\documentclass[UTF8,12pt,a4paper,fontset=fandol]{ctexbook}

% 支持从项目根目录或当前笔记目录编译。
\makeatletter
\def\input@path{{../}{../../}{../../../}{../../../../}}
\makeatother
\usepackage{researchnotes}

% 编译章节列表：{0} 编译全文；{3} 只编译第三章；{2,3} 或 {2,5} 编译多章。
% 正文章号范围为 1--17，多个章号用半角逗号分隔，仍按原书顺序输出。
% 列表中含 0 时编译全文。日常只需修改下面这一行。
\newcommand{\CompileChapter}{5}

\title{数学分析教材（伍胜健）笔记}
\author{}
\date{\today}

\begin{document}
\ifcompileallchapters

\maketitle

\frontmatter
\tableofcontents

\chapter{绪论}

\fi

\mainmatter

% 以下章、节、小节及章末习题栏目按教材目录顺序整理。
% 第一至六章目录来源：数学分析/textbooks/Real-Analysis (1)/Real-Analysis (1).pdf，第 I--III 页。
% 后续可在对应标题下补充学习思考与习题解答。

\selectchapter{1}
\ifcompilechapter
\chapter{函数}
\label{chap:functions}

\section{实数}
\label{sec:real-numbers}

\subsection{数集}
\label{subsec:number-sets}

\subsection{实数系的连续性}
\label{subsec:continuity-of-real-numbers}

\subsection{有界集与确界}
\label{subsec:bounded-sets-and-bounds}

以下均设 $E$ 为非空实数集，即 $\varnothing\neq E\subseteq\mathbb{R}$，所讨论的界均为实数。

\begin{definition}
\label{def:upper-and-lower-bounds}
若存在 $M\in\mathbb{R}$，使得
\[
    \forall x\in E,\quad x\leq M,
\]
则称 $E$ \emph{有上界}（bounded above），并称 $M$ 为 $E$ 的一个\emph{上界}（upper bound）。若存在 $m\in\mathbb{R}$，使得
\[
    \forall x\in E,\quad m\leq x,
\]
则称 $E$ \emph{有下界}（bounded below），并称 $m$ 为 $E$ 的一个\emph{下界}（lower bound）。
\end{definition}

\emph[caution]{上界和下界都不要求属于 $E$，也通常不唯一}：若 $M$ 是上界，则任意 $M'\geq M$ 也是上界；若 $m$ 是下界，则任意 $m'\leq m$ 也是下界。

\begin{definition}
\label{def:bounded-and-unbounded-sets}
若 $E$ 既有上界又有下界，则称 $E$ 为\emph{有界集}（bounded set），即存在实数 $m,M$，使得
\[
    \forall x\in E,\quad m\leq x\leq M.
\]
若 $E$ 不是有界集，则称 $E$ 为\emph{无界集}（unbounded set）。具体地，$E$ 无上界（unbounded above）是指
\[
    \forall M\in\mathbb{R},\quad \exists x\in E,\quad x>M;
\]
$E$ 无下界（unbounded below）是指
\[
    \forall m\in\mathbb{R},\quad \exists x\in E,\quad x<m.
\]
\emph[caution]{因此，$E$ 无界意味着它至少在一个方向上无界，并不要求它同时无上界和无下界。}
\end{definition}

\begin{definition}
\label{def:maximum-and-minimum}
若 $a\in E$ 且对所有 $x\in E$ 都有 $x\leq a$，则称 $a$ 为 $E$ 的\emph{最大值}（maximum），记为 $a=\max E$。若 $b\in E$ 且对所有 $x\in E$ 都有 $b\leq x$，则称 $b$ 为 $E$ 的\emph{最小值}（minimum），记为 $b=\min E$。
\end{definition}

\begin{definition}
\label{def:supremum}
若 $\beta\in\mathbb{R}$ 是 $E$ 的上界，且对 $E$ 的任意上界 $M$ 都有 $\beta\leq M$，则称 $\beta$ 为 $E$ 的\emph{上确界}（supremum），记为 $\beta=\sup E$。也就是说，上确界是\emph{最小的上界}（least upper bound）。等价地，$\beta=\sup E$ 当且仅当
\[
    \begin{cases}
        \forall x\in E,\quad x\leq\beta,\\
        \forall\varepsilon>0,\quad \exists x_{\varepsilon}\in E,\quad
        \beta-\varepsilon<x_{\varepsilon}\leq\beta.
    \end{cases}
\]
第一条说明 $\beta$ 是上界，\emph[tip]{第二条说明无论把 $\beta$ 减小多少，所得的数都不再是上界}。事实上，任意小于 $\beta$ 的实数都可写成 $\beta-\varepsilon$，其中 $\varepsilon>0$，故第二条恰好排除了所有更小的上界。
\end{definition}

\begin{definition}
\label{def:infimum}
若 $\alpha\in\mathbb{R}$ 是 $E$ 的下界，且对 $E$ 的任意下界 $m$ 都有 $m\leq\alpha$，则称 $\alpha$ 为 $E$ 的\emph{下确界}（infimum），记为 $\alpha=\inf E$。也就是说，下确界是\emph{最大的下界}（greatest lower bound）。等价地，$\alpha=\inf E$ 当且仅当
\[
    \begin{cases}
        \forall x\in E,\quad \alpha\leq x,\\
        \forall\varepsilon>0,\quad \exists y_{\varepsilon}\in E,\quad
        \alpha\leq y_{\varepsilon}<\alpha+\varepsilon.
    \end{cases}
\]
第一条说明 $\alpha$ 是下界，第二条说明无论把 $\alpha$ 增大多少，所得的数都不再是下界；因为任意大于 $\alpha$ 的实数都可写成 $\alpha+\varepsilon$，其中 $\varepsilon>0$。
\end{definition}

确界若存在，则必唯一。例如，若 $\beta_1,\beta_2$ 都是 $E$ 的上确界，由最小上界的定义，分别有 $\beta_1\leq\beta_2$ 和 $\beta_2\leq\beta_1$，故 $\beta_1=\beta_2$；下确界的唯一性同理。\emph[caution]{确界不一定属于原集合，而最大值和最小值必须属于原集合。}\emph{当上确界属于 $E$ 时，它就是最大值；当下确界属于 $E$ 时，它就是最小值}。反过来，最大值必为上确界，最小值必为下确界：例如，若 $a=\max E$，则 $a$ 是上界，且由 $a\in E$ 可知每个上界都不小于 $a$，所以 $a=\sup E$。

例如，开区间 $(0,1)$ 和闭区间 $[0,1]$ 的上确界都是 $1$，下确界都是 $0$。对任意 $\varepsilon>0$，取 $\delta=\min\{\varepsilon/2,1/2\}$，则 $\delta$ 与 $1-\delta$ 都属于 $(0,1)$，且 $0<\delta<\varepsilon$、$1-\varepsilon<1-\delta<1$，因而两个区间都满足上述确界刻画。对于 $(0,1)$，两个确界均不属于集合，所以它既无最大值也无最小值；对于 $[0,1]$，则有 $\max[0,1]=1$、$\min[0,1]=0$。

\subsection{几个常用不等式}
\label{subsec:common-inequalities}

\begin{definition}
\label{def:real-absolute-value}
实数 $x$ 的\emph{绝对值}（absolute value）定义为
\[
    |x|=
    \begin{cases}
        x, & x\geq 0,\\
        -x, & x<0,
    \end{cases}
\]
它表示数轴上 $x$ 到原点的距离，也可写成 $|x|=\sqrt{x^2}$。
\end{definition}

由定义可知，$-|x|\leq x\leq |x|$；对任意 $c\geq 0$，$|x|\leq c$ 等价于 $-c\leq x\leq c$，也等价于 $x\leq c$ 与 $-x\leq c$ 同时成立。这些性质使我们能够用绝对值统一估计正、负两种情形下的大小。

\begin{proposition}[三角不等式]
\label{prop:real-triangle-inequality}
对任意实数 $a,b$，都有
\begin{equation}
    |a+b|\leq |a|+|b|.
    \label{eq:real-triangle-inequality}
\end{equation}
式~\eqref{eq:real-triangle-inequality} 称为\emph{三角不等式}（triangle inequality），即两数之和的绝对值不超过它们绝对值之和；等号成立当且仅当 $ab\geq 0$。
\end{proposition}

\begin{proof}
由绝对值的性质，$-|a|\leq a\leq |a|$ 且 $-|b|\leq b\leq |b|$，两式相加，得到
\[
    -(|a|+|b|)\leq a+b\leq |a|+|b|.
\]
由于 $|a|+|b|\geq 0$，这等价于 $|a+b|\leq |a|+|b|$。

等号成立当且仅当 $a+b=|a|+|b|$ 或 $a+b=-(|a|+|b|)$。前一种情形等价于 $(|a|-a)+(|b|-b)=0$，由于两项均非负，又等价于 $a\geq 0$ 且 $b\geq 0$；后一种情形同理等价于 $a\leq 0$ 且 $b\leq 0$。因此，等号成立当且仅当 $ab\geq 0$，也就是 $a,b$ 同号或至少有一个为零。
\end{proof}

\emph[caution]{三角不等式给出 $|a-b|\leq |a|+|b|$ 这一上界；利用它还可得到 $|a-b|$ 的下界，即下面的反向三角不等式。}

\begin{corollary}[反向三角不等式]
\label{cor:real-reverse-triangle-inequality}
对任意实数 $a,b$，都有\emph{反向三角不等式}（reverse triangle inequality）
\begin{equation}
    \bigl||a|-|b|\bigr|\leq |a-b|.
    \label{eq:real-reverse-triangle-inequality}
\end{equation}
即两数绝对值之差的绝对值不超过它们之差的绝对值。
\end{corollary}

\begin{proof}
将 $a$ 写成 $(a-b)+b$，将 $b$ 写成 $(b-a)+a$，分别应用命题~\ref{prop:real-triangle-inequality}，得到
\[
    \begin{aligned}
        |a|&\leq |a-b|+|b|,\\
        |b|&\leq |b-a|+|a|=|a-b|+|a|.
    \end{aligned}
\]
移项后分别有 $|a|-|b|\leq |a-b|$ 和 $|b|-|a|\leq |a-b|$，合并即得
\[
    -|a-b|\leq |a|-|b|\leq |a-b|.
\]
由于 $|a-b|\geq 0$，根据绝对值的性质，这等价于 $\bigl||a|-|b|\bigr|\leq |a-b|$。
\end{proof}

\begin{definition}
\label{def:arithmetic-and-geometric-means}
设 $n$ 为正整数，$a_1,\ldots,a_n$ 为非负实数。它们的\emph{算术平均}（arithmetic mean）与\emph{几何平均}（geometric mean）分别定义为
\[
    A_n=\frac{a_1+\cdots+a_n}{n},
    \qquad
    G_n=\sqrt[n]{a_1a_2\cdots a_n}.
\]
其中，几何平均取非负的 $n$ 次方根。下面的不等式说明，这两个平均值之间存在确定的大小关系。
\end{definition}

\begin{proposition}[算术平均--几何平均不等式]
\label{prop:am-gm-inequality}
对任意正整数 $n$ 及非负实数 $a_1,\ldots,a_n$，都有
\begin{equation}
    \sqrt[n]{a_1a_2\cdots a_n}
    \leq \frac{a_1+\cdots+a_n}{n}.
    \label{eq:am-gm-inequality}
\end{equation}
式~\eqref{eq:am-gm-inequality} 称为\emph{算术平均--几何平均不等式}（arithmetic mean--geometric mean inequality，简称 AM--GM 不等式），即 $G_n\leq A_n$；\emph[caution]{等号成立当且仅当 $a_1=\cdots=a_n$。}
\end{proposition}

\begin{proof}
对 $n$ 作数学归纳。$n=1$ 时两边相等；$n=2$ 时，由
\[
    \frac{a_1+a_2}{2}-\sqrt{a_1a_2}
    =\frac{(\sqrt{a_1}-\sqrt{a_2})^2}{2}\geq 0
\]
可得结论，且等号成立当且仅当 $a_1=a_2$。

设 $n\geq 3$，并假设结论对 $n-1$ 个非负实数成立。若 $a_1,\ldots,a_n$ 中有零，则几何平均为零，不等式成立；此时等号成立当且仅当算术平均也为零，即所有 $a_i$ 均为零。以下设各 $a_i>0$，记
\[
    A=\frac{a_1+\cdots+a_n}{n}>0.
\]
若所有 $a_i$ 相等，则等号成立。若它们不全相等，由 $\sum_{i=1}^{n}a_i=nA$ 可知，必有一项小于 $A$，另一项大于 $A$；否则各项都在 $A$ 的同一侧且至少一项严格不等，其和便不等于 $nA$。重新编号后可设 $a_1<A<a_2$。

\emph[tip]{将 $a_1,a_2$ 替换为 $A$ 和 $b=a_1+a_2-A$，保持两项之和不变，同时增大两项之积。}事实上，$b=a_1+(a_2-A)>0$，且
\[
    Ab-a_1a_2=(A-a_1)(a_2-A)>0.
\]
其余 $n-1$ 个数 $b,a_3,\ldots,a_n$ 的和为 $(n-1)A$，由归纳假设，
\[
    ba_3\cdots a_n\leq A^{n-1}.
\]
由于各项均为正，因而
\[
    a_1a_2\cdots a_n
    <Ab\,a_3\cdots a_n
    \leq A\cdot A^{n-1}=A^n.
\]
两边取非负的 $n$ 次方根，便得 $G_n<A=A_n$。因此，各数不全相等时为严格不等式，全相等时为等式，归纳完成。
\end{proof}

下面讨论幂与一次式之间的比较。先回顾两个代数恒等式：设 $a,b\in\mathbb{R}$，$m$ 为正整数，则有\emph{幂差公式}
\begin{equation}
    a^m-b^m=(a-b)\sum_{k=0}^{m-1}a^{m-1-k}b^k;
    \label{eq:integer-power-difference}
\end{equation}
当 $m$ 为奇数时，还有\emph{幂和公式}
\begin{equation}
    a^m+b^m=(a+b)\sum_{k=0}^{m-1}(-1)^k a^{m-1-k}b^k.
    \label{eq:odd-integer-power-sum}
\end{equation}
在幂差公式中，将右端乘开后，中间项 $a^{m-k}b^k$（$1\leq k\leq m-1$）分别出现一次正号和一次负号，抵消后只剩 $a^m-b^m$；幂和公式则由幂差公式中以 $-b$ 代替 $b$ 得到，其中用到 $m$ 为奇数时 $(-b)^m=-b^m$。下面利用幂差公式证明伯努利不等式。

\begin{proposition}[伯努利不等式]
\label{prop:bernoulli-inequality}
对任意实数 $x\geq-1$ 和任意正整数 $n$，都有
\begin{equation}
    (1+x)^n\geq1+nx.
    \label{eq:bernoulli-inequality}
\end{equation}
式~\eqref{eq:bernoulli-inequality} 称为\emph{伯努利不等式}（Bernoulli's inequality）；等号成立当且仅当 $n=1$ 或 $x=0$。
\end{proposition}

\begin{proof}
当 $n=1$ 或 $x=0$ 时，两边相等。以下设 $n\geq2$ 且 $x\neq0$。若 $x=-1$，则 $(1+x)^n=0>1-n=1+nx$。以下只需讨论 $x>-1$ 的情形，由幂差公式~\eqref{eq:integer-power-difference}，
\[
    (1+x)^n-1
    =x\bigl[1+(1+x)+(1+x)^2+\cdots+(1+x)^{n-1}\bigr].
\]
\emph[tip]{先比较方括号内的 $n$ 项之和与 $n$ 的大小，再根据 $x$ 的正负确定乘积的不等号方向。}当 $x>0$ 时，$1+x>1$，方括号内除首项等于 $1$ 外，其余 $n-1$ 项均大于 $1$，故其和大于 $n$；乘以 $x>0$，得到 $(1+x)^n-1>nx$。当 $-1<x<0$ 时，$0<1+x<1$，方括号内除首项外，其余各项均小于 $1$，故其和小于 $n$；由于 $x<0$，乘以 $x$ 后不等号反向，仍得到 $(1+x)^n-1>nx$。因此不等式在全部条件下成立，且 $n\geq2$、$x\neq0$ 时均为严格不等式，这也给出了所述等号条件。
\end{proof}

为了便于应用，还可通过代换 $x=B/A$，将伯努利不等式写成双参数形式。

\begin{corollary}[伯努利不等式的双参数形式]
\label{cor:two-parameter-bernoulli}
设 $A,B\in\mathbb{R}$，满足 $A>0$ 且 $A+B\geq0$。对任意正整数 $n$，都有
\begin{equation}
    (A+B)^n\geq A^n+nA^{n-1}B.
    \label{eq:two-parameter-bernoulli}
\end{equation}
当 $n=1$ 时恒取等号；当 $n>1$ 时，等号成立当且仅当 $B=0$。
\end{corollary}

\begin{proof}
由 $A>0$ 且 $A+B\geq0$，有 $1+B/A\geq0$，即 $B/A\geq-1$。在式~\eqref{eq:bernoulli-inequality} 中取 $x=B/A$，再将两边乘以正数 $A^n$，得到
\[
    (A+B)^n=A^n\left(1+\frac BA\right)^n
    \geq A^n\left(1+n\frac BA\right)
    =A^n+nA^{n-1}B.
\]
上述变形保持等号条件，而 $B/A=0$ 等价于 $B=0$，故等号条件随之成立。
\end{proof}

伯努利不等式把幂与一次式联系起来。下面两个例题分别展示换元与取商的用法：前者用于证明根式不等式，后者用于比较底数与指数同时变化的两个幂。

\begin{example}
\label{ex:bernoulli-root-bound}
设 $x\geq-1$，$n$ 为正整数。证明
\begin{equation}
    \sqrt[n]{1+x}\leq1+\frac xn,
    \label{eq:bernoulli-root-bound}
\end{equation}
并确定等号条件。这里根式取非负的 $n$ 次方根。
\end{example}

\begin{solve}
\emph[tip]{把根式设为 $1+t$，就能使原来的正整数次幂不等式适用。}令
\[
    t=\sqrt[n]{1+x}-1\geq-1.
\]
由伯努利不等式，
\[
    1+x=(1+t)^n\geq1+nt.
\]
移项并除以正整数 $n$，得到 $t\leq x/n$，即式~\eqref{eq:bernoulli-root-bound}。等号成立当且仅当 $n=1$ 或 $t=0$，而 $t=0$ 等价于 $x=0$。
\end{solve}

\begin{example}
\label{ex:bernoulli-adjacent-powers}
对任意正整数 $n$，证明
\[
    \left(1+\frac1{n+1}\right)^{n+1}
    >\left(1+\frac1n\right)^n.
\]
\end{example}

\begin{solve}
两边的底数与指数同时变化，\emph[tip]{取商后将相同指数的部分合并，可以得到适合伯努利不等式的形式。}具体地，
\[
    \frac{\left(1+\frac1{n+1}\right)^{n+1}}
         {\left(1+\frac1n\right)^n}
    =\frac{n+2}{n+1}
      \left(\frac{n(n+2)}{(n+1)^2}\right)^n
    =\frac{n+2}{n+1}
      \left(1-\frac1{(n+1)^2}\right)^n.
\]
由于 $-1/(n+1)^2\geq-1$，由伯努利不等式可得
\[
    \begin{aligned}
        \frac{\left(1+\frac1{n+1}\right)^{n+1}}
             {\left(1+\frac1n\right)^n}
        &\geq\frac{n+2}{n+1}
          \left(1-\frac{n}{(n+1)^2}\right)\\
        &=\frac{(n+2)(n^2+n+1)}{(n+1)^3}
          =1+\frac1{(n+1)^3}>1.
    \end{aligned}
\]
分母为正，故原不等式成立。这一比较说明，$(1+1/n)^n$ 随正整数 $n$ 严格增大。
\end{solve}

\subsection{常用记号}
\label{subsec:common-notation}

\section{函数的概念}
\label{sec:concept-of-functions}

\subsection{函数的定义}
\label{subsec:definition-of-functions}

\subsection{由已知函数构造新函数的方法}
\label{subsec:constructing-functions}

\section{函数的性质}
\label{sec:properties-of-functions}

以下讨论实值函数 $f:D\to\mathbb{R}$，其中 $\varnothing\neq D\subseteq\mathbb{R}$ 为定义域。有界性描述函数值的范围，单调性描述函数值随自变量变化的方向，周期性与奇偶性则分别描述平移和反射所保持的关系。讨论函数性质时，必须同时明确函数的对应法则与讨论范围。

\subsection{函数的有界性}
\label{subsec:boundedness-of-functions}

设 $E$ 是 $D$ 的非空子集，记 $f(E)=\{f(x):x\in E\}$。函数 $f$ 在 $E$ 上的有界性，就是函数值集合 $f(E)$ 的有界性，因而可沿用第~\ref{subsec:bounded-sets-and-bounds} 小节中的概念。

\begin{definition}
\label{def:bounded-functions}
若存在实数 $U$，使得对所有 $x\in E$ 都有 $f(x)\leq U$，则称 $f$ 在 $E$ 上\emph{有上界}（bounded above）；若存在实数 $L$，使得对所有 $x\in E$ 都有 $L\leq f(x)$，则称 $f$ 在 $E$ 上\emph{有下界}（bounded below）。若 $f$ 在 $E$ 上既有上界又有下界，则称它在 $E$ 上\emph{有界}（bounded），等价地，
\[
    \exists M\geq 0,\quad \forall x\in E,\quad |f(x)|\leq M.
\]
若不存在这样的 $M$，则称 $f$ 在 $E$ 上\emph{无界}（unbounded），即
\[
    \forall M>0,\quad \exists x\in E,\quad |f(x)|>M.
\]
当 $E=D$ 时，简称 $f$ 为有界函数或无界函数。
\end{definition}

上述等价性来自绝对值的性质：若 $L\leq f(x)\leq U$，取 $M=\max\{|L|,|U|\}$，便有 $|f(x)|\leq M$；反之，$|f(x)|\leq M$ 给出 $-M\leq f(x)\leq M$。\emph[caution]{这里的界 $M$ 必须对所有 $x\in E$ 同时有效，不能随 $x$ 改变。}无界并不意味着某一点的函数值是无穷大，而是函数值的绝对值能够超过任意预先给定的正数。

\begin{example}
\label{ex:function-boundedness-domain}
判断函数 $f(x)=1/x$ 在 $(0,1)$ 与 $[1,+\infty)$ 上的有界性。
\end{example}

\begin{solve}
在 $(0,1)$ 上，$f(x)>1$，所以 $f$ 有下界。对任意 $M>0$，取 $x=1/(M+1)\in(0,1)$，便有 $f(x)=M+1>M$，故 $f$ 无上界，从而无界。在 $[1,+\infty)$ 上，由 $0<f(x)\leq 1$ 可知 $f$ 有界。这也说明，定义域有界并不保证函数有界，定义域无界也不妨碍函数有界。
\end{solve}

有界性只要求存在统一的界，不要求函数达到这些界。例如，$f(x)=x$ 在 $(0,1)$ 上的函数值集合仍为 $(0,1)$，其上确界为 $1$、下确界为 $0$，但函数没有最大值和最小值。

\subsection{函数的单调性}
\label{subsec:monotonicity-of-functions}

\begin{definition}
\label{def:monotone-functions}
设 $\varnothing\neq E\subseteq D$。若对任意 $x_1,x_2\in E$，由 $x_1<x_2$ 总能推出 $f(x_1)\leq f(x_2)$，则称 $f$ 在 $E$ 上\emph{单调增加}或\emph{不减}（nondecreasing）；若总能推出 $f(x_1)\geq f(x_2)$，则称 $f$ 在 $E$ 上\emph{单调减少}或\emph{不增}（nonincreasing）。这两类函数统称为\emph{单调函数}（monotone function）。

若对任意 $x_1<x_2$ 都有 $f(x_1)<f(x_2)$，则称 $f$ 在 $E$ 上\emph{严格单调增加}（strictly increasing）；若都有 $f(x_1)>f(x_2)$，则称它在 $E$ 上\emph{严格单调减少}（strictly decreasing），二者统称为严格单调函数。这里各条件中的 $x_1,x_2$ 均取自 $E$。
\end{definition}

本笔记用“单调增加、单调减少”表示允许相等的情形，用“严格”排除相等。因此，常值函数既不减又不增；严格单调函数在不同自变量处的函数值必不相同。\emph[caution]{单调性要求比较讨论范围内的任意两点，不能仅比较若干点，也不能把各个区间上的单调性直接拼接为整个定义域上的单调性。}

\begin{proposition}
\label{prop:monotone-function-endpoint-bounds}
若 $a<b$ 且 $f$ 在闭区间 $[a,b]$ 上单调，则 $f$ 在该区间上有界。具体地，若 $f$ 不减，则
\[
    f(a)\leq f(x)\leq f(b),\qquad x\in[a,b];
\]
若 $f$ 不增，则
\[
    f(b)\leq f(x)\leq f(a),\qquad x\in[a,b].
\]
因此，它在该区间上的最大值和最小值都可在端点取得。
\end{proposition}

\begin{proof}
对每个 $x\in[a,b]$，都有 $a\leq x\leq b$；分别对 $a,x$ 与 $x,b$ 应用单调性的定义，即得相应不等式，端点处允许等号。因此，两个端点值分别给出函数的最大值和最小值。
\end{proof}

\begin{example}
\label{ex:monotonicity-on-subsets}
讨论 $f(x)=x^2$ 与 $g(x)=1/x$ 的单调性。
\end{example}

\begin{solve}
对 $x_1<x_2$，有
\[
    x_2^2-x_1^2=(x_2-x_1)(x_2+x_1).
\]
若 $0\leq x_1<x_2$，右端为正；若 $x_1<x_2\leq 0$，右端为负。因此，$x^2$ 在 $[0,+\infty)$ 上严格增加，在 $(-\infty,0]$ 上严格减少，在整个 $\mathbb{R}$ 上不单调。

对于 $g$，若 $x_1<x_2$ 同在 $(-\infty,0)$ 或同在 $(0,+\infty)$，则
\[
    g(x_2)-g(x_1)=\frac{x_1-x_2}{x_1x_2}<0.
\]
所以 $g$ 在这两个区间上分别严格减少。但 $-1<1$ 而 $g(-1)=-1<1=g(1)$，不符合整个定义域上单调减少的要求；结合它在 $(0,+\infty)$ 上严格减少可知，$g$ 在 $\mathbb{R}\setminus\{0\}$ 上并不单调。
\end{solve}

\subsection{函数的周期性}
\label{subsec:periodicity-of-functions}

\begin{definition}
\label{def:periodic-functions}
若存在常数 $T>0$，使得对每个 $x\in D$ 都有
\[
    x+T\in D,\qquad x-T\in D,\qquad f(x+T)=f(x),
\]
则称 $f$ 为\emph{周期函数}（periodic function），称 $T$ 为它的一个\emph{周期}（period）。本节只讨论正周期；若所有正周期中存在最小者，则称它为 $f$ 的\emph{最小正周期}或\emph{基本周期}（fundamental period）。
\end{definition}

周期性表示函数的整个图像沿水平方向平移 $T$ 后保持不变。定义中要求定义域对正、负两个方向的平移都保持不变，不能仅在 $x$ 与 $x+T$ 恰好都属于定义域时比较函数值。由于把 $x-T$ 代入周期关系可得 $f(x-T)=f(x)$，反复使用该关系便有
\[
    f(x+kT)=f(x),\qquad x\in D,\quad k\in\mathbb{Z}.
\]
特别地，$T$ 的任意正整数倍仍是周期，但周期不一定是最小正周期。例如，定义在 $\mathbb{R}$ 上的常值函数以每个正数为周期，因而没有最小正周期。

\begin{example}
\label{ex:periods-of-sine-functions}
求 $f(x)=\sin x$ 与 $g(x)=\sin^2 x$ 在 $\mathbb{R}$ 上的最小正周期。
\end{example}

\begin{solve}
由 $\sin(x+2\pi)=\sin x$ 可知 $2\pi$ 是 $f$ 的周期。若 $T>0$ 是任意周期，令 $x=0$，得 $\sin T=0$，所以 $T=k\pi$，其中 $k$ 为正整数；再令 $x=\pi/2$，得 $\cos T=1$，所以 $k$ 必为偶数。因此，$f$ 的最小正周期为 $2\pi$。

由 $\sin(x+\pi)=-\sin x$ 可知 $\sin^2(x+\pi)=\sin^2 x$，所以 $\pi$ 是 $g$ 的周期。若 $T>0$ 是 $g$ 的周期，令 $x=0$，便有 $\sin^2 T=0$，从而 $T=k\pi$，其中 $k$ 为正整数。因此，$g$ 的最小正周期为 $\pi$。
\end{solve}

\begin{proposition}
\label{prop:periodic-function-boundedness}
设 $f:\mathbb{R}\to\mathbb{R}$ 以 $T>0$ 为周期，则 $f$ 在 $\mathbb{R}$ 上有界，当且仅当它在 $[0,T)$ 上有界。
\end{proposition}

\begin{proof}
若 $f$ 在 $\mathbb{R}$ 上有界，同一个界也适用于 $[0,T)$。反之，设 $|f(t)|\leq M$ 对所有 $t\in[0,T)$ 成立。任意实数 $x$ 都可写成 $x=kT+t$，其中 $k\in\mathbb{Z}$、$0\leq t<T$；由周期性，$|f(x)|=|f(t)|\leq M$，故 $f$ 在 $\mathbb{R}$ 上有界。
\end{proof}

\subsection{函数的奇偶性}
\label{subsec:parity-of-functions}

\begin{definition}
\label{def:even-and-odd-functions}
设定义域 $D$ \emph{关于原点对称}，即 $x\in D$ 时总有 $-x\in D$。若
\[
    \forall x\in D,\quad f(-x)=f(x),
\]
则称 $f$ 为\emph{偶函数}（even function）；若
\[
    \forall x\in D,\quad f(-x)=-f(x),
\]
则称 $f$ 为\emph{奇函数}（odd function）。若上述两类条件都不满足，则称 $f$ 为非奇非偶函数；定义域不关于原点对称的函数也不属于奇函数或偶函数。
\end{definition}

偶函数的图像关于 $y$ 轴对称，奇函数的图像关于原点对称，因为图像上的点 $(x,f(x))$ 分别对应于 $(-x,f(x))$ 与 $(-x,-f(x))$。\emph[caution]{判断奇偶性应先检查定义域的对称性，再比较 $f(-x)$ 与 $f(x)$。}若奇函数在 $0$ 处有定义，则 $f(0)=-f(0)$，所以 $f(0)=0$；但奇函数的定义域不一定包含 $0$，例如 $f(x)=1/x$。在同一非空对称定义域上，一个函数既奇又偶，当且仅当它恒为零，因为这时 $f(x)=f(-x)=-f(x)$。

\begin{example}
\label{ex:parity-and-domain}
判断定义在 $\mathbb{R}$ 上的 $f(x)=x^2$、$g(x)=x^3$、$h(x)=x^2+x$ 的奇偶性，并讨论将 $f$ 限制在 $[0,+\infty)$ 上后的情形。
\end{example}

\begin{solve}
三个函数的定义域 $\mathbb{R}$ 都关于原点对称。由 $f(-x)=x^2=f(x)$ 与 $g(-x)=-x^3=-g(x)$，可知 $f$ 是偶函数，$g$ 是奇函数。对于 $h$，有 $h(-1)=0$、$h(1)=2$，既不满足 $h(-1)=h(1)$，也不满足 $h(-1)=-h(1)$，所以它非奇非偶。将 $f$ 限制在 $[0,+\infty)$ 上后，定义域不再关于原点对称，因此该限制函数非奇非偶，尽管它仍由公式 $x^2$ 给出。
\end{solve}

\begin{proposition}[奇偶分解]
\label{prop:even-odd-decomposition}
设 $D$ 非空且关于原点对称，则任意函数 $f:D\to\mathbb{R}$ 都能唯一写成一个偶函数与一个奇函数之和：
\begin{equation}
    f(x)=f_{\mathrm e}(x)+f_{\mathrm o}(x),\qquad
    f_{\mathrm e}(x)=\frac{f(x)+f(-x)}{2},\qquad
    f_{\mathrm o}(x)=\frac{f(x)-f(-x)}{2}.
    \label{eq:even-odd-decomposition}
\end{equation}
其中两个分量的定义域均为 $D$。
\end{proposition}

\begin{proof}
由 $D$ 的对称性，式~\eqref{eq:even-odd-decomposition} 中各函数均有定义。直接计算得
\[
    f_{\mathrm e}(-x)=\frac{f(-x)+f(x)}{2}=f_{\mathrm e}(x),
    \qquad
    f_{\mathrm o}(-x)=\frac{f(-x)-f(x)}{2}=-f_{\mathrm o}(x),
\]
故 $f_{\mathrm e}$ 为偶函数，$f_{\mathrm o}$ 为奇函数，且两式相加得到 $f=f_{\mathrm e}+f_{\mathrm o}$。若另有 $f=u+v$，其中 $u$ 为偶函数、$v$ 为奇函数，则
\[
    f(x)=u(x)+v(x),\qquad f(-x)=u(x)-v(x).
\]
相加与相减分别给出 $u=f_{\mathrm e}$、$v=f_{\mathrm o}$，所以分解唯一。
\end{proof}

\section{初等函数}
\label{sec:elementary-functions}

\section*{习题一}
\addcontentsline{toc}{section}{习题一}
\fi

\selectchapter{2}
\ifcompilechapter
\chapter{序列的极限}
\label{chap:sequence-limits}

\section{序列极限的定义}
\label{sec:definition-of-sequence-limits}

本节研究一列实数随下标增大时的变化趋势。\emph{序列趋近某个实数，意味着无论给定多小的正数作为允许的误差，都能找到一项，使这一项及其后面所有项与该实数的距离都小于给定的误差。}本节还将介绍无穷小量与无穷大量，分别讨论序列趋于零和绝对值无限增大的情形。以下下标均取正整数，并约定 $\mathbb{N}=\{1,2,\ldots\}$。

\subsection{序列}
\label{subsec:sequences}

\begin{definition}
\label{def:real-sequence}
从正整数集到实数集的函数 $x:\mathbb{N}\to\mathbb{R}$ 称为\emph{实数序列}（real sequence），简称序列或数列。记 $x_n=x(n)$，通常把它写成
\[
    x_1,x_2,\ldots,x_n,\ldots,
    \qquad\text{或简记为 }\{x_n\},
\]
其中 $x_n$ 称为第 $n$ 项，也称\emph{通项}（general term）。
\end{definition}

给出序列不要求一定有简单的通项公式，也可以用递推关系或其他确定规则指定每一项。例如，$1,1/2,1/3,\ldots$ 的通项为 $1/n$；将 $\pi$ 依次截断到小数点后 $n-1$ 位得到 $3,3.1,3.14,\ldots$，也确定了一个序列。区间 $[0,1)$ 中的有理数还能不重复地排成序列：先列出 $0$，再将既约分数 $p/q$ 按分母 $q=2,3,\ldots$ 递增排列，同一分母内按分子递增排列，得到
\[
    0,\ \frac12,\ \frac13,\ \frac23,\ \frac14,\ \frac34,\ \frac15,\ \frac25,\ldots.
\]
每个有理数都有唯一的既约表示，且排在任一固定分母之前的分数只有有限个，因而每个数都会在某个有限位置被列到。

\subsection{序列极限的定义}
\label{subsec:definition-of-sequence-limits}

对于 $x_n=1/n$，各项始终在零的右侧；对于 $x_n=(-1)^n/n$，各项在零的两侧交替出现；若令 $x_{2k-1}=1/k$、$x_{2k}=0$，序列还会反复取到零。三者的变化方式不同，但其全部充分靠后的项都能任意接近零。可见，趋近极限并不要求始终从同一侧逼近，也不要求误差逐项减小或各项始终异于极限。

\begin{definition}
\label{def:sequence-limit}
设 $\{x_n\}$ 为实数列。若存在实数 $a$，使
\begin{equation}
    \forall\varepsilon>0,\quad\exists N\in\mathbb{N},\quad
    \forall n>N,\quad |x_n-a|<\varepsilon,
    \label{eq:epsilon-N-sequence-limit}
\end{equation}
则称序列\emph{收敛}（convergent），称 $a$ 为它的\emph{极限}（limit），记为
\[
    \lim_{n\to\infty}x_n=a,
    \qquad\text{或}\qquad x_n\to a\quad(n\to\infty).
\]
若不存在这样的实数 $a$，则称序列\emph{发散}（divergent）。
\end{definition}

这里 $\varepsilon$ 是事先任意给定的误差，$N$ 可以依赖 $\varepsilon$，但选定后必须对全部 $n>N$ 同时有效。把 $n>N$ 改成 $n\geq N$ 得到等价定义，因为可以把阈值增减一个整数。若某个 $N$ 已经适用，更大的阈值也适用；因此证明不需要找最小的 $N$，适当放大估计往往更方便。

从数轴上看，式~\eqref{eq:epsilon-N-sequence-limit} 表示：对每个 $\varepsilon>0$，开区间 $(a-\varepsilon,a+\varepsilon)$ 包含从某项起的全部项，等价地，落在这个区间外的项只有有限项。\emph[caution]{这里按下标计数，要求的是“最终所有项都接近”，而不只是“有无穷多项接近”。}例如，$(-1)^n$ 有无穷多项等于 $1$，却不趋于 $1$。把序列画成离散点 $(n,x_n)$，收敛则表示足够靠右的所有点都位于水平带 $a-\varepsilon<y<a+\varepsilon$ 内。

\begin{example}
\label{ex:geometric-sequence-limit}
证明 $|q|<1$ 时 $q^n\to0$。
\end{example}

\begin{solve}
若 $q=0$，则 $q^n=0$，结论成立。以下设 $0<|q|<1$。给定 $\varepsilon>0$，由对数函数的严格单调递增性，
\[
    |q^n-0|=|q|^n<\varepsilon
    \quad\Longleftrightarrow\quad
    n\ln|q|<\ln\varepsilon
    \quad\Longleftrightarrow\quad
    n>\frac{\ln\varepsilon}{\ln|q|}.
\]
其中最后一步用到 $\ln|q|<0$，除以它时不等号反向。因此，取正整数 $N>\ln\varepsilon/\ln|q|$，则当 $n>N$ 时，就有 $|q^n-0|<\varepsilon$，故 $q^n\to0$。
\end{solve}

\begin{example}
\label{ex:fixed-root-sequence-limit}
证明对任意固定的 $a>0$，都有 $\sqrt[n]{a}\to1$。
\end{example}

\begin{solve}
若 $a>1$，令 $t_n=\sqrt[n]{a}-1>0$，由\crosschapterref{1}{伯努利不等式~\eqref{eq:bernoulli-inequality}}{伯努利不等式}，得到 $a=(1+t_n)^n\geq1+nt_n$，故
\[
    0<t_n\leq\frac{a-1}{n}.
\]
取 $N>(a-1)/\varepsilon$ 就能使 $|\sqrt[n]{a}-1|<\varepsilon$。$a=1$ 时序列恒为 $1$；$0<a<1$ 时令 $b=1/a>1$，则
\[
    0<1-\sqrt[n]{a}
    =\frac{\sqrt[n]{b}-1}{\sqrt[n]{b}}
    \leq\sqrt[n]{b}-1\leq\frac{b-1}{n},
\]
同样按定义得到极限为 $1$。
\end{solve}

为了否定以某个给定的实数 $a$ 为极限，必须逐个否定量词：
\begin{equation}
    x_n\not\to a
    \quad\Longleftrightarrow\quad
    \exists\varepsilon_0>0,\quad\forall N\in\mathbb{N},\quad
    \exists n>N,\quad |x_n-a|\geq\varepsilon_0.
    \label{eq:negation-of-sequence-limit}
\end{equation}
所以序列发散的完整表述是：\emph{对每个 $a\in\mathbb{R}$}，都满足式~\eqref{eq:negation-of-sequence-limit}；其中 $\varepsilon_0$ 可以随 $a$ 改变，但不能随 $N$ 改变。只证明不趋于一个候选值，还没有排除其他实数成为极限的可能。

\begin{example}
\label{ex:alternating-sequence-divergence-by-definition}
用定义证明 $x_n=(-1)^n$ 发散。
\end{example}

\begin{solve}
任取 $a\in\mathbb{R}$。由 $2=|1-(-1)|\leq|1-a|+|-1-a|$，至少有 $|1-a|\geq1$ 或 $|-1-a|\geq1$。取 $\varepsilon_0=1$；前一种情形可在任意 $N$ 后选偶数下标，后一种情形可选奇数下标，都得到 $|x_n-a|\geq1$。因此任意实数 $a$ 都不是极限，序列发散。
\end{solve}

\subsection{无穷小量}
\label{subsec:infinitesimal-sequences}

\begin{definition}
\label{def:infinitesimal-sequence}
若实数列 $\{\alpha_n\}$ 满足 $\alpha_n\to0$，则称它为\emph{无穷小量}（infinitesimal），也记为 $\alpha_n=o(1)$（$n\to\infty$）。
\end{definition}

无穷小量是变化过程的性质，不是某个固定实数的大小。非零常数即使很小，组成的常值序列也不是无穷小量；恒为零的序列则是无穷小量。由极限定义中的绝对值可直接得到
\begin{equation}
    \alpha_n\to0\ \Longleftrightarrow\ |\alpha_n|\to0,
    \qquad
    x_n\to a\ \Longleftrightarrow\ x_n=a+\alpha_n,\quad\alpha_n\to0.
    \label{eq:sequence-limit-and-infinitesimal-error}
\end{equation}
后一关系把一般收敛问题转化为误差是否为无穷小量的问题。若 $c\in\mathbb{R}$ 固定，则 $c\alpha_n\to0$：$c=0$ 时直接成立，$c\neq0$ 时把原来的误差容限改成 $\varepsilon/|c|$ 即可。无穷小量的更多运算性质将在第~\ref{sec:properties-of-sequence-limits} 节给出。

\begin{example}
\label{ex:factorial-dominates-geometric-sequence}
设 $a>1$，证明 $a^n/n!$ 是无穷小量，其中 $n!=1\cdot2\cdots n$。
\end{example}

\begin{solve}
记 $u_n=a^n/n!>0$。选正整数 $K$ 使 $K+1\geq2a$，则对 $n\geq K$，
\[
    \frac{u_{n+1}}{u_n}=\frac{a}{n+1}\leq\frac12.
\]
反复使用这一估计，得 $0<u_n\leq u_K(1/2)^{n-K}$（$n\geq K$）。由例~\ref{ex:geometric-sequence-limit}，右端的几何因子趋于零，因此给定 $\varepsilon>0$，可选 $N\geq K$，使所有 $n>N$ 都满足 $u_K(1/2)^{n-K}<\varepsilon$，于是 $u_n<\varepsilon$。这直接按定义证明了 $u_n\to0$，也说明比较相邻项有时能把复杂表达式的误差控制转化为几何衰减。
\end{solve}

\subsection{无穷大量}
\label{subsec:infinite-sequences}

并非所有发散序列都具有相同的变化方式。$(-1)^n$ 始终在两个值之间振荡，$n$ 则最终超过每一个固定的界，$(-1)^n n$ 虽然符号交替，绝对值也最终超过每一个界。后两类变化由无穷大量描述。

\begin{definition}
\label{def:infinite-sequence}
若实数列 $\{x_n\}$ 满足
\begin{equation}
    \forall M>0,\quad\exists N\in\mathbb{N},\quad
    n>N\ \Longrightarrow\ x_n>M,
    \label{eq:positive-infinite-sequence}
\end{equation}
则称它为\emph{正无穷大量}，记为 $x_n\to+\infty$。将结论换成 $x_n<-M$，定义\emph{负无穷大量}，记为 $x_n\to-\infty$；换成 $|x_n|>M$，定义\emph{无穷大量}（infinitely large sequence），记为 $x_n\to\infty$，即 $|x_n|\to+\infty$。其中不带正负号的 $\infty$ 表示绝对值趋于无穷，不要求符号最终固定。
\end{definition}

正、负无穷大量都是无穷大量，反过来不一定成立，如 $(-1)^n n$。趋于无穷的序列在有限极限意义下仍然发散；若允许有限极限以及上述无穷趋向，则讨论的是\emph{广义极限}（generalized limit）。在采用扩展实数系 $\mathbb{R}\cup\{-\infty,+\infty\}$ 时，还要注意 $(-1)^n n$ 没有单一的扩展实数极限，尽管它满足这里的绝对值趋于无穷。

\emph[caution]{无界只要求能找到任意大的绝对值，无穷大量则要求所有充分靠后的项都具有任意大的绝对值。}例如，令 $x_{2k}=2k$、$x_{2k-1}=1$，则序列无界，也不能收敛于任何实数：对任意候选值 $a$，偶数项最终与它的距离超过 $1$。但奇数项始终为 $1$，所以它不是无穷大量。无界序列必发散还将在下一节由“收敛必有界”统一说明。

\begin{example}
\label{ex:harmonic-sequence-positive-infinity}
证明 $s_n=\sum_{j=1}^n1/j\to+\infty$。
\end{example}

\begin{solve}
按下标位于 $2$ 的相邻幂之间分组。第 $j$ 组从 $2^{j-1}+1$ 到 $2^j$，有 $2^{j-1}$ 项，每项至少为 $1/2^j$，因此
\[
    s_{2^k}=1+\sum_{j=1}^k\sum_{r=2^{j-1}+1}^{2^j}\frac1r
    \geq1+\frac{k}{2}.
\]
任给 $M>0$，选正整数 $k$ 使 $1+k/2>M$，再令 $N=2^k$。由于所有加项为正，对每个 $n>N$ 都有 $s_n\geq s_{2^k}>M$，故 $s_n\to+\infty$。证明不仅找到了任意大的若干项，还利用正项累加控制了它们后面的全部项。
\end{solve}

\begin{proposition}[无穷小量与倒数的关系]
\label{prop:infinitesimal-sequence-reciprocal}
若 $x_n$ 从某项起恒不为零，则
\begin{equation}
    x_n\to0\quad\Longleftrightarrow\quad
    \left|\frac1{x_n}\right|\to+\infty,
    \qquad
    |x_n|\to+\infty\quad\Longleftrightarrow\quad\frac1{x_n}\to0.
    \label{eq:sequence-infinitesimal-reciprocal}
\end{equation}
只需对倒数有定义的尾部讨论；有限个起始位置不影响这些极限。
\end{proposition}

\begin{proof}
对任意 $M>0$，不等式 $0<|x_n|<1/M$ 等价于 $|1/x_n|>M$，给出第一条等价关系。第二条同理由 $|x_n|>1/\varepsilon$ 与 $|1/x_n|<\varepsilon$ 的等价性得到。趋于无穷时，取 $M=1$ 即能保证原序列最终非零。
\end{proof}

\begin{theorem}[算术平均保持极限]
\label{thm:cesaro-mean-of-convergent-sequence}
设 $\{x_n\}$ 为实数列，记前 $n$ 项的算术平均为
\[
    A_n=\frac{x_1+\cdots+x_n}{n}.
\]
若 $x_n\to a\in\mathbb{R}$，则 $A_n\to a$；若 $x_n\to+\infty$ 或 $x_n\to-\infty$，则分别有 $A_n\to+\infty$ 或 $A_n\to-\infty$。有限极限的这一结论也称\emph{切萨罗平均定理}（Cesàro mean theorem）。
\end{theorem}

\begin{proof}
先设 $a$ 有限。给定 $\varepsilon>0$，选正整数 $K$ 使 $k>K$ 时 $|x_k-a|<\varepsilon/2$，并记固定的有限和 $C=\sum_{k=1}^K|x_k-a|$。对 $n>K$，
\[
    |A_n-a|\leq\frac{C}{n}+\frac1n\sum_{k=K+1}^n|x_k-a|
    \leq\frac{C}{n}+\frac{\varepsilon}{2}.
\]
再选正整数 $N>\max\{K,2C/\varepsilon\}$，便有 $n>N$ 时 $|A_n-a|<\varepsilon$。证明将有限个可能较大的起始误差与统一较小的尾部误差分开控制。

若 $x_n\to+\infty$，给定 $M>0$，选 $K$ 使 $k>K$ 时 $x_k>4M$，记 $S=x_1+\cdots+x_K$。当 $n>\max\{2K,|S|/M\}$ 时，
\[
    A_n=\frac{S}{n}+\frac1n\sum_{k=K+1}^n x_k
    >-M+4M\frac{n-K}{n}>M,
\]
故 $A_n\to+\infty$。若 $x_n\to-\infty$，对 $-x_n$ 使用已证结论即可。
\end{proof}

这个定理的逆命题不成立：$x_n=(-1)^n$ 不收敛，但其算术平均的绝对值不超过 $1/n$，所以 $A_n\to0$。此外，条件 $|x_n|\to+\infty$ 不能替代具有确定符号的无穷趋向；例如 $x_n=(-1)^n n$ 的平均满足
\[
    A_{2k}=\frac12,\qquad A_{2k-1}=-\frac{k}{2k-1},
\]
始终有界，因此并非无穷大量。

\section{序列极限的性质}
\label{sec:properties-of-sequence-limits}

% 根据教材第 2.2 节整理，第 40--49 页；分章 PDF 第 12--21 页。
极限定义适合直接估计误差，但面对复杂表达式，每次重新寻找 $N$ 会使计算繁琐。本节从定义推出极限的基本性质与运算法则，把新的极限问题分解为已知问题；使用这些法则时，必须先核对有关极限是否存在，以及分母是否满足条件。

\begin{definition}
\label{def:bounded-sequence}
若实数列 $\{x_n\}$ 满足
\[
    \exists M>0,\quad\forall n\in\mathbb{N},\quad |x_n|\leq M,
\]
则称它为\emph{有界序列}（bounded sequence），也记为 $x_n=O(1)$（$n\to\infty$）。这等价于取值集合 $\{x_n:n\in\mathbb{N}\}$ 有界。上界、下界及相应的无界性均按\crosschapterref{1}{定义~\ref{def:upper-and-lower-bounds}}{上、下界的定义} 和\crosschapterref{1}{定义~\ref{def:bounded-and-unbounded-sets}}{有界集与无界集的定义} 理解。
\end{definition}

\begin{theorem}[有限项不变性、唯一性与有界性]
\label{thm:basic-properties-of-sequence-limits}
改变序列的有限多项不改变其敛散性，收敛时也不改变极限。每个收敛实数列的极限唯一，并且序列有界。
\end{theorem}

\begin{proof}
若两个序列从第 $K$ 项后相同，只需把原来的误差阈值 $N$ 换成 $\max\{N,K\}$，两者便满足同一个极限估计，给出有限项不变性。

若同一序列同时趋于 $a,b$ 且 $a\neq b$，取 $\varepsilon=|a-b|/3$，使充分大的 $n$ 同时满足 $|x_n-a|<\varepsilon$、$|x_n-b|<\varepsilon$，则
\[
    |a-b|\leq|a-x_n|+|x_n-b|<\frac23|a-b|,
\]
矛盾，故极限唯一。

若 $x_n\to a$，在定义中取 $\varepsilon=1$，得到正整数 $N$，使 $n>N$ 时 $|x_n|\leq|a|+|x_n-a|<|a|+1$。令
\[
    M=\max\{|x_1|,\ldots,|x_N|,|a|+1\},
\]
便有 $|x_n|\leq M$ 对所有 $n$ 成立。尾部的统一界与有限个起始值共同给出了整个序列的界。
\end{proof}

因此，删去或增加有限个首项也不会改变极限；特别地，对固定非负整数 $r$，若 $x_n\to a$，则 $x_{n+r}\to a$，反之亦然。这是后文在递推式两边取极限时需要核对的一点。\emph[caution]{有界性是收敛的必要条件，不是充分条件。}例如，$(-1)^n$ 有界而发散；无界序列则由上述定理必定发散。有限项不变性还说明，只要尾部有界，补上前面有限项后仍然有界。

\begin{theorem}[保序性与最终保号性]
\label{thm:sequence-limit-order-properties}
设 $x_n\to a$、$y_n\to b$，其中 $a,b\in\mathbb{R}$。若 $a<b$，则对任意 $c\in(a,b)$，存在正整数 $N$，使 $n>N$ 时 $x_n<c<y_n$。若从某项起恒有 $x_n\leq y_n$，则 $a\leq b$。特别地，若 $a>0$，则对每个 $c\in(0,a)$，最终有 $x_n>c$；若 $a<0$，则有对应的负值结论。
\end{theorem}

\begin{proof}
当 $a<b$ 时，取 $\varepsilon=\min\{c-a,b-c\}>0$，并使两个序列在同一阈值后均满足各自的误差估计，便得 $x_n<a+\varepsilon\leq c$、$y_n>b-\varepsilon\geq c$。若最终有 $x_n\leq y_n$ 而 $a>b$，将已证结论用于 $y_n,x_n$ 及中间值 $(a+b)/2$，就会得到最终 $y_n<x_n$，矛盾。取一个适当的常值序列作比较，便得到保号结论。
\end{proof}

\emph[caution]{取极限通常只能保留非严格不等式。}例如，$1-1/n<1+1/n$ 对每个 $n$ 成立，但两侧都趋于 $1$。反过来，若极限 $a\neq0$，则最终可使 $|x_n-a|<|a|/2$，从而
\begin{equation}
    \frac{|a|}{2}<|x_n|<\frac{3|a|}{2}.
    \label{eq:nonzero-sequence-limit-away-from-zero}
\end{equation}
这比“最终非零”更强：它给出了绝对值的统一正下界，是控制倒数及商的关键。

\begin{theorem}[极限的四则运算]
\label{thm:sequence-limit-arithmetic}
设 $x_n\to a$、$y_n\to b$，其中 $a,b\in\mathbb{R}$，则
\begin{equation}
    \begin{aligned}
        \lim_{n\to\infty}(x_n\pm y_n)&=a\pm b,\\
        \lim_{n\to\infty}x_ny_n&=ab,\\
        \lim_{n\to\infty}\frac{x_n}{y_n}&=\frac ab\qquad(b\neq0).
    \end{aligned}
    \label{eq:sequence-limit-arithmetic}
\end{equation}
商在 $y_n\neq0$ 的尾部上讨论，有限个分母为零的起始位置不影响尾部极限。
\end{theorem}

\begin{proof}
对于和、差，利用
\[
    |(x_n\pm y_n)-(a\pm b)|\leq|x_n-a|+|y_n-b|,
\]
使右端两项分别小于 $\varepsilon/2$ 即可。

对于乘积，\emph[tip]{通过加减中间项，把乘积误差分解成已知的两个收敛误差。}由 $\{y_n\}$ 有界，可取 $M\geq1$ 使所有 $|y_n|\leq M$，于是
\[
    |x_ny_n-ab|
    =|y_n(x_n-a)+a(y_n-b)|
    \leq M|x_n-a|+|a|\,|y_n-b|.
\]
使 $|x_n-a|<\varepsilon/(2M)$ 且 $|y_n-b|<\varepsilon/(2(|a|+1))$ 同时成立，右端即小于 $\varepsilon$。

若 $b\neq0$，由式~\eqref{eq:nonzero-sequence-limit-away-from-zero}，最终有 $|y_n|>|b|/2$，故
\[
    \left|\frac1{y_n}-\frac1b\right|
    =\frac{|y_n-b|}{|b|\,|y_n|}
    \leq\frac2{|b|^2}|y_n-b|\to0.
\]
再对 $x_n\cdot(1/y_n)$ 使用乘积法则，便得到商的结论。
\end{proof}

四则运算可以反复用于有限个收敛序列，特别地，$c x_n\to ca$、$x_n^m\to a^m$（$m$ 为固定正整数）。由反向三角不等式，还有
\[
    \bigl||x_n|-|a|\bigr|\leq|x_n-a|,
\]
所以 $|x_n|\to|a|$；反向不能直接成立，例如 $|(-1)^n|\to1$ 而原序列发散。\emph[caution]{这些运算法则的前提是各部分已有相应的有限极限；不能用“拆开后出现无穷”的形式计算替代条件核对。}

\begin{example}
\label{ex:sequence-limit-arithmetic-applications}
求
\[
    \lim_{n\to\infty}\frac{3n^2+4n-100}{4n^2+5n+105},
    \qquad
    \lim_{n\to\infty}(1+q+\cdots+q^n)\quad(|q|<1).
\]
\end{example}

\begin{solve}
第一问将分子、分母同除以 $n^2$，得到
\[
    \frac{3+4/n-100/n^2}{4+5/n+105/n^2}\to\frac34,
\]
此时分母的极限为 $4\neq0$，商的法则才可使用。第二问先用有限几何和的恒等式，再由 $q^{n+1}\to0$，得到
\[
    1+q+\cdots+q^n=\frac{1-q^{n+1}}{1-q}\to\frac1{1-q}.
\]
这里处理的是含有 $n+1$ 项的整个有限和，不能将每一项分别取极限后任意交换无限次求和与极限。
\end{solve}

\begin{theorem}[夹逼收敛原理]
\label{thm:squeeze-principle-for-sequences}
设三个实数列从某项起满足 $x_n\leq z_n\leq y_n$。若 $x_n\to a$、$y_n\to a$，其中 $a\in\mathbb{R}$，则 $z_n\to a$。
\end{theorem}

\begin{proof}
给定 $\varepsilon>0$，取足够大的阈值，使其后既满足夹逼不等式，又满足 $|x_n-a|<\varepsilon$ 与 $|y_n-a|<\varepsilon$，于是
\[
    a-\varepsilon<x_n\leq z_n\leq y_n<a+\varepsilon,
\]
即 $|z_n-a|<\varepsilon$。
\end{proof}

夹逼的关键在于两端趋于同一个数，而不在于给中间序列找到任意上、下界。常用的误差形式是：若 $|u_n-a|\leq v_n$ 最终成立且 $v_n\to0$，则 $u_n\to a$。特别地，\emph{无穷小量与有界序列的乘积仍为无穷小量}：若 $\alpha_n\to0$、$|b_n|\leq M$，则 $|\alpha_n b_n|\leq M|\alpha_n|\to0$，这里 $\{b_n\}$ 本身不必收敛。

\begin{example}
\label{ex:nth-root-limits-by-squeeze}
证明
\[
    \lim_{n\to\infty}\sqrt[n]{n}=1,
    \qquad
    \lim_{n\to\infty}\left(a_1^n+\cdots+a_m^n\right)^{1/n}
    =\max\{a_1,\ldots,a_m\},
\]
其中 $m$ 是固定正整数，$a_1,\ldots,a_m$ 是固定正实数。
\end{example}

\begin{solve}
对 $n\geq2$，令 $h_n=\sqrt[n]{n}-1>0$。由二项式展开中各项非负，
\[
    n=(1+h_n)^n\geq\frac{n(n-1)}2h_n^2,
    \qquad 0<h_n\leq\sqrt{\frac2{n-1}}.
\]
右端趋于零可由定义核实：给定 $\varepsilon>0$，只需 $n>1+2/\varepsilon^2$。故夹逼原理给出 $h_n\to0$，即 $\sqrt[n]{n}\to1$。例~\ref{ex:fixed-root-sequence-limit} 只讨论固定底数，不能直接把其中的 $a$ 换成随 $n$ 变化的 $n$。

记 $A=\max\{a_1,\ldots,a_m\}>0$，则
\[
    A^n\leq a_1^n+\cdots+a_m^n\leq mA^n,
\]
从而
\[
    A\leq\left(a_1^n+\cdots+a_m^n\right)^{1/n}\leq A\sqrt[n]{m}.
\]
因为 $m$ 固定，由 $\sqrt[n]{m}\to1$ 及夹逼原理，得到所求极限。该结论表示，取很高的正整数幂后，最大项控制了这个有限和的 $n$ 次方根。
\end{solve}

另一个常用结论是：若 $u_n\to0$，$p$ 为固定正整数，则 $\sqrt[p]{1+u_n}\to1$，其中只需考察 $|u_n|<1$ 的尾部，使被开方数为正。事实上，正数开方保持大小关系，且 $0\leq t\leq1$ 时 $t^p\leq t$、$t\geq1$ 时 $t^p\geq t$，所以
\[
    1-|u_n|\leq\sqrt[p]{1-|u_n|}
    \leq\sqrt[p]{1+u_n}\leq\sqrt[p]{1+|u_n|}\leq1+|u_n|.
\]
两端都趋于 $1$，故结论成立。这一证明不需要预先引用后续的函数连续性理论。

\begin{example}
\label{ex:recurrence-sequence-geometric-error-bound}
设 $a_1=2$，$a_{n+1}=2+1/a_n$（$n\geq1$）。证明序列收敛并求其极限。
\end{example}

\begin{analyze}
若序列收敛于非零数 $L$，四则运算法则要求 $L=2+1/L$，其正根为 $L=1+\sqrt2$。这一计算只找到候选值，还没有证明收敛；下面直接估计 $|a_n-L|$，使误差被一个趋于零的几何序列控制。
\end{analyze}

\begin{solve}
由递推式归纳可得 $a_n\geq2$，故分母始终为正。令 $L=1+\sqrt2$，利用 $L=2+1/L$，有
\[
    a_{n+1}-L=\frac1{a_n}-\frac1L=-\frac{a_n-L}{La_n}.
\]
因此，令 $r=1/(2L)\in(0,1)$，得到
\[
    |a_{n+1}-L|\leq r|a_n-L|,
    \qquad |a_n-L|\leq r^{n-1}|a_1-L|.
\]
右端趋于零，由夹逼原理得 $a_n\to1+\sqrt2$。证明把候选值用于构造误差，没有把“假设收敛后满足的方程”当成收敛性的依据。
\end{solve}

由四则运算法则，有限个无穷小量的和、差与积仍为无穷小量，但商不能只由“分子、分母都趋于零”来判断。以共同分母 $y_n=1/n$ 为例，分别取分子 $1/n^2$、$1/n$、$1/\sqrt n$、$(-1)^n/n$，它们都趋于零，对应的商却分别趋于 $0$、$1$、$+\infty$，或振荡而无极限。因而 $0/0$ 只是\emph{不定式}（indeterminate form）的形式记号，并不是数值。

无穷大量的运算也须区分确定情形与不定情形。同趋于 $+\infty$ 的两列相加仍趋于 $+\infty$，同趋于 $-\infty$ 时相加仍趋于 $-\infty$；两个绝对值趋于无穷的序列相乘，乘积的绝对值仍趋于无穷。若 $|x_n|\to+\infty$ 且最终 $|y_n|\geq c>0$，则 $|x_ny_n|\geq c|x_n|\to+\infty$，这里不必要求 $y_n$ 另外有上界。反之，$\infty-\infty$、$0\cdot\infty$、$\infty/\infty$ 等形式一般不能确定结果；例如，$n-(n-1)=1$ 与 $n-2n=-n$ 已给出不同趋向。所有这些记号都应还原为具体序列及其极限条件，不能把 $\infty$ 当作实数参与运算。

\begin{proposition}[多项式序列由最高次项决定无穷趋向]
\label{prop:polynomial-sequence-leading-term-infinity}
设 $P(n)=a_kn^k+a_{k-1}n^{k-1}+\cdots+a_0$，其中 $k\geq1$、$a_k\neq0$。则 $a_k>0$ 时 $P(n)\to+\infty$，$a_k<0$ 时 $P(n)\to-\infty$。
\end{proposition}

\begin{proof}
将最高次幂提出，得到
\[
    P(n)=n^k\left(a_k+\frac{a_{k-1}}n+\cdots+\frac{a_0}{n^k}\right).
\]
括号中的序列由有限次运算法则趋于 $a_k$。若 $a_k>0$，则最终括号大于 $a_k/2>0$，而 $n^k\to+\infty$，故 $P(n)\to+\infty$；若 $a_k<0$，对 $-P(n)$ 使用同一结论即可。这个证明避免了把各项分别记成无穷后相加，因为低次项的系数可能与最高次项异号，直接相加会出现尚未确定的无穷差。
\end{proof}

\section{单调收敛原理}
\label{sec:monotone-convergence}

\subsection{单调收敛原理}
\label{subsec:monotone-convergence-principle}

% 根据教材第 2.3.1 小节整理，第 49--52 页；分章 PDF 第 21--24 页。
用极限定义验证收敛，通常要先找到候选极限；但对递推给出的序列，极限值往往并不容易预先看出。单调收敛原理提供了另一条途径：先利用单调性和有界性证明极限存在，再借助递推关系求出极限。以下讨论的都是实数列，下标取正整数。

\begin{definition}
\label{def:monotone-sequences}
若序列 $\{x_n\}$ 满足
\[
    x_n\leq x_{n+1},\qquad \forall n\geq1,
\]
则称它为\emph{单调递增序列}（monotonically increasing sequence），也称单调上升序列或非减序列；若满足
\[
    x_n\geq x_{n+1},\qquad \forall n\geq1,
\]
则称它为\emph{单调递减序列}（monotonically decreasing sequence），也称单调下降序列或非增序列。这两类序列统称为\emph{单调序列}（monotone sequence）。若相应的不等式对每个 $n$ 都严格成立，则分别称为严格递增序列或严格递减序列。
\end{definition}

这里的递增、递减都允许相邻项相等，常值序列同时属于这两类。将序列看作定义在正整数集上的函数，序列的单调性就是该函数的单调性；比较相邻两项已经足够，因为连续使用 $x_k\leq x_{k+1}$，便得到 $m<n$ 时 $x_m\leq x_n$。

\begin{theorem}[单调收敛原理]
\label{thm:monotone-convergence-principle}
单调有界的实数列必收敛。更具体地，若 $\{x_n\}$ 单调递增且有上界，则
\begin{equation}
    \lim_{n\to\infty}x_n=\sup\{x_n:n\geq1\};
    \label{eq:increasing-sequence-limit-supremum}
\end{equation}
若 $\{x_n\}$ 单调递减且有下界，则
\begin{equation}
    \lim_{n\to\infty}x_n=\inf\{x_n:n\geq1\}.
    \label{eq:decreasing-sequence-limit-infimum}
\end{equation}
这里的极限都是有限实数。
\end{theorem}

\begin{proof}
先设 $\{x_n\}$ 单调递增且有上界。由实数的确界定理，非空有上界的集合 $E=\{x_n:n\geq1\}$ 存在上确界，记为 $a=\sup E$。\emph[tip]{先用上确界的逼近性质找到一项靠近 $a$，再用单调性控制它后面的全部项。}事实上，$x_n\leq a$ 对每个 $n$ 成立；又由\crosschapterref{1}{定义~\ref{def:supremum}}{上确界定义}，任给 $\varepsilon>0$，存在正整数 $N$，使得
\[
    a-\varepsilon<x_N\leq a.
\]
当 $n\geq N$ 时，由单调性得
\[
    a-\varepsilon<x_N\leq x_n\leq a,
\]
故 $|x_n-a|=a-x_n<\varepsilon$。这就证明了 $x_n\to a$。

再设 $\{x_n\}$ 单调递减且有下界，令 $b=\inf E$。由\crosschapterref{1}{定义~\ref{def:infimum}}{下确界定义}，任给 $\varepsilon>0$，存在正整数 $N$，使 $b\leq x_N<b+\varepsilon$。当 $n\geq N$ 时，
\[
    b\leq x_n\leq x_N<b+\varepsilon,
\]
所以 $|x_n-b|=x_n-b<\varepsilon$，即 $x_n\to b$。
\end{proof}

递增序列自动以首项 $x_1$ 为下界，因而只需另外寻找上界；递减序列自动以首项为上界，只需另外寻找下界。证明中的确界定理体现了实数系的完备性：单调性保证序列始终朝同一方向变化，而有限的确界为它提供了最终趋近的位置。\emph[caution]{一个任意选取的上界不一定是递增序列的极限，极限对应的是其取值集合的上确界。}例如，$x_n=1-1/n$ 有上界 $2$，但它的上确界及极限都是 $1$。

若单调递增序列收敛于 $a$，则它的每一项都不超过 $a$；若某一项 $x_N=a$，则对所有 $n\geq N$，都有 $a=x_N\leq x_n\leq a$，所以此后的各项恒等于 $a$。因此，这类序列或者始终严格小于极限，或者从某一项起恒等于极限；递减序列有对应的结论。

\begin{corollary}
\label{cor:monotone-sequence-generalized-limit}
每个单调实数列都有广义极限：单调递增且无上界时，$x_n\to+\infty$；单调递减且无下界时，$x_n\to-\infty$。若允许确界取扩展实数值，则式~\eqref{eq:increasing-sequence-limit-supremum} 和式~\eqref{eq:decreasing-sequence-limit-infimum} 对无界的相应情形也成立。
\end{corollary}

\begin{proof}
若 $\{x_n\}$ 单调递增且无上界，则任给 $M>0$，存在正整数 $N$，使 $x_N>M$。当 $n\geq N$ 时，$x_n\geq x_N>M$，故 $x_n\to+\infty$。若序列单调递减且无下界，则任给 $M>0$，可选 $N$ 使 $x_N<-M$；对所有 $n\geq N$，有 $x_n\leq x_N<-M$，故 $x_n\to-\infty$。
\end{proof}

\emph[caution]{趋于无穷不属于通常所说的收敛，单调性和有界性也不能单独保证收敛。}例如，$x_n=n$ 单调递增却无上界，$x_n=(-1)^n$ 有界却不收敛；反过来，收敛序列也不一定单调，例如 $x_n=(-1)^n/n\to0$。若序列只是从某一项起单调且有界，仍可对这一尾部应用单调收敛原理，因为删去有限个首项不改变收敛性及极限。

\begin{example}
\label{ex:monotone-square-root-iteration}
% 教材例 2.3.1，第 50--51 页。
设 $A>0$、$x_1>0$，递推定义
\begin{equation}
    x_{n+1}=\frac12\left(x_n+\frac{A}{x_n}\right),\qquad n\geq1.
    \label{eq:monotone-square-root-iteration}
\end{equation}
证明 $\{x_n\}$ 收敛，并求出它的极限。
\end{example}

\begin{solve}
由 $x_1>0$ 和递推式归纳可知，每一项都为正，故递推过程始终有定义。又因为
\[
    x_{n+1}-\sqrt A
    =\frac{x_n^2-2\sqrt A\,x_n+A}{2x_n}
    =\frac{(x_n-\sqrt A)^2}{2x_n}\geq0,
\]
所以 $x_n\geq\sqrt A$ 对所有 $n\geq2$ 成立。在此基础上，对 $n\geq2$ 有
\[
    x_{n+1}-x_n
    =\frac{A-x_n^2}{2x_n}\leq0.
\]
因此，尾部序列 $\{x_n\}_{n\geq2}$ 单调递减且有下界 $\sqrt A$。由定理~\ref{thm:monotone-convergence-principle}，原序列收敛；记极限为 $a$，则 $a\geq\sqrt A>0$。

现在才可在式~\eqref{eq:monotone-square-root-iteration} 两边取极限。由于 $a>0$，倒数的极限运算法则适用，得到
\[
    a=\frac12\left(a+\frac{A}{a}\right),
    \qquad a^2=A.
\]
结合 $a>0$，可知 $\lim_{n\to\infty}x_n=\sqrt A$。注意，当 $x_1<\sqrt A$ 时，第一步会使 $x_2>x_1$，所以这里应证明从第二项起递减，不能直接声称整个序列递减。
\end{solve}

这个递推过程给出了求平方根的近似方法。例如，取 $A=2$、$x_1=2$，依次得到
\[
    x_2=\frac32,\qquad x_3=\frac{17}{12},\qquad
    x_4=\frac{577}{408}\approx1.41421569.
\]
上述恒等式还给出了误差关系 $x_{n+1}-\sqrt A=(x_n-\sqrt A)^2/(2x_n)$，说明每一步近似误差如何由前一步的误差决定。

\begin{example}
\label{ex:recurrence-fixed-point-without-convergence}
% 教材例 2.3.2，第 51 页。
设 $x_1=1$、$x_{n+1}=-x_n$。说明为什么不能仅凭递推式中形式上的极限方程判断 $\{x_n\}$ 的极限。
\end{example}

\begin{solve}
由递推式得到 $x_n=(-1)^{n-1}$，其奇数项恒为 $1$，偶数项恒为 $-1$，因而序列不收敛。若未经证明便设 $x_n\to a$，再在递推式两边取极限，就会得到 $a=-a$，进而算出 $a=0$；但这个数并不是该序列的极限。方程 $a=-a$ 只说明：\emph{如果序列有有限极限，那么该极限只能是 $0$}，并没有证明极限存在。
\end{solve}

一般地，对递推式 $x_{n+1}=f(x_n)$，若已证明 $x_n\to a$，还须确认 $f(x_n)\to f(a)$，才能得到极限方程 $a=f(a)$；例如可核对所用极限运算法则的条件，或在后续连续性理论中核对 $f$ 在 $a$ 处连续。仅求出方程的解，不能替代对序列收敛性的证明。

\begin{example}
\label{ex:monotone-odd-even-golden-ratio}
% 教材例 2.3.3，第 51--52 页。
设 $a_1=1$，并且
\begin{equation}
    a_{n+1}=1+\frac1{a_n},\qquad n\geq1.
    \label{eq:golden-ratio-recurrence}
\end{equation}
证明 $\{a_n\}$ 收敛，并求其极限。
\end{example}

\begin{analyze}
前五项为 $1,2,3/2,5/3,8/5$，它们满足 $a_1<a_3<a_5<a_4<a_2$。原序列并不单调，但奇数项可能递增、偶数项可能递减。\emph[tip]{对隔一项的递推关系应用单调收敛原理，再证明两个子列趋于同一个数。}
\end{analyze}

\begin{solve}
首先归纳证明 $1\leq a_n\leq2$：首项满足此条件；若 $1\leq a_n\leq2$，则 $3/2\leq1+1/a_n\leq2$，故下一项仍满足此条件。特别地，所有分母均为正，序列有界。

连续使用两次式~\eqref{eq:golden-ratio-recurrence}，得到
\begin{equation}
    a_{n+2}=1+\frac1{1+1/a_n}
    =2-\frac1{a_n+1},\qquad n\geq1.
    \label{eq:golden-ratio-two-step-recurrence}
\end{equation}
因此对每个 $n\geq1$，有
\[
    a_{n+4}-a_{n+2}
    =\frac1{a_n+1}-\frac1{a_{n+2}+1}
    =\frac{a_{n+2}-a_n}{(a_n+1)(a_{n+2}+1)}.
\]
分母为正，所以相邻两次的同奇偶项之差同号。由 $a_3-a_1=1/2>0$ 和 $a_4-a_2=-1/3<0$ 归纳可知，奇数项子列 $\{a_{2k-1}\}$ 严格递增，偶数项子列 $\{a_{2k}\}$ 严格递减。两个子列均有界，由单调收敛原理，存在 $\alpha,\beta\in[1,2]$，使
\[
    \lim_{k\to\infty}a_{2k-1}=\alpha,\qquad
    \lim_{k\to\infty}a_{2k}=\beta.
\]
分别沿奇数、偶数下标在式~\eqref{eq:golden-ratio-two-step-recurrence} 中取极限，得到
\[
    \alpha=2-\frac1{\alpha+1},\qquad
    \beta=2-\frac1{\beta+1}.
\]
两者都满足 $t^2-t-1=0$。由于 $\alpha,\beta>0$，负根被排除，故
\[
    \alpha=\beta=\frac{1+\sqrt5}{2}=:L.
\]

最后验证整个序列也趋于 $L$。任给 $\varepsilon>0$，由两个子列都趋于 $L$，存在正整数 $K$，使 $k\geq K$ 时同时有 $|a_{2k-1}-L|<\varepsilon$ 和 $|a_{2k}-L|<\varepsilon$。每个 $n\geq2K$ 都属于这两类下标之一，因而均满足 $|a_n-L|<\varepsilon$。所以
\[
    \lim_{n\to\infty}a_n=\frac{1+\sqrt5}{2}.\qedhere
\]
\end{solve}

这三个例子体现了单调收敛原理的使用次序：先检查递推式是否始终有定义，再寻找单调性与适当方向的界；原序列不单调时，可以考察从某一项起的尾部，或分别考察奇、偶子列。确认有关序列收敛后，才利用递推式求候选值，并结合取值范围排除不合条件的根；若拆成子列，还必须说明这些子列的极限相同且能控制原序列的全部充分靠后的项。

\subsection{无理数 \texorpdfstring{$e$}{e} 和欧拉常数 \texorpdfstring{$\gamma$}{gamma}}
\label{subsec:e-and-euler-constant}

\section{实数系连续性的基本定理}
\label{sec:fundamental-theorems-of-real-continuity}

实数系的连续性可以通过闭区间套、有限覆盖、聚点及柯西序列等不同形式来表述，这些结论各有适合的应用场合。本节沿用教材的论证顺序：由单调收敛原理得到闭区间套定理，再依次证明有限覆盖定理、聚点原理、有界序列的子列定理和柯西收敛准则。以下序列的下标均取正整数。

\subsection{闭区间套定理}
\label{subsec:nested-interval-theorem}

不断缩小一个非空闭区间，并要求后一个区间包含在前一个区间中，直观上应当能保留至少一个公共点。闭区间套定理把这一事实表述为严格的存在性结论；若区间长度还趋于零，则公共点唯一。

\begin{definition}
\label{def:nested-closed-intervals}
设 $I_n=[a_n,b_n]$ 为非空有界闭区间，即 $a_n,b_n\in\mathbb{R}$ 且 $a_n\leq b_n$。若
\[
    I_{n+1}\subseteq I_n\qquad(n=1,2,\ldots),
\]
则称 $\{I_n\}$ 为\emph{闭区间套}（nested closed intervals）。这一包含关系等价于
\[
    a_n\leq a_{n+1}\leq b_{n+1}\leq b_n.
\]
下面的定理在嵌套条件之外，另要求区间长度趋于零。
\end{definition}

\begin{theorem}[闭区间套定理]
\label{thm:nested-closed-intervals}
设 $\{[a_n,b_n]\}$ 为一列非空有界闭区间，满足
\[
    [a_{n+1},b_{n+1}]\subseteq[a_n,b_n]\quad(n=1,2,\ldots),
    \qquad
    \lim_{n\to\infty}(b_n-a_n)=0.
\]
则存在唯一实数 $c$ 属于每个区间，即
\[
    \bigcap_{n=1}^{\infty}[a_n,b_n]=\{c\},
\]
并且
\[
    \lim_{n\to\infty}a_n=\lim_{n\to\infty}b_n=c.
\]
\end{theorem}

\begin{proof}
由嵌套条件，有 $a_n\leq a_{n+1}\leq b_{n+1}\leq b_n$，所以 $\{a_n\}$ 单调上升且以 $b_1$ 为上界，$\{b_n\}$ 单调下降且以 $a_1$ 为下界。\emph[tip]{分别对左右端点应用单调收敛原理。}记
\[
    a=\lim_{n\to\infty}a_n=\sup\{a_n:n\geq1\},
    \qquad
    b=\lim_{n\to\infty}b_n=\inf\{b_n:n\geq1\}.
\]
对任意 $m,n$，若 $m\leq n$，则 $a_m\leq a_n\leq b_n$；若 $m>n$，则 $a_m\leq b_m\leq b_n$。因此每个右端点都是左端点集合的上界，每个左端点都是右端点集合的下界，从而
\[
    a_n\leq a\leq b\leq b_n\qquad(n=1,2,\ldots).
\]
又由长度趋于零及极限的运算法则，
\[
    b-a=\lim_{n\to\infty}(b_n-a_n)=0.
\]
令 $c=a=b$，便有 $a_n\leq c\leq b_n$，所以 $c$ 属于每个区间。若另有 $c'$ 也属于所有区间，则 $a_n\leq c'\leq b_n$ 对每个 $n$ 成立；由夹逼收敛原理，$c'=c$，故公共点唯一。
\end{proof}

\begin{example}
\label{ex:sqrt-two-bisection}
利用\emph{二分法}（bisection method）构造一列闭区间，证明存在唯一的正实数 $c$ 满足 $c^2=2$，并给出用区间中点近似 $c$ 时的误差估计。
\end{example}

\begin{solve}
取 $a_1=1$、$b_1=2$，令 $I_1=[a_1,b_1]$。已得到区间 $I_n=[a_n,b_n]$ 后，取其中点 $m_n=(a_n+b_n)/2$，根据 $m_n^2$ 与 $2$ 的大小关系定义
\begin{equation}
    I_{n+1}=[a_{n+1},b_{n+1}]=
    \begin{cases}
        [m_n,b_n], & m_n^2\leq2,\\
        [a_n,m_n], & m_n^2>2.
    \end{cases}
    \label{eq:sqrt-two-bisection-update}
\end{equation}
\emph[tip]{每次保留端点平方仍分处 $2$ 两侧的半区间。}初始时 $a_1^2\leq2\leq b_1^2$；若 $a_n^2\leq2\leq b_n^2$，则第一种情形有 $m_n^2\leq2\leq b_n^2$，第二种情形有 $a_n^2\leq2<m_n^2$。因此由归纳可知，对所有 $n$ 都有
\[
    1\leq a_n<b_n\leq2,
    \qquad a_n^2\leq2\leq b_n^2.
\]
同时，由式~\eqref{eq:sqrt-two-bisection-update}，这些非空闭区间逐个嵌套，且每次长度减半，故
\[
    b_n-a_n=\frac{1}{2^{n-1}}\longrightarrow0.
\]
由闭区间套定理~\ref{thm:nested-closed-intervals}，存在唯一公共点 $c\in[1,2]$，并且 $a_n\to c$、$b_n\to c$。利用极限的乘法法则，在 $a_n^2\leq2\leq b_n^2$ 中取极限，得到 $c^2\leq2\leq c^2$，所以 $c^2=2$，且 $c>0$。若另有正实数 $d$ 满足 $d^2=2$，则 $(c-d)(c+d)=c^2-d^2=0$；由于 $c+d>0$，必有 $c=d$。因此满足该方程的正实数存在且唯一，将它记为 $\sqrt2$。

由于 $c\in[a_n,b_n]$，用区间中点 $m_n$ 近似 $c$ 时，绝对误差不超过区间长度的一半，即
\begin{equation}
    \left|\sqrt2-m_n\right|
    =\left|c-\frac{a_n+b_n}{2}\right|
    \leq\frac{b_n-a_n}{2}
    =\frac1{2^n}.
    \label{eq:sqrt-two-bisection-error}
\end{equation}
这里 $I_n$ 由初始区间经过 $n-1$ 次二分得到。给定允许的绝对误差 $\varepsilon>0$，取满足 $2^{-n}\leq\varepsilon$ 的正整数 $n$，构造出 $I_n$ 后输出 $m_n$，便可保证误差不超过 $\varepsilon$。因此，闭区间套定理既保证了所构造实数的存在，也使二分过程获得了明确的近似误差估计。
\end{solve}

\begin{remark}
\label{rem:nested-interval-assumptions}
\emph[caution]{闭性保证极限位置不会被排除，长度趋于零保证公共点唯一。}例如，开区间 $(0,1/n)$ 逐个嵌套且长度趋于零，却没有公共点。若进一步要求开区间严格嵌套，即 $a_n<a_{n+1}<b_{n+1}<b_n$，并且 $b_n-a_n\to0$，则结论仍成立：对相应的闭区间套应用定理，所得公共点 $c$ 满足 $a_n<a_{n+1}\leq c\leq b_{n+1}<b_n$，因而属于每个开区间。

若保留闭区间的嵌套条件而不要求长度趋于零，则同样由单调收敛原理得到 $\alpha=\lim a_n=\sup\{a_n:n\geq1\}$、$\beta=\lim b_n=\inf\{b_n:n\geq1\}$，且
\begin{equation}
    \bigcap_{n=1}^{\infty}[a_n,b_n]=[\alpha,\beta].
    \label{eq:nested-interval-intersection}
\end{equation}
这是因为 $a_n\leq\alpha\leq\beta\leq b_n$，而任一公共点 $x$ 满足 $a_n\leq x\leq b_n$，取极限便得 $\alpha\leq x\leq\beta$。例如，$[0,1+1/n]$ 的交集为 $[0,1]$。另外，$[n,+\infty)$ 虽逐个嵌套且均为非空闭集，交集却为空；它们不是有界闭区间，不满足定理的假设。
\end{remark}

\begin{example}
% 教材例 2.4.1，第 58 页。
设 $\{[a_n,b_n]\}$ 为一列非空有界闭区间，且任意两个区间都有公共点。证明这些区间有公共点，即存在 $c\in\mathbb{R}$ 使 $c\in[a_n,b_n]$ 对每个 $n$ 成立。
\end{example}

\begin{solve}
由 $[a_i,b_i]\cap[a_j,b_j]\neq\varnothing$，可知 $a_i\leq b_j$ 对任意 $i,j$ 成立。因此每个 $b_j$ 都是集合 $\{a_n:n\geq1\}$ 的上界，令 $c=\sup\{a_n:n\geq1\}$，便有 $a_n\leq c\leq b_n$ 对每个 $n$ 成立。这里仅要求区间两两相交，并不要求它们构成区间套。
\end{solve}

\subsection{有限覆盖定理}
\label{subsec:finite-covering-theorem}

有限覆盖定理研究用许多开区间覆盖一个闭区间时，能否从中选出有限个仍覆盖它。它常用于把函数在各点附近的局部性质推广为整个闭区间上的性质。

\begin{definition}
\label{def:real-neighborhood}
设 $x\in\mathbb{R}$，$\delta>0$。开区间
\[
    U(x,\delta)=(x-\delta,x+\delta)
\]
称为 $x$ 的一个\emph{邻域}（neighborhood），相应的\emph{去心邻域}（punctured neighborhood）记为
\[
    U_0(x,\delta)=U(x,\delta)\setminus\{x\}
    =\{y\in\mathbb{R}:0<|y-x|<\delta\}.
\]
\end{definition}

\begin{definition}
\label{def:open-cover-and-finite-subcover}
设 $A\subseteq\mathbb{R}$，$\{E_\lambda\}_{\lambda\in\Lambda}$ 为一族实数集，其中指标集 $\Lambda$ 不要求有限或可数。若
\[
    A\subseteq\bigcup_{\lambda\in\Lambda}E_\lambda,
\]
则称这族集合为 $A$ 的一个\emph{覆盖}（cover）。若其中每个 $E_\lambda$ 都是开区间，则称它为 $A$ 的一个\emph{开覆盖}（open cover）；若指标集 $\Lambda$ 只有有限多个元素，则称它为一个\emph{有限覆盖}（finite cover）。这里采用教材本节的表述，以开区间作为开覆盖的成员。

若能从原覆盖中选出有限个集合 $E_{\lambda_1},\ldots,E_{\lambda_N}$，仍满足
\[
    A\subseteq E_{\lambda_1}\cup\cdots\cup E_{\lambda_N},
\]
则称它们组成一个\emph{有限子覆盖}（finite subcover）。
\end{definition}

\begin{theorem}[有限覆盖定理]
\label{thm:finite-covering}
设 $a,b\in\mathbb{R}$ 且 $a\leq b$，则闭区间 $[a,b]$ 的任意开覆盖都有有限子覆盖。这是\emph{海涅--博雷尔定理}（Heine--Borel theorem）的闭区间形式。
\end{theorem}

\begin{proof}
设给定的开覆盖为 $\{E_\lambda\}_{\lambda\in\Lambda}$。当 $a=b$ 时，选取其中任意一个包含 $a$ 的开区间即可。以下设 $a<b$，并反设不能从该覆盖中选出 $[a,b]$ 的有限子覆盖。\emph[tip]{将区间二等分，始终保留不能被原覆盖中的有限个开区间覆盖的那一半。}因为若左右两个闭半区间各有有限子覆盖，将二者合并便得到整个区间的有限子覆盖，所以至少一个半区间仍不能被有限覆盖。

将第一次选出的半区间记为 $[a_1,b_1]$，反复二等分并作同样选择，得到闭区间套 $\{[a_n,b_n]\}$，其中每个区间都不能被原覆盖中的有限个开区间覆盖，且
\[
    b_n-a_n=\frac{b-a}{2^n}\longrightarrow0.
\]
由定理~\ref{thm:nested-closed-intervals}，存在唯一公共点 $\xi$，并且 $a_n\to\xi$、$b_n\to\xi$。由于 $\xi\in[a,b]$，原覆盖中存在开区间 $E_{\lambda_0}=(c,d)$ 包含 $\xi$，即 $c<\xi<d$。由端点序列的收敛性，充分大的 $n$ 满足
\[
    c<a_n\leq\xi\leq b_n<d,
\]
于是 $[a_n,b_n]\subseteq E_{\lambda_0}$。这说明只用原覆盖中的一个开区间就能覆盖 $[a_n,b_n]$，与其选择方式矛盾。
\end{proof}

\begin{remark}
\emph[caution]{有限覆盖定理要求从给定的任意开覆盖中选取有限子覆盖。}如果只要求另找有限个开区间覆盖 $[a,b]$，那么一个开区间 $(a-1,b+1)$ 就已经足够，这并不是定理要说明的内容。

开区间 $(0,1)$ 有开覆盖 $\{(1/n,1):n\geq2\}$。任取有限个这样的区间，记其下标的最大值为 $N$，则它们的并集为 $(1/N,1)$，不能覆盖 $(0,1)$。无界闭集 $\mathbb{R}$ 有开覆盖 $\{(-n,n):n\geq1\}$，任意有限个成员的并集都是有界的，也不能覆盖 $\mathbb{R}$。这说明若去掉被覆盖集合的闭性或有界性，结论未必成立。

另一方面，单点集族 $\{\{x\}:x\in[0,1]\}$ 覆盖闭区间 $[0,1]$，但任取有限个成员只能覆盖有限个点，所以没有有限子覆盖。这些单点集不是开区间，说明不能将定理中的开覆盖任意换成一般的覆盖。
\end{remark}

\begin{example}
% 教材例 2.4.2，第 60--61 页。
设实函数 $f$ 定义在 $[a,b]$ 上，且对任意 $x\in[a,b]$，都存在 $\delta(x)>0$，使 $f$ 在 $U(x,\delta(x))\cap[a,b]$ 上有界。证明 $f$ 在整个 $[a,b]$ 上有界。
\end{example}

\begin{solve}
对每个 $x\in[a,b]$，取上述邻域及实数 $M_x>0$，使得当 $y\in U(x,\delta(x))\cap[a,b]$ 时，$|f(y)|\leq M_x$。这些邻域组成 $[a,b]$ 的一个开覆盖，故由定理~\ref{thm:finite-covering}，可选出有限个邻域 $U(x_1,\delta(x_1)),\ldots,U(x_N,\delta(x_N))$，仍覆盖 $[a,b]$。令
\[
    M=\max\{M_{x_1},\ldots,M_{x_N}\}.
\]
任取 $y\in[a,b]$，总有某个 $k$ 使 $y\in U(x_k,\delta(x_k))$，从而 $|f(y)|\leq M_{x_k}\leq M$。因此 $f$ 在 $[a,b]$ 上有统一的界。有限子覆盖的作用在于把随点变化的局部界归结为有限多个数，再取它们的最大值。
\end{solve}

\begin{example}
\label{ex:continuous-function-bounded-by-finite-cover}
设 $a,b\in\mathbb{R}$ 且 $a<b$，实函数 $f:[a,b]\to\mathbb{R}$ 在 $[a,b]$ 上连续。利用有限覆盖定理~\ref{thm:finite-covering} 证明 $f$ 在 $[a,b]$ 上有界，即存在 $M>0$，使得对任意 $x\in[a,b]$ 都有 $|f(x)|\leq M$。
\end{example}

\subsection{聚点原理}
\label{subsec:accumulation-point-principle}

本小节先讨论实数集的聚点，再通过聚点原理研究序列的子列。集合的聚点描述的是集合中的点在某个位置附近不断聚集；子列则保留原序列中所选各项的先后顺序，二者需要加以区分。

\begin{definition}
\label{def:accumulation-point-of-real-set}
设 $E\subseteq\mathbb{R}$。若 $x_0\in\mathbb{R}$ 满足
\[
    \forall\delta>0,\quad U_0(x_0,\delta)\cap E\neq\varnothing,
\]
则称 $x_0$ 为 $E$ 的一个\emph{聚点}（accumulation point）。等价地，对每个 $\delta>0$，存在 $x\in E$ 使 $0<|x-x_0|<\delta$。这里 $x_0$ 不一定属于 $E$。
\end{definition}

教材给出了聚点的另外两个等价刻画：$x_0$ 的每个邻域都含有 $E$ 中无穷多个不同的点；存在由 $E$ 中互不相同的点组成的序列 $\{x_n\}$，使 $x_n\to x_0$。为说明等价性，先设 $x_0$ 为聚点。若某个邻域只含有限个异于 $x_0$ 的点，便可缩小半径以排除这些点，与定义矛盾，故每个邻域都含有无穷多个点。进而可逐次选取 $x_n\in E\cap U(x_0,1/n)$，并排除 $x_0$ 和此前选过的有限个点，使各 $x_n$ 互不相同；由 $|x_n-x_0|<1/n$ 得 $x_n\to x_0$。反过来，若互不相同的点列 $\{x_n\}\subseteq E$ 趋于 $x_0$，则任一邻域内都有该序列充分靠后的项，其中至多一项等于 $x_0$，故该去心邻域与 $E$ 相交，$x_0$ 是聚点。

若 $x_0\in E$，但 $x_0$ 不是 $E$ 的聚点，则称它为 $E$ 的\emph{孤立点}（isolated point）。由定义，此时存在 $\delta>0$，使 $U(x_0,\delta)\cap E=\{x_0\}$。例如，$E=\{1/n:n\geq1\}$ 的唯一聚点是 $0$，且 $0\notin E$，而 $E$ 的每个点都是孤立点。$0$ 是聚点可由 $1/n\to0$ 得到；对任意 $x_0\neq0$，先取不含 $0$ 的充分小邻域，其中仅有有限个 $1/n$，再缩小邻域便能排除异于 $x_0$ 的这些点，故 $x_0$ 不是聚点。

\begin{theorem}[聚点原理]
\label{thm:accumulation-point-principle}
实数集中的任何有界无限子集至少有一个聚点。
\end{theorem}

\begin{proof}
设 $E\subseteq\mathbb{R}$ 有界且含有无穷多个不同的点，取闭区间 $[a,b]$ 使 $E\subseteq[a,b]$。反设 $E$ 没有聚点，则对每个 $x\in[a,b]$，存在 $\delta_x>0$，使
\[
    U_0(x,\delta_x)\cap E=\varnothing.
\]
因此 $U(x,\delta_x)$ 中至多有 $E$ 的一个点，即可能存在的中心点 $x$。这些邻域组成 $[a,b]$ 的一个开覆盖，由定理~\ref{thm:finite-covering}，可选取有限个邻域，使
\[
    [a,b]\subseteq\bigcup_{j=1}^{m}U(x_j,\delta_{x_j}).
\]
由于每个邻域至多含有 $E$ 的一个点，而它们覆盖了 $E$，故 $E$ 至多有 $m$ 个点，与 $E$ 为无限集矛盾。因此 $E$ 必有聚点。
\end{proof}

\begin{definition}
\label{def:subsequence-and-sequence-cluster-point}
设 $\{x_n\}$ 为实数列，$n_1<n_2<\cdots$ 为严格递增的正整数列，则 $\{x_{n_k}\}$ 称为 $\{x_n\}$ 的一个\emph{子列}（subsequence）。若某个子列满足 $x_{n_k}\to\xi$，则称 $\xi$ 为原序列的一个\emph{聚点}（cluster point），也称\emph{子列极限}（subsequential limit）。这里 $n_k\geq k$，因而子列的下标必趋于无穷。
\end{definition}

\begin{theorem}[收敛序列的子列]
\label{thm:subsequence-of-convergent-sequence}
若 $x_n\to a$，则 $\{x_n\}$ 的任意子列 $\{x_{n_k}\}$ 都满足 $x_{n_k}\to a$。
\end{theorem}

\begin{proof}
对任意 $\varepsilon>0$，存在正整数 $N$，使 $n>N$ 时 $|x_n-a|<\varepsilon$。当 $k>N$ 时，由 $n_k\geq k>N$，得到 $|x_{n_k}-a|<\varepsilon$，所以 $x_{n_k}\to a$。
\end{proof}

这一结论给出常用的发散判据：若一个序列存在两个收敛子列，且它们的极限不同，则原序列必发散。

\begin{theorem}[波尔查诺--魏尔斯特拉斯定理]
\label{thm:bolzano-weierstrass-sequence}
每个有界实数列都至少有一个收敛子列。这一结论称为\emph{波尔查诺--魏尔斯特拉斯定理}（Bolzano--Weierstrass theorem），是聚点原理在序列问题中的应用。
\end{theorem}

\begin{proof}
设 $\{x_n\}$ 有界。若它只取有限多个不同的值，则其中至少有一个值出现无穷多次；否则每个值都只出现有限次，整个序列便只有有限项，产生矛盾。选出等于该值的各项，即得一个收敛的常值子列。

若它取无穷多个不同的值，则取值集合 $E=\{x_n:n\geq1\}$ 为有界无限实数集。由定理~\ref{thm:accumulation-point-principle}，$E$ 至少有一个聚点，记为 $a$。对每个正整数 $k$，邻域 $U(a,1/k)$ 内都有 $E$ 的无穷多个不同的点，因此也有原序列的无穷多项。先取 $x_{n_1}\in U(a,1)$；选定 $n_{k-1}$ 后，再取 $n_k>n_{k-1}$，使 $x_{n_k}\in U(a,1/k)$。这样便得到子列 $\{x_{n_k}\}$，且
\[
    |x_{n_k}-a|<\frac1k\longrightarrow0.
\]
所以 $x_{n_k}\to a$，定理得证。
\end{proof}

\begin{notealert}[title={波尔查诺--魏尔斯特拉斯定理的直观理解}]
定理~\ref{thm:bolzano-weierstrass-sequence} 说明，一串数即使一直在不同位置之间跳动，只要始终被限制在某个有限区间内，就一定能按原来的先后顺序挑出无穷多项，使这些项趋于某个固定的实数。这里，\emph{有界只限制取值范围}，并不要求数列逐渐稳定；选取子列时可以跳过一些项，但必须保留所选各项在原数列中的先后顺序。

例如，数列 $x_n=(-1)^n$ 的各项依次为
\[
    -1,\ 1,\ -1,\ 1,\ -1,\ 1,\ \ldots,
\]
始终位于 $[-1,1]$ 内。挑出偶数位置的项，得到常值子列 $x_{2k}=1$，它收敛于 $1$；挑出奇数位置的项，得到常值子列 $x_{2k-1}=-1$，它收敛于 $-1$。两个子列的极限不同，由定理~\ref{thm:subsequence-of-convergent-sequence} 可知原数列不收敛。

因此，该定理保证的是\emph{至少能选出一个收敛子列}，并不保证原数列收敛，也不保证所有收敛子列具有相同的极限。
\end{notealert}

\begin{remark}
\label{rem:set-and-sequence-cluster-points}
\emph[caution]{序列的聚点不能简单地等同于其取值集合的聚点。}例如，常值序列 $x_n=c$ 的每个子列都趋于 $c$，但取值集合 $\{c\}$ 没有聚点。这也是上述证明先区分取值集合有限与无限两种情况的原因；若取值集合有限，不能直接对它应用集合的聚点原理。
\end{remark}

\subsection{柯西收敛准则}
\label{subsec:cauchy-convergence-criterion}

用极限定义证明 $x_n\to a$，需要确定候选极限 $a$。柯西收敛准则改为比较序列后面的项彼此之间的距离，因此可以在尚不知道极限值时判断收敛性。

\begin{definition}
\label{def:cauchy-sequence}
若实数列 $\{x_n\}$ 满足
\begin{equation}
    \forall\varepsilon>0,\quad\exists N\geq1,\quad
    \forall m,n>N,\quad |x_m-x_n|<\varepsilon,
    \label{eq:cauchy-sequence-condition}
\end{equation}
其中 $N,m,n$ 均为正整数，则称它为\emph{柯西序列}（Cauchy sequence），简称柯西列，也称基本列。等价地，对每个 $\varepsilon>0$，存在正整数 $N$，使得对所有 $n>N$ 及所有正整数 $p$ 都有 $|x_{n+p}-x_n|<\varepsilon$。
\end{definition}

条件~\eqref{eq:cauchy-sequence-condition} 要求足够靠后的\emph{任意两项}都彼此接近。$N$ 可以依赖于 $\varepsilon$，但一旦选定，就必须同时适用于所有 $m,n>N$；在等价写法中，$N$ 也不能随 $p$ 改变。

\begin{theorem}[柯西收敛准则]
\label{thm:cauchy-convergence-criterion}
实数列收敛的充分必要条件是它为柯西列。
\end{theorem}

\begin{proof}
先证必要性。若 $x_n\to a\in\mathbb{R}$，则对任意 $\varepsilon>0$，存在正整数 $N$，使得 $n>N$ 时 $|x_n-a|<\varepsilon/2$。于是对任意 $m,n>N$，由三角不等式得
\[
    |x_m-x_n|\leq|x_m-a|+|x_n-a|<\varepsilon,
\]
故 $\{x_n\}$ 是柯西列。

再证充分性。\emph[tip]{先证明柯西列有界，再取收敛子列，最后利用柯西条件把子列的收敛推广到整个序列。}在式~\eqref{eq:cauchy-sequence-condition} 中取 $\varepsilon=1$，得到正整数 $N_0$，使 $n>N_0$ 时
\[
    |x_n-x_{N_0+1}|<1,
    \qquad |x_n|<1+|x_{N_0+1}|.
\]
令 $M=\max\{|x_1|,\ldots,|x_{N_0}|,1+|x_{N_0+1}|\}$，便有 $|x_n|\leq M$ 对所有 $n$ 成立，所以该序列有界。由定理~\ref{thm:bolzano-weierstrass-sequence}，存在子列 $x_{n_k}\to a\in\mathbb{R}$。

给定 $\varepsilon>0$，由子列收敛，存在正整数 $K$，使 $k>K$ 时 $|x_{n_k}-a|<\varepsilon/2$；由柯西条件，存在正整数 $N$，使 $m,n>N$ 时 $|x_m-x_n|<\varepsilon/2$。固定一个 $k>K$ 且 $n_k>N$，则对所有 $n>N$，都有
\[
    |x_n-a|\leq|x_n-x_{n_k}|+|x_{n_k}-a|<\varepsilon.
\]
因此 $x_n\to a$。
\end{proof}

\begin{example}
% 教材例 2.4.8，第 66--67 页；同时说明相邻项之差趋于零并不足以保证收敛。
设 $s_n=\sum_{k=1}^{n}1/k$。证明相邻项之差 $s_{n+1}-s_n\to0$，但 $\{s_n\}$ 不收敛。
\end{example}

\begin{solve}
由定义，$s_{n+1}-s_n=1/(n+1)\to0$。然而对任意正整数 $n$，都有
\[
    s_{2n}-s_n=\sum_{k=n+1}^{2n}\frac1k
    \geq n\cdot\frac1{2n}=\frac12.
\]
取 $\varepsilon_0=1/2$，无论给定多大的正整数 $N$，总可选取 $n>N$，使下标 $n,2n$ 都大于 $N$，而 $|s_{2n}-s_n|\geq\varepsilon_0$。因此它不满足柯西条件，由定理~\ref{thm:cauchy-convergence-criterion} 可知该序列不收敛。\emph[caution]{相邻两项的距离趋于零，不能替代尾部任意两项的距离一致变小。}
\end{solve}

\begin{example}
% 教材例 2.4.9，第 67 页。
证明序列 $y_n=\sum_{k=1}^{n}1/k^2$ 收敛。
\end{example}

\begin{solve}
对正整数 $n>m$，由 $1/k^2<1/((k-1)k)$，有
\[
    |y_n-y_m|=\sum_{k=m+1}^{n}\frac1{k^2}
    <\sum_{k=m+1}^{n}\left(\frac1{k-1}-\frac1k\right)
    =\frac1m-\frac1n<\frac1m.
\]
给定 $\varepsilon>0$，取正整数 $N>1/\varepsilon$。当 $n,m>N$ 且 $n>m$ 时，上式给出 $|y_n-y_m|<1/m<\varepsilon$；$m>n$ 时交换两者即可，$m=n$ 时差为零。因此 $\{y_n\}$ 是柯西序列，由定理~\ref{thm:cauchy-convergence-criterion} 可知它收敛。这个证明不需要预先确定极限值。
\end{solve}

\begin{remark}
\label{rem:completeness-of-real-numbers}
设 $K\subseteq\mathbb{R}$，若由 $K$ 中的点组成的任意柯西序列都收敛到 $K$ 中某一点，则称 $K$ 是\emph{完备的}（complete），这称为 $K$ 的\emph{完备性}（completeness）。柯西收敛准则说明 $\mathbb{R}$ 是完备的。$\mathbb{Q}$ 则不完备：设 $q_n$ 为 $\sqrt2$ 保留小数点后 $n$ 位得到的截断小数，由 $0\leq\sqrt2-q_n<10^{-n}$ 可知 $q_n\to\sqrt2$，所以 $\{q_n\}$ 是有理数柯西序列，却没有有理数极限。

任意非空有界开区间 $(a,b)$ 也不完备。例如，$x_n=b-(b-a)/(n+1)$ 的每一项都在 $(a,b)$ 中，并且 $x_n\to b$，所以它是柯西序列，但极限 $b$ 不属于 $(a,b)$。因此，讨论完备性时必须明确要求极限存在于哪个集合中。
\end{remark}

\section{序列的上下极限}
\label{sec:upper-and-lower-sequence-limits}

% 根据《数学分析》第一册第 2.5 节整理，教材第 68--76 页；分章 PDF 第 40--48 页。
普通极限要求序列的所有充分靠后的项都趋近同一个数，但有界序列也可能不断振荡。上一节的波尔查诺--魏尔斯特拉斯定理保证有界序列至少有一个收敛子列，却不保证这些子列具有相同的极限。为了描述序列在尾部的取值范围，本节引入上下极限：先求每一段尾部的上、下确界，再考察这些确界随起始位置后移而趋向何处。以下 $\{x_n\}$ 均为实数列，下标取正整数。

\subsection{从尾部确界定义上下极限}
\label{subsec:liminf-limsup-definitions}

先设 $\{x_n\}$ 有界，记它从第 $n$ 项开始的取值集合为 $E_n=\{x_k:k\geq n\}$，并令
\begin{equation}
    l_n=\inf E_n=\inf_{k\geq n}x_k,
    \qquad h_n=\sup E_n=\sup_{k\geq n}x_k.
    \label{eq:sequence-tail-bounds}
\end{equation}
由于 $E_{n+1}\subseteq E_n$，缩小取值集合不会使下确界减小，也不会使上确界增大，因此
\[
    l_n\leq l_{n+1}\leq h_{n+1}\leq h_n,
    \qquad l_n\leq x_n\leq h_n.
\]
于是 $\{l_n\}$ 单调增加，$\{h_n\}$ 单调减少，并且两者都有界。由单调收敛原理，它们分别收敛于 $\sup_{n\geq1}l_n$ 和 $\inf_{n\geq1}h_n$。

\begin{definition}
\label{def:liminf-limsup}
设 $\{x_n\}$ 是有界实数列，$l_n,h_n$ 由式~\eqref{eq:sequence-tail-bounds} 定义。称
\begin{equation}
    \begin{aligned}
        \liminf_{n\to\infty}x_n
        &=\lim_{n\to\infty}l_n
        =\sup_{n\geq1}\inf_{k\geq n}x_k,\\
        \limsup_{n\to\infty}x_n
        &=\lim_{n\to\infty}h_n
        =\inf_{n\geq1}\sup_{k\geq n}x_k
    \end{aligned}
    \label{eq:liminf-limsup-tail-formulas}
\end{equation}
分别为序列的\emph{下极限}（limit inferior）和\emph{上极限}（limit superior），也分别记为 $\underline{\lim}_{n\to\infty}x_n$ 和 $\overline{\lim}_{n\to\infty}x_n$。
\end{definition}

上述定义可以推广到无界实数列。记 $\overline{\mathbb{R}}=\mathbb{R}\cup\{-\infty,+\infty\}$ 为\emph{扩展实数系}（extended real number system），约定非空无上界实数集的上确界为 $+\infty$，非空无下界实数集的下确界为 $-\infty$。在这一约定下，式~\eqref{eq:sequence-tail-bounds} 对任意实数列都有意义，而式~\eqref{eq:liminf-limsup-tail-formulas} 中的极限按广义极限理解。具体地，若原序列无上界，删去有限项后仍无上界，故所有 $h_n=+\infty$，上极限为 $+\infty$；若原序列无下界，则所有 $l_n=-\infty$，下极限为 $-\infty$。若原序列有下界，则 $l_n$ 为单调增加的实数列，其极限可能有限，也可能为 $+\infty$；有上界时，$h_n$ 的极限同理可能有限，也可能为 $-\infty$。因此，\emph{每个实数列的上下极限都在扩展实数系中存在}，并且恒有
\begin{equation}
    \liminf_{n\to\infty}x_n\leq\limsup_{n\to\infty}x_n.
    \label{eq:liminf-le-limsup}
\end{equation}

\begin{notealert}[title={上下极限描述尾部的边界}]
下极限和上极限考察的是不断舍去前面有限项之后，剩余各项的下、上确界最终趋向哪里，而不是原序列的最小值和最大值。例如，$x_n=1/n$ 的上确界为 $1$，但其尾部上确界 $h_n=1/n\to0$，所以上极限为 $0$。改变、增加或删去开头有限项都不会改变上下极限；特别地，对固定非负整数 $r$，$\{x_{n+r}\}$ 与 $\{x_n\}$ 具有相同的上下极限，因为前者的尾部确界就是后者的 $l_{n+r}$ 与 $h_{n+r}$。
\end{notealert}

\begin{example}
% 参照教材例 2.5.1，第 69--70 页。
求下列四个序列的上下极限：
\[
    \sin\frac{n\pi}{2},\qquad n^{(-1)^n},\qquad (-1)^n n,\qquad n.
\]
\end{example}

\begin{solve}
对于 $\sin(n\pi/2)$，每个尾部都含有 $1$ 和 $-1$，且所有项均在 $[-1,1]$ 内，故 $l_n=-1$、$h_n=1$。对于 $n^{(-1)^n}$，偶数位置的项等于 $n$，奇数位置的项等于 $1/n$；每个尾部都无上界，同时含有任意接近 $0$ 的正数，故 $l_n=0$、$h_n=+\infty$。因此
\[
    \begin{aligned}
        \liminf_{n\to\infty}\sin\frac{n\pi}{2}&=-1,
        &\limsup_{n\to\infty}\sin\frac{n\pi}{2}&=1,\\
        \liminf_{n\to\infty}n^{(-1)^n}&=0,
        &\limsup_{n\to\infty}n^{(-1)^n}&=+\infty.
    \end{aligned}
\]
对于 $(-1)^n n$，每个尾部既无上界又无下界，故 $l_n=-\infty$、$h_n=+\infty$。对于 $x_n=n$，有 $l_n=n$、$h_n=+\infty$。于是
\[
    \begin{aligned}
        \liminf_{n\to\infty}(-1)^n n&=-\infty,
        &\limsup_{n\to\infty}(-1)^n n&=+\infty,\\
        \liminf_{n\to\infty}n&=+\infty,
        &\limsup_{n\to\infty}n&=+\infty.
    \end{aligned}
\]
最后一个例子说明，上下极限相等时，其共同值仍可能是无穷大；这时原序列具有广义极限，但不在通常的有限极限意义下收敛。
\end{solve}

\subsection{等价刻画、子列与收敛判据}
\label{subsec:liminf-limsup-characterizations}

\begin{theorem}[上下极限的等价刻画]
\label{thm:liminf-limsup-characterizations}
设 $\{x_n\}$ 为有界实数列，$h\in\mathbb{R}$。则 $h=\limsup_{n\to\infty}x_n$ 当且仅当对任意 $\varepsilon>0$，以下两个条件同时成立：
\begin{equation}
    \begin{aligned}
        &\exists N\geq1,\quad\forall n\geq N,\quad x_n<h+\varepsilon;\\
        &\forall K\geq1,\quad\exists n\geq K,\quad x_n>h-\varepsilon.
    \end{aligned}
    \label{eq:limsup-epsilon-characterization}
\end{equation}
这里 $N,K,n$ 均为正整数。等价地，存在子列趋于 $h$，且任何收敛子列的极限都不超过 $h$，即上极限是所有子列极限中的最大值。

相应地，对实数 $l$，$l=\liminf_{n\to\infty}x_n$ 当且仅当对任意 $\varepsilon>0$，同时有
\begin{equation}
    \begin{aligned}
        &\exists N\geq1,\quad\forall n\geq N,\quad x_n>l-\varepsilon;\\
        &\forall K\geq1,\quad\exists n\geq K,\quad x_n<l+\varepsilon.
    \end{aligned}
    \label{eq:liminf-epsilon-characterization}
\end{equation}
等价地，下极限是所有子列极限中的最小值。
\end{theorem}

\begin{proof}
先证上极限的刻画。若 $h_n\to h$，则给定 $\varepsilon>0$，充分大的 $n$ 满足 $x_n\leq h_n<h+\varepsilon$。另一方面，对任意正整数 $K$，有 $h_K\geq h>h-\varepsilon$；由上确界的定义，必存在 $n\geq K$ 使 $x_n>h-\varepsilon$，否则 $h-\varepsilon$ 就是该尾部的上界，产生矛盾。因此式~\eqref{eq:limsup-epsilon-characterization} 成立。

反过来，若式~\eqref{eq:limsup-epsilon-characterization} 成立，则第一条给出充分大的 $n$ 满足 $h_n\leq h+\varepsilon$，第二条给出每个 $h_n\geq h-\varepsilon$。记 $H=\lim h_n$，取极限得到 $h-\varepsilon\leq H\leq h+\varepsilon$；由 $\varepsilon$ 的任意性，$H=h$。

下面证明子列极限的刻画。在式~\eqref{eq:limsup-epsilon-characterization} 中依次取 $\varepsilon=1/k$，由第一条选取正整数 $N_k$，使 $n\geq N_k$ 时 $x_n<h+1/k$。令 $n_0=0$，再由第二条递归选取 $n_k\geq\max\{N_k,n_{k-1}+1\}$，使 $x_{n_k}>h-1/k$，便得到下标严格递增的子列，且有
\[
    h-\frac1k<x_{n_k}<h+\frac1k.
\]
故 $x_{n_k}\to h$。若另一子列 $x_{m_k}\to b$，则第一条保证它充分靠后的项小于 $h+\varepsilon$，取极限得 $b\leq h+\varepsilon$，再由 $\varepsilon$ 的任意性得 $b\leq h$。这说明实际的上极限是子列极限中的最大值；反之，若某个数也具有这一最大值性质，它必等于实际的上极限。

由于 $\sup_{k\geq n}(-x_k)=-\inf_{k\geq n}x_k=-l_n$，对序列 $\{-x_n\}$ 应用已经证明的结论，便得到下极限的两个等价刻画。
\end{proof}

\emph[caution]{“充分靠后的所有项”与“无穷多项”不能互换。}式~\eqref{eq:limsup-epsilon-characterization} 的第一条说，最终所有项都低于 $h+\varepsilon$；第二条说，无论下标走到多远，后面仍有项高于 $h-\varepsilon$。它不要求最终所有项都高于 $h-\varepsilon$，因此有上极限并不意味着原序列收敛。

\begin{remark}
\label{rem:liminf-limsup-and-set-cluster-points}
按定义~\ref{def:subsequence-and-sequence-cluster-point}，对任意有界序列，上、下极限分别是它的最大、最小序列聚点。如果进一步要求各项互不相同，那么它们也是取值集合 $E=\{x_n:n\geq1\}$ 的最大、最小聚点：趋于上下极限的子列由互不相同的数组成，故其极限是 $E$ 的聚点；反过来，$E$ 的任一聚点的每个邻域都含有无穷多个不同的点，可以逐次选取下标递增的项趋于该点，所以它也是子列极限。教材对取值集合聚点所作的互异性假设不能直接省略，常值序列便是反例，参见备注~\ref{rem:set-and-sequence-cluster-points}。
\end{remark}

\begin{theorem}[子列的上下极限与收敛判据]
\label{thm:subsequence-limits-and-convergence}
设 $\{x_n\}$ 为任意实数列，$\{x_{n_j}\}$ 为其任一子列，则在扩展实数系中有
\begin{equation}
    \liminf_{n\to\infty}x_n
    \leq\liminf_{j\to\infty}x_{n_j}
    \leq\limsup_{j\to\infty}x_{n_j}
    \leq\limsup_{n\to\infty}x_n.
    \label{eq:subsequence-liminf-limsup-bounds}
\end{equation}
此外，对任意 $a\in\overline{\mathbb{R}}$，都有
\begin{equation}
    \lim_{n\to\infty}x_n=a
    \quad\Longleftrightarrow\quad
    \liminf_{n\to\infty}x_n=\limsup_{n\to\infty}x_n=a.
    \label{eq:convergence-via-liminf-limsup}
\end{equation}
当 $a=\pm\infty$ 时，左端按广义极限理解。
\end{theorem}

\begin{proof}
对子列从第 $j$ 项开始的尾部，记
\[
    \widetilde l_j=\inf_{i\geq j}x_{n_i},
    \qquad \widetilde h_j=\sup_{i\geq j}x_{n_i}.
\]
该尾部取值集合包含于原序列从第 $n_j$ 项开始的尾部，所以
\[
    l_{n_j}\leq\widetilde l_j\leq\widetilde h_j\leq h_{n_j}.
\]
因 $n_j\to\infty$，单调尾部确界序列的这两个子列仍趋于原来的广义极限。令 $j\to\infty$，即得式~\eqref{eq:subsequence-liminf-limsup-bounds}。

对于收敛判据，先设 $a\in\mathbb{R}$。若上下极限都等于 $a$，则由 $l_n\leq x_n\leq h_n$ 和夹逼收敛原理得到 $x_n\to a$。反之，若 $x_n\to a$，则对任意 $\varepsilon>0$，充分大的 $n$ 及所有 $k\geq n$ 都满足 $a-\varepsilon<x_k<a+\varepsilon$，从而 $a-\varepsilon\leq l_n\leq h_n\leq a+\varepsilon$，这说明 $l_n,h_n$ 都趋于 $a$。

若 $a=+\infty$，则 $x_n\to+\infty$ 意味着对每个实数 $M$，最终所有项都大于 $M$，因而最终 $l_n\geq M$，所以 $l_n\to+\infty$；由 $h_n\geq l_n$，也有 $h_n\to+\infty$。反之，由 $x_n\geq l_n\to+\infty$ 得 $x_n\to+\infty$。当 $a=-\infty$ 时，利用尾部上确界 $h_n$ 作对称论证即可。
\end{proof}

\begin{example}
% 教材例 2.5.2，第 74 页。
将 $[0,1]\cap\mathbb{Q}$ 中的所有有理数不重复地排成序列 $\{x_n\}$，求它的上下极限。
\end{example}

\begin{solve}
任意非空开区间 $(a,b)\subseteq(0,1)$ 都含有无穷多个有理数：由有理数的稠密性，其中有一个有理数；若其中只有有限个有理数，则可在去掉这些点后选出一个非空开子区间，与稠密性矛盾。因此，删去开头有限项后，剩余项仍在 $[0,1]$ 中稠密，尤其仍有任意接近 $0$ 和 $1$ 的项。于是对每个 $n$，都有 $l_n=0$、$h_n=1$，故
\[
    \liminf_{n\to\infty}x_n=0,
    \qquad \limsup_{n\to\infty}x_n=1.
\]
也可从聚点看：上述性质保证 $[0,1]$ 中每一点都是取值集合的聚点，区间外的点则不是；其最小、最大聚点分别为 $0$ 和 $1$，由备注~\ref{rem:liminf-limsup-and-set-cluster-points} 得到同一结论。由上下极限不相等可知，该序列不收敛。
\end{solve}

\subsection{上下极限的保序性与运算性质}
\label{subsec:liminf-limsup-operations}

\begin{theorem}
\label{thm:liminf-limsup-operations}
设 $\{x_n\}$、$\{y_n\}$ 为有界实数列，简记
\[
    \underline x=\liminf_{n\to\infty}x_n,\quad
    \overline x=\limsup_{n\to\infty}x_n,\quad
    \underline y=\liminf_{n\to\infty}y_n,\quad
    \overline y=\limsup_{n\to\infty}y_n.
\]
若充分大的 $n$ 都满足 $x_n\leq y_n$，则 $\underline x\leq\underline y$ 且 $\overline x\leq\overline y$。取相反数时，上下极限交换位置并变号：
\begin{equation}
    \liminf_{n\to\infty}(-x_n)=-\overline x,
    \qquad \limsup_{n\to\infty}(-x_n)=-\underline x.
    \label{eq:liminf-limsup-negation}
\end{equation}
对于和，有
\begin{equation}
    \begin{aligned}
        \underline x+\underline y
        &\leq\liminf_{n\to\infty}(x_n+y_n)
        \leq\underline x+\overline y,\\
        \underline x+\overline y
        &\leq\limsup_{n\to\infty}(x_n+y_n)
        \leq\overline x+\overline y.
    \end{aligned}
    \label{eq:liminf-limsup-sum}
\end{equation}
交换 $x_n,y_n$ 后，同样得到含 $\overline x+\underline y$ 的两个中间界。若进一步有 $x_n\geq0$、$y_n\geq0$ 对所有 $n$ 成立，则
\begin{equation}
    \begin{aligned}
        \underline x\,\underline y
        &\leq\liminf_{n\to\infty}(x_ny_n)
        \leq\underline x\,\overline y,\\
        \underline x\,\overline y
        &\leq\limsup_{n\to\infty}(x_ny_n)
        \leq\overline x\,\overline y.
    \end{aligned}
    \label{eq:liminf-limsup-product}
\end{equation}
交换两列后，也得到含 $\overline x\,\underline y$ 的相应界。
\end{theorem}

\begin{proof}
保序性由充分靠后的每个尾部上的下确界、上确界比较得到；有限个起始项不影响结论。取相反数的公式则由 $\inf_{k\geq n}(-x_k)=-\sup_{k\geq n}x_k$ 及其对偶式取极限得到。

为证明和的不等式，分别用 $l_n^x,h_n^x$ 与 $l_n^y,h_n^y$ 表示两列的尾部确界。对所有 $k\geq n$，有
\[
    l_n^x+l_n^y\leq x_k+y_k\leq h_n^x+h_n^y.
\]
取尾部确界，再令 $n\to\infty$，得到式~\eqref{eq:liminf-limsup-sum} 最左端和最右端的界。为证明 $\liminf(x_n+y_n)\leq\underline x+\overline y$，由定理~\ref{thm:liminf-limsup-characterizations} 选取 $x_{n_j}\to\underline x$；由于 $\{y_{n_j}\}$ 有界，再由定理~\ref{thm:bolzano-weierstrass-sequence} 选取子列 $y_{n_{j_i}}\to b$。这时 $b\leq\overline y$，且
\[
    x_{n_{j_i}}+y_{n_{j_i}}\longrightarrow\underline x+b
    \leq\underline x+\overline y.
\]
由定理~\ref{thm:subsequence-limits-and-convergence}，原序列的下极限不超过该子列极限，所需不等式成立。再选取 $y_{m_j}\to\overline y$，并从有界的 $\{x_{m_j}\}$ 中选出趋于某个 $c\geq\underline x$ 的子列，便得到和的一个子列极限 $c+\overline y\geq\underline x+\overline y$，故 $\limsup(x_n+y_n)\geq\underline x+\overline y$。

当两列非负时，有 $0\leq l_n^x,l_n^y$，从而
\[
    l_n^x l_n^y\leq x_k y_k\leq h_n^x h_n^y
    \qquad(k\geq n),
\]
取尾部确界并令 $n\to\infty$，得到乘积的两个外侧界。沿用上一段选出的子列，分别得到乘积的子列极限 $\underline x\,b\leq\underline x\,\overline y$ 与 $c\,\overline y\geq\underline x\,\overline y$；这里乘法保序依赖于 $\underline x,\overline y\geq0$。再应用式~\eqref{eq:subsequence-liminf-limsup-bounds}，即得两个中间界。
\end{proof}

\emph[caution]{一般不能把上下极限直接拆成两个上下极限的和或积。}例如，取 $x_n=(-1)^n$、$y_n=-(-1)^n$，则两列的上极限都为 $1$，但 $x_n+y_n=0$ 的上极限为 $0$。再取非负序列 $u_n=1+(-1)^n$、$v_n=1-(-1)^n$，则两列的上极限都为 $2$，但 $u_nv_n=0$。不同序列接近各自上极限的项未必出现在相同下标处。以上运算定理针对有界序列，不能在无界情形下不加条件地使用含 $+\infty-\infty$ 或 $0\cdot\infty$ 的表达式。

\begin{corollary}[收敛项与常数因子]
\label{cor:liminf-limsup-affine-rules}
设 $a_n\to a\in\mathbb{R}$，$\{x_n\}$ 有界，则
\[
    \liminf_{n\to\infty}(a_n+x_n)=a+\liminf_{n\to\infty}x_n,
    \qquad
    \limsup_{n\to\infty}(a_n+x_n)=a+\limsup_{n\to\infty}x_n.
\]
若 $c>0$，则乘以 $c$ 时上下极限分别乘以 $c$；若 $c<0$，则
\[
    \liminf_{n\to\infty}(cx_n)=c\limsup_{n\to\infty}x_n,
    \qquad
    \limsup_{n\to\infty}(cx_n)=c\liminf_{n\to\infty}x_n.
\]
当 $c=0$ 时，两者均为 $0$。
\end{corollary}

\begin{proof}
由 $\liminf a_n=\limsup a_n=a$，在式~\eqref{eq:liminf-limsup-sum} 中交换两列后使用中间界，即得
\[
    a+\liminf x_n\leq\liminf(a_n+x_n)\leq a+\liminf x_n,
\]
上极限同理。对于 $c>0$，由 $\inf_{k\geq n}(cx_k)=c\inf_{k\geq n}x_k$ 及对应的上确界等式取极限即可；$c<0$ 时，再结合式~\eqref{eq:liminf-limsup-negation} 得到结论。
\end{proof}

\subsection{用上下极限研究序列的敛散性}
\label{subsec:liminf-limsup-applications}

\begin{example}
% 教材例 2.5.3，第 75--76 页。
设 $\{x_n\}$ 为实数列，$\alpha\geq2$ 为常数，令 $y_n=x_n+\alpha x_{n+1}$。证明：$\{y_n\}$ 收敛的充分必要条件是 $\{x_n\}$ 收敛，并求两者极限之间的关系。
\end{example}

\begin{solve}
若 $x_n\to a$，则由极限的运算法则，$y_n\to(1+\alpha)a$。反过来，设 $y_n\to b$。\emph[tip]{先证明原序列有界，再由关系式求出上下极限之间的约束。}由 $\{y_n\}$ 有界，可取 $M>0$，使 $|x_1|\leq M$ 且所有 $|y_n|\leq M$。若 $|x_n|\leq M$，则
\[
    |x_{n+1}|=\frac1\alpha|y_n-x_n|
    \leq\frac{|y_n|+|x_n|}{\alpha}
    \leq\frac{2M}{\alpha}\leq M.
\]
由数学归纳法，$|x_n|\leq M$ 对所有 $n$ 成立，故原序列有界。

记 $l=\liminf x_n$、$h=\limsup x_n$。因为移去首项不改变上下极限，在 $x_n=y_n-\alpha x_{n+1}$ 两边应用推论~\ref{cor:liminf-limsup-affine-rules}，得到
\[
    h=b-\alpha l,\qquad l=b-\alpha h.
\]
两式相减，得 $h-l=\alpha(h-l)$。由于 $\alpha\geq2$，故 $h=l$，再代回即得 $h=l=b/(1+\alpha)$。由定理~\ref{thm:subsequence-limits-and-convergence}，$x_n\to b/(1+\alpha)$。条件不能任意去掉：若 $\alpha=1$ 且 $x_n=(-1)^n$，则 $y_n=0$ 收敛，但 $\{x_n\}$ 不收敛。
\end{solve}

\begin{example}
% 参照教材例 2.5.4，第 76 页。
利用上下极限证明序列 $\{\sin n\}$ 发散，其中角度以弧度计。
\end{example}

\begin{solve}
目标是找到两组下标，使相应的正弦值分别保持在一个正数以上和一个负数以下。记 $\lfloor t\rfloor$ 为不超过 $t$ 的最大整数，对正整数 $k$ 取
\[
    n_k=\left\lfloor2k\pi+\frac\pi4\right\rfloor+1,
    \qquad
    m_k=\left\lfloor(2k+1)\pi+\frac\pi4\right\rfloor+1.
\]
两列下标都严格增加。利用 $t<\lfloor t\rfloor+1\leq t+1$ 以及 $1<\pi/2$，得到
\[
    2k\pi+\frac\pi4<n_k<2k\pi+\frac{3\pi}4,
    \qquad
    (2k+1)\pi+\frac\pi4<m_k<(2k+1)\pi+\frac{3\pi}4.
\]
因此 $\sin n_k>\sqrt2/2$，而 $\sin m_k<-\sqrt2/2$。由式~\eqref{eq:subsequence-liminf-limsup-bounds}，
\[
    \liminf_{n\to\infty}\sin n
    \leq\liminf_{k\to\infty}\sin m_k
    \leq-\frac{\sqrt2}{2}
    <\frac{\sqrt2}{2}
    \leq\limsup_{k\to\infty}\sin n_k
    \leq\limsup_{n\to\infty}\sin n.
\]
所以上下极限不同，原序列发散。这里仅需将上下极限分开，无须求出它们的精确值。
\end{solve}

\section*{习题二}
\addcontentsline{toc}{section}{习题二}
\fi

\selectchapter{3}
\ifcompilechapter
\chapter{函数的极限与连续性}
\label{chap:function-limits-and-continuity}

序列极限研究下标趋于无穷时各项的变化，函数极限则允许自变量从不同方向趋于有限点或无穷远。\emph{二者的共同核心是：通过限制自变量的变化范围，使函数值与目标值的误差任意小。}本章先建立函数极限及其判别方法，再用极限刻画连续性；随后借助实数系的连续性，把逐点的局部结论推广为闭区间上的整体结论，最后用无穷小量与无穷大量的阶描述不同变化速度。

\section{函数的极限}
\label{sec:function-limits}

\subsection{函数极限的定义}
\label{subsec:definition-of-function-limits}

函数在某点可能没有定义，但在该点附近仍可能表现出确定的变化趋势。例如，$x\sin(1/x)$ 在 $x=0$ 没有定义，而 $|x\sin(1/x)|\leq|x|$ 表明它的值随着 $x\to0$ 趋近于零。这里所要刻画的是附近的函数值，而不是中心点的取值。以下沿用\crosschapterref{2}{定义~\ref{def:real-neighborhood}}{实数邻域的定义} 中的邻域 $U(x_0,\delta)$ 与去心邻域 $U_0(x_0,\delta)$。

\begin{definition}
\label{def:finite-function-limit}
设实函数 $f$ 在某个 $U_0(x_0,\delta_0)$ 上有定义，其中 $\delta_0>0$。若存在实数 $A$，使得
\begin{equation}
    \forall\varepsilon>0,\quad\exists\delta\in(0,\delta_0],\quad
    0<|x-x_0|<\delta\ \Longrightarrow\ |f(x)-A|<\varepsilon,
    \label{eq:epsilon-delta-function-limit}
\end{equation}
则称 $x$ 趋于 $x_0$ 时 $f(x)$ 的\emph{极限}（limit）为 $A$，记为
\[
    \lim_{x\to x_0}f(x)=A,
    \qquad\text{或}\qquad f(x)\to A\quad(x\to x_0).
\]
没有特别说明时，“极限存在”指有限极限存在。
\end{definition}

式~\eqref{eq:epsilon-delta-function-limit} 中，$\varepsilon$ 是任意给定的函数值误差，$\delta$ 则是为达到这一误差而选取的自变量范围。$\delta$ 可以依赖 $\varepsilon$、$x_0$ 与函数，但不能在选定之后再依赖邻域中的具体点 $x$。\emph[caution]{条件 $0<|x-x_0|$ 排除了中心点，所以极限与 $f(x_0)$ 是否有定义、取何值无关。}从几何上看，图像在足够小的去心邻域内全部落入水平带 $A-\varepsilon<y<A+\varepsilon$，而不只是有若干点落入其中。

\begin{example}
\label{ex:epsilon-delta-function-limits}
用定义证明
\[
    \lim_{x\to1}\frac{\sqrt{x}-1}{x-1}=\frac12,
    \qquad
    \lim_{x\to0}x\sin\frac1x=0.
\]
\end{example}

\begin{solve}
第一式先把差化为含 $|x-1|$ 的表达式。当 $0<|x-1|<1$ 时，$x>0$，从而
\[
    \left|\frac{\sqrt{x}-1}{x-1}-\frac12\right|
    =\left|\frac1{\sqrt{x}+1}-\frac12\right|
    =\frac{|x-1|}{2(\sqrt{x}+1)^2}
    \leq\frac{|x-1|}{2}.
\]
给定 $\varepsilon>0$，取 $\delta=\min\{1,2\varepsilon\}$，便有所需误差小于 $\varepsilon$。第二式由
\[
    \left|x\sin\frac1x-0\right|\leq|x|
\]
可直接取 $\delta=\varepsilon$。两式都不需要在极限点补定义函数值。
\end{solve}

\emph[tip]{用定义证明极限时，通常先限制 $|x-x_0|<\delta_1$，控制其余因子的大小，再把 $|f(x)-A|$ 估计为 $C|x-x_0|$ 或另一容易控制的量，最后取若干限制的最小值。}例如，$|x-2|<1$ 时 $|x|<3$，故
\[
    |x^4-16|=|x-2|\,|x+2|\,(x^2+4)\leq65|x-2|,
\]
从而证明 $\lim_{x\to2}x^4=16$ 时可取 $\delta=\min\{1,\varepsilon/65\}$。

\subsection{函数极限的性质}
\label{subsec:properties-of-function-limits}

函数极限只由任意充分小的去心邻域内的取值决定：如果两个函数在某个去心邻域内相同，它们在中心点的极限情况就相同。与收敛序列整体有界不同，函数有有限极限只能推出相应的\emph{局部有界性}，不能推出整个定义域上的有界性。

\begin{theorem}[唯一性、局部有界性与保序性]
\label{thm:local-properties-of-function-limits}
设 $\lim_{x\to x_0}f(x)=A\in\mathbb{R}$，则该极限唯一，且 $f$ 在 $x_0$ 的某个去心邻域内有界。若另有 $\lim_{x\to x_0}g(x)=B\in\mathbb{R}$，则有以下结论：若某个去心邻域内恒有 $f(x)\leq g(x)$，则 $A\leq B$；若 $A<B$，则对每个 $C\in(A,B)$，存在去心邻域使其中恒有 $f(x)<C<g(x)$。特别地，$A>0$ 时充分靠近 $x_0$ 的函数值为正，$A<0$ 时为负。
\end{theorem}

\begin{proof}
若 $A,B$ 都是 $f$ 的极限，则对任意 $\varepsilon>0$，在足够小的去心邻域内同时有 $|f(x)-A|<\varepsilon/2$ 与 $|f(x)-B|<\varepsilon/2$，故 $|A-B|<\varepsilon$。由 $\varepsilon$ 的任意性得 $A=B$。在极限定义中取 $\varepsilon=1$，得到局部估计 $|f(x)|<|A|+1$。

若 $A<B$，取 $\varepsilon=\tfrac12\min\{C-A,B-C\}$，并使两个误差同时小于 $\varepsilon$，便有 $f(x)<C<g(x)$。如果 $f(x)\leq g(x)$ 在某个去心邻域内成立但 $A>B$，将这一结论用于 $g,f$，就会在更小的去心邻域内得到 $g(x)<f(x)$，矛盾。因此 $A\leq B$。
\end{proof}

\emph[caution]{严格不等号在取极限后可能变成等号。}例如，$x^2>0$ 对所有 $x\neq0$ 成立，但 $\lim_{x\to0}x^2=0$。此外，$1/x$ 在每个 $x_0>0$ 处都有有限极限，却在 $(0,+\infty)$ 上无界。

\begin{theorem}[四则运算与夹逼原理]
\label{thm:function-limit-arithmetic-and-squeeze}
设 $f(x)\to A$、$g(x)\to B$（$x\to x_0$），其中 $A,B\in\mathbb{R}$，则
\begin{equation}
    \begin{aligned}
        \lim_{x\to x_0}(f(x)\pm g(x))&=A\pm B,\\
        \lim_{x\to x_0}f(x)g(x)&=AB,\\
        \lim_{x\to x_0}\frac{f(x)}{g(x)}&=\frac AB\qquad(B\neq0).
    \end{aligned}
    \label{eq:function-limit-arithmetic}
\end{equation}
商的条件 $B\neq0$ 保证 $g(x)$ 在某个去心邻域内不为零。若 $f(x)\leq h(x)\leq g(x)$ 在某个去心邻域内成立，且 $A=B$，则 $h(x)\to A$。
\end{theorem}

\begin{proof}
和与差的结论由误差不超过 $|f(x)-A|+|g(x)-B|$ 得到。乘积利用 $f$ 的局部有界性 $|f(x)|\leq M$，将误差估计为
\[
    |f(x)g(x)-AB|\leq M|g(x)-B|+|B|\,|f(x)-A|,
\]
令右端两个误差分别足够小即可。若 $B\neq0$，先限制 $|g(x)-B|<|B|/2$，则 $|g(x)|>|B|/2$，因而
\[
    \left|\frac1{g(x)}-\frac1B\right|
    \leq\frac{2}{|B|^2}|g(x)-B|\to0.
\]
再用乘积法则得到商的结论。对于夹逼原理，给定 $\varepsilon>0$，使两个外侧函数同时满足 $A-\varepsilon<f(x)$ 与 $g(x)<A+\varepsilon$，便有 $|h(x)-A|<\varepsilon$。
\end{proof}

\begin{theorem}[复合函数的极限]
\label{thm:composition-of-function-limits}
设 $g(x)\to u_0$（$x\to x_0$），$f(u)\to A$（$u\to u_0$），且复合函数在 $x_0$ 的某个去心邻域内有定义。如果该邻域内还有 $g(x)\neq u_0$，则
\begin{equation}
    \lim_{x\to x_0}f(g(x))=A.
    \label{eq:composition-of-function-limits}
\end{equation}
若 $f$ 在 $u_0$ 处有定义且 $f(u_0)=A$，则可以取消 $g(x)\neq u_0$ 这一条件。
\end{theorem}

\begin{proof}
给定 $\varepsilon>0$，由外层函数的极限选取 $\eta>0$，使 $0<|u-u_0|<\eta$ 时 $|f(u)-A|<\varepsilon$。再由内层函数的极限选取 $\delta>0$，使 $0<|x-x_0|<\delta$ 时 $|g(x)-u_0|<\eta$。在 $g(x)\neq u_0$ 的条件下，可以将 $u=g(x)$ 代入前一估计。若 $f(u_0)=A$，则 $g(x)=u_0$ 时误差本来就是零，结论仍成立。
\end{proof}

\begin{example}
\label{ex:composition-limit-exception}
说明复合函数极限定理中排除 $g(x)=u_0$ 的条件不能无条件删去。
\end{example}

\begin{solve}
设
\[
    f(u)=\begin{cases}1,&u\neq0,\\0,&u=0,\end{cases}
    \qquad
    g(x)=\begin{cases}x,&x\in\mathbb{Q},\\0,&x\notin\mathbb{Q}.\end{cases}
\]
由 $|g(x)|\leq|x|$ 得 $g(x)\to0$，且 $\lim_{u\to0}f(u)=1$。然而在非零有理数点，$f(g(x))=1$；在无理数点，$f(g(x))=0$。有理数与无理数都在实数中稠密，所以复合函数在 $0$ 附近始终取这两个不同的值，其极限不存在。问题在于内层函数反复取到 $0$，而 $f(0)$ 与外层函数的极限不同。
\end{solve}

四则运算应在核对各部分极限及分母条件后使用。例如，对正整数 $n$，先在 $x\neq1$ 时作恒等变形
\[
    \frac{x^n-1}{x-1}=1+x+\cdots+x^{n-1},
\]
再取极限，得到 $\lim_{x\to1}(x^n-1)/(x-1)=n$；原式中分母趋于零，不能直接套用商的法则。

\subsection{函数极限概念的推广}
\label{subsec:generalizations-of-function-limits}

自变量除了从两侧趋于一个有限点，还可以只从一侧趋近，或趋于无穷远；函数值也可能不趋于有限数而在绝对值上无限增大。推广定义时，关键是明确自变量最终进入什么集合，以及函数值需要满足什么条件。

\begin{definition}
\label{def:one-sided-function-limits}
设 $f$ 在 $(x_0,x_0+\delta_0)$ 上有定义。若
\[
    \forall\varepsilon>0,\quad\exists\delta\in(0,\delta_0],\quad
    0<x-x_0<\delta\ \Longrightarrow\ |f(x)-A|<\varepsilon,
\]
则称 $A$ 为 $f$ 在 $x_0$ 的\emph{右极限}（right-hand limit），记为 $\lim_{x\to x_0+}f(x)=A$，也记为 $f(x_0+0)=A$。将条件换成 $0<x_0-x<\delta$，便得到\emph{左极限}（left-hand limit），记为 $\lim_{x\to x_0-}f(x)=A$ 或 $f(x_0-0)=A$。
\end{definition}

\begin{proposition}[双侧极限与单侧极限]
\label{prop:two-sided-and-one-sided-limits}
设 $f$ 在 $x_0$ 的某个去心邻域内有定义，则
\[
    \lim_{x\to x_0}f(x)=A
    \quad\Longleftrightarrow\quad
    f(x_0-0)=f(x_0+0)=A.
\]
\end{proposition}

\begin{proof}
双侧极限的误差估计分别限制在左右两侧便给出必要性。反之，给定 $\varepsilon>0$，分别选取左右两侧的半径 $\delta_-,\delta_+$，再取 $\delta=\min\{\delta_-,\delta_+\}$，便在整个去心邻域内得到同一误差估计。
\end{proof}

例如，取整函数 $\lfloor x\rfloor$ 在整数 $k$ 处的左右极限分别为 $k-1,k$，所以双侧极限不存在。

\begin{definition}
\label{def:function-limits-at-infinity}
设 $f$ 在某个右侧无界区间上有定义。若
\[
    \forall\varepsilon>0,\quad\exists X>0,\quad
    x>X\ \Longrightarrow\ |f(x)-A|<\varepsilon,
\]
则记 $\lim_{x\to+\infty}f(x)=A$。相应地，将 $x>X$ 换成 $x<-X$ 或 $|x|>X$，分别定义 $\lim_{x\to-\infty}f(x)=A$ 与 $\lim_{x\to\infty}f(x)=A$，并要求函数在对应的尾部有定义。这里不带正负号的 $x\to\infty$ 表示 $|x|\to+\infty$，同时考察正、负两个方向。
\end{definition}

因此，$\lim_{x\to\infty}f(x)=A$ 当且仅当 $x\to+\infty$ 与 $x\to-\infty$ 的两个极限均为 $A$。上述有限极限仍有唯一性、相应尾部上的有界性、保序性、四则运算与夹逼性质，只需在证明中把去心邻域换成对应的自变量范围。

\begin{definition}
\label{def:generalized-function-limits}
设 $f$ 在 $x_0$ 的某个去心邻域内有定义。若
\[
    \forall M>0,\quad\exists\delta>0,\quad
    0<|x-x_0|<\delta\ \Longrightarrow\ f(x)>M,
\]
则称 $f(x)$ 趋于正无穷，记为 $\lim_{x\to x_0}f(x)=+\infty$。将结论换成 $f(x)<-M$ 或 $|f(x)|>M$，分别定义 $f(x)\to-\infty$ 与 $f(x)\to\infty$。这些称为\emph{广义极限}（generalized limits），最后一种也称函数为该过程中的\emph{无穷大量}（infinitely large quantity）。自变量趋向单侧或无穷远时，按相同方式替换自变量条件。
\end{definition}

\emph[caution]{函数值无界不等于函数趋于无穷大。}无界只要求能找到任意大的绝对值，而趋于无穷大要求自变量充分接近目标以后，所有函数值的绝对值都超过给定的界。例如，$x^2\sin x$ 在 $x\to+\infty$ 时无界，但在 $x=n\pi$ 处恒为零，故不是无穷大量。另一个区别是：$1/x$ 在 $x\to0$ 时满足 $|1/x|\to+\infty$，却既不趋于 $+\infty$，也不趋于 $-\infty$；其左右广义极限分别为 $-\infty,+\infty$。无穷符号不是实数，$\infty-\infty$、$0\cdot\infty$、$0/0$、$\infty/\infty$ 等形式不能按实数四则运算处理。

\begin{example}
\label{ex:radical-limits-at-infinity}
设 $F(x)=\sqrt{4x^2+2x+3}$，求 $x\to\pm\infty$ 时 $F(x)/x$ 的极限，以及
\[
    \lim_{x\to+\infty}(F(x)-2x),
    \qquad \lim_{x\to-\infty}(F(x)+2x).
\]
\end{example}

\begin{solve}
由 $\sqrt{x^2}=|x|$，有
\[
    \frac{F(x)}x=\frac{|x|}x\sqrt{4+\frac2x+\frac3{x^2}},
\]
故正、负方向的极限分别为 $2,-2$。对两个差式作有理化，得到
\[
    F(x)-2x=\frac{2x+3}{F(x)+2x},
    \qquad
    F(x)+2x=\frac{2x+3}{F(x)-2x}.
\]
在对应方向上将分子、分母同除以 $x$，分别得到 $2/(2+2)=1/2$ 与 $2/(-2-2)=-1/2$。这里既要消去无穷大量相减造成的不定形式，也要保留负方向上 $|x|/x=-1$ 的符号。
\end{solve}

\subsection{序列极限与函数极限的关系}
\label{subsec:sequence-and-function-limits}

要严格说明 $f(x)$ 在 $x\to x_0$ 时不趋于一个给定的实数 $A$，应将式~\eqref{eq:epsilon-delta-function-limit} 的量词逐个否定：
\begin{equation}
    \exists\varepsilon_0>0,\quad\forall\delta>0,\quad
    \exists x\in U_0(x_0,\delta),\quad |f(x)-A|\geq\varepsilon_0,
    \label{eq:negation-of-function-limit}
\end{equation}
其中只需考虑足够小且位于定义范围内的邻域。这表示任意靠近 $x_0$ 都仍能找到误差较大的点，并不要求整个邻域内的点都具有较大误差；它也只否定以 $A$ 为极限，尚未排除其他候选极限。

\begin{theorem}[海涅定理]
\label{thm:heine-sequential-criterion}
设 $f$ 在 $U_0(x_0,\delta_0)$ 上有定义，$A\in\mathbb{R}$。则 $\lim_{x\to x_0}f(x)=A$ 的充分必要条件是：对每一个满足
\[
    x_n\in U_0(x_0,\delta_0),\qquad x_n\to x_0
\]
的序列，都有 $f(x_n)\to A$。这一结论称为\emph{海涅定理}（Heine's theorem），也称函数极限的\emph{序列判别法}（sequential criterion）。
\end{theorem}

\begin{proof}
若函数极限为 $A$，给定 $\varepsilon>0$，先由函数极限得到 $\delta>0$，再由 $x_n\to x_0$ 得到充分大的 $n$ 满足 $0<|x_n-x_0|<\delta$，于是 $|f(x_n)-A|<\varepsilon$。

反之，若函数不以 $A$ 为极限，由式~\eqref{eq:negation-of-function-limit}，存在固定的 $\varepsilon_0>0$，对每个正整数 $n$ 都可选出 $x_n$ 满足
\[
    0<|x_n-x_0|<\min\{\delta_0,1/n\},
    \qquad |f(x_n)-A|\geq\varepsilon_0.
\]
这样 $x_n\to x_0$，而 $f(x_n)$ 不趋于 $A$，与假设矛盾。
\end{proof}

定理同样适用于单侧极限、无穷远处的极限以及广义极限，只需要求测试序列遵守相应的趋近方向。\emph[caution]{必须检查所有允许的序列；只检查一个序列，即使是 $x_n=x_0+1/n$，也不足以确认函数极限。}反过来，要否定极限存在，只需找出两条允许的序列，使对应函数值趋于不同的数，或者找出一条使函数值不收敛的序列。

\begin{example}
\label{ex:oscillatory-function-no-limit}
证明 $\lim_{x\to0}\sin(1/x)$ 不存在。
\end{example}

\begin{solve}
取 $x_n=1/(n\pi)$ 与 $y_n=1/(2n\pi+\pi/2)$，则两者都从右侧趋于零，而
\[
    \sin\frac1{x_n}=0,\qquad \sin\frac1{y_n}=1.
\]
由定理~\ref{thm:heine-sequential-criterion}，右极限已不存在，因而双侧极限也不存在。这与 $x\sin(1/x)\to0$ 并不矛盾：后者的振幅被趋于零的因子 $x$ 控制住了。
\end{solve}

\subsection{极限存在性定理和两个重要极限}
\label{subsec:limit-existence-and-two-important-limits}

极限的定义适合验证候选值，存在性定理则可以在尚未知道极限值时先判断它是否存在。单调性用确界确定单侧极限，柯西条件则通过比较邻近点的函数值来避免预先指定极限。

\begin{theorem}[单调函数的单侧极限]
\label{thm:monotone-function-one-sided-limits}
设 $f$ 在 $(a,b)$ 上单调递增，其中 $a,b$ 可以是无穷端点，则在扩展实数意义下
\begin{equation}
    \lim_{x\to a+}f(x)=\inf_{a<x<b}f(x),
    \qquad
    \lim_{x\to b-}f(x)=\sup_{a<x<b}f(x).
    \label{eq:monotone-function-endpoint-limits}
\end{equation}
无穷端点处分别理解为 $x\to-\infty$ 或 $x\to+\infty$。单调递减时，两个公式中的下确界、上确界互换。特别地，实值单调函数在区间每个内点都有有限的左右极限。
\end{theorem}

\begin{proof}
以递增函数的左端极限为例。记 $L=\inf f((a,b))$。若 $L$ 有限，由下确界性质，对任意 $\varepsilon>0$，存在 $c\in(a,b)$ 使 $L\leq f(c)<L+\varepsilon$。当 $a<x<c$ 时，由单调性有 $L\leq f(x)\leq f(c)<L+\varepsilon$，故左端极限为 $L$。若 $L=-\infty$，则对每个 $M>0$，可选 $c$ 使 $f(c)<-M$，于是 $a<x<c$ 时均有 $f(x)<-M$。右端极限及递减情形同理。在任一内点附近，用该点两旁固定位置的函数值夹住函数，便排除了无穷单侧极限。
\end{proof}

\begin{theorem}[函数极限的柯西收敛准则]
\label{thm:cauchy-criterion-for-function-limits}
设 $f$ 在 $x_0$ 的某个去心邻域内有定义，则 $\lim_{x\to x_0}f(x)$ 存在且有限，当且仅当
\begin{equation}
    \forall\varepsilon>0,\quad\exists\delta>0,\quad
    x,y\in U_0(x_0,\delta)\ \Longrightarrow\ |f(x)-f(y)|<\varepsilon.
    \label{eq:cauchy-condition-for-function-limits}
\end{equation}
\end{theorem}

\begin{proof}
若极限为 $A$，使 $|f(x)-A|$ 与 $|f(y)-A|$ 都小于 $\varepsilon/2$ 即得必要性。反之，固定一条位于定义范围内且趋于 $x_0$ 的序列 $\{x_n\}$，由式~\eqref{eq:cauchy-condition-for-function-limits}，$\{f(x_n)\}$ 是柯西列。由\crosschapterref{2}{定理~\ref{thm:cauchy-convergence-criterion}}{序列的柯西收敛准则}，存在 $A\in\mathbb{R}$ 使 $f(x_n)\to A$。给定 $\varepsilon>0$，先将柯西条件中的误差取为 $\varepsilon/2$，得到 $\delta$；再固定足够大的 $n$，使 $x_n\in U_0(x_0,\delta)$ 且 $|f(x_n)-A|<\varepsilon/2$。于是对所有 $x\in U_0(x_0,\delta)$，
\[
    |f(x)-A|\leq|f(x)-f(x_n)|+|f(x_n)-A|<\varepsilon,
\]
所以 $f(x)\to A$。此证明只需选取一条固定序列，再利用柯西条件控制全部邻近点。
\end{proof}

\begin{theorem}[两个重要极限]
\label{thm:two-fundamental-function-limits}
当角度采用弧度制时，有
\begin{equation}
    \lim_{x\to0}\frac{\sin x}{x}=1,
    \qquad
    \lim_{t\to\infty}\left(1+\frac1t\right)^t=e.
    \label{eq:two-fundamental-function-limits}
\end{equation}
第二式在 $t\to+\infty$ 和 $t\to-\infty$ 两个方向均成立，且在对应尾部底数为正；等价地，
\begin{equation}
    \lim_{x\to0}(1+x)^{1/x}=e,
    \label{eq:fundamental-exponential-limit}
\end{equation}
其中取 $x>-1$ 且 $x\neq0$。
\end{theorem}

\begin{proof}
由单位圆中的三角形、扇形与外切三角形的面积比较，$0<x<\pi/2$ 时有 $\sin x<x<\tan x$，因此
\[
    \cos x<\frac{\sin x}{x}<1.
\]
这一比较也给出 $|\sin x|\leq|x|$；结合 $1-\cos x=2\sin^2(x/2)\leq x^2/2$，得到 $\cos x\to1$。由夹逼原理及 $\sin x/x$ 的偶性，第一个极限成立。弧度制使扇形面积恰为 $x/2$，所以该单位约定不能省略。

第二式以序列极限 $e=\lim_{n\to\infty}(1+1/n)^n$ 为基础。这个序列的收敛性也可由二项式展开核实：其第 $k$ 个展开项为 $\frac1{k!}\prod_{j=0}^{k-1}(1-j/n)$，随 $n$ 增大而不减，且由 $k!\geq2^{k-1}$（$k\geq1$），各项之和不超过 $\sum_{k=0}^n1/k!<3$，因而由单调收敛原理收敛。以下只用正底数实幂的代数性质与单调性，其构造与连续性见命题~\ref{prop:continuity-of-exponential-logarithmic-power-functions}。对实数 $t>1$，令 $n=\lfloor t\rfloor$，则
\[
    \left(1+\frac1{n+1}\right)^n
    \leq\left(1+\frac1t\right)^t
    \leq\left(1+\frac1n\right)^{n+1}.
\]
两端分别为 $\frac{(1+1/(n+1))^{n+1}}{1+1/(n+1)}$ 与 $(1+1/n)^n(1+1/n)$，均趋于 $e$，故正方向的极限为 $e$。当 $t\to-\infty$ 时，令 $s=-t-1\to+\infty$，则
\[
    \left(1+\frac1t\right)^t
    =\left(1+\frac1s\right)^{s+1}
    =\left(1+\frac1s\right)^s\left(1+\frac1s\right)\to e.
\]
最后令 $t=1/x$，便得到式~\eqref{eq:fundamental-exponential-limit}。
\end{proof}

\begin{example}
\label{ex:fundamental-trigonometric-limit-use}
求 $\lim_{x\to0}\sin(\sin x)/x$。
\end{example}

\begin{solve}
在充分小的去心邻域内 $\sin x\neq0$，并且 $\sin x\to0$，故
\[
    \frac{\sin(\sin x)}x
    =\frac{\sin(\sin x)}{\sin x}\frac{\sin x}x\to1\cdot1=1.
    \qedhere
\]
\end{solve}

普通极限不存在时，还可像序列一样描述邻近取值的上下边界。设 $f$ 在 $U_0(x_0,\delta_0)$ 内有定义，对 $0<\delta<\delta_0$，令
\[
    l(\delta)=\inf_{0<|x-x_0|<\delta}f(x),
    \qquad h(\delta)=\sup_{0<|x-x_0|<\delta}f(x).
\]
确界允许取扩展实数。随着 $\delta$ 缩小，$l(\delta)$ 不减，$h(\delta)$ 不增，因而可以定义\emph{函数的下极限与上极限}（limit inferior and limit superior）
\begin{equation}
    \liminf_{x\to x_0}f(x)=\lim_{\delta\to0+}l(\delta),
    \qquad
    \limsup_{x\to x_0}f(x)=\lim_{\delta\to0+}h(\delta).
    \label{eq:function-liminf-limsup}
\end{equation}
由 $l(\delta)\leq f(x)\leq h(\delta)$ 及确界定义可得：有限极限 $\lim_{x\to x_0}f(x)=A$ 存在，当且仅当上下极限都等于同一个\emph{有限数} $A$。具体地，极限存在时，任给 $\varepsilon>0$，充分小的邻域满足 $A-\varepsilon\leq l(\delta)\leq h(\delta)\leq A+\varepsilon$；反之，上下界同时趋于 $A$ 就夹住了全部邻近函数值。若上下极限共同等于 $+\infty$ 或 $-\infty$，得到的是相应的广义极限。例如，$\sin(1/x)$ 在 $0$ 的上下极限分别为 $-1,1$，而 $1/x^2$ 的上下极限均为 $+\infty$。

\section{函数的连续与间断}
\label{sec:continuity-and-discontinuity}

\subsection{函数的连续与间断}
\label{subsec:continuity-and-discontinuity}

% 教材第 103--109 页；分章 PDF 第 21--27 页。
极限只描述附近的变化，连续性还要求这种变化与该点的函数值衔接。讨论连续性时必须把定义域和端点包括在内；同一个表达式限制在不同定义域上时，允许的趋近方向也会改变。

\begin{definition}
\label{def:continuity-at-a-point-and-on-intervals}
设 $f:I\to\mathbb{R}$，其中 $I$ 为区间，$x_0\in I$。若
\begin{equation}
    \forall\varepsilon>0,\quad\exists\delta>0,\quad
    x\in I,\ |x-x_0|<\delta
    \ \Longrightarrow\ |f(x)-f(x_0)|<\varepsilon,
    \label{eq:epsilon-delta-continuity}
\end{equation}
则称 $f$ 在 $x_0$ \emph{连续}（continuous），否则称它在该点\emph{间断}（discontinuous）。当 $x_0$ 为内点时，这等价于 $\lim_{x\to x_0}f(x)=f(x_0)$。若只要求右侧或左侧的点满足上述条件，分别称为\emph{右连续}（right-continuous）或\emph{左连续}（left-continuous），相应条件是 $f(x_0+0)=f(x_0)$ 或 $f(x_0-0)=f(x_0)$。

若 $f$ 在区间 $I$ 的每一点都连续，则称 $f$ 在 $I$ 上连续，记为 $f\in C(I)$。特别地，$f\in C([a,b])$ 要求它在 $(a,b)$ 上连续、在 $a$ 处右连续且在 $b$ 处左连续。
\end{definition}

连续性也可以写成增量形式：令 $\Delta x=x-x_0$、$\Delta y=f(x_0+\Delta x)-f(x_0)$，则连续就是 $\Delta x\to0$ 时 $\Delta y\to0$，其中 $x_0+\Delta x$ 必须在定义域内。由海涅定理~\ref{thm:heine-sequential-criterion}，它还等价于：对任意 $\{x_n\}\subseteq I$，只要 $x_n\to x_0$，就有 $f(x_n)\to f(x_0)$。这里允许 $x_n=x_0$，因为这时函数值误差已经为零。

\begin{definition}
\label{def:classification-of-discontinuities}
设 $x_0$ 是 $f$ 的内点间断点。若左右极限都存在且有限，则称 $x_0$ 为\emph{第一类间断点}（discontinuity of the first kind）；其中，左右极限相等而不等于 $f(x_0)$ 时，称为\emph{可去间断点}（removable discontinuity），左右极限不等时，称为\emph{跳跃间断点}（jump discontinuity）。若至少一个单侧有限极限不存在，则称为\emph{第二类间断点}（discontinuity of the second kind）；其中包括单侧趋于无穷以及振荡而无极限的情形。
\end{definition}

可去间断点可以且只能通过把该点函数值改为共同的单侧极限来消除，跳跃与第二类间断不能靠改变一个点的值消除。若原函数在 $x_0$ 没有定义，则严格说来 $x_0$ 不属于其定义域；此时若有限极限 $A$ 存在，可补定义 $f(x_0)=A$，得到该点的\emph{连续延拓}（continuous extension）。类似地，$f\in C((a,b))$ 能连续延拓到 $[a,b]$，当且仅当两个端点的内侧极限均存在且有限，且端点值由这两个极限唯一确定。

\begin{example}
\label{ex:types-of-discontinuities}
讨论下列函数在 $0$ 处的连续性与间断类型：
\[
    f(x)=\begin{cases}x\sin(1/x),&x\neq0,\\1,&x=0,\end{cases}
    \qquad g(x)=\lfloor x\rfloor,
\]
\[
    h(x)=\begin{cases}e^{1/x},&x\neq0,\\0,&x=0,\end{cases}
    \qquad q(x)=\begin{cases}\sin(1/x),&x\neq0,\\0,&x=0.\end{cases}
\]
\end{example}

\begin{solve}
$f$ 的左右极限都为 $0$，而 $f(0)=1$，故 $0$ 是可去间断点，改取 $f(0)=0$ 即可连续。$g(0-0)=-1$、$g(0+0)=g(0)=0$，故 $g$ 在 $0$ 右连续，但该点是跳跃间断点。$h(0-0)=0$、$h(0+0)=+\infty$，故它在 $0$ 左连续，但该点属于第二类间断点。$q$ 的两侧极限均因振荡而不存在，故 $0$ 也是第二类间断点。后二者说明，第二类间断既可能伴随无界性，也可能发生在有界函数中；关于指数函数的连续性及尾部极限，见第~\ref{subsec:continuity-of-elementary-functions} 小节。
\end{solve}

\begin{example}
\label{ex:riemann-function-continuity}
在 $(0,1)$ 上定义\emph{黎曼函数}（Riemann function，也称 Thomae function）
\[
    R(x)=\begin{cases}
        1/q,&x=p/q\in\mathbb{Q},\ p,q\text{ 为互素正整数},\\
        0,&x\notin\mathbb{Q}.
    \end{cases}
\]
证明它在每个无理数点连续，在每个有理数点具有可去间断。
\end{example}

\begin{solve}
固定 $x_0\in(0,1)$。给定 $\varepsilon>0$，取正整数 $N$ 使 $1/N<\varepsilon$。区间内分母不超过 $N$ 的既约分数只有有限个，故可选 $\delta>0$，使 $U_0(x_0,\delta)\subseteq(0,1)$ 且排除这些有限个分数中所有异于 $x_0$ 的点。于是在该去心邻域内，有理数点的分母均大于 $N$，无理数点的值为零，因此
\[
    0\leq R(x)<1/N<\varepsilon.
\]
所以对每个 $x_0$ 都有 $\lim_{x\to x_0}R(x)=0$。无理数点满足 $R(x_0)=0$，故连续；有理数点满足 $R(x_0)>0$，故为可去间断点。这一例子表明，函数可以在稠密的点集上连续，同时在另一个稠密点集上间断。
\end{solve}

\begin{proposition}[单调函数的间断类型]
\label{prop:monotone-function-discontinuities}
实值单调函数在区间内部只可能有跳跃间断点。
\end{proposition}

\begin{proof}
以递增函数为例。由定理~\ref{thm:monotone-function-one-sided-limits}，内点 $x_0$ 的两个单侧极限均存在且有限，并由单调性有
\[
    f(x_0-0)\leq f(x_0)\leq f(x_0+0).
\]
若左右极限相等，则三者相等，函数在该点连续。因此只要它不连续，就必有 $f(x_0-0)<f(x_0+0)$，即发生跳跃。递减情形同理。
\end{proof}

\subsection{连续函数的性质}
\label{subsec:properties-of-continuous-functions}

% 教材第 109--111 页；分章 PDF 第 27--29 页。
连续函数在一点附近的性质可以直接由函数极限导出，但此时中心点也必须计入。连续性还允许把函数运算与极限过程交换，这是许多极限计算和后续微积分论证的基础。

\begin{proposition}[连续性的局部性质与运算]
\label{prop:local-properties-and-operations-of-continuity}
若 $f$ 在 $x_0$ 连续，则它在该点某个相对邻域内有界；若 $f(x_0)>0$，则对任意 $c\in(0,f(x_0))$，存在相对邻域使其中恒有 $f(x)>c$，负值情形类似。若 $f,g$ 都在 $x_0$ 连续，则 $f\pm g$、$fg$、$|f|$ 在该点连续；$g(x_0)\neq0$ 时，$f/g$ 也在该点连续。
\end{proposition}

\begin{proof}
局部有界性取式~\eqref{eq:epsilon-delta-continuity} 中的误差为 $1$ 即得。保号性取误差为 $f(x_0)-c$ 即得。四则运算由式~\eqref{eq:function-limit-arithmetic} 导出；绝对值的连续性则由反向三角不等式
\[
    \bigl||f(x)|-|f(x_0)|\bigr|\leq|f(x)-f(x_0)|
\]
得到。这些论证也适用于端点的单侧连续性。
\end{proof}

因此，连续函数的有限和、有限积仍连续，且
\[
    \max\{f,g\}=\frac{f+g+|f-g|}{2},
    \qquad \min\{f,g\}=\frac{f+g-|f-g|}{2}
\]
也连续。这里的运算次数有限；无限个连续函数的极限是否连续，需要后续函数列理论中的额外条件。

\begin{theorem}[复合函数的连续性]
\label{thm:composition-of-continuous-functions}
设 $g$ 在 $x_0$ 连续，$f$ 在 $u_0=g(x_0)$ 连续，且 $g$ 在 $x_0$ 附近的取值属于 $f$ 的定义域，则 $f\circ g$ 在 $x_0$ 连续。
\end{theorem}

\begin{proof}
给定 $\varepsilon>0$，先由 $f$ 在 $u_0$ 连续选取 $\eta>0$，使 $|u-u_0|<\eta$ 时 $|f(u)-f(u_0)|<\varepsilon$；再由 $g$ 在 $x_0$ 连续选取 $\delta>0$，使 $|x-x_0|<\delta$ 时 $|g(x)-g(x_0)|<\eta$。代入即得 $|f(g(x))-f(g(x_0))|<\varepsilon$。由于外层函数在中心点也满足这一估计，不需要排除 $g(x)=u_0$。
\end{proof}

\begin{theorem}[反函数的连续性]
\label{thm:continuity-of-inverse-functions}
设 $I$ 为非退化区间，$f:I\to\mathbb{R}$ 连续且严格单调，则 $f$ 的值域 $J=f(I)$ 是区间，反函数 $f^{-1}:J\to I$ 连续且具有与 $f$ 相同的严格单调方向。端点处的连续性按相对定义域理解。
\end{theorem}

\begin{proof}
严格单调性保证单射性，从而反函数在 $J$ 上存在，且具有相同的单调方向。$J$ 是区间将由下一节的介值定理~\ref{thm:intermediate-value-property} 得到；下面的连续性证明只需对 $y\in J$ 考察，不依赖尚未证明的介值性。

不妨设 $f$ 严格递增。固定 $y_0=f(x_0)$，先设 $x_0$ 是 $I$ 的内点。给定 $\varepsilon>0$，在 $I$ 内选 $x_-,x_+$ 使
\[
    x_0-\varepsilon<x_-<x_0<x_+<x_0+\varepsilon.
\]
令 $\delta=\min\{y_0-f(x_-),f(x_+)-y_0\}>0$。若 $y\in J$ 且 $|y-y_0|<\delta$，由严格单调性得
\[
    x_-<f^{-1}(y)<x_+,
\]
从而 $|f^{-1}(y)-x_0|<\varepsilon$。若 $x_0$ 是属于 $I$ 的左端点，只选 $x_+$ 并利用 $f^{-1}(y)\geq x_0$；右端点类似。递减情形可对 $-f$ 作同样论证。
\end{proof}

\subsection{初等函数的连续性}
\label{subsec:continuity-of-elementary-functions}

% 教材第 111--113 页；分章 PDF 第 29--31 页。
初等函数的连续性由基本初等函数逐步建立，再通过有限次四则运算与复合传递。多项式的连续性来自 $x^n$ 的极限；有理函数在分母非零处连续；三角函数则可先由和差公式得到
\[
    |\sin x-\sin y|\leq|x-y|,
    \qquad |\cos x-\cos y|\leq|x-y|,
\]
从而正弦、余弦处处连续，其余三角函数在相应分母非零处连续。证明这些结论不需要预先使用导数。

对于实指数，先说明无理指数的含义。给定 $a>1$，有理数 $r=p/q$ 的幂由正数开方定义，负指数通过倒数定义；对任意实数 $x$，可统一写成
\begin{equation}
    a^x=\sup\{a^r:r\in\mathbb{Q},\ r<x\}.
    \label{eq:real-exponent-via-supremum}
\end{equation}
这个集合非空且有上界，因为可以选整数 $m<x<n$，使 $a^m$ 属于集合，而 $a^n$ 是上界。该定义与有理指数一致：由\crosschapterref{1}{伯努利不等式~\eqref{eq:bernoulli-inequality}}{伯努利不等式}，$a^{1/N}\to1$，因为若 $a^{1/N}\geq1+\eta$，则 $a\geq(1+\eta)^N\geq1+N\eta$，对固定 $\eta>0$ 和充分大 $N$ 不可能成立。对有理数 $x$，从下方用 $x-1/N$ 逼近即可验证式~\eqref{eq:real-exponent-via-supremum} 的确等于原来的 $a^x$。

\begin{proposition}[指数、对数与幂函数的连续性]
\label{prop:continuity-of-exponential-logarithmic-power-functions}
对 $a>0$，指数函数 $a^x$ 在 $\mathbb{R}$ 上连续；$a\neq1$ 时，其反函数 $\log_a x$ 在 $(0,+\infty)$ 上连续。对任意 $\alpha\in\mathbb{R}$，幂函数 $x^\alpha=e^{\alpha\ln x}$ 在 $(0,+\infty)$ 上连续。
\end{proposition}

\begin{proof}
先设 $a>1$。由确界定义及有理数稠密性，若 $x_1<x_2$，可选有理数 $r,s$ 满足 $x_1<r<s<x_2$，得到 $a^{x_1}\leq a^r<a^s\leq a^{x_2}$，故指数函数严格递增。固定 $x_0\in\mathbb{R}$，对正整数 $N$，令
\[
    r_N=\frac{\lfloor Nx_0\rfloor-1}{N},
    \qquad s_N=\frac{\lfloor Nx_0\rfloor+2}{N}.
\]
则 $r_N<x_0<s_N$、$s_N-r_N=3/N$。当 $x$ 足够靠近 $x_0$，使 $r_N<x<s_N$ 时，$a^x,a^{x_0}$ 均在 $[a^{r_N},a^{s_N}]$ 内，因而
\[
    |a^x-a^{x_0}|\leq a^{s_N}-a^{r_N}
    =a^{r_N}(a^{3/N}-1)\leq a^m(a^{3/N}-1),
\]
其中 $m>x_0$ 为固定整数。已知 $a^{1/N}\to1$，故先选充分大的 $N$，再选相应邻域，就能使此误差任意小，证明连续性。$a=1$ 时为常函数；$0<a<1$ 时定义 $a^x=1/(1/a)^x$，由商的连续性得到结论。

当 $a>1$ 时，$a^n\to+\infty$、$a^{-n}\to0$，结合单调性和介值性，指数函数的值域为 $(0,+\infty)$；$0<a<1$ 时同理。因此由反函数连续性定理，对数函数连续。幂函数的连续性再由指数、对数的复合得到。
\end{proof}

严格单调的三角函数分支的反函数也连续：例如，$\arcsin x$ 在 $[-1,1]$ 上连续，端点处按单侧理解；$\arccos x$ 在 $[-1,1]$ 上连续，$\arctan x$ 在 $\mathbb{R}$ 上连续。于是由有限次运算得到的\emph{初等函数在其定义域内连续}。\emph[caution]{“在定义域内”是结论的一部分：例如，$\ln x$ 不在 $0$ 有定义，$1/x$ 的定义域不含 $0$，不能通过直接代入这些点求有限极限。}

\begin{example}
\label{ex:variable-exponent-function-limit}
设 $u(x)\to a>0$、$v(x)\to b\in\mathbb{R}$（$x\to x_0$），证明 $u(x)^{v(x)}\to a^b$，并求 $\lim_{x\to0}(1-2x)^{1/x}$。
\end{example}

\begin{solve}
由 $a>0$ 及局部保号性，$u(x)$ 在充分小的去心邻域内为正，可写成
\[
    u(x)^{v(x)}=e^{v(x)\ln u(x)}\to e^{b\ln a}=a^b.
\]
第二问的指数趋于无穷，不能使用刚证明的有限指数结论。令 $t=-2x$，则
\[
    (1-2x)^{1/x}=\bigl((1+t)^{1/t}\bigr)^{-2}\to e^{-2},
\]
其中先用式~\eqref{eq:fundamental-exponential-limit}，再用正数上的幂函数连续性。这说明处理变指数极限时，应先判断指数是否有有限极限。
\end{solve}

\section{闭区间上连续函数的基本性质}
\label{sec:continuous-functions-on-closed-intervals}

% 教材第 113--121 页；分章 PDF 第 31--39 页。
逐点连续给出的邻域大小和局部界可以随点变化，而闭区间的有限覆盖性质与有界序列的子列性质，使这些局部信息能够推出整体结论。以下除另有说明外，均设 $a,b\in\mathbb{R}$ 且 $a<b$；闭性与有界性都是区间条件中不可随意删去的部分。

\begin{theorem}[有界性与最值定理]
\label{thm:boundedness-and-extreme-values}
若 $f\in C([a,b])$，则 $f$ 在 $[a,b]$ 上有界，并存在 $\xi,\eta\in[a,b]$，使
\begin{equation}
    f(\xi)=\min_{x\in[a,b]}f(x)=m,
    \qquad f(\eta)=\max_{x\in[a,b]}f(x)=M.
    \label{eq:extreme-values-on-closed-interval}
\end{equation}
\end{theorem}

\begin{proof}
先证有界性。由连续性的局部有界性，每个 $x\in[a,b]$ 都有邻域 $U(x,\delta_x)$，使 $f$ 在该邻域与 $[a,b]$ 的交集上有界。对这些邻域应用\crosschapterref{2}{有限覆盖定理~\ref{thm:finite-covering}}{有限覆盖定理}，取有限子覆盖，再取相应有限多个局部界的最大值，就得到整个区间上的界。\crosschapterref{2}{这也完成了例~\ref{ex:continuous-function-bounded-by-finite-cover} 中提出的证明。}{}

由有界性，$M=\sup\{f(x):x\in[a,b]\}$ 是实数。由上确界定义，对每个正整数 $n$ 可选 $x_n\in[a,b]$，使 $M-1/n<f(x_n)\leq M$，于是 $f(x_n)\to M$。由\crosschapterref{2}{定理~\ref{thm:bolzano-weierstrass-sequence}}{波尔查诺--魏尔斯特拉斯定理}，存在子列 $x_{n_k}\to\eta$；闭区间条件保证 $\eta\in[a,b]$。再由连续性，
\[
    f(\eta)=\lim_{k\to\infty}f(x_{n_k})=M.
\]
对 $-f$ 应用同一论证，便得到最小值。
\end{proof}

\emph[caution]{有界只保证确界有限，最值定理还保证确界实际取到。}例如，$1/x$ 在 $(0,1)$ 上连续却无界，$x$ 在 $(0,1)$ 上连续且有界却没有最大值与最小值；$x$ 在 $\mathbb{R}$ 上连续也不能推出有界性。这些例子分别说明开区间和无界区间不足以替代有界闭区间。

\begin{corollary}[严格正的连续函数有统一正下界]
\label{cor:positive-continuous-function-lower-bound}
若 $f\in C([a,b])$ 且 $f(x)>0$ 对每个 $x\in[a,b]$ 成立，则存在 $c>0$ 使 $f(x)\geq c$ 对整个区间成立。特别地，若连续函数在闭区间上处处非零，则 $|f|$ 有正下界，$1/f$ 连续且有界。
\end{corollary}

\begin{proof}
由最值定理，存在 $\xi\in[a,b]$ 使 $c=f(\xi)=\min f$。由于最小值在区间内取到，且函数处处为正，故 $c>0$。对 $|f|$ 应用这一结论即得后半部分。
\end{proof}

\begin{theorem}[介值定理]
\label{thm:intermediate-value-property}
设 $f$ 在区间 $I$ 上连续，$x_1,x_2\in I$。若实数 $c$ 介于 $f(x_1)$ 与 $f(x_2)$ 之间，则存在介于 $x_1,x_2$ 之间的点 $\xi\in I$ 使 $f(\xi)=c$；当 $c$ 严格介于两个函数值之间时，可以取 $\xi$ 严格介于两个点之间。这一性质称为\emph{介值性}（intermediate value property）。因此 $f(I)$ 是区间；特别地，若 $f\in C([a,b])$，则
\begin{equation}
    f([a,b])=[m,M],
    \label{eq:continuous-image-of-closed-interval}
\end{equation}
其中 $m,M$ 由式~\eqref{eq:extreme-values-on-closed-interval} 给出。
\end{theorem}

\begin{proof}
端点值相等或 $c$ 等于端点值时直接取相应点。对于严格介于的情况，通过交换两点或考虑 $-f$，只需证明 $u<v$ 且 $f(u)<c<f(v)$ 的情形。令
\[
    E=\{x\in[u,v]:f(x)\leq c\},\qquad \xi=\sup E.
\]
$E$ 非空且有上界。由 $u,v$ 处的连续性与严格不等式，$u$ 右侧有一段仍满足 $f(x)<c$，$v$ 左侧有一段满足 $f(x)>c$，故 $u<\xi<v$。由确界性质，存在 $t_n\in E$ 使 $t_n\to\xi$，因此连续性给出 $f(\xi)=\lim f(t_n)\leq c$。若 $f(\xi)<c$，则由连续性，$\xi$ 右侧足够靠近的点仍满足 $f(x)<c$，从而属于 $E$，与 $\xi$ 为上界矛盾。故 $f(\xi)=c$。

任取 $f(I)$ 中两个值，它们之间的值都由上述结论取到，所以 $f(I)$ 是区间。闭区间情形再结合最大值、最小值的存在性，就得到式~\eqref{eq:continuous-image-of-closed-interval}。
\end{proof}

\begin{corollary}[零点存在定理]
\label{cor:bolzano-zero-existence}
若 $f\in C([a,b])$ 且 $f(a)f(b)<0$，则存在 $\xi\in(a,b)$ 使 $f(\xi)=0$。此结论也称\emph{波尔查诺零点定理}（Bolzano's theorem），由介值定理取中间值 $0$ 得到。
\end{corollary}

\emph[caution]{端点异号保证至少有一个零点，不保证唯一。}例如，$x(x^2-1)$ 在 $[-2,2]$ 上满足端点异号，却有三个零点。若另有严格单调性，则同一个函数值至多取到一次，可以把存在性与唯一性结合起来。介值定理还表明，对连续函数的有限个取值 $f(x_1),\ldots,f(x_n)$，其算术平均介于这些值的最小值与最大值之间，因而也必是区间内某点的函数值。

\begin{proposition}[连续函数的单射性与严格单调性]
\label{prop:continuous-injective-function-strict-monotonicity}
设 $f$ 在非退化区间 $I$ 上连续，则它为单射的充分必要条件是它严格单调。因此，区间上的连续函数在其值域上存在反函数，当且仅当它严格单调。
\end{proposition}

\begin{proof}
严格单调必为单射。反之，若 $f$ 连续且为单射，对任意 $x_1<x_2<x_3$，$f(x_2)$ 必严格介于 $f(x_1),f(x_3)$ 之间。否则它比两端值都大或都小，在中间函数值与两端函数值之间选取一个适当的值，由介值定理，该值会在 $(x_1,x_2)$ 和 $(x_2,x_3)$ 内各取到一次，违反单射性。

固定 $u<v$。若 $f(u)<f(v)$，任取 $s<t$，把 $u,v,s,t$ 的不同点按大小排列。上述三点性质保证其中任意相邻两段的函数值变化方向一致，而 $f(u)<f(v)$ 又确定这一方向只能递增，故 $f(s)<f(t)$。若 $f(u)>f(v)$，同理得到严格递减。
\end{proof}

对于一般的单调函数，连续性的另一判据是其值域是否为区间：连续性由介值定理保证值域为区间；反之，单调函数若发生间断，有限的单侧极限之间必有跳跃，从而遗漏一段值中的点，值域就不可能是区间。端点处同样考察函数值与内侧极限之间的跳跃。

\begin{example}
\label{ex:continuous-zero-and-bisection}
证明方程 $x^4+2x-1=0$ 在 $(0,1)$ 内有唯一解，并说明如何用二分法给出带误差保证的近似解。
\end{example}

\begin{solve}
令 $F(x)=x^4+2x-1$。它在 $[0,1]$ 上连续，且 $F(0)=-1<0<F(1)=2$，故存在零点。若 $0\leq x<y\leq1$，则
\[
    F(y)-F(x)=(y-x)(y^3+y^2x+yx^2+x^3+2)>0,
\]
故 $F$ 严格递增，零点 $\xi$ 唯一。

取初始区间 $[a_0,b_0]=[0,1]$。每步令 $c_n=(a_n+b_n)/2$；若 $F(c_n)=0$，便得到精确解并终止；否则保留使端点函数值异号的半区间。零点始终在保留区间内，进行 $n$ 次二分后 $b_n-a_n=2^{-n}$，所以其中点 $c_n$ 满足
\[
    |c_n-\xi|\leq\frac{b_n-a_n}{2}=2^{-(n+1)}.
\]
给定误差容限 $\varepsilon>0$，选取 $2^{-(n+1)}\leq\varepsilon$ 即可停止并输出 $c_n$。例如，八次二分后得到区间 $[121/256,122/256]$，其中点 $243/512=0.474609375$ 的误差不超过 $1/512$。二分法的根存在性来自连续与异号条件，误差保证来自区间长度；唯一性则由另外证明的严格单调性提供。
\end{solve}

逐点连续仍允许不同位置需要不同的邻域半径。为了在整个区间内统一控制函数值的变化，需要加强量词要求。

\begin{definition}
\label{def:uniform-continuity}
设 $f:I\to\mathbb{R}$。若
\begin{equation}
    \forall\varepsilon>0,\quad\exists\delta>0,\quad
    \forall x,y\in I,\quad
    |x-y|<\delta\ \Longrightarrow\ |f(x)-f(y)|<\varepsilon,
    \label{eq:uniform-continuity}
\end{equation}
则称 $f$ 在 $I$ 上\emph{一致连续}（uniformly continuous）。
\end{definition}

逐点连续的次序是“对每个 $x_0$，对每个 $\varepsilon$，存在 $\delta$”，而一致连续要求“对每个 $\varepsilon$，存在同时适用于所有点对的 $\delta$”。因此一致连续蕴含连续，但反向一般不成立。否定式为
\begin{equation}
    \exists\varepsilon_0>0,\quad\forall\delta>0,\quad
    \exists x,y\in I,\quad |x-y|<\delta,
    \quad |f(x)-f(y)|\geq\varepsilon_0.
    \label{eq:negation-of-uniform-continuity}
\end{equation}

\begin{theorem}[一致连续的序列判据]
\label{thm:sequential-criterion-uniform-continuity}
函数 $f$ 在区间 $I$ 上一致连续，当且仅当对任意两列 $x_n,y_n\in I$，只要 $|x_n-y_n|\to0$，就有 $|f(x_n)-f(y_n)|\to0$。这里不要求 $x_n,y_n$ 各自收敛。
\end{theorem}

\begin{proof}
必要性直接将式~\eqref{eq:uniform-continuity} 用于充分大的 $n$。若一致连续不成立，在式~\eqref{eq:negation-of-uniform-continuity} 中依次取 $\delta=1/n$，得到 $|x_n-y_n|<1/n$，而函数值之差始终不小于固定的 $\varepsilon_0$，与序列条件矛盾。
\end{proof}

\begin{example}
\label{ex:uniform-continuity-examples}
讨论 $\sin x$ 在 $\mathbb{R}$ 上、$1/x$ 在 $(0,+\infty)$ 与 $[c,+\infty)$（$c>0$）上，以及 $\sqrt{x}$ 在 $[0,+\infty)$ 上的一致连续性。
\end{example}

\begin{solve}
由 $|\sin x-\sin y|\leq|x-y|$，对任意 $\varepsilon>0$ 取 $\delta=\varepsilon$，得正弦函数在整条实轴上一致连续。对 $1/x$，取 $x_n=1/n$、$y_n=1/(n+1)$，则两点距离趋于零，但
\[
    \left|\frac1{x_n}-\frac1{y_n}\right|=1,
\]
故它在 $(0,+\infty)$ 上不一致连续；在 $[c,+\infty)$ 上则有
\[
    \left|\frac1x-\frac1y\right|=\frac{|x-y|}{xy}\leq\frac{|x-y|}{c^2},
\]
取 $\delta=c^2\varepsilon$ 即可。最后，对 $x,y\geq0$，不妨设 $x\geq y$，则
\[
    (\sqrt{x}-\sqrt{y})^2\leq(\sqrt{x}-\sqrt{y})(\sqrt{x}+\sqrt{y})=x-y,
\]
所以 $|\sqrt{x}-\sqrt{y}|\leq\sqrt{|x-y|}$，取 $\delta=\varepsilon^2$ 即得一致连续。这也说明在无界区间上一致连续的函数可以无界。
\end{solve}

\begin{theorem}[康托尔定理]
\label{thm:heine-cantor-uniform-continuity}
若 $f\in C([a,b])$，则 $f$ 在 $[a,b]$ 上一致连续。这一结论称为\emph{康托尔定理}（Heine--Cantor theorem）。
\end{theorem}

\begin{proof}
反设不一致连续。由式~\eqref{eq:negation-of-uniform-continuity}，可选两列 $x_n,y_n\in[a,b]$，使 $|x_n-y_n|<1/n$ 且 $|f(x_n)-f(y_n)|\geq\varepsilon_0>0$。由有界序列的子列定理，可选子列 $x_{n_k}\to\xi\in[a,b]$。又因
\[
    |y_{n_k}-\xi|\leq|y_{n_k}-x_{n_k}|+|x_{n_k}-\xi|\to0,
\]
故 $y_{n_k}\to\xi$。在 $\xi$ 处的连续性给出 $f(x_{n_k}),f(y_{n_k})\to f(\xi)$，于是两者之差趋于零，与固定的正下界矛盾。
\end{proof}

\begin{proposition}[有穷开区间上的一致连续性]
\label{prop:uniform-continuity-and-endpoint-extension}
设 $a,b$ 为有限实数，$f\in C((a,b))$。则 $f$ 在 $(a,b)$ 上一致连续，当且仅当两个单侧极限 $f(a+0),f(b-0)$ 都存在且有限，也即 $f$ 可以连续延拓到 $[a,b]$。
\end{proposition}

\begin{proof}
若一致连续，给定 $\varepsilon>0$，取相应的 $\delta>0$。任意两个位于 $(a,a+\eta)$ 的点，在 $0<\eta<\min\{\delta,b-a\}$ 时距离小于 $\delta$，其函数值之差便小于 $\varepsilon$。这正是 $x\to a+$ 的柯西条件，故由单侧版本的定理~\ref{thm:cauchy-criterion-for-function-limits}，右极限存在且有限；右端点同理。反之，以这两个极限补上端点值，得到闭区间上的连续函数，由康托尔定理知其一致连续，限制回 $(a,b)$ 后结论仍成立。
\end{proof}

因此，有穷区间上的一致连续函数必有界：开区间情形可通过上述延拓及最值定理看出，其他端点类型同理。但有界连续仍不保证一致连续，例如 $\sin(1/x)$ 在 $(0,1)$ 上有界连续，却没有有限的左端点右极限。一致连续函数的和与差仍一致连续，乘积则需要额外控制；两个有界的一致连续函数的乘积由
\[
    |f(x)g(x)-f(y)g(y)|
    \leq\|f\|_\infty|g(x)-g(y)|+\|g\|_\infty|f(x)-f(y)|
\]
可知一致连续，其中 $\|f\|_\infty=\sup_{x\in I}|f(x)|$，$\|g\|_\infty$ 同理。不加有界条件时，$f(x)=g(x)=x$ 在 $\mathbb{R}$ 上均一致连续，但 $x^2$ 不一致连续：取 $x_n=n$、$y_n=n+1/n$，则两点距离趋于零，平方之差却趋于 $2$。

\section{无穷小量与无穷大量的阶}
\label{sec:orders-of-infinitesimals-and-infinities}

% 教材第 121--127 页；分章 PDF 第 39--45 页。
本节比较的对象是同一极限过程中变化的量，而不是某个固定的“小数”或“大数”。以下以 $x\to x_0$ 为规范；换成单侧趋近或趋于无穷远时，只需相应改变定义中的邻域条件。没有注明极限过程的阶比较是不完整的，例如 $x^2$ 与 $x$ 在 $x\to0$ 和 $x\to+\infty$ 时的相对大小恰好相反。

\begin{definition}
\label{def:infinitesimal-and-infinite-functions}
若 $f(x)\to0$，则称 $f$ 为该过程中的\emph{无穷小量}（infinitesimal）；若 $|f(x)|\to+\infty$，则称为\emph{无穷大量}（infinitely large quantity）。特别地，
\[
    f(x)\to A\in\mathbb{R}
    \quad\Longleftrightarrow\quad f(x)=A+\alpha(x),\quad\alpha(x)\to0.
\]
若 $f$ 在相应的去心邻域内恒不为零，则
\begin{equation}
    f(x)\to0
    \quad\Longleftrightarrow\quad
    \left|\frac1{f(x)}\right|\to+\infty.
    \label{eq:infinitesimal-and-reciprocal-infinity}
\end{equation}
\end{definition}

式~\eqref{eq:infinitesimal-and-reciprocal-infinity} 来自 $0<|f(x)|<1/M$ 与 $|1/f(x)|>M$ 的等价性，其中“恒不为零”保证倒数有定义。有限个无穷小量的和仍是无穷小量，无穷小量与局部有界函数的乘积也是无穷小量；这些均由极限运算或夹逼原理得到。无限多个无穷小量的相加则不能直接套用有限次运算的结论。

\begin{definition}
\label{def:little-o-big-o-and-asymptotic-equivalence}
设 $f,g$ 在相应去心邻域内有定义，且 $g$ 恒不为零。定义\emph{小 $o$ 记号}（little-o notation）、\emph{大 $O$ 记号}（big-O notation）与\emph{渐近等价}（asymptotic equivalence）如下：
\begin{equation}
    \begin{aligned}
        f(x)=o(g(x))
        &\quad\Longleftrightarrow\quad \frac{f(x)}{g(x)}\to0,\\
        f(x)=O(g(x))
        &\quad\Longleftrightarrow\quad
        \exists C>0,\ \exists\delta>0,\quad
        0<|x-x_0|<\delta\ \Longrightarrow\ |f(x)|\leq C|g(x)|,\\
        f(x)\sim g(x)
        &\quad\Longleftrightarrow\quad \frac{f(x)}{g(x)}\to1.
    \end{aligned}
    \label{eq:asymptotic-comparison-notations}
\end{equation}
若 $f,g$ 都是无穷小量且 $f/g\to0$，称 $f$ 是比 $g$ \emph{高阶的无穷小量}（higher-order infinitesimal）；若 $f/g\to c\in\mathbb{R}\setminus\{0\}$，称二者为\emph{同阶无穷小量}（infinitesimals of the same order），$c=1$ 时称为\emph{等价无穷小量}（equivalent infinitesimals）。若二者都是无穷大量，非零有限比值和比值为 $1$ 分别定义同阶与等价无穷大量，而 $f/g\to0$ 表示 $f$ 是比 $g$ \emph{低阶的无穷大量}。
\end{definition}

小 $o$ 与大 $O$ 本身不要求 $f,g$ 都趋于零或无穷。例如，$f=o(1)$ 就是 $f\to0$，$f=O(1)$ 表示局部有界。$f=o(g)$ 必推出 $f=O(g)$，反向不成立；按照本节采用的教材定义，同阶要求比值具有非零有限极限，比值只是在两个正常数之间波动并不足以称为同阶。例如，$f(x)=x(2+\sin(1/x))$ 与 $g(x)=x$ 在 $x\to0$ 时互为大 $O$，但比值没有极限。若 $f\sim g$，则 $f$ 也在充分小的去心邻域内不为零，且
\begin{equation}
    f\sim g
    \quad\Longleftrightarrow\quad
    f=g+o(g).
    \label{eq:asymptotic-equivalence-remainder}
\end{equation}

\emph[caution]{阶记号中的等号表达一种渐近关系，不能当作两个固定函数之间可任意倒转的恒等式。}同一行中不同的 $o(g)$ 可以代表不同函数。例如，$x^2=o(x)$ 也满足 $x^2=O(x)$，只能说明小 $o$ 类包含于大 $O$ 类，不能据此写成 $o(x)=O(x)$；$x=O(x)$ 却不满足 $x=o(x)$。当 $x\to0$ 时，正整数幂的指数越大，无穷小量阶越高；对倒数 $1/x^n$ 而言，指数越大则无穷大量阶越高。若需要非整数幂的双侧比较，使用 $|x-x_0|^\alpha$ 可避免负底数实幂的定义问题。

\begin{proposition}[对数、幂与指数的增长次序]
\label{prop:logarithm-power-exponential-orders}
设 $a>1$、$\alpha>0$。当 $x\to+\infty$ 时，
\begin{equation}
    \ln x=o(x^\alpha),
    \qquad x^\alpha=o(a^x).
    \label{eq:logarithm-power-exponential-orders}
\end{equation}
即对数增长慢于任意正幂，任意正幂增长慢于底数大于 $1$ 的指数函数。
\end{proposition}

\begin{proof}
先证第二式，取整数 $m>\alpha$，令 $n=\lfloor x\rfloor$、$h=a-1>0$。当 $n\geq2m$ 时，由二项式定理，
\[
    a^x\geq a^n=(1+h)^n\geq\binom nm h^m
    \geq\frac{h^m}{2^m m!}n^m.
\]
又 $x<n+1\leq2n$，所以
\[
    0<\frac{x^\alpha}{a^x}
    \leq\frac{2^{m+\alpha}m!}{h^m}n^{\alpha-m}\to0.
\]
对于第一式，令 $t=\ln x\to+\infty$，则
\[
    \frac{\ln x}{x^\alpha}=\frac{t}{(e^\alpha)^t}\to0,
\]
其中使用已经证明的幂与指数的比较，且 $e^\alpha>1$。
\end{proof}

\begin{theorem}[常用等价无穷小]
\label{thm:standard-equivalent-infinitesimals}
当 $x\to0$ 时，有
\begin{equation}
    \begin{aligned}
        \sin x&\sim x,&\tan x&\sim x,&1-\cos x&\sim\frac{x^2}{2},\\
        \arcsin x&\sim x,&\arctan x&\sim x,&\ln(1+x)&\sim x,\\
        e^x-1&\sim x,&a^x-1&\sim x\ln a,&(1+x)^\alpha-1&\sim\alpha x.
    \end{aligned}
    \label{eq:standard-equivalent-infinitesimals}
\end{equation}
其中 $a>0$、$a\neq1$，$\alpha\in\mathbb{R}\setminus\{0\}$；反三角函数取通常的主值分支，对数与实幂均在零点附近取其有定义的范围。
\end{theorem}

\begin{proof}
$\sin x\sim x$ 是第一个重要极限；由 $\cos x\to1$，得到 $\tan x/x=(\sin x/x)/\cos x\to1$。又
\[
    \frac{1-\cos x}{x^2/2}
    =\left(\frac{\sin(x/2)}{x/2}\right)^2\to1.
\]
令 $t=\arcsin x$，由反函数连续性知 $t\to0$，且 $x=\sin t$，故 $\arcsin x/x=t/\sin t\to1$；反正切的结论同理用 $x=\tan t$ 得到。

对第二个重要极限取连续的自然对数，得到
\[
    \frac{\ln(1+x)}x
    =\ln\bigl((1+x)^{1/x}\bigr)\to\ln e=1.
\]
令 $t=e^x-1$，则 $t\to0$ 且 $x=\ln(1+t)$，从而 $(e^x-1)/x=t/\ln(1+t)\to1$。再分别将 $x\ln a$ 和 $\alpha\ln(1+x)$ 代入指数等价式，得到
\[
    \frac{a^x-1}{x\ln a}\to1,
    \qquad
    \frac{(1+x)^\alpha-1}{\alpha x}
    =\frac{e^{\alpha\ln(1+x)}-1}{\alpha\ln(1+x)}
      \frac{\ln(1+x)}x\to1.
\]
参数非零保证所写比较量在去心邻域内不为零。
\end{proof}

结合式~\eqref{eq:asymptotic-equivalence-remainder}，这些等价式也可写成带余量的表达，例如
\[
    \sin x=x+o(x),\qquad
    \cos x=1-\frac{x^2}{2}+o(x^2),\qquad
    e^x=1+x+o(x),\qquad
    (1+x)^\alpha=1+\alpha x+o(x).
\]
最后一个展开在 $\alpha=0$ 时也成立，但此时不能把零函数写成与 $\alpha x$ 等价，因为比较的分母恒为零。所有这些式子都来自本章的极限与连续性，不需要提前使用泰勒公式。

\begin{theorem}[等价因子的替换]
\label{thm:replacement-of-equivalent-factors}
在同一极限过程中，设 $f\sim\widetilde f$、$g\sim\widetilde g$，所有函数在相应去心邻域内有定义，且商中的分母恒不为零。若右端极限存在，则
\begin{equation}
    \lim f(x)h(x)=\lim\widetilde f(x)h(x),
    \qquad
    \lim\frac{f(x)}{g(x)}=\lim\frac{\widetilde f(x)}{\widetilde g(x)}.
    \label{eq:replacement-of-equivalent-factors}
\end{equation}
这些结论对有限极限以及 $+\infty,-\infty$ 和绝对值趋于无穷的情形均成立。
\end{theorem}

\begin{proof}
由等价关系，$f/\widetilde f\to1$、$\widetilde g/g\to1$，因而
\[
    fh=\frac f{\widetilde f}\,(\widetilde f h),
    \qquad
    \frac fg=\frac f{\widetilde f}\frac{\widetilde g}g
        \frac{\widetilde f}{\widetilde g}.
\]
右端中额外的因子都趋于 $1$，有限极限情形由乘积法则成立。对无穷极限，这些因子最终为正，且可以限制在任意接近 $1$ 的正数区间内，所以也保持原来的无穷趋向。
\end{proof}

\begin{example}
\label{ex:equivalent-infinitesimal-calculations}
求 $\ln\cos2x$ 在 $x\to0$ 时形如 $cx^\alpha$ 的等价无穷小，并计算
\[
    \lim_{x\to0}\frac{\tan x-\sin x}{x^3}.
\]
\end{example}

\begin{solve}
在零点附近 $\cos2x>0$，且 $\cos2x=1-2\sin^2x$，故由 $\ln(1+t)\sim t$ 与 $\sin x\sim x$，
\[
    \ln\cos2x=\ln(1-2\sin^2x)\sim-2\sin^2x\sim-2x^2.
\]
第二问先作精确变形，使可能相消的差变成乘积：
\[
    \tan x-\sin x=\frac{\sin x(1-\cos x)}{\cos x}.
\]
再对乘积中的因子作等价替换，得到
\[
    \frac{\tan x-\sin x}{x^3}
    =\frac{\sin x}{x}\frac{1-\cos x}{x^2}\frac1{\cos x}
    \to1\cdot\frac12\cdot1=\frac12.
    \qedhere
\]
\end{solve}

\emph[caution]{等价替换可以用于乘除中的因子，不能无条件用于和差中的各项。}虽然 $\tan x\sim x$ 且 $\sin x\sim x$，但把上一例的分子直接替成 $x-x=0$ 就会丢失决定结果的三阶小量。等价只保证每项的相对误差趋于零；当主项相消时，这些误差可能与剩余量同阶甚至更大。处理和差应先恒等变形，或把余项写明并核对其相对于整个和差的阶。

\section*{习题三}
\addcontentsline{toc}{section}{习题三}
\fi

\selectchapter{4}
\ifcompilechapter
\chapter{导数与微分}
\label{chap:derivatives-and-differentials}

% 依据教材第 134--170 页；分章 PDF 第 1--37 页的正文，不含习题四。
上一章用极限描述函数值的趋近过程，本章进一步研究函数在一点附近如何变化。导数刻画局部变化率，微分把这一变化写成线性部分与高阶余项之和；二者在一元函数中等价，却侧重不同的问题。以下先从速度与切线引出导数，再建立各类表示下的求导法则，最后讨论微分近似和高阶运算。本章未特别说明时，可导均指\emph[caution]{有限导数存在}，自变量的增量可以取正值，也可以取负值。

\section{导数}
\label{sec:derivatives}

\subsection{导数概念的引入}
\label{subsec:introduction-to-derivatives}

设物体沿一条直线运动，用带方向的位置函数 $s=s(t)$ 描述其运动。在时刻 $t_0$ 与 $t_0+\Delta t$ 之间，平均速度为
\[
    \frac{\Delta s}{\Delta t}
    =\frac{s(t_0+\Delta t)-s(t_0)}{\Delta t},\qquad \Delta t\ne0.
\]
若此比值在 $\Delta t\to0$ 时趋于有限值，就用该极限定义时刻 $t_0$ 的\emph{瞬时速度}（instantaneous velocity）：
\[
    v(t_0)=\lim_{\Delta t\to0}
    \frac{s(t_0+\Delta t)-s(t_0)}{\Delta t}.
\]
位置采用带方向的坐标，因此速度也带符号，其绝对值才是速率。若物体始终朝同一方向运动，也可以用所走路程来表示上述变化。这里取极限的理由是：一个固定时间段的平均量无法刻画某个时刻的局部运动状态，必须让时间段缩短到该时刻附近。

几何中的切线问题具有相同的数学结构。对曲线 $y=f(x)$ 上的点 $P=(x_0,f(x_0))$，取相邻点 $Q=(x_0+h,f(x_0+h))$，割线 $PQ$ 的斜率为
\[
    \frac{f(x_0+h)-f(x_0)}h.
\]
若它在 $h\to0$ 时趋于有限值 $k$，则将过 $P$、斜率为 $k$ 的直线定义为曲线在 $P$ 处的切线，其方程为
\begin{equation}
\label{eq:tangent-line-from-derivative}
    y-f(x_0)=k(x-x_0).
\end{equation}
这一切线概念描述割线的极限位置，不要求切线只与曲线相交一次，也不要求曲线位于切线的一侧。例如 $y=x^3$ 在原点的切线为 $y=0$，曲线在原点穿过切线。瞬时速度与切线斜率虽然背景不同，都归结为\emph{函数增量与自变量增量之比的极限}；密度、增长率等局部量也可由同一结构刻画。

\subsection{导数的定义}
\label{subsec:definition-of-derivatives}

\begin{definition}
\label{def:derivative-at-a-point}
设函数 $f$ 在 $x_0$ 的某邻域内有定义。若有限极限
\begin{equation}
\label{eq:definition-of-derivative}
    f'(x_0)=\lim_{h\to0}\frac{f(x_0+h)-f(x_0)}h
    =\lim_{x\to x_0}\frac{f(x)-f(x_0)}{x-x_0}
\end{equation}
存在，则称 $f$ 在 $x_0$ \emph{可导}（differentiable），称 $f'(x_0)$ 为 $f$ 在 $x_0$ 的\emph{导数}（derivative）。若 $y=f(x)$，还可记为
\[
    \left.\frac{\mathrm dy}{\mathrm dx}\right|_{x=x_0},\qquad
    \left.\frac{\mathrm df}{\mathrm dx}\right|_{x=x_0}.
\]
若 $f$ 在开区间 $I$ 的每一点均可导，则对应 $x\mapsto f'(x)$ 构成 $I$ 上的\emph{导函数}（derivative function），通常也简称导数。
\end{definition}

在一点的导数 $f'(x_0)$ 是一个数，导函数 $f'$ 才是函数。差商必须保留中心点的函数值 $f(x_0)$；改变这一个值便可能改变可导性。式~\eqref{eq:definition-of-derivative} 中 $h$ 的两侧趋近也不可省略，差商趋于 $+\infty$ 或 $-\infty$ 不属于本章定义的可导情形。几何上，此时可能出现竖直切线，但不能把无穷大当作有限斜率代入式~\eqref{eq:tangent-line-from-derivative}。

\begin{example}
\label{ex:derivatives-of-constant-and-square-root}
用定义求常函数 $f(x)=C$ 的导数，以及 $g(x)=\sqrt{x}$ 在 $x_0>0$ 处的导数。
\end{example}
\begin{solve}
常函数的差商恒为零，故 $f'(x)=0$。对平方根函数，利用有理化得
\[
    \frac{\sqrt{x_0+h}-\sqrt{x_0}}h
    =\frac1{\sqrt{x_0+h}+\sqrt{x_0}}
    \longrightarrow\frac1{2\sqrt{x_0}},
\]
所以 $g'(x_0)=1/(2\sqrt{x_0})$。这里 $x_0>0$ 保证充分小的正、负增量都在定义域内；端点 $0$ 处只能先讨论右侧差商，其值为 $1/\sqrt{h}\to+\infty$，并无有限右导数。
\end{solve}

基本初等函数的导数可直接由上一章的基本极限推出。对任意实数 $\alpha$，在 $x>0$ 时令 $t=h/x$，便有
\[
    \frac{(x+h)^\alpha-x^\alpha}h
    =x^{\alpha-1}\frac{(1+t)^\alpha-1}t
    \longrightarrow\alpha x^{\alpha-1}.
\]
同样，利用指数、对数的等价无穷小，有
\[
    \frac{a^{x+h}-a^x}h=a^x\frac{a^h-1}h
    \longrightarrow a^x\ln a\quad(a>0),
\]
以及
\[
    \frac{\log_a(x+h)-\log_a x}h
    =\frac1{x\ln a}\frac{\ln(1+h/x)}{h/x}
    \longrightarrow\frac1{x\ln a}
    \quad(x>0,\ a>0,\ a\ne1).
\]
对三角函数，用和差化积及 $\sin t/t\to1$，得到
\[
    \frac{\sin(x+h)-\sin x}h
    =\cos\left(x+\frac h2\right)\frac{\sin(h/2)}{h/2}
    \longrightarrow\cos x,
\]
余弦的计算类似。因此可集中记为
\begin{equation}
\label{eq:basic-power-exponential-trigonometric-derivatives}
\begin{gathered}
    (x^\alpha)'=\alpha x^{\alpha-1}\quad(x>0),\qquad
    (a^x)'=a^x\ln a\quad(a>0),\qquad (e^x)'=e^x,\\
    (\log_a x)'=\frac1{x\ln a},\qquad (\ln x)'=\frac1x\quad(x>0),\\
    (\sin x)'=\cos x,\qquad (\cos x)'=-\sin x.
\end{gathered}
\end{equation}
其中对数底数仍须满足 $a>0$ 且 $a\ne1$。幂函数公式在负半轴或原点能否使用，须另核对幂的定义域及可导性；例如非负整数次幂在全实轴上可导，不能因此把任意实指数的定义域也扩展到全实轴。

\subsection{单侧导数}
\label{subsec:one-sided-derivatives}

\begin{definition}
\label{def:one-sided-derivatives}
设 $f$ 在 $[x_0,x_0+\delta)$ 上有定义。若有限极限
\[
    f'_+(x_0)=\lim_{h\to0+}\frac{f(x_0+h)-f(x_0)}h
\]
存在，则称其为 $f$ 在 $x_0$ 的\emph{右导数}（right derivative）。若 $f$ 在 $(x_0-\delta,x_0]$ 上有定义，则以 $h\to0-$ 定义\emph{左导数}（left derivative）$f'_-(x_0)$。函数在闭区间上可导，通常指内点均可导且两个端点具有相应的有限单侧导数。
\end{definition}

\begin{proposition}
\label{prop:derivative-one-sided-and-continuity}
若 $f$ 在 $x_0$ 的邻域内有定义，则它在 $x_0$ 可导，当且仅当 $f'_-(x_0)$、$f'_+(x_0)$ 均存在且相等。可导必连续；左可导必左连续，右可导必右连续。
\end{proposition}
\begin{proof}
第一项是双侧极限存在的单侧刻画。若 $f'(x_0)$ 存在，则
\[
    f(x_0+h)-f(x_0)
    =h\,\frac{f(x_0+h)-f(x_0)}h\longrightarrow0,
\]
故 $f$ 在 $x_0$ 连续。只让 $h$ 从一侧趋于零，就得到相应的单侧连续性。
\end{proof}

\begin{example}
\label{ex:absolute-value-and-oscillatory-differentiability}
比较 $f(x)=|x|$ 及由振荡函数
\[
    S(x)=\begin{cases}\sin(1/x),&x\ne0,\\0,&x=0\end{cases}
\]
构造的 $S(x)$、$xS(x)$、$x^2S(x)$ 在原点的连续性与可导性。
\end{example}
\begin{solve}
对 $|x|$，有 $f'_+(0)=1$、$f'_-(0)=-1$，故其在原点连续却不可导。对 $S$，取趋于零且使 $\sin(1/x)$ 分别等于 $1$、$-1$ 的两列点，便知它在原点没有极限，是第二类间断。由 $|xS(x)|\leq|x|$，$xS(x)$ 在原点连续，但它的差商为 $S(h)$，即使单侧极限也不存在，因而不可导。再乘一个 $x$ 后，$x^2S(x)$ 的差商为 $h\sin(1/h)\to0$，故它在原点可导且导数为零。
\end{solve}

可导性要求函数增量在一阶尺度上稳定；仅有连续性只要求增量趋于零，并不控制它与自变量增量的比值。上述 $x^2S(x)$ 还说明切线可以在任意小邻域内与曲线反复相交：原点处的切线是 $y=0$，而 $x=1/(k\pi)$ 给出趋向原点的一列非零交点。后面求得其非零点处的导数后，还会看到\emph[caution]{可导也不意味着导函数连续}。

\section{求导数的方法}
\label{sec:differentiation-methods}

通过定义推导基本函数的导数以后，应建立与函数构造方式对应的运算法则。四则运算、反函数与复合函数的法则构成基础；隐式、参数式和极坐标式求导则是这些法则在不同表示下的应用。运算前须先确认函数的定义域、所用表示对应的局部分支以及分母非零等条件。

\subsection{函数四则运算的导数}
\label{subsec:arithmetic-differentiation-rules}

\begin{theorem}
\label{thm:arithmetic-differentiation-rules}
设 $f,g$ 在 $x_0$ 可导，$a,b$ 为常数，则 $af+bg$、$fg$ 在 $x_0$ 可导，且
\begin{align}
    (af+bg)'(x_0)&=af'(x_0)+bg'(x_0),\label{eq:derivative-linearity}\\
    (fg)'(x_0)&=f'(x_0)g(x_0)+f(x_0)g'(x_0).\label{eq:derivative-product-rule}
\end{align}
若另有 $g(x_0)\ne0$，则 $f/g$ 在 $x_0$ 的某邻域内有定义并可导，且
\begin{equation}
\label{eq:derivative-quotient-rule}
    \left(\frac fg\right)'(x_0)
    =\frac{f'(x_0)g(x_0)-f(x_0)g'(x_0)}{g(x_0)^2}.
\end{equation}
\end{theorem}
\begin{proof}
线性法则直接来自差商的线性运算。对乘积，写成
\[
    \frac{f(x_0+h)g(x_0+h)-f(x_0)g(x_0)}h
    =\frac{f(x_0+h)-f(x_0)}h\,g(x_0+h)
    +f(x_0)\frac{g(x_0+h)-g(x_0)}h.
\]
由可导必连续，$g(x_0+h)\to g(x_0)$，取极限即得乘积法则。若 $g(x_0)\ne0$，连续性保证充分靠近 $x_0$ 时 $g$ 仍非零，并且
\[
    \frac{1/g(x_0+h)-1/g(x_0)}h
    =-\frac1{g(x_0+h)g(x_0)}
    \frac{g(x_0+h)-g(x_0)}h
    \longrightarrow-\frac{g'(x_0)}{g(x_0)^2}.
\]
将 $f/g$ 写成 $f\cdot(1/g)$，再用乘积法则，得到商法则。
\end{proof}

有限个可导函数的线性组合可逐项求导，有限乘积则每次选取一个因子求导后相加：
\[
    \left(\prod_{j=1}^m f_j\right)'
    =\sum_{k=1}^m f'_k\prod_{\substack{1\leq j\leq m\\j\ne k}}f_j.
\]
这里强调有限性，无穷项的求和与求导能否交换，需要后续另加条件。

\begin{example}
\label{ex:derivatives-of-tangent-and-cotangent}
求 $\tan x$、$\cot x$ 的导数。
\end{example}
\begin{solve}
在 $\cos x\ne0$ 处，由商法则，
\[
    (\tan x)'=\left(\frac{\sin x}{\cos x}\right)'
    =\frac{\cos^2x+\sin^2x}{\cos^2x}=\frac1{\cos^2x}.
\]
同理，在 $\sin x\ne0$ 处，$(\cot x)'=-1/\sin^2x$。这些条件正是原函数有定义的条件；不能在三角函数的极点处套用求导公式。
\end{solve}

\subsection{反函数的求导法则}
\label{subsec:inverse-function-rule}

\begin{theorem}
\label{thm:inverse-function-derivative}
设 $f$ 在开区间 $I$ 上连续且严格单调，$J=f(I)$，反函数记为 $g=f^{-1}$。若 $x_0\in I$，$y_0=f(x_0)$，且 $f'(x_0)$ 存在并非零，则 $g$ 在 $y_0$ 可导，且
\begin{equation}
\label{eq:inverse-function-derivative}
    g'(y_0)=\frac1{f'(x_0)}=\frac1{f'(g(y_0))}.
\end{equation}
特别地，若 $f$ 在 $I$ 上处处可导且导数非零，则其反函数在 $J$ 上处处可导。
\end{theorem}
\begin{proof}
由\crosschapterref{3}{定理~\ref{thm:continuity-of-inverse-functions}}{反函数的连续性定理}，反函数连续，故 $y\to y_0$ 时 $x=g(y)\to x_0$。又因单射性，$y\ne y_0$ 时 $x\ne x_0$，故
\[
    \frac{g(y)-g(y_0)}{y-y_0}
    =\frac{x-x_0}{f(x)-f(x_0)}
    \longrightarrow\frac1{f'(x_0)}.
\]
分母极限非零正是 $f'(x_0)\ne0$ 的作用。
\end{proof}

记成 $\mathrm dx/\mathrm dy=1/(\mathrm dy/\mathrm dx)$ 时，两边的导数必须取在同一对应点。反函数存在并不足以保证它可导，例如 $y=x^3$ 在原点的导数为零，其反函数 $x=\sqrt[3]{y}$ 在原点没有有限导数。

\begin{example}
\label{ex:derivatives-of-inverse-trigonometric-functions}
求四个常用反三角函数的导数，并说明定义分支与适用区间。
\end{example}
\begin{solve}
若 $y=\arctan x\in(-\pi/2,\pi/2)$，则 $x=\tan y$，从而
\[
    \frac{\mathrm dy}{\mathrm dx}
    =\frac1{1+\tan^2y}=\frac1{1+x^2}.
\]
若 $y=\arcsin x\in[-\pi/2,\pi/2]$，则在 $|x|<1$ 时 $\cos y>0$，因此
\[
    \frac{\mathrm dy}{\mathrm dx}
    =\frac1{\cos y}=\frac1{\sqrt{1-x^2}}.
\]
相同方法用于值域为 $[0,\pi]$ 的 $\arccos$，以及值域为 $(0,\pi)$ 的 $\operatorname{arccot}$，给出
\begin{equation}
\label{eq:inverse-trigonometric-derivatives}
\begin{gathered}
    (\arcsin x)'=\frac1{\sqrt{1-x^2}},\qquad
    (\arccos x)'=-\frac1{\sqrt{1-x^2}}\quad(|x|<1),\\
    (\arctan x)'=\frac1{1+x^2},\qquad
    (\operatorname{arccot}x)'=-\frac1{1+x^2}\quad(x\in\mathbb R).
\end{gathered}
\end{equation}
反正弦与反余弦在 $\pm1$ 处有定义，但相应的单侧差商无有限极限，因此公式中的区间不能改成闭区间。
\end{solve}

\subsection{复合函数的求导法则}
\label{subsec:chain-rule}

\begin{theorem}[链式法则]
\label{thm:one-variable-chain-rule}
设 $g$ 在 $x_0$ 的邻域内有定义且在 $x_0$ 可导，$u_0=g(x_0)$；设 $f$ 在 $u_0$ 的邻域内有定义且在 $u_0$ 可导，并且复合函数 $F=f\circ g$ 在 $x_0$ 的某邻域内有定义。则 $F$ 在 $x_0$ 可导，且
\begin{equation}
\label{eq:one-variable-chain-rule}
    F'(x_0)=f'(g(x_0))g'(x_0).
\end{equation}
这称为\emph{链式法则}（chain rule），亦称链锁法则。若 $y=f(u)$、$u=g(x)$，则常记成
\[
    \frac{\mathrm dy}{\mathrm dx}
    =\frac{\mathrm dy}{\mathrm du}\frac{\mathrm du}{\mathrm dx}.
\]
\end{theorem}
\begin{proof}
令
\[
    A(u)=\begin{cases}
        \dfrac{f(u)-f(u_0)}{u-u_0},&u\ne u_0,\\[4pt]
        f'(u_0),&u=u_0.
    \end{cases}
\]
由可导性，$A$ 在 $u_0$ 连续，且对所有邻近 $u$ 都有 $f(u)-f(u_0)=A(u)(u-u_0)$。于是
\[
    \frac{F(x)-F(x_0)}{x-x_0}
    =A(g(x))\frac{g(x)-g(x_0)}{x-x_0}.
\]
因 $g$ 在 $x_0$ 连续，$A(g(x))\to f'(u_0)$，取极限即可。
\end{proof}

证明中把差商在 $u=u_0$ 处连续补值，是为了容许 $g(x)=g(x_0)$ 在邻近点反复发生。若直接除以 $g(x)-g(x_0)$，就可能在这些点除以零，因此不能未经说明地把复合函数的差商拆成两个普通商的乘积。多层复合可逐层使用链式法则，顺序是外层在内层函数值处求导，再乘内层的导数。

\begin{example}
\label{ex:chain-rule-and-logarithm-derivatives}
比较 $\sin(x^2)$ 与 $\sin^2x$ 的导数，并求 $\ln|x|$ 的导数。
\end{example}
\begin{solve}
前两者的外层函数分别为正弦与平方，故
\[
    (\sin(x^2))'=2x\cos(x^2),\qquad
    (\sin^2x)'=2\sin x\cos x=\sin2x.
\]
对 $\ln|x|$，在 $x>0$ 时导数为 $1/x$；在 $x<0$ 时写成 $\ln(-x)$，由链式法则仍得 $(-1)/(-x)=1/x$。因此
\[
    (\ln|x|)'=\frac1x\qquad(x\ne0).
\]
同理，由 $x^\alpha=e^{\alpha\ln x}$ 可重新得到 $x>0$ 上的幂函数求导公式，这也说明同一函数可以用不同的分解方式求导。
\end{solve}

\begin{example}
\label{ex:binomial-square-sum-by-differentiation}
设 $n\geq1$ 为整数，利用求导证明
\[
    \sum_{k=0}^n\binom nk k^2=n(n+1)2^{n-2}.
\]
\end{example}
\begin{solve}
对有限和恒等式 $(1+x)^n=\sum_{k=0}^n\binom nk x^k$ 求导，再乘 $x$ 后求导，得到
\[
    n(1+x)^{n-1}+n(n-1)x(1+x)^{n-2}
    =\sum_{k=1}^n\binom nk k^2x^{k-1}.
\]
当 $n\geq2$ 时令 $x=1$，左边为 $n2^{n-1}+n(n-1)2^{n-2}=n(n+1)2^{n-2}$。$n=1$ 时直接验得两边均为 $1$。这里先乘 $x$ 再求导，是为了使单项式的系数连续获得两次因子 $k$。
\end{solve}

\subsection{隐函数的求导法}
\label{subsec:implicit-differentiation}

若变量 $x,y$ 通过方程 $F(x,y)=0$ 联系，且在某个指定范围及分支上，每个 $x$ 唯一对应一个 $y=y(x)$，则称这个函数为方程所确定的\emph{隐函数}（implicit function）。方程本身未必在整个范围内确定一个函数，例如 $x^2+y^2=1$ 对 $|x|<1$ 同时给出正、负两个分支。\emph[caution]{对方程形式求导不能代替隐函数存在性与可导性的证明}；除可直接化为反函数等情形外，本节先假定所选分支已经可导，一般的隐函数存在定理留待多元微分学讨论。

隐函数求导的操作是：在恒等式 $F(x,y(x))=0$ 两边对 $x$ 求导，凡含 $y$ 的部分均用链式法则，再从所得等式解出 $y'(x)$。解出的公式只能在推导中所除因子非零的点使用；无法解出时，应重新判断分支、改用 $x=x(y)$，或检查曲线是否存在竖直切线或奇点。

\begin{example}
\label{ex:kepler-equation-implicit-derivative}
设 $0<e<1$ 为常数，证明\emph{开普勒方程}（Kepler's equation）
\[
    y-x-e\sin y=0
\]
在全实轴上确定唯一可导函数 $y=y(x)$，并求其导数。
\end{example}
\begin{solve}
令 $\varphi(y)=y-e\sin y$。它连续，且因 $e\sin y$ 有界，$\varphi(y)\to\pm\infty$ 当 $y\to\pm\infty$。对 $y_2>y_1$，由 $|\sin y_2-\sin y_1|\leq y_2-y_1$，有
\[
    \varphi(y_2)-\varphi(y_1)
    \geq(1-e)(y_2-y_1)>0.
\]
所以 $\varphi$ 严格递增，由介值定理，它的值域为 $\mathbb R$。于是方程唯一确定 $y=\varphi^{-1}(x)$。又有 $\varphi'(y)=1-e\cos y\geq1-e>0$，可用反函数求导定理保证 $y$ 可导。对方程求导得
\[
    y'-1-e\cos y\,y'=0,
    \qquad y'(x)=\frac1{1-e\cos y(x)}.
\]
虽然未得到 $y(x)$ 的显式表达式，导数却已由该方程及其唯一分支确定。
\end{solve}

\begin{example}
\label{ex:circle-tangents-by-implicit-differentiation}
求单位圆 $x^2+y^2=1$ 上斜率为 $-1/2$ 的切线。
\end{example}
\begin{solve}
在 $y\ne0$ 的上、下半圆分支上，$2x+2yy'=0$，故 $y'=-x/y$。设切点为 $(x_0,y_0)$，则 $y_0=2x_0$，结合圆的方程得到
\[
    (x_0,y_0)=\left(\frac1{\sqrt5},\frac2{\sqrt5}\right)
    \quad\text{或}\quad
    \left(-\frac1{\sqrt5},-\frac2{\sqrt5}\right).
\]
将两点分别代入切线方程，得 $y=-x/2+\sqrt5/2$ 与 $y=-x/2-\sqrt5/2$。圆上 $y=0$ 的两点不能使用 $y'=-x/y$，它们的切线分别为 $x=1$、$x=-1$。
\end{solve}

对乘积、商或变指数幂函数，先取对数往往比直接展开简洁，这称为\emph{对数求导法}（logarithmic differentiation）。若 $y=u(x)^{v(x)}$，其中 $u>0$ 且 $u,v$ 均可导，则 $\ln y=v\ln u$，从而
\begin{equation}
\label{eq:logarithmic-differentiation-variable-power}
    y'=u^v\left(v'\ln u+v\frac{u'}u\right).
\end{equation}
这里 $u>0$ 使任意实指数的幂及对数都有定义。对其他非零函数的乘积或商，可以在其非零邻域内使用 $\ln|y|$，而不能在零点直接除以 $y$。

\begin{example}
\label{ex:logarithmic-differentiation-examples}
求 $y=x^{e^x}$（$x>0$）及 $z=(1+x^{-2})^{x^2}$（$x\ne0$）的导数。
\end{example}
\begin{solve}
第一式取对数得 $\ln y=e^x\ln x$，所以
\[
    y'=x^{e^x}e^x\left(\ln x+\frac1x\right).
\]
第二式有 $\ln z=x^2\ln(1+x^{-2})$，故
\[
    \frac{z'}z
    =2x\ln(1+x^{-2})-\frac{2x}{1+x^2},
\]
即
\[
    z'=2x(1+x^{-2})^{x^2}
    \left(\ln(1+x^{-2})-\frac1{1+x^2}\right).
\]
两个例子中底数与指数均随 $x$ 改变，不能只套用常指数幂函数或常底数指数函数的求导公式。
\end{solve}

\subsection{参数式函数的求导法}
\label{subsec:parametric-differentiation}

设曲线由\emph{参数方程}（parametric equations）$x=\varphi(t)$、$y=\psi(t)$ 给出，$t_0$ 对应点 $(x_0,y_0)$。为保证它在这一点附近确实确定可导函数 $y=y(x)$，可假设 $\varphi,\psi$ 在 $t_0$ 的某邻域内可导，$\varphi$ 在该邻域内严格单调，且 $\varphi'(t_0)\ne0$。反函数与链式法则于是给出
\begin{equation}
\label{eq:parametric-first-derivative}
    \left.\frac{\mathrm dy}{\mathrm dx}\right|_{x=x_0}
    =\frac{\psi'(t_0)}{\varphi'(t_0)}.
\end{equation}
其依据是 $y(x)=\psi(\varphi^{-1}(x))$，而不是把两个无穷小增量形式地约掉。一个常用的充分条件是 $\varphi$ 在 $t_0$ 附近连续可导且 $\varphi'(t_0)\ne0$：连续性使导数在充分小邻域同号，再由下一章的\crosschapterref{5}{定理~\ref{thm:derivative-monotonicity-test}}{单调性的导数判别定理} 得局部严格单调。

若 $\varphi'(t_0)=0$ 而 $\psi'(t_0)\ne0$，可在相应局部可逆条件下改以 $y$ 为自变量，此时曲线的切线竖直。若二者同时为零，上述公式没有结论，需要直接考察曲线；例如 $x=t^3,y=t^3$ 在 $t=0$ 的两个参数导数均为零，但曲线仍是直线 $y=x$。

\begin{example}
\label{ex:ellipse-tangent-parametric-and-implicit}
设 $a,b>0$，求椭圆 $x^2/a^2+y^2/b^2=1$ 在任意点 $(x_0,y_0)$ 处的切线方程。
\end{example}
\begin{solve}
当 $y_0\ne0$ 时，椭圆在切点附近选定的上、下分支可导，由隐函数求导得
\[
    \frac{2x}{a^2}+\frac{2yy'}{b^2}=0,
    \qquad y'=-\frac{b^2x}{a^2y}.
\]
在切点代入直线方程，并利用 $x_0^2/a^2+y_0^2/b^2=1$，得到
\begin{equation}
\label{eq:ellipse-tangent-line}
    \frac{x_0x}{a^2}+\frac{y_0y}{b^2}=1.
\end{equation}
也可用参数式 $x=a\cos t,y=b\sin t$，在 $\sin t_0\ne0$ 时求得 $\mathrm dy/\mathrm dx=-b\cos t/(a\sin t)$，化简后得到同一方程。当 $y_0=0$ 时，$x_0=\pm a$，可将 $x$ 看成 $y$ 的函数，求得 $\mathrm dx/\mathrm dy=0$，切线为 $x=\pm a$；这两条竖直切线也由式~\eqref{eq:ellipse-tangent-line} 统一表达。
\end{solve}

\subsection{极坐标式函数的求导法}
\label{subsec:polar-differentiation}

若曲线由\emph{极坐标方程}（polar equation）$r=r(\theta)$ 给出，其中 $r$ 为极径、$\theta$ 为极角，则可转为参数式
\[
    x=r(\theta)\cos\theta,\qquad y=r(\theta)\sin\theta.
\]
若 $r$ 在所考虑点附近连续可导，且 $r'\cos\theta-r\sin\theta\ne0$，则参数式求导公式给出
\begin{equation}
\label{eq:polar-curve-derivative}
    \frac{\mathrm dy}{\mathrm dx}
    =\frac{r'(\theta)\sin\theta+r(\theta)\cos\theta}
    {r'(\theta)\cos\theta-r(\theta)\sin\theta}.
\end{equation}
分母为零而分子非零时切线竖直。事实上，
\[
    \left(\frac{\mathrm dx}{\mathrm d\theta}\right)^2
    +\left(\frac{\mathrm dy}{\mathrm d\theta}\right)^2
    =(r')^2+r^2,
\]
因此只要 $r,r'$ 不同时为零，参数曲线就有非零的切向方向；若二者同时为零，应另作分析。

当 $r>0$ 时，把切向方向相对于向径分解，其沿向径和垂直向径的分量分别为 $r'$ 与 $r$。若 $\alpha$ 表示切线与 $x$ 轴的夹角，$\beta$ 表示从向径逆时针转到切线的角，并按模 $\pi$ 选取在 $[0,\pi)$ 内，则 $\beta\equiv\alpha-\theta\pmod\pi$，且在 $r'\ne0$ 时
\begin{equation}
\label{eq:polar-radius-tangent-angle}
    \tan\beta=\frac r{r'}.
\end{equation}
当 $r'=0$ 而 $r>0$ 时，切线与向径垂直。用正切求角只能先得到模 $\pi$ 的关系，最终角度必须结合取值区间判断。

\begin{example}
\label{ex:lemniscate-radius-tangent-angle}
设 $a>0$，双纽线在 $0<\theta<\pi/4$ 的分支由 $r^2=a^2\cos2\theta$、$r>0$ 给出，求向径与切线间的上述夹角 $\beta$。
\end{example}
\begin{solve}
对极角求导，得 $2rr'=-2a^2\sin2\theta$，从而
\[
    \tan\beta=\frac r{r'}
    =-\frac{r^2}{a^2\sin2\theta}
    =-\cot2\theta
    =\tan\left(2\theta+\frac\pi2\right).
\]
因 $2\theta+\pi/2\in(\pi/2,\pi)$，与所选角度范围一致，故 $\beta=2\theta+\pi/2$。
\end{solve}

\section{微分}
\label{sec:differentials}

\subsection{微分的定义}
\label{subsec:definition-of-differentials}

导数由差商极限定义，微分则直接从函数增量中提取线性部分。例如圆面积 $S(r)=\pi r^2$ 满足
\[
    S(r+h)-S(r)=2\pi rh+\pi h^2.
\]
当 $h\to0$ 时，第一项关于 $h$ 线性，第二项为 $o(|h|)$。这使我们能够用容易计算的线性项近似真实增量，同时用余项描述误差尺度。

\begin{definition}
\label{def:one-variable-differential}
设 $f$ 在 $x_0$ 的某邻域内有定义。若存在常数 $A$，使
\begin{equation}
\label{eq:differentiability-linear-part}
    \Delta y=f(x_0+h)-f(x_0)=Ah+o(|h|)\qquad(h\to0),
\end{equation}
则称 $f$ 在 $x_0$ \emph{可微}（differentiable），称线性映射 $h\mapsto Ah$ 为 $f$ 在 $x_0$ 的\emph{微分}（differential），其在增量 $h$ 上的值记为 $\mathrm df(x_0)=Ah$，或记作 $\mathrm dy=Ah$。这里 $o(|h|)$ 指余项与 $|h|$ 之比趋于零，也常简写为 $o(h)$。
\end{definition}

\begin{theorem}
\label{thm:one-variable-differentiable-iff-derivable}
一元函数 $f$ 在 $x_0$ 可微，当且仅当它在 $x_0$ 可导；微分的系数唯一，且等于 $f'(x_0)$。因此
\begin{equation}
\label{eq:one-variable-differential}
    \mathrm dy=f'(x_0)\,\mathrm dx,\qquad \mathrm dx=h.
\end{equation}
\end{theorem}
\begin{proof}
若式~\eqref{eq:differentiability-linear-part} 成立，对 $h\ne0$ 除以 $h$，余项与 $h$ 的比仍趋于零，所以差商趋于 $A$，即 $f'(x_0)=A$。反之，若导数存在，记
\[
    \varepsilon(h)=\frac{f(x_0+h)-f(x_0)}h-f'(x_0)\longrightarrow0,
\]
便有 $f(x_0+h)-f(x_0)=f'(x_0)h+h\varepsilon(h)$，最后一项为 $o(|h|)$，故 $f$ 可微。因为极限唯一，系数也唯一。
\end{proof}

对独立变量本身，恒等函数的微分是 $\mathrm dx=h=\Delta x$；对因变量，$\mathrm dy$ 一般只等于增量的线性部分，并不等于 $\Delta y$。几何上，真实增量是曲线从 $x_0$ 到 $x_0+h$ 的纵向变化，微分是切线在同一横向增量下的纵向变化，二者相差 $o(|h|)$。于是
\begin{equation}
\label{eq:differential-linear-approximation}
    f(x_0+h)=f(x_0)+f'(x_0)h+o(|h|)
\end{equation}
表达了函数在一点附近的线性近似。

\emph[caution]{微分首先保证相对于自变量增量的一阶误差趋于零，不自动保证相对于函数增量的误差趋于零。}若 $f'(x_0)\ne0$，则充分小的非零 $h$ 满足 $\Delta y\ne0$，并且
\[
    \frac{\Delta y-\mathrm dy}{\Delta y}
    =\frac{o(|h|)/h}{f'(x_0)+o(|h|)/h}\longrightarrow0.
\]
若 $f'(x_0)=0$，则线性部分为零，这一相对误差结论可能失效。例如 $f(x)=x^2$ 在原点有 $\mathrm dy=0$，而 $\Delta y=h^2$，用微分近似函数增量的相对误差恒为 $1$。

\begin{example}
\label{ex:fourth-root-linear-approximation}
利用微分估算 $\sqrt[4]{80}$。
\end{example}
\begin{solve}
取 $f(x)=x^{1/4}$、$x_0=81$、$h=-1$。因 $f(81)=3$，$f'(81)=1/(4\cdot81^{3/4})=1/108$，一阶近似给出
\[
    \sqrt[4]{80}\approx3-\frac1{108}\approx2.99074.
\]
选择 $81$ 的原因是函数值和导数都容易计算。这个数值来自一阶线性近似；若要保证具体误差上界，还须对余项作定量估计，不能只由 $o(|h|)$ 推断所写各位均准确。
\end{solve}

\begin{example}
\label{ex:square-area-relative-error-differential}
正方形面积由边长测量值计算。若希望面积的相对误差不超过 $1\%$，边长允许多大的相对误差？比较微分估计与严格保证。
\end{example}
\begin{solve}
设真实边长为 $a>0$，测量增量为 $\Delta a$，$S=a^2$。由 $\mathrm dS=2a\,\mathrm da$，一阶估计为
\[
    \frac{|\Delta S|}{S}\approx\frac{|\mathrm dS|}{S}
    =2\frac{|\Delta a|}a.
\]
因此边长相对误差约为 $0.5\%$ 时，面积相对误差约为 $1\%$。若要求对误差正负两种情形都严格保证，则令 $\varepsilon=\Delta a/a$，有精确关系
\[
    \frac{\Delta S}{S}=2\varepsilon+\varepsilon^2.
\]
当 $|\varepsilon|\leq\eta<1$ 时，其绝对值的最大值为 $2\eta+\eta^2$，故应取
\[
    \eta\leq\sqrt{1.01}-1\approx0.00498756,
\]
即约 $0.498756\%$。这与一阶估计接近，但揭示了近似精度与严格误差保证的区别。
\end{solve}

\subsection{一阶微分的形式不变性}
\label{subsec:invariance-of-first-differentials}

设 $y=f(u)$、$u=g(x)$，且 $g$ 在 $x$ 处可导、$f$ 在 $u=g(x)$ 处可导。以 $x$ 为独立变量时，由链式法则，
\[
    \mathrm dy=f'(u)g'(x)\,\mathrm dx,
    \qquad \mathrm du=g'(x)\,\mathrm dx,
\]
因此仍有
\begin{equation}
\label{eq:invariance-of-first-differential}
    \mathrm dy=f'(u)\,\mathrm du.
\end{equation}
若 $u$ 本身是独立变量，微分也是这一形式；若 $u$ 依赖于 $x$，形式仍相同。这称为\emph{一阶微分的形式不变性}（invariance of the first differential）。两种情形中 $\mathrm du$ 的含义不同：前者是自变量增量 $\Delta u$，后者是 $g'(x)\Delta x$，一般只与实际增量 $g(x+\Delta x)-g(x)$ 相差 $o(|\Delta x|)$。

形式不变性使微分运算可以按函数本身的表达式展开，而不必在每一步写出最底层自变量。若 $u,v$ 都是同一自变量的可微函数，则
\begin{equation}
\label{eq:differential-arithmetic-rules}
\begin{gathered}
    \mathrm d(au+bv)=a\,\mathrm du+b\,\mathrm dv,\qquad
    \mathrm d(uv)=v\,\mathrm du+u\,\mathrm dv,\\
    \mathrm d\left(\frac uv\right)
    =\frac{v\,\mathrm du-u\,\mathrm dv}{v^2}\quad(v\ne0),
\end{gathered}
\end{equation}
其中 $a,b$ 为常数。基本函数的微分由其导数直接得到，例如
\[
\begin{gathered}
    \mathrm d(u^\alpha)=\alpha u^{\alpha-1}\,\mathrm du\quad(u>0),\qquad
    \mathrm d(\ln|u|)=\frac{\mathrm du}u\quad(u\ne0),\\
    \mathrm d(e^u)=e^u\,\mathrm du,\qquad
    \mathrm d(\sin u)=\cos u\,\mathrm du,\qquad
    \mathrm d(\arctan u)=\frac{\mathrm du}{1+u^2}.
\end{gathered}
\]
在相应局部可逆、分母非零的条件下，从 $\mathrm dy=f'(x)\mathrm dx$ 可以得到反函数求导公式；从参数式的 $\mathrm dx=\varphi'(t)\mathrm dt$、$\mathrm dy=\psi'(t)\mathrm dt$ 也可以读出式~\eqref{eq:parametric-first-derivative}。这些记号简化了运算，成立依据仍是前面证明过的求导法则。

\begin{example}
\label{ex:differentials-of-composite-and-fractional-functions}
求 $y=\arctan(e^x)$ 的微分；再设 $|a|<1$、$|x|<1$，$z=(x-a)/(1-ax)$，证明
\[
    \frac{\mathrm dz}{1-z^2}=\frac{\mathrm dx}{1-x^2}.
\]
\end{example}
\begin{solve}
对第一式逐层微分，得
\[
    \mathrm dy=\frac1{1+e^{2x}}\,\mathrm d(e^x)
    =\frac{e^x}{1+e^{2x}}\,\mathrm dx
    =\frac{\mathrm dx}{e^x+e^{-x}}.
\]
对第二式，$1-ax>0$，商的微分给出
\[
    \mathrm dz=\frac{1-a^2}{(1-ax)^2}\,\mathrm dx,
    \qquad
    1-z^2=\frac{(1-a^2)(1-x^2)}{(1-ax)^2}>0.
\]
将二式相除即得所求恒等式，且上述条件同时保证各分母非零。
\end{solve}

\begin{example}
\label{ex:implicit-differentiation-using-differentials}
设 $x>0$，可导分支 $y=y(x)$ 满足 $\sin(y^2)=\cos\sqrt{x}$，求其导数。
\end{example}
\begin{solve}
对方程微分，得
\[
    2y\cos(y^2)\,\mathrm dy
    =-\frac{\sin\sqrt{x}}{2\sqrt{x}}\,\mathrm dx.
\]
因此在 $y\cos(y^2)\ne0$ 的点，
\[
    y'=-\frac{\sin\sqrt{x}}{4y\sqrt{x}\cos(y^2)}.
\]
此处只用了给定可导分支满足的恒等式；若某点的系数 $y\cos(y^2)$ 为零，就不能从该等式除法解出导数，也不能把 $y$ 在整个范围内一律写成同一个反正弦显式分支。
\end{solve}

\section{高阶导数与高阶微分}
\label{sec:higher-derivatives-and-differentials}

\subsection{高阶导数}
\label{subsec:higher-derivatives}

位置函数 $s(t)$ 的一阶导数是速度，速度的导数则是加速度，即 $s''(t)$。这说明变化率本身也可以作为新的研究对象，从而引出高阶导数。

\begin{definition}
\label{def:higher-order-derivatives}
若 $f$ 在开区间上可导，且导函数 $f'$ 在 $x$ 处可导，则称 $(f')'(x)$ 为 $f$ 在 $x$ 处的\emph{二阶导数}（second derivative），记为 $f''(x)$。递归地，若 $f^{(n-1)}$ 在 $x$ 的某邻域内存在且在 $x$ 处可导，则定义
\[
    f^{(n)}(x)=\bigl(f^{(n-1)}\bigr)'(x),\qquad n\geq2,
\]
称为 $n$ 阶导数，也记为 $y^{(n)}$ 或 $\mathrm d^ny/\mathrm dx^n$；约定 $f^{(0)}=f$。若在区间上直至 $n$ 阶导数都存在，则称 $f$ 在该区间 $n$ 次可导；若这些导数还连续，则称 $f$ 为 $C^n$ 函数。
\end{definition}

高阶导数必须逐阶存在，$n$ 次可导与属于 $C^n$ 不应混为一谈。例如前面构造的 $f(x)=x^2\sin(1/x)$（$x\ne0$）、$f(0)=0$ 在全实轴可导，但当 $x\ne0$ 时
\[
    f'(x)=2x\sin(1/x)-\cos(1/x),\qquad f'(0)=0,
\]
导函数在原点没有极限，因而 $f\notin C^1(\mathbb R)$；也由于 $f'$ 在原点不连续，$f''(0)$ 不存在。

\begin{proposition}
\label{prop:standard-higher-derivative-formulas}
对正整数 $n$，下列公式在所列定义域内成立：
\begin{align}
    (e^{ax})^{(n)}&=a^ne^{ax}\quad(a\in\mathbb R),\label{eq:exponential-higher-derivatives}\\
    (x^\alpha)^{(n)}&=\alpha(\alpha-1)\cdots(\alpha-n+1)x^{\alpha-n}
    \quad(x>0),\label{eq:power-higher-derivatives}\\
    (\ln(1+x))^{(n)}&=(-1)^{n-1}\frac{(n-1)!}{(1+x)^n}
    \quad(x>-1),\label{eq:logarithm-higher-derivatives}\\
    (\sin x)^{(n)}&=\sin\left(x+\frac{n\pi}2\right),\qquad
    (\cos x)^{(n)}=\cos\left(x+\frac{n\pi}2\right).
    \label{eq:trigonometric-higher-derivatives}
\end{align}
其中 $0!=1$。若 $P$ 是 $m$ 次多项式，则在整个实轴上有 $P^{(n)}\equiv0$（$n>m$）。
\end{proposition}
\begin{proof}
以上等式均以一阶公式为起点作数学归纳。指数函数每求导一次获得因子 $a$；幂函数每求导一次，系数多乘当前指数，指数减一。对对数，若公式对 $n$ 成立，则再求导得到
\[
    (-1)^{n-1}(n-1)!\,[-n(1+x)^{-n-1}]
    =(-1)^n\frac{n!}{(1+x)^{n+1}},
\]
所以对 $n+1$ 也成立。正弦每求导一次相位增加 $\pi/2$，余弦亦然，归纳即得三角函数公式。多项式是有限个非负整数次幂的线性组合，每个单项式在求导次数超过其次数后都为零，故得到最后一项。
\end{proof}

求简单函数的高阶导数时，先算若干低阶结果以发现规律，再用归纳法验证，不能把有限次观察直接当作一般证明。三角函数的高阶公式也可由欧拉公式作形式上的提示，但上述归纳证明只使用实函数求导，不依赖尚未建立的复函数微分理论。

\begin{example}
\label{ex:arctangent-higher-derivatives}
设 $y=\arctan x$，求 $y^{(n)}(x)$ 及它们在原点的值。
\end{example}
\begin{solve}
由 $x=\tan y$，$y'=\cos^2y$。前几次求导提示下式：
\begin{equation}
\label{eq:arctangent-higher-derivatives}
    y^{(n)}=(n-1)!\cos^ny\,
    \sin\left(n\left(y+\frac\pi2\right)\right),\qquad n\geq1.
\end{equation}
当 $n=1$ 时，右边为 $\cos^2y$。若对 $n$ 成立，记 $\theta=y+\pi/2$，对 $x$ 再求导得
\[
\begin{aligned}
    y^{(n+1)}
    &=n!\cos^{n+1}y\,
      [-\sin y\sin(n\theta)+\cos y\cos(n\theta)]\\
    &=n!\cos^{n+1}y\cos(n\theta+y)
     =n!\cos^{n+1}y\sin((n+1)\theta),
\end{aligned}
\]
这正是 $n+1$ 的公式。由于 $y(0)=0$，故
\[
    y^{(2m)}(0)=0\quad(m\geq1),\qquad
    y^{(2m+1)}(0)=(-1)^m(2m)!\quad(m\geq0).
\]
式~\eqref{eq:arctangent-higher-derivatives} 中的 $y$ 只是 $\arctan x$ 的简记；若需纯粹关于 $x$ 的表达式，可用 $\cos y=1/\sqrt{1+x^2}$ 代回。
\end{solve}

\subsection{莱布尼茨公式}
\label{subsec:leibniz-formula}

有限线性组合的高阶导数仍可逐项计算，即 $(af+bg)^{(n)}=af^{(n)}+bg^{(n)}$。乘积的高阶导数则涉及两因子之间的求导次数分配：总共求导 $n$ 次，其中 $k$ 次落在一个因子上，另外 $n-k$ 次落在另一个因子上。

\begin{theorem}[莱布尼茨公式]
\label{thm:leibniz-higher-product-rule}
设 $u,v$ 在开区间 $I$ 上均 $n$ 次可导，则它们的乘积也 $n$ 次可导，且
\begin{equation}
\label{eq:leibniz-higher-product-rule}
    (uv)^{(n)}(x)=\sum_{k=0}^n\binom nk
    u^{(k)}(x)v^{(n-k)}(x),\qquad x\in I,
\end{equation}
其中 $u^{(0)}=u$、$v^{(0)}=v$，$\binom nk=n!/[k!(n-k)!]$。这称为\emph{莱布尼茨公式}（Leibniz formula）。
\end{theorem}
\begin{proof}
对 $n$ 归纳。$n=1$ 时是乘积法则。假设结论对 $n$ 成立，且两函数均有 $n+1$ 阶导数，则
\[
\begin{aligned}
    (uv)^{(n+1)}
    &=\sum_{k=0}^n\binom nk u^{(k+1)}v^{(n-k)}
      +\sum_{k=0}^n\binom nk u^{(k)}v^{(n+1-k)}\\
    &=u^{(n+1)}v+uv^{(n+1)}
      +\sum_{k=1}^n
       \left(\binom n{k-1}+\binom nk\right)
       u^{(k)}v^{(n+1-k)}.
\end{aligned}
\]
利用 $\binom n{k-1}+\binom nk=\binom{n+1}k$，即得 $n+1$ 的结论。此证明只要求所讨论的有限阶导数存在，无须假定两函数具有任意阶导数。
\end{proof}

若只需要原点或其他指定点的高阶导数，直接逐阶求完整表达式常会产生大量无用项。更有效的办法是寻找函数满足的低阶恒等式，使用莱布尼茨公式求导后再代入指定点，建立递推关系。

\begin{example}
\label{ex:arcsine-derivative-recurrence-at-zero}
设 $y=\arcsin x$，求 $y^{(n)}(0)$。
\end{example}
\begin{solve}
在 $|x|<1$ 内，$y'=1/\sqrt{1-x^2}$，且各阶导数均存在。由此可得
\[
    (1-x^2)y''-xy'=0.
\]
对该式求 $n-2$ 阶导数，并用莱布尼茨公式处理二次多项式与一次多项式，得到对 $n\geq2$ 成立的关系
\[
    (1-x^2)y^{(n)}-(2n-3)xy^{(n-1)}
    -(n-2)^2y^{(n-2)}=0.
\]
令 $x=0$，得到递推式
\begin{equation}
\label{eq:arcsine-derivatives-zero-recurrence}
    y^{(n)}(0)=(n-2)^2y^{(n-2)}(0),\qquad n\geq2.
\end{equation}
再用初值 $y(0)=0$、$y'(0)=1$，可得
\[
    y^{(2m)}(0)=0\quad(m\geq0),\qquad
    y^{(2m+1)}(0)=\bigl[(2m-1)!!\bigr]^2\quad(m\geq0).
\]
这里 $(2m-1)!!=1\cdot3\cdots(2m-1)$（$m\geq1$），并约定 $(-1)!!=1$。由于多项式因子在少数几次求导后即消失，递推式比完整高阶导数表达式简洁得多。
\end{solve}

\begin{example}
\label{ex:tangent-derivatives-zero-through-eight}
设 $y=\tan x$，求 $y^{(n)}(0)$，$1\leq n\leq8$。
\end{example}
\begin{solve}
在原点附近，$y$ 有任意阶导数。由 $y(-x)=-y(x)$ 逐次求导并令 $x=0$，得到所有偶数阶导数在原点均为零。令 $a_j=y^{(j)}(0)$，对恒等式 $y\cos x=\sin x$ 求 $2m+1$ 阶导数，再令 $x=0$，由莱布尼茨公式得到
\[
    \sum_{j=0}^m(-1)^j\binom{2m+1}{2j}a_{2m+1-2j}
    =(-1)^m.
\]
$m=0$ 时 $a_1=1$。依次取 $m=1,2,3$，有
\[
\begin{aligned}
    a_3-3a_1&=-1,\\
    a_5-10a_3+5a_1&=1,\\
    a_7-21a_5+35a_3-7a_1&=-1.
\end{aligned}
\]
所以
\[
    a_1=1,\quad a_3=2,\quad a_5=16,\quad a_7=272,
    \qquad a_2=a_4=a_6=a_8=0.
\]
这些数值随后可以用于泰勒多项式的构造，但本例的计算本身只依赖求导法则与恒等式。
\end{solve}

\subsection{一般函数的高阶导数}
\label{subsec:higher-derivatives-of-general-functions}

一般复合函数、反函数、隐函数或参数式函数未必具有简洁的任意阶公式，通常应从低阶求导关系逐次计算，并始终注明对哪个变量求导。例如在 $f,g$ 相应的二阶导数存在时，
\begin{equation}
\label{eq:second-derivative-of-composite-function}
    \frac{\mathrm d^2}{\mathrm dx^2}f(g(x))
    =f''(g(x))(g'(x))^2+f'(g(x))g''(x).
\end{equation}
这是先对 $f'(g(x))g'(x)$ 使用乘积法则，再对第一因子使用链式法则的结果。反函数也须再次求导：若 $g=f^{-1}$，$f$ 在开区间上连续严格单调、二次可导且 $f'\ne0$，则
\begin{equation}
\label{eq:second-derivative-of-inverse-function}
    g''(y)=-\frac{f''(g(y))}{(f'(g(y)))^3}.
\end{equation}
因而二阶导数之间并无一阶导数那样的简单倒数关系。

若参数式 $x=\varphi(t)$、$y=\psi(t)$ 在所考虑的邻域内二次连续可导，且 $\varphi'\ne0$，则局部可把 $y$ 看成 $x$ 的函数。对任何可导的参数表达式，关于 $x$ 求导相当于先关于 $t$ 求导再除以 $\varphi'(t)$，所以
\begin{equation}
\label{eq:parametric-second-derivative}
    \frac{\mathrm d^2y}{\mathrm dx^2}
    =\frac1{\varphi'(t)}\frac{\mathrm d}{\mathrm dt}
      \left(\frac{\psi'(t)}{\varphi'(t)}\right)
    =\frac{\psi''(t)\varphi'(t)-\psi'(t)\varphi''(t)}{(\varphi'(t))^3}.
\end{equation}
不能把 $\mathrm d^2y/\mathrm dx^2$ 写成 $\psi''/\varphi''$；每增加一次关于 $x$ 的求导，都必须重新使用一次 $\mathrm d/\mathrm dx=(1/\varphi')\,\mathrm d/\mathrm dt$。

\begin{example}
\label{ex:ellipse-second-derivative-three-methods}
设 $a,b>0$，上半椭圆由 $x^2/a^2+y^2/b^2=1$、$y>0$ 给出，分别用显式、参数式与隐式表示求 $y''$。
\end{example}
\begin{solve}
显式表示为 $y=(b/a)\sqrt{a^2-x^2}$（$|x|<a$），直接求导得到
\[
    y'=-\frac{bx}{a\sqrt{a^2-x^2}},\qquad
    y''=-\frac{ab}{(a^2-x^2)^{3/2}}
        =-\frac{b^4}{a^2y^3}.
\]
参数表示为 $x=a\cos t$、$y=b\sin t$（$0<t<\pi$），于是
\[
    \frac{\mathrm dy}{\mathrm dx}=-\frac ba\cot t,\qquad
    \frac{\mathrm d^2y}{\mathrm dx^2}
    =\frac1{-a\sin t}\frac{\mathrm d}{\mathrm dt}
      \left(-\frac ba\cot t\right)
    =-\frac b{a^2\sin^3t}
    =-\frac{b^4}{a^2y^3}.
\]
最后，不解出 $y$，直接把椭圆方程对 $x$ 连续求导两次，得到
\[
    \frac x{a^2}+\frac{yy'}{b^2}=0,\qquad
    \frac1{a^2}+\frac{(y')^2+yy''}{b^2}=0.
\]
代入 $y'=-b^2x/(a^2y)$，并利用椭圆方程化简，仍得 $y''=-b^4/(a^2y^3)$。三种计算的定义域一致；本例中隐式表示可以避免反复处理根式。
\end{solve}

\subsection{高阶微分}
\label{subsec:higher-differentials}

以 $x$ 为独立变量时，把 $\mathrm dx$ 看成求微分过程中保持不变的增量。若 $f$ 二次可导，则对 $\mathrm df=f'(x)\mathrm dx$ 再求微分，得到\emph{二阶微分}（second differential）
\[
    \mathrm d^2f=\mathrm d(\mathrm df)=f''(x)(\mathrm dx)^2.
\]
递归地，当所需阶数的导数存在时，定义 $n$ 阶微分为 $\mathrm d^nf=\mathrm d(\mathrm d^{n-1}f)$，从而
\begin{equation}
\label{eq:higher-differentials-independent-variable}
    \mathrm d^nf=f^{(n)}(x)(\mathrm dx)^n.
\end{equation}
常把 $(\mathrm dx)^n$ 简记为 $\mathrm dx^n$，因此高阶导数可记成 $\mathrm d^nf/\mathrm dx^n$。这里的 $\mathrm dx^2$ 表示 $(\mathrm dx)^2$，不能与 $\mathrm d(x^2)=2x\mathrm dx$ 或 $\mathrm d^2x$ 混淆；以 $x$ 为独立变量时，$\mathrm d^2x=0$。

\emph[caution]{高阶微分一般不具有一阶微分的形式不变性。}设 $y=f(u)$、$u=g(x)$，并假定两函数在相应邻域内二次可导。若 $u$ 是独立变量，则 $\mathrm d^2y=f''(u)(\mathrm du)^2$；但以 $x$ 为独立变量时，由式~\eqref{eq:second-derivative-of-composite-function}，
\begin{equation}
\label{eq:second-differential-dependent-inner-variable}
\begin{aligned}
    \mathrm d^2y
    &=\left[f''(u)(g'(x))^2+f'(u)g''(x)\right](\mathrm dx)^2\\
    &=f''(u)(\mathrm du)^2+f'(u)\,\mathrm d^2u.
\end{aligned}
\end{equation}
与把 $u$ 看作独立变量时相比，一般多出 $f'(u)\mathrm d^2u$。若 $g$ 为仿射函数，这个附加项为零；在一般非线性换元下则不能略去。

\begin{example}
\label{ex:failure-of-invariance-second-differential}
设 $y=u^2$、$u=x^2$，比较以 $u$ 和以 $x$ 为独立变量时的二阶微分。
\end{example}
\begin{solve}
以 $u$ 为独立变量，$\mathrm d^2y=2(\mathrm du)^2$。以 $x$ 为独立变量，则 $y=x^4$，所以 $\mathrm d^2y=12x^2(\mathrm dx)^2$。若只把 $\mathrm du=2x\mathrm dx$ 代入前一个形式，只能得到 $8x^2(\mathrm dx)^2$，遗漏的正是
\[
    f'(u)\mathrm d^2u=2u\cdot2(\mathrm dx)^2
    =4x^2(\mathrm dx)^2.
\]
因此高阶微分计算必须先明确独立变量；其求值虽归结为高阶导数，变换自变量时却不能照搬一阶微分的形式不变性。
\end{solve}

\section*{习题四}
\addcontentsline{toc}{section}{习题四}
\fi

\selectchapter{5}
\ifcompilechapter
\chapter{导数的应用}
\label{chap:applications-of-derivatives}

% 本章正文依据教材第 177--229 页；分章 PDF 第 1--53 页，习题五不在本次整理范围内。
导数描述函数在一点的局部变化，而函数的单调性、极值和图像形状涉及一段区间上的整体性质。连接二者的基本工具是微分中值定理：它把区间两端的函数值之差表示为某个内点处导数与自变量增量的乘积。本章先建立这一联系，再讨论用导数求极限、用多项式逼近函数，以及通过导数确定函数的升降、凹凸与渐近行为。以下所说的可导均指有限导数存在，闭区间端点若涉及导数，则另行说明单侧导数。

\section{微分中值定理}
\label{sec:differential-mean-value-theorems}

% 教材第 177--186 页；分章 PDF 第 1--10 页。
微分中值定理的基本证明路线是：由极值点处的必要条件得到罗尔定理，再通过减去适当的辅助函数，依次得到拉格朗日定理和柯西定理。证明中的辅助函数不是额外假设，而是为了把待证结论转化成某个函数导数为零的问题。

\subsection{费马定理}
\label{subsec:fermat-theorem}

\begin{definition}
\label{def:local-extrema-and-stationary-points}
设 $f$ 在 $x_0$ 的某邻域内有定义。若存在 $\delta>0$，使得 $|x-x_0|<\delta$ 时恒有 $f(x)\leq f(x_0)$，则称 $x_0$ 为 $f$ 的\emph{极大值点}（local maximum point），$f(x_0)$ 为相应的极大值；将不等号反向便得到\emph{极小值点}（local minimum point）和极小值。若在 $0<|x-x_0|<\delta$ 时相应不等式严格成立，则称为严格极大值或严格极小值。二者统称\emph{极值}（local extremum）。若 $f'(x_0)=0$，则称 $x_0$ 为\emph{驻点}（stationary point）。
\end{definition}

\begin{theorem}[费马定理]
\label{thm:fermat-stationary-point}
设 $f$ 在 $x_0$ 的某邻域内有定义。若 $x_0$ 是 $f$ 的极值点，且 $f$ 在 $x_0$ 可导，则 $f'(x_0)=0$。这一结论称为\emph{费马定理}（Fermat's theorem）。
\end{theorem}

\begin{proof}
以极大值点 $x_0$ 为例，令 $h=x-x_0$。由极大值点的定义，存在 $\delta>0$，使得当 $0<|h|<\delta$ 时，附近点 $x_0+h$ 的函数值不超过 $f(x_0)$，即
\[
    f(x_0+h)\leq f(x_0)
    \quad\Longrightarrow\quad
    f(x_0+h)-f(x_0)\leq0.
\]
这个结论对正、负两种 $h$ 都成立，因此无论从左侧还是右侧接近 $x_0$，差商的分子始终非正。

当 $0<h<\delta$ 时，分子非正、分母为正，所以
\[
    \frac{f(x_0+h)-f(x_0)}{h}\leq0.
\]
令 $h\to0^+$，由于 $f$ 在 $x_0$ 可导，这些差商的极限就是 $f'(x_0)$；非正差商的极限仍然非正，因此
\[
    f'(x_0)
    =\lim_{h\to0^+}\frac{f(x_0+h)-f(x_0)}{h}
    \leq0.
\]

当 $-\delta<h<0$ 时，分子仍然非正，但分母为负，除以负数时不等号反向，于是
\[
    \frac{f(x_0+h)-f(x_0)}{h}\geq0.
\]
再令 $h\to0^-$，同样由可导性及极限的保序性，得到
\[
    f'(x_0)
    =\lim_{h\to0^-}\frac{f(x_0+h)-f(x_0)}{h}
    \geq0.
\]
上述两个单侧极限都等于同一个 $f'(x_0)$，故 $f'(x_0)$ 同时非正、非负，只能等于零。极小值情形将上述不等号反向即可。
\emph[tip]{极大值条件控制差商分子的符号，$h$ 的正负决定差商的符号，而可导性保证左右两侧的差商趋向同一个数。}
\end{proof}

费马定理给出的是可导内点取得极值的必要条件。$f(x)=x^3$ 在 $0$ 处为驻点，却不取极值；$f(x)=|x|$ 在 $0$ 处取得严格极小值，却不可导。闭区间上的最值还可能出现在端点，例如 $f(x)=x$ 在 $[0,1]$ 上的最大值出现在 $1$，因而不能将费马定理直接用于端点。

\subsection{罗尔微分中值定理}
\label{subsec:rolle-theorem}

\begin{theorem}[罗尔微分中值定理]
\label{thm:rolle-mean-value}
设 $a<b$，$f\in C([a,b])$，$f$ 在 $(a,b)$ 内可导，且 $f(a)=f(b)$。则至少存在 $\xi\in(a,b)$，使
\[
    f'(\xi)=0.
\]
这一结论称为\emph{罗尔微分中值定理}（Rolle's theorem）。
\end{theorem}

\begin{proof}
由闭区间上连续函数的最值定理，$f$ 在 $[a,b]$ 上取得最大值 $M$ 和最小值 $m$。若 $M=m$，则 $f$ 为常函数，任取 $\xi\in(a,b)$ 即可。若 $M>m$，则 $M,m$ 中至少有一个不等于共同的端点值 $f(a)=f(b)$，这一最值必在内点 $\xi$ 取得。由费马定理，$f'(\xi)=0$。
\end{proof}

几何上，端点等高的可导曲线至少有一条水平切线。定理保证中值点存在，但通常既不能确定其具体位置，也不保证唯一。三个条件各有作用：$f(x)=x$ 不满足端点等值，$f(x)=|x|$ 在 $[-1,1]$ 内不可导，二者都没有导数为零的内点；在 $[0,1]$ 上定义 $f(0)=1$、$f(x)=x$（$x>0$），则端点等值且内点可导，但端点不连续，结论同样不成立。

\begin{example}
\label{ex:rolle-auxiliary-function-and-root-count}
设 $f$ 在 $(a,b)$ 内可导，证明：$f$ 的任意两个不同零点之间都有 $f'+f$ 的零点。
\end{example}

\begin{solve}
设 $u<v$ 是 $f$ 的两个零点。取 $F(x)=e^xf(x)$，则 $F$ 在 $[u,v]$ 上连续、在 $(u,v)$ 内可导，且 $F(u)=F(v)=0$。罗尔定理给出 $\xi\in(u,v)$，使
\[
    0=F'(\xi)=e^\xi\bigl(f'(\xi)+f(\xi)\bigr),
\]
故 $f'(\xi)+f(\xi)=0$。这里乘以 $e^x$ 的目的，是让乘积的导数产生所需的 $f'+f$。
\end{solve}

\begin{example}
\label{ex:rolle-exponential-polynomial-root-count}
设 $n$ 为非负整数，$P_n$ 为次数不超过 $n$ 的实多项式。证明：方程 $e^x-P_n(x)=0$ 至多有 $n+1$ 个不同实根。
\end{example}

\begin{solve}
令 $F(x)=e^x-P_n(x)$，则 $F$ 在 $\mathbb R$ 上任意阶可导。反设方程至少有 $n+2$ 个不同实根，取其中 $n+2$ 个，按大小排列为 $x_1<\cdots<x_{n+2}$，则 $F(x_i)=0$（$i=1,\ldots,n+2$）。

对每个 $i=1,\ldots,n+1$，$F$ 在 $[x_i,x_{i+1}]$ 上连续、在 $(x_i,x_{i+1})$ 内可导，且两端点函数值均为零。由罗尔定理，存在 $\xi_i\in(x_i,x_{i+1})$，使 $F'(\xi_i)=0$。这些开区间两两不交，故 $\xi_1,\ldots,\xi_{n+1}$ 互不相同，因而 $F'$ 至少有 $n+1$ 个不同零点。对各阶导数重复这一过程，每应用一次罗尔定理，所能保证的不同零点个数减少一个；于是 $F^{(k)}$ 至少有 $n+2-k$ 个不同零点（$k=1,\ldots,n+1$），特别地，$F^{(n+1)}$ 至少有一个零点。

然而，$P_n$ 的次数不超过 $n$，故 $P_n^{(n+1)}(x)=0$，从而
\[
    F^{(n+1)}(x)=e^x-P_n^{(n+1)}(x)=e^x>0,
    \qquad x\in\mathbb R.
\]
这与 $F^{(n+1)}$ 存在零点矛盾。因此原方程至多有 $n+1$ 个不同实根。
\end{solve}

\begin{example}
\label{ex:rolle-with-prescribed-interior-value}
设 $f\in C([0,1])$，在 $(0,1)$ 内可导，$f(0)=f(1)=0$，且某个 $t_0\in(0,1)$ 满足 $f(t_0)=A$。证明存在 $\xi\in(0,1)$，使 $f'(\xi)=A$。
\end{example}

\begin{solve}
若 $A=0$，直接使用罗尔定理。若 $A\neq0$，令 $F(x)=f(x)-Ax$，则
\[
    F(0)=0,\qquad F(t_0)=(1-t_0)A,\qquad F(1)=-A.
\]
后两个值异号，故由介值定理存在 $\eta\in(t_0,1)$，使 $F(\eta)=0$。在 $[0,\eta]$ 上应用罗尔定理，得到 $F'(\xi)=f'(\xi)-A=0$。
\end{solve}

\subsection{拉格朗日微分中值定理}
\label{subsec:lagrange-mean-value-theorem}

当两端点的函数值不相等时，将曲线减去通过两端点的直线，便可恢复罗尔定理所需的端点等值条件。

\begin{theorem}[拉格朗日微分中值定理]
\label{thm:lagrange-mean-value}
设 $a<b$，$f\in C([a,b])$，且 $f$ 在 $(a,b)$ 内可导。则存在 $\xi\in(a,b)$，使
\begin{equation}
    f(b)-f(a)=f'(\xi)(b-a).
    \label{eq:lagrange-mean-value}
\end{equation}
这一结论称为\emph{拉格朗日微分中值定理}（Lagrange mean value theorem）。
\end{theorem}

\begin{proof}
令
\[
    F(x)=f(x)-\left[f(a)+\frac{f(b)-f(a)}{b-a}(x-a)\right].
\]
则 $F$ 在闭区间连续、开区间可导，且 $F(a)=F(b)=0$。罗尔定理给出 $F'(\xi)=0$，即 $f'(\xi)=(f(b)-f(a))/(b-a)$。
\end{proof}

式~\eqref{eq:lagrange-mean-value} 表示某条切线平行于端点之间的割线，即某个瞬时变化率等于区间上的平均变化率。若 $h\neq0$，且 $f$ 在以 $x_0$、$x_0+h$ 为端点的闭区间上连续、内部可导，则也可写成
\begin{equation}
    f(x_0+h)=f(x_0)+f'(x_0+\theta h)h,
    \qquad 0<\theta<1.
    \label{eq:lagrange-increment-form}
\end{equation}
此式同时适用于正、负增量，但 $\theta$ 通常依赖于 $x_0,h$，不能当作固定常数。

\begin{corollary}[导数恒为零]
\label{cor:zero-derivative-constant}
设 $I$ 为非空区间，$f$ 在 $I$ 上连续、内部可导。若 $I$ 的内部每一点都满足 $f'(x)=0$，则存在常数 $C$，使 $f(x)=C$ 对所有 $x\in I$ 成立。
\end{corollary}

\begin{proof}
任取 $x,y\in I$ 且 $x<y$。因为 $I$ 是区间，所以 $[x,y]\subseteq I$，从而 $f$ 在 $[x,y]$ 上连续、在 $(x,y)$ 内可导。由拉格朗日中值定理，存在 $\xi\in(x,y)$，使
\[
    f(y)-f(x)=f'(\xi)(y-x)=0.
\]
因此任意两个不同点的函数值均相等。固定 $x_0\in I$，取 $C=f(x_0)$，便有 $f(x)=C$ 对所有 $x\in I$ 成立；$I$ 只有一个点时，结论也成立。
\end{proof}

这里的区间条件不能删除：在 $\mathbb{R}\setminus\{0\}$ 上分别取 $f(x)=0$（$x<0$）和 $f(x)=1$（$x>0$），虽处处有 $f'=0$，却不是整个定义域上的常函数。

\begin{corollary}[导数相等]
\label{cor:equal-derivatives-constant-difference}
设 $I$ 为非空区间，$f,g$ 都在 $I$ 上连续、内部可导。若 $I$ 的内部每一点都满足 $f'(x)=g'(x)$，则存在常数 $C$，使 $f(x)=g(x)+C$ 对所有 $x\in I$ 成立，即 $f$ 与 $g$ 只相差一个常数。
\end{corollary}

\begin{proof}
令 $H(x)=f(x)-g(x)$，则 $H$ 在 $I$ 上连续、内部可导，且内部每一点都满足
\[
    H'(x)=f'(x)-g'(x)=0.
\]
由推论~\ref{cor:zero-derivative-constant}，$H$ 在 $I$ 上为常数。因此固定 $x_0\in I$，取 $C=H(x_0)=f(x_0)-g(x_0)$，便有 $f(x)-g(x)=C$ 对所有 $x\in I$ 成立。
\end{proof}

\begin{corollary}[导数有界]
\label{cor:bounded-derivative-lipschitz}
设 $I$ 为非空区间，$f$ 在 $I$ 上连续、内部可导。若存在常数 $M\geq0$，使 $I$ 的内部每一点都满足 $|f'(x)|\leq M$，则
\begin{equation}
    |f(x)-f(y)|\leq M|x-y|,\qquad x,y\in I.
    \label{eq:bounded-derivative-lipschitz}
\end{equation}
特别地，$f$ 在 $I$ 上一致连续。这里的 $M$ 是在整个区间内部都适用的同一个上界，不能随所选的点而改变。
\end{corollary}

\begin{proof}
任取 $x,y\in I$ 且 $x<y$。因为 $[x,y]\subseteq I$，所以 $f$ 在 $[x,y]$ 上连续、在 $(x,y)$ 内可导。由拉格朗日中值定理，存在 $\xi\in(x,y)$，使 $f(y)-f(x)=f'(\xi)(y-x)$。结合 $|f'(\xi)|\leq M$，得到
\[
    |f(y)-f(x)|=|f'(\xi)|\,|y-x|\leq M|y-x|.
\]
交换 $x,y$ 不改变这个不等式，$x=y$ 时两边均为零，因此式~\eqref{eq:bounded-derivative-lipschitz} 对任意 $x,y\in I$ 都成立。

由式~\eqref{eq:bounded-derivative-lipschitz} 及下文命题~\ref{prop:lipschitz-implies-uniform-continuity}，可知 $f$ 在 $I$ 上一致连续。
\end{proof}

\begin{notealert}[title={什么是一致连续}]
\emph{一致连续}（uniform continuity）要求在整个区间内用同一个距离标准控制函数值的变化。为理解这一要求，先回顾普通连续性：$f$ 在区间 $I$ 上连续，是指它在每一点 $x_0\in I$ 都连续，即给定任意 $\varepsilon>0$，总能找到 $\delta>0$，使对任意 $x\in I$ 都有
\[
    |x-x_0|<\delta\quad\Longrightarrow\quad|f(x)-f(x_0)|<\varepsilon.
\]
这里先固定位置 $x_0$，再寻找合适的 $\delta$，所以 $\delta$ 可以同时依赖于 $\varepsilon$ 和 $x_0$，允许不同位置使用不同的距离标准。

一致连续则要求：给定任意 $\varepsilon>0$，都能先找到一个统一的 $\delta>0$，使对任意两点 $x,y\in I$ 都有
\[
    |x-y|<\delta\quad\Longrightarrow\quad|f(x)-f(y)|<\varepsilon.
\]
对于固定的函数 $f$ 和区间 $I$，这个 $\delta$ 可以依赖于 $\varepsilon$，但\emph{不能依赖于所选点的位置}。直观地说，普通连续是“每个位置都能找到适合自己的距离标准”，一致连续是“所有位置都能共用一个距离标准”；“一致”指这个标准统一，并不是说函数各处的变化率相同。将 $y$ 固定为任意一点 $x_0$，便得到 $f$ 在 $x_0$ 处连续，因此\emph{一致连续蕴含连续}；反向一般不成立，例如 $x^2$ 在 $\mathbb{R}$ 上连续却不一致连续，具体见下文例~\ref{ex:lipschitz-absolute-value-and-square}。

在本推论中，由于 $M\geq0$ 是整个区间共用的常数，给定任意 $\varepsilon>0$，统一取 $\delta=\varepsilon/(M+1)>0$，则任意满足 $|x-y|<\delta$ 的 $x,y\in I$ 都有
\[
    |f(x)-f(y)|\leq M|x-y|\leq(M+1)|x-y|<\varepsilon.
\]
这个 $\delta$ 与两点的位置无关，并且上述取法也适用于 $M=0$。因此，推论中的“特别地”表示：式~\eqref{eq:bounded-derivative-lipschitz} 给出了整个区间统一适用的变化量估计，从而保证一致连续。
\end{notealert}

\emph[tip]{导数 $f'(x)$ 描述某一点附近的变化率，而式~\eqref{eq:bounded-derivative-lipschitz} 把这种局部控制推广到了区间内任意两点之间。几何上，$|f'(x)|\leq M$ 表示区间内部各点处切线的斜率绝对值都不超过 $M$，结论则说明连接图像上任意两个不同点的割线，其斜率绝对值也不超过 $M$。}这种用同一个常数控制所有点对的性质，可以脱离可导性单独定义。

\begin{definition}
\label{def:lipschitz-continuity}
设 $I$ 为非空区间，$f:I\to\mathbb{R}$。若存在常数 $L\geq0$，使对任意 $x,y\in I$ 都有
\[
    |f(x)-f(y)|\leq L|x-y|,
\]
则称 $f$ 在 $I$ 上满足\emph{利普希茨条件}（Lipschitz condition），也称 $f$ 在 $I$ 上\emph{利普希茨连续}（Lipschitz continuous）；这样的 $L$ 称为 $f$ 在 $I$ 上的一个\emph{利普希茨常数}（Lipschitz constant）。
\end{definition}

这个定义的关键是\emph{先选定一个有限常数 $L$，再要求它对区间内所有点对都有效}。当 $x\neq y$ 时，不等式等价于
\[
    \left|\frac{f(x)-f(y)}{x-y}\right|\leq L,
\]
即函数图像上任意一条割线的斜率绝对值都不超过 $L$。利普希茨常数通常不唯一：若 $L$ 可用，则任意 $L'\geq L$ 也可用；$L=0$ 对应常函数。这里限制的是函数值的变化速度，并不要求函数值本身有界，例如 $f(x)=x$ 在 $\mathbb{R}$ 上可取 $L=1$，却是无界函数。

\begin{proposition}[利普希茨连续蕴含一致连续]
\label{prop:lipschitz-implies-uniform-continuity}
设 $I$ 为非空区间，$f:I\to\mathbb{R}$。若 $f$ 在 $I$ 上利普希茨连续，则 $f$ 在 $I$ 上一致连续。
\end{proposition}

\begin{proof}
由利普希茨连续的定义，存在固定常数 $L\geq0$，使对任意 $x,y\in I$ 都有
\[
    |f(x)-f(y)|\leq L|x-y|.
\]
若 $L>0$，给定任意 $\varepsilon>0$，取 $\delta=\varepsilon/L>0$。任取满足 $|x-y|<\delta$ 的 $x,y\in I$，便有
\[
    |f(x)-f(y)|\leq L|x-y|<L\delta=\varepsilon.
\]
这里的 $\delta$ 只依赖于 $\varepsilon$ 和固定的利普希茨常数 $L$，与点对 $x,y$ 的位置无关，因此同时适用于区间内所有点对。

若 $L=0$，则对任意 $x,y\in I$ 都有 $|f(x)-f(y)|=0$，所以 $f$ 是常函数。给定任意 $\varepsilon>0$，例如取 $\delta=1$，所有满足 $|x-y|<\delta$ 的点对都满足 $|f(x)-f(y)|=0<\varepsilon$。因此，无论 $L>0$ 还是 $L=0$，都满足一致连续的定义，故 $f$ 在 $I$ 上一致连续。
\end{proof}

将上述证明用于推论~\ref{cor:bounded-derivative-lipschitz} 时，只需令 $L=M$：当 $M>0$ 时，统一取 $\delta=\varepsilon/M$ 即可；当 $M=0$ 时，函数为常函数。这个证明只使用函数增量的不等式，不需要预先假设 $f$ 可导。由此得到\emph{利普希茨连续蕴含一致连续，进而蕴含连续}。利普希茨条件给出了明确的线性控制，而一致连续只要求对每个 $\varepsilon>0$ 能找到统一的 $\delta>0$，不要求这种线性关系；后面的例~\ref{ex:power-function-uniform-continuity-by-mean-value} 将说明一致连续不能反推出利普希茨连续。

\begin{proposition}
\label{prop:lipschitz-derivative-criterion}
设 $f$ 在非空区间 $I$ 上连续、内部可导，$L\geq0$ 为给定常数。则 $f$ 在 $I$ 上以 $L$ 为利普希茨常数，当且仅当 $I$ 的内部每一点都满足 $|f'(x)|\leq L$。
\end{proposition}

\begin{proof}
若内部恒有 $|f'|\leq L$，由推论~\ref{cor:bounded-derivative-lipschitz} 即得利普希茨条件。反之，设该条件成立。对任意内部点 $x$，当 $h\neq0$ 且 $|h|$ 充分小时，$x+h\in I$，故
\[
    \left|\frac{f(x+h)-f(x)}{h}\right|\leq L.
\]
令 $h\to0$，利用 $f$ 在 $x$ 处可导以及绝对值函数的连续性，得到 $|f'(x)|\leq L$。
\end{proof}

\begin{example}
\label{ex:lipschitz-absolute-value-and-square}
判断 $f(x)=|x|$ 与 $g(x)=x^2$ 在 $\mathbb{R}$ 上是否满足利普希茨条件，证明 $g$ 在 $\mathbb{R}$ 上连续但不一致连续，并讨论 $g$ 限制在 $[-R,R]$（$R>0$）以及 $[0,1]$ 上时的一致连续性。
\end{example}

\begin{solve}
由三角不等式 $|x|\leq|x-y|+|y|$ 及交换 $x,y$ 后的不等式，得到
\[
    \bigl||x|-|y|\bigr|\leq|x-y|.
\]
因此 $f(x)=|x|$ 在 $\mathbb{R}$ 上满足利普希茨条件，可取 $L=1$。但它在零点的左右导数分别为 $-1$ 和 $1$，所以在零点不可导。这说明利普希茨条件本身不保证处处可导；命题~\ref{prop:lipschitz-derivative-criterion} 中的可导性是预先给定的假设。

对于 $g(x)=x^2$，先考察它在一个固定位置 $x_0\in\mathbb{R}$ 处的连续性。由
\[
    |g(x)-g(x_0)|=|x^2-x_0^2|=|x-x_0|\,|x+x_0|,
\]
给定任意 $\varepsilon>0$，取
\[
    \delta=\min\left\{1,\frac{\varepsilon}{2|x_0|+1}\right\}>0.
\]
当 $|x-x_0|<\delta$ 时，$|x+x_0|\leq|x-x_0|+2|x_0|<2|x_0|+1$，从而
\[
    |g(x)-g(x_0)|\leq|x-x_0|(2|x_0|+1)
    <\delta(2|x_0|+1)\leq\varepsilon.
\]
因此 $g$ 在 $x_0$ 处连续。由于 $x_0$ 任意，$g$ 在整个 $\mathbb{R}$ 上连续。这里的 $\delta$ 允许依赖于位置 $x_0$；对固定的 $\varepsilon>0$，上述取法在 $|x_0|\to+\infty$ 时趋于零，但仅凭某一种取法依赖于 $x_0$，还不能断定不存在统一的 $\delta$，需要进一步验证。

为了否定一致连续性，取
\[
    x_n=n,\qquad y_n=n+\frac1n,\qquad n=1,2,\ldots.
\]
这两点的距离满足 $|y_n-x_n|=1/n\to0$，但函数值的距离为
\[
    |g(y_n)-g(x_n)|=\left(n+\frac1n\right)^2-n^2
    =2+\frac1{n^2}>2.
\]
具体地，固定 $\varepsilon_0=1$。无论给出多小的 $\delta>0$，都能取正整数 $n>1/\delta$，使 $|y_n-x_n|=1/n<\delta$，而 $|g(y_n)-g(x_n)|>2>\varepsilon_0$。因此不存在对整个实轴上所有点对都有效的 $\delta$，$g$ 在 $\mathbb{R}$ 上不一致连续；由命题~\ref{prop:lipschitz-implies-uniform-continuity}，它也不满足整个实轴上的利普希茨条件。

若将定义域限制在 $[-R,R]$，则任意 $x,y\in[-R,R]$ 都满足
\[
    |x^2-y^2|=|x-y|\,|x+y|\leq2R|x-y|,
\]
所以可取利普希茨常数 $L=2R$，并由上述命题得到一致连续性。特别地，在 $[0,1]$ 上有 $|x^2-y^2|\leq2|x-y|$，给定任意 $\varepsilon>0$，统一取 $\delta=\varepsilon/2$，就能保证 $|x-y|<\delta$ 时 $|x^2-y^2|<\varepsilon$。

从图像上看，$x^2$ 在 $|x|$ 越大的位置越陡，为保持同样的函数值误差，需要越来越小的自变量变化尺度；上述点列把这一现象转化为了严格的反例。将定义域限制在有界区间后，变化速度可以用一个有限常数统一控制。因此，连续、一致连续及利普希茨条件的区别，既体现在量词和增量估计上，也与所考察的定义域有关。
\end{solve}

\begin{example}
\label{ex:mean-value-logarithm-and-identity}
证明当 $x>-1$ 且 $x\neq0$ 时，
\[
    \frac{x}{1+x}<\ln(1+x)<x,
\]
并证明对任意 $x\in\mathbb{R}$，
\[
    \arctan x=\arcsin\frac{x}{\sqrt{1+x^2}}.
\]
\end{example}

\begin{solve}
对 $\ln(1+t)$ 在以 $0,x$ 为端点的区间上应用中值定理，得 $\ln(1+x)=x/(1+\xi)$，其中 $\xi$ 严格介于 $0,x$ 之间。若 $x>0$，则 $1/(1+x)<1/(1+\xi)<1$；若 $-1<x<0$，倒数的次序反向，而乘以负数 $x$ 后再次反向。两种情形都给出所需不等式。

第二式中，$u(x)=x/\sqrt{1+x^2}$ 满足 $|u(x)|<1$，且
\[
    u'(x)=\frac1{(1+x^2)^{3/2}},\qquad
    \frac{d}{dx}\arcsin u(x)
    =\frac{u'(x)}{\sqrt{1-u(x)^2}}=\frac1{1+x^2}.
\]
这与 $\arctan x$ 的导数相同，故二者之差为常数；再代入 $x=0$，知常数为零。
\end{solve}

\begin{example}
\label{ex:power-function-uniform-continuity-by-mean-value}
证明当 $0<\alpha<1$ 时，$f(x)=x^\alpha$ 在 $[0,+\infty)$ 上一致连续，但在该区间上不满足利普希茨条件。
\end{example}

\begin{solve}
$f$ 在 $[0,2]$ 上连续，故由康托尔定理在该区间上一致连续。在 $[1,+\infty)$ 上，$|f'(x)|=\alpha x^{\alpha-1}\leq\alpha<1$，因此 $|f(x)-f(y)|\leq|x-y|$。给定 $\varepsilon>0$，取 $[0,2]$ 上对应的 $\delta_1>0$，并令 $\delta=\min\{\delta_1,\varepsilon,1\}$。任意 $x,y\geq0$ 若满足 $|x-y|<\delta$，则二者或者都属于 $[0,2]$，或者都属于 $[1,+\infty)$：若有一个点大于 $2$，另一个点与它的距离小于 $1$，便也大于 $1$。在第一种情形下使用 $\delta_1$，在第二种情形下使用导数给出的增量估计，均得 $|f(x)-f(y)|<\varepsilon$，故 $f$ 在整个 $[0,+\infty)$ 上一致连续。

另一方面，取 $y=0$、$x>0$，则
\[
    \frac{|f(x)-f(0)|}{|x-0|}=x^{\alpha-1}\longrightarrow+\infty
    \qquad(x\to0^+).
\]
因此不存在有限常数 $L$ 能控制所有这样的点对，$f$ 在 $[0,+\infty)$ 上不满足利普希茨条件。这说明一致连续不能反推出利普希茨连续。与此同时，$f'(x)=\alpha x^{\alpha-1}$ 在 $x\to0^+$ 时无界，也说明导数有界是区间上一致连续的充分条件，却不是必要条件。
\end{solve}

\subsection{柯西微分中值定理}
\label{subsec:cauchy-mean-value-theorem}

拉格朗日公式中的 $b-a$ 是函数 $g(x)=x$ 的端点增量。将它换成一般函数的增量，就得到比较两个函数变化率的工具。

\begin{theorem}[柯西微分中值定理]
\label{thm:cauchy-mean-value}
设 $a<b$，$f,g\in C([a,b])$，且二者在 $(a,b)$ 内可导。则存在 $\xi\in(a,b)$，使
\begin{equation}
    \bigl(f(b)-f(a)\bigr)g'(\xi)
    =\bigl(g(b)-g(a)\bigr)f'(\xi).
    \label{eq:cauchy-mean-value-product}
\end{equation}
若另有 $g'(x)\neq0$ 对所有 $x\in(a,b)$ 成立，则 $g(b)\neq g(a)$，并有
\begin{equation}
    \frac{f(b)-f(a)}{g(b)-g(a)}=\frac{f'(\xi)}{g'(\xi)}.
    \label{eq:cauchy-mean-value-quotient}
\end{equation}
这一结论称为\emph{柯西微分中值定理}（Cauchy mean value theorem）。
\end{theorem}

\begin{proof}
先从待证结论寻找辅助函数。记 $f,g$ 在固定区间 $[a,b]$ 上的增量为
\[
    \Delta f=f(b)-f(a),\qquad \Delta g=g(b)-g(a).
\]
式~\eqref{eq:cauchy-mean-value-product} 等价于存在 $\xi\in(a,b)$，使
\[
    \Delta g\,f'(\xi)-\Delta f\,g'(\xi)=0.
\]
这提示我们：若能找到一个函数，其导数恰好是 $\Delta g\,f'(x)-\Delta f\,g'(x)$，就可以尝试用罗尔定理证明这个导数在某处等于零。

由于 $\Delta f,\Delta g$ 都是常数，直接取
\[
    F_0(x)=\Delta g\,f(x)-\Delta f\,g(x),
\]
便有
\[
    F_0'(x)=\Delta g\,f'(x)-\Delta f\,g'(x).
\]
这样，待证结论就转化为存在 $\xi\in(a,b)$ 使 $F_0'(\xi)=0$。

下面检查罗尔定理的条件。由 $f,g$ 的连续性与可导性，$F_0$ 在 $[a,b]$ 上连续，在 $(a,b)$ 内可导。罗尔定理只要求两个端点的函数值相等，因此计算
\[
    \begin{aligned}
        F_0(b)-F_0(a)
        &=\Delta g\,[f(b)-f(a)]-\Delta f\,[g(b)-g(a)]\\
        &=\Delta g\,\Delta f-\Delta f\,\Delta g=0.
    \end{aligned}
\]
于是 $F_0(a)=F_0(b)$，端点条件也得到满足，无须再通过减去常数将端点函数值调整为零。

由罗尔定理~\ref{thm:rolle-mean-value}，存在 $\xi\in(a,b)$，使
\[
    0=F_0'(\xi)=\Delta g\,f'(\xi)-\Delta f\,g'(\xi).
\]
移项并代回 $\Delta f,\Delta g$ 的定义，即得式~\eqref{eq:cauchy-mean-value-product}。这一推导没有使用除法，因此不要求 $g(b)-g(a)$ 或 $g'(\xi)$ 非零。

若进一步假设 $g'(x)\neq0$ 对所有 $x\in(a,b)$ 成立，则必有 $g(b)\neq g(a)$。否则，由 $g(a)=g(b)$ 及 $g$ 的连续性、可导性，对 $g$ 应用罗尔定理，便存在 $\eta\in(a,b)$ 使 $g'(\eta)=0$，与假设矛盾。又因 $\xi\in(a,b)$，有 $g'(\xi)\neq0$，故可将式~\eqref{eq:cauchy-mean-value-product} 两边同时除以非零量 $\bigl(g(b)-g(a)\bigr)g'(\xi)$，得到商式~\eqref{eq:cauchy-mean-value-quotient}。
\end{proof}

分别对 $f,g$ 使用拉格朗日公式，只会得到 $f'(\xi_1)/g'(\xi_2)$，两个中值点未必相同，不能据此推出式~\eqref{eq:cauchy-mean-value-quotient}。几何上，把 $x$ 看作参数，则该式表示参数曲线 $(g(x),f(x))$ 的某条切线平行于端点连线。

\begin{example}
\label{ex:cauchy-mean-value-weighted-derivative}
设 $f$ 在 $(0,1]$ 上连续、在 $(0,1)$ 内可导，且 $\lim_{x\to0+}\sqrt{x}\,f'(x)$ 存在且有限。证明 $f$ 在 $(0,1]$ 上一致连续。
\end{example}

\begin{solve}
由极限的局部有界性，存在 $M>0$ 和 $c\in(0,1)$，使 $0<x<c$ 时有 $|\sqrt{x}f'(x)|\leq M$。在任意不同的 $x,y\in(0,c)$ 之间，对 $f(t)$ 和 $g(t)=\sqrt{t}$ 使用柯西定理，得到
\[
    \frac{f(x)-f(y)}{\sqrt{x}-\sqrt{y}}=2\sqrt{\xi}\,f'(\xi),
\]
从而
\[
    |f(x)-f(y)|\leq2M|\sqrt{x}-\sqrt{y}|
    \leq2M\sqrt{|x-y|}.
\]
相同点的情形也满足该估计。于是当 $x,y\to0+$ 时，函数值之差趋于零，由函数极限的柯西准则，$f(0+)$ 存在且有限。以此值补定义 $f(0)$，得到 $[0,1]$ 上的连续函数，康托尔定理给出其一致连续性，限制回 $(0,1]$ 后结论仍成立。
\end{solve}

\section{洛必达法则}
\label{sec:lhopital-rule}

% 教材第 186--196 页；分章 PDF 第 10--20 页。
当分子、分母分别趋于零或同时无界时，仅凭它们各自的极限不能确定商的极限；$0/0$、$\infty/\infty$ 因而称为\emph{不定式}（indeterminate form）。它们是极限过程的类型标记，不是可以进行普通算术的数。\emph{洛必达法则}（L'H\^opital's rule）利用柯西中值定理，将某些函数商的极限化为导数商的极限。以下的无穷极限允许取 $+\infty$ 或 $-\infty$，双侧结论须在两侧分别满足条件并取得相同的极限。

\subsection{\texorpdfstring{$\frac{0}{0}$}{0/0}型不定式}
\label{subsec:zero-over-zero-indeterminate-form}

\begin{theorem}[$0/0$ 型洛必达法则]
\label{thm:lhopital-zero-over-zero}
设 $f,g$ 在有限点 $a$ 的某个右侧去心邻域 $(a,a+\delta)$ 内可导，并满足
\[
    f(x)\to0,\qquad g(x)\to0,\qquad g'(x)\neq0.
\]
若
\[
    \lim_{x\to a+}\frac{f'(x)}{g'(x)}=L,
    \qquad L\in\mathbb{R}\cup\{-\infty,+\infty\},
\]
则 $g$ 在该邻域内不为零，且
\begin{equation}
    \lim_{x\to a+}\frac{f(x)}{g(x)}=L.
    \label{eq:lhopital-zero-over-zero}
\end{equation}
同一结论适用于 $x\to a-$，并适用于满足相应条件的 $x\to+\infty$ 或 $x\to-\infty$。
\end{theorem}

\begin{proof}
补定义 $f(a)=g(a)=0$，使二者在 $a$ 处连续。若某个 $x>a$ 满足 $g(x)=0$，罗尔定理会给出 $(a,x)$ 内的 $g'$ 零点，与假设矛盾。对 $[a,x]$ 使用柯西中值定理，得到
\[
    \frac{f(x)}{g(x)}=\frac{f'(\xi_x)}{g'(\xi_x)},
    \qquad a<\xi_x<x.
\]
当 $x\to a+$ 时，$\xi_x\to a+$，故由导数商的极限得到结论，$L$ 有限或为有符号无穷时均成立。左侧极限的证明相同。对于 $x\to+\infty$，令 $t=1/x\to0+$，则
\[
    \frac{\frac{d}{dt}f(1/t)}{\frac{d}{dt}g(1/t)}
    =\frac{f'(1/t)(-1/t^2)}{g'(1/t)(-1/t^2)}
    =\frac{f'(1/t)}{g'(1/t)},
\]
因而转化为已证的有限端点情形；$x\to-\infty$ 用 $t\to0-$ 同理。
\end{proof}

\begin{example}
\label{ex:lhopital-zero-over-zero-calculations}
计算
\[
    \lim_{x\to0}\frac{\sinh x-x}{\sin x-x},
    \qquad
    \lim_{x\to0+}\frac{\sqrt{x}}{1-e^{2\sqrt{x}}},
\]
其中 $\sinh x=(e^x-e^{-x})/2$、$\cosh x=(e^x+e^{-x})/2$。
\end{example}

\begin{solve}
第一式在每次求导后仍为 $0/0$ 型，且各步分母的导数在充分小的去心邻域内均不为零。依次使用洛必达法则，得
\[
    \lim_{x\to0}\frac{\sinh x-x}{\sin x-x}
    =\lim_{x\to0}\frac{\cosh x-1}{\cos x-1}
    =\lim_{x\to0}\frac{\sinh x}{-\sin x}
    =\lim_{x\to0}\frac{\cosh x}{-\cos x}=-1.
\]
第二式先令 $t=\sqrt{x}\to0+$，避免对根式反复求导，再得
\[
    \lim_{t\to0+}\frac{t}{1-e^{2t}}
    =\lim_{t\to0+}\frac1{-2e^{2t}}=-\frac12.
    \qedhere
\]
\end{solve}

\subsection{\texorpdfstring{$\frac{\infty}{\infty}$}{infinity/infinity}型不定式}
\label{subsec:infinity-over-infinity-indeterminate-form}

\begin{theorem}[分母无界时的洛必达法则]
\label{thm:lhopital-infinite-denominator}
设 $f,g$ 在某个单侧去心邻域内可导，或在某条趋向无穷远的半直线上可导，且在相应极限过程中
\[
    |g(x)|\to+\infty,\qquad g'(x)\neq0,
    \qquad \frac{f'(x)}{g'(x)}\to L
    \in\mathbb{R}\cup\{-\infty,+\infty\}.
\]
则
\begin{equation}
    \frac{f(x)}{g(x)}\to L.
    \label{eq:lhopital-infinite-denominator}
\end{equation}
特别地，$\infty/\infty$ 型不定式可用此定理处理；定理本身\emph[caution]{只要求分母绝对值趋于无穷，不要求分子也趋于无穷，甚至不要求分子预先存在极限}。
\end{theorem}

\begin{proof}
先设 $L$ 有限。给定 $\varepsilon>0$，在足够靠近所考察端点的区间内，有 $|f'/g'-L|<\varepsilon$。在该区间固定一点 $y$，对更靠近端点的 $x$，将柯西定理用于 $f-Lg$ 与 $g$，得到
\[
    |f(x)-Lg(x)-f(y)+Lg(y)|
    \leq\varepsilon|g(x)-g(y)|.
\]
由于 $|g(x)|\to+\infty$，除以 $|g(x)|$ 后有
\[
    \left|\frac{f(x)}{g(x)}-L\right|
    \leq\varepsilon\left|1-\frac{g(y)}{g(x)}\right|
       +\frac{|f(y)-Lg(y)|}{|g(x)|}.
\]
右端的极限为 $\varepsilon$，由 $\varepsilon$ 的任意性得到结论。这里先按 $\varepsilon$ 选取区间和固定点 $y$，再令 $x$ 趋近端点，不能误认为相应的中值点必自动趋于端点。

若 $L=+\infty$，给定任意 $A>0$，可在足够靠近端点的区间内使 $f'/g'>2A$。同样固定 $y$，柯西定理给出
\[
    Q_x:=\frac{f(x)-f(y)}{g(x)-g(y)}>2A,
    \qquad
    \frac{f(x)}{g(x)}
    =Q_x\left(1-\frac{g(y)}{g(x)}\right)+\frac{f(y)}{g(x)}.
\]
括号内的因子趋于 $1$ 且最终为正，末项趋于零，故最终 $f(x)/g(x)>A$，即商趋于 $+\infty$。$L=-\infty$ 时对 $-f$ 使用此结论即可。上述论证同时适用于有限端点和无穷远处。
\end{proof}

\begin{example}
\label{ex:lhopital-growth-rates}
设 $\alpha>0$，用洛必达法则验证
\[
    \lim_{x\to+\infty}\frac{\ln x}{x^\alpha}=0,
    \qquad
    \lim_{x\to+\infty}\frac{x^\alpha}{e^x}=0.
\]
\end{example}

\begin{solve}
第一式的导数商为 $1/(\alpha x^\alpha)\to0$，故结论成立。第二式取正整数 $m>\alpha$，先对 $x^m/e^x$ 连续应用 $m$ 次洛必达法则，最后得到 $m!/e^x\to0$；每次使用前的商都是 $\infty/\infty$ 型，分母导数始终非零。当 $x\geq1$ 时，$0\leq x^\alpha/e^x\leq x^m/e^x$，再由夹逼原理得到所需极限。这为\crosschapterref{3}{命题~\ref{prop:logarithm-power-exponential-orders} 的增长次序}{对数、幂与指数的增长次序}提供了另一种证明。
\end{solve}

以下用 $C^n(I)$ 表示在区间 $I$ 上具有 $n$ 阶连续导数的函数类，约定 $C^0(I)=C(I)$；用 $C^\infty(I)$ 表示任意阶连续可导函数类，其元素称为\emph{光滑函数}（smooth function）。闭区间上的 $C^n$ 要求内部的各阶导数连续延拓至端点，端点导数按单侧意义理解。

\begin{example}
\label{ex:smooth-flat-exponential-function}
定义
\[
    f(x)=
    \begin{cases}
        e^{-1/x^2},&x\neq0,\\
        0,&x=0.
    \end{cases}
\]
证明 $f$ 在 $\mathbb{R}$ 上具有任意阶连续导数，且 $f^{(n)}(0)=0$ 对一切非负整数 $n$ 成立。
\end{example}

\begin{solve}
先注意，对任意 $m\geq0$，令 $t=1/x^2\to+\infty$，由幂与指数的增长比较得到
\[
    |x|^{-m}e^{-1/x^2}=t^{m/2}e^{-t}\to0.
\]
归纳证明：存在多项式 $P_n$，使 $x\neq0$ 时有 $f^{(n)}(x)=P_n(1/x)e^{-1/x^2}$，而 $f^{(n)}(0)=0$ 且 $f^{(n)}$ 在零点连续。$n=0$ 时取 $P_0=1$ 即可。若结论对 $n$ 成立，则
\[
    f^{(n+1)}(0)=\lim_{x\to0}\frac{P_n(1/x)e^{-1/x^2}}x=0.
\]
在非零点直接求导，则所得多项式满足
\[
    P_{n+1}(t)=-t^2P_n'(t)+2t^3P_n(t),
\]
故 $f^{(n+1)}(x)\to0=f^{(n+1)}(0)$，归纳完成，因而 $f\in C^\infty(\mathbb{R})$。
\end{solve}

\subsection{其他类型的不定式}
\label{subsec:other-indeterminate-forms}

$0\cdot\infty$ 型可改写成 $u/(1/v)$ 或 $v/(1/u)$；$\infty-\infty$ 型可通分、有理化或提取主项；对 $u(x)^{v(x)}$，若 $u(x)>0$，则先研究
\[
    \ln\bigl(u(x)^{v(x)}\bigr)=v(x)\ln u(x).
\]
这样，$1^\infty$、$0^0$ 和 $\infty^0$ 型通常转为乘积，再转为商。所作变形须在相应去心邻域内有意义，最后通过指数函数还原极限。选择变形的标准是让导数商更简单，而不是机械地把每个表达式都交给洛必达法则。

\begin{example}
\label{ex:other-indeterminate-form-transformations}
设 $\alpha>0$，计算
\[
    \lim_{x\to0+}x^\alpha\ln x,
    \qquad \lim_{x\to0+}x^x,
    \qquad \lim_{x\to+\infty}\bigl(\sqrt{x^2+x}-x\bigr).
\]
\end{example}

\begin{solve}
第一式写成 $\ln x/x^{-\alpha}$，是 $\infty/\infty$ 型，导数商为 $-x^\alpha/\alpha\to0$，故极限为零。第二式取对数得 $\ln(x^x)=x\ln x\to0$，故 $x^x\to e^0=1$。第三式直接有理化便有
\[
    \sqrt{x^2+x}-x
    =\frac{x}{\sqrt{x^2+x}+x}
    =\frac1{\sqrt{1+1/x}+1}\to\frac12.
    \qedhere
\]
\end{solve}

\begin{example}
\label{ex:power-means-parameter-limits}
设 $a_1,\ldots,a_n>0$，对 $r\neq0$ 定义\emph{幂平均}（power mean）
\[
    M_r=\left(\frac1n\sum_{i=1}^n a_i^r\right)^{1/r}.
\]
求 $r\to0$、$r\to+\infty$ 和 $r\to-\infty$ 时的极限。
\end{example}

\begin{solve}
当 $r\to0$ 时，
\[
    \ln M_r=\frac{\ln\bigl(\frac1n\sum_i a_i^r\bigr)}r
    \longrightarrow\frac1n\sum_{i=1}^n\ln a_i,
\]
其中使用了 $0/0$ 型洛必达法则。因此可连续补定义
\begin{equation}
    M_0=\left(\prod_{i=1}^n a_i\right)^{1/n}.
    \label{eq:power-mean-zero-parameter}
\end{equation}
记 $A=\max_i a_i$。当 $r>0$ 时，$A^r/n\leq\frac1n\sum_i a_i^r\leq A^r$，故
\[
    A n^{-1/r}\leq M_r\leq A,
\]
夹逼得到 $M_r\to A$。再利用 $M_r(a_1,\ldots,a_n)=1/M_{-r}(1/a_1,\ldots,1/a_n)$，得到
\begin{equation}
    \lim_{r\to-\infty}M_r=\min_i a_i,
    \qquad \lim_{r\to+\infty}M_r=\max_i a_i.
    \label{eq:power-mean-infinite-parameter-limits}
\end{equation}
这些极限说明几何平均和最大、最小值都可纳入同一参数族；其单调性将在第~\ref{subsec:convexity-and-concavity} 小节中证明。
\end{solve}

\begin{example}
\label{ex:iterated-sine-asymptotic-rate}
设 $0<x_1<\pi$，$x_{n+1}=\sin x_n$。证明 $\sqrt{n}\,x_n\to\sqrt{3}$。
\end{example}

\begin{solve}
由 $0<\sin x<x$（$0<x<\pi$），序列正且严格递减，故收敛。设极限为 $\ell\geq0$，由连续性得 $\ell=\sin\ell$，因而 $\ell=0$。为了求收敛速度，先由洛必达法则和 $\sin t/t\to1$ 得到
\[
    \lim_{t\to0}\frac{t-\sin t}{t^3}
    =\lim_{t\to0}\frac{1-\cos t}{3t^2}
    =\lim_{t\to0}\frac{\sin t}{6t}=\frac16.
\]
于是，令 $u_n=1/x_n^2$，有
\[
    u_{n+1}-u_n
    =\frac{x_n^2-\sin^2x_n}{x_n^2\sin^2x_n}
    \longrightarrow\frac13.
\]
由\crosschapterref{2}{定理~\ref{thm:cesaro-mean-of-convergent-sequence}}{算术平均保持极限的定理}，这些增量的算术平均也趋于 $1/3$，故
\[
    \frac{u_n}n=\frac{u_1}n+
    \frac1n\sum_{k=1}^{n-1}(u_{k+1}-u_k)\longrightarrow\frac13.
\]
因此 $nx_n^2\to3$，再由 $x_n>0$ 得到结论。
\end{solve}

\emph[caution]{洛必达法则是有条件的充分判据，不能由函数商有极限反推导数商有极限。}例如，$x\to+\infty$ 时 $(x+\sin x)/x\to1$，但导数商 $1+\cos x$ 没有极限。每次重复使用法则，都须重新核对不定式类型、可导性、分母导数非零及导数商的极限；若这些条件不满足，当前的求导链便不能作为证明。例如 $(x-\sin x)/(x+\sin x)\to1$，但第一次形式求导后的分母导数 $1+\cos x$ 有任意大的零点，导数商本身也无极限，继续求导得出 $-1$ 是无效推理。

\section{泰勒公式}
\label{sec:taylor-formula}

% 教材第 197--210 页；分章 PDF 第 21--34 页。
函数在 $x_0$ 处连续，只能保证 $f(x)=f(x_0)+o(1)$；可导则进一步给出 $f(x)=f(x_0)+f'(x_0)(x-x_0)+o(x-x_0)$。泰勒公式将这一思想推广到高阶：通过匹配函数及其若干阶导数，用多项式逼近函数，并单独研究逼近误差。佩亚诺余项描述靠近展开点时误差的阶，拉格朗日余项则通过高阶导数给出可计算的误差界；插值多项式处理的是在多个指定节点上匹配函数值的另一类逼近问题。

\subsection{带佩亚诺余项的泰勒公式}
\label{subsec:taylor-formula-with-peano-remainder}

\begin{theorem}[带佩亚诺余项的泰勒公式]
\label{thm:taylor-peano}
设 $n\geq1$，$f$ 在 $x_0$ 处具有 $n$ 阶导数；这里指 $f$ 在 $x_0$ 附近具有直到 $n-1$ 阶的导数，且 $f^{(n-1)}$ 在 $x_0$ 可导，不要求 $f^{(n)}$ 在邻域内连续。则当 $x\to x_0$ 时，
\begin{equation}
    f(x)=\sum_{k=0}^n\frac{f^{(k)}(x_0)}{k!}(x-x_0)^k
         +o\bigl((x-x_0)^n\bigr),
    \label{eq:taylor-peano}
\end{equation}
其中约定 $f^{(0)}=f$、$0!=1$。多项式
\[
    T_n(x)=\sum_{k=0}^n\frac{f^{(k)}(x_0)}{k!}(x-x_0)^k
\]
称为 $f$ 在 $x_0$ 处的 $n$ 阶\emph{泰勒多项式}（Taylor polynomial）；$R_n(x)=f(x)-T_n(x)$ 称为余项，式~\eqref{eq:taylor-peano} 中的余项称为\emph{佩亚诺余项}（Peano remainder）。当 $x_0=0$ 时，称为\emph{麦克劳林公式}（Maclaurin formula）。
\end{theorem}

\begin{proof}
当 $n=1$ 时就是导数定义。对 $n\geq2$，只需证明 $R_n(x)/(x-x_0)^n\to0$。由于各阶导数在 $x_0$ 的相应连续性，逐次求导后直到第 $n-2$ 次的商均为 $0/0$ 型，分母的下一阶导数在去心邻域内非零。连续使用 $n-1$ 次洛必达法则，原极限归结为
\[
    \lim_{x\to x_0}
    \frac{f^{(n-1)}(x)-f^{(n-1)}(x_0)-f^{(n)}(x_0)(x-x_0)}
         {n!(x-x_0)}=0.
\]
最后一步由 $f^{(n-1)}$ 在 $x_0$ 的可导性成立，再沿各步洛必达法则反推即得结论。
\end{proof}

\begin{proposition}[多项式逼近的唯一性]
\label{prop:taylor-polynomial-uniqueness}
若次数不超过 $n$ 的多项式 $P,Q$ 都满足
\[
    f(x)-P(x)=o\bigl((x-x_0)^n\bigr),\qquad
    f(x)-Q(x)=o\bigl((x-x_0)^n\bigr),
\]
则 $P=Q$。因此在定理~\ref{thm:taylor-peano} 的条件下，任何具有这一余项阶的多项式逼近，其系数必为 $f^{(k)}(x_0)/k!$。
\end{proposition}

\begin{proof}
相减得到 $P-Q=o((x-x_0)^n)$。若 $P-Q$ 不恒为零，设其关于 $x-x_0$ 的最低非零项为 $c_m(x-x_0)^m$，其中 $m\leq n$、$c_m\neq0$。除以 $(x-x_0)^m$ 并令 $x\to x_0$，左端趋于 $c_m$，右端趋于零，矛盾。
\end{proof}

由唯一性，可以通过恒等变形、有限展开的相乘与复合来求泰勒多项式，无须总是直接计算所有高阶导数；每一步仍须核对余项的阶。\emph[caution]{不能在没有额外条件时逐项微分一个带小 $o$ 余项的等式，因为小 $o$ 关系本身不控制余项的导数。}

\begin{proposition}[常用麦克劳林公式]
\label{prop:standard-maclaurin-formulas}
对固定的正整数 $n$，当 $x\to0$ 时，有
\begin{align}
    e^x&=\sum_{k=0}^n\frac{x^k}{k!}+o(x^n),
        \label{eq:maclaurin-exponential}\\
    \sin x&=\sum_{k=0}^n\frac{(-1)^kx^{2k+1}}{(2k+1)!}
              +o(x^{2n+1}),
        \label{eq:maclaurin-sine}\\
    \cos x&=\sum_{k=0}^n\frac{(-1)^kx^{2k}}{(2k)!}
              +o(x^{2n}),
        \label{eq:maclaurin-cosine}\\
    \ln(1+x)&=\sum_{k=1}^n\frac{(-1)^{k-1}}k x^k+o(x^n),
        \label{eq:maclaurin-logarithm}\\
    (1+x)^\alpha&=\sum_{k=0}^n\binom\alpha k x^k+o(x^n),
        \label{eq:maclaurin-real-power}
\end{align}
其中 $\alpha\in\mathbb{R}$，$\binom\alpha0=1$，且
\[
    \binom\alpha k=\frac{\alpha(\alpha-1)\cdots(\alpha-k+1)}{k!}
    \quad(k\geq1).
\]
对数和实幂取 $x>-1$ 的零点邻域。此外，
\begin{align}
    \arctan x&=\sum_{k=0}^n\frac{(-1)^k}{2k+1}x^{2k+1}
                  +o(x^{2n+1}),
        \label{eq:maclaurin-arctangent}\\
    \arcsin x&=\sum_{k=0}^n
        \frac{(2k-1)!!}{(2k)!!(2k+1)}x^{2k+1}
                  +o(x^{2n+1}),
        \label{eq:maclaurin-arcsine}
\end{align}
其中 $(2k)!!=2\cdot4\cdots(2k)$、$(2k-1)!!=1\cdot3\cdots(2k-1)$，并约定 $0!!=(-1)!!=1$。
\end{proposition}

\begin{proof}
$e^x$ 的各阶导数在零点均为 $1$；正弦与余弦的导数循环给出交替系数及相应的零系数。又对 $k\geq1$，
\[
    \frac{d^k}{dx^k}\ln(1+x)
    =\frac{(-1)^{k-1}(k-1)!}{(1+x)^k},\qquad
    \frac{d^k}{dx^k}(1+x)^\alpha
    =\alpha(\alpha-1)\cdots(\alpha-k+1)(1+x)^{\alpha-k}.
\]
在零点代入泰勒公式，便得前五式。反正切和反正弦在 $(-1,1)$ 内光滑，它们的导函数分别为 $(1+x^2)^{-1}$ 与 $(1-x^2)^{-1/2}$。由实幂公式及唯一性，这两个导函数的 $x^{2k}$ 项系数分别为 $(-1)^k$ 和 $(2k-1)!!/(2k)!!$，奇次项系数为零。由于导函数的这些泰勒系数分别等于原函数相应导数除以 $(2k)!$，原函数的 $x^{2k+1}$ 项系数还须除以 $2k+1$。再结合两函数在零点的值为零，即得后两式；此处使用的是光滑函数的泰勒系数关系，没有对未知的余项求导。
\end{proof}

\begin{example}
\label{ex:taylor-composition-and-product}
求 $e^{\cos x}$、$e^x\ln(1+x)$ 和 $\sin^2(1+x^2)$ 在零点展开到四次项、余项为 $o(x^4)$ 的公式。
\end{example}

\begin{solve}
令 $u=\cos x-1=-x^2/2+x^4/24+o(x^4)$，则 $u=O(x^2)$、$u^2=x^4/4+o(x^4)$，所以
\[
    e^{\cos x}=e\,e^u
    =e\left(1+u+\frac{u^2}{2}+o(u^2)\right)
    =e\left(1-\frac{x^2}{2}+\frac{x^4}{6}+o(x^4)\right).
\]
第二式将两个有限展开相乘，保留总次数不超过 $4$ 的项，得
\[
    e^x\ln(1+x)=x+\frac{x^2}{2}+\frac{x^3}{3}+0\cdot x^4+o(x^4).
\]
四次项系数为 $-1/4+1/3-1/4+1/6=0$，不能未经计算就将它漏去。第三式利用二倍角公式，得
\[
    \sin^2(1+x^2)
    =\frac{1-\cos2\cos(2x^2)+\sin2\sin(2x^2)}2
    =\sin^21+\sin2\,x^2+\cos2\,x^4+o(x^4).
    \qedhere
\]
\end{solve}

\begin{example}
\label{ex:taylor-peano-sequence-limit}
计算 $\lim_{n\to\infty}n^2(1-n\sin(1/n))$。
\end{example}

\begin{solve}
由 $\sin t=t-t^3/6+o(t^3)$，令 $t=1/n$，得到
\[
    n^2\left(1-n\sin\frac1n\right)
    =n^2\left(\frac1{6n^2}+o(n^{-2})\right)\longrightarrow\frac16.
    \qedhere
\]
\end{solve}

佩亚诺公式是固定阶数 $n$、令 $x\to x_0$ 的局部结论，不直接给出某个指定 $x$ 上的数值误差，也不允许未经证明地令 $n\to\infty$。例~\ref{ex:smooth-flat-exponential-function} 中的函数在零点的每个泰勒多项式都恒为零，但任意 $x\neq0$ 处函数值为正，说明即使具有任意阶连续导数，也未必等于由其导数组成的无限展开。

\subsection{带拉格朗日余项的泰勒公式}
\label{subsec:taylor-formula-with-lagrange-remainder}

\begin{theorem}[带拉格朗日余项的泰勒公式]
\label{thm:taylor-lagrange}
设 $n\geq0$，$f\in C^n([a,b])$，且 $f^{(n+1)}$ 在 $(a,b)$ 内存在。则对任意不同的 $x_0,x\in[a,b]$，存在严格介于二者之间的 $\xi$，使
\begin{equation}
    f(x)=\sum_{k=0}^n\frac{f^{(k)}(x_0)}{k!}(x-x_0)^k
       +\frac{f^{(n+1)}(\xi)}{(n+1)!}(x-x_0)^{n+1}.
    \label{eq:taylor-lagrange}
\end{equation}
最后一项称为\emph{拉格朗日余项}（Lagrange remainder），其中 $\xi=x_0+\theta(x-x_0)$、$0<\theta<1$。
\end{theorem}

\begin{proof}
固定 $x$，以 $t$ 为变量，在 $x_0,x$ 之间定义
\[
    F(t)=f(x)-\sum_{k=0}^n\frac{f^{(k)}(t)}{k!}(x-t)^k,
    \qquad G(t)=(x-t)^{n+1}.
\]
二者在相应闭区间连续、内部可导，且 $F(x)=G(x)=0$。求导时相邻项抵消，得到
\[
    F'(t)=-\frac{f^{(n+1)}(t)}{n!}(x-t)^n,
    \qquad G'(t)=-(n+1)(x-t)^n.
\]
$G'$ 在该开区间内非零，由柯西中值定理，
\[
    \frac{F(x_0)}{G(x_0)}
    =\frac{F(x_0)-F(x)}{G(x_0)-G(x)}
    =\frac{F'(\xi)}{G'(\xi)}
    =\frac{f^{(n+1)}(\xi)}{(n+1)!},
\]
整理即得所需公式。
\end{proof}

\begin{corollary}[泰勒逼近的误差界]
\label{cor:taylor-lagrange-error-bound}
在定理~\ref{thm:taylor-lagrange} 的条件下，若 $|f^{(n+1)}(t)|\leq M$ 对 $a<t<b$ 成立，则
\begin{equation}
    |f(x)-T_n(x)|\leq\frac{M}{(n+1)!}|x-x_0|^{n+1},
    \qquad x\in[a,b].
    \label{eq:taylor-lagrange-error-bound}
\end{equation}
当 $x=x_0$ 时，两端均为零。因此只要控制相应区间上的高阶导数，就能按给定精度选择多项式次数。
\end{corollary}

常用麦克劳林公式的余项也可写成这一形式。例如对 $n\geq1$，
\begin{align*}
    e^x&=\sum_{k=0}^n\frac{x^k}{k!}
        +\frac{e^{\theta_1x}}{(n+1)!}x^{n+1},\\
    \ln(1+x)&=\sum_{k=1}^n\frac{(-1)^{k-1}}k x^k
        +\frac{(-1)^nx^{n+1}}{(n+1)(1+\theta_2x)^{n+1}},\\
    (1+x)^\alpha&=\sum_{k=0}^n\binom\alpha k x^k
        +\binom\alpha{n+1}(1+\theta_3x)^{\alpha-n-1}x^{n+1},
\end{align*}
其中各 $\theta_i\in(0,1)$ 可互不相同；指数公式适用于任意实数 $x$，后两式在 $x>-1$ 时适用。对正弦展开到 $x^{2m+1}$ 项时，可在泰勒多项式中同时保留系数为零的 $x^{2m+2}$ 项，从而得到绝对值不超过 $|x|^{2m+3}/(2m+3)!$ 的余项；余弦展开到 $x^{2m}$ 项时，同理得到绝对值不超过 $|x|^{2m+2}/(2m+2)!$ 的余项。上述两项估计对 $m\geq0$ 和任意实数 $x$ 成立。

\begin{example}
\label{ex:taylor-approximation-and-irrationality-of-e}
用有限和近似 $e$，使绝对误差小于 $10^{-5}$，并用泰勒余项证明 $e$ 为无理数。
\end{example}

\begin{solve}
将指数函数在零点展开并令 $x=1$，有
\[
    e=\sum_{k=0}^n\frac1{k!}+\frac{e^\theta}{(n+1)!},
    \qquad 0<\theta<1.
\]
这里可先取 $n=2$，由 $e^\theta<e$ 得 $e<5/2+e/6$，从而 $e<3$。因此取 $n=8$ 时，
\[
    0<e-\sum_{k=0}^8\frac1{k!}<\frac3{9!}<10^{-5}.
\]
有限和为 $2.7182787698\ldots$；若交付小数近似，还须把舍入误差计入总误差。

若 $e=p/q$ 为有理数，其中 $p,q$ 为正整数，取整数 $n>\max\{q,3\}$，将上述等式乘以 $n!$，得到
\[
    n!e-\sum_{k=0}^n\frac{n!}{k!}=\frac{e^\theta}{n+1}.
\]
左端为整数，右端却严格介于 $0$ 与 $3/(n+1)<1$ 之间，矛盾。因此 $e$ 为无理数。
\end{solve}

\begin{example}
\label{ex:taylor-factorial-exponential-limit}
计算 $\lim_{n\to\infty}n\sin(2\pi n!e)$。
\end{example}

\begin{solve}
将 $e^x$ 展开到 $n+1$ 次并令 $x=1$，乘以 $n!$，可得
\[
    n!e=N_n+d_n,\qquad
    N_n=\sum_{k=0}^n\frac{n!}{k!}\in\mathbb{Z},\qquad
    d_n=\frac1{n+1}+\frac{e^{\theta_n}}{(n+1)(n+2)},
\]
其中 $0<\theta_n<1$。由 $1<e^{\theta_n}<3$，得 $d_n\to0$ 且 $nd_n\to1$。利用正弦的周期性及 $\sin t/t\to1$，得到
\[
    n\sin(2\pi n!e)=n\sin(2\pi d_n)
    =2\pi nd_n\frac{\sin(2\pi d_n)}{2\pi d_n}\longrightarrow2\pi.
    \qedhere
\]
\end{solve}

\subsection{拉格朗日插值多项式}
\label{subsec:lagrange-interpolation-polynomial}

泰勒多项式匹配一个点上的多阶导数；若已知的是不同节点上的函数值，则应改用\emph{多项式插值}（polynomial interpolation）。设 $n\geq0$ 为整数，互不相同的节点为 $x_0,\ldots,x_n$，给定数据 $y_0,\ldots,y_n$，目标是寻找次数不超过 $n$ 的多项式 $L_n$，使 $L_n(x_i)=y_i$。

\begin{theorem}[拉格朗日插值公式]
\label{thm:lagrange-interpolation}
上述插值问题存在唯一解，且
\begin{equation}
    L_n(x)=\sum_{i=0}^n y_i\ell_i(x),\qquad
    \ell_i(x)=\prod_{\substack{0\leq j\leq n\\j\neq i}}
                  \frac{x-x_j}{x_i-x_j}.
    \label{eq:lagrange-interpolation}
\end{equation}
这里约定空乘积为 $1$。$\ell_i$ 称为\emph{拉格朗日插值基函数}（Lagrange basis polynomial），满足 $\ell_i(x_j)=1$（$i=j$）或 $0$（$i\neq j$）。$L_n$ 称为\emph{拉格朗日插值多项式}（Lagrange interpolation polynomial）；其实际次数可以小于 $n$。
\end{theorem}

\begin{proof}
节点互异保证各分母非零。基函数在本节点取 $1$，在其他节点取 $0$，因而式~\eqref{eq:lagrange-interpolation} 满足全部插值条件且次数不超过 $n$，给出存在性。若另有同次数上界的多项式 $P_n$ 满足条件，则 $P_n-L_n$ 有 $n+1$ 个不同零点。次数不超过 $n$ 的非零多项式不可能有这么多零点，故 $P_n=L_n$。
\end{proof}

\begin{theorem}[插值余项]
\label{thm:lagrange-interpolation-remainder}
设 $f\in C^{n+1}([a,b])$，$x_0,\ldots,x_n\in[a,b]$ 互不相同，且 $L_n(x_i)=f(x_i)$。对任意不是插值节点的 $x\in[a,b]$，存在位于 $x,x_0,\ldots,x_n$ 的最小值与最大值之间的 $\xi$，使
\begin{equation}
    f(x)-L_n(x)=\frac{f^{(n+1)}(\xi)}{(n+1)!}\omega_{n+1}(x),
    \qquad \omega_{n+1}(x)=\prod_{i=0}^n(x-x_i).
    \label{eq:lagrange-interpolation-remainder}
\end{equation}
当 $x$ 是节点时误差为零。若 $|f^{(n+1)}|\leq M$，则全部 $x\in[a,b]$ 均满足
\[
    |f(x)-L_n(x)|\leq\frac{M}{(n+1)!}
        \prod_{i=0}^n|x-x_i|.
\]
\end{theorem}

\begin{proof}
固定非节点 $x$，令 $K=(f(x)-L_n(x))/\omega_{n+1}(x)$，并构造关于 $t$ 的函数
\[
    F(t)=f(t)-L_n(t)-K\omega_{n+1}(t).
\]
$F$ 在互不相同的 $x,x_0,\ldots,x_n$ 上都为零，故连续使用 $n+1$ 次罗尔定理，存在严格位于这些点的最小值与最大值之间的 $\xi$，使 $F^{(n+1)}(\xi)=0$。由于 $L_n^{(n+1)}=0$，而首一多项式 $\omega_{n+1}$ 的 $n+1$ 阶导数为 $(n+1)!$，得 $K=f^{(n+1)}(\xi)/(n+1)!$，代回即得余项公式。节点处单独使用插值条件，不能在那里直接除以 $\omega_{n+1}(x)$。
\end{proof}

\begin{example}
\label{ex:quadratic-interpolation-data}
给定三组数据 $(1,0.0000)$、$(1.5,0.4055)$、$(2,0.6931)$，求相应的二次插值多项式，并估算 $f(1.2)$。
\end{example}

\begin{solve}
第一组数据的函数值为零，故
\[
    L_2(x)=-1.6220(x-1)(x-2)+1.3862(x-1)(x-1.5),
\]
从而 $L_2(1.2)=0.176348$，按四位小数可写为 $f(1.2)\approx0.1763$。这个数是根据所给数据得到的插值估计；仅有三组数据不能给出真实函数的误差保证。若这些数据还经过舍入，则总误差除式~\eqref{eq:lagrange-interpolation-remainder} 的函数插值误差外，还包括数据误差。具体地，节点数据误差为 $\varepsilon_i$ 时，多项式的变化为 $\sum_i\varepsilon_i\ell_i(x)$，其绝对值不超过 $\sum_i|\varepsilon_i\ell_i(x)|$。
\end{solve}

\section{利用导数研究函数}
\label{sec:studying-functions-with-derivatives}

% 教材第 210--229 页；分章 PDF 第 34--53 页。
中值定理把导数的符号转化为函数增量的符号，泰勒公式则用局部主项决定函数在一点附近的行为。因此，研究函数时应先确定定义域和连续性，再在适用的区间上考察一阶、二阶或更高阶导数；跨越不连续点或定义域空缺时，不能直接沿用同一个区间上的判别。

\subsection{函数的单调性}
\label{subsec:derivative-tests-for-monotonicity}

\begin{theorem}[单调性的导数判别]
\label{thm:derivative-monotonicity-test}
设 $I$ 为非退化区间，$f$ 在 $I$ 上连续、内部可导。则 $f$ 在 $I$ 上单调递增（即非减）的充分必要条件是内部恒有 $f'\geq0$；单调递减（即非增）的充分必要条件是内部恒有 $f'\leq0$。
\end{theorem}

\begin{proof}
若 $f$ 非减，则任意内点 $x$ 处充分小的非零增量 $h$ 满足 $(f(x+h)-f(x))/h\geq0$，取极限得 $f'(x)\geq0$。反之，若内部 $f'\geq0$，则对任意 $u<v$、$u,v\in I$，由拉格朗日定理，
\[
    f(v)-f(u)=f'(\xi)(v-u)\geq0,
\]
故 $f$ 非减。对 $-f$ 使用该结论即可得到非增情形。
\end{proof}

\begin{corollary}[严格单调的判据]
\label{cor:strict-monotonicity-test}
在定理~\ref{thm:derivative-monotonicity-test} 的条件下，$f$ 严格递增的充分必要条件是 $f'\geq0$，且 $f'$ 在内部任意非退化子区间上都不恒为零。相应地，严格递减等价于 $f'\leq0$ 且 $f'$ 在任何此类子区间上不恒为零。特别地，内部处处有 $f'>0$ 或 $f'<0$ 是相应严格单调性的充分条件。
\end{corollary}

\begin{proof}
严格递增必非减，故 $f'\geq0$；若 $f'$ 在某个子区间上恒为零，则 $f$ 在其上为常数，矛盾。反之，由 $f'\geq0$ 知 $f$ 非减，若存在 $u<v$ 而 $f(u)=f(v)$，则对 $u<t<v$ 有 $f(u)\leq f(t)\leq f(v)$，故 $f$ 在 $(u,v)$ 上为常数，违反第二个条件。严格递减的情形同理。
\end{proof}

$f(x)=x^3$ 在实轴上严格递增，但 $f'(0)=0$，说明严格递增不等于导数处处为正。利用单调性也可证明不等式：对 $0<x<\pi/2$，$x-\sin x$ 的导数为 $1-\cos x>0$，$\tan x-x$ 的导数为 $\tan^2x>0$，二者在零点的值均为零，因此 $\sin x<x<\tan x$。这种方法的关键是把待证不等式改写成一个函数与其已知初值的比较。

\subsection{函数的极值}
\label{subsec:extrema-of-functions}

费马定理只筛选可导极值点的候选位置；要确定候选点是否真正取得极值，还需比较其两侧函数值。最直接的方法是通过导数符号确定两侧的升降。

\begin{theorem}[极值的第一导数判别]
\label{thm:first-derivative-extremum-test}
设 $f$ 在 $x_0$ 处连续，在某个去心邻域 $0<|x-x_0|<\delta$ 内可导。若 $f'$ 在 $x_0$ 左侧为正、右侧为负，则 $x_0$ 为严格极大值点；若左侧为负、右侧为正，则为严格极小值点；若两侧均为正，或两侧均为负，则 $x_0$ 不是极值点。这里不要求 $f'(x_0)$ 存在。
\end{theorem}

\begin{proof}
对邻域内任意 $x<x_0$，可在 $[x,x_0]$ 上使用中值定理；对 $x>x_0$，可在 $[x_0,x]$ 上使用。若左侧导数为正、右侧为负，则分别得到 $f(x)<f(x_0)$，故为严格极大值。反向变号时得到两侧均有 $f(x)>f(x_0)$。两侧同为正时，左侧函数值小于 $f(x_0)$、右侧大于 $f(x_0)$；同为负时次序相反，均不可能取得极值。
\end{proof}

本判据的同号、异号描述指在某个完整的单侧邻域内保持所列符号，并不保证导数任意振荡时也可用简单的符号表判定。连续性也不能省略，否则 $x_0$ 的函数值可能脱离两侧的升降趋势。

\begin{example}
\label{ex:polynomial-stationary-point-classification}
求 $f(x)=(x+2)^2(x-1)^3$ 在实轴上的极值。
\end{example}

\begin{solve}
求导得
\[
    f'(x)=(x+2)(x-1)^2(5x+4).
\]
驻点为 $-2,-4/5,1$。导数在由它们分出的四个开区间上依次为正、负、正、正，故 $x=-2$ 为严格极大值点，极大值为 $0$；$x=-4/5$ 为严格极小值点，极小值为
\[
    \left(\frac65\right)^2\left(-\frac95\right)^3
    =-\frac{26244}{3125}.
\]
$x=1$ 两侧导数均为正，不是极值点。
\end{solve}

\begin{theorem}[极值的高阶导数判别]
\label{thm:higher-derivative-extremum-test}
设 $n\geq2$，$f$ 在 $x_0$ 处具有 $n$ 阶导数，且
\[
    f'(x_0)=\cdots=f^{(n-1)}(x_0)=0,
    \qquad f^{(n)}(x_0)\neq0.
\]
若 $n$ 为偶数，则 $f^{(n)}(x_0)>0$ 时取得严格极小值，$f^{(n)}(x_0)<0$ 时取得严格极大值；若 $n$ 为奇数，则 $x_0$ 不是极值点。特别地，$f'(x_0)=0$ 且 $f''(x_0)\neq0$ 时，二阶导数的正、负分别对应严格极小、极大值。
\end{theorem}

\begin{proof}
由佩亚诺公式，
\[
    f(x)-f(x_0)
    =\left(\frac{f^{(n)}(x_0)}{n!}+o(1)\right)(x-x_0)^n.
\]
括号内的量在充分小的去心邻域内与 $f^{(n)}(x_0)$ 同号。偶次幂在两侧均为正，故差值保持相应符号；奇次幂在左右两侧异号，故差值变号，不可能取得极值。
\end{proof}

若 $f''(x_0)=0$，二阶判据没有结论：$x^4,-x^4,x^3$ 在零点分别有极小值、极大值和无极值。若所有阶导数都为零，高阶判据也可能无法使用，此时仍可回到定义或考察导数符号。

\begin{theorem}[达布定理]
\label{thm:darboux-derivative-intermediate-value}
设 $f$ 在 $[a,b]$ 上连续、在 $(a,b)$ 内可导，有限单侧导数 $f'_+(a)$、$f'_-(b)$ 存在。则对严格介于这两个导数之间的任意实数 $\lambda$，存在 $\xi\in(a,b)$，使 $f'(\xi)=\lambda$。在开区间内任取两个点作端点，便知\emph{导函数具有介值性}；这一结论称为\emph{达布定理}（Darboux's theorem）。
\end{theorem}

\begin{proof}
令 $F(x)=f(x)-\lambda x$。若 $f'_+(a)<\lambda<f'_-(b)$，则 $F'_+(a)<0$，故在 $a$ 的右侧存在函数值严格小于 $F(a)$ 的点；又 $F'_-(b)>0$，故在 $b$ 的左侧存在函数值严格小于 $F(b)$ 的点。因此连续函数 $F$ 在闭区间上取得的最小值不可能位于任一端点，必在内点 $\xi$ 取得。费马定理给出 $F'(\xi)=0$，即 $f'(\xi)=\lambda$。若端点导数的大小次序相反，则改取内点最大值，证明相同。
\end{proof}

达布定理不要求导函数连续。例如 $f(x)=x^2\sin(1/x)$（$x\neq0$）、$f(0)=0$ 在实轴上处处可导，但
\[
    f'(x)=2x\sin(1/x)-\cos(1/x)\quad(x\neq0),\qquad f'(0)=0,
\]
在零点不连续，仍然具有介值性。因此若函数在一个区间内处处可导且 $f'$ 从不为零，则 $f'$ 不能由正变负，必保持同一个符号，函数便严格单调；不能仅凭“导函数连续”来解释这一事实，因为连续性未必成立。

局部极值与整个定义范围上的最大、最小值需要区分。若 $f\in C([a,b])$，则最值存在，且最值点必属于端点、内部驻点或内部不可导点；这是由费马定理排除其余内点得到的。求最值时应检查全部候选点，而不能只解 $f'=0$；若候选点无穷多，还须进一步分析，不能预设总能作有限次比较。

\begin{example}
\label{ex:weighted-cylinder-material-minimization}
在薄壁近似下制作无盖圆柱形容器，设侧壁厚度为固定的 $\delta>0$，底厚为 $3\delta$，给定内部容积 $V>0$。忽略边缘及厚度对半径、容积关系的修正，求用料最省的底面半径与高度关系。
\end{example}

\begin{solve}
设底面半径为 $r>0$、高度为 $h>0$，以 $V=\pi r^2h$ 为体积约束。用料体积在此模型中为
\[
    S=3\delta\pi r^2+2\delta\pi rh,
\]
故消去 $h$ 后，问题化为在 $r>0$ 上最小化
\[
    S(r)=\delta\left(3\pi r^2+\frac{2V}{r}\right),\qquad
    S'(r)=\delta\left(6\pi r-\frac{2V}{r^2}\right).
\]
唯一驻点为 $r_0=(V/(3\pi))^{1/3}$，其左侧导数为负、右侧为正，故是整个正半轴上的唯一最小值点。相应高度满足
\[
    h_0=\frac{V}{\pi r_0^2}=3r_0.
\]
这一结论来自所列薄壁模型；改变底与侧壁的厚度比例时，目标函数及最优比例也会随之改变。
\end{solve}

\subsection{函数的凹凸性}
\label{subsec:convexity-and-concavity}

单调性描述函数值怎样变化，凹凸性则描述割线和切线斜率怎样变化。本节沿用教材及凸分析中的约定：\emph{图像位于弦线下方的函数称为凸函数}，以免与仅凭“向上凹”“向下凹”等图形用语产生混淆。

\begin{definition}
\label{def:convex-and-concave-functions}
设 $I$ 为区间，$f:I\to\mathbb{R}$。若对任意 $x,y\in I$ 和 $t\in(0,1)$ 均有
\begin{equation}
    f\bigl(tx+(1-t)y\bigr)\leq tf(x)+(1-t)f(y),
    \label{eq:convex-function-definition}
\end{equation}
则称 $f$ 为 $I$ 上的\emph{凸函数}（convex function）；若 $x\neq y$ 时不等式总严格成立，则称为\emph{严格凸函数}（strictly convex function）。将不等号反向，得到\emph{凹函数}（concave function）及严格凹函数。$f$ 为凹函数等价于 $-f$ 为凸函数。
\end{definition}

仿射函数 $f(x)=Ax+B$ 同时为凸函数和凹函数；反之，在非退化区间上同时凸、凹的函数，由式~\eqref{eq:convex-function-definition} 处处取等号可知，其任意三点共线，故必为仿射函数。二次函数 $ax^2+bx+c$ 在 $a>0$ 时严格凸，在 $a<0$ 时严格凹。定义本身不以可导为前提，例如 $|x|$ 是实轴上的凸函数，却在零点不可导。

\begin{lemma}[割线斜率的刻画]
\label{lem:convex-secant-slopes}
$f$ 在区间 $I$ 上为凸函数，当且仅当对任意 $u<v<w$、$u,v,w\in I$，均有
\begin{equation}
    \frac{f(v)-f(u)}{v-u}
    \leq\frac{f(w)-f(u)}{w-u}
    \leq\frac{f(w)-f(v)}{w-v}.
    \label{eq:convex-secant-slopes}
\end{equation}
严格凸对应于上述不等式对任意这类三点均严格成立。
\end{lemma}

\begin{proof}
将 $v$ 写为 $v=\frac{w-v}{w-u}u+\frac{v-u}{w-u}w$。代入凸函数定义，移项即得式~\eqref{eq:convex-secant-slopes} 中的两个不等式；每一步都是等价变形，故也给出充分性及严格情形。
\end{proof}

\begin{theorem}[凹凸性的导数判别]
\label{thm:convexity-derivative-tests}
设 $f\in C([a,b])$，且在 $(a,b)$ 内可导。则 $f$ 在 $[a,b]$ 上为凸函数，当且仅当 $f'$ 在 $(a,b)$ 内非减；$f$ 严格凸，当且仅当 $f'$ 严格递增。若 $f$ 在 $(a,b)$ 内二阶可导，则凸性等价于 $f''\geq0$；严格凸等价于 $f''\geq0$，且 $f''$ 在任意非退化子区间上都不恒为零。凹性和严格凹性分别对应于导数非增、严格递减，或相应的二阶导数反向不等式。
\end{theorem}

\begin{proof}
若 $f$ 凸，对任意内点 $u<v$，由割线斜率的比较并令相邻端点趋向 $u,v$，得
\begin{equation}
    f'(u)\leq\frac{f(v)-f(u)}{v-u}\leq f'(v),
    \label{eq:convex-derivative-secant-bounds}
\end{equation}
因此 $f'$ 非减。若 $f$ 严格凸而某两个内点有 $f'(u)=f'(v)$，则非减性迫使 $f'$ 在 $[u,v]$ 上为常数，中值定理又迫使 $f$ 在其上为仿射函数，与严格凸矛盾，故 $f'$ 严格递增。

反之，对任意 $u<v<w$、$u,v,w\in[a,b]$，在 $[u,v]$ 和 $[v,w]$ 上分别使用中值定理，得到两个相邻割线斜率为 $f'(\xi_1)$ 与 $f'(\xi_2)$，其中 $\xi_1<v<\xi_2$。若 $f'$ 非减，则前者不大于后者，由引理~\ref{lem:convex-secant-slopes} 得凸性；若 $f'$ 严格递增，则得到严格凸性。二阶导数的判据再对 $f'$ 使用定理~\ref{thm:derivative-monotonicity-test} 及推论~\ref{cor:strict-monotonicity-test} 即得。
\end{proof}

在二阶可导情形，$f''>0$ 是严格凸的充分条件，却不是必要条件；例如 $x^4$ 严格凸，但在零点有 $f''(0)=0$。另一个等价刻画是切线支撑不等式：在上述可导条件下，凸性等价于对任意内点 $u$ 和任意 $x\in[a,b]$ 有
\begin{equation}
    f(x)\geq f(u)+f'(u)(x-u).
    \label{eq:convex-supporting-tangent}
\end{equation}
凸性推出该式可由割线斜率取极限得到。反之，令 $z=tx+(1-t)y$ 为内点，将以 $z$ 为切点、分别取 $x,y$ 的两个不等式加权相加，便得到式~\eqref{eq:convex-function-definition}。因此凸函数的图像位于每条切线上方，同时位于相应的弦线下方。

\begin{theorem}[詹森不等式]
\label{thm:jensen-finite-convex-combinations}
设 $f$ 为区间 $I$ 上的凸函数，$x_1,\ldots,x_n\in I$，权重 $t_i\geq0$ 满足 $\sum_{i=1}^n t_i=1$。则
\begin{equation}
    f\left(\sum_{i=1}^n t_ix_i\right)
    \leq\sum_{i=1}^n t_if(x_i).
    \label{eq:jensen-finite-convex-combinations}
\end{equation}
这称为\emph{詹森不等式}（Jensen's inequality）。若 $f$ 严格凸，则等号成立当且仅当所有正权重所对应的 $x_i$ 都相等。凹函数满足反向不等式。
\end{theorem}

\begin{proof}
删除零权重项，以下设权重均为正。$n=1$ 无须证明，$n=2$ 为凸性的定义。若结论对 $n-1$ 成立，令 $s=1-t_n$，则当 $n\geq2$ 时 $0<s<1$，且
\[
    \sum_{i=1}^nt_ix_i
    =s\left(\sum_{i=1}^{n-1}\frac{t_i}s x_i\right)+t_nx_n.
\]
先对这两个点使用凸性，再对括号内的凸组合使用归纳假设，得到所需不等式。严格凸时，若前 $n-1$ 个点不全相等，归纳步骤已经严格；若它们全相等但与 $x_n$ 不同，最外层的二点不等式严格。因此只有全部正权重点相等时才能取等号。凹函数对 $-f$ 使用结论即可。
\end{proof}

\begin{example}
\label{ex:power-mean-monotonicity-by-convexity}
对例~\ref{ex:power-means-parameter-limits} 的幂平均，证明 $M_r$ 关于实参数 $r$ 非减；若 $a_i$ 不全相等，则严格递增，并导出调和、几何与算术平均不等式。
\end{example}

\begin{solve}
若 $a_i$ 全相等，则所有幂平均都等于这个共同值。以下设不全相等，记 $S_r=\sum_i a_i^r$。当 $r\neq0$ 时，直接求导得
\[
    \frac{d}{dr}\ln M_r
    =\frac1{r^2}\left(
       r\frac{\sum_i a_i^r\ln a_i}{S_r}-\ln\frac{S_r}n\right).
\]
令 $u_i=na_i^r/S_r$，则 $u_i>0$、$\frac1n\sum_i u_i=1$，且因 $r\neq0$，这些 $u_i$ 不全相等。函数 $\varphi(u)=u\ln u$ 在 $(0,+\infty)$ 上满足 $\varphi''(u)=1/u>0$，为严格凸函数。由严格詹森不等式，
\[
    r\frac{\sum_i a_i^r\ln a_i}{S_r}-\ln\frac{S_r}n
    =\frac1n\sum_i u_i\ln u_i>\varphi(1)=0.
\]
因此 $M_r$ 在两个半轴上分别严格递增。由式~\eqref{eq:power-mean-zero-parameter} 的连续补定义，严格递增性也跨过 $r=0$：例如对 $r<0$，有 $M_r<M_{r/2}\leq M_0$，正半轴同理。

取 $r=-1,0,1$ 并结合式~\eqref{eq:power-mean-infinite-parameter-limits}，得到
\begin{equation}
    \min_i a_i
    \leq\frac{n}{\sum_i1/a_i}
    \leq\left(\prod_i a_i\right)^{1/n}
    \leq\frac1n\sum_i a_i
    \leq\max_i a_i.
    \label{eq:harmonic-geometric-arithmetic-mean-chain}
\end{equation}
当 $a_i$ 不全相等时，三个平均之间均严格不等；最小、最大值与这些平均之间也严格不等，因为有限参数下正数的平均只有在各项相等时才达到相应极端值。全相等时各式同时取等号。
\end{solve}

\begin{example}
\label{ex:holder-inequality-from-concavity}
设 $a_i,b_i>0$（$1\leq i\leq n$），$p,q>1$ 且 $1/p+1/q=1$。证明\emph{赫尔德不等式}（H\"older's inequality）
\begin{equation}
    \sum_{i=1}^n a_ib_i
    \leq\left(\sum_{i=1}^n a_i^p\right)^{1/p}
         \left(\sum_{i=1}^n b_i^q\right)^{1/q},
    \label{eq:holder-inequality-finite}
\end{equation}
并确定等号条件。
\end{example}

\begin{solve}
记 $A=\sum_i a_i^p$、$B=\sum_i b_i^q$，取权重 $t_i=b_i^q/B$，以及 $z_i=a_i^p/b_i^q$。函数 $\varphi(z)=z^{1/p}$ 在正半轴上满足 $\varphi''(z)<0$，因而严格凹。对它使用反向詹森不等式，得
\[
    \frac1B\sum_i a_ib_i
    =\sum_i t_i z_i^{1/p}
    \leq\left(\sum_i t_i z_i\right)^{1/p}
    =\left(\frac AB\right)^{1/p},
\]
其中使用了 $q-q/p=1$。两边乘以 $B$，再由 $1-1/p=1/q$ 得到式~\eqref{eq:holder-inequality-finite}。等号成立当且仅当所有 $z_i$ 相等，即存在常数 $c>0$，使 $a_i^p=cb_i^q$ 对全部 $i$ 成立。取 $p=q=2$ 即得柯西--施瓦茨不等式；先对正数使用结论，再由连续性允许零项，并以绝对值替代实数项，可得常用的实向量形式
\[
    \left|\sum_i a_ib_i\right|
    \leq\left(\sum_i a_i^2\right)^{1/2}
         \left(\sum_i b_i^2\right)^{1/2}.
    \qedhere
\]
\end{solve}

\subsection{拐点}
\label{subsec:inflection-points}

\begin{definition}
\label{def:inflection-point}
设 $f$ 在 $x_0$ 的某邻域内连续。若在 $x_0$ 的某个左、右邻域内，函数分别为凸、凹，或分别为凹、凸，并在 $x_0$ 处实际发生凹凸性的改变，则称 $(x_0,f(x_0))$ 为图像的\emph{拐点}（inflection point），也常简称 $x_0$ 为拐点横坐标。这里“实际改变”要求对任意充分小的 $\eta>0$，$f$ 在 $(x_0-\eta,x_0)$ 和 $(x_0,x_0+\eta)$ 上均不恒为仿射函数，以排除没有实际凹凸变化的退化情形；一侧严格凸、另一侧严格凹时自然满足此要求。
\end{definition}

定义关注的是凹凸性的改变，本身不要求二阶导数存在。若 $f$ 在去心邻域内二阶可导，且 $f''$ 在 $x_0$ 两侧分别严格为正、严格为负，或次序相反，则由凹凸性判别立即知 $x_0$ 对应一个拐点。

\begin{theorem}[拐点的必要条件与一个充分条件]
\label{thm:inflection-derivative-tests}
若 $x_0$ 为拐点横坐标，且 $f''(x_0)$ 存在，则 $f''(x_0)=0$。若 $f$ 在 $x_0$ 的某邻域内二阶可导，且
\[
    f''(x_0)=0,\qquad f'''(x_0)\neq0,
\]
其中三阶导数只需在 $x_0$ 存在，则 $x_0$ 对应一个拐点。
\end{theorem}

\begin{proof}
先证必要性。$f''(x_0)$ 存在保证 $f'$ 在附近有定义并在 $x_0$ 连续。若 $f$ 左凸右凹，则 $f'$ 左侧非减、右侧非增，再结合其在 $x_0$ 的连续性，知 $f'$ 在 $x_0$ 取得局部极大值；左凹右凸时则取得局部极小值。将费马定理用于 $f'$，得到 $f''(x_0)=0$。

对充分性，由 $f''$ 在 $x_0$ 可导且 $f''(x_0)=0$，有
\[
    f''(x)=f'''(x_0)(x-x_0)+o(x-x_0)
           =\bigl(f'''(x_0)+o(1)\bigr)(x-x_0).
\]
由于 $f'''(x_0)\neq0$，$f''$ 在充分小的左右邻域内严格异号，因此凹凸性在该点改变。
\end{proof}

$f''(x_0)=0$ 仍只是候选条件：$x^4$ 在零点二阶导数为零，却始终凸，没有拐点；$x^3$ 在零点确实改变凹凸性，因而有拐点。实立方根函数 $x^{1/3}$ 在零点连续，左侧严格凸、右侧严格凹，也有拐点，但在零点没有有限的一阶导数，更没有二阶导数。因此寻找拐点时，除 $f''=0$ 的点外，还应检查函数连续而二阶导数不存在的点，并逐一核对两侧的凹凸性。导函数取得局部极值是上述必要性证明中的结论，不能不加单侧单调性条件就反过来认定原函数有拐点。

\subsection{渐近线}
\label{subsec:asymptotes}

函数图像的一条分支向无穷远延伸时，若其上的点到某条直线的距离趋于零，则称这条直线为该分支的\emph{渐近线}（asymptote）。这一概念只限制远端行为，不要求图像永远位于直线的一侧，也不禁止二者相交。

\begin{definition}
\label{def:vertical-horizontal-oblique-asymptotes}
若 $x\to x_0+$ 或 $x\to x_0-$ 时有 $|f(x)|\to+\infty$，则直线 $x=x_0$ 称为相应分支的\emph{垂直渐近线}（vertical asymptote）。若在 $x\to+\infty$ 或 $x\to-\infty$ 的某一个过程中有
\begin{equation}
    f(x)-(kx+b)\to0,\qquad k,b\in\mathbb{R},
    \label{eq:nonvertical-asymptote-definition}
\end{equation}
则 $y=kx+b$ 是该方向上的非垂直渐近线；$k=0$ 时称为\emph{水平渐近线}（horizontal asymptote），$k\neq0$ 时称为\emph{斜渐近线}（oblique asymptote）。
\end{definition}

图像上点 $(x,f(x))$ 到 $y=kx+b$ 的距离为 $|f(x)-kx-b|/\sqrt{1+k^2}$，因此式~\eqref{eq:nonvertical-asymptote-definition} 正好等价于几何定义。垂直渐近线通常只需有一侧无穷分支，不要求两侧都发散，更不要求两侧有相同的无穷符号。

\begin{proposition}[非垂直渐近线的计算]
\label{prop:asymptote-slope-and-intercept}
在固定的 $x\to+\infty$ 或 $x\to-\infty$ 方向上，非垂直渐近线存在，当且仅当以下两个极限依次存在且有限：
\begin{equation}
    k=\lim\frac{f(x)}x,\qquad
    b=\lim\bigl(f(x)-kx\bigr).
    \label{eq:asymptote-slope-and-intercept}
\end{equation}
此时渐近线唯一，为 $y=kx+b$；特别地，若 $f(x)\to b\in\mathbb{R}$，则水平渐近线为 $y=b$。
\end{proposition}

\begin{proof}
若式~\eqref{eq:nonvertical-asymptote-definition} 成立，则 $f(x)=kx+b+o(1)$，除以 $x$ 得 $f(x)/x\to k$，再相减得 $f(x)-kx\to b$。反之，第二个极限直接给出 $f(x)-(kx+b)\to0$。两个系数均由极限唯一确定，故渐近线唯一。
\end{proof}

必须分别研究两个无穷方向；例如 $f(x)=\sqrt{x^2+1}$ 在 $+\infty$ 处的渐近线是 $y=x$，在 $-\infty$ 处是 $y=-x$。仅有有限的斜率极限也不足以保证存在渐近线：$f(x)=x+\sqrt{x}$ 在 $x\to+\infty$ 时有 $f(x)/x\to1$，但 $f(x)-x\to+\infty$，故无斜渐近线。另一方面，$f(x)=\sin x/x$ 在 $+\infty$ 处以 $y=0$ 为渐近线，却与它相交无穷多次。

\subsection{函数的作图}
\label{subsec:graphing-functions}

作图的目标是把已经证明的性质组合成一致的图像描述。先求定义域，检查奇偶性、周期性、间断点与边界极限；再用 $f'$ 的零点和不可导点划分区间，确定升降及极值；用 $f''$ 的零点和不存在点划分凹凸区间，核实拐点；最后确定各个方向的渐近线，并计算截距和必要的特殊点。列表时必须把定义域的断点列入分界，函数不存在的点不能作为极值点或拐点。有限个采样点不能代替这些区间性质的证明。

\begin{example}
\label{ex:graphing-gaussian-function}
研究 $f(x)=e^{-x^2}$ 的图像特征。
\end{example}

\begin{solve}
定义域为 $\mathbb{R}$，$f$ 为偶函数，且 $0<f(x)\leq1$，故图像关于纵轴对称，与纵轴交于 $(0,1)$，不与横轴相交。由 $x\to\pm\infty$ 时 $f(x)\to0$，知 $y=0$ 是两个方向的水平渐近线。求导得
\[
    f'(x)=-2xe^{-x^2},\qquad
    f''(x)=(4x^2-2)e^{-x^2}.
\]
因此 $f$ 在 $(-\infty,0)$ 上严格递增，在 $(0,+\infty)$ 上严格递减，$x=0$ 是唯一极大值点，也是全局最大值点，最大值为 $1$。函数在 $(-\infty,-1/\sqrt2)$ 和 $(1/\sqrt2,+\infty)$ 上严格凸，在 $(-1/\sqrt2,1/\sqrt2)$ 上严格凹，两个拐点为
\[
    \left(-\frac1{\sqrt2},e^{-1/2}\right),\qquad
    \left(\frac1{\sqrt2},e^{-1/2}\right).
\]
据此可画出从两端趋近横轴、在中央达到峰值的光滑曲线。零是函数的下确界，但从不取得，因此函数在实轴上没有最小值。
\end{solve}

\begin{example}
\label{ex:graphing-rational-function}
研究 $g(x)=(x-1)^3/(x+1)^2$ 的图像特征。
\end{example}

\begin{solve}
定义域为 $\mathbb{R}\setminus\{-1\}$，横、纵轴截距分别为 $(1,0)$ 与 $(0,-1)$。当 $x\to-1\pm$ 时，分子趋于 $-8$、分母从正侧趋于零，故两侧均有 $g(x)\to-\infty$，垂直渐近线为 $x=-1$。多项式除法给出
\[
    g(x)=x-5+\frac{12x+4}{(x+1)^2},
\]
因此两个无穷方向均以 $y=x-5$ 为斜渐近线。导数为
\[
    g'(x)=\frac{(x-1)^2(x+5)}{(x+1)^3},\qquad
    g''(x)=\frac{24(x-1)}{(x+1)^4}.
\]
一阶导数的符号说明，$g$ 在 $(-\infty,-5)$ 上严格递增，在 $(-5,-1)$ 上严格递减，在 $(-1,+\infty)$ 上严格递增；最后一个区间中 $g'(1)=0$ 只是一个孤立零点，不妨碍严格递增。因此 $x=-5$ 是唯一极大值点，极大值为 $g(-5)=-27/2$；$x=1$ 是驻点，却不是极值点。

二阶导数在 $(-\infty,-1)$ 和 $(-1,1)$ 上为负，在 $(1,+\infty)$ 上为正，故图像在前两个区间严格凹，在最后一个区间严格凸，$(1,0)$ 为具有水平切线的拐点。$x=-1$ 不在定义域内，不能作为拐点。将这些区间性质与渐近线、截距结合，即可确定两个分支的变化趋势；两个分支分别向无穷延伸，函数在整个定义域上没有最大值或最小值。
\end{solve}

\section*{习题五}
\addcontentsline{toc}{section}{习题五}
\fi

\selectchapter{6}
\ifcompilechapter
\chapter{不定积分}
\label{chap:indefinite-integrals}

% 本章正文依据教材第 241--270 页；分章 PDF 第 1--30 页，习题六不在本次整理范围内。
求导从函数得到它的变化率，求原函数则由变化率反求函数。例如，直线运动中由位置 $s(t)$ 求速度是计算 $s'(t)$，由速度 $v(t)$ 求位置则要解 $s'(t)=v(t)$；仅有速度还不能确定初始位置，因此结果中自然出现任意常数。本章先说明原函数及其全体组成的不定积分，再由复合函数和乘积的求导法则导出换元法与分部积分法，最后讨论若干能够系统求解的函数类型。以下计算均在被积函数有定义的实开区间上进行，参数视为常数，根式取所声明的实分支。

\section{原函数与不定积分}
\label{sec:antiderivatives-and-indefinite-integrals}

\subsection{原函数与不定积分的概念}
\label{subsec:concepts-of-antiderivatives-and-indefinite-integrals}

\begin{definition}
\label{def:antiderivative-and-indefinite-integral}
设 $f$ 定义在开区间 $I$ 上。若函数 $F$ 在 $I$ 上可导，且对每个 $x\in I$ 都有 $F'(x)=f(x)$，则称 $F$ 为 $f$ 在 $I$ 上的一个\emph{原函数}（antiderivative）。若这样的函数存在，则称 $f$ 在 $I$ 上的全体原函数为它的\emph{不定积分}（indefinite integral），记为
\[
    \int f(x)\,dx.
\]
其中 $f(x)$ 为被积函数，$x$ 为积分变量，$f(x)\,dx$ 为被积表达式。若讨论闭区间，端点处的导数相应理解为内侧单侧导数；本章主要使用开区间形式。
\end{definition}

\begin{theorem}[原函数之间的关系]
\label{thm:antiderivatives-differ-by-constant}
若 $F$ 是 $f$ 在区间 $I$ 上的一个原函数，则 $f$ 在 $I$ 上的全体原函数恰为
\begin{equation}
    \int f(x)\,dx=F(x)+C,\qquad C\in\mathbb{R}.
    \label{eq:indefinite-integral-family}
\end{equation}
这里 $C$ 称为\emph{积分常数}（constant of integration）。
\end{theorem}

\begin{proof}
对任意常数 $C$，有 $(F+C)'=f$，故 $F+C$ 都是原函数。反过来，若 $G$ 也是原函数，则 $(G-F)'=0$。任取 $x_1<x_2$ 且 $x_1,x_2\in I$，函数 $G-F$ 在 $[x_1,x_2]$ 上连续、在 $(x_1,x_2)$ 内可导，由拉格朗日中值定理得
\[
    (G-F)(x_2)-(G-F)(x_1)=0.
\]
因此 $G-F$ 在整个区间上为常数，结论成立。
\end{proof}

几何上，$F(x)+C$ 的图像由同一条曲线上下平移得到，它们在相同横坐标处具有相同的切线斜率。若再指定 $F(x_0)=y_0$，则积分常数由 $C=y_0-F_0(x_0)$ 唯一确定，其中 $F_0$ 是任取的一个原函数。\emph[caution]{定理中的区间条件保证积分常数在所讨论范围内统一；定义域不连通时，不同区间上的常数可以彼此独立。}例如，$1/x$ 在 $(-\infty,0)$ 和 $(0,\infty)$ 上分别有原函数 $\ln|x|+C_-$ 与 $\ln|x|+C_+$，没有理由要求 $C_-=C_+$。

不定积分与微分互为逆运算，但两种运算顺序的结果应加以区别。对任意一个选定的原函数，均有
\[
    \frac{d}{dx}\left(\int f(x)\,dx\right)=f(x),
    \qquad d\left(\int f(x)\,dx\right)=f(x)\,dx;
\]
而若 $F$ 可导，则
\[
    \int F'(x)\,dx=F(x)+C,
    \qquad \int dF(x)=F(x)+C.
\]
前一组公式表示任取积分族中的一个成员再求导，都会还原被积函数；后一组公式表示由导数只能恢复函数到一个加法常数。\emph[tip]{核查一个不定积分结果的直接方法，是在声明的区间上对结果求导。}原函数的存在性也不能省略：下一章将证明连续函数在区间上必有原函数，而有跳跃间断的函数不可能在跨越跳跃点的区间上处处等于某个导函数，因为导函数具有达布介值性。

\subsection{基本不定积分表和不定积分的线性性质}
\label{subsec:basic-integrals-and-linearity}

将基本求导公式反向使用，便得到基本积分表。下式中的 $C$ 各自独立；实幂公式先取 $x>0$，若某个指数允许在更大的实区间上定义且可导，则可在相应区间上使用。
\begin{align}
    \int x^\alpha\,dx&=\frac{x^{\alpha+1}}{\alpha+1}+C
        &&(\alpha\neq-1),\label{eq:basic-power-integral}\\
    \int\frac{dx}{x}&=\ln|x|+C
        &&(x\neq0),\label{eq:basic-logarithmic-integral}\\
    \int e^x\,dx&=e^x+C,\qquad
    \int a^x\,dx=\frac{a^x}{\ln a}+C
        &&(a>0,\ a\neq1),\notag\\
    \int\cos x\,dx&=\sin x+C,\qquad
    \int\sin x\,dx=-\cos x+C,\notag\\
    \int\frac{dx}{\cos^2x}&=\tan x+C,\qquad
    \int\frac{dx}{\sin^2x}=-\cot x+C,\notag\\
    \int\frac{dx}{1+x^2}&=\arctan x+C,\qquad
    \int\frac{dx}{\sqrt{1-x^2}}=\arcsin x+C
        &&(|x|<1),\notag\\
    \int\sinh x\,dx&=\cosh x+C,\qquad
    \int\cosh x\,dx=\sinh x+C,\notag\\
    \int\frac{dx}{\cosh^2x}&=\tanh x+C.
        \label{eq:basic-hyperbolic-integrals}
\end{align}
其中 $\sinh x=(e^x-e^{-x})/2$、$\cosh x=(e^x+e^{-x})/2$、$\tanh x=\sinh x/\cosh x$。含 $\sin x$ 或 $\cos x$ 的分母不得为零，积分公式分别在去掉这些零点后的每个区间上理解。式~\eqref{eq:basic-power-integral} 不能代入 $\alpha=-1$，这一情形由式~\eqref{eq:basic-logarithmic-integral} 单独处理；$\ln|x|$ 的绝对值使同一表达式可以分别用于正、负半轴。

\begin{proposition}[不定积分的线性性质]
\label{prop:linearity-of-indefinite-integrals}
设 $F'=f$、$G'=g$ 在同一区间 $I$ 上成立，$a,b\in\mathbb{R}$。则
\begin{equation}
    \int\bigl(af(x)+bg(x)\bigr)\,dx=aF(x)+bG(x)+C.
    \label{eq:linearity-of-indefinite-integrals}
\end{equation}
特别地，和、差可以逐项积分，非零常数因子可以移到积分号外。
\end{proposition}

\begin{proof}
由求导的线性性质，$(aF+bG)'=af+bg$，再由定理~\ref{thm:antiderivatives-differ-by-constant} 即得式~\eqref{eq:linearity-of-indefinite-integrals}。
\end{proof}

通常将线性性质简写为 $\int(af+bg)\,dx=a\int f\,dx+b\int g\,dx$，其含义是将右侧积分常数合并为左侧的一个任意常数。若把不定积分严格视为函数族，这一简写在 $a=b=0$ 时需要另加任意常数；直接使用式~\eqref{eq:linearity-of-indefinite-integrals} 则对所有 $a,b$ 都成立。线性性质只适用于线性组合，不能据此将乘积或商的不定积分拆成原函数的乘积或商。

\begin{example}
\label{ex:basic-indefinite-integral-transformations}
求 $\displaystyle\int\frac{x^4}{1+x^2}\,dx$、$\displaystyle\int\tan^2x\,dx$ 和 $\displaystyle\int\frac{dx}{\sin^2 2x}$。
\end{example}

\begin{solve}
积分表的使用往往先要做恒等变形。多项式除法给出
\[
    \frac{x^4}{1+x^2}=x^2-1+\frac1{1+x^2},
\]
故第一个积分为 $x^3/3-x+\arctan x+C$。利用 $\tan^2x=1/\cos^2x-1$，第二个积分为 $\tan x-x+C$。对第三个积分，由 $\sin2x=2\sin x\cos x$ 及 $\sin^2x+\cos^2x=1$，有
\[
    \frac1{\sin^2 2x}=\frac14\left(\frac1{\cos^2x}+\frac1{\sin^2x}\right),
\]
因而
\[
    \int\frac{dx}{\sin^2 2x}=\frac14(\tan x-\cot x)+C.
    \qedhere
\]
\end{solve}

\section{换元法与分部积分法}
\label{sec:substitution-and-integration-by-parts}

复合函数的链式法则对应换元积分，乘积的求导法则对应分部积分。换元法中，变量和微分必须同时替换；分部积分中，则通过重新分配微分使剩余积分更简单。这两类方法的有效性都来自求导公式，因而最终结果可以逐项求导核对。

\subsection{第一换元法}
\label{subsec:first-substitution-method}

\begin{theorem}[第一换元积分法]
\label{thm:first-substitution-indefinite-integral}
设 $F'(u)=f(u)$ 在区间 $J$ 上成立，函数 $u=\varphi(x)$ 在区间 $I$ 上可导，且 $\varphi(I)\subseteq J$。则
\begin{equation}
    \int f(\varphi(x))\varphi'(x)\,dx
       =F(\varphi(x))+C.
    \label{eq:first-substitution-indefinite-integral}
\end{equation}
此法称为\emph{第一换元法}（first substitution method），也常称凑微分法。
\end{theorem}

\begin{proof}
由链式法则，$\bigl(F(\varphi(x))\bigr)'=F'(\varphi(x))\varphi'(x)=f(\varphi(x))\varphi'(x)$，故右侧给出了被积函数的全部原函数。
\end{proof}

令 $u=\varphi(x)$，将 $du=\varphi'(x)\,dx$ 一并代入，公式便写成
\[
    \int f(\varphi(x))\varphi'(x)\,dx
       =\int f(u)\,du=F(u)+C=F(\varphi(x))+C.
\]
这里中间的积分以 $u$ 为变量，最后须代回 $x$。这种写法依据的是复合函数微分公式，无须把 $dx$ 当成可以任意约去的普通数。\emph[tip]{选择内层函数后，要检查剩余因子是否恰好是它的导数，或只差一个常数倍。}第一换元法不要求 $\varphi$ 单调，也不要求它有反函数，例如 $u=x^2$ 可用于整个实轴上的 $\int 2x\cos(x^2)\,dx=\sin(x^2)+C$。

分别取 $u=ax+b$（$a\neq0$）和 $u=f(x)$，可得
\begin{align*}
    \int\frac{dx}{ax+b}&=\frac1a\ln|ax+b|+C,\\
    \int(ax+b)^\alpha\,dx
        &=\frac{(ax+b)^{\alpha+1}}{a(\alpha+1)}+C
        \qquad(\alpha\neq-1),\\
    \int\frac{f'(x)}{f(x)}\,dx&=\ln|f(x)|+C
        \qquad(f(x)\neq0).
\end{align*}
前两式中的 $1/a$ 正是由 $du=a\,dx$ 产生的因子；第三式要求 $f$ 在所讨论区间可导且处处不为零。对于 $a>0$，将积分表中的自变量缩放，还得到
\begin{align}
    \int\frac{dx}{a^2+x^2}&=\frac1a\arctan\frac xa+C,
       \label{eq:scaled-arctangent-integral}\\
    \int\frac{dx}{a^2-x^2}&=\frac1{2a}\ln\left|\frac{a+x}{a-x}\right|+C
       \qquad(x\neq\pm a),\label{eq:difference-of-squares-integral}\\
    \int\frac{dx}{\sqrt{a^2-x^2}}&=\arcsin\frac xa+C
       \qquad(|x|<a).\label{eq:scaled-arcsine-integral}
\end{align}
其中第二式也可先用 $1/(a^2-x^2)=\bigl(1/(a+x)+1/(a-x)\bigr)/(2a)$ 分解后积分；两个途径都要求分别在避开 $\pm a$ 的区间上使用。

\begin{example}
\label{ex:two-substitutions-for-square-root-integral}
求 $\displaystyle\int\frac{dx}{\sqrt{x-x^2}}$，其中 $0<x<1$，并比较两种表达式。
\end{example}

\begin{solve}
配方 $x-x^2=\bigl(1-(2x-1)^2\bigr)/4$，令 $u=2x-1$，得到
\[
    \int\frac{dx}{\sqrt{x-x^2}}=\arcsin(2x-1)+C.
\]
另一方面，将 $\sqrt{x-x^2}$ 写成 $\sqrt{x}\sqrt{1-x}$，由 $dx/\sqrt{x}=2\,d\sqrt{x}$，得
\[
    \int\frac{dx}{\sqrt{x-x^2}}=2\arcsin\sqrt{x}+C.
\]
令 $\theta=\arcsin\sqrt{x}\in(0,\pi/2)$，则 $\sin(2\theta-\pi/2)=2x-1$，且 $2\theta-\pi/2\in(-\pi/2,\pi/2)$，故
\[
    2\arcsin\sqrt{x}=\arcsin(2x-1)+\frac\pi2.
\]
两种结果相差常数，因而表示同一原函数族。
\end{solve}

\begin{example}
\label{ex:tangent-secant-integrals-and-recurrence}
求 $\displaystyle\int\tan x\,dx$、$\displaystyle\int\frac{dx}{\cos x}$，并建立 $T_n=\int\tan^n x\,dx$ 的递推公式。
\end{example}

\begin{solve}
在 $\cos x\neq0$ 的每个区间上，由 $d(\cos x)=-\sin x\,dx$，有
\[
    \int\tan x\,dx=-\ln|\cos x|+C.
\]
对第二个积分，用 $u=\sin x$ 凑微分，得
\[
    \int\frac{dx}{\cos x}
       =\int\frac{d(\sin x)}{1-\sin^2x}
       =\frac12\ln\left|\frac{1+\sin x}{1-\sin x}\right|+C
       =\ln|\sec x+\tan x|+C,
\]
其中 $\sec x=1/\cos x$，最后一步使用 $1-\sin^2x=\cos^2x$。若 $n\geq2$ 为整数，由 $\tan^2x=\sec^2x-1$，得到
\begin{equation}
    T_n=\frac{\tan^{n-1}x}{n-1}-T_{n-2},
    \qquad T_0=x,\quad T_1=-\ln|\cos x|.
    \label{eq:tangent-power-integral-recurrence}
\end{equation}
这里及下文的递推式均按相差一个常数的意义理解，基式取一个代表原函数，最终统一加 $C$。例如 $T_4=\tan^3x/3-\tan x+x+C$。
\end{solve}

\subsection{第二换元法}
\label{subsec:second-substitution-method}

\begin{theorem}[第二换元积分法]
\label{thm:second-substitution-indefinite-integral}
设 $x=\varphi(t)$ 在开区间 $J$ 上可导，且 $\varphi'(t)\neq0$。记 $I=\varphi(J)$，$t=\psi(x)$ 为其反函数。若 $f$ 在 $I$ 上有定义，且
\[
    \int f(\varphi(t))\varphi'(t)\,dt=G(t)+C,
\]
则
\begin{equation}
    \int f(x)\,dx=G(\psi(x))+C.
    \label{eq:second-substitution-indefinite-integral}
\end{equation}
此法称为\emph{第二换元法}（second substitution method）。
\end{theorem}

\begin{proof}
导函数具有达布介值性，而 $\varphi'$ 从不为零，因此它在 $J$ 上不能变号。由拉格朗日中值定理，$\varphi$ 严格单调；它又连续，故 $I$ 是开区间，反函数 $\psi$ 存在且可导，满足 $\psi'(x)=1/\varphi'(\psi(x))$。于是
\[
    \frac{d}{dx}G(\psi(x))
       =f(\varphi(\psi(x)))\varphi'(\psi(x))
          \frac1{\varphi'(\psi(x))}=f(x).
    \qedhere
\]
\end{proof}

第一换元法从已有的复合结构中提取内层变量，第二换元法则主动选择 $x=\varphi(t)$，把原积分改成较容易处理的 $t$ 的积分，再通过反函数代回。\emph[caution]{第二换元法必须说明所取区间，特别是根式化简时不能漏掉绝对值。}例如 $\sqrt{a^2\cos^2t}=a|\cos t|$，只有先限制 $\cos t>0$ 才能写成 $a\cos t$。

\begin{example}
\label{ex:trigonometric-substitution-for-semicircle}
求 $\displaystyle\int\sqrt{a^2-x^2}\,dx$，其中 $a>0$、$|x|<a$。
\end{example}

\begin{solve}
令 $x=a\sin t$，$-\pi/2<t<\pi/2$。此时 $\cos t>0$，故 $dx=a\cos t\,dt$、$\sqrt{a^2-x^2}=a\cos t$，从而
\[
    \int\sqrt{a^2-x^2}\,dx
      =a^2\int\cos^2t\,dt
      =\frac{a^2}{2}(t+\sin t\cos t)+C.
\]
代回 $t=\arcsin(x/a)$、$\sin t=x/a$ 和 $\cos t=\sqrt{a^2-x^2}/a$，得
\begin{equation}
    \int\sqrt{a^2-x^2}\,dx
      =\frac12\left(x\sqrt{a^2-x^2}+a^2\arcsin\frac xa\right)+C.
    \label{eq:semicircle-antiderivative}
\end{equation}
\end{solve}

\begin{example}
\label{ex:logarithmic-integrals-of-quadratic-radicals}
求 $\displaystyle\int\frac{dx}{\sqrt{x^2-a^2}}$ 和 $\displaystyle\int\frac{dx}{\sqrt{x^2+a^2}}$，其中 $a>0$。
\end{example}

\begin{solve}
先在 $x>a$ 上令 $x=a\sec t$，$0<t<\pi/2$。由 $dx=a\sec t\tan t\,dt$、$\sqrt{x^2-a^2}=a\tan t$，有
\[
    \int\frac{dx}{\sqrt{x^2-a^2}}
       =\int\sec t\,dt
       =\ln|\sec t+\tan t|+C
       =\ln|x+\sqrt{x^2-a^2}|+C.
\]
在 $x<-a$ 上可令 $x=-a\cosh t$（$t>0$），或直接对最后的表达式求导，同样得到
\begin{equation}
    \int\frac{dx}{\sqrt{x^2-a^2}}
       =\ln|x+\sqrt{x^2-a^2}|+C,\qquad |x|>a,
    \label{eq:reciprocal-difference-radical-integral}
\end{equation}
但两侧区间的常数可以不同。对第二个积分，令 $x=a\sinh t$，$t\in\mathbb{R}$，则 $\sqrt{x^2+a^2}=a\cosh t>0$，故积分为 $t+C$。解出 $t=\ln\bigl((x+\sqrt{x^2+a^2})/a\bigr)$，并将 $-\ln a$ 并入常数，得
\begin{equation}
    \int\frac{dx}{\sqrt{x^2+a^2}}
       =\ln\bigl(x+\sqrt{x^2+a^2}\bigr)+C,
       \qquad x\in\mathbb{R}.
    \label{eq:reciprocal-sum-radical-integral}
\end{equation}
这里对数的真数恒正，因为 $\sqrt{x^2+a^2}>|x|$。也可用 $x=a\tan t$、$-\pi/2<t<\pi/2$ 将第二个积分化为 $\int\sec t\,dt$。
\end{solve}

上述换元分别利用 $1-\sin^2t=\cos^2t$、$\sec^2t-1=\tan^2t$ 和 $\cosh^2t-\sinh^2t=1$ 消去根号。对同一变量的多个有理次幂，则可按指数分母的最小公倍数选择换元。

\begin{example}
\label{ex:common-denominator-substitution-for-radicals}
求 $\displaystyle\int\frac{dx}{\sqrt{x}+\sqrt[3]{x}}$，其中 $x>0$。
\end{example}

\begin{solve}
令 $x=t^6$，$t>0$，则 $\sqrt{x}=t^3$、$\sqrt[3]{x}=t^2$，因此
\begin{align*}
    \int\frac{dx}{\sqrt{x}+\sqrt[3]{x}}
       &=6\int\frac{t^3}{1+t}\,dt
        =6\int\left(t^2-t+1-\frac1{1+t}\right)\,dt\\
       &=2t^3-3t^2+6t-6\ln(1+t)+C\\
       &=2\sqrt{x}-3\sqrt[3]{x}+6\sqrt[6]{x}
           -6\ln\bigl(1+\sqrt[6]{x}\bigr)+C.
\end{align*}
这表明第二换元法除三角换元外，也可通过统一根指数将无理表达式化成有理函数。
\end{solve}

\subsection{分部积分法}
\label{subsec:integration-by-parts}

\begin{theorem}[分部积分公式]
\label{thm:integration-by-parts-indefinite}
设 $u,v$ 在区间 $I$ 上可导，且 $u'v$ 在 $I$ 上有原函数，则 $uv'$ 也有原函数，并有
\begin{equation}
    \int u(x)v'(x)\,dx=u(x)v(x)-\int u'(x)v(x)\,dx.
    \label{eq:integration-by-parts-indefinite}
\end{equation}
这称为\emph{分部积分法}（integration by parts），通常简记为
\[
    \int u\,dv=uv-\int v\,du.
\]
式中的任意常数统一包含在右侧的不定积分中。
\end{theorem}

\begin{proof}
若 $H'=u'v$，由乘积求导公式，$(uv-H)'=u'v+uv'-u'v=uv'$，所以 $uv-H$ 是 $uv'$ 的一个原函数，公式成立。
\end{proof}

分部积分的选择同时涉及两个目标：$dv$ 应当容易求出一个原函数 $v$，而 $v\,du$ 应当比原被积表达式简单，或能与原积分形成可解的线性关系。通常多项式求导会降低次数，对数和反三角函数求导会消去这类函数，指数与正弦、余弦则在反复求导、积分后保持同类形式；这些特点提供选择依据，但不构成无条件适用的固定次序。

\begin{example}
\label{ex:integration-by-parts-basic-examples}
求 $\displaystyle\int\ln x\,dx$（$x>0$）、$\displaystyle\int\arctan x\,dx$ 和 $\displaystyle\int x^2\sin x\,dx$。
\end{example}

\begin{solve}
将 $\ln x$ 看成 $\ln x\cdot1$，取 $u=\ln x$、$dv=dx$，有
\[
    \int\ln x\,dx=x\ln x-\int x\frac{dx}{x}
       =x\ln x-x+C.
\]
类似地，取 $u=\arctan x$、$dv=dx$，有
\[
    \int\arctan x\,dx
       =x\arctan x-\int\frac{x}{1+x^2}\,dx
       =x\arctan x-\frac12\ln(1+x^2)+C.
\]
对最后一个积分，先取 $u=x^2$、$dv=\sin x\,dx$，再对剩余积分分部一次，得到
\begin{align*}
    \int x^2\sin x\,dx
       &=-x^2\cos x+2\int x\cos x\,dx\\
       &=-x^2\cos x+2x\sin x+2\cos x+C.
\end{align*}
多项式的次数每次降低一次，因此这里反复分部能够在有限步后结束。
\end{solve}

\begin{example}
\label{ex:integration-by-parts-radical}
求 $\displaystyle I=\int\sqrt{x^2-a^2}\,dx$，其中 $a>0$、$|x|>a$。
\end{example}

\begin{solve}
取 $u=\sqrt{x^2-a^2}$、$dv=dx$，并将 $x^2$ 写成 $(x^2-a^2)+a^2$，得到
\[
    I=x\sqrt{x^2-a^2}-\int\frac{x^2}{\sqrt{x^2-a^2}}\,dx
      =x\sqrt{x^2-a^2}-I-a^2\int\frac{dx}{\sqrt{x^2-a^2}}.
\]
在相差常数的意义下移项，再使用式~\eqref{eq:reciprocal-difference-radical-integral}，得
\begin{equation}
    \int\sqrt{x^2-a^2}\,dx
      =\frac12\left(x\sqrt{x^2-a^2}
             -a^2\ln|x+\sqrt{x^2-a^2}|\right)+C.
    \label{eq:difference-radical-integral}
\end{equation}
这类计算并不是把未知积分当成已知量，而是通过分部积分建立了关于原函数的线性关系。
\end{solve}

\begin{example}
\label{ex:cyclic-exponential-trigonometric-integrals}
求 $\displaystyle I=\int e^{ax}\cos bx\,dx$ 与 $\displaystyle J=\int e^{ax}\sin bx\,dx$，其中 $a,b\in\mathbb{R}$、$a^2+b^2>0$。
\end{example}

\begin{solve}
当 $a\neq0$ 时，对三角函数求导、对指数函数积分，分部公式给出 $aI=e^{ax}\cos bx+bJ$ 和 $aJ=e^{ax}\sin bx-bI$，均在相差常数的意义下成立。联立求解，并将常数合并，得到
\begin{align}
    \int e^{ax}\cos bx\,dx
      &=\frac{e^{ax}}{a^2+b^2}(a\cos bx+b\sin bx)+C,
        \label{eq:exponential-cosine-integral}\\
    \int e^{ax}\sin bx\,dx
      &=\frac{e^{ax}}{a^2+b^2}(a\sin bx-b\cos bx)+C.
        \label{eq:exponential-sine-integral}
\end{align}
当 $a=0$ 时，条件保证 $b\neq0$，直接积分也得到同样的公式。若 $a=b=0$，则两个积分分别为 $x+C$ 与任意常数，应单独处理。
\end{solve}

联立积分的思路也适用于分子与分母有简单导数关系的表达式。记 $D(x)=a\cos x+b\sin x$，其中 $a^2+b^2>0$，在 $D(x)\neq0$ 的区间上令
\[
    I=\int\frac{\sin x}{D(x)}\,dx,\qquad
    J=\int\frac{\cos x}{D(x)}\,dx.
\]
由 $bI+aJ=x$、$-aI+bJ=\ln|D(x)|$（相差常数意义下），得
\[
    I=\frac{bx-a\ln|D(x)|}{a^2+b^2}+C,
    \qquad J=\frac{ax+b\ln|D(x)|}{a^2+b^2}+C.
\]
这里建立关系的依据是分子可分别组合成 $D$ 与 $D'$，无须机械地反复分部。

\begin{proposition}[常用积分递推公式]
\label{prop:power-integral-recurrences}
对整数 $n\geq2$，记 $K_n=\int\cos^n x\,dx$、$S_n=\int\sin^n x\,dx$，则
\begin{align}
    K_n&=\frac{\sin x\cos^{n-1}x}{n}+\frac{n-1}{n}K_{n-2},
       \label{eq:cosine-power-integral-recurrence}\\
    S_n&=-\frac{\sin^{n-1}x\cos x}{n}+\frac{n-1}{n}S_{n-2}.
       \label{eq:sine-power-integral-recurrence}
\end{align}
可取基式 $K_0=S_0=x$、$K_1=\sin x$、$S_1=-\cos x$。此外，对 $a>0$、整数 $n\geq1$，记 $H_n(x)=\int dx/(x^2+a^2)^n$，有
\begin{equation}
    H_{n+1}(x)=\frac{x}{2na^2(x^2+a^2)^n}
          +\frac{2n-1}{2na^2}H_n(x),
    \qquad H_1(x)=\frac1a\arctan\frac xa.
    \label{eq:quadratic-denominator-integral-recurrence}
\end{equation}
这些等式均允许最后添加一个任意常数。
\end{proposition}

\begin{proof}
在 $K_n=\int\cos^{n-1}x\,d(\sin x)$ 中分部，得到
\[
    K_n=\sin x\cos^{n-1}x
          +(n-1)\int\sin^2x\cos^{n-2}x\,dx
       =\sin x\cos^{n-1}x+(n-1)(K_{n-2}-K_n),
\]
移项即得第一式。同理，从 $S_n=-\int\sin^{n-1}x\,d(\cos x)$ 得到
\[
    S_n=-\sin^{n-1}x\cos x+(n-1)(S_{n-2}-S_n),
\]
移项即得第二式。最后，直接计算
\[
    \left(\frac{x}{(x^2+a^2)^n}\right)'
       =\frac{2na^2}{(x^2+a^2)^{n+1}}
           -\frac{2n-1}{(x^2+a^2)^n},
\]
移项并积分，即得式~\eqref{eq:quadratic-denominator-integral-recurrence}。
\end{proof}

例如，$K_2=x/2+\sin2x/4+C$、$S_2=x/2-\sin2x/4+C$，而
\[
    \int\frac{dx}{(x^2+a^2)^2}
       =\frac{x}{2a^2(x^2+a^2)}+\frac1{2a^3}\arctan\frac xa+C.
\]
递推的作用是逐步降低正整数指数，因此需要同时给出起始积分。另一方面，\emph[caution]{存在原函数并不意味着原函数能用初等函数表示。}连续的初等函数在其定义区间上有原函数，但这些原函数未必仍是初等函数；这一存在性结论将在定积分理论中证明，不能因上述技巧未能奏效就断言原函数不存在。

\section{其他类型函数的不定积分}
\label{sec:indefinite-integrals-of-other-functions}

这一节的共同思路是先辨认函数的代数结构，再通过分解或换元化成已知类型。有理函数的不定积分可以由有限个代数步骤求出；三角函数有理式以及某些无理函数则可经适当换元化成有理函数。这里的可求解性指原函数可以用初等函数表示，所采用的表达式仍须在相应的定义区间上理解。

\subsection{有理函数的不定积分}
\label{subsec:integrals-of-rational-functions}

\begin{definition}
\label{def:rational-function-and-proper-fraction}
若 $P,Q$ 是实系数多项式且 $Q$ 不是零多项式，则 $R(x)=P(x)/Q(x)$ 称为\emph{有理函数}（rational function），定义域为 $Q(x)\neq0$ 的点集。约去公因子后，非零分式若满足 $\deg P<\deg Q$，称为\emph{真分式}（proper rational fraction）；否则称为假分式。
\end{definition}

多项式除法将任意有理函数写成 $P/Q=S+T/Q$，其中 $S,T$ 为多项式，$T=0$ 或 $\deg T<\deg Q$。多项式可以逐项积分，因此关键在于处理真分式。根据实系数多项式的因式分解，忽略非零常数因子，分母可写成
\[
    Q(x)=\prod_{i=1}^r(x-a_i)^{m_i}
           \prod_{j=1}^s(x^2+p_jx+q_j)^{n_j},
    \qquad p_j^2-4q_j<0,
\]
其中各一次因子互不相同，各二次因子互不相同，$m_i,n_j$ 为正整数；没有某类因子时，相应乘积省略。

\begin{theorem}[部分分式分解]
\label{thm:real-partial-fraction-decomposition}
若 $T/Q$ 是上述形式的真分式，则存在唯一的实常数 $A_{ik},B_{jk},C_{jk}$，使
\begin{equation}
    \frac{T(x)}{Q(x)}
       =\sum_{i=1}^r\sum_{k=1}^{m_i}\frac{A_{ik}}{(x-a_i)^k}
        +\sum_{j=1}^s\sum_{k=1}^{n_j}
           \frac{B_{jk}x+C_{jk}}{(x^2+p_jx+q_j)^k}.
    \label{eq:real-partial-fraction-decomposition}
\end{equation}
这种表示称为\emph{部分分式分解}（partial fraction decomposition）。
\end{theorem}

\begin{proof}
先说明存在性。若 $U,V$ 为互素多项式，多项式的贝祖恒等式给出 $AU+BV=1$，故
\[
    \frac{T}{UV}=\frac{AT}{V}+\frac{BT}{U}.
\]
逐次使用此式，可将分母拆为各个互素不可约因子的幂；再对分子作除法，并反复按对应的一次或二次因子展开余式，就得到分子次数分别小于 $1$ 或 $2$ 的各项。过程中可能产生多项式项，但原式和所有真分式项在无穷远处都趋于零，因此总多项式项只能为零。这便给出式~\eqref{eq:real-partial-fraction-decomposition}。

为说明唯一性，将两个分解相减并乘以 $Q$。对任一不可约因子 $h$，设它在 $Q$ 中的次数为 $m$；所得多项式恒等式模 $h$ 化简后，只有分母为 $h^m$ 的项可能保留。其余因子与 $h$ 互素，故该项的分子被 $h$ 整除，而其次数小于 $\deg h$，因此分子必为零。将恒等式约去一个 $h$ 后重复此过程，再逐个处理其他因子，得所有系数均为零。
\end{proof}

\emph[caution]{重因子的各次幂必须全部列出；不可约二次因子上方要设一次多项式，不能一律设成常数。}将式~\eqref{eq:real-partial-fraction-decomposition} 通分后比较系数，可得到线性方程组；对一次因子，也可先代入其零点确定最高幂项的系数，再确定其余系数。分解之后只需处理以下四种基本形式：
\[
    \frac{A}{x-a},\qquad \frac{A}{(x-a)^k}\ (k\geq2),\qquad
    \frac{Bx+C}{x^2+px+q},\qquad
    \frac{Bx+C}{(x^2+px+q)^k}\ (k\geq2),
\]
其中 $p^2<4q$。前两类分别给出
\[
    \int\frac{A}{x-a}\,dx=A\ln|x-a|+C_0,
    \qquad
    \int\frac{A}{(x-a)^k}\,dx
       =\frac{A}{1-k}(x-a)^{1-k}+C_0.
\]
这里用 $C_0$ 表示积分常数，以免与分子系数 $C$ 混淆。对后二类，记
\[
    Q_2(x)=x^2+px+q=\left(x+\frac p2\right)^2+\delta^2,
    \qquad \delta=\sqrt{q-\frac{p^2}4}>0,
\]
并将分子写成 $Bx+C=(B/2)Q_2'(x)+(C-Bp/2)$，从而
\begin{equation}
    \int\frac{Bx+C}{Q_2(x)}\,dx
       =\frac B2\ln Q_2(x)
        +\frac{C-Bp/2}{\delta}\arctan\frac{x+p/2}{\delta}+C_0.
    \label{eq:irreducible-quadratic-linear-numerator-integral}
\end{equation}
当 $k\geq2$ 时，同样的分子分解给出
\[
    \int\frac{Bx+C}{Q_2(x)^k}\,dx
      =\frac{B}{2(1-k)}Q_2(x)^{1-k}
        +\left(C-\frac{Bp}2\right)\int\frac{dy}{(y^2+\delta^2)^k},
    \qquad y=x+\frac p2.
\]
最后的积分由式~\eqref{eq:quadratic-denominator-integral-recurrence} 递推到反正切积分。因此，\emph{任何实有理函数在每个无极点区间上的原函数，都可以写成有理函数、对数和反正切的有限组合。}

\begin{example}
\label{ex:partial-fractions-simple-factors}
求 $\displaystyle\int\frac{x}{(x^2+1)(x-1)}\,dx$。
\end{example}

\begin{solve}
设
\[
    \frac{x}{(x^2+1)(x-1)}=\frac{Ax+B}{x^2+1}+\frac D{x-1}.
\]
通分比较系数，得 $A+D=0$、$B-A=1$、$D-B=0$，故 $A=-1/2$、$B=D=1/2$。于是
\begin{align*}
    \int\frac{x}{(x^2+1)(x-1)}\,dx
      &=-\frac12\int\frac{x}{x^2+1}\,dx
          +\frac12\int\frac{dx}{x^2+1}+\frac12\int\frac{dx}{x-1}\\
      &=-\frac14\ln(1+x^2)+\frac12\arctan x
          +\frac12\ln|x-1|+C.
\end{align*}
结果分别在 $(-\infty,1)$ 和 $(1,\infty)$ 上成立。
\end{solve}

\begin{example}
\label{ex:partial-fractions-repeated-factor}
求 $\displaystyle\int\frac{x^3+1}{x(x-1)^3}\,dx$。
\end{example}

\begin{solve}
对重因子 $(x-1)^3$ 列出全部三项，得到分解
\[
    \frac{x^3+1}{x(x-1)^3}
      =-\frac1x+\frac2{x-1}+\frac1{(x-1)^2}+\frac2{(x-1)^3}.
\]
该恒等式可通过通分直接核对，故
\[
    \int\frac{x^3+1}{x(x-1)^3}\,dx
      =-\ln|x|+2\ln|x-1|-\frac1{x-1}-\frac1{(x-1)^2}+C.
\]
这里分母的零点把定义域分为三个区间，积分常数应分别理解。
\end{solve}

\subsection{三角函数有理式的不定积分}
\label{subsec:integrals-of-rational-trigonometric-expressions}

设 $R(u,v)$ 是二元有理函数，即两个二元多项式之比。形如 $R(\sin x,\cos x)$ 的表达式称为\emph{三角函数有理式}（rational expression in trigonometric functions）；由 $\sin x,\cos x,\tan x,\cot x$ 等经有限次四则运算得到的表达式都属此类，而 $\sqrt{1+\sin^2x}$ 或 $\sin(x^2)$ 不属于这一类型。这里始终在原表达式的分母及所用三角函数均有定义的区间上积分。

令
\begin{equation}
    t=\tan\frac x2,\qquad
    \sin x=\frac{2t}{1+t^2},\qquad
    \cos x=\frac{1-t^2}{1+t^2},\qquad
    dx=\frac{2\,dt}{1+t^2}.
    \label{eq:tangent-half-angle-substitution}
\end{equation}
在相邻奇数倍 $\pi$ 之间，这一换元可逆，反函数须按所在区间写成 $x=2\arctan t+2k\pi$，其中 $k\in\mathbb{Z}$。因此
\begin{equation}
    \int R(\sin x,\cos x)\,dx
      =\int R\left(\frac{2t}{1+t^2},\frac{1-t^2}{1+t^2}\right)
             \frac{2\,dt}{1+t^2},
    \label{eq:rationalization-of-trigonometric-integral}
\end{equation}
右端为有理函数的不定积分。式~\eqref{eq:tangent-half-angle-substitution} 称为\emph{万能代换}（tangent half-angle substitution）；它说明三角函数有理式在各定义区间上都有初等原函数。不过，它会在 $x=(2k+1)\pi$ 处引入换元的分界；原被积函数若在这些点仍有定义，所得分段表达式需通过调整积分常数衔接，不能把换元失效误认为原函数在该处不存在。

\begin{example}
\label{ex:tangent-half-angle-integral}
求 $\displaystyle\int\frac{dx}{5+4\sin x}$。
\end{example}

\begin{solve}
令 $t=\tan(x/2)$，则
\[
    \int\frac{dx}{5+4\sin x}
      =\int\frac{2\,dt}{5t^2+8t+5}
      =\frac23\arctan\frac{5t+4}{3}+C.
\]
因此，在每个不跨越奇数倍 $\pi$ 的区间上，可以写成
\[
    \frac23\arctan\frac{5\tan(x/2)+4}{3}+C.
\]
由于 $5+4\sin x\geq1$，原被积函数在实轴上连续；上式中出现的分段来自换元和反正切主值分支，应按前述方式衔接。
\end{solve}

万能代换保证了可处理性，但具体计算常可利用奇偶结构选择更简单的变量。以下等式均指有理函数恒等式，并限于相应分母非零处。若 $R(u,-v)=-R(u,v)$，则可写成 $R(u,v)=vS(u,v^2)$，其中 $S$ 是有理函数；这是因为 $R(u,v)/v$ 关于 $v$ 为偶有理函数，故可写成 $v^2$ 的有理式。令 $t=\sin x$，有
\[
    \int R(\sin x,\cos x)\,dx
       =\int S(t,1-t^2)\,dt.
\]
若 $R(-u,v)=-R(u,v)$，同理可写成 $R(u,v)=uS(u^2,v)$，令 $t=\cos x$，有
\[
    \int R(\sin x,\cos x)\,dx
       =-\int S(1-t^2,t)\,dt.
\]
若 $R(-u,-v)=R(u,v)$，则在 $v\neq0$ 处可表示为 $S(u/v,v^2)$；令 $t=\tan x$，有
\[
    \int R(\sin x,\cos x)\,dx
       =\int S\left(t,\frac1{1+t^2}\right)\frac{dt}{1+t^2}.
\]
最后一条只要求两个变量\emph{同时变号}时不变，不要求对每个变量分别为偶函数。这三种换元的共同作用，是利用一个微分因子或平方关系消去三角函数，从而减少有理式的次数。

\begin{example}
\label{ex:parity-substitutions-for-trigonometric-integrals}
求 $\displaystyle\int\frac{\cos^3x}{1+\sin^2x}\,dx$ 和 $\displaystyle\int\frac{dx}{a^2\sin^2x+b^2\cos^2x}$，其中 $a,b>0$。
\end{example}

\begin{solve}
第一个被积函数关于 $\cos x$ 为奇式，故用 $t=\sin x$，将 $\cos^2x$ 写成 $1-t^2$，得
\[
    \int\frac{\cos^3x}{1+\sin^2x}\,dx
      =\int\frac{1-t^2}{1+t^2}\,dt
      =2\arctan t-t+C
      =2\arctan(\sin x)-\sin x+C.
\]
第二个被积函数同时变号不变，在 $\cos x\neq0$ 的区间上，将分子分母同除以 $\cos^2x$，再令 $t=\tan x$，得
\[
    \int\frac{dx}{a^2\sin^2x+b^2\cos^2x}
       =\int\frac{dt}{a^2t^2+b^2}
       =\frac1{ab}\arctan\left(\frac ab\tan x\right)+C.
\]
原被积函数处处有定义，后一表达式仍须在 $\cos x=0$ 处通过分支常数衔接。两例都比直接使用万能代换简洁。
\end{solve}

\subsection{无理函数的不定积分}
\label{subsec:integrals-of-irrational-functions}

本小节讨论教材中的两类无理函数：由线性分式的有理次根组成的有理式，以及二项式微分式。其核心都是选择新变量，使根式化为整数次幂；这并不意味着任意含根式的积分都能用同样方法求出初等原函数。

先考虑
\begin{equation}
    \int R\left(x,\sqrt[m]{\frac{ax+b}{cx+d}}\right)\,dx,
    \qquad m\in\mathbb{Z}_{>0},\quad ad-bc\neq0,
    \label{eq:linear-fraction-radical-integral}
\end{equation}
其中 $R$ 为二元有理函数。令 $t=\sqrt[m]{(ax+b)/(cx+d)}$，则由 $t^m=(ax+b)/(cx+d)$ 解得
\begin{equation}
    x=\frac{dt^m-b}{a-ct^m},\qquad
    dx=\frac{m(ad-bc)t^{m-1}}{(a-ct^m)^2}\,dt.
    \label{eq:linear-fraction-radical-substitution}
\end{equation}
将式~\eqref{eq:linear-fraction-radical-substitution} 代入式~\eqref{eq:linear-fraction-radical-integral}，得到
\[
    \int R\left(\frac{dt^m-b}{a-ct^m},t\right)
        \frac{m(ad-bc)t^{m-1}}{(a-ct^m)^2}\,dt,
\]
它是有理函数的不定积分。条件 $ad-bc\neq0$ 排除了线性分式为常数的退化情形，也保证逆变换的非退化性。实际应用中还须排除 $cx+d=0$、$a-ct^m=0$ 以及 $R$ 的其他分母零点；偶次根取非负实根，并在所选开区间上固定分支。若 $t=0$ 使逆变换导数为零，可先在其两侧适用的区间上计算，再核查结果能否延拓。

若同时出现同一个线性分式的若干有理次幂，可取这些指数分母的最小公倍数 $M$，令 $t=((ax+b)/(cx+d))^{1/M}$，将它们同时化成 $t$ 的整数次幂。例~\ref{ex:common-denominator-substitution-for-radicals} 中的 $x=t^6$ 正是这一原则的简单情形。

\begin{example}
\label{ex:linear-fraction-radical-rationalization}
求 $\displaystyle\int\sqrt{\frac{x-1}{x+1}}\,dx$，其中 $x>1$。
\end{example}

\begin{solve}
令 $t=\sqrt{(x-1)/(x+1)}$，则 $0<t<1$，且
\[
    x=\frac{1+t^2}{1-t^2},\qquad dx=\frac{4t}{(1-t^2)^2}\,dt.
\]
因此
\begin{align*}
    \int\sqrt{\frac{x-1}{x+1}}\,dx
       &=\int\frac{4t^2}{(1-t^2)^2}\,dt\\
       &=\frac{2t}{1-t^2}-\int\frac{2\,dt}{1-t^2}
        =\frac{2t}{1-t^2}-\ln\frac{1+t}{1-t}+C.
\end{align*}
由 $2t/(1-t^2)=\sqrt{x^2-1}$、$(1+t)/(1-t)=x+\sqrt{x^2-1}$，代回得
\[
    \int\sqrt{\frac{x-1}{x+1}}\,dx
       =\sqrt{x^2-1}-\ln\bigl(x+\sqrt{x^2-1}\bigr)+C.
    \qedhere
\]
\end{solve}

第二类是\emph{二项式微分式}（binomial differential）
\begin{equation}
    x^m(a+bx^n)^p\,dx,
    \qquad m,n,p\in\mathbb{Q},\quad n\neq0,\quad ab\neq0.
    \label{eq:binomial-differential}
\end{equation}
为避免实有理幂的分支歧义，先限制在 $x>0$、$a+bx^n>0$ 的开区间上；其他有定义的实分支需另外核对。令 $t=x^n$，则
\begin{equation}
    \int x^m(a+bx^n)^p\,dx
       =\frac1n\int t^q(a+bt)^p\,dt,
    \qquad q=\frac{m+1}{n}-1.
    \label{eq:binomial-differential-standard-form}
\end{equation}
这里 $q$ 表示换元后 $t$ 的指数；也有资料用 $(m+1)/n$ 作为参数，二者相差 $1$，判定整性时等价。

\begin{theorem}[二项式微分式的初等可积条件]
\label{thm:chebyshev-binomial-integrals}
在式~\eqref{eq:binomial-differential} 的非退化条件下，二项式微分式有初等原函数，当且仅当
\begin{equation}
    p\in\mathbb{Z},\qquad q\in\mathbb{Z},\qquad
    p+q\in\mathbb{Z}
    \label{eq:binomial-elementary-integrability-conditions}
\end{equation}
三者至少有一个成立。这一结论称为\emph{切比雪夫二项式微分式定理}（Chebyshev's theorem on binomial differentials）。
\end{theorem}

\begin{proof}[充分性证明]
若 $p$ 为整数，将 $q$ 写成 $s/k$，其中 $s\in\mathbb{Z}$、$k\in\mathbb{Z}_{>0}$，令 $t=u^k$，则
\[
    t^q(a+bt)^p\,dt=k u^{s+k-1}(a+bu^k)^p\,du,
\]
为有理表达式。若 $q$ 为整数，将 $p$ 写成 $s/k$，令 $a+bt=u^k$，则 $t=(u^k-a)/b$、$dt=(k/b)u^{k-1}\,du$，而 $(a+bt)^p=u^s$，所以也得到有理函数的积分。

最后，若 $r=p+q$ 为整数，仍写 $p=s/k$，把被积式改写为 $t^r(b+a/t)^p\,dt$，令
\[
    b+\frac at=u^k,\qquad
    t=\frac{a}{u^k-b},\qquad
    dt=-\frac{ak u^{k-1}}{(u^k-b)^2}\,du.
\]
于是被积式成为
\[
    -\frac{k a^{r+1}u^{s+k-1}}{(u^k-b)^{r+2}}\,du,
\]
仍为有理表达式。这就证明了三种情形均可化为有理函数积分。必要性的证明超出本章范围，此处引用教材给出的切比雪夫结论。
\end{proof}

若 $a=0$、$b=0$ 或 $n=0$，原式通常直接退化为常数倍的幂函数，应先单独处理，不能套用定理中的非退化判据。三种换元也允许同时适用，应以所得有理式最简单为准。

\begin{example}
\label{ex:binomial-integrability-and-nonintegrability}
比较 $\displaystyle\int\frac{dx}{\sqrt{1+x^2}}$ 与 $\displaystyle\int\frac{dx}{\sqrt{x+x^3}}$ 在 $x>0$ 上的初等可积性。
\end{example}

\begin{solve}
第一个积分中 $m=0$、$n=2$、$p=-1/2$，故 $q=-1/2$、$p+q=-1\in\mathbb{Z}$，符合第三种情形。依照该情形令 $u=\sqrt{1+x^{-2}}>1$，也就是 $x=(u^2-1)^{-1/2}$，则
\[
    \int\frac{dx}{\sqrt{1+x^2}}
      =-\int\frac{du}{u^2-1}
      =\frac12\ln\frac{u+1}{u-1}+C
      =\ln\bigl(x+\sqrt{1+x^2}\bigr)+C.
\]
这与式~\eqref{eq:reciprocal-sum-radical-integral} 在 $a=1$ 时一致。对第二个积分，先写成
\[
    \int x^{-1/2}(1+x^2)^{-1/2}\,dx.
\]
此时 $m=-1/2$、$n=2$、$p=-1/2$，故 $q=-3/4$、$p+q=-5/4$，三者均非整数。由定理的必要性，这个积分不能用初等函数表示；但被积函数在 $(0,\infty)$ 连续，因而仍有原函数。这里区分的是原函数的表达形式，而非它的存在与否。
\end{solve}

\section*{习题六}
\addcontentsline{toc}{section}{习题六}

% 第七至十七章目录根据用户提供的 8 张教材目录图片整理。
\fi

\selectchapter{7}
\ifcompilechapter
\chapter{定积分}
\label{chap:definite-integrals}

本章研究由“分割、近似求和、取极限”得到的累积量。面积、路程和质量虽有不同背景，却具有共同的数学结构：先把整体分成小块，用局部的简单量近似各块，再证明近似和在分割充分细时趋于同一个数。以下首先建立黎曼积分，继而讨论可积性、积分性质与原函数，最后将这一思想用于几何和物理问题。除另有说明外，均设 $a<b$；$C([a,b])$ 表示闭区间上的连续函数类，$R[a,b]$ 表示其上的黎曼可积函数类。
% 教材依据：数学分析/textbooks/数学分析（分章原文）/数学分析-7.pdf，正文第 1--82 页；不含习题七。

\section{定积分的概念与微积分基本定理}
\label{sec:definite-integrals-and-fundamental-theorem}

\subsection{曲边梯形的面积}
\label{subsec:area-of-curvilinear-trapezoids}

设 $f\in C([a,b])$ 且 $f\geq0$，考虑曲线 $y=f(x)$、$x$ 轴与直线 $x=a,x=b$ 围成的曲边梯形。取分点 $a=x_0<x_1<\cdots<x_n=b$，记 $\Delta x_i=x_i-x_{i-1}$，在第 $i$ 个小区间任取 $\xi_i$。以高 $f(\xi_i)$ 的矩形近似相应部分，所得面积近似值为 $\sum_{i=1}^n f(\xi_i)\Delta x_i$。若记
\[
    m_i=\min_{[x_{i-1},x_i]}f,\qquad
    M_i=\max_{[x_{i-1},x_i]}f,
\]
则上下两组矩形夹住曲边梯形，且两组面积之差为 $\sum_i(M_i-m_i)\Delta x_i$。连续函数在闭区间上一致连续，因而只要所有小区间足够短，就能使这一差小于任意给定的正数。这既解释了矩形近似为何有效，也提示我们以矩形和的共同极限刻画面积；这里的图形直观用于引出定义，后面的可积性证明不预先假定该面积已经存在。

\begin{example}
\label{ex:riemann-area-square-function}
用矩形和求曲线 $y=x^2$ 在 $[0,1]$ 上的曲边梯形面积。
\end{example}
\begin{solve}
将 $[0,1]$ 等分为 $n$ 份。由于 $x^2$ 单调增加，任意取点的矩形和均夹在
\[
    \frac1{n^3}\sum_{i=1}^n(i-1)^2
    \quad\text{与}\quad
    \frac1{n^3}\sum_{i=1}^n i^2
\]
之间。利用 $\sum_{i=1}^n i^2=n(n+1)(2n+1)/6$，两端均趋于 $1/3$，故所求面积为 $1/3$。对于任意足够细的分割，同样可用上述上下矩形之差控制误差；等分计算给出了极限的具体值。
\end{solve}

若质点沿直线同向运动，连续速度为 $v(t)\geq0$，则小时间段上的路程近似为 $v(\xi_i)\Delta t_i$，总路程也由同类和式的极限给出。允许改变运动方向时，速度可以取负；速度的积分对应位移，速率 $|v|$ 的积分才对应路程。

\subsection{定积分的定义}
\label{subsec:definition-of-definite-integrals}

\begin{definition}
\label{def:partition-and-riemann-sum}
区间 $[a,b]$ 的一个\emph{分割}（partition）是有限点集
\[
    P:\quad a=x_0<x_1<\cdots<x_n=b,
    \qquad \|P\|=\max_{1\leq i\leq n}\Delta x_i.
\]
$\|P\|$ 称为分割的\emph{网格宽度}（mesh）。对于定义在 $[a,b]$ 上的实值函数 $f$，在每个 $[x_{i-1},x_i]$ 中任取标记点 $\xi_i$，称
\[
    S(f;P,\xi)=\sum_{i=1}^n f(\xi_i)\Delta x_i
\]
为相应的\emph{黎曼和}（Riemann sum）。
\end{definition}

\begin{definition}
\label{def:riemann-definite-integral}
若存在 $I\in\mathbb R$，使对每个 $\varepsilon>0$，都存在 $\delta>0$，对 $[a,b]$ 的任何分割 $P$ 和任何标记点选取 $\xi$，只要 $\|P\|<\delta$，便有
\[
    |S(f;P,\xi)-I|<\varepsilon,
\]
则称 $f$ 在 $[a,b]$ 上\emph{黎曼可积}（Riemann integrable），并将 $I$ 称为它的\emph{定积分}（definite integral），记作
\begin{equation}
    \int_a^b f(x)\,dx=\lim_{\|P\|\to0}S(f;P,\xi).
    \label{eq:riemann-integral-definition}
\end{equation}
其中 $a,b$ 分别是积分下限、上限，$f$ 是被积函数，$x$ 是积分变量。
\end{definition}

定义同时要求\emph[caution]{分割与标记点两方面的任意性}；仅证明某一列特殊黎曼和收敛，尚不足以证明可积。积分变量是求和过程中使用的哑变量，故 $\int_a^b f(x)\,dx=\int_a^b f(t)\,dt$，但更换变量字母与换元公式所做的变量变换是不同的操作。

积分值具有唯一性。若 $I,J$ 都满足定义，对给定 $\varepsilon>0$ 取一份同时满足两个误差要求的足够细分割，则 $|I-J|\leq|I-S|+|S-J|<2\varepsilon$；由 $\varepsilon$ 任意可知 $I=J$。

\subsection{定积分的几何意义}
\label{subsec:geometric-meaning-of-definite-integrals}

对于连续的非负函数，$\int_a^b f(x)\,dx$ 给出相应曲边梯形的面积；对于非正函数，其值等于面积的相反数。一般情况下，定积分表示\emph{有向面积}（signed area），即轴上方的面积减去轴下方的面积。当零点很多时也不必把图形预先分成有限块，可令 $f_+(x)=\max\{f(x),0\}$、$f_-(x)=\max\{-f(x),0\}$，将两部分面积分别表示为 $\int_a^b f_+(x)\,dx$ 和 $\int_a^b f_-(x)\,dx$。

因此 $\int_a^b f(x)\,dx$ 可能为负，也可能因正负抵消而为零；图像与 $x$ 轴之间的总面积应写为 $\int_a^b|f(x)|\,dx$。例如 $\int_{-1}^1x\,dx=0$，而 $\int_{-1}^1|x|\,dx=1$。

\subsection{连续函数的可积性}
\label{subsec:integrability-of-continuous-functions}

\begin{theorem}
\label{thm:continuous-functions-riemann-integrable}
每个 $f\in C([a,b])$ 都在 $[a,b]$ 上黎曼可积。
\end{theorem}
\begin{proof}
记 $f$ 在 $[a,b]$ 上的连续模为
\[
    \omega_f(h)=\sup\{|f(x)-f(y)|:x,y\in[a,b],\ |x-y|\leq h\}.
\]
一致连续性给出 $\omega_f(h)\to0$（$h\to0^+$）。若 $P'$ 是 $P$ 的细分，则任意两个分别对应 $P,P'$ 的黎曼和满足
\[
    |S(f;P,\xi)-S(f;P',\eta)|
    \leq\omega_f(\|P\|)(b-a),
\]
因为细分前后的标记点均位于同一个原小区间内。对任意两个分割 $P,Q$，以它们的公共细分 $P\cup Q$ 作中介，得到
\[
    |S(f;P,\xi)-S(f;Q,\eta)|
    \leq[\omega_f(\|P\|)+\omega_f(\|Q\|)](b-a).
\]
取一列网格宽度趋于零的等分分割及任意标记点，上式说明相应和式是柯西序列，因实数完备性而收敛，设极限为 $I$。再固定任意 $P,\xi$，使该等分序列趋于极限，便有
\[
    |S(f;P,\xi)-I|\leq\omega_f(\|P\|)(b-a).
\]
右边随 $\|P\|\to0$ 趋于零，故 $I$ 满足定积分定义。
\end{proof}

证明中的关键是统一控制所有小区间上的振幅，而非分别对每个点使用连续性。它还给出连续函数矩形求积的误差估计：若 $\|P\|\leq h$，则误差至多为 $(b-a)\omega_f(h)$。

\subsection{微积分基本定理}
\label{subsec:fundamental-theorem-of-calculus}

积分定义提供累积量的理论含义，但直接计算黎曼和通常很繁琐。微积分基本定理把积分与原函数联系起来，从而将求和极限转化为端点值之差。

\begin{theorem}[牛顿--莱布尼茨公式]
\label{thm:newton-leibniz-integrable-antiderivative}
设 $f\in R[a,b]$，$F$ 在 $[a,b]$ 上连续、在 $(a,b)$ 内可导，且 $F'(x)=f(x)$（$a<x<b$），则
\begin{equation}
    \int_a^b f(x)\,dx=F(b)-F(a)
    =\left.F(x)\right|_a^b.
    \label{eq:newton-leibniz-formula}
\end{equation}
这一公式称为\emph{牛顿--莱布尼茨公式}（Newton--Leibniz formula），也是\emph{微积分基本定理}（fundamental theorem of calculus）的一个形式。
\end{theorem}
\begin{proof}
对任意分割 $P$，在每个 $[x_{i-1},x_i]$ 上用拉格朗日中值定理，存在 $\xi_i\in(x_{i-1},x_i)$，使
\[
    F(x_i)-F(x_{i-1})=F'(\xi_i)\Delta x_i=f(\xi_i)\Delta x_i.
\]
将这些等式相加，得到 $F(b)-F(a)=S(f;P,\xi)$。令 $\|P\|\to0$，由 $f$ 可积即得结论。
\end{proof}

定理同时要求可积性与原函数；目前尚未证明任意连续函数都有原函数，这将在第~\ref{subsec:definite-integrals-with-variable-limits} 小节由变限积分构造出来。定积分是一个数，不定积分则表示原函数族，不能混淆两者。

\begin{example}
\label{ex:power-sum-limit-by-integral}
设 $\alpha>0$，求 $\displaystyle\lim_{n\to\infty}\frac{1^\alpha+2^\alpha+\cdots+n^\alpha}{n^{\alpha+1}}$。
\end{example}
\begin{solve}
原式是 $x^\alpha$ 在 $[0,1]$ 上的等分右端点黎曼和：
\[
    \frac1n\sum_{k=1}^n\left(\frac kn\right)^\alpha
    \longrightarrow\int_0^1x^\alpha\,dx
    =\left.\frac{x^{\alpha+1}}{\alpha+1}\right|_0^1
    =\frac1{\alpha+1}.
\]
这里先由连续性确认可积，再识别黎曼和并用原函数求值；若负指数导致函数在零点无界，就须使用下一章的广义积分理论。
\end{solve}

\section{可积性问题}
\label{sec:integrability}

可积性研究要回答哪些函数具有定义~\ref{def:riemann-definite-integral} 中的共同极限。连续性只是充分条件，许多不连续函数仍然可积；达布理论通过上下确界消去标记点的任意性，使这一问题转化为可控制的振幅估计。

\subsection{可积的必要条件}
\label{subsec:necessary-conditions-for-integrability}

\begin{theorem}
\label{thm:riemann-integrable-implies-bounded}
若 $f\in R[a,b]$，则 $f$ 在 $[a,b]$ 上有界。
\end{theorem}
\begin{proof}
若 $f$ 无界，则任意有限分割中至少有一个小区间上的 $f$ 无界，否则有限个局部界会给出整体界。由可积定义取 $\varepsilon=1$ 所对应的 $\delta>0$，固定一份 $\|P\|<\delta$ 的分割，再固定其他小区间的标记点。在无界的那个小区间内改变标记点，可使相应项的绝对值任意大，因而使整个黎曼和的绝对值任意大。这与所有这些黎曼和都位于 $(I-1,I+1)$ 中矛盾。
\end{proof}

\begin{example}
\label{ex:dirichlet-function-not-riemann-integrable}
狄利克雷函数 $D(x)=1$（$x\in\mathbb Q$）、$D(x)=0$（$x\notin\mathbb Q$）在任何非退化闭区间上都不可积。
\end{example}
\begin{solve}
每个正长度区间内都有有理数与无理数。对于同一个分割，全部选有理标记点得到黎曼和 $b-a$，全部选无理标记点得到 $0$。这两个值不随网格变细而接近，故不可能存在共同极限。特别地，在 $[0,1]$ 上，等分分割的右端点和恒为 $1$，仍不能据此判定可积。
\end{solve}

\subsection{达布理论}
\label{subsec:darboux-theory}

\begin{definition}
\label{def:darboux-sums-and-integrals}
设 $f$ 在 $[a,b]$ 上有界。对分割 $P$ 记
\[
    M_i=\sup_{[x_{i-1},x_i]}f,\qquad
    m_i=\inf_{[x_{i-1},x_i]}f,\qquad \omega_i=M_i-m_i.
\]
称 $\omega_i$ 为该小区间上的\emph{振幅}（oscillation），称
\[
    U(f,P)=\sum_{i=1}^nM_i\Delta x_i,\qquad
    L(f,P)=\sum_{i=1}^nm_i\Delta x_i
\]
分别为\emph{达布上和、下和}（upper and lower Darboux sums）。再定义\emph{上积分、下积分}为
\[
    \overline{\int_a^b}f=\inf_P U(f,P),\qquad
    \underline{\int_a^b}f=\sup_P L(f,P),
\]
其中上下确界遍取 $[a,b]$ 的全部分割。
\end{definition}

每个黎曼和都满足 $L(f,P)\leq S(f;P,\xi)\leq U(f,P)$。若 $f$ 不连续，上下确界未必能取到，达布和因而未必是某一组标记点给出的黎曼和，但可以由适当选点的黎曼和任意逼近。

\begin{proposition}
\label{prop:darboux-refinement}
若 $P'$ 是 $P$ 的细分，则
\[
    L(f,P)\leq L(f,P')\leq U(f,P')\leq U(f,P).
\]
对任意两个分割 $P,Q$，都有 $L(f,P)\leq U(f,Q)$，从而下积分不大于上积分。
\end{proposition}
\begin{proof}
在小区间内增加分点后，各新小区间的上确界不超过原上确界，下确界不小于原下确界，而长度之和不变，故细分使上和不增、下和不减。任意 $P,Q$ 的公共细分给出
\[
    L(f,P)\leq L(f,P\cup Q)\leq U(f,P\cup Q)\leq U(f,Q).
\]
再分别取上下确界即得最后的结论。
\end{proof}

\begin{theorem}[达布定理]
\label{thm:darboux-sum-limits}
对 $[a,b]$ 上的任意有界函数 $f$，都有
\[
    \lim_{\|P\|\to0}U(f,P)=\overline{\int_a^b}f,\qquad
    \lim_{\|P\|\to0}L(f,P)=\underline{\int_a^b}f.
\]
\end{theorem}
\begin{proof}
只证上和。记 $m=\inf f,M=\sup f$，$I^*=\overline{\int_a^b}f$。给定 $\varepsilon>0$，选一份固定分割 $Q$，使 $U(f,Q)<I^*+\varepsilon/2$，并设 $Q$ 有 $N$ 个内分点。任取 $P$，加入 $Q$ 的内分点时，只有内部含这些分点的小区间会改变；它们的总长度不超过 $N\|P\|$，故
\[
    0\leq U(f,P)-U(f,P\cup Q)\leq(M-m)N\|P\|.
\]
利用 $U(f,P\cup Q)\leq U(f,Q)$，得到
\[
    0\leq U(f,P)-I^*
    <\frac{\varepsilon}{2}+(M-m)N\|P\|.
\]
取网格足够小即可。$M=m$ 或 $N=0$ 时后一项为零；下和可对 $-f$ 应用同一论证。
\end{proof}

\begin{theorem}[黎曼可积判据]
\label{thm:darboux-integrability-criteria}
设 $f$ 在 $[a,b]$ 上有界，则以下条件等价：$f\in R[a,b]$；上积分等于下积分；对每个 $\varepsilon>0$，存在分割 $P$，使
\begin{equation}
    U(f,P)-L(f,P)=\sum_{i=1}^n\omega_i\Delta x_i<\varepsilon.
    \label{eq:darboux-oscillation-criterion}
\end{equation}
这些条件又等价于：对任意 $\eta>0,\rho>0$，存在分割，使振幅 $\omega_i\geq\eta$ 的那些小区间的长度总和小于 $\rho$。
\end{theorem}
\begin{proof}
若 $f$ 可积且积分为 $I$，则对足够细的分割，所有黎曼和都在 $I$ 附近。利用上下确界定义逐个选点，可使黎曼和任意逼近上和或下和，故它们也都在 $I$ 附近，上下积分相等。反之，若二者同为 $I$，则由定理~\ref{thm:darboux-sum-limits} 及 $L\leq S\leq U$，任何足够细分割的任意黎曼和都趋于 $I$。

若上下积分相等，分别选下和、上和逼近共同值，再取公共细分，便得到式~\eqref{eq:darboux-oscillation-criterion}。反之，任意分割均有
\[
    0\leq\overline{\int_a^b}f-\underline{\int_a^b}f
    \leq U(f,P)-L(f,P),
\]
故任意小的右边迫使两个积分相等。

若振幅和能任意小，取其小于 $\eta\rho$，则 $\omega_i\geq\eta$ 的区间总长度至多为 $\eta^{-1}\sum_i\omega_i\Delta x_i<\rho$。反之，记整体振幅为 $\Omega$，把和式按 $\omega_i<\eta$ 与 $\omega_i\geq\eta$ 分成两部分，便有
\[
    \sum_i\omega_i\Delta x_i
    \leq\eta(b-a)+\Omega\sum_{\omega_i\geq\eta}\Delta x_i
    <\eta(b-a)+\Omega\rho.
\]
先取足够小的 $\eta,\rho$ 即得所需估计；$\Omega=0$ 的常数函数情形直接成立。
\end{proof}

判据只要求\emph[tip]{对每个精度找到一份振幅和足够小的分割}，不必一开始就检查所有分割。例如若 $\|P_n\|\to0$，且这些分割上的所有标记点和一致趋于同一个数，则 $U(f,P_n)-L(f,P_n)\to0$，足以推出可积；这里“所有标记点和”不能替换成“各取一组标记点的和”。

\subsection{可积函数类}
\label{subsec:classes-of-integrable-functions}

\begin{theorem}
\label{thm:finite-discontinuities-integrable}
闭区间上的有界函数若只有有限个间断点，则黎曼可积。
\end{theorem}
\begin{proof}
设 $|f|\leq B$，可取 $B>0$。用有限个小区间包住所有间断点，使它们的总长度小于 $\varepsilon/(4B)$，并使内部间断点落在所选区间的内部。去掉这些区间内部后，剩下有限个闭区间，$f$ 在其上连续，因而可分别细分，使这些部分上的振幅和小于 $\varepsilon/2$。合并分点得到的总振幅和小于
\[
    2B\frac{\varepsilon}{4B}+\frac{\varepsilon}{2}=\varepsilon.
\]
由达布判据可积。
\end{proof}

在有限个点上改变可积函数的值，仍然得到可积函数，而且积分不变。事实上，两函数之差 $h$ 只在有限个点非零，是一个有界函数；用总长度任意小的区间包住这些点，便可使 $|h|$ 的上和任意小，从而 $\int_a^b h=0$。同样的包围方法还表明：有界函数的间断点集若只有有限多个聚点，也可积；先用总长度很小的区间包住聚点，剩余部分的间断点必为有限个，否则紧致性会在该部分产生新的聚点。

\begin{theorem}
\label{thm:monotone-functions-riemann-integrable}
闭区间上的单调函数黎曼可积。
\end{theorem}
\begin{proof}
以单调增加为例，函数由两个端点值夹住，故有界。对于 $n$ 等分分割，$\omega_i=f(x_i)-f(x_{i-1})$，于是
\[
    \sum_i\omega_i\Delta x_i
    =\frac{b-a}{n}\sum_i[f(x_i)-f(x_{i-1})]
    =\frac{(b-a)[f(b)-f(a)]}{n}\longrightarrow0.
\]
单调减少时对 $-f$ 应用同一结论。
\end{proof}

\begin{example}
\label{ex:thomae-function-riemann-integrable}
在 $[0,1]$ 上定义\emph{黎曼函数}（Thomae's function）
\[
    T(x)=
    \begin{cases}
        1/q,&x=p/q\in(0,1]\cap\mathbb Q,\quad p,q\in\mathbb N,\ (p,q)=1,\\
        1,&x=0,\\
        0,&x\notin\mathbb Q.
    \end{cases}
\]
证明它可积且积分为零。
\end{example}
\begin{solve}
给定 $\eta>0$，满足 $T(x)\geq\eta$ 的点只有有限个，因为这些有理数的既约分母满足 $q\leq1/\eta$。用总长度小于任给 $\rho>0$ 的有限个小区间将它们包住，并使这些点不落在内部小区间的边界上。取包含这些端点的分割后，振幅不小于 $\eta$ 的小区间只能出现在被包住的部分，其长度总和小于 $\rho$，故 $T$ 可积。每个小区间均可取无理标记点，所得黎曼和恒为零，故 $\int_0^1T(x)\,dx=0$。每个有理点都因附近存在无理数而不连续；每个无理点则可避开上述有限个高取值点而得到连续性。这说明稠密的间断点并不必然妨碍可积。
\end{solve}

\begin{example}
\label{ex:oscillatory-sine-riemann-integrable}
函数 $g(x)=\sin(1/x)$（$x\neq0$）、$g(0)=0$ 在每个有限闭区间上可积。
\end{example}
\begin{solve}
$|g|\leq1$，且除零点外连续，因此至多有一个间断点，由定理~\ref{thm:finite-discontinuities-integrable} 即可。零点附近振荡很快，但这一部分可被总长度任意小的区间包住，其振幅和贡献因有界性而可控制。
\end{solve}

\begin{theorem}[勒贝格可积性判据]
\label{thm:lebesgue-criterion-riemann-integrability}
设 $f$ 在 $[a,b]$ 上有界，则 $f$ 黎曼可积的充分必要条件是其间断点集 $E$ 为\emph{零测集}（set of measure zero）：对每个 $\varepsilon>0$，都存在有限或可数个开区间 $(\alpha_k,\beta_k)$ 覆盖 $E$，且 $\sum_k(\beta_k-\alpha_k)<\varepsilon$。
\end{theorem}

这里给出该判据作为可积函数类的整体刻画，证明留待实变函数论。有限集、可数集都是零测集，但零测集也可能不可数；不能把“间断点有限”误当作黎曼可积的必要条件。

\section{定积分的性质}
\label{sec:properties-of-definite-integrals}

为统一端点次序，约定 $\int_a^a f(x)\,dx=0$，以及 $\int_b^a f(x)\,dx=-\int_a^b f(x)\,dx$。以下涉及某个区间时均要求被积函数在该区间可积；积分保序等不等式首先按 $a<b$ 陈述，交换上下限时须相应改变不等号方向。

\begin{theorem}[线性、区间可加性与保序性]
\label{thm:basic-properties-riemann-integral}
若 $f,g\in R[a,b]$，$\alpha,\beta\in\mathbb R$，则 $\alpha f+\beta g\in R[a,b]$，且
\begin{equation}
    \int_a^b(\alpha f+\beta g)\,dx
    =\alpha\int_a^b f\,dx+\beta\int_a^b g\,dx.
    \label{eq:integral-linearity}
\end{equation}
若 $a<c<b$，则 $f\in R[a,b]$ 当且仅当它在 $[a,c]$、$[c,b]$ 上都可积，此时
\begin{equation}
    \int_a^b f\,dx=\int_a^c f\,dx+\int_c^b f\,dx.
    \label{eq:integral-interval-additivity}
\end{equation}
若 $f(x)\leq g(x)$ 对所有 $x\in[a,b]$ 成立，则 $\int_a^b f\leq\int_a^b g$。特别地，若 $m\leq f\leq M$，则
\begin{equation}
    m(b-a)\leq\int_a^b f(x)\,dx\leq M(b-a).
    \label{eq:integral-bounds}
\end{equation}
\end{theorem}
\begin{proof}
线性由同一分割、同一标记点下的黎曼和恒等式及极限的线性得到。对区间可加性，若整体可积，则先把 $c$ 加入振幅和足够小的分割，再限制到两个子区间，各自的振幅和都不超过整体振幅和。反之，将两个子区间上振幅和分别小于 $\varepsilon/2$ 的分割合并，便得到整体可积性。在包含 $c$ 的分割上，整体黎曼和恰为两个子区间黎曼和之和，取极限即得式~\eqref{eq:integral-interval-additivity}。保序性来自每一项 $f(\xi_i)\Delta x_i\leq g(\xi_i)\Delta x_i$；取常数函数比较即得上下界估计。
\end{proof}

式~\eqref{eq:integral-interval-additivity} 在 $a,b,c$ 次序任意时仍成立，只须三个端点所涉及的区间上可积，并使用有向积分的约定。

\begin{theorem}[绝对值、乘积与连续复合]
\label{thm:riemann-integrable-closure}
若 $f,g\in R[a,b]$，则 $|f|,fg\in R[a,b]$，且
\begin{equation}
    \left|\int_a^b f(x)\,dx\right|
    \leq\int_a^b|f(x)|\,dx.
    \label{eq:integral-triangle-inequality}
\end{equation}
若 $m\leq f\leq M$，$\varphi$ 在 $[m,M]$ 上连续，则 $\varphi\circ f\in R[a,b]$。因此若再有 $|g(x)|\geq c>0$，则 $f/g\in R[a,b]$。
\end{theorem}
\begin{proof}
由 $||f(x)|-|f(y)||\leq|f(x)-f(y)|$，绝对值函数在每个小区间上的振幅不超过 $f$ 的振幅，故 $|f|$ 可积。对 $-|f|\leq f\leq|f|$ 积分即可得到式~\eqref{eq:integral-triangle-inequality}。设 $|f|,|g|\leq B$，则
\[
    |f(x)g(x)-f(y)g(y)|
    \leq B|f(x)-f(y)|+B|g(x)-g(y)|,
\]
故 $\omega_i(fg)\leq B[\omega_i(f)+\omega_i(g)]$。分别取两函数振幅和足够小的分割，再取公共细分，便知乘积可积。

对于连续复合，$\varphi$ 在紧区间 $[m,M]$ 上有界且一致连续。给定 $\varepsilon>0$，取 $\eta>0$，使 $|u-v|<\eta$ 时 $|\varphi(u)-\varphi(v)|<\varepsilon/[2(b-a)]$。由 $f$ 可积，选分割使振幅不小于 $\eta$ 的小区间总长度足够小；在其余小区间上，$\varphi\circ f$ 的振幅不超过 $\varepsilon/[2(b-a)]$，而坏区间的贡献由 $\varphi$ 的整体振幅乘以它们的总长度控制。两部分之和可小于 $\varepsilon$，故复合可积。最后，若 $|g|\geq c$，则
\[
    \left|\frac1{g(x)}-\frac1{g(y)}\right|
    \leq c^{-2}|g(x)-g(y)|,
\]
所以 $1/g$ 可积，再用乘积结论即可。
\end{proof}

商的分母仅仅“不为零”还不够：在 $[0,1]$ 上令 $g(x)=x$（$x>0$）、$g(0)=1$，则 $g$ 可积且处处非零，但 $1/g$ 无界。两个可积函数的复合也未必可积：取例~\ref{ex:thomae-function-riemann-integrable} 中的 $T$，再令 $\varphi(0)=0$、$\varphi(u)=1$（$0<u\leq1$），则 $\varphi$ 可积，$\varphi(T(x))$ 却是狄利克雷函数。

\begin{example}
\label{ex:step-and-continuous-integral-approximation}
设 $f\in R[a,b]$。证明对每个 $\varepsilon>0$，存在阶梯函数 $h$ 与连续函数 $g$，使
\[
    \int_a^b|f-h|\,dx<\varepsilon,\qquad
    \int_a^b|f-g|\,dx<\varepsilon.
\]
还可以使 $\int_a^b|f-g|^2\,dx$ 任意小。
\end{example}
\begin{solve}
\emph{阶梯函数}（step function）是由有限个分点分段、在各开小区间内为常数的函数，分点处可另行赋值。选取满足 $\sum_i\omega_i\Delta x_i<\varepsilon$ 的分割，在每个开小区间令 $h=m_i$，在分点任意赋有限值，则除有限个点外 $|f-h|\leq\omega_i$，故其积分小于 $\varepsilon$。

再令 $g$ 为依次连接 $(x_i,f(x_i))$ 的分段线性函数。对于 $x\in[x_{i-1},x_i]$，线性插值 $g(x)$ 与 $f(x)$ 都在 $[m_i,M_i]$ 内，故 $|f(x)-g(x)|\leq\omega_i$，得到相同的积分估计。若 $|f|\leq B$，则构造还保证 $|g|\leq B$，因而
\[
    \int_a^b|f-g|^2\,dx\leq2B\int_a^b|f-g|\,dx.
\]
把右边的近似精度进一步缩小即可；$B=0$ 时取 $g=0$。这里得到的是积分意义下的近似，不声称 $\max|f-g|$ 能任意小。
\end{solve}

\begin{proposition}
\label{prop:zero-integral-continuous-functions}
设 $f\in C([a,b])$。若 $f\geq0$ 且 $\int_a^b f=0$，则 $f\equiv0$；若 $f$ 在每个子区间上的积分都是零，也有 $f\equiv0$。
\end{proposition}
\begin{proof}
若存在 $x_0$ 使 $f(x_0)>0$，连续性保证在一个正长度的子区间上 $f(x)\geq f(x_0)/2$，其积分严格为正。这在第一种情形下与整体积分为零矛盾，在第二种情形下与子区间积分为零矛盾。第二种情形若只有负值点，对 $-f$ 作同一论证即可；端点处非零时使用单侧邻域。
\end{proof}

非负条件排除了正负抵消，连续性排除了仅在孤立点取正值的情形，两者的作用不同。例如只有一个点取值为 $1$、其余点为 $0$ 的函数积分为零，却不恒为零。

\begin{example}
\label{ex:logarithm-integral-inequality}
设 $p,q>0$，证明 $\displaystyle\frac{p-q}{p}\leq\ln\frac pq\leq\frac{p-q}{q}$。
\end{example}
\begin{solve}
若 $q<p$，则在 $[q,p]$ 上有 $1/p\leq1/x\leq1/q$，积分得到结论。若 $p<q$，则先在 $[p,q]$ 上积分 $1/q\leq1/x\leq1/p$，再改变符号，也得到同一结论。$p=q$ 时两边均为零；$p\neq q$ 时，由 $1/x$ 的严格单调性，两处不等号均严格。
\end{solve}

\section{原函数的存在性与定积分的计算}
\label{sec:antiderivative-existence-and-definite-integration}

\subsection{变限定积分}
\label{subsec:definite-integrals-with-variable-limits}

设 $f\in R[a,b]$，定义\emph{变上限积分}（integral with variable upper limit）与变下限积分
\[
    \Phi(x)=\int_a^x f(t)\,dt,\qquad
    \Psi(x)=\int_x^b f(t)\,dt,\qquad x\in[a,b].
\]
这里 $x$ 是函数的自变量，$t$ 是积分内的哑变量，不能因积分变量改名而抹去上限对 $x$ 的依赖。

\begin{theorem}[变限积分与原函数]
\label{thm:variable-integral-and-antiderivative}
若 $f\in R[a,b]$ 且 $|f|\leq B$，则 $\Phi$ 满足
\[
    |\Phi(x)-\Phi(y)|\leq B|x-y|,
\]
故它是\emph{利普希茨连续}（Lipschitz continuous）的，特别地连续。若 $f$ 在 $x_0$ 连续，则 $\Phi'(x_0)=f(x_0)$，端点导数按单侧理解。因此每个连续函数都有原函数，且 $\Psi'=-f$。
\end{theorem}
\begin{proof}
由区间可加性与积分估计，
\[
    |\Phi(x)-\Phi(y)|=\left|\int_y^x f(t)\,dt\right|\leq B|x-y|.
\]
若 $f$ 在 $x_0$ 连续，则对充分小且使 $x_0+h\in[a,b]$ 的非零 $h$，
\[
    \left|\frac{\Phi(x_0+h)-\Phi(x_0)}h-f(x_0)\right|
    =\left|\frac1h\int_{x_0}^{x_0+h}[f(t)-f(x_0)]\,dt\right|
    \leq\sup_{t\text{ 在 }x_0\text{ 与 }x_0+h\text{ 之间}}|f(t)-f(x_0)|
    \longrightarrow0.
\]
最后利用 $\Psi(x)=\int_a^b f-\Phi(x)$ 即得下限求导公式。
\end{proof}

若 $f$ 连续，任一原函数 $F$ 与 $\Phi$ 的导数相同，故 $F-\Phi$ 为常数，由此也能推出牛顿--莱布尼茨公式。反之，可积函数未必存在原函数，例如阶跃函数违背\crosschapterref{5}{定理~\ref{thm:darboux-derivative-intermediate-value}}{达布定理} 中导数的介值性；存在原函数也未必黎曼可积，例如
\[
    F(x)=
    \begin{cases}x^2\sin(1/x^2),&x\neq0,\\0,&x=0,\end{cases}
    \quad
    F'(x)=
    \begin{cases}2x\sin(1/x^2)-\dfrac2x\cos(1/x^2),&x\neq0,\\0,&x=0.\end{cases}
\]
$F'$ 在零点附近无界，因而在包含零点的非退化闭区间上不可积。可积性、存在原函数和连续性是三个不同层次的性质。

设 $f$ 连续，$u,v$ 在所论区间可导，且它们的值均落在 $f$ 的定义区间内。利用固定下限的 $\Phi$ 和链式法则，
\begin{equation}
    \frac{d}{dx}\int_{u(x)}^{v(x)}f(t)\,dt
    =f(v(x))v'(x)-f(u(x))u'(x).
    \label{eq:variable-integral-chain-rule}
\end{equation}
这一公式适用于被积函数仅依赖 $t$ 的情形；若被积函数还含参数 $x$，则还须处理其对参数的变化，不能直接套用。

\begin{example}
\label{ex:variable-gaussian-integral-derivative}
求 $\displaystyle\frac{d}{dx}\int_x^{x^2}e^{-t^2}\,dt$。
\end{example}
\begin{solve}
由式~\eqref{eq:variable-integral-chain-rule}，
\[
    \frac{d}{dx}\int_x^{x^2}e^{-t^2}\,dt
    =2xe^{-x^4}-e^{-x^2}.
\]
即使 $e^{-t^2}$ 的原函数不能用初等函数表示，变限积分仍可直接求导。
\end{solve}

\begin{example}
\label{ex:variable-integral-limit}
求 $\displaystyle\lim_{x\to0}\frac{\int_0^{x^2}\sin t\,dt}{x^4}$。
\end{example}
\begin{solve}
由 $\int_0^{x^2}\sin t\,dt=1-\cos(x^2)$ 及 $1-\cos u\sim u^2/2$，极限为 $1/2$。也可先确认其为 $0/0$ 型，再用洛必达法则得到
\[
    \lim_{x\to0}\frac{2x\sin(x^2)}{4x^3}
    =\frac12\lim_{x\to0}\frac{\sin(x^2)}{x^2}=\frac12.
\]
\end{solve}

\subsection{定积分的计算}
\label{subsec:evaluation-of-definite-integrals}

计算时先检查可积性，再按表达式的分段点、对称性或适当换元化简。对于分段函数，分别使用各段的原函数即可，不必强求一个统一表达式。例如令 $f(x)=x$（$x<0$）、$f(x)=-\sin x$（$x\geq0$），则对任意 $c>0$，
\[
    \int_{-c}^{\pi}f(x)\,dx
    =\int_{-c}^0x\,dx-\int_0^\pi\sin x\,dx
    =-\frac{c^2}{2}-2.
\]

\begin{theorem}[定积分换元公式]
\label{thm:definite-integral-substitution}
设 $f\in C([a,b])$，$\varphi\in C^1([\alpha,\beta])$，且 $\varphi([\alpha,\beta])\subseteq[a,b]$，则
\begin{equation}
    \int_\alpha^\beta f(\varphi(t))\varphi'(t)\,dt
    =\int_{\varphi(\alpha)}^{\varphi(\beta)}f(x)\,dx.
    \label{eq:definite-integral-substitution}
\end{equation}
此公式不要求 $\varphi$ 单调；若将 $\varphi^{-1}$ 作为计算工具，则需另行保证相应区间上的可逆性。
\end{theorem}
\begin{proof}
连续函数 $f$ 有原函数 $F$，而 $F\circ\varphi$ 的导数是 $f(\varphi(t))\varphi'(t)$。两边分别使用牛顿--莱布尼茨公式，均得到 $F(\varphi(\beta))-F(\varphi(\alpha))$。
\end{proof}

若 $\varphi$ 在闭区间上连续、内部可导，且 $\varphi'\in R[\alpha,\beta]$，则其余条件不变时同一换元公式仍成立：$f\circ\varphi$ 连续，与可积的 $\varphi'$ 相乘仍可积，因而上述原函数论证仍适用。

换元时须同时变换被积式、微分和上下限；新上下限已经表示同一对端点，不必再将结果换回旧变量。例如 $a>0$ 时，令 $x=a\sin t$，可得
\[
    \int_0^a\sqrt{a^2-x^2}\,dx
    =a^2\int_0^{\pi/2}\cos^2t\,dt=\frac{\pi a^2}{4}.
\]

\begin{proposition}[对称性与周期性]
\label{prop:integral-symmetry-and-periodicity}
若 $f\in R[-a,a]$ 且 $a>0$，则偶函数满足 $\int_{-a}^a f=2\int_0^a f$，奇函数满足 $\int_{-a}^a f=0$。若 $f$ 连续且以 $T>0$ 为周期，则对任意 $c\in\mathbb R$，
\[
    \int_c^{c+T}f(x)\,dx=\int_0^T f(x)\,dx.
\]
\end{proposition}
\begin{proof}
在负半段作反射变换，得到 $\int_{-a}^0f(x)\,dx=\int_0^af(-t)\,dt$，分别代入奇偶关系即得前两式；反射关系对可积函数也可直接由反向排列分割证明。对于周期函数，将积分拆成 $\int_c^T f+\int_T^{c+T}f$，后一项令 $x=t+T$ 得 $\int_0^cf$，相加即得；有向积分约定保证任意 $c$ 时都有效。
\end{proof}

\begin{example}
\label{ex:reflection-integral-x-sine}
求 $\displaystyle I=\int_0^\pi\frac{x\sin x}{1+\cos^2x}\,dx$。
\end{example}
\begin{solve}
令 $x=\pi-t$，得到 $I=\int_0^\pi(\pi-x)\sin x/(1+\cos^2x)\,dx$。将两式相加，
\[
    2I=\pi\int_0^\pi\frac{\sin x}{1+\cos^2x}\,dx
    =-\pi\left.\arctan(\cos x)\right|_0^\pi=\frac{\pi^2}{2},
\]
故 $I=\pi^2/4$。反射对称性消去了原式中难处理的因子 $x$。
\end{solve}

\begin{example}
\label{ex:hyperbolic-distance-on-interval}
在 $(-1,1)$ 内定义
\[
    d(x,y)=\left|\int_x^y\frac{dt}{1-t^2}\right|.
\]
验证它是距离，并说明变换 $T_c(x)=(x-c)/(1-cx)$（$|c|<1$）保持此距离。
\end{example}
\begin{solve}
令 $H(x)=\frac12\ln\frac{1+x}{1-x}$，则 $H'(x)=1/(1-x^2)>0$，且 $H$ 将 $(-1,1)$ 严格递增地映到 $\mathbb R$ 上。于是 $d(x,y)=|H(y)-H(x)|$，非负性、正定性、对称性及三角不等式都由实数上的绝对值性质得到。又有
\[
    1-T_c(x)^2=\frac{(1-c^2)(1-x^2)}{(1-cx)^2}>0,\qquad
    T_c'(x)=\frac{1-c^2}{(1-cx)^2},
\]
故 $T_c(x)\in(-1,1)$，且 $T_c'(x)/(1-T_c(x)^2)=1/(1-x^2)$。用换元公式可得 $d(T_c(x),T_c(y))=d(x,y)$。

固定 $x_0$ 时，$x\to1^-$ 或 $x\to-1^+$ 均有 $d(x_0,x)\to+\infty$。若 $(x_n)$ 在此距离下为柯西序列，则 $(H(x_n))$ 是实数柯西序列，收敛于某个 $u\in\mathbb R$，从而 $x_n$ 在 $d$ 下收敛于 $H^{-1}(u)\in(-1,1)$。因此这个度量空间完备。这是区间上的\emph{双曲距离}（hyperbolic distance）的一种归一化。
\end{solve}

\begin{theorem}[定积分分部积分公式]
\label{thm:definite-integration-by-parts}
设 $u,v$ 在 $[a,b]$ 上连续、在内部可导，且 $u',v'$ 具有端点赋值后属于 $R[a,b]$，则
\begin{equation}
    \int_a^b u(x)v'(x)\,dx
    =\left.u(x)v(x)\right|_a^b-\int_a^b u'(x)v(x)\,dx.
    \label{eq:definite-integration-by-parts}
\end{equation}
特别地，$u,v\in C^1([a,b])$ 时适用。
\end{theorem}
\begin{proof}
由乘积可积性，$u'v$、$uv'$ 可积，而 $(uv)'=u'v+uv'$。对该等式积分，并对左边用牛顿--莱布尼茨公式，即得结论。
\end{proof}

\begin{example}
\label{ex:integration-by-parts-logarithm}
求 $\displaystyle\int_0^1x\ln(1+x)\,dx$。
\end{example}
\begin{solve}
取 $u=\ln(1+x)$、$dv=x\,dx$，则
\[
    \int_0^1x\ln(1+x)\,dx
    =\frac{\ln2}{2}-\frac12\int_0^1\frac{x^2}{1+x}\,dx
    =\frac{\ln2}{2}-\frac12\int_0^1\left(x-1+\frac1{1+x}\right)\,dx
    =\frac14.
\]
\end{solve}

\begin{example}
\label{ex:wallis-recurrence-and-product}
求 $I_n=\int_0^{\pi/2}\sin^nx\,dx$（$n=0,1,2,\ldots$），并推导瓦利斯公式。
\end{example}
\begin{solve}
反射换元 $x=\pi/2-t$ 给出 $I_n=\int_0^{\pi/2}\cos^nx\,dx$。$I_0=\pi/2$、$I_1=1$；当 $n\geq2$ 时，分部积分得到
\[
    I_n=\int_0^{\pi/2}\sin^{n-1}x\,d(-\cos x)
    =(n-1)\int_0^{\pi/2}\sin^{n-2}x\cos^2x\,dx
    =(n-1)(I_{n-2}-I_n).
\]
因此 $I_n=(n-1)I_{n-2}/n$。记双阶乘为同奇偶正整数的连乘，约定 $0!!=(-1)!!=1$，则
\[
    I_{2k}=\frac{(2k-1)!!}{(2k)!!}\frac{\pi}{2},\qquad
    I_{2k+1}=\frac{(2k)!!}{(2k+1)!!}.
\]
由于 $0\leq\sin x\leq1$，有 $I_{2n+1}\leq I_{2n}\leq I_{2n-1}$；又由递推式，$I_{2n-1}/I_{2n+1}=(2n+1)/(2n)\to1$，故夹逼得到 $I_{2n}/I_{2n+1}\to1$。代入显式表达式，得到\emph{瓦利斯公式}（Wallis formula）
\begin{equation}
    \lim_{n\to\infty}
    \frac1{2n+1}\left(\frac{(2n)!!}{(2n-1)!!}\right)^2
    =\frac{\pi}{2}.
    \label{eq:wallis-product-limit}
\end{equation}
\end{solve}

\section{定积分中值定理}
\label{sec:integral-mean-value-theorems}

\subsection{定积分第一中值定理}
\label{subsec:first-integral-mean-value-theorem}

\begin{theorem}[定积分第一中值定理]
\label{thm:first-integral-mean-value}
设 $f\in C([a,b])$，$g\in R[a,b]$，且 $g$ 在 $[a,b]$ 上不变号，则存在 $\xi\in[a,b]$，使
\begin{equation}
    \int_a^b f(x)g(x)\,dx=f(\xi)\int_a^b g(x)\,dx.
    \label{eq:first-integral-mean-value}
\end{equation}
特别地，取 $g\equiv1$，有 $\int_a^b f(x)\,dx=f(\xi)(b-a)$。
\end{theorem}
\begin{proof}
不妨设 $g\geq0$，记 $m=\min f,M=\max f$。由 $mg\leq fg\leq Mg$，积分得到
\[
    m\int_a^b g\leq\int_a^b fg\leq M\int_a^b g.
\]
若 $\int g=0$，则 $\int fg=0$，任意 $\xi$ 都适用；否则 $\int fg/\int g\in[m,M]$，由 $f$ 的连续性和介值性可取到相应的 $f(\xi)$。$g\leq0$ 时对 $-g$ 应用结论。
\end{proof}

第一中值定理将加权平均表示为某点的函数值，$g$ 不变号是保证平均值落在 $f$ 的取值范围内的关键。若只假定 $f$ 可积，仍可得到 $\int fg=\mu\int g$，其中 $\inf f\leq\mu\leq\sup f$，但未必有 $\mu=f(\xi)$。若 $f,g$ 都连续且 $g$ 不变号，则还可选 $\xi\in(a,b)$：若 $f$ 在内部已经取到该平均值即无须再证；否则 $f$ 与平均值之差在内部连续且严格同号，当 $g$ 不恒为零时，其乘积的积分也严格同号，与平均值定义矛盾。

\begin{example}
\label{ex:concentrating-integral-kernel}
设 $f\in C([0,1])$，求 $\displaystyle\lim_{n\to\infty}\int_0^1\frac{nf(x)}{1+n^2x^2}\,dx$。
\end{example}
\begin{solve}
设 $|f|\leq B$，在 $r_n=n^{-1/2}$ 处分段。远离零点的部分满足
\[
    \left|\int_{r_n}^1\frac{nf(x)}{1+n^2x^2}\,dx\right|
    \leq\frac Bn\int_{r_n}^1\frac{dx}{x^2}
    \leq\frac B{\sqrt n}\longrightarrow0.
\]
近零点部分的权函数非负，由第一中值定理，存在 $\xi_n\in[0,r_n]$，使
\[
    \int_0^{r_n}\frac{nf(x)}{1+n^2x^2}\,dx
    =f(\xi_n)\arctan(nr_n)
    =f(\xi_n)\arctan\sqrt n
    \longrightarrow\frac{\pi}{2}f(0).
\]
所以原极限为 $\pi f(0)/2$。此处权函数在零点附近集中，不能仅凭它在每个 $x>0$ 处趋于零就把极限移入积分号。
\end{solve}

\begin{theorem}[带积分余项的泰勒公式]
\label{thm:taylor-integral-remainder}
设 $n\geq0$，$f$ 在包含 $x_0,x$ 的区间上具有 $n+1$ 阶连续导数，则
\begin{equation}
    f(x)=\sum_{k=0}^n\frac{f^{(k)}(x_0)}{k!}(x-x_0)^k+R_n(x),
    \qquad
    R_n(x)=\frac1{n!}\int_{x_0}^x(x-t)^n f^{(n+1)}(t)\,dt.
    \label{eq:taylor-integral-remainder}
\end{equation}
\end{theorem}
\begin{proof}
当 $n=0$ 时，这就是 $f(x)-f(x_0)=\int_{x_0}^x f'(t)\,dt$。当 $n\geq1$ 时，对余项中的积分分部积分，
\[
    R_n(x)
    =-\frac{f^{(n)}(x_0)}{n!}(x-x_0)^n
      +\frac1{(n-1)!}\int_{x_0}^x(x-t)^{n-1}f^{(n)}(t)\,dt
    =R_{n-1}(x)-\frac{f^{(n)}(x_0)}{n!}(x-x_0)^n.
\]
从 $R_0=f(x)-f(x_0)$ 逐次代入即得公式，积分方向的约定保证 $x<x_0$ 时也成立。
\end{proof}

令 $t=x_0+s(x-x_0)$，可把余项写成
\[
    R_n(x)=\frac{(x-x_0)^{n+1}}{n!}
    \int_0^1(1-s)^n f^{(n+1)}(x_0+s(x-x_0))\,ds.
\]
以 $(1-s)^n$ 为权函数用第一中值定理，便得到某个 $\xi$ 严格位于 $x_0,x$ 之间的拉格朗日余项
\[
    R_n(x)=\frac{f^{(n+1)}(\xi)}{(n+1)!}(x-x_0)^{n+1}.
\]
改对整个连续被积函数取通常的积分平均值，则对某个 $\theta\in(0,1)$ 得到柯西余项
\[
    R_n(x)=\frac{(1-\theta)^n}{n!}
    f^{(n+1)}(x_0+\theta(x-x_0))(x-x_0)^{n+1}.
\]
这里默认 $x\neq x_0$，相等时余项为零。不同余项公式中的中间点一般不同，不能当作同一个点；它们在本定理较强的连续可导条件下都可由积分余项推出。

\subsection{定积分第二中值定理}
\label{subsec:second-integral-mean-value-theorem}

第一中值定理适合把连续因子移出积分；第二中值定理则利用单调性，把一个单调因子替换成端点值，并允许另一个因子变号。证明的核心是有限和的分部求和，而非对单调函数直接求导。

\begin{lemma}[阿贝尔变换]
\label{lem:abel-summation-finite}
设 $a_1,\ldots,a_n$、$b_1,\ldots,b_n$ 为实数，记 $B_0=0$、$B_k=\sum_{i=1}^kb_i$，则
\begin{equation}
    \sum_{i=1}^na_ib_i
    =\sum_{i=1}^{n-1}(a_i-a_{i+1})B_i+a_nB_n.
    \label{eq:abel-summation-finite}
\end{equation}
\end{lemma}
\begin{proof}
将 $b_i=B_i-B_{i-1}$ 代入左边，再将两组有限和错位相减，即得右边；$B_0=0$ 消去了起始项。
\end{proof}

\begin{theorem}[定积分第二中值定理]
\label{thm:second-integral-mean-value}
设 $g\in R[a,b]$。若 $f$ 非负且单调增加，则存在 $\xi_1\in[a,b]$，使
\[
    \int_a^b f(x)g(x)\,dx=f(b)\int_{\xi_1}^b g(x)\,dx;
\]
若 $f$ 非负且单调减少，则存在 $\xi_2\in[a,b]$，使
\[
    \int_a^b f(x)g(x)\,dx=f(a)\int_a^{\xi_2}g(x)\,dx.
\]
一般地，只要 $f$ 单调，就存在 $\xi\in[a,b]$，使
\begin{equation}
    \int_a^b f(x)g(x)\,dx
    =f(a)\int_a^\xi g(x)\,dx+f(b)\int_\xi^b g(x)\,dx.
    \label{eq:second-integral-mean-value}
\end{equation}
\end{theorem}
\begin{proof}
先证非负单调减少的情形。令 $G(x)=\int_a^x g(t)\,dt$，则 $G$ 连续；记 $m=\min G,M=\max G$。对分割 $P$ 设
\[
    Q_P=\sum_{i=1}^n f(x_{i-1})\int_{x_{i-1}}^{x_i}g(t)\,dt.
\]
取 $a_i=f(x_{i-1})$、$b_i=G(x_i)-G(x_{i-1})$，由阿贝尔变换，
\[
    Q_P=\sum_{i=1}^{n-1}[f(x_{i-1})-f(x_i)]G(x_i)
          +f(x_{n-1})G(b).
\]
所有系数非负，且总和为 $f(a)$，故 $mf(a)\leq Q_P\leq Mf(a)$。设 $|g|\leq B$，由 $f$ 单调可得
\[
    \left|\int_a^bfg-Q_P\right|
    \leq B\sum_i[f(x_{i-1})-f(x_i)]\Delta x_i
    \leq B\|P\|[f(a)-f(b)]\longrightarrow0.
\]
所以 $mf(a)\leq\int fg\leq Mf(a)$。若 $f(a)=0$，则 $f\equiv0$；否则 $\int fg/f(a)$ 是 $G$ 的某个函数值，第二式得证。反射变量 $x=a+b-t$ 即把非负单调增加情形化为已证情形，得到第一式。

若 $f$ 单调增加，则 $F=f-f(a)$ 非负单调增加，因而对某个 $\xi$，
\[
    \int_a^bfg=f(a)\int_a^bg+[f(b)-f(a)]\int_\xi^bg,
\]
整理即得式~\eqref{eq:second-integral-mean-value}。单调减少时改用 $F=f-f(b)$ 及非负单调减少的结论，同样得到该式。
\end{proof}

\begin{example}
\label{ex:monotone-weighted-first-moment}
设 $f$ 在 $[a,b]$ 上单调增加，证明
\[
    \int_a^bxf(x)\,dx\geq\frac{a+b}{2}\int_a^bf(x)\,dx.
\]
\end{example}
\begin{solve}
令 $c=(a+b)/2$，对 $f(x)$ 和 $g(x)=x-c$ 用式~\eqref{eq:second-integral-mean-value}。由于 $\int_a^b(x-c)\,dx=0$，有
\[
    \int_a^bf(x)(x-c)\,dx
    =[f(b)-f(a)]\int_\xi^b(x-c)\,dx
    =\frac{f(b)-f(a)}2(\xi-a)(b-\xi)\geq0,
\]
整理即得结论。此结论不要求 $f$ 非负。
\end{solve}

\begin{example}
\label{ex:oscillatory-integral-at-zero}
令 $g(t)=\sin(1/t)$（$t\neq0$）、$g(0)=0$，$\Phi(x)=\int_0^xg(t)\,dt$。估计零点附近的 $\Phi(x)$，并考察
\[
    f(x)=
    \begin{cases}|x|^{-1/2}\Phi(x),&x\neq0,\\0,&x=0.\end{cases}
\]
在零点的可导性。
\end{example}
\begin{solve}
先设 $0<\delta<x$，换元 $u=1/t$ 得
\[
    \int_\delta^x\sin(1/t)\,dt
    =\int_{1/x}^{1/\delta}\frac{\sin u}{u^2}\,du.
\]
$u^{-2}$ 非负单调减少，由第二中值定理，存在 $\eta\in[1/x,1/\delta]$，使右边等于 $x^2\int_{1/x}^{\eta}\sin u\,du$，其绝对值不超过 $2x^2$。由 $|\int_0^\delta g|\leq\delta$，令 $\delta\to0^+$ 得 $|\Phi(x)|\leq2x^2$。$g$ 为奇函数，故 $\Phi$ 为偶函数，这一估计对 $x<0$ 也成立。

于是 $\Phi'(0)=0$，而
\[
    \left|\frac{f(x)-f(0)}x\right|
    \leq2|x|^{1/2}\longrightarrow0,
\]
故 $f'(0)=0$。尽管 $g$ 在零点不连续，其变限积分仍可能在那里可导；连续性是定理~\ref{thm:variable-integral-and-antiderivative} 的充分条件，不能反向推断为必要条件。
\end{solve}

\section{定积分在几何学中的应用}
\label{sec:geometric-applications-of-definite-integrals}

\subsection{直角坐标系下平面图形的面积}
\label{subsec:plane-areas-in-cartesian-coordinates}

设 $f,g\in C([a,b])$ 且 $g\leq f$。对于\emph{$X$ 型区域}（$x$-simple region）
\[
    D=\{(x,y):a\leq x\leq b,\ g(x)\leq y\leq f(x)\},
\]
每个竖直细条的高度为上边界减下边界，故
\begin{equation}
    A(D)=\int_a^b[f(x)-g(x)]\,dx.
    \label{eq:cartesian-area-x-simple}
\end{equation}
同理，若 $D=\{(x,y):c\leq y\leq d,\ \varphi(y)\leq x\leq\psi(y)\}$ 为\emph{$Y$ 型区域}（$y$-simple region），则
\[
    A(D)=\int_c^d[\psi(y)-\varphi(y)]\,dy.
\]
一般图形可在内部不重叠的前提下分块求和；交点会影响上下或左右边界的次序，应先解出交点再决定分段。

\begin{example}
\label{ex:area-line-hyperbola}
求直线 $y=2,y=x$ 与第一象限内的双曲线 $xy=1$ 所围图形的面积。
\end{example}
\begin{solve}
沿水平方向切片时，$1\leq y\leq2$，左边界 $x=1/y$、右边界 $x=y$，所以
\[
    A=\int_1^2\left(y-\frac1y\right)\,dy=\frac32-\ln2.
\]
若沿竖直方向切片，则须在 $x=1$ 处分段：
\[
    A=\int_{1/2}^1\left(2-\frac1x\right)\,dx+\int_1^2(2-x)\,dx.
\]
两者结果相同，选取切片方向能减少分段和计算。
\end{solve}

\begin{example}
\label{ex:young-inequality-integral-form}
设 $f:[0,\infty)\to[0,\infty)$ 连续且严格增加，$f(0)=0$。对 $a\geq0$ 及属于 $f$ 值域的 $b\geq0$，证明\emph{杨氏不等式}（Young's inequality）
\begin{equation}
    ab\leq\int_0^af(x)\,dx+\int_0^bf^{-1}(y)\,dy,
    \label{eq:young-integral-inequality}
\end{equation}
等号成立当且仅当 $b=f(a)$。若要求对任意 $b\geq0$ 使用此式，还须假定 $f(x)\to+\infty$。
\end{example}
\begin{solve}
令 $c=f^{-1}(b)$。矩形 $[0,c]\times[0,b]$ 被曲线 $y=f(x)$ 分成两部分，按竖直与水平切片分别计算，得到
\[
    \int_0^cf(x)\,dx+\int_0^bf^{-1}(y)\,dy=bc.
\]
因此式~\eqref{eq:young-integral-inequality} 右边减去左边恰为
\[
    \int_c^a[f(x)-b]\,dx.
\]
若 $a>c$，被积函数在 $(c,a]$ 上严格为正；若 $a<c$，则在 $[a,c)$ 上严格为负，但积分方向相反，结果仍严格为正。只有 $a=c$ 时为零，故结论成立。反函数仅在原函数值域上定义，不能省略这一范围条件。
\end{solve}

\subsection{参数方程表示的曲线所围平面图形的面积}
\label{subsec:plane-areas-enclosed-by-parametric-curves}

设简单闭曲线由分段 $C^1$ 参数方程 $x=x(t),y=y(t)$（$\alpha\leq t\leq\beta$）描绘一周，除首尾重合外不自交，并按\emph{正向}（positive orientation）行进，即所围区域始终在前进方向的左侧。对于可分成有限个 $X$ 型或 $Y$ 型区域的图形，由直角坐标面积公式换元可得
\begin{equation}
    A=\int_\alpha^\beta x(t)y'(t)\,dt
     =-\int_\alpha^\beta y(t)x'(t)\,dt
     =\frac12\int_\alpha^\beta[x(t)y'(t)-y(t)x'(t)]\,dt.
    \label{eq:parametric-enclosed-area}
\end{equation}
具体地，每个 $X$ 型区域的下边界由左向右、上边界由右向左，故 $-\int y\,dx$ 等于上、下边界高度差的积分。分块时内部公共边界被反向经过两次，其贡献相消；闭曲线上又有 $\int_\alpha^\beta(xy)'\,dt=0$，得到另外两个表达式。这里的 $\int y\,dx$ 只是 $\int y(t)x'(t)\,dt$ 的简写，曲线积分的系统定义留待后文。

若参数反向，公式右边变号；若重复描绘，则会重复累计面积。对一条描绘一次的简单闭曲线，未知方向时可取右边的绝对值；对自交曲线，直接取绝对值未必得到各部分几何面积之和。

\begin{example}
\label{ex:ellipse-area-by-parameter}
求椭圆 $x^2/a^2+y^2/b^2=1$（$a,b>0$）所围面积。
\end{example}
\begin{solve}
取 $x=a\cos t,y=b\sin t$，$0\leq t\leq2\pi$，该参数按正向描绘一周。因此
\[
    A=\int_0^{2\pi}ab\cos^2t\,dt=\pi ab.
\]
也可把椭圆视为 $X$ 型区域，计算 $\int_{-a}^a2b\sqrt{1-x^2/a^2}\,dx$；两种方法表达的是同一面积。
\end{solve}

\subsection{微元法}
\label{subsec:differential-element-method}

\emph{微元法}（differential element method）是把可加的整体量写成局部近似量之和，再取极限的方法。设 $Q([a,b])$ 对子区间可加，且对任意足够细的分割，每个部分满足
\[
    Q([x_{i-1},x_i])=q(\xi_i)\Delta x_i+r_i,\qquad
    |r_i|\leq\eta(\|P\|)\Delta x_i,\qquad \eta(h)\to0,
\]
其中 $q$ 可积，则
\[
    \left|Q([a,b])-\sum_iq(\xi_i)\Delta x_i\right|
    \leq\eta(\|P\|)(b-a)\longrightarrow0,
\]
从而 $Q([a,b])=\int_a^bq(x)\,dx$。常将局部的一阶主部记作 $dQ=q(x)\,dx$，称为该量的微元。

这一记法不是把有限增量 $\Delta Q$ 与 $q(x)\Delta x$ 无条件画等号，而是概括有误差控制的极限过程；仅有逐点的近似，尚不能保证任意细分后误差总和趋于零。例如曲边梯形从 $a$ 累积到 $x$ 的面积 $A(x)=\int_a^xf(t)\,dt$ 在 $f$ 连续时满足 $A'(x)=f(x)$，因而 $dA=f(x)\,dx$。实际使用时应先确认量的可加性、选定自变量、求出微元，并检查它覆盖的范围。

\subsection{极坐标方程表示的曲线所围平面图形的面积}
\label{subsec:plane-areas-enclosed-by-polar-curves}

设 $r=r(\theta)\geq0$ 在 $[\alpha,\beta]$ 上连续，$0<\beta-\alpha\leq2\pi$，曲线与两条边界射线围成区域，每条中间射线上的部分是从原点到 $r(\theta)$ 的线段，且内部不被重复扫过。一个小角段的内、外扇形面积分别为 $\frac12m_i^2\Delta\theta_i$、$\frac12M_i^2\Delta\theta_i$，其中 $m_i,M_i$ 是该段半径的最小值、最大值。由一致连续性，上下扇形面积和之差趋于零，故
\begin{equation}
    A=\frac12\int_\alpha^\beta r(\theta)^2\,d\theta,
    \qquad dA=\frac12r(\theta)^2\,d\theta.
    \label{eq:polar-area}
\end{equation}
若同一角度下区域位于内、外半径 $r_1,r_2$ 之间，则面积为 $\frac12\int(r_2^2-r_1^2)\,d\theta$。应用时先检查描绘区间，尤其多叶曲线可能在一个参数周期内重复描绘。

\begin{example}
\label{ex:cardioid-polar-area}
求心形线 $r=a(1+\cos\theta)$（$a>0$）所围面积。
\end{example}
\begin{solve}
$0\leq\theta\leq2\pi$ 描绘其边界一次，故
\[
    A=\frac{a^2}{2}\int_0^{2\pi}(1+2\cos\theta+\cos^2\theta)\,d\theta
    =\frac{3\pi a^2}{2}.
\]
\end{solve}

\begin{example}
\label{ex:circular-arc-area-two-methods}
单位圆弧 $x=\cos t,y=\sin t$（$0\leq t\leq3\pi/4$）与 $x$ 轴、直线 $x=-1/\sqrt2$ 围成图形，求其面积。
\end{example}
\begin{solve}
在正向边界的另外两条线段上，分别有 $dx=0$ 或 $y=0$，所以
\[
    A=-\int_0^{3\pi/4}\sin t(-\sin t)\,dt
    =\frac{3\pi}{8}+\frac14.
\]
也可将图形分成张角 $3\pi/4$ 的单位扇形与两直角边均为 $1/\sqrt2$ 的直角三角形，分别得到 $3\pi/8$ 与 $1/4$。
\end{solve}

\subsection{平行截面面积为已知的立体的体积}
\label{subsec:volumes-from-parallel-cross-sections}

设边界分片光滑的有界立体介于平面 $x=a,x=b$ 之间，垂直于 $x$ 轴的截面面积为连续函数 $A(x)\geq0$。在通常的分片光滑几何情形下，薄片体积的一阶主部为 $A(x)\,dx$，故由微元法得到
\begin{equation}
    V=\int_a^b A(x)\,dx.
    \label{eq:volume-by-cross-sections}
\end{equation}
这里使用了薄片体积可由相邻柱体近似的几何性质；严格的一般体积定义与截面积分原理将在重积分中建立，不能仅凭一条截面积函数就省略对立体本身的可测性条件。

若 $0\leq y\leq f(x)$ 绕 $x$ 轴旋转一周，截面为半径 $f(x)$ 的圆盘，故
\[
    V=\pi\int_a^b f(x)^2\,dx.
\]
若区域为 $0\leq r(x)\leq y\leq R(x)$，则截面为圆环，$V=\pi\int_a^b[R(x)^2-r(x)^2]\,dx$。半径是到旋转轴的距离，图形跨过旋转轴时应先确定实际截面，不能机械相减。

\begin{example}
\label{ex:ellipsoid-volume-cross-sections}
求椭球体 $x^2/a^2+y^2/b^2+z^2/c^2\leq1$（$a,b,c>0$）的体积。
\end{example}
\begin{solve}
固定 $x\in[-a,a]$，截面为半轴 $b\sqrt{1-x^2/a^2}$、$c\sqrt{1-x^2/a^2}$ 的椭圆，故
\[
    V=\int_{-a}^a\pi bc\left(1-\frac{x^2}{a^2}\right)\,dx
    =\frac{4\pi abc}{3}.
\]
\end{solve}

\begin{example}
\label{ex:astroid-solid-volume}
星形线 $x^{2/3}+y^{2/3}=a^{2/3}$（$a>0$）所围区域绕 $x$ 轴旋转一周，求所得立体体积。
\end{example}
\begin{solve}
上半支为 $y=(a^{2/3}-|x|^{2/3})^{3/2}$，截面为圆盘。利用对称性并令 $x=at$，
\[
    V=2\pi\int_0^a(a^{2/3}-x^{2/3})^3\,dx
    =2\pi a^3\int_0^1(1-t^{2/3})^3\,dt
    =\frac{32\pi a^3}{105},
\]
其中最后一个积分展开为 $1-9/5+9/7-1/3=16/105$。
\end{solve}

\subsection{曲线的弧长}
\label{subsec:arc-length-of-curves}

\begin{definition}
\label{def:rectifiable-curve-length}
设 $\gamma:[\alpha,\beta]\to\mathbb R^m$ 为连续参数曲线。对参数分割 $P:\alpha=t_0<\cdots<t_n=\beta$，作依次连接 $\gamma(t_i)$ 的内接折线，定义其长度为
\[
    \ell(P)=\sum_{i=1}^n\|\gamma(t_i)-\gamma(t_{i-1})\|.
\]
若 $\sup_P\ell(P)<\infty$，则称曲线\emph{可求长}（rectifiable），并定义\emph{弧长}（arc length）$L=\sup_P\ell(P)$；否则称其长度无穷。这里 $\|\cdot\|$ 是欧氏范数，曲线若被重复经过，相应路段按经过次数计入。
\end{definition}

细分使折线长不减，但连续曲线未必可求长。例如 $\gamma(t)=(t,t\sin(1/t))$（$t>0$）、$\gamma(0)=(0,0)$ 在 $[0,1]$ 上连续。取 $t_n^+=1/(2n\pi+\pi/2)$、$t_n^-=1/(2n\pi+3\pi/2)$，对应纵坐标为 $t_n^+$、$-t_n^-$；连接这一对点的弦长至少为 $t_n^++t_n^-$。把前 $N$ 对点按参数次序列入同一分割，折线长至少为 $\sum_{n=1}^N(t_n^++t_n^-)$，此和随 $N$ 发散，故曲线不可求长。

\begin{theorem}[弧长积分公式]
\label{thm:arc-length-integral}
若 $\gamma\in C^1([\alpha,\beta];\mathbb R^m)$，则曲线可求长，且
\begin{equation}
    L=\int_\alpha^\beta\|\gamma'(t)\|\,dt.
    \label{eq:parametric-arc-length}
\end{equation}
分段 $C^1$ 曲线可逐段求和。若还要求 $\gamma'(t)\neq0$，则称参数曲线\emph{正则}（regular）；公式本身也允许导数在某些点为零。
\end{theorem}
\begin{proof}
记 $\gamma'$ 的连续模为 $\omega(h)$。由各坐标上的牛顿--莱布尼茨公式，对任意 $\tau_i\in[t_{i-1},t_i]$，
\[
    \|\gamma(t_i)-\gamma(t_{i-1})-\gamma'(\tau_i)\Delta t_i\|
    \leq\omega(\|P\|)\Delta t_i.
\]
利用范数的反三角不等式，得到
\[
    \left|\ell(P)-\sum_i\|\gamma'(\tau_i)\|\Delta t_i\right|
    \leq\omega(\|P\|)(\beta-\alpha).
\]
故网格趋于零时，折线长趋于式~\eqref{eq:parametric-arc-length} 的右边。对任意固定分割 $Q$，取包含 $Q$ 且网格趋于零的细分，因折线长随细分不减，知 $\ell(Q)$ 不超过该积分；另一方面，足够细的折线长能任意逼近该积分，故其上确界恰为该积分。
\end{proof}

平面参数曲线的弧长微元为 $ds=\sqrt{x'(t)^2+y'(t)^2}\,dt$，空间曲线则为 $ds=\sqrt{x'^2+y'^2+z'^2}\,dt$，此处参数按增加方向积分。对于函数图像 $y=f(x)$，$f\in C^1([a,b])$，有
\[
    L=\int_a^b\sqrt{1+f'(x)^2}\,dx.
\]
极坐标下 $x=r(\theta)\cos\theta,y=r(\theta)\sin\theta$，交叉项消去，得到
\begin{equation}
    L=\int_\alpha^\beta\sqrt{r(\theta)^2+r'(\theta)^2}\,d\theta.
    \label{eq:polar-arc-length}
\end{equation}

\begin{example}
\label{ex:parabola-arc-length}
求抛物线 $y=x^2$ 在 $0\leq x\leq1$ 上的弧长。
\end{example}
\begin{solve}
由图像弧长公式，
\[
    L=\int_0^1\sqrt{1+4x^2}\,dx
    =\left.\left(\frac{x}{2}\sqrt{1+4x^2}
       +\frac14\ln(2x+\sqrt{1+4x^2})\right)\right|_0^1
    =\frac{\sqrt5}{2}+\frac14\ln(2+\sqrt5).
\]
\end{solve}

\begin{example}
\label{ex:lemniscate-elliptic-arc-length}
设 $a>0$、$0<\beta<\pi/4$，双纽线 $r^2=2a^2\cos2\theta$ 在 $0\leq\theta\leq\beta$ 上的弧长可表示为什么积分？
\end{example}
\begin{solve}
该区间上 $r>0$ 且光滑，由 $2rr'=-4a^2\sin2\theta$ 得
\[
    \sqrt{r^2+r'^2}
    =\frac{\sqrt2a}{\sqrt{\cos2\theta}}
    =\frac{\sqrt2a}{\sqrt{1-2\sin^2\theta}},
    \qquad
    L=\int_0^\beta\frac{\sqrt2a}{\sqrt{1-2\sin^2\theta}}\,d\theta.
\]
这属于\emph{椭圆积分}（elliptic integral）。建立积分表示已经解决了弧长的表达问题，实际数值可再用数值积分求出；不能把“没有初等形式”理解为该积分不存在。
\end{solve}

\subsection{旋转体的侧面积}
\label{subsec:lateral-areas-of-solids-of-revolution}

设曲线 $x=x(t),y=y(t)\geq0$（$\alpha\leq t\leq\beta$）分段 $C^1$，绕 $x$ 轴旋转后不重复覆盖曲面内部。一个小弧段可用弦段近似，弦段旋转所得圆台的侧面积为 $\pi(y_i+y_{i+1})\ell_i$；在足够细的分割下，以连续性及弧长公式控制误差，得到微元 $dS=2\pi y\,ds$，从而
\begin{equation}
    S=2\pi\int_\alpha^\beta
      y(t)\sqrt{x'(t)^2+y'(t)^2}\,dt.
    \label{eq:surface-of-revolution-area}
\end{equation}
这是对分片光滑旋转曲面的面积公式，面积的系统定义将在曲面积分中给出。特别地，$y=f(x)\geq0$ 时，
\[
    S=2\pi\int_a^b f(x)\sqrt{1+f'(x)^2}\,dx.
\]
绕 $y$ 轴时，把半径换成 $|x(t)|$；绕 $x$ 轴但 $y$ 可负时使用 $|y(t)|$，并先检查是否有不同弧段旋转后覆盖同一片曲面。侧面积不包括端面圆盘，曲线重复描绘会使该积分重复计数。

\begin{example}
\label{ex:spherical-zone-volume-and-area}
半径为 $R>0$ 的球在 $x=a$ 与 $x=a+h$ 之间的部分，其中 $-R\leq a<a+h\leq R$，求其体积与球面部分的面积。
\end{example}
\begin{solve}
母线为 $y=\sqrt{R^2-x^2}$，截面面积是 $\pi(R^2-x^2)$，故
\[
    V=\pi\int_a^{a+h}(R^2-x^2)\,dx
    =\pi h\left(R^2-a^2-ah-\frac{h^2}{3}\right).
\]
在 $|x|<R$ 时，$y\sqrt{1+y'^2}=R$，所以球面部分面积为 $S=2\pi Rh$。若端点达到 $\pm R$，可取正则参数 $x=R\sin t,y=R\cos t$ 直接应用式~\eqref{eq:surface-of-revolution-area}，得到同一结果，无须在导数发散的端点强用图像形式的 $C^1$ 条件。取 $a=-R,h=2R$，分别得到球体体积 $4\pi R^3/3$ 与球面积 $4\pi R^2$。球面带面积只由高度 $h$ 决定，而体积还与所在位置 $a$ 有关。
\end{solve}

\section{定积分在物理学中的应用}
\label{sec:physical-applications-of-definite-integrals}

质量、静力矩与转动惯量都是可加量，因此可从有限个质点的定义出发，借助微元法推广到连续分布。设平面内有限个质点的位置为 $(x_i,y_i)$、质量为 $m_i\geq0$，总质量 $M=\sum_i m_i>0$，其关于坐标轴的\emph{静力矩}（first moment of mass）、\emph{质心}（center of mass）坐标与\emph{转动惯量}（moment of inertia）分别定义为
\[
    M_x=\sum_i m_i y_i,\qquad M_y=\sum_i m_i x_i,\qquad
    \bar x=\frac{M_y}{M},\qquad \bar y=\frac{M_x}{M},
\]
\[
    I_x=\sum_i m_i y_i^2,\qquad I_y=\sum_i m_i x_i^2.
\]
静力矩使用到轴的有向坐标，可能为负；转动惯量使用到轴距离的平方，必非负。静力矩的量纲为质量乘长度，转动惯量的量纲为质量乘长度平方，这也可用于检查微元是否写对。

考虑分段 $C^1$ 的物质曲线 $\gamma(t)=(x(t),y(t))$（$\alpha\leq t\leq\beta$），假定按参数描绘一次，其\emph{线密度}（linear mass density）$\rho(t)\geq0$ 连续。线密度表示单位弧长的质量，记速度大小 $v_\gamma(t)=\sqrt{x'(t)^2+y'(t)^2}$，则
\[
    ds=v_\gamma(t)\,dt,\qquad dm=\rho(t)v_\gamma(t)\,dt.
\]
因此质量与静力矩为
\begin{equation}
    \begin{aligned}
        M&=\int_\alpha^\beta\rho(t)v_\gamma(t)\,dt,\\
        M_x&=\int_\alpha^\beta y(t)\rho(t)v_\gamma(t)\,dt,\qquad
        M_y=\int_\alpha^\beta x(t)\rho(t)v_\gamma(t)\,dt.
    \end{aligned}
    \label{eq:curve-mass-and-first-moments}
\end{equation}
当 $M>0$ 时，$\bar x=M_y/M$、$\bar y=M_x/M$；转动惯量则为
\begin{equation}
    I_x=\int_\alpha^\beta y(t)^2\rho(t)v_\gamma(t)\,dt,\qquad
    I_y=\int_\alpha^\beta x(t)^2\rho(t)v_\gamma(t)\,dt.
    \label{eq:curve-moments-of-inertia}
\end{equation}
这些公式是把小弧段近似为集中在段内一点的质点，再取分割极限得到的；由于被积量连续，局部近似误差可统一控制。特别要区分单位弧长的密度与单位参数的质量变化率，后者为 $\rho(t)v_\gamma(t)$，一般不等于 $\rho(t)$。

若 $\rho$ 为正常数，则它在质心公式中约去，质心就是曲线的\emph{形心}（centroid）；取 $\rho=1$，总质量等于弧长 $L$，有
\[
    \bar x=\frac1L\int_\alpha^\beta x(t)v_\gamma(t)\,dt,\qquad
    \bar y=\frac1L\int_\alpha^\beta y(t)v_\gamma(t)\,dt.
\]
形心只取决于几何形状；质量分布不均匀时，质心还取决于密度。

\begin{theorem}[古鲁金第一定理]
\label{thm:guldin-surface-area}
设可求长的分段光滑平面曲线长为 $L>0$，均匀分布的形心到旋转轴的距离为 $d$。若曲线位于轴的同一侧，旋转一周所得曲面没有重复覆盖，则其面积为
\[
    S=2\pi dL.
\]
这一结论也称\emph{帕普斯--古鲁金定理}（Pappus--Guldinus theorem）的曲面面积形式。
\end{theorem}
\begin{proof}
取旋转轴为 $x$ 轴，曲线位于 $y\geq0$ 一侧，则 $d=\bar y$。由式~\eqref{eq:surface-of-revolution-area} 及均匀曲线的形心公式，
\[
    S=2\pi\int_\alpha^\beta y(t)v_\gamma(t)\,dt=2\pi\bar yL.
\]
\end{proof}

\begin{example}
\label{ex:torus-surface-area-guldin}
圆周 $x^2+(y-a)^2=R^2$（$a>R>0$）绕 $x$ 轴旋转一周，求环面的面积。
\end{example}
\begin{solve}
圆周均匀分布的形心在圆心 $(0,a)$，长为 $2\pi R$，且圆周不与旋转轴相交，旋转生成的曲面没有重复覆盖。由古鲁金第一定理，
\[
    S=2\pi a\cdot2\pi R=4\pi^2aR.
\]
\end{solve}

再考虑平面薄片 $D=\{(x,y):a\leq x\leq b,\ y_1(x)\leq y\leq y_2(x)\}$，其中 $y_1,y_2$ 连续，质量均匀分布，取单位面密度。设面积 $A=\int_a^b(y_2-y_1)\,dx>0$，则总质量为 $M=A$。一条竖直薄片的质量微元为 $dm=(y_2-y_1)\,dx$，其形心的一阶近似坐标为 $(x,(y_1+y_2)/2)$，故
\[
    dM_y=x(y_2-y_1)\,dx,\qquad
    dM_x=\frac{y_1+y_2}{2}(y_2-y_1)\,dx.
\]
由此得到形心公式
\begin{equation}
    \bar x=\frac{\int_a^bx[y_2(x)-y_1(x)]\,dx}{A},\qquad
    \bar y=\frac{\frac12\int_a^b[y_2(x)^2-y_1(x)^2]\,dx}{A}.
    \label{eq:plane-region-centroid}
\end{equation}
若面密度为其他正常数，同一因子在分子分母约去，形心不变。沿竖条再对 $y^2$ 累加，还可得单位面密度薄片的转动惯量
\[
    I_x=\frac13\int_a^b[y_2(x)^3-y_1(x)^3]\,dx,\qquad
    I_y=\int_a^bx^2[y_2(x)-y_1(x)]\,dx.
\]
计算 $I_x$ 时不能把整条竖条的质量都放在其形心而忽略竖条本身沿 $y$ 方向的分布；静力矩允许这一替换，转动惯量一般不允许。

\begin{theorem}[古鲁金第二定理]
\label{thm:guldin-volume}
设面积为 $A>0$、可分为有限个上述类型区域的平面图形位于旋转轴的一侧，形心到轴的距离为 $d$，绕轴旋转一周不重复计入体积，则所得旋转体的体积为
\[
    V=2\pi dA.
\]
这是帕普斯--古鲁金定理的体积形式。
\end{theorem}
\begin{proof}
先对满足 $0\leq y_1\leq y_2$ 的 $X$ 型区域，利用圆环截面公式与式~\eqref{eq:plane-region-centroid}，
\[
    V=\pi\int_a^b(y_2^2-y_1^2)\,dx=2\pi\bar yA.
\]
对于可由有限块此类区域拼成、旋转后各块内部不重叠的区域，由体积与静力矩的可加性求和，得到同一等式。
\end{proof}

\begin{example}
\label{ex:semidisk-centroid}
求半圆盘 $D=\{(x,y):x^2+y^2\leq R^2,\ y\geq0\}$ 的形心。
\end{example}
\begin{solve}
由关于 $y$ 轴的对称性，$\bar x=0$。半圆盘面积为 $\pi R^2/2$，绕 $x$ 轴旋转一周得到体积 $4\pi R^3/3$ 的球。古鲁金第二定理给出
\[
    \frac{4\pi R^3}{3}
    =2\pi\bar y\frac{\pi R^2}{2},
    \qquad
    (\bar x,\bar y)=\left(0,\frac{4R}{3\pi}\right).
\]
这是均匀半圆盘的形心；弧线与薄片的质量分布不同，其形心也不同。
\end{solve}

\begin{example}
\label{ex:polar-wire-moment-of-inertia}
设平面曲线由 $r=r(\theta)\geq0$（$\alpha\leq\theta\leq\beta$）给出，$r\in C^1$，线密度为 $1$，求它关于极轴的转动惯量。
\end{example}
\begin{solve}
极轴即 $x$ 轴，点到轴的距离为 $|r(\theta)\sin\theta|$，而 $dm=ds=\sqrt{r(\theta)^2+r'(\theta)^2}\,d\theta$。因此
\[
    I_x=\int_\alpha^\beta
    r(\theta)^2\sin^2\theta
    \sqrt{r(\theta)^2+r'(\theta)^2}\,d\theta.
\]
若线密度为非负连续函数 $\rho(\theta)$，在被积式中再乘 $\rho(\theta)$ 即可；若改求关于原点的极转动惯量，则距离平方改为 $r(\theta)^2$。
\end{solve}

\section*{习题七}
\addcontentsline{toc}{section}{习题七}
\fi

\selectchapter{8}
\ifcompilechapter
\chapter{广义积分}
\label{chap:improper-integrals}

% 本章依据“数学分析/textbooks/数学分析（分章原文）/数学分析-8.pdf”正文第 1--26 页整理，不含习题八。
定积分要求积分区间有限且被积函数有界。本章把这两个限制分别放宽：在无穷区间上以截断区间的积分极限定义\emph{无穷积分}（improper integral over an infinite interval），在有限区间上以避开无界点后的积分极限定义\emph{瑕积分}（improper integral with a singularity），二者统称\emph{广义积分}（improper integral）。两类问题都先在正常的闭区间上积分，再研究截断参数的极限；因此，写出广义积分的符号并不表示这个极限已经存在，也不能把定积分的运算法则直接用于尚未证明收敛的表达式。

\section{无穷积分的基本概念与性质}
\label{sec:infinite-integrals-concepts-and-properties}

考虑一个理想化的引力做功问题：设地球质量为 $M$、半径为 $R$，质量固定为 $m$ 的物体沿径向从地表升至高度 $H$，忽略空气阻力、其他天体及地球自转的影响。以 $G$ 表示万有引力常数，则克服引力所做的功为
\[
    W(H)=\int_R^{R+H}\frac{GMm}{r^2}\,\mathrm{d}r
        =GMm\left(\frac1R-\frac1{R+H}\right).
\]
每一个有限的 $H$ 对应普通定积分；物体远离地球时的总功则由 $H\to+\infty$ 的极限给出，等于 $GMm/R$。在同一模型下，把这个极限与初动能 $mv^2/2$ 比较，得到逃逸速度 $v=\sqrt{2GM/R}=\sqrt{2gR}$，其中 $g=GM/R^2$。无穷积分描述的正是这种截断位置不断远移时的累积量。

\begin{definition}[无穷积分及其收敛性]
\label{def:improper-integrals-at-infinity}
设实函数 $f$ 在 $[a,+\infty)$ 上\emph{局部黎曼可积}（locally Riemann integrable），即对每个 $X>a$，$f$ 在 $[a,X]$ 上黎曼可积。若存在有限实数极限
\begin{equation}
    \lim_{X\to+\infty}\int_a^X f(x)\,\mathrm{d}x=I,
    \label{eq:improper-integral-right-limit}
\end{equation}
则称无穷积分 $\int_a^{+\infty}f(x)\,\mathrm{d}x$ \emph{收敛}（convergent），其值定义为 $I$；若无有限极限，则称其\emph{发散}（divergent）。类似地，若 $f$ 在 $(-\infty,b]$ 的每个有限闭子区间上可积，则定义
\[
    \int_{-\infty}^b f(x)\,\mathrm{d}x
    =\lim_{X\to-\infty}\int_X^b f(x)\,\mathrm{d}x,
\]
并以右端有限极限的存在性定义收敛。
\end{definition}

极限等于 $+\infty$ 或 $-\infty$ 仍属于发散；必要时可以用这两个符号说明发散方向，但它们不是收敛积分的值。后文的“局部可积”均指在所讨论区间内每个避开瑕点的有限闭子区间上黎曼可积。

\begin{definition}[整个实轴上的无穷积分]
\label{def:improper-integral-whole-line}
设 $f$ 在实轴的每个有限闭区间上可积。固定 $c\in\mathbb{R}$，只有当 $\int_{-\infty}^c f(x)\,\mathrm{d}x$ 与 $\int_c^{+\infty}f(x)\,\mathrm{d}x$ 分别收敛时，才称 $\int_{-\infty}^{+\infty}f(x)\,\mathrm{d}x$ 收敛，并定义
\begin{equation}
    \int_{-\infty}^{+\infty}f(x)\,\mathrm{d}x
    =\int_{-\infty}^c f(x)\,\mathrm{d}x
     +\int_c^{+\infty}f(x)\,\mathrm{d}x.
    \label{eq:improper-integral-whole-line-split}
\end{equation}
若其中至少一个积分发散，则整个实轴上的积分发散。
\end{definition}

分点 $c$ 的选取不影响敛散性，也不影响收敛时的积分值，因为改变分点只是在两部分之间转移一个有限区间上的定积分。更一般地，对任意 $b>a$，由
\[
    \int_a^X f(x)\,\mathrm{d}x
    =\int_a^b f(x)\,\mathrm{d}x+\int_b^X f(x)\,\mathrm{d}x
    \qquad(X>b)
\]
可知，$\int_a^{+\infty}f$ 与 $\int_b^{+\infty}f$ 同敛散。因此无穷积分的敛散性取决于尾部，而有限初始区间只改变收敛后的数值。若 $f\geq0$，截断积分表示曲线下的面积；若 $f$ 变号，则表示正、负部分面积的代数和，后者可以因抵消而收敛。

\begin{proposition}[收敛积分的线性与原函数计算]
\label{prop:improper-integral-linearity-and-primitives}
若 $\int_a^{+\infty}f$、$\int_a^{+\infty}g$ 均收敛，$\lambda,\mu\in\mathbb{R}$，则
\[
    \int_a^{+\infty}(\lambda f+\mu g)\,\mathrm{d}x
    =\lambda\int_a^{+\infty}f\,\mathrm{d}x
     +\mu\int_a^{+\infty}g\,\mathrm{d}x.
\]
若 $f$ 局部可积且有原函数 $F$，则 $\int_a^{+\infty}f$ 收敛当且仅当 $\lim_{X\to+\infty}F(X)$ 为有限实数，此时
\begin{equation}
    \int_a^{+\infty}f(x)\,\mathrm{d}x
    =\lim_{X\to+\infty}F(X)-F(a).
    \label{eq:improper-newton-leibniz}
\end{equation}
左侧无穷区间上的积分有对应结论；整个实轴上的积分则要求原函数在两端分别有有限极限。
\end{proposition}

\begin{proof}
在每个有限截断区间上，分别使用定积分的线性与牛顿--莱布尼茨公式，再取相应极限即可。线性公式中两项的收敛性保证可以把和的极限拆开；不能把两个发散积分之差当成通常的实数运算。
\end{proof}

\begin{example}
\label{ex:improper-integral-arctangent}
计算 $\int_{-\infty}^{+\infty}\frac{\mathrm{d}x}{1+x^2}$。
\end{example}

\begin{solve}
取分点 $0$，以 $\arctan x$ 为原函数，分别得到
\[
    \int_{-\infty}^0\frac{\mathrm{d}x}{1+x^2}=\frac\pi2,
    \qquad
    \int_0^{+\infty}\frac{\mathrm{d}x}{1+x^2}=\frac\pi2.
\]
两部分均有有限值，故原积分收敛，值为 $\pi$。
\end{solve}

\begin{example}
\label{ex:improper-power-and-log-benchmarks}
讨论幂函数、对数函数构成的基本比较积分。
\end{example}

\begin{solve}
对 $p\in\mathbb{R}$，直接计算截断积分得
\begin{equation}
    \int_1^{+\infty}\frac{\mathrm{d}x}{x^p}
    \begin{cases}
        =\dfrac1{p-1},&p>1,\\
        \text{发散至 }+\infty,&p\leq1.
    \end{cases}
    \label{eq:infinite-power-integral-benchmark}
\end{equation}
这里 $p=1$ 时截断积分为 $\ln X$，$p\neq1$ 时为 $(X^{1-p}-1)/(1-p)$。临界指数 $p=1$ 必须单独检验，不能从 $p\neq1$ 的表达式直接代入。

换元 $t=\ln x$ 后，
\[
    \int_2^{+\infty}\frac{\mathrm{d}x}{x(\ln x)^p}
    =\int_{\ln2}^{+\infty}\frac{\mathrm{d}t}{t^p},
\]
因此仍当且仅当 $p>1$ 时收敛，收敛时的值为 $(\ln2)^{1-p}/(p-1)$。再令 $t=\ln\ln x$，得到
\[
    \int_{e^e}^{+\infty}
        \frac{\mathrm{d}x}{x\ln x(\ln\ln x)^p}
    =\int_1^{+\infty}\frac{\mathrm{d}t}{t^p},
\]
其收敛范围同样为 $p>1$。这些积分说明：幂次达到临界值后，还需由下一层对数因子决定敛散性。
\end{solve}

\begin{proposition}[换元与分部积分]
\label{prop:improper-substitution-and-parts}
设 $f$ 在 $[a,+\infty)$ 上局部可积，$\varphi$ 在 $[\alpha,\beta)$ 上连续可微且严格递增，其中 $\beta$ 可以为 $+\infty$，并满足 $\varphi(\alpha)=a$、$\lim_{t\to\beta-}\varphi(t)=+\infty$。则
\begin{equation}
    \int_a^{+\infty}f(x)\,\mathrm{d}x
    =\int_\alpha^\beta f(\varphi(t))\varphi'(t)\,\mathrm{d}t;
    \label{eq:improper-substitution}
\end{equation}
等式指两边同时收敛或同时发散，收敛时值相等。递减换元应同时调整积分上下限。

若 $u,v\in C^1([a,+\infty))$，且 $L=\lim_{X\to+\infty}u(X)v(X)$ 为有限实数，则 $\int_a^{+\infty}uv'$ 与 $\int_a^{+\infty}u'v$ 同敛散，收敛时
\begin{equation}
    \int_a^{+\infty}u(x)v'(x)\,\mathrm{d}x
    =L-u(a)v(a)-\int_a^{+\infty}u'(x)v(x)\,\mathrm{d}x.
    \label{eq:improper-integration-by-parts}
\end{equation}
这里要求的是乘积 $uv$ 的极限存在，不要求两个因子各有有限极限。
\end{proposition}

\begin{proof}
在 $[\alpha,T]$ 上使用定积分换元公式，右端等于 $\int_a^{\varphi(T)}f$；$T\to\beta-$ 时，$\varphi(T)$ 遍历趋于无穷的截断位置，故两边的极限存在性等价。分部积分则先写出有限区间上的恒等式
\[
    \int_a^X uv'\,\mathrm{d}x
    =u(X)v(X)-u(a)v(a)-\int_a^X u'v\,\mathrm{d}x.
\]
由于边界乘积已有有限极限，右端剩下的积分极限与左端同时存在，从而得到结论。
\end{proof}

\begin{example}
\label{ex:improper-substitution-and-parts-computations}
计算 $\int_1^{+\infty}\frac{\mathrm{d}x}{x\sqrt{1+x^2}}$ 与 $\int_1^{+\infty}\frac{\arctan x}{x^3}\,\mathrm{d}x$。
\end{example}

\begin{solve}
第一个积分令 $t=1/x$，将无穷端点移到 $t=0$，得到
\[
    \int_1^{+\infty}\frac{\mathrm{d}x}{x\sqrt{1+x^2}}
    =\int_0^1\frac{\mathrm{d}t}{\sqrt{1+t^2}}
    =\ln(1+\sqrt2).
\]
换元后的被积函数在 $0$ 处可连续延拓，故右端为普通定积分，这也同时证明原积分收敛。

第二个积分取 $u=\arctan x$、$v=-1/(2x^2)$，则 $uv\to0$，且
\begin{align*}
    \int_1^{+\infty}\frac{\arctan x}{x^3}\,\mathrm{d}x
    &=\frac\pi8+\frac12\int_1^{+\infty}
        \frac{\mathrm{d}x}{x^2(1+x^2)}\\
    &=\frac\pi8+\frac12\int_1^{+\infty}
        \left(\frac1{x^2}-\frac1{1+x^2}\right)\,\mathrm{d}x
      =\frac12.
\end{align*}
右侧各积分均已证明收敛，因此分部积分的使用有依据。
\end{solve}

\begin{example}
\label{ex:exponentially-damped-trigonometric-integrals}
设 $a>0$、$b\in\mathbb{R}$，计算
\[
    I_s=\int_0^{+\infty}e^{-ax}\sin bx\,\mathrm{d}x,
    \qquad
    I_c=\int_0^{+\infty}e^{-ax}\cos bx\,\mathrm{d}x.
\]
\end{example}

\begin{solve}
先验证收敛性：对 $Y>X\geq0$，两个被积函数的尾部积分绝对值均不超过
\[
    \int_X^Y e^{-ax}\,\mathrm{d}x\leq\frac{e^{-aX}}a\longrightarrow0.
\]
因此由函数极限的柯西准则，两个截断积分都有有限极限；下一节的绝对收敛与比较判别法也可直接说明这一点。分别分部积分，边界项因 $e^{-ax}\to0$ 而存在，得
\[
    I_s=\frac ba I_c,\qquad
    I_c=\frac1a-\frac ba I_s.
\]
解这两个线性方程可得
\[
    I_s=\frac{b}{a^2+b^2},\qquad
    I_c=\frac{a}{a^2+b^2}.
    \qedhere
\]
\end{solve}

\begin{definition}[无穷积分的柯西主值]
\label{def:cauchy-principal-value-at-infinity}
设 $f$ 在实轴的每个有限闭区间上可积。若对称截断的有限极限存在，则称
\begin{equation}
    \operatorname{PV}\int_{-\infty}^{+\infty}f(x)\,\mathrm{d}x
    =\lim_{X\to+\infty}\int_{-X}^X f(x)\,\mathrm{d}x
    \label{eq:cauchy-principal-value-infinity}
\end{equation}
为该积分的\emph{柯西主值}（Cauchy principal value）。
\end{definition}

通常收敛要求左右两个极限独立存在，故一定蕴含主值存在且与通常积分相等；主值只考察一种对称截断方式，反向推论不成立。例如 $f(x)=x/(1+x^2)$ 为奇函数，每个对称截断积分都为零，故主值为零，但
\[
    \int_0^X\frac{x}{1+x^2}\,\mathrm{d}x=\frac12\ln(1+X^2)\to+\infty,
\]
所以通常的整个实轴积分发散。\emph[caution]{主值中的对称抵消不能替代两端各自的收敛性检验。}

\section{无穷积分敛散性的判别法}
\label{sec:convergence-tests-for-infinite-integrals}

定义把敛散性归结为 $F(X)=\int_a^X f$ 是否有有限极限，但一般不能求出 $F$ 的初等表达式。判别法的作用是利用被积函数的大小、符号或振荡性质直接控制尾部积分。除特别说明外，本节出现的函数均在 $[a,+\infty)$ 上局部可积；改变有限初始区间不影响结论。

\begin{theorem}[无穷积分的柯西准则]
\label{thm:cauchy-criterion-improper-integrals}
$\int_a^{+\infty}f(x)\,\mathrm{d}x$ 收敛的充分必要条件是：对每个 $\varepsilon>0$，存在 $A>a$，使任意 $v>u>A$ 均满足
\begin{equation}
    \left|\int_u^v f(x)\,\mathrm{d}x\right|<\varepsilon.
    \label{eq:improper-integral-cauchy-tail}
\end{equation}
\end{theorem}

\begin{proof}
令 $F(X)=\int_a^X f(x)\,\mathrm{d}x$，则 $F(v)-F(u)=\int_u^v f$。将函数极限的柯西准则\crosschapterref{3}{，即定理~\ref{thm:cauchy-criterion-for-function-limits}}{}，用于 $X\to+\infty$，即得结论。
\end{proof}

这里要求\emph{任意两个足够远的截断点之间的积分一致地小}，不是只检验某一列区间或某种配对方式的极限。这个量词要求排除了利用特殊截断掩盖发散的情形。

\begin{definition}[绝对收敛与条件收敛]
\label{def:absolute-and-conditional-improper-integrals}
若 $\int_a^{+\infty}|f(x)|\,\mathrm{d}x$ 收敛，则称 $\int_a^{+\infty}f(x)\,\mathrm{d}x$ \emph{绝对收敛}（absolutely convergent）；若后者收敛而前者发散，则称后者\emph{条件收敛}（conditionally convergent）。
\end{definition}

\begin{theorem}[绝对收敛蕴含收敛]
\label{thm:absolute-improper-integrability-implies-convergence}
若 $\int_a^{+\infty}|f(x)|\,\mathrm{d}x$ 收敛，则 $\int_a^{+\infty}f(x)\,\mathrm{d}x$ 收敛，且
\[
    \left|\int_a^{+\infty}f(x)\,\mathrm{d}x\right|
    \leq\int_a^{+\infty}|f(x)|\,\mathrm{d}x.
\]
\end{theorem}

\begin{proof}
由 $f$ 的局部可积性，$|f|$ 也局部可积。有限区间上的不等式
\[
    \left|\int_u^v f(x)\,\mathrm{d}x\right|
    \leq\int_u^v|f(x)|\,\mathrm{d}x
\]
与 $|f|$ 的柯西准则一起给出 $f$ 的柯西准则，故原积分收敛。再在 $[a,X]$ 上使用同一不等式并令 $X\to+\infty$，得到结论中的估计。
\end{proof}

\begin{theorem}[非负函数积分的有界判据]
\label{thm:nonnegative-improper-integral-boundedness}
若 $f(x)\geq0$，则 $\int_a^{+\infty}f$ 收敛的充分必要条件是其截断积分有共同上界，即存在常数 $C$，使
\[
    \int_a^X f(x)\,\mathrm{d}x\leq C\qquad(X>a).
\]
若无此上界，则截断积分趋于 $+\infty$。
\end{theorem}

\begin{proof}
由定积分对积分区间的可加性，$F(X)=\int_a^X f$ 单调非减。单调函数的极限存在性定理表明，它有有限极限当且仅当有上界；无上界时极限为 $+\infty$。
\end{proof}

\begin{theorem}[比较判别法及其极限形式]
\label{thm:comparison-tests-improper-integrals}
\emph{比较判别法}（comparison test）及其\emph{极限比较形式}（limit comparison test）如下。设 $f,g\geq0$，且存在 $A\geq a$、$C>0$，使 $x\geq A$ 时有 $f(x)\leq Cg(x)$。若 $\int_a^{+\infty}g$ 收敛，则 $\int_a^{+\infty}f$ 收敛；若 $\int_a^{+\infty}f$ 发散，则 $\int_a^{+\infty}g$ 发散。

若进一步有 $g(x)>0$，且 $\lim_{x\to+\infty}f(x)/g(x)=\ell$，则：$0<\ell<+\infty$ 时两积分同敛散；$\ell=0$ 时由 $\int_a^{+\infty}g$ 收敛可推出 $\int_a^{+\infty}f$ 收敛；$\ell=+\infty$ 时由 $\int_a^{+\infty}g$ 发散可推出 $\int_a^{+\infty}f$ 发散。
\end{theorem}

\begin{proof}
若 $\int_a^{+\infty}g$ 收敛，则 $\int_A^Xf\leq C\int_A^Xg$ 有共同上界，由非负函数积分的有界判据得 $\int_A^{+\infty}f$ 收敛，加回有限初始区间即可。发散结论为此命题的逆否命题。若 $0<\ell<+\infty$，充分大时有 $(\ell/2)g\leq f\leq(3\ell/2)g$，故可双向比较；$\ell=0$ 时最终有 $f\leq g$，$\ell=+\infty$ 时最终有 $f\geq g$，分别使用单向比较即可。
\end{proof}

比较的方向可以概括为“大者收敛则小者收敛，小者发散则大者发散”，但前提是被积函数非负。对变号函数，可以先比较绝对值以证明绝对收敛；仅有 $f\sim g$ 而不能控制符号与绝对值时，不能据此照搬上述敛散性结论。极限为零或无穷时也只有定理给出的单向推论。

以例~\ref{ex:improper-power-and-log-benchmarks} 为基准，并利用对数慢于任意正幂增长，可得对 $p,q\in\mathbb{R}$，
\begin{equation}
    \int_e^{+\infty}\frac{\mathrm{d}x}{x^p(\ln x)^q}
    \text{ 收敛}\quad\Longleftrightarrow\quad
    p>1\ \text{或}\ (p=1\text{ 且 }q>1).
    \label{eq:power-log-infinite-integral-test}
\end{equation}
具体地，$p>1$ 时选 $1<r<p$，有 $x^{-p}(\ln x)^{-q}=o(x^{-r})$；$p<1$ 时选 $p<r<1$，则该被积函数与 $x^{-r}$ 的比值趋于 $+\infty$；$p=1$ 时令 $t=\ln x$。同理，对任意多项式 $P$ 和 $a>0$，$|P(x)|e^{-ax}\leq Ce^{-ax/2}$ 在充分大时成立，所以 $\int_0^{+\infty}P(x)e^{-ax}\,\mathrm{d}x$ 绝对收敛；这也给出例~\ref{ex:exponentially-damped-trigonometric-integrals} 中两个积分的另一种收敛性证明。

\begin{example}
\label{ex:rational-infinite-integral-degree-test}
设 $n$ 为正整数，讨论 $\int_0^{+\infty}\frac{x^2}{x^n+x^2+1}\,\mathrm{d}x$ 的敛散性，并说明有理函数的一般判据。
\end{example}

\begin{solve}
被积函数在每个有限闭区间上连续。$n>2$ 时，它在无穷远处等价于 $x^{2-n}$，故当且仅当 $n-2>1$，即 $n>3$ 时收敛。$n=1$、$n=2$ 时，被积函数分别趋于 $1$、$1/2$，故均发散；这两种情形的分母最高次项不同，不能直接套用 $n>2$ 时的等价式。

一般地，设非零实多项式 $P,Q$ 的次数分别为 $m,n$，首项系数分别为 $p_m,q_n$，且 $Q$ 在 $[a,+\infty)$ 上无零点，则
\[
    \frac{P(x)}{Q(x)}\sim\frac{p_m}{q_n}x^{m-n}
    \qquad(x\to+\infty).
\]
有理函数在充分远处符号固定，故其积分收敛等价于绝对收敛，再由幂积分判据可知，这当且仅当 $n-m\geq2$。若有限区间中还有未消去的极点，则必须另行检验，次数判据只处理无穷端。
\end{solve}

\begin{example}
\label{ex:absolute-comparison-and-log-cancellation}
证明 $\int_1^{+\infty}\frac{\cos x\sin(1/x)}x\,\mathrm{d}x$ 绝对收敛，并讨论 $p>0$ 时 $\int_1^{+\infty}[x^{-p}-\ln(1+x^{-p})]\,\mathrm{d}x$ 的敛散性。
\end{example}

\begin{solve}
由 $|\sin t|\leq|t|$，第一问的被积函数绝对值不超过 $x^{-2}$，故绝对收敛。第二问利用 $t-\ln(1+t)=t^2/2+o(t^2)$，令 $t=x^{-p}$，得到
\[
    x^{-p}-\ln(1+x^{-p})\sim\frac1{2x^{2p}}.
\]
被积函数非负，故极限比较合法，原积分当且仅当 $2p>1$，即 $p>1/2$ 时收敛。这里首阶项相消，必须展开到二阶；若只把 $\ln(1+x^{-p})$ 换成 $x^{-p}$，便会丢失决定敛散性的主项。
\end{solve}

\begin{example}
\label{ex:continuous-integrable-spikes}
构造连续正函数 $f$，使 $\int_1^{+\infty}f(x)\,\mathrm{d}x$ 收敛，但 $f(x)$ 在无穷远处不趋于零，且无界。
\end{example}

\begin{solve}
用越来越窄的三角形尖峰控制总面积。记 $r_+=\max\{r,0\}$，在每个整数 $n\geq2$ 附近放置高为 $2^n$、半底宽为 $3^{-n}$ 的尖峰，定义
\[
    f(x)=\frac1{x^2}
       +\sum_{n=2}^{\infty}2^n\bigl(1-3^n|x-n|\bigr)_+,
    \qquad x\geq1.
\]
各尖峰支撑区间 $[n-3^{-n},n+3^{-n}]$ 互不相交，且任意有限闭区间只遇到有限个尖峰，所以这里的和在局部实际为有限和，$f$ 连续且严格为正。第 $n$ 个三角形的面积为 $(2/3)^n$，故对任意 $X>1$，
\[
    \int_1^X f(x)\,\mathrm{d}x
    \leq1+\sum_{n=2}^{\infty}\left(\frac23\right)^n
    =\frac73.
\]
由非负函数积分的有界判据，广义积分收敛；但 $f(n)=n^{-2}+2^n\to+\infty$，所以函数值并不趋于零。这一论证只用了有限区间上有限多个尖峰的面积估计，没有交换无穷求和与积分。
\end{solve}

\emph[caution]{积分收敛控制的是尾部累积量，一般不强迫被积函数逐点趋于零，即使函数连续且非负也不例外。}反过来，$f(x)\to0$ 也不足以保证收敛，$f(x)=1/x$ 即为反例。若另外知道 $f$ 在尾部单调，则积分收敛必推出 $f(x)\to0$：单调性保证扩充实数极限存在，若极限非零，则尾部最终符号固定且绝对值有正下界，其积分不可能有有限极限。

\begin{theorem}[狄利克雷判别法]
\label{thm:dirichlet-test-infinite-integrals}
设 $f$ 局部可积，其累积积分一致有界，即存在 $M>0$，使
\[
    \left|\int_a^X f(x)\,\mathrm{d}x\right|\leq M
    \qquad(X\geq a).
\]
若 $g$ 在 $[a,+\infty)$ 上单调且 $\lim_{x\to+\infty}g(x)=0$，则 $\int_a^{+\infty}f(x)g(x)\,\mathrm{d}x$ 收敛。这称为\emph{狄利克雷判别法}（Dirichlet's test）。
\end{theorem}

\begin{proof}
记 $F(X)=\int_a^X f$，则任意 $v>u\geq a$ 有 $|\int_u^v f|=|F(v)-F(u)|\leq2M$。在有限区间 $[u,v]$ 上，$g$ 单调，故第二积分中值定理给出某个 $\xi\in[u,v]$，使
\[
    \int_u^v fg\,\mathrm{d}x
    =g(u)\int_u^\xi f\,\mathrm{d}x
     +g(v)\int_\xi^v f\,\mathrm{d}x.
\]
因此
\[
    \left|\int_u^v fg\,\mathrm{d}x\right|
    \leq2M\bigl(|g(u)|+|g(v)|\bigr).
\]
当 $u,v$ 都足够大时，右端可任意小；由柯西准则，乘积积分收敛。单调函数在有限区间上可积，故 $fg$ 的局部可积性也有保证。
\end{proof}

\begin{theorem}[阿贝尔判别法]
\label{thm:abel-test-infinite-integrals}
若 $\int_a^{+\infty}f(x)\,\mathrm{d}x$ 收敛，而 $g$ 在 $[a,+\infty)$ 上单调有界，则 $\int_a^{+\infty}f(x)g(x)\,\mathrm{d}x$ 收敛。这称为\emph{阿贝尔判别法}（Abel's test）。
\end{theorem}

\begin{proof}
由 $g$ 单调有界，存在有限极限 $\ell=\lim_{x\to+\infty}g(x)$。又因 $\int_a^Xf$ 收敛于有限值，其在尾部有界，在每个有限初始区间上也有界，故可将狄利克雷判别法用于 $f$ 与 $g-\ell$，得到 $\int_a^{+\infty}f(g-\ell)$ 收敛。再加上已收敛的 $\ell\int_a^{+\infty}f$ 即得结论。
\end{proof}

两种判别法的分工是：狄利克雷判别法只要求一个因子的累积积分有界，但要求另一个因子单调趋于零；阿贝尔判别法要求第一个因子的积分已经收敛，另一个因子只需单调有界。例如，在下例证明 $\int_1^{+\infty}\sin x/x\,\mathrm{d}x$ 收敛后，再乘以单调有界的 $\arctan x$，即可用阿贝尔判别法判断所得积分收敛。两种判别法通常给出收敛性，绝对收敛仍需另行检验。

\begin{example}
\label{ex:sinc-integral-conditional-convergence}
证明 $\int_{-\infty}^{+\infty}\frac{\sin x}{x}\,\mathrm{d}x$ 条件收敛，其中在 $x=0$ 处把被积函数连续延拓为 $1$。
\end{example}

\begin{solve}
函数在有限区间上连续，只需讨论无穷尾部。在 $[1,+\infty)$ 上，$\int_1^X\sin x\,\mathrm{d}x=\cos1-\cos X$ 一致有界，而 $1/x$ 单调趋于零，故由狄利克雷判别法，$\int_1^{+\infty}\sin x/x\,\mathrm{d}x$ 收敛。被积函数为偶函数，负半轴上的积分也收敛，因而整个实轴上的通常积分收敛。

另一方面，对 $X>1$，利用 $|\sin x|\geq\sin^2x$，有
\[
    \int_1^X\frac{|\sin x|}{x}\,\mathrm{d}x
    \geq\int_1^X\frac{\sin^2x}{x}\,\mathrm{d}x
    =\frac12\ln X-\frac12\int_1^X\frac{\cos2x}{x}\,\mathrm{d}x.
\]
最后一个振荡积分由狄利克雷判别法收敛，右端因 $\ln X\to+\infty$ 而趋于 $+\infty$。所以绝对值积分发散，原积分条件收敛。这里先判定收敛而未求其值，说明判别法可以处理没有初等原函数的积分。
\end{solve}

\begin{example}
\label{ex:oscillatory-log-improper-integral}
固定 $A>1$，讨论参数 $p>0$ 时
\[
    I_p=\int_A^{+\infty}\ln\left(1+\frac{\sin x}{x^p}\right)\,\mathrm{d}x
\]
的收敛性与绝对收敛性。
\end{example}

\begin{solve}
由于 $|\sin x|/x^p<1$，对数在整个积分区间上有定义。令 $u(x)=\sin x/x^p$，并设
\[
    d(x)=u(x)-\ln(1+u(x)).
\]
由 $\ln(1+u)\leq u$ 以及泰勒展开，在充分大的 $x$ 上有
\[
    d(x)=\frac{u(x)^2}{2}\bigl(1+O(u(x))\bigr),\qquad
    \frac14\frac{\sin^2x}{x^{2p}}
    \leq d(x)\leq\frac34\frac{\sin^2x}{x^{2p}}.
\]
在 $u=0$ 的点两侧均为零，上述不等式仍成立，因而不必在那里作除法。对任意 $p>0$，$\int_A^{+\infty}u(x)\,\mathrm{d}x$ 均由狄利克雷判别法收敛，而
\[
    \int_A^X\frac{\sin^2x}{x^{2p}}\,\mathrm{d}x
    =\frac12\int_A^X\frac{\mathrm{d}x}{x^{2p}}
     -\frac12\int_A^X\frac{\cos2x}{x^{2p}}\,\mathrm{d}x.
\]
振荡项对全部 $p>0$ 收敛，故非负积分 $\int_A^{+\infty}d$ 当且仅当 $2p>1$ 时收敛。于是 $0<p\leq1/2$ 时，$I_p$ 的截断积分趋于 $-\infty$；$p>1/2$ 时，$I_p$ 收敛。

为区分绝对收敛，注意 $|\ln(1+u)|$ 在 $u\to0$ 时与 $|u|$ 可双向作常数倍比较。$p>1$ 时，$|u(x)|\leq x^{-p}$，故绝对收敛。$0<p\leq1$ 时，在每个区间 $[n\pi+\pi/6,n\pi+5\pi/6]$ 上均有 $|\sin x|\geq1/2$，所以
\[
    \int_{n\pi}^{(n+1)\pi}\frac{|\sin x|}{x^p}\,\mathrm{d}x
    \geq\frac{\pi}{3((n+1)\pi)^p}.
\]
对充分大的连续整数 $n$ 求和，并与 $\int t^{-p}\,\mathrm{d}t$ 比较，得到绝对值积分发散。因此最终分类为：$p>1$ 时绝对收敛，$1/2<p\leq1$ 时条件收敛，$0<p\leq1/2$ 时发散。

尽管在 $\sin x\neq0$ 且 $x\to+\infty$ 时有 $\ln(1+u(x))/u(x)\to1$，当 $p\leq1/2$ 时它的积分与 $u$ 的积分却不同敛散。首阶振荡项的抵消使二阶的非负累积量成为关键，这正是变号情形不能直接套用等价比较的原因。
\end{solve}

\section{瑕积分}
\label{sec:singular-integrals}

\subsection{瑕积分的概念}
\label{subsec:concept-of-singular-integrals}

有些函数只在有限区间的端点或少数内点附近无界，在避开这些点的闭区间上仍可积。无界性使它不满足普通黎曼积分的必要条件，却不排除截断积分有有限极限。若函数在某点某一侧的每个充分小去心邻域内均无界，就称该点为该侧的\emph{瑕点}（singular point）；瑕积分考察的正是逼近瑕点时能否得到确定的有限累积量。

\begin{definition}[端点处的瑕积分]
\label{def:singular-integrals-at-endpoints}
设 $f$ 在 $(a,b]$ 上有定义，在每个 $[a+\varepsilon,b]$ 上可积，其中 $0<\varepsilon<b-a$。若有限极限存在，则定义
\begin{equation}
    \int_a^b f(x)\,\mathrm{d}x
    =\lim_{\varepsilon\downarrow0}
        \int_{a+\varepsilon}^b f(x)\,\mathrm{d}x,
    \label{eq:singular-integral-left-endpoint}
\end{equation}
并称相应瑕积分收敛；否则称其发散。若无界端点为 $b$，则相应定义为
\begin{equation}
    \int_a^b f(x)\,\mathrm{d}x
    =\lim_{\varepsilon\downarrow0}
        \int_a^{b-\varepsilon} f(x)\,\mathrm{d}x.
    \label{eq:singular-integral-right-endpoint}
\end{equation}
\end{definition}

瑕点处是否定义函数值及如何赋值不影响上述极限。另一方面，单纯在端点缺少定义却可有界延拓的函数，可能直接补成普通可积函数，并不一定构成真正的无界瑕点。

若两端都是瑕点，应在 $c\in(a,b)$ 处分成 $\int_a^c f+\int_c^b f$，要求两部分分别收敛。若内部点 $c$ 为瑕点，则同样定义
\[
    \int_a^b f(x)\,\mathrm{d}x
    =\lim_{\varepsilon\downarrow0}\int_a^{c-\varepsilon}f(x)\,\mathrm{d}x
     +\lim_{\eta\downarrow0}\int_{c+\eta}^b f(x)\,\mathrm{d}x,
\]
两个极限必须独立存在；有限多个瑕点时依此逐段处理。若另有无穷端点，也须分开检验每个异常端点，不能让一个瑕点附近的发散与另一个端点的发散互相抵消。

对内部瑕点，也可另外定义\emph{柯西主值}
\[
    \operatorname{PV}\int_a^b f(x)\,\mathrm{d}x
    =\lim_{\varepsilon\downarrow0}
       \left(\int_a^{c-\varepsilon}f(x)\,\mathrm{d}x
             +\int_{c+\varepsilon}^b f(x)\,\mathrm{d}x\right),
\]
这里假定其余端点正常。它与定义~\ref{def:cauchy-principal-value-at-infinity} 一样，只考察特定的同步截断。例如 $\operatorname{PV}\int_{-1}^1\mathrm{d}x/x=0$，但两侧积分分别发散，所以通常的瑕积分不存在。

若 $f$ 在 $(a,b]$ 上局部可积且有原函数 $F$，则逐个截断后使用牛顿--莱布尼茨公式得到
\begin{equation}
    \int_a^b f(x)\,\mathrm{d}x
    =F(b)-\lim_{x\to a+}F(x),
    \label{eq:singular-newton-leibniz}
\end{equation}
其中右端要求原函数的单侧极限为有限实数；右端点为瑕点时作相应修改。原函数有界本身并不足以保证该极限存在，关键仍是极限而非有界性。

\begin{example}
\label{ex:unbounded-derivative-singular-integral}
设 $F(0)=0$，$F(x)=x^2\sin(1/x^2)$（$0<|x|\leq1$），讨论其导函数的积分。
\end{example}

\begin{solve}
由导数定义有 $F'(0)=\lim_{h\to0}h\sin(1/h^2)=0$，而对 $x>0$，
\[
    f(x)=F'(x)=2x\sin\frac1{x^2}-\frac2x\cos\frac1{x^2}.
\]
取 $x_n=(2\pi n)^{-1/2}$，便有 $f(x_n)=-2/x_n\to-\infty$，故 $f$ 在零点附近无界，不能在 $[0,1]$ 上作普通黎曼积分。但它在每个 $[\varepsilon,1]$ 上连续，且 $F(\varepsilon)\to0$，所以
\[
    \int_0^1 f(x)\,\mathrm{d}x
    =\lim_{\varepsilon\downarrow0}
        \bigl(F(1)-F(\varepsilon)\bigr)=\sin1.
\]
这说明有原函数并不保证普通黎曼可积，却可以借原函数的端点极限计算瑕积分。
\end{solve}

\begin{example}
\label{ex:singular-arccosine-and-logarithm}
计算 $\int_{-1}^1\frac{\arccos x}{\sqrt{1-x^2}}\,\mathrm{d}x$ 与 $\int_0^1\ln x\,\mathrm{d}x$。
\end{example}

\begin{solve}
第一问的被积函数在 $x\to1-$ 时极限为 $1$，可作连续延拓；在 $x=-1$ 附近才有无界瑕点。以 $-\tfrac12(\arccos x)^2$ 为原函数，逐端取极限可得
\[
    \int_{-1}^1\frac{\arccos x}{\sqrt{1-x^2}}\,\mathrm{d}x
    =0-\left(-\frac{\pi^2}{2}\right)=\frac{\pi^2}{2}.
\]
第二问在 $x=0$ 附近无界。先在 $[\varepsilon,1]$ 上分部积分，再利用 $\varepsilon\ln\varepsilon\to0$，得到
\[
    \int_0^1\ln x\,\mathrm{d}x
    =\lim_{\varepsilon\downarrow0}
        [x\ln x-x]_\varepsilon^1=-1.
    \qedhere
\]
\end{solve}

\begin{example}
\label{ex:singular-power-integral-benchmark}
设 $a<b$、$p\in\mathbb{R}$，讨论 $\int_a^b(b-x)^{-p}\,\mathrm{d}x$ 的敛散性。
\end{example}

\begin{solve}
若 $p=1$，截断积分为 $\ln(b-a)-\ln\varepsilon\to+\infty$；若 $p\neq1$，则
\[
    \int_a^{b-\varepsilon}\frac{\mathrm{d}x}{(b-x)^p}
    =\frac{(b-a)^{1-p}-\varepsilon^{1-p}}{1-p}.
\]
因此
\begin{equation}
    \int_a^b\frac{\mathrm{d}x}{(b-x)^p}
    \begin{cases}
        =\dfrac{(b-a)^{1-p}}{1-p},&p<1,\\
        \text{发散至 }+\infty,&p\geq1.
    \end{cases}
    \label{eq:singular-power-integral-benchmark}
\end{equation}
$p\leq0$ 时被积函数在端点可连续延拓，实际为普通定积分；$0<p<1$ 时虽无界却可积。左端点的 $\int_a^b(x-a)^{-p}\,\mathrm{d}x$ 有相同结论。与式~\eqref{eq:infinite-power-integral-benchmark} 对照，无穷远处要求 $p>1$，有限瑕点附近要求 $p<1$，不能混用。
\end{solve}

两类广义积分可以通过换元相互转化。对右端点 $b$，令 $t=1/(b-x)$，先在截断区间上换元再取极限，有
\begin{equation}
    \int_a^b f(x)\,\mathrm{d}x
    =\int_{1/(b-a)}^{+\infty}
        f\left(b-\frac1t\right)\frac{\mathrm{d}t}{t^2}.
    \label{eq:singular-to-infinite-integral}
\end{equation}
相应地，左端点的瑕积分可写成 $\int_{1/(b-a)}^{+\infty}f(a+1/t)t^{-2}\,\mathrm{d}t$。反过来，对 $a>0$，令 $t=1/x$ 可得
\[
    \int_a^{+\infty}f(x)\,\mathrm{d}x
    =\int_0^{1/a}f(1/t)\frac{\mathrm{d}t}{t^2}.
\]
这些等式均表示两端极限的等价性；右端可能是瑕积分，也可能在补定义后成为普通定积分。尤其对幂函数，式~\eqref{eq:singular-to-infinite-integral} 把 $(b-x)^{-p}$ 化为 $t^{p-2}$，无穷积分条件 $2-p>1$ 正好等价于瑕积分条件 $p<1$。瑕积分的换元和分部积分也应遵循同一原则：先在避开瑕点的闭区间上作普通运算，再核对换元后的端点极限；分部积分还需检验边界乘积的有限极限，不能直接在瑕点处代入无定义的量。

\subsection{瑕积分敛散性的判别法}
\label{subsec:convergence-tests-for-singular-integrals}

以下以 $b$ 为唯一可能的无界端点，设所有函数均在 $[a,b)$ 上局部可积。左端点的结论可由反射换元 $x=a+b-t$ 得到。若不便直接求原函数或换元后的表达式更复杂，可以在原区间上使用与无穷积分平行的尾部判据。

\begin{theorem}[瑕积分的柯西准则与绝对收敛]
\label{thm:cauchy-and-absolute-singular-integrals}
$\int_a^b f(x)\,\mathrm{d}x$ 收敛的充分必要条件是：对任意 $\varepsilon>0$，存在 $\delta\in(0,b-a)$，使任意 $b-\delta<u<v<b$ 均有
\begin{equation}
    \left|\int_u^v f(x)\,\mathrm{d}x\right|<\varepsilon.
    \label{eq:singular-integral-cauchy-tail}
\end{equation}
若 $\int_a^b|f(x)|\,\mathrm{d}x$ 收敛，称原瑕积分绝对收敛，此时原积分必收敛；若原积分收敛而绝对值积分发散，则称其条件收敛。
\end{theorem}

\begin{proof}
令 $F(t)=\int_a^t f$，其在 $t\to b-$ 时的函数极限柯西准则恰为式~\eqref{eq:singular-integral-cauchy-tail}。绝对收敛蕴含收敛仍由 $|\int_u^v f|\leq\int_u^v|f|$ 与柯西准则得到。
\end{proof}

\begin{theorem}[瑕积分的比较判别法]
\label{thm:comparison-tests-singular-integrals}
设 $f,g\geq0$。若存在 $\delta_0\in(0,b-a)$、$C>0$，使 $b-\delta_0<x<b$ 时有 $f(x)\leq Cg(x)$，则 $\int_a^b g$ 收敛蕴含 $\int_a^b f$ 收敛，$\int_a^b f$ 发散蕴含 $\int_a^b g$ 发散。

若 $g>0$ 且 $\lim_{x\to b-}f(x)/g(x)=\ell$，则 $0<\ell<+\infty$ 时两积分同敛散；$\ell=0$ 时由 $\int_a^b g$ 收敛可推得 $\int_a^b f$ 收敛；$\ell=+\infty$ 时由 $\int_a^b g$ 发散可推得 $\int_a^b f$ 发散。
\end{theorem}

\begin{proof}
对非负函数，$\int_a^t f$ 随 $t\uparrow b$ 单调非减，有限极限的存在仍等价于这些截断积分有共同上界。因此邻近端点的常数倍比较可转化为截断积分的比较，而远离端点的部分只是有限定积分。极限形式通过在端点附近建立常数倍不等式推出，与定理~\ref{thm:comparison-tests-improper-integrals} 的证明相同。
\end{proof}

最基本的比较对象是 $(b-x)^{-p}$，其收敛阈值见式~\eqref{eq:singular-power-integral-benchmark}。例如 $f\geq0$ 且 $f(x)\sim C(b-x)^{-p}$、$C>0$ 时，积分当且仅当 $p<1$ 时收敛。比较前应先确定全部瑕点，并在每个瑕点附近分别找主项；函数在某一端的等价式不能说明另一端的行为。

\begin{example}
\label{ex:two-endpoint-singular-power-integral}
讨论 $\int_{-1}^1(1-x^2)^{-p}\,\mathrm{d}x$ 的敛散性，其中 $p\in\mathbb{R}$。
\end{example}

\begin{solve}
在 $x\to1-$ 与 $x\to-1+$ 时分别有
\[
    (1-x^2)^{-p}\sim2^{-p}(1-x)^{-p},\qquad
    (1-x^2)^{-p}\sim2^{-p}(1+x)^{-p}.
\]
在 $0$ 处分开后，两部分由同一个幂函数判据控制，故原积分当且仅当 $p<1$ 时收敛；$p\geq1$ 时两端均发散。$p\leq0$ 时端点可作有界连续延拓，属于普通可积情形，不能不加区分地说所有参数下都有两个无界瑕点。
\end{solve}

\begin{theorem}[瑕积分的狄利克雷与阿贝尔判别法]
\label{thm:dirichlet-abel-singular-integrals}
若存在 $M>0$，使任意 $a\leq t<b$ 均有 $|\int_a^t f(x)\,\mathrm{d}x|\leq M$，而 $g$ 在 $[a,b)$ 上单调且 $\lim_{x\to b-}g(x)=0$，则 $\int_a^b f(x)g(x)\,\mathrm{d}x$ 收敛；这是瑕积分的狄利克雷判别法。若 $\int_a^b f(x)\,\mathrm{d}x$ 已收敛，而 $g$ 在 $[a,b)$ 上单调有界，则乘积积分也收敛；这是瑕积分的阿贝尔判别法。
\end{theorem}

\begin{proof}
第一种情形下，对 $a\leq u<v<b$，有限区间上的第二积分中值定理与累积积分有界性给出
\[
    \left|\int_u^v fg\,\mathrm{d}x\right|
    \leq2M\bigl(|g(u)|+|g(v)|\bigr).
\]
当 $u,v\to b-$ 时右端一致趋于零，故由瑕积分的柯西准则收敛。第二种情形下，$g$ 有有限端点极限 $\ell$，而 $\int_a^t f$ 有界；先把第一种情形用于 $f(g-\ell)$，再加回已收敛的 $\ell\int_a^b f$ 即可。
\end{proof}

\begin{example}
\label{ex:mixed-logarithmic-improper-integral}
讨论含无穷端点与有限异常端点的积分
\[
    \int_0^{+\infty}\frac{\ln(1+x)}{x^\alpha}\,\mathrm{d}x,
    \qquad\alpha\in\mathbb{R}.
\]
\end{example}

\begin{solve}
在 $1$ 处分成 $\int_0^1$ 与 $\int_1^{+\infty}$。当 $x\to0+$ 时，$\ln(1+x)\sim x$，故
\[
    \frac{\ln(1+x)}{x^\alpha}\sim x^{1-\alpha}.
\]
由左端点的幂函数判据，零点附近的积分当且仅当 $\alpha<2$ 时收敛；其中 $\alpha\leq1$ 时可连续延拓，$1<\alpha<2$ 时虽无界却可积。当 $x\to+\infty$ 时，$\ln(1+x)\sim\ln x$，由式~\eqref{eq:power-log-infinite-integral-test}，无穷尾部当且仅当 $\alpha>1$ 时收敛。两项必须同时满足，故原积分收敛的充分必要条件为
\[
    1<\alpha<2.
\]
临界值 $\alpha=1$ 在无穷端发散，$\alpha=2$ 在零点端发散；将两类端点混在同一个截断过程中计算，可能掩盖其中一个条件。
\end{solve}

\section*{习题八}
\addcontentsline{toc}{section}{习题八}
\fi

\selectchapter{9}
\ifcompilechapter
\chapter{数项级数}
\label{chap:numerical-series}

% 本章依据《数学分析》第九章正文整理：数学分析-9.pdf，第 1--42 页（书内第 123--164 页）；不含习题九。
有限个数相加只涉及代数运算，无穷多个数相加则必须借助极限：先取有限项的和，再研究这些和是否趋向一个确定的实数。因此，数项级数本质上是部分和序列的极限问题，却又因通项的结构而具有专门的判别方法。本章先讨论收敛的含义及正项级数，再研究正负项抵消所产生的条件收敛，最后说明分组、重排、相乘及无穷乘积的适用条件；这些内容也是以后研究函数项级数的基础。

\section{数项级数的基本概念}
\label{sec:basic-concepts-of-numerical-series}

\subsection{数项级数的基本概念}
\label{subsec:basic-concepts-of-numerical-series}

\begin{definition}
设 $\{a_n\}_{n\geq1}$ 为实数列，形式上的无穷和
\[
    a_1+a_2+\cdots+a_n+\cdots=\sum_{n=1}^{\infty}a_n
\]
称为\emph{数项级数}（numerical series），$a_n$ 称为\emph{通项}（general term）。令
\begin{equation}
    S_n=\sum_{k=1}^n a_k,\qquad n\geq1,
    \label{eq:numerical-series-partial-sums}
\end{equation}
称 $S_n$ 为第 $n$ 个\emph{部分和}（partial sum）。若 $S_n\to S\in\mathbb R$，则称级数\emph{收敛}（convergent），并记 $\sum_{n=1}^{\infty}a_n=S$，其中 $S$ 为级数的和；否则称级数\emph{发散}（divergent）。
\end{definition}

$\sum_{n=1}^{\infty}a_n$ 既可表示待研究的级数，也可在收敛已经确立后表示其和，使用时须区分。发散包含部分和趋于 $+\infty$、趋于 $-\infty$ 以及没有极限等情形；只有前两种情形才分别称为发散到 $+\infty$、$-\infty$。级数的通项序列与部分和序列也不能混淆：前者描述每次增加多少，后者才决定无穷和是否存在。

\begin{example}
\label{ex:geometric-numerical-series}
讨论\emph{几何级数}（geometric series，也称等比级数）$\sum_{n=0}^{\infty}q^n$ 的敛散性，其中 $q\in\mathbb R$，约定零次幂为 $1$。
\end{example}

\begin{solve}
取前 $n$ 项之和。当 $q\neq1$ 时，
\[
    \sum_{k=0}^{n-1}q^k=\frac{1-q^n}{1-q}.
\]
故 $|q|<1$ 时级数收敛，且和为 $1/(1-q)$。$q=1$ 时部分和为 $n$；$q=-1$ 时部分和在 $1$ 与 $0$ 之间交替；$q>1$ 时部分和趋于 $+\infty$；$q<-1$ 时部分和的两个奇偶子列分别趋于相反的无穷。因此收敛的充分必要条件是 $|q|<1$，而 $q<-1$ 时不能把发散误写为趋于某一个带符号的无穷。
\end{solve}

\begin{example}
\label{ex:telescoping-shifted-harmonic-series}
设 $m$ 为正整数，证明 $\sum_{n=1}^{\infty}1/(n(n+m))$ 收敛，并求其和。
\end{example}

\begin{solve}
把通项拆为 $m^{-1}(1/n-1/(n+m))$，有限求和后相消，得
\[
    \sum_{k=1}^{n}\frac1{k(k+m)}
    =\frac1m\left(\sum_{k=1}^{m}\frac1k
                     -\sum_{k=n+1}^{n+m}\frac1k\right).
\]
右侧最后一个有限和介于 $0$ 与 $m/(n+1)$ 之间，趋于零，因而
\[
    \sum_{n=1}^{\infty}\frac1{n(n+m)}=\frac1m\sum_{k=1}^{m}\frac1k.
    \qedhere
\]
\end{solve}

\begin{example}
\label{ex:factorial-series-for-e}
由 $e=\lim_{m\to\infty}(1+1/m)^m$ 出发，证明 $e=\sum_{k=0}^{\infty}1/k!$，并估计截断误差。
\end{example}

\begin{solve}
令 $T_n=\sum_{k=0}^n1/k!$。由 $k!\geq2^{k-1}$（$k\geq1$）知 $T_n\leq3$，故其单调递增而有有限极限 $T$。对固定的 $n$ 及 $m\geq n$，二项式展开给出
\[
    \sum_{k=0}^n\frac1{k!}
       \prod_{j=0}^{k-1}\left(1-\frac jm\right)
    \leq\left(1+\frac1m\right)^m
    \leq\sum_{k=0}^m\frac1{k!}=T_m,
\]
其中空乘积取 $1$。先令 $m\to\infty$，得 $T_n\leq e\leq T$；再令 $n\to\infty$，即得 $T=e$。对 $n\geq1$，利用 $(n+j)!\geq n!(n+1)^j$，且 $j\geq2$ 时严格大于，可得
\begin{equation}
    0<e-T_n=\sum_{j=1}^{\infty}\frac1{(n+j)!}
      <\frac1{n!}\sum_{j=1}^{\infty}\frac1{(n+1)^j}
      =\frac1{n\,n!}.
    \label{eq:factorial-series-tail-bound}
\end{equation}
这里对无穷尾和的估计可先对有限尾和成立，再取极限。这一估计也可用于证明 $e$ 无理：若 $e=p/q$，取整数 $n\geq\max\{q,2\}$，则 $n!(e-T_n)$ 为整数，却严格介于 $0$ 与 $1/n$ 之间，矛盾。\crosschapterref{5}{例~\ref{ex:taylor-approximation-and-irrationality-of-e}}{用泰勒公式近似 $e$ 的例题} 给出了利用泰勒余项的另一种证明。
\end{solve}

\begin{proposition}[有限项改变与线性运算]
\label{prop:series-finite-changes-and-linearity}
改变级数有限项的值不改变其敛散性，但一般会改变其和。对非零常数 $c$，$\sum a_n$ 与 $\sum ca_n$ 同敛散。若 $\sum a_n=A$、$\sum b_n=B$ 都收敛，则对任意 $\alpha,\beta\in\mathbb R$，
\[
    \sum_{n=1}^{\infty}(\alpha a_n+\beta b_n)=\alpha A+\beta B.
\]
\end{proposition}

\begin{proof}
改变有限项后，充分靠后的部分和只相差一个固定常数；乘以非零常数后，部分和也乘以该常数。这就证明前两项结论。最后一项来自有限和的线性及收敛序列的极限运算法则。
\end{proof}

线性公式要求相关级数已经收敛，不能把两个发散级数的形式相减当作两个无穷和之差。有限项改变的性质也说明，判别敛散性时往往只需核对充分靠后的项；但求和时必须把被删去的初始项补回。

\subsection{柯西准则}
\label{subsec:cauchy-criterion-for-series}

\begin{theorem}[级数的柯西准则]
\label{thm:cauchy-criterion-for-numerical-series}
实数项级数 $\sum_{n=1}^{\infty}a_n$ 收敛的充分必要条件是
\begin{equation}
    \forall\varepsilon>0,\quad\exists N\in\mathbb N,\quad
    \forall n\geq N,\ \forall p\geq1,\quad
    \left|\sum_{k=n+1}^{n+p}a_k\right|<\varepsilon,
    \label{eq:cauchy-condition-for-numerical-series}
\end{equation}
其中 $n,p$ 均为整数。
\end{theorem}

\begin{proof}
式~\eqref{eq:cauchy-condition-for-numerical-series} 中的尾段和等于 $S_{n+p}-S_n$，故这一条件恰好是部分和序列的柯西条件。应用\crosschapterref{2}{定理~\ref{thm:cauchy-convergence-criterion}}{序列的柯西收敛准则} 即得结论。
\end{proof}

\emph[caution]{这里的 $N$ 必须同时适用于所有尾段长度 $p$，不能在固定 $p$ 后才选取 $N$。}特别取 $p=1$，得到收敛的必要条件 $a_n\to0$；其逆命题不成立。例如调和级数 $\sum1/n$ 的通项趋于零，但
\[
    S_{2n}-S_n=\sum_{k=n+1}^{2n}\frac1k\geq\frac12,
\]
违反柯西条件，故发散。由于其部分和单调递增，实际上发散到 $+\infty$。

\begin{example}
\label{ex:series-term-test-divergence}
证明 $\sum_{n=1}^{\infty}n\ln(1+1/n)$ 发散；并讨论 $\sum_{n=1}^{\infty}\sin(n\theta)$，其中 $\theta\in\mathbb R$。
\end{example}

\begin{solve}
由 $\ln(1+t)\sim t$（$t\to0$），有 $n\ln(1+1/n)\to1$，故第一个级数发散。对第二个级数，若 $\theta\in\pi\mathbb Z$，各项均为零，级数收敛；若 $\theta\notin\pi\mathbb Z$ 而级数收敛，则必有 $\sin(n\theta)\to0$。由
\[
    \sin((n+1)\theta)
    =\cos\theta\sin(n\theta)+\sin\theta\cos(n\theta)
\]
及 $\sin\theta\neq0$ 推出 $\cos(n\theta)\to0$，与 $\sin^2(n\theta)+\cos^2(n\theta)=1$ 矛盾。因此第二个级数仅在 $\theta\in\pi\mathbb Z$ 时收敛。
\end{solve}

\section{正项级数}
\label{sec:positive-term-series}

\subsection{比较判别法}
\label{subsec:comparison-tests-for-series}

\begin{definition}
若 $\sum_{n=1}^{\infty}|a_n|$ 收敛，称 $\sum a_n$ \emph{绝对收敛}（absolutely convergent）；若 $\sum a_n$ 收敛而 $\sum|a_n|$ 发散，称其\emph{条件收敛}（conditionally convergent）。若每项 $a_n\geq0$，称 $\sum a_n$ 为\emph{正项级数}（positive-term series）；本书的这一名称允许通项为零，要求严格为正时另行说明。
\end{definition}

\begin{proposition}
\label{prop:absolute-convergence-implies-convergence}
绝对收敛级数必收敛。
\end{proposition}

\begin{proof}
由三角不等式，
\[
    \left|\sum_{k=n+1}^{n+p}a_k\right|
    \leq\sum_{k=n+1}^{n+p}|a_k|.
\]
右端由 $\sum|a_n|$ 的柯西条件一致趋于零，故 $\sum a_n$ 也满足式~\eqref{eq:cauchy-condition-for-numerical-series}。
\end{proof}

\begin{theorem}[正项级数的有界性判据]
\label{thm:positive-series-bounded-partial-sums}
正项级数收敛的充分必要条件是其部分和序列有上界；若发散，则部分和趋于 $+\infty$。
\end{theorem}

\begin{proof}
正项级数的部分和单调递增，故由单调收敛原理，有上界时收敛，无上界时趋于 $+\infty$。反向结论来自收敛序列必有界。
\end{proof}

\begin{theorem}[比较判别法]
\label{thm:series-comparison-test}
设 $a_n,b_n\geq0$，并且存在 $C>0$ 及正整数 $N$，使 $n\geq N$ 时 $a_n\leq Cb_n$。若 $\sum b_n$ 收敛，则 $\sum a_n$ 收敛；若 $\sum a_n$ 发散，则 $\sum b_n$ 发散。
\end{theorem}

\begin{proof}
忽略有限个初始项后，部分和满足 $\sum_{k=N}^n a_k\leq C\sum_{k=N}^n b_k$。由定理~\ref{thm:positive-series-bounded-partial-sums}，右侧有界蕴含左侧有界；第二个结论是第一个结论的逆否命题。
\end{proof}

\begin{corollary}[极限比较判别法]
\label{cor:series-limit-comparison-test}
设 $a_n\geq0$、$b_n>0$，且 $a_n/b_n\to\ell\in[0,+\infty]$。若 $0<\ell<+\infty$，则两级数同敛散；若 $\ell=0$，则 $\sum b_n$ 收敛可推出 $\sum a_n$ 收敛；若 $\ell=+\infty$，则 $\sum b_n$ 发散可推出 $\sum a_n$ 发散。
\end{corollary}

\begin{proof}
$0<\ell<+\infty$ 时，充分大的 $n$ 满足 $(\ell/2)b_n\leq a_n\leq(3\ell/2)b_n$，可向两个方向比较。$\ell=0$ 时最终有 $a_n\leq b_n$；$\ell=+\infty$ 时最终有 $a_n\geq b_n$，分别应用比较判别法。
\end{proof}

\emph[caution]{极限为零或无穷时只得到所述的单向结论。}例如 $1/n^2=o(1/n)$，却不能从调和级数发散推出 $\sum1/n^2$ 发散。比较判别法中的非负条件也至关重要：一般任意项级数的抵消不能仅由通项大小控制。

\begin{example}
\label{ex:p-series-convergence}
讨论\emph{$p$ 级数}（$p$-series）$\sum_{n=1}^{\infty}n^{-p}$，其中 $p\in\mathbb R$。
\end{example}

\begin{solve}
$p\leq1$ 时，$n^{-p}\geq1/n$，故发散。$p>1$ 时按二进区间估计，每个 $j\geq0$ 满足
\[
    \sum_{n=2^j}^{2^{j+1}-1}\frac1{n^p}
    \leq 2^j\frac1{2^{jp}}=2^{-j(p-1)}.
\]
右侧关于 $j$ 构成收敛的几何级数，因而所有原级数部分和都有同一个上界，故收敛。因此 $p$ 级数恰在 $p>1$ 时收敛。这里使用的是有限分组估计与部分和有界性，不预先假定无穷分组可以任意拆合。
\end{solve}

\begin{example}
\label{ex:comparison-series-fibonacci-and-logcos}
设 $F_1=1$、$F_2=2$、$F_{n+2}=F_{n+1}+F_n$，证明 $\sum1/F_n$ 收敛；并判断 $\sum_{n=1}^{\infty}[-\ln\cos(1/n)]$ 的敛散性。
\end{example}

\begin{solve}
$F_n$ 为递增正数列，故 $F_{n+2}\geq2F_n$。由此 $F_{2k-1}\geq2^{k-1}$、$F_{2k}\geq2^k$，两个奇偶子列的倒数都可用几何级数控制，原正项级数的部分和因而有界。对第二个级数，$0<\cos(1/n)<1$，并且
\[
    -\ln\cos\frac1n
    \sim1-\cos\frac1n\sim\frac1{2n^2}.
\]
由极限比较判别法，该级数收敛。识别等价主项后与几何级数、$p$ 级数比较，是处理正项级数的基本方法。
\end{solve}

\subsection{达朗贝尔判别法与柯西判别法}
\label{subsec:dalembert-and-cauchy-tests}

几何级数的通项之比及通项的 $n$ 次方根都反映其指数衰减率。因此，考察一般正项级数的相邻项比值或根值，可以把它与几何级数联系起来。以下判别式中的上、下极限均允许取扩充实数值。

\begin{theorem}[达朗贝尔比值判别法]
\label{thm:dalembert-ratio-test}
设 $a_n>0$，记
\[
    r^*=\limsup_{n\to\infty}\frac{a_{n+1}}{a_n},\qquad
    r_*=\liminf_{n\to\infty}\frac{a_{n+1}}{a_n}.
\]
若 $r^*<1$，则 $\sum a_n$ 收敛；若 $r_*>1$，则发散。特别地，若比值有极限 $r$，则 $r<1$ 时收敛，$r>1$ 时发散；$r=1$ 时本判别法没有结论。此法称为\emph{达朗贝尔判别法}（d'Alembert's ratio test）。
\end{theorem}

\begin{proof}
若 $r^*<1$，取 $r^*<q<1$，则从某个 $N$ 开始有 $a_{n+1}\leq qa_n$，从而 $a_n\leq a_Nq^{n-N}$，可与收敛几何级数比较。若 $r_*>1$，取 $1<q<r_*$，则最终有 $a_n\geq a_Nq^{n-N}$，通项不趋于零，故发散。
\end{proof}

\begin{theorem}[柯西根值判别法]
\label{thm:cauchy-root-test}
设 $a_n\geq0$，记 $\rho=\limsup_{n\to\infty}\sqrt[n]{a_n}$。若 $\rho<1$，则 $\sum a_n$ 收敛；若 $\rho>1$，则发散；$\rho=1$ 时没有结论。此法称为\emph{柯西判别法}（Cauchy's root test）。
\end{theorem}

\begin{proof}
若 $\rho<1$，取 $\rho<q<1$，则最终有 $\sqrt[n]{a_n}<q$，即 $a_n<q^n$。若 $\rho>1$，取 $1<q<\rho$，则有无穷多个 $n$ 满足 $a_n>q^n$，通项不可能趋于零。注意发散部分只需这样的下标无穷多，不要求所有充分大的下标都满足该不等式。
\end{proof}

对任意实数 $p$，$a_n=n^{-p}$ 的比值与根值都趋于 $1$，但其级数可收敛也可发散，因此两个判别法的临界值都不能补上等号。任意项级数应用这两种方法时，通常先对 $|a_n|$ 作判断：得到收敛即为绝对收敛；因根值或比值大于 $1$ 而得到的发散，则来自原通项也不趋于零。

对严格正数列还有
\begin{equation}
    \liminf_{n\to\infty}\frac{a_{n+1}}{a_n}
    \leq\liminf_{n\to\infty}\sqrt[n]{a_n}
    \leq\limsup_{n\to\infty}\sqrt[n]{a_n}
    \leq\limsup_{n\to\infty}\frac{a_{n+1}}{a_n}.
    \label{eq:series-ratio-root-limits}
\end{equation}
证明时，任取大于最右端上极限的有限正数 $q$，最终有 $a_{n+1}\leq qa_n$，故 $\limsup\sqrt[n]{a_n}\leq q$，再让 $q$ 下降到该上极限；最右端为 $+\infty$ 时此侧不等式自动成立。左侧同理，对小于比值下极限的正数作几何下界；下极限为零时直接利用根值非负。因而比值有极限时根值也有相同极限，根值判别法能够覆盖比值判别法的上述有效情形。

\begin{example}
\label{ex:factorial-power-ratio-test}
讨论 $\sum_{n=1}^{\infty}z^nn!/n^n$，其中 $z\geq0$。
\end{example}

\begin{solve}
$z=0$ 时级数各项为零。$z>0$ 时记通项为 $a_n$，则
\[
    \frac{a_{n+1}}{a_n}=z\left(\frac n{n+1}\right)^n\longrightarrow\frac ze.
\]
所以 $0<z<e$ 时收敛，$z>e$ 时发散。$z=e$ 时比值判别法失效，须另行分析；由 $(1+1/n)^n<e$，有 $a_{n+1}/a_n>1$，故 $a_n\geq a_1>0$，仍发散。因此原级数恰在 $0\leq z<e$ 时收敛。
\end{solve}

\begin{example}
\label{ex:root-test-with-oscillating-ratios}
设 $0<a<b<1$，奇数 $n$ 时令 $c_n=a^n$，偶数 $n$ 时令 $c_n=b^n$。判断 $\sum c_n$ 的敛散性。
\end{example}

\begin{solve}
$\sqrt[n]{c_n}$ 按奇偶分别为 $a,b$，故根值上极限为 $b<1$，级数收敛。另一方面，$c_{n+1}/c_n$ 沿奇数下标趋于 $+\infty$，沿偶数下标趋于零，故比值判别法的两个充分条件均不满足。这说明根值可以保留被相邻项比值的振荡掩盖的整体衰减信息。
\end{solve}

\subsection{拉贝判别法}
\label{subsec:raabe-test}

当相邻项比值趋于 $1$ 时，几何级数的比较通常不再敏感。对参考通项 $a_n=n^{-p}$，有
\[
    n\left(\frac{a_n}{a_{n+1}}-1\right)
    =n\left[\left(1+\frac1n\right)^p-1\right]\longrightarrow p.
\]
因此，考察比值偏离 $1$ 的一阶程度，可以转而与 $p$ 级数比较。

\begin{theorem}[拉贝判别法]
\label{thm:raabe-series-test}
设 $a_n>0$，令 $R_n=n(a_n/a_{n+1}-1)$。若 $\liminf R_n>1$，则 $\sum a_n$ 收敛；若 $\limsup R_n<1$，则发散。特别地，$R_n\to r$ 时，$r>1$ 可判收敛，$r<1$ 可判发散，而 $r=1$ 一般没有结论。此法称为\emph{拉贝判别法}（Raabe's test）。另外，若最终恒有 $R_n\leq1$，即使 $R_n\to1$，仍可判定级数发散。
\end{theorem}

\begin{proof}
若 $\liminf R_n>1$，选取实数 $1<p<c<\liminf R_n$。由 $(1+1/n)^p=1+p/n+o(1/n)$，充分大的 $n$ 满足
\[
    \frac{a_n}{a_{n+1}}\geq1+\frac cn
      >\left(1+\frac1n\right)^p,
\]
即 $(n+1)^pa_{n+1}<n^pa_n$。所以 $n^pa_n$ 从某项起有上界，$a_n\leq Cn^{-p}$，级数收敛。

若最终有 $R_n\leq1$，则 $a_n/a_{n+1}\leq(n+1)/n$，即 $(n+1)a_{n+1}\geq na_n$。于是最终有 $a_n\geq C/n$，其中 $C>0$，故级数发散。$\limsup R_n<1$ 蕴含这一最终不等式，发散部分随即成立。
\end{proof}

\begin{example}
\label{ex:raabe-double-factorial-series}
设 $p>0$，讨论
\[
    \sum_{n=1}^{\infty}\left[\frac{(2n-1)!!}{(2n)!!}\right]^p,
\]
其中 $(2n-1)!!=1\cdot3\cdots(2n-1)$，$(2n)!!=2\cdot4\cdots2n$。
\end{example}

\begin{solve}
记通项为 $a_n$，则
\[
    R_n=n\left[\left(\frac{2n+2}{2n+1}\right)^p-1\right]
        \longrightarrow\frac p2.
\]
故 $p>2$ 时收敛，$0<p<2$ 时发散。$p=2$ 时极限为 $1$，但进一步计算有
\[
    R_n=\frac{n(4n+3)}{(2n+1)^2}<1,
\]
可由定理~\ref{thm:raabe-series-test} 的最终不等式判定发散。因此恰在 $p>2$ 时收敛。临界情形的结论来自额外的不等式，而不是把极限判据擅自延伸到 $r=1$。
\end{solve}

\subsection{柯西积分判别法}
\label{subsec:cauchy-integral-test}

\begin{theorem}[柯西积分判别法]
\label{thm:cauchy-integral-test-for-series}
设 $f:[1,+\infty)\to[0,+\infty)$ 单调不增，则正项级数 $\sum_{n=1}^{\infty}f(n)$ 与广义积分 $\int_1^{\infty}f(x)\,dx$ 同敛散。这称为\emph{柯西积分判别法}（Cauchy's integral test）。这里 $f$ 在每个有限闭区间上单调有界，故黎曼可积。
\end{theorem}

\begin{proof}
由 $f(n+1)\leq f(x)\leq f(n)$（$n\leq x\leq n+1$），得到
\[
    f(n+1)\leq\int_n^{n+1}f(x)\,dx\leq f(n).
\]
有限求和后有
\begin{equation}
    \sum_{k=2}^{n+1}f(k)
    \leq\int_1^{n+1}f(x)\,dx
    \leq\sum_{k=1}^n f(k).
    \label{eq:series-integral-comparison}
\end{equation}
两边都是非负量累积而成的单调序列，一边有界时另一边也有界，故级数与积分同时收敛。反之任一无界则另一也无界，均趋于 $+\infty$。
\end{proof}

积分判别法给出的是敛散性的等价，级数的和一般不等于积分值。若已知收敛，记 $r_n=\sum_{k=n+1}^{\infty}f(k)$，相同的区间估计还给出
\begin{equation}
    \int_{n+1}^{\infty}f(x)\,dx\leq r_n
       \leq\int_n^{\infty}f(x)\,dx,
    \label{eq:integral-test-series-tail-bound}
\end{equation}
可用于估计截断误差。单调性只需从某点起成立，有限个初始项不影响判别。

\begin{example}
\label{ex:logarithmic-p-series}
讨论 $\sum_{n=2}^{\infty}1/(n^p(\ln n)^q)$，其中 $p,q\in\mathbb R$。
\end{example}

\begin{solve}
$p>1$ 时，选取 $1<r<p$，由 $n^{r-p}(\ln n)^{-q}\to0$，可与收敛的 $\sum n^{-r}$ 比较。$p<1$ 时，选取 $p<r<1$，有 $n^{r-p}(\ln n)^{-q}\to+\infty$，可与发散的 $\sum n^{-r}$ 比较。这里使用了任意固定的对数幂都慢于正幂的增长。

$p=1$ 时，$f(x)=1/(x(\ln x)^q)$ 最终递减，令 $t=\ln x$，得
\[
    \int_2^{\infty}\frac{dx}{x(\ln x)^q}
      =\int_{\ln2}^{\infty}t^{-q}\,dt.
\]
后者恰在 $q>1$ 时收敛。因此原级数收敛的充分必要条件是 $p>1$，或 $p=1$ 且 $q>1$；其余情形均发散。特别地，$\sum1/(n(\ln n)^q)$ 的拉贝量趋于 $1$，仍可有两种不同的敛散性。
\end{solve}

\begin{example}
\label{ex:positive-series-tail-powered-denominator}
设 $a_n>0$ 且 $\sum a_n$ 收敛，令 $r_n=\sum_{k=n}^{\infty}a_k$。证明当 $p<1$ 时，$\sum a_n/r_n^p$ 收敛。
\end{example}

\begin{solve}
有 $r_n\downarrow0$、$a_n=r_n-r_{n+1}$。$0\leq p<1$ 时，因 $t^{-p}$ 单调不增，
\[
    \frac{a_n}{r_n^p}
      \leq\int_{r_{n+1}}^{r_n}t^{-p}\,dt.
\]
对 $n=1,\ldots,N$ 求和，得到部分和不超过
\[
    \int_{r_{N+1}}^{r_1}t^{-p}\,dt
      \leq\frac{r_1^{1-p}}{1-p}.
\]
故级数收敛。$p<0$ 时，$r_n^{-p}\leq r_1^{-p}$，直接与 $\sum a_n$ 比较即可。取 $0<p<1$ 时，新通项与 $a_n$ 的比值为 $r_n^{-p}\to+\infty$，说明在一个收敛正项级数之上仍可构造通项相对更大而保持收敛的级数。
\end{solve}

\section{任意项级数}
\label{sec:general-term-series}

对于正负项兼有的级数，绝对收敛仍可用正项级数的方法判定；但若绝对值级数发散，原级数还可能因抵消而收敛。研究这种情形，既要控制通项趋零，也要控制正负项累积抵消的方式。

\subsection{交错级数的敛散性}
\label{subsec:convergence-of-alternating-series}

先说明一种可以从分组级数返回原级数的情形。一般而言，分组后的部分和仅是原部分和的一个子列，其收敛不足以保证整个序列收敛；若每组长度有统一上界，且原通项趋零，则可控制组内遗漏的部分和。

\begin{theorem}[有界长度分组的逆向结论]
\label{thm:bounded-block-grouping-series}
设 $a_n\to0$，$0=n_0<n_1<n_2<\cdots$ 为整数，且存在正整数 $M$ 使 $n_k-n_{k-1}\leq M$。令
\[
    b_k=\sum_{j=n_{k-1}+1}^{n_k}a_j.
\]
若 $\sum b_k=S$，则 $\sum a_n=S$。
\end{theorem}

\begin{proof}
记原部分和为 $S_n$，则 $S_{n_k}=\sum_{j=1}^k b_j\to S$。给定 $n$，取 $k$ 使 $n_k\leq n<n_{k+1}$，便有
\[
    |S_n-S_{n_k}|
    \leq M\sup_{j\geq n_k+1}|a_j|\longrightarrow0.
\]
因 $n\to\infty$ 时 $k\to\infty$，故 $S_n\to S$。
\end{proof}

\begin{definition}
设 $b_n\geq0$，形如 $\sum_{n=1}^{\infty}(-1)^{n-1}b_n$ 或其相反数的级数称为\emph{交错级数}（alternating series）。
\end{definition}

\begin{theorem}[莱布尼茨判别法及余项估计]
\label{thm:leibniz-alternating-series-test}
设 $b_n\geq0$、$b_{n+1}\leq b_n$ 且 $b_n\to0$，则交错级数 $\sum_{n=1}^{\infty}(-1)^{n-1}b_n$ 收敛。若其和为 $S$，前 $n$ 项和为 $S_n$，则
\begin{equation}
    0\leq(-1)^n(S-S_n)\leq b_{n+1},
    \qquad |S-S_n|\leq b_{n+1}.
    \label{eq:leibniz-series-remainder}
\end{equation}
这称为\emph{莱布尼茨判别法}（Leibniz's alternating series test）。
\end{theorem}

\begin{proof}
有
\[
    S_{2n}=\sum_{k=1}^n(b_{2k-1}-b_{2k})\geq0,
\]
并且 $S_{2n}$ 随 $n$ 单调递增。另一方面，
\[
    S_{2n}=b_1-\sum_{k=1}^{n-1}(b_{2k}-b_{2k+1})-b_{2n}\leq b_1.
\]
故 $S_{2n}$ 收敛；又 $S_{2n+1}-S_{2n}=b_{2n+1}\to0$，所以奇偶部分和有相同极限 $S$。奇数部分和单调递减，偶数部分和单调递增，因而 $S$ 位于任意相邻的两个部分和之间。$S_{n+1}-S_n=(-1)^nb_{n+1}$，即得式~\eqref{eq:leibniz-series-remainder}。
\end{proof}

单调性只须从某一项起成立，此时余项估计对相应靠后的截断仍成立。仅有符号交错和 $b_n\to0$ 不能保证收敛，例如令 $b_{2k-1}=1/\sqrt{k}$、$b_{2k}=1/(2\sqrt{k})$，则 $b_n\to0$，但偶数部分和为 $\frac12\sum_{k=1}^n k^{-1/2}\to+\infty$。同时，莱布尼茨判别法只先给出收敛，是否绝对收敛还需另行判断 $\sum b_n$。

\begin{example}
\label{ex:alternating-logarithmic-p-series}
讨论 $\sum_{n=2}^{\infty}(-1)^n/(n^p(\ln n)^q)$ 的收敛性与绝对收敛性，其中 $p,q\in\mathbb R$。
\end{example}

\begin{solve}
记 $b(x)=x^{-p}(\ln x)^{-q}$，则
\[
    \frac{b'(x)}{b(x)}=-\frac1x\left(p+\frac q{\ln x}\right).
\]
当 $p>0$，或 $p=0$ 且 $q>0$ 时，$b(n)\to0$ 且最终递减，故交错级数收敛。其余情形通项绝对值不趋于零，级数发散。由例~\ref{ex:logarithmic-p-series}，其中 $p>1$，或 $p=1$ 且 $q>1$ 时绝对收敛；其余已经判为收敛的情形均为条件收敛。特别地，交错调和级数条件收敛，而 $\sum(-1)^n/n^2$ 绝对收敛。
\end{solve}

\subsection{狄利克雷判别法和阿贝尔判别法}
\label{subsec:dirichlet-and-abel-tests}

交错符号的作用在于其部分和有界。更一般地，如果一列数的部分和始终受控，再乘上单调趋零的权重，其尾段和也能受控。实现这一想法的工具是离散的分部积分公式。

\begin{lemma}[阿贝尔变换]
\label{lem:abel-summation-by-parts}
设 $A_0=0$、$A_k=\sum_{j=1}^k a_j$，则对整数 $1\leq m\leq n$，
\begin{equation}
    \sum_{k=m}^n a_kb_k
    =A_nb_n-A_{m-1}b_m
      +\sum_{k=m}^{n-1}A_k(b_k-b_{k+1}).
    \label{eq:abel-summation-by-parts}
\end{equation}
该恒等式称为\emph{阿贝尔变换}（Abel transformation），也称\emph{分部求和}（summation by parts）。
\end{lemma}

\begin{proof}
代入 $a_k=A_k-A_{k-1}$，分别把两个有限和的下标平移，再合并同一 $A_k$ 的系数，即得等式。整个过程仅涉及有限和，不需要任何收敛假设。
\end{proof}

\begin{theorem}[狄利克雷判别法]
\label{thm:dirichlet-test-for-series}
若 $\sum a_n$ 的部分和序列有界，$\{b_n\}$ 单调且 $b_n\to0$，则 $\sum a_nb_n$ 收敛。这称为\emph{狄利克雷判别法}（Dirichlet's test）。
\end{theorem}

\begin{proof}
设 $|A_n|\leq M$。对 $m\leq k\leq n$，令 $C_k=\sum_{j=m}^k a_j=A_k-A_{m-1}$，则 $|C_k|\leq2M$。在式~\eqref{eq:abel-summation-by-parts} 中从 $m$ 重新累加，得到
\[
    \sum_{k=m}^n a_kb_k=C_nb_n+\sum_{k=m}^{n-1}C_k(b_k-b_{k+1}).
\]
单调趋零的实数列保持同一弱符号，且
$\sum_{k=m}^{n-1}|b_k-b_{k+1}|=|b_m|-|b_n|$，因此
\[
    \left|\sum_{k=m}^n a_kb_k\right|
      \leq2M\left(|b_n|+\sum_{k=m}^{n-1}|b_k-b_{k+1}|\right)
      =2M|b_m|\longrightarrow0.
\]
这一估计对任意 $n\geq m$ 一致成立，故由柯西准则，级数收敛。
\end{proof}

\begin{theorem}[阿贝尔判别法]
\label{thm:abel-test-for-series}
若 $\sum a_n$ 收敛，$\{b_n\}$ 单调有界，则 $\sum a_nb_n$ 收敛。这称为\emph{阿贝尔判别法}（Abel's test）。
\end{theorem}

\begin{proof}
单调有界保证 $b_n\to b\in\mathbb R$，故 $b_n-b$ 单调趋零。收敛级数 $\sum a_n$ 的部分和有界，由狄利克雷判别法，$\sum a_n(b_n-b)$ 收敛；而 $\sum ba_n$ 也收敛。二者相加，便得 $\sum a_nb_n$ 收敛。
\end{proof}

两种判别法的假设不能混用：狄利克雷判别法对第一因子只要求部分和有界，但第二因子必须趋零；阿贝尔判别法允许第二因子趋于非零常数，却要求第一因子本身构成收敛级数。取 $a_n=(-1)^{n-1}$，其部分和有界，便可由狄利克雷判别法重新得到莱布尼茨判别法。

\begin{example}
\label{ex:dirichlet-trigonometric-series}
设 $a_n\geq0$ 单调趋零，讨论 $\sum a_n\sin(nx)$ 与 $\sum a_n\cos(nx)$。若另有 $\sum a_n=+\infty$，说明何时为条件收敛。
\end{example}

\begin{solve}
当 $x\notin2\pi\mathbb Z$ 时，由三角恒等式的有限求和，
\begin{align*}
    \sum_{k=1}^n\sin(kx)
       &=\frac{\cos(x/2)-\cos((n+1/2)x)}{2\sin(x/2)},\\
    \sum_{k=1}^n\cos(kx)
       &=\frac{\sin((n+1/2)x)-\sin(x/2)}{2\sin(x/2)}.
\end{align*}
对固定的 $x$，两个部分和均有与 $n$ 无关的上界，故由狄利克雷判别法，两级数都收敛。$x\in2\pi\mathbb Z$ 时，正弦级数恒为零，余弦级数等于 $\sum a_n$。若 $\sum a_n$ 收敛，则 $|a_n\sin(nx)|,|a_n\cos(nx)|\leq a_n$，所以两级数对一切 $x$ 均绝对收敛。

以下设 $\sum a_n=+\infty$。当 $x\notin\pi\mathbb Z$ 时，$2x\notin2\pi\mathbb Z$，故 $\sum a_n\cos(2nx)$ 收敛。利用 $|\sin t|\geq\sin^2t$、$|\cos t|\geq\cos^2t$，在有限和层面有
\begin{align*}
    \sum_{k=1}^n a_k|\sin(kx)|
       &\geq\frac12\sum_{k=1}^n a_k
             -\frac12\sum_{k=1}^n a_k\cos(2kx),\\
    \sum_{k=1}^n a_k|\cos(kx)|
       &\geq\frac12\sum_{k=1}^n a_k
             +\frac12\sum_{k=1}^n a_k\cos(2kx).
\end{align*}
两个右端均趋于 $+\infty$，因此两级数均条件收敛。$x\in\pi\mathbb Z$ 时，正弦级数恒为零，绝对收敛；若 $x$ 为 $\pi$ 的奇数倍，余弦级数为 $\sum(-1)^na_n$，收敛但不绝对收敛；若为 $2\pi$ 的整数倍，余弦级数发散。于是，在 $\sum a_n$ 发散的假设下，正弦级数恰在 $x\notin\pi\mathbb Z$ 时条件收敛，余弦级数恰在 $x\notin2\pi\mathbb Z$ 时条件收敛。
\end{solve}

\section{数项级数的性质}
\label{sec:properties-of-numerical-series}

有限和的结合律、交换律及分配律可以自由使用，但级数还包含取极限的步骤。分组会选取部分和子列，重排会改变累积抵消的次序，乘法则要规定怎样累加所有两两乘积。因此，这些运算能否保持收敛及其和，必须分别讨论。

\subsection{结合律}
\label{subsec:associativity-of-series}

\begin{theorem}[收敛级数的分组]
\label{thm:grouping-convergent-series}
设 $\sum a_n=S$，任取整数 $0=n_0<n_1<n_2<\cdots$，令
\[
    b_k=a_{n_{k-1}+1}+\cdots+a_{n_k}.
\]
则 $\sum b_k$ 收敛且和为 $S$。也就是说，保持原次序，把相邻的有限项依次分组，不改变收敛级数的和。
\end{theorem}

\begin{proof}
分组级数的第 $k$ 个部分和为 $\sum_{j=1}^k b_j=S_{n_k}$，是原部分和序列的子列，故同样趋于 $S$。
\end{proof}

逆命题一般不成立：$1-1+1-1+\cdots$ 发散，分为 $(1-1)+(1-1)+\cdots$ 却得到和为零的级数。这也说明，不能从一个已经收敛的分组级数出发任意拆开括号，除非另有控制组内部分和的条件。定理~\ref{thm:bounded-block-grouping-series} 给出了一种这样的条件。

对正项级数，分组收敛可以推出原级数收敛，因为原部分和被相邻的分组部分和夹住。更一般地，若每一组内各项同为非负或同为非正，则组内部分和单调，仍有
\[
    \min\{S_{n_{k-1}},S_{n_k}\}\leq S_n
       \leq\max\{S_{n_{k-1}},S_{n_k}\},
    \qquad n_{k-1}\leq n\leq n_k.
\]
若分组部分和趋于 $S$，夹逼即得原部分和也趋于 $S$。这个结论依赖组内同号，不要求组长有统一上界。

\subsection{交换律}
\label{subsec:commutativity-of-series}

\begin{definition}
若 $\pi:\mathbb N\to\mathbb N$ 为双射，则称 $\sum_{n=1}^{\infty}a_{\pi(n)}$ 为 $\sum a_n$ 的一个\emph{重排}（rearrangement）。双射保证每个原下标恰好出现一次，既不遗漏也不重复。
\end{definition}

\begin{theorem}[有界位移重排]
\label{thm:bounded-displacement-rearrangement}
若存在正整数 $M$ 使 $|\pi(n)-n|\leq M$ 对所有 $n$ 成立，则 $\sum a_n$ 与 $\sum a_{\pi(n)}$ 同敛散，并且收敛时和相同。
\end{theorem}

\begin{proof}
先设原级数收敛，则 $a_n\to0$。记两级数部分和分别为 $S_n,T_n$。$\{\pi(1),\ldots,\pi(n)\}$ 包含所有不超过 $n-M$ 的正整数，且不含大于 $n+M$ 的整数，故 $n>M$ 时
\[
    |T_n-S_n|\leq\sum_{k=n-M+1}^{n+M}|a_k|\longrightarrow0.
\]
右侧只有固定的 $2M$ 项，故两者有相同极限。反之，逆双射也满足 $|\pi^{-1}(n)-n|\leq M$，对重排后的级数应用刚证明的方向即可。
\end{proof}

\begin{theorem}[绝对收敛级数的重排]
\label{thm:absolutely-convergent-series-rearrangement}
若 $\sum a_n$ 绝对收敛，则任意重排 $\sum a_{\pi(n)}$ 也绝对收敛，并且
\[
    \sum_{n=1}^{\infty}a_{\pi(n)}=\sum_{n=1}^{\infty}a_n,
    \qquad
    \sum_{n=1}^{\infty}|a_{\pi(n)}|=\sum_{n=1}^{\infty}|a_n|.
\]
\end{theorem}

\begin{proof}
先考察非负项。$\sum_{k=1}^n|a_{\pi(k)}|\leq\sum_{k=1}^{\infty}|a_k|$，故重排后的绝对值级数收敛，且其和不超过原和；再用逆双射，得到反向不等式，从而绝对值级数的和相同。

记 $\sum a_n=S$。给定 $\varepsilon>0$，取 $K$ 使 $\sum_{k>K}|a_k|<\varepsilon$；由双射性，存在 $N$，使 $n\geq N$ 时重排的前 $n$ 项已经包含原级数的前 $K$ 项。于是
\[
    |T_n-S|
    \leq|T_n-S_K|+|S_K-S|
    \leq2\sum_{k>K}|a_k|<2\varepsilon.
\]
故 $T_n\to S$，任意重排不改变原和。
\end{proof}

\begin{example}
\label{ex:alternating-harmonic-rearrangement}
证明交错调和级数的和为 $\ln2$，而把它重排为每两个正项后接一个负项时，和变为 $\frac32\ln2$。
\end{example}

\begin{solve}
令 $H_n=\sum_{k=1}^n1/k$。对固定正整数 $m$，
\[
    H_{mn}-H_n=\frac1n\sum_{k=n+1}^{mn}\frac1{k/n}
       \longrightarrow\int_1^m\frac{dx}{x}=\ln m,
\]
这是连续函数 $1/x$ 在 $[1,m]$ 上的黎曼和。原级数的偶数部分和为 $H_{2n}-H_n\to\ln2$，奇数部分和与相邻偶数部分和之差趋零，因此原和为 $\ln2$。

重排后的级数为
\[
    1+\frac13-\frac12+\frac15+\frac17-\frac14+\cdots.
\]
它的前 $3n$ 项之和为
\[
    T_{3n}=H_{4n}-\frac12H_{2n}-\frac12H_n
          =(H_{4n}-H_{2n})+\frac12(H_{2n}-H_n)
          \longrightarrow\frac32\ln2.
\]
$T_{3n+1}-T_{3n}=1/(4n+1)$，$T_{3n+2}-T_{3n+1}=1/(4n+3)$，均趋于零，故整个重排级数收敛到这一新值。该重排每项都出现一次；改变和的原因是条件收敛，不能据此否认它是合法的重排。
\end{solve}

\begin{theorem}[黎曼重排定理]
\label{thm:riemann-rearrangement}
若实数项级数 $\sum a_n$ 条件收敛，则对任意 $s\in\mathbb R$，存在一个重排收敛到 $s$；也分别存在重排使部分和趋于 $+\infty$ 或 $-\infty$。这称为\emph{黎曼重排定理}（Riemann rearrangement theorem）。
\end{theorem}

\begin{proof}
定义正部与负部的大小为
\[
    a_n^+=\max\{a_n,0\},\qquad a_n^-=\max\{-a_n,0\}.
\]
于是 $a_n=a_n^+-a_n^-$、$|a_n|=a_n^++a_n^-$。条件收敛蕴含
\begin{equation}
    \sum a_n^+=+\infty,\qquad \sum a_n^-=+\infty,
    \qquad a_n\to0.
    \label{eq:conditional-series-positive-negative-parts}
\end{equation}
事实上，若两个非负级数都收敛，则原级数绝对收敛；若仅一个收敛，另一个发散，则原部分和趋于一个带符号的无穷，与收敛矛盾。

把正项按原次序列为 $p_1,p_2,\ldots$，把负项按原次序列为 $-q_1,-q_2,\ldots$，其中 $p_j,q_j>0$ 且均趋零。给定 $s\in\mathbb R$，先依次取正项，直到部分和首次大于 $s$；若初始值 $0$ 已大于 $s$，这一首次正项阶段可为空。再依次取负项，直到首次小于 $s$；此后正负阶段交替进行。每个阶段都能在有限步结束，因为两类项的总量都为 $+\infty$。除初始阶段外，每次向上越过 $s$ 的超出量不超过最后加入的正项，每次向下越过的不足量不超过最后加入的负项，故阶段端点的部分和都趋于 $s$。阶段内部分和单调地位于相邻端点之间，也趋于 $s$，因此整个重排收敛到 $s$。每类项都按原次序不断取用，故每一非零项最终都会出现；若有零项，按原次序在各阶段间逐个插入即可。

若要求趋于 $+\infty$，第 $k$ 阶段先至少取一个尚未使用的正项，并继续取正项直到部分和大于 $k+q_k$，再加入负项 $-q_k$。阶段末部分和大于 $k$，阶段内的正项累加不低于上一阶段末值，最后加入负项后也仍大于 $k$，故全部部分和趋于 $+\infty$。正负两类项都被穷尽，零项仍可依次插入。对 $-a_n$ 作同样构造，便得到趋于 $-\infty$ 的重排。
\end{proof}

因此，对已经收敛的实数项级数，绝对收敛保证所有重排同和；若只是条件收敛，则任意重排同和的结论不能成立。使用交换律前，应先核对绝对收敛，或核对具体重排是否满足定理~\ref{thm:bounded-displacement-rearrangement} 等额外条件。

\subsection{级数的乘法（分配律）}
\label{subsec:multiplication-and-distributivity-of-series}

设 $\sum a_n$、$\sum b_n$ 为两个级数，其所有两两乘积 $a_ib_j$ 构成以 $(i,j)\in\mathbb N^2$ 为下标的乘积矩阵。有限分配律并不直接规定无穷矩阵的累加顺序，须先选定排序或分组方式。

\begin{definition}
令
\begin{equation}
    c_n=\sum_{k=1}^n a_kb_{n+1-k},\qquad n\geq1,
    \label{eq:cauchy-product-coefficients}
\end{equation}
则级数 $\sum_{n=1}^{\infty}c_n$ 称为两个级数的\emph{柯西乘积}（Cauchy product）。它按乘积矩阵的斜对角线求和，把下标满足 $i+j=n+1$ 的各项归为一组。
\end{definition}

另一种方式是按扩大的正方形边界分组。记 $A_n=\sum_{k=1}^n a_k$、$B_n=\sum_{k=1}^n b_k$，并令 $A_0=B_0=0$，则第 $n$ 组为
\[
    d_n=A_nb_n+a_nB_{n-1}=A_nB_n-A_{n-1}B_{n-1}.
\]
所以 $\sum_{k=1}^n d_k=A_nB_n$。只要两个原级数分别收敛到 $A,B$，这一\emph{按正方形分组的级数}就收敛到 $AB$；但分组的组长为 $2n-1$，并无统一上界，故不能仅凭此式推出拆开每组之后的级数也收敛，更不能推出柯西乘积收敛。

\begin{theorem}[绝对收敛级数的乘法]
\label{thm:absolutely-convergent-series-product}
若 $\sum a_n=A$、$\sum b_n=B$ 都绝对收敛，则把所有 $a_ib_j$ 按任意次序各取一次所成的级数均绝对收敛，且和为 $AB$。特别地，柯西乘积也绝对收敛，且
\begin{equation}
    \sum_{n=1}^{\infty}\sum_{k=1}^n a_kb_{n+1-k}
      =\left(\sum_{n=1}^{\infty}a_n\right)
       \left(\sum_{n=1}^{\infty}b_n\right).
    \label{eq:absolute-series-cauchy-product}
\end{equation}
\end{theorem}

\begin{proof}
记 $U=\sum|a_n|$、$V=\sum|b_n|$。任意有限个乘积项的绝对值之和都不超过一个包含它们的有限矩形的绝对值和，而后者不超过 $UV$。故任一枚举对应的绝对值级数部分和有界，级数绝对收敛。由定理~\ref{thm:absolutely-convergent-series-rearrangement}，所有枚举的和相同。

选定一种逐层列出正方形边界元素的枚举，它在列完 $n\times n$ 正方形时的部分和为 $A_nB_n\to AB$。由于这一枚举已经证明收敛，其和就等于该子列的极限 $AB$，于是所有枚举均同和。最后，把斜对角线枚举分组即得柯西乘积；又有
\[
    \sum_{n=1}^N|c_n|
    \leq\sum_{\substack{i,j\geq1\\i+j\leq N+1}}|a_i b_j|
    \leq UV,
\]
故它也绝对收敛，且由分组定理保持和 $AB$。
\end{proof}

\begin{example}
\label{ex:conditional-series-cauchy-square-divergence}
设 $0<\alpha\leq1/2$，$a_n=(-1)^{n-1}/n^\alpha$。证明 $\sum a_n$ 收敛，但它与自身的柯西乘积发散。
\end{example}

\begin{solve}
由莱布尼茨判别法，原级数收敛；因 $\sum n^{-\alpha}$ 发散，它为条件收敛。自乘的柯西乘积通项为
\[
    c_n=(-1)^{n-1}\sum_{k=1}^n
        \frac1{[k(n+1-k)]^\alpha}.
\]
由 $k(n+1-k)\leq(n+1)^2/4$，得到
\[
    |c_n|\geq\frac{4^\alpha n}{(n+1)^{2\alpha}}.
\]
$0<\alpha<1/2$ 时该下界趋于 $+\infty$；$\alpha=1/2$ 时趋于 $2$。所以 $c_n$ 不趋于零，柯西乘积发散。相同两个原级数的正方形分组乘积却仍收敛到原和的平方，说明排序与分组方式在缺少绝对收敛条件时确实会影响结论。
\end{solve}

\section{无穷乘积}
\label{sec:infinite-products}

无穷乘积同样由有限运算的极限定义，但收敛约定与级数不同。如果允许部分积趋于零也算收敛，那么单个零因子就能使后续全部部分积变为零，无法反映其余因子的累积性质。因此，本节按照教材约定，只把具有有限非零极限的部分积称为收敛。

\begin{definition}
设 $\{v_n\}_{n\geq1}$ 为实数列，称 $\prod_{n=1}^{\infty}v_n$ 为\emph{无穷乘积}（infinite product），并称
\[
    P_n=\prod_{k=1}^n v_k
\]
为第 $n$ 个\emph{部分积}（partial product）。若 $P_n\to P\in\mathbb R\setminus\{0\}$，则称该无穷乘积\emph{收敛}，并记 $\prod_{n=1}^{\infty}v_n=P$；否则称为发散。特别地，部分积趋于零也属于这里的发散情形。
\end{definition}

\begin{proposition}
\label{prop:infinite-product-necessary-condition}
无穷乘积 $\prod v_n$ 收敛的必要条件是每个因子非零，并且 $v_n\to1$。所以收敛乘积的因子从某项起均为正数；对因子全非零的乘积，改变有限个因子而仍保持这些因子非零，不改变敛散性。
\end{proposition}

\begin{proof}
若有零因子，则充分靠后的 $P_n$ 均为零，不满足定义。若 $P_n\to P\neq0$，则
\[
    v_n=\frac{P_n}{P_{n-1}}\longrightarrow1.
\]
因此充分大的 $n$ 满足 $v_n>0$。有限个非零因子的替换，只使充分靠后的部分积乘以一个固定非零常数，故不改变是否具有有限非零极限。
\end{proof}

$v_n\to1$ 不是充分条件。例如
$\prod_{k=1}^n(1+1/k)=n+1\to+\infty$；而 $\prod_{k=2}^n(1-1/k)=1/n\to0$，同样发散。对最终为正的因子，可以把有限个初始非零因子提出，只研究正因子尾积；这样得到的尾积若收敛，其极限必为正数，而整个乘积的符号由初始因子决定。

\begin{theorem}[无穷乘积的对数判据]
\label{thm:infinite-product-logarithmic-criterion}
设 $v_n>0$，则 $\prod v_n$ 收敛的充分必要条件是实数项级数 $\sum\ln v_n$ 收敛；收敛时
\begin{equation}
    \prod_{n=1}^{\infty}v_n
      =\exp\left(\sum_{n=1}^{\infty}\ln v_n\right).
    \label{eq:infinite-product-logarithm}
\end{equation}
特别地，对 $u_n>-1$，可取 $v_n=1+u_n$ 来研究 $\prod(1+u_n)$。
\end{theorem}

\begin{proof}
有限部分积与对数部分和满足
\[
    P_n=\exp\left(\sum_{k=1}^n\ln v_k\right),\qquad
    \sum_{k=1}^n\ln v_k=\ln P_n.
\]
对数部分和趋于有限实数 $L$ 时，由指数函数连续性，$P_n\to e^L>0$。反之，若乘积收敛，则正部分积的极限 $P$ 非零，故 $P>0$，由对数函数连续性得到对数部分和趋于 $\ln P$。
\end{proof}

若对数部分和趋于 $-\infty$，部分积趋于零；若趋于 $+\infty$，部分积趋于 $+\infty$。这两种情形均不满足乘积收敛的定义。对一般有正有负的 $u_n$，不能仅凭 $\ln(1+u_n)\sim u_n$ 就断言两个级数同敛散，因为任意项级数不适用正项级数的极限比较判别法。

\begin{corollary}
\label{cor:nonnegative-infinite-product-test}
若 $u_n\geq0$，则 $\prod(1+u_n)$ 收敛的充分必要条件是 $\sum u_n$ 收敛。
\end{corollary}

\begin{proof}
若 $\sum u_n$ 收敛，则 $u_n\to0$，利用 $0\leq\ln(1+u_n)\leq u_n$ 可知对数级数收敛，进而乘积收敛。反之，乘积收敛首先蕴含 $u_n\to0$，从而最终有 $0\leq u_n\leq1$，且
\[
    \ln(1+u_n)=\int_0^{u_n}\frac{dt}{1+t}\geq\frac{u_n}{2}.
\]
对数级数收敛，由比较判别法知 $\sum u_n$ 收敛。
\end{proof}

\begin{example}
\label{ex:infinite-product-p-power}
讨论 $\prod_{n=1}^{\infty}(1+n^{-p})$，其中 $p\in\mathbb R$。
\end{example}

\begin{solve}
当 $p>1$ 时，$\sum n^{-p}$ 收敛，由推论~\ref{cor:nonnegative-infinite-product-test}，乘积收敛。$0<p\leq1$ 时对数级数与发散的 $\sum n^{-p}$ 同敛散，故部分积趋于 $+\infty$。$p\leq0$ 时，各因子均不小于 $2$，部分积也趋于 $+\infty$。因此原乘积恰在 $p>1$ 时收敛。
\end{solve}

\begin{example}
\label{ex:wallis-infinite-product}
证明
\begin{equation}
    \prod_{n=1}^{\infty}\left(1-\frac1{4n^2}\right)=\frac2\pi,
    \qquad
    \prod_{n=1}^{\infty}\frac{(2n)^2}{(2n-1)(2n+1)}=\frac\pi2.
    \label{eq:wallis-infinite-product}
\end{equation}
第二个等式通常称为\emph{沃利斯乘积}（Wallis product）。
\end{example}

\begin{solve}
记 $r_n=(2n-1)!!/(2n)!!$，则有限部分积为
\[
    P_n=\prod_{k=1}^n\frac{(2k-1)(2k+1)}{(2k)^2}
       =(2n+1)r_n^2.
\]
为求极限，令 $I_m=\int_0^{\pi/2}\sin^m x\,dx$。对整数 $m\geq2$，分部积分给出 $I_m=(m-1)(I_{m-2}-I_m)$，即
\[
    I_m=\frac{m-1}{m}I_{m-2},\qquad I_0=\frac\pi2,\quad I_1=1.
\]
递推得到 $I_{2n}=(\pi/2)r_n$、$I_{2n+1}=1/((2n+1)r_n)$。又由 $0\leq\sin x\leq1$，
\[
    \frac{2n+1}{2n+2}
       =\frac{I_{2n+2}}{I_{2n}}
       \leq\frac{I_{2n+1}}{I_{2n}}\leq1,
\]
所以 $I_{2n+1}/I_{2n}\to1$，并有
\[
    P_n=\frac2\pi\frac{I_{2n}}{I_{2n+1}}\longrightarrow\frac2\pi.
\]
极限非零，故乘积收敛；取有限部分积的倒数，再取极限，就得到第二个等式。
\end{solve}

\begin{example}
\label{ex:euler-product-zeta}
对实数 $s>1$，定义 $\zeta(s)=\sum_{n=1}^{\infty}n^{-s}$。将素数按递增顺序记为 $p_1,p_2,\ldots$，证明\emph{欧拉乘积公式}（Euler product formula）
\begin{equation}
    \zeta(s)=\prod_{j=1}^{\infty}(1-p_j^{-s})^{-1},
    \qquad
    \prod_{j=1}^{\infty}(1-p_j^{-s})=\frac1{\zeta(s)}.
    \label{eq:euler-product-zeta}
\end{equation}
\end{example}

\begin{solve}
$s>1$ 保证 $p$ 级数 $\zeta(s)$ 有限。对前 $N$ 个素数，分别使用几何级数，并利用正整数的素因数分解唯一性，得
\[
    Q_N=\prod_{j=1}^N(1-p_j^{-s})^{-1}
       =\sum_{\substack{m\geq1\\m\text{ 的素因子均属于 }\{p_1,\ldots,p_N\}}}\frac1{m^s},
\]
其中整数 $1$ 也包括在右侧。这个展开可先把每个几何级数截为有限和，使用有限分配律，再让各截断长度趋于无穷；所有项非负，故有限和单调增加，不涉及未经证明的条件收敛级数重排。唯一分解保证每个符合条件的整数恰好出现一次。

右侧包括所有 $m\leq p_N$，而未出现的整数都大于 $p_N$，因此
\[
    0\leq\zeta(s)-Q_N\leq\sum_{m>p_N}\frac1{m^s}\longrightarrow0.
\]
故第一项乘积收敛到 $\zeta(s)>0$；取倒数得到第二项乘积的非零极限。这说明无穷乘积可以把整数的乘法分解结构与级数联系起来。

这个关系也能反证素数只有有限个的假设。若全部素数仅为 $p_1,\ldots,p_N$，相同的有限素数展开对每个 $s>1$ 都给出
\[
    \zeta(s)=\prod_{j=1}^N(1-p_j^{-s})^{-1}.
\]
当 $s\downarrow1$ 时，右侧趋于有限正数 $\prod_{j=1}^N(1-p_j^{-1})^{-1}$，左侧却无界。事实上，任意给定 $M>0$，可先取整数 $K$ 使 $\sum_{n=1}^K1/n>2M$，再由这个有限和关于 $s$ 的连续性，令 $s>1$ 充分接近 $1$，使 $\zeta(s)\geq\sum_{n=1}^K n^{-s}>M$。矛盾说明素数不可能只有有限个。这里先固定有限的 $K$ 再取 $s\downarrow1$，没有把极限直接交换进一个尚未核实的无穷和。
\end{solve}

\section*{习题九}
\addcontentsline{toc}{section}{习题九}
\fi

\selectchapter{10}
\ifcompilechapter
\chapter{函数序列与函数项级数}
\label{chap:function-sequences-and-series}

% 本章正文依据数学分析-10.pdf 第 1--36 页及第 37 页习题前的结束语，
% 对应教材第 171--207 页习题十之前的正文；习题十不在本次整理范围内。
用多项式、三角函数等简单函数逼近较复杂的函数，会自然地得到函数序列及函数项级数。这里不仅要对每个自变量求一个数列极限，还要研究所得函数能否继承连续性、可积性和可导性。困难在于：对每个固定点都收敛，并不能控制整个区间上的逼近误差。本章以\emph{一致收敛}为主要工具，依次讨论收敛的概念与判别、极限与其他运算的交换，并说明如何用级数构造具有指定性质的函数。以下主要讨论定义在非空集合上的实值函数；积分均指黎曼积分，除另有说明外，区间端点的连续性、导数和极限均按相应单侧理解。

\section{函数序列与函数项级数的基本问题}
\label{sec:basic-problems-of-function-sequences-and-series}

\begin{definition}
\label{def:pointwise-function-sequence-convergence}
设 $E_0\subseteq\mathbb R$ 非空，对每个 $n\in\mathbb N$，$f_n:E_0\to\mathbb R$。按正整数排列的函数 $f_1,f_2,\ldots$ 称为\emph{函数序列}（sequence of functions）。若数列 $\{f_n(x)\}$ 在 $x\in E_0$ 处收敛，则称 $x$ 为该函数序列的收敛点；全部收敛点组成的集合 $E$ 称为\emph{收敛域}（domain of convergence）。在 $E$ 上定义
\[
    f(x)=\lim_{n\to\infty}f_n(x),
\]
称 $f$ 为\emph{极限函数}（limit function），并称 $f_n$ 在 $E$ 上\emph{逐点收敛}（pointwise convergence）于 $f$。
\end{definition}

\begin{definition}
\label{def:function-series-and-sum-function}
设 $u_n:E_0\to\mathbb R$，形式和 $\sum_{n=1}^\infty u_n(x)$ 称为\emph{函数项级数}（series of functions）。记
\[
    S_N(x)=\sum_{n=1}^N u_n(x).
\]
若数列 $\{S_N(x)\}$ 收敛，则称级数在 $x$ 处收敛；全部收敛点组成其收敛域。在收敛域上，
\[
    S(x)=\lim_{N\to\infty}S_N(x)=\sum_{n=1}^\infty u_n(x)
\]
称为\emph{和函数}（sum function）。若 $\sum_{n=1}^\infty|u_n(x)|$ 收敛，则称原级数在该点绝对收敛。
\end{definition}

固定 $x$ 后，这些问题分别退化为数列与数项级数问题，因此可先用数列极限与数项级数的方法求收敛域与极限函数。函数项级数的收敛等价于其部分和序列的收敛；反过来，约定 $f_0=0$ 后，恒有
\[
    f_N(x)=\sum_{n=1}^N\bigl(f_n(x)-f_{n-1}(x)\bigr),
\]
故函数序列也可写成一个级数的部分和序列。但从逐点收敛进一步推断函数性质时，需要新的条件。

\begin{example}
\label{ex:pointwise-limits-and-operation-failures}
考察 $f_n(x)=x^n$（$x\in\mathbb R$）、$g_n(x)=nx(1-x^2)^n$（$0\leq x\leq1$）、$h_n(x)=\sin(nx)/\sqrt n$（$x\in\mathbb R$）及 $p_n(x)=\frac{2x}{\pi}\arctan(nx)$（$x\in\mathbb R$）的极限及其连续性、积分或导数，说明逐点收敛本身不能保证这些运算与极限交换。
\end{example}

\begin{solve}
首先，$f_n(x)=x^n$ 在实轴上的收敛域为 $(-1,1]$，且
\[
    f(x)=\begin{cases}0,&-1<x<1,\\1,&x=1.\end{cases}
\]
每个 $f_n$ 都连续，但极限函数在 $1$ 处不左连续。特别地，先令 $n\to\infty$ 再令 $x\to1-$，与按相反次序取极限，分别得到 $0$ 与 $1$。

其次，在 $[0,1]$ 上令 $g_n(x)=nx(1-x^2)^n$。端点处恒为零；固定 $0<x<1$ 后，几何衰减快于 $n$ 的增长，故 $g_n(x)\to0$。另一方面，换元 $t=1-x^2$ 得
\[
    \int_0^1g_n(x)\,dx=\frac{n}{2(n+1)}\longrightarrow\frac12,
    \qquad \int_0^1\lim_{n\to\infty}g_n(x)\,dx=0.
\]
这里极限函数甚至是连续函数，但积分仍不能与极限交换；靠近零点的狭窄区域保留了不能忽略的积分贡献。

再次，令 $h_n(x)=\sin(nx)/\sqrt n$，则实轴上恒有 $|h_n(x)|\leq1/\sqrt n$，故极限函数为零；但 $h_n'(0)=\sqrt n\to+\infty$，因而极限的导数不等于导数的极限。最后，令
\[
    p_n(x)=\frac{2x}{\pi}\arctan(nx),\qquad
    p_n'(x)=\frac2\pi\arctan(nx)+\frac{2nx}{\pi(1+n^2x^2)}.
\]
有 $p_n(x)\to|x|$，且导函数在每一点都有极限：$x>0$ 时为 $1$，$x<0$ 时为 $-1$，$x=0$ 时为 $0$。然而 $|x|$ 在零点不可导。这说明即使导函数序列也逐点收敛，仍不足以保证极限函数可导。
\end{solve}

上述例子提出三个需要分别回答的问题：连续函数的极限是否连续，可积函数的极限是否可积并允许交换积分与极限，可导函数的极限是否可导并允许交换微分与极限。第三个问题所需的控制比前两个更强，不能把同一条收敛条件机械地用于所有运算。

\section{一致收敛的概念}
\label{sec:concept-of-uniform-convergence}

逐点收敛的量词次序是
\[
    \forall x\in E\ \forall\varepsilon>0\ \exists N=N(\varepsilon,x)
    \ \forall n\geq N:\quad |f_n(x)-f(x)|<\varepsilon.
\]
不同点可以需要不同的项数。若希望一个足够靠后的函数在整个集合上同时逼近 $f$，就必须使这里的 $N$ 与 $x$ 无关。

\begin{definition}
\label{def:uniform-function-sequence-and-series-convergence}
设 $f_n,f:E\to\mathbb R$，$E\neq\varnothing$。若
\begin{equation}
    \forall\varepsilon>0\ \exists N=N(\varepsilon)\ \forall n\geq N
    \ \forall x\in E:\quad |f_n(x)-f(x)|<\varepsilon,
    \label{eq:uniform-function-sequence-convergence}
\end{equation}
则称 $f_n$ 在 $E$ 上\emph{一致收敛}（uniform convergence）于 $f$，记为 $f_n\rightrightarrows f$。若级数 $\sum_{n=1}^\infty u_n$ 的部分和序列在 $E$ 上一致收敛于 $S$，则称该级数在 $E$ 上一致收敛。写 $R_N=S-S_N$ 时，这等价于余项 $R_N\rightrightarrows0$。
\end{definition}

一致收敛必逐点收敛，因而极限函数唯一。几何上，对给定的 $\varepsilon$，充分靠后的每条曲线 $y=f_n(x)$ 在全部 $x\in E$ 处，都落在 $f(x)-\varepsilon<y<f(x)+\varepsilon$ 的带形区域中。这里不要求 $E$ 闭、有界，也不要求函数连续。与式~\eqref{eq:uniform-function-sequence-convergence} 对应，\emph[caution]{不一致收敛}的准确表述是
\begin{equation}
    \exists\varepsilon_0>0\ \forall N\in\mathbb N\ \exists n\geq N
    \ \exists x\in E:\quad |f_n(x)-f(x)|\geq\varepsilon_0.
    \label{eq:negation-of-uniform-function-convergence}
\end{equation}
因此，只要能找到 $n_j\to\infty$ 和随之移动的点 $x_j\in E$，使误差始终不小于一个固定正数，就能否定一致收敛；这里并不要求全部 $x_j$ 是同一个点。

\begin{example}
\label{ex:uniform-powers-and-endpoint-factor}
讨论 $x^n$ 在不同区间上的一致收敛性，并证明：若 $g\in C([0,1])$ 且 $g(1)=0$，则 $x^ng(x)\rightrightarrows0$ 于 $[0,1]$。
\end{example}

\begin{solve}
对 $0<r<1$，在 $[-r,r]$ 或 $(-r,r)$ 上均有 $|x^n|\leq r^n\to0$，故一致收敛于零。在 $(-1,1)$ 或 $[0,1)$ 上，取 $x_n=2^{-1/n}\to1$，则 $x_n^n=1/2$，故不一致收敛于零。同样的点列也说明在 $[0,1]$ 上不一致收敛于其分段定义的极限函数。可见极限函数在一个集合上连续，也不意味着原序列一定在该集合上一致收敛。

对于第二个结论，记 $M=\max_{[0,1]}|g|$。给定 $\varepsilon>0$，由 $g$ 在 $1$ 处连续且 $g(1)=0$，可取 $0<\delta<1$，使 $1-\delta<x\leq1$ 时 $|g(x)|<\varepsilon$。在这一小段上，全部 $n$ 都满足 $|x^ng(x)|<\varepsilon$；在其余部分 $0\leq x\leq1-\delta$，有
\[
    |x^ng(x)|\leq M(1-\delta)^n\longrightarrow0.
\]
取与 $x$ 无关的充分大 $N$，便可同时控制两部分。这个证明体现了常用方法：靠近端点时利用系数函数变小，离开端点后利用几何因子的统一衰减。
\end{solve}

\begin{example}
\label{ex:uniform-telescoping-rational-series}
讨论 $x>0$ 时级数
\[
    \sum_{n=1}^\infty
    \frac{(2n+1)x}{(1+n^2x)(1+(n+1)^2x)}
\]
在 $[a,+\infty)$（$a>0$）及 $(0,+\infty)$ 上的一致收敛性。
\end{example}

\begin{solve}
通项可写为 $1/(1+n^2x)-1/(1+(n+1)^2x)$，故
\[
    S_N(x)=\frac1{1+x}-\frac1{1+(N+1)^2x},\qquad
    S(x)=\frac1{1+x}.
\]
在 $[a,+\infty)$ 上，余项绝对值不超过 $1/(1+(N+1)^2a)\to0$，所以一致收敛。在 $(0,+\infty)$ 上取 $x_N=(N+1)^{-2}$，则余项恒为 $1/2$，所以不一致收敛。这也说明无界区间上可以一致收敛，而有问题的区域可能是一个未包含的有限端点附近。
\end{solve}

\begin{proposition}
\label{prop:elementary-uniform-convergence-properties}
一致收敛具有以下性质：在 $E$ 上一致收敛的序列，在其任意非空子集上也一致收敛；若同一序列在 $E_1,E_2$ 上分别一致收敛于同一个函数的限制，则在 $E_1\cup E_2$ 上也一致收敛。若 $f_n\rightrightarrows f$、$g_n\rightrightarrows g$ 于 $E$，则对固定实数 $\alpha,\beta$，有 $\alpha f_n+\beta g_n\rightrightarrows\alpha f+\beta g$。若还存在常数 $M_1,M_2$，使全部 $n$、$x\in E$ 满足 $|f_n(x)|\leq M_1$、$|g_n(x)|\leq M_2$，则 $f_ng_n\rightrightarrows fg$。
\end{proposition}

\begin{proof}
限制到子集时沿用原来的 $N$；有限并集取两组 $N$ 的最大值。线性组合使用
\[
    |\alpha f_n+\beta g_n-\alpha f-\beta g|
    \leq|\alpha|\,|f_n-f|+|\beta|\,|g_n-g|.
\]
最后一项中，取逐点极限知 $|f|\leq M_1$、$|g|\leq M_2$，再由
\[
    |f_ng_n-fg|\leq M_1|g_n-g|+M_2|f_n-f|
\]
得到结论。
\end{proof}

上述对全部 $n$、$x$ 同时成立的界称为函数序列的\emph{一致有界性}（uniform boundedness）。它比每一个函数分别有界更强。乘积的结论不能无条件使用：在 $(0,1)$ 上，$(1-x)x^n\rightrightarrows0$，常值序列 $g_n(x)=1/(1-x)$ 一致收敛于自身，但乘积 $x^n$ 不一致收敛于零。这里 $g_n$ 没有关于 $x$ 的统一界。有限并集的性质也不能直接推广到可数并集，例如 $[0,1)=\bigcup_{m\geq2}[0,1-1/m]$，$x^n$ 在每个小区间上一致收敛，在并集上却不一致收敛。

\section{函数序列与函数项级数一致收敛的判别法}
\label{sec:uniform-convergence-tests}

\subsection{柯西准则}
\label{subsec:cauchy-criterion-for-uniform-convergence}

\begin{theorem}[一致收敛的柯西准则]
\label{thm:uniform-convergence-cauchy-criterion}
设 $f_n:E\to\mathbb R$，$E\neq\varnothing$。则 $\{f_n\}$ 在 $E$ 上一致收敛的充分必要条件是
\begin{equation}
    \forall\varepsilon>0\ \exists N\ \forall m,n\geq N
    \ \forall x\in E:\quad |f_n(x)-f_m(x)|<\varepsilon.
    \label{eq:uniform-cauchy-criterion}
\end{equation}
因此，$\sum_{n=1}^\infty u_n$ 在 $E$ 上一致收敛，等价于
\begin{equation}
    \forall\varepsilon>0\ \exists N\ \forall q\geq p\geq N
    \ \forall x\in E:\quad
    \left|\sum_{k=p}^q u_k(x)\right|<\varepsilon.
    \label{eq:uniform-series-cauchy-criterion}
\end{equation}
\end{theorem}

\begin{proof}
若 $f_n\rightrightarrows f$，将 $|f_n-f_m|$ 用 $|f_n-f|+|f_m-f|$ 控制即可。反之，式~\eqref{eq:uniform-cauchy-criterion} 对每个固定 $x$ 给出实数柯西列，故由实数系的完备性存在 $f(x)=\lim_m f_m(x)$。给定 $\varepsilon>0$，先把柯西条件中的误差取为 $\varepsilon/2$，固定 $n\geq N$ 后令 $m\to\infty$，得到对全部 $x\in E$ 成立的
\[
    |f_n(x)-f(x)|\leq\frac\varepsilon2<\varepsilon.
\]
于是得到一致收敛。将该准则用于部分和序列，就得到级数形式；有限个初始项不影响结论。
\end{proof}

\begin{corollary}
\label{cor:uniform-series-term-necessary-condition}
若 $\sum_{n=1}^\infty u_n$ 在 $E$ 上一致收敛，则 $u_n\rightrightarrows0$ 于 $E$。
\end{corollary}

\begin{proof}
在式~\eqref{eq:uniform-series-cauchy-criterion} 中取 $p=q=n$。
\end{proof}

通项一致趋于零仍只是必要条件。例如 $u_n(x)=1/n$ 在任何非空 $E$ 上都一致趋于零，但相应级数在每一点都发散。与定义相比，柯西准则的主要作用是无需预先求出极限函数或和函数。

\begin{proposition}[补入端点]
\label{prop:uniform-convergence-extension-to-endpoints}
设 $a<b$，每个 $f_n\in C([a,b])$，且 $\{f_n\}$ 在 $(a,b)$ 上一致收敛，则它在 $[a,b]$ 上也一致收敛。相应地，若级数各项在 $[a,b]$ 上连续，且级数在 $(a,b)$ 上一致收敛，则在 $[a,b]$ 上也一致收敛。
\end{proposition}

\begin{proof}
在开区间上用柯西条件将 $|f_n-f_m|$ 控制在 $\varepsilon/2$ 内。固定充分大的 $n,m$，分别令 $x\to a+$、$x\to b-$，由连续性得到端点处的差不超过 $\varepsilon/2$，故在整个闭区间上小于 $\varepsilon$。再用柯西准则即得结论。对级数只需将 $f_n$ 换成连续的部分和。
\end{proof}

例如 $\sum_{n=1}^\infty(-1)^n/(1+nx)$ 在每个 $x>0$ 处由交错级数判别法收敛，在 $[a,+\infty)$（$a>0$）上的余项绝对值不超过 $1/(1+(N+1)a)$，所以一致收敛。但它在 $(0,+\infty)$ 上不一致收敛，因为取 $x_n=1/n$ 时通项绝对值恒为 $1/2$，不满足推论~\ref{cor:uniform-series-term-necessary-condition}。也可反证：若在 $(0,+\infty)$ 上一致收敛，则限制到 $(0,1)$ 并补入端点后，在 $x=0$ 也应收敛，与 $\sum(-1)^n$ 发散矛盾。

\subsection{一致收敛的判别法}
\label{subsec:uniform-convergence-tests}

\begin{theorem}[上确界判别法]
\label{thm:uniform-convergence-supremum-test}
设 $f_n,f:E\to\mathbb R$，记扩展实数
\[
    d_n=\sup_{x\in E}|f_n(x)-f(x)|\in[0,+\infty].
\]
则 $f_n\rightrightarrows f$ 于 $E$ 的充分必要条件是 $d_n\to0$。级数一致收敛于 $S$ 的相应条件为 $\sup_{x\in E}|S_N(x)-S(x)|\to0$。
\end{theorem}

\begin{proof}
若一致收敛，则对给定 $\varepsilon>0$，充分大的 $n$ 满足全部误差小于 $\varepsilon/2$，从而 $d_n\leq\varepsilon/2<\varepsilon$。反之，$d_n<\varepsilon$ 直接给出对全部 $x$ 成立的误差界。
\end{proof}

在误差有界时，可将该上确界记为一致范数 $\|f_n-f\|_{\infty,E}$。判别法只需上确界，不要求最大值能够取到；“最值判别法”是其常用名称，不能因此遗漏开区间或无界区间。例如在实轴上，
\[
    f_n(x)=\frac{x}{1+n^2x^2},\qquad
    \sup_{x\in\mathbb R}|f_n(x)|=\frac1{2n}\longrightarrow0.
\]
这个上界可直接由 $1+n^2x^2\geq2n|x|$ 得到，并在 $x=\pm1/n$ 处取到，故 $f_n\rightrightarrows0$。

\begin{example}
\label{ex:uniform-composite-exponential-limit}
设 $a>0$，证明 $f_n(x)=n^2(e^{1/(nx)}-1)\sin(1/(nx))$ 在 $[a,+\infty)$ 上一致收敛于 $1/x^2$。
\end{example}

\begin{solve}
令 $H(t)=((e^t-1)/t)(\sin t/t)$（$t\neq0$），并令 $H(0)=1$。由两个基本极限，$H$ 在零点连续。令 $t=1/(nx)$，则
\[
    \left|f_n(x)-\frac1{x^2}\right|
       =\frac1{x^2}|H(t)-1|
       \leq\frac1{a^2}\sup_{0\leq t\leq1/(na)}|H(t)-1|
       \longrightarrow0.
\]
最后的极限来自零点连续性的定义，且该界与 $x$ 无关。证明一致收敛时，不能仅将对每个固定 $x$ 成立的小 $o$ 余项直接当作关于全部 $x$ 的统一估计；这里通过 $0<t\leq1/(na)$ 补足了这一控制。
\end{solve}

\begin{definition}
\label{def:absolute-uniform-convergence}
若 $\sum_{n=1}^\infty|u_n(x)|$ 在 $E$ 上一致收敛，则称 $\sum_{n=1}^\infty u_n(x)$ 在 $E$ 上\emph{绝对一致收敛}（absolute uniform convergence）。
\end{definition}

由柯西准则和三角不等式，绝对一致收敛必推出一致收敛及每一点的绝对收敛。反向一般不成立：\emph[caution]{“一致收敛且逐点绝对收敛”并不等于“绝对一致收敛”}。此外，$\sum_n\sup_{x\in E}|u_n(x)|<\infty$ 是绝对一致收敛的充分条件，不能将它与定义混同。

\begin{theorem}[魏尔斯特拉斯 $M$ 判别法]
\label{thm:weierstrass-uniform-m-test}
若存在非负实数 $M_n$，使全部 $x\in E$ 满足 $|u_n(x)|\leq M_n$，且数项级数 $\sum_{n=1}^\infty M_n$ 收敛，则 $\sum_{n=1}^\infty u_n$ 在 $E$ 上绝对一致收敛。这称为\emph{魏尔斯特拉斯 $M$ 判别法}（Weierstrass $M$-test）。
\end{theorem}

\begin{proof}
对任意 $q\geq p$，有
\[
    \sup_{x\in E}\sum_{k=p}^q|u_k(x)|\leq\sum_{k=p}^qM_k.
\]
右边由数项级数的柯西准则可一致控制，再用式~\eqref{eq:uniform-series-cauchy-criterion} 即可。
\end{proof}

\begin{example}
\label{ex:uniform-exponential-function-series}
设实数 $\alpha>1$，证明 $\sum_{n=1}^\infty x^\alpha e^{-nx}$ 在 $[0,+\infty)$ 上绝对一致收敛，其中 $x=0$ 时每项定义为零。
\end{example}

\begin{solve}
对固定 $n$，函数在正半轴上的导数为 $x^{\alpha-1}e^{-nx}(\alpha-nx)$，因此最大值在 $x=\alpha/n$ 处取得，且
\[
    \sup_{x\geq0}x^\alpha e^{-nx}
    =\left(\frac\alpha e\right)^\alpha\frac1{n^\alpha}.
\]
由 $p$ 级数收敛及 $M$ 判别法得到结论。这里控制的是整个正半轴上的最大误差，而不是仅对固定 $x$ 使用指数衰减。
\end{solve}

当正负抵消对收敛起主要作用时，应使用狄利克雷或阿贝尔判别法。两者的共同工具是\emph{阿贝尔变换}（Abel transformation），也称分部求和。对固定 $x$ 和 $q\geq p$，记 $A_k(x)=\sum_{j=p}^ku_j(x)$，直接展开可得
\begin{equation}
    \sum_{k=p}^qu_kv_k
    =A_qv_q+\sum_{k=p}^{q-1}A_k(v_k-v_{k+1}).
    \label{eq:uniform-abel-summation}
\end{equation}
若对这一段全部 $k$ 有 $|A_k|\leq B$，且 $v_k$ 关于 $k$ 单调，则
\begin{equation}
    \left|\sum_{k=p}^qu_kv_k\right|
    \leq B\bigl(|v_q|+|v_p-v_q|\bigr).
    \label{eq:uniform-abel-summation-bound}
\end{equation}
单调性使差的绝对值之和可以望远镜求和；这里及以下判别法中的单调性都指项指标 $n$ 的单调性。

\begin{theorem}[狄利克雷判别法]
\label{thm:uniform-dirichlet-test}
设 $u_n,v_n:E\to\mathbb R$。若部分和一致有界，即存在 $M>0$，使
\[
    \left|\sum_{k=1}^nu_k(x)\right|\leq M,
    \qquad n\geq1,\ x\in E,
\]
并且对每个固定 $x$，$v_n(x)$ 关于 $n$ 单调，且 $v_n\rightrightarrows0$ 于 $E$，则 $\sum_{n=1}^\infty u_nv_n$ 在 $E$ 上一致收敛。这称为\emph{狄利克雷判别法}（Dirichlet's test）。
\end{theorem}

\begin{proof}
任意尾段的部分和 $A_k$ 满足 $|A_k|\leq2M$。固定 $x$ 后，单调趋于零的实数列 $v_n(x)$ 不越过零，故
\[
    |v_q(x)|+|v_p(x)-v_q(x)|=|v_p(x)|.
\]
由式~\eqref{eq:uniform-abel-summation-bound}，任意 $q\geq p$ 满足
\[
    \sup_{x\in E}\left|\sum_{k=p}^qu_k(x)v_k(x)\right|
    \leq2M\sup_{x\in E}|v_p(x)|\longrightarrow0.
\]
于是级数满足一致柯西条件。
\end{proof}

\begin{theorem}[阿贝尔判别法]
\label{thm:uniform-abel-test}
设 $\sum_{n=1}^\infty u_n$ 在 $E$ 上一致收敛；对每个固定 $x\in E$，$v_n(x)$ 关于 $n$ 单调，且存在 $C>0$，使全部 $n,x$ 满足 $|v_n(x)|\leq C$。则 $\sum_{n=1}^\infty u_nv_n$ 在 $E$ 上一致收敛。这称为\emph{阿贝尔判别法}（Abel's test）。
\end{theorem}

\begin{proof}
给定 $\varepsilon>0$，由 $\sum u_n$ 的一致柯西条件，可选 $N$，使 $q\geq k\geq p\geq N$ 时，全部 $x$ 都有 $|A_k(x)|<\varepsilon/(3C+1)$。式~\eqref{eq:uniform-abel-summation-bound} 中的另一因子满足
\[
    |v_q|+|v_p-v_q|\leq3C,
\]
故任意充分靠后的尾段均小于 $\varepsilon$，再用柯西准则即可。
\end{proof}

狄利克雷判别法要求一个因子的部分和一致有界，另一个因子一致趋于零；阿贝尔判别法要求前一个级数已经一致收敛，另一个因子只需单调且一致有界。阿贝尔判别法中虽有逐点极限 $v(x)=\lim_n v_n(x)$，却不能擅自假设 $v_n\rightrightarrows v$，再套用狄利克雷判别法；$v_n(x)=x^n$ 在 $[0,1]$ 上就是反例。

\begin{example}
\label{ex:uniform-dirichlet-and-abel-applications}
应用上述判别法讨论 $\sum_{n=1}^\infty(-1)^n/(n+x)$，以及给定实数列 $\{a_n\}$ 后的 $\sum_{n=1}^\infty a_n/n^x$。
\end{example}

\begin{solve}
对于第一个级数，在 $[0,+\infty)$ 上取 $u_n=(-1)^n$、$v_n(x)=1/(n+x)$。前者的部分和绝对值不超过 $1$，后者关于 $n$ 递减且 $0<v_n(x)\leq1/n$，所以级数一致收敛。若 $I$ 是不含负整数的任意有限闭区间，则从某个指标起，对全部 $x\in I$ 都有 $n+x\geq n/2>0$；舍去有限个初始项后仍可用相同判别法，故在 $I$ 上也一致收敛。对任意固定的非极点 $x$，$1/|n+x|\sim1/n$，所以该级数在每一点都不绝对收敛。

对于第二个级数，它在 $[1,+\infty)$ 上一致收敛的充分必要条件是 $\sum_{n=1}^\infty a_n/n$ 收敛。必要性取 $x=1$ 即得。充分性将通项写为
\[
    \frac{a_n}{n^x}=\frac{a_n}{n}\,\frac1{n^{x-1}}.
\]
与 $x$ 无关的级数 $\sum a_n/n$ 在该集合上一致收敛，而 $n^{1-x}$ 对每个 $x\geq1$ 关于 $n$ 单调，且始终介于 $0$ 与 $1$ 之间，故由阿贝尔判别法得到结论。
\end{solve}

\begin{example}
\label{ex:uniform-and-pointwise-absolute-not-absolute-uniform}
证明 $\sum_{n=1}^\infty(-1)^n(1-x)x^n$ 在 $[0,1]$ 上一致收敛，且每一点都绝对收敛，但不绝对一致收敛。
\end{example}

\begin{solve}
取 $u_n=(-1)^n$、$v_n=(1-x)x^n$。对每个固定 $x\in[0,1]$，$v_n$ 关于 $n$ 非增；由例~\ref{ex:uniform-powers-and-endpoint-factor}，$v_n\rightrightarrows0$，故狄利克雷判别法给出一致收敛。每点的绝对值级数是几何级数，其和为
\[
    A(x)=\sum_{n=1}^\infty(1-x)x^n
    =\begin{cases}x,&0\leq x<1,\\0,&x=1.\end{cases}
\]
在 $x<1$ 处，绝对值级数的第 $N$ 个余项为 $x^{N+1}$，其在 $[0,1)$ 上的上确界始终为 $1$，故绝对值级数不一致收敛。
\end{solve}

\section{一致收敛的函数序列和函数项级数}
\label{sec:uniformly-convergent-function-sequences-and-series}

一致收敛使误差估计能够先对所有自变量统一成立，再固定某个充分靠后的函数使用其性质。下面连续性证明中的三项分解、积分中的统一误差界，以及可导性证明中的差商控制，都是这一原则的具体形式。

\subsection{极限函数的连续性}
\label{subsec:continuity-of-limit-functions}

\begin{theorem}[连续性与一致极限]
\label{thm:continuity-of-uniform-limit}
设 $E\subseteq\mathbb R$，$x_0\in E$，每个 $f_n$ 在 $x_0$ 处相对于 $E$ 连续，且 $f_n\rightrightarrows f$ 于 $E$。则 $f$ 在 $x_0$ 处连续。特别地，连续函数序列的一致极限仍连续；各项连续且一致收敛的函数项级数，其和函数连续。
\end{theorem}

\begin{proof}
给定 $\varepsilon>0$，先固定一个充分大的 $n_0$，使全部 $x\in E$ 满足 $|f_{n_0}(x)-f(x)|<\varepsilon/3$。再由这个固定函数在 $x_0$ 的连续性选取 $\delta>0$，使 $x\in E$、$|x-x_0|<\delta$ 时，$|f_{n_0}(x)-f_{n_0}(x_0)|<\varepsilon/3$。于是
\[
    |f(x)-f(x_0)|
    \leq |f(x)-f_{n_0}(x)|
       +|f_{n_0}(x)-f_{n_0}(x_0)|
       +|f_{n_0}(x_0)-f(x_0)|<\varepsilon.
\]
对级数应用于连续的部分和即可。
\end{proof}

\begin{theorem}[一致收敛下交换两次极限]
\label{thm:exchange-limits-under-uniform-convergence}
设 $x_0$ 为 $E$ 的聚点，$f_n\rightrightarrows f$ 于 $E\setminus\{x_0\}$，且每个有限极限 $a_n=\lim_{x\to x_0,\,x\in E}f_n(x)$ 存在。则 $\{a_n\}$ 收敛，$\lim_{x\to x_0,\,x\in E}f(x)$ 存在，并且
\begin{equation}
    \lim_{x\to x_0}\lim_{n\to\infty}f_n(x)
    =\lim_{n\to\infty}\lim_{x\to x_0}f_n(x).
    \label{eq:exchange-uniform-double-limits}
\end{equation}
式中的 $x$ 均限制在 $E\setminus\{x_0\}$ 内。若级数一致收敛且每项在 $x_0$ 有有限极限，则相应地可交换求和与该点极限。
\end{theorem}

\begin{proof}
一致柯西条件中先固定充分大的 $m,n$，再令 $x\to x_0$，得 $|a_n-a_m|\leq\varepsilon$，故 $a_n\to a\in\mathbb R$。为证明 $f(x)\to a$，将
\[
    |f(x)-a|\leq|f(x)-f_{n_0}(x)|
          +|f_{n_0}(x)-a_{n_0}|+|a_{n_0}-a|
\]
的首尾两项通过固定充分大的 $n_0$ 控制，再通过 $x\to x_0$ 控制中间一项。对级数，将 $f_n$ 换为部分和，$a_n$ 换为各项极限的有限和即可。
\end{proof}

连续性是局部性质，因而只在 $x_0$ 的一个相对邻域内一致收敛，已经足够保证该点连续。对开区间 $I$，若每个 $x_0\in I$ 都有一个邻域，使 $f_n$ 在其中一致收敛，则称\emph{局部一致收敛}（locally uniform convergence）；若在任意有限闭子区间 $[c,d]\subset I$ 上都一致收敛，则称\emph{内闭一致收敛}（uniform convergence on compact subintervals）。二者在开区间上等价：一个方向直接在点的附近取闭子区间；另一个方向由闭子区间的有限覆盖性，选出有限个局部邻域，并取所需项数的最大值。函数项级数使用其部分和作相同定义。因此，开区间上连续函数的内闭一致极限仍连续。

\begin{example}
\label{ex:locally-uniform-trigonometric-series}
若实数列 $\{a_n\}$ 单调且 $a_n\to0$，证明 $\sum_{n=1}^\infty a_n\cos(nx)$ 与 $\sum_{n=1}^\infty a_n\sin(nx)$ 的和函数在 $(0,2\pi)$ 内连续。
\end{example}

\begin{solve}
对 $0<x<2\pi$，有限和公式为
\[
    \sum_{k=1}^N\cos(kx)
      =\frac{\sin((N+\tfrac12)x)-\sin(x/2)}{2\sin(x/2)},
    \qquad
    \sum_{k=1}^N\sin(kx)
      =\frac{\cos(x/2)-\cos((N+\tfrac12)x)}{2\sin(x/2)}.
\]
在任意 $[\delta,2\pi-\delta]$（$0<\delta<\pi$）上，两者绝对值都不超过 $1/\sin(\delta/2)$；$a_n$ 与 $x$ 无关，故一致趋于零。由狄利克雷判别法，两级数均内闭一致收敛；各项连续，所以和函数在整个开区间内连续。靠近 $0$ 或 $2\pi$ 时，上述部分和界可能增大，不能据此声称在整个开区间上一致收敛。
\end{solve}

\begin{theorem}[狄尼定理]
\label{thm:dini-uniform-convergence}
设 $a<b$，$f_n\in C([a,b])$，在 $[a,b]$ 上逐点收敛于 $f$，且对每个固定 $x$，$f_n(x)$ 关于 $n$ 单调。则 $f\in C([a,b])$ 的充分必要条件是 $f_n\rightrightarrows f$ 于 $[a,b]$。这一结论称为\emph{狄尼定理}（Dini's theorem）。
\end{theorem}

\begin{proof}
一致收敛推出连续性已证。反过来，若 $f$ 连续，令 $g_n=|f_n-f|$，则每个 $g_n$ 连续，且 $g_n(x)$ 对每点单调下降到零。给定 $\varepsilon>0$，记相对开集
\[
    U_n=\{x\in[a,b]:g_n(x)<\varepsilon\}.
\]
有 $U_n\subseteq U_{n+1}$，且这些集合覆盖 $[a,b]$。由有限覆盖定理，可取有限子覆盖；设其中最大的指标为 $N$，由嵌套关系便有 $U_N=[a,b]$。所以全部 $n\geq N$、$x\in[a,b]$ 都满足 $g_n(x)<\varepsilon$，即一致收敛。
\end{proof}

\begin{corollary}[非负项级数的狄尼定理]
\label{cor:dini-nonnegative-function-series}
设 $u_n\in C([a,b])$、$u_n\geq0$，且 $\sum_{n=1}^\infty u_n$ 在 $[a,b]$ 上逐点收敛于 $S$。则该级数一致收敛的充分必要条件是 $S$ 连续。
\end{corollary}

\begin{proof}
连续的部分和 $S_N$ 关于 $N$ 逐点非减，直接应用狄尼定理。
\end{proof}

狄尼定理的条件各有作用。$x^n$ 在 $[0,1]$ 上单调且逐点收敛，但极限不连续，故不一致收敛；在 $[0,1)$ 上其极限恒为零且连续，却仍不一致收敛，说明紧的闭区间条件不能省去。另一方面，在 $[0,1]$ 上，$h_n(x)=nx^2e^{-nx^2}$ 逐点趋于连续函数零，但 $h_n(1/\sqrt n)=1/e$，故不一致收敛；这里缺少关于 $n$ 的单调性。判定极限是否连续时，一致收敛是充分条件，而不是一般的必要条件。

\subsection{极限函数的积分}
\label{subsec:integration-of-limit-functions}

\begin{theorem}[一致极限的可积性与积分极限]
\label{thm:integration-of-uniform-limit}
设 $a<b$，每个 $f_n$ 在 $[a,b]$ 上黎曼可积，且 $f_n\rightrightarrows f$ 于 $[a,b]$。则 $f$ 在 $[a,b]$ 上黎曼可积，且
\begin{equation}
    \int_a^bf(x)\,dx
       =\lim_{n\to\infty}\int_a^bf_n(x)\,dx.
    \label{eq:uniform-limit-integration}
\end{equation}
此外，若 $F_n(x)=\int_a^xf_n(t)\,dt$、$F(x)=\int_a^xf(t)\,dt$，则 $F_n\rightrightarrows F$ 于 $[a,b]$。
\end{theorem}

\begin{proof}
由一致收敛，固定充分大的 $n_0$ 可使 $|f-f_{n_0}|\leq1$，故 $f$ 有界。给定 $\varepsilon>0$，再选 $n_0$ 使 $\|f-f_{n_0}\|_\infty<\eta=\varepsilon/(4(b-a))$。对任意子区间 $J$，记 $\omega_J(h)=\sup_Jh-\inf_Jh$ 为振幅，则
\[
    \omega_J(f)\leq\omega_J(f_{n_0})+2\eta.
\]
由于 $f_{n_0}$ 可积，可选分割 $P$，使其上、下达布和之差小于 $\varepsilon/2$，于是
\[
    U(f,P)-L(f,P)
    \leq U(f_{n_0},P)-L(f_{n_0},P)+2\eta(b-a)<\varepsilon.
\]
由达布可积准则，$f$ 可积。再由积分的保序性及绝对值不等式，
\begin{equation}
    \left|\int_a^b(f_n-f)\,dx\right|
       \leq(b-a)\|f_n-f\|_\infty\longrightarrow0.
    \label{eq:uniform-integral-error-bound}
\end{equation}
同样有 $|F_n(x)-F(x)|\leq(b-a)\|f_n-f\|_\infty$，此界与 $x$ 无关，故变上限积分序列也一致收敛。
\end{proof}

\begin{corollary}[逐项积分]
\label{cor:termwise-integration-uniform-series}
若每个 $u_n$ 在有限闭区间 $[a,b]$ 上可积，且 $\sum_{n=1}^\infty u_n$ 在该区间上一致收敛于 $S$，则 $S$ 可积，且
\begin{equation}
    \int_a^b\left(\sum_{n=1}^\infty u_n(x)\right)\,dx
       =\sum_{n=1}^\infty\int_a^bu_n(x)\,dx.
    \label{eq:uniform-series-termwise-integration}
\end{equation}
这种交换称为\emph{逐项积分}（termwise integration）。
\end{corollary}

\begin{proof}
将定理用于部分和 $S_N$，再使用有限和的积分线性性质。
\end{proof}

这些结论不要求每项连续，但要求在同一个有限闭区间上一致收敛。内闭一致收敛只保证在各个有限闭子区间内可交换积分；对无穷区间、瑕积分或趋向收敛域边界的积分上限，还需额外控制积分尾部，不能直接沿用式~\eqref{eq:uniform-integral-error-bound}。

\begin{example}
\label{ex:logarithm-series-by-termwise-integration}
证明 $|x|<1$ 时
\begin{equation}
    \frac12\ln\frac{1+x}{1-x}
       =\sum_{n=0}^\infty\frac{x^{2n+1}}{2n+1}.
    \label{eq:logarithm-odd-power-series}
\end{equation}
\end{example}

\begin{solve}
由几何级数，$1/(1-t^2)=\sum_{n=0}^\infty t^{2n}$。固定 $|x|<1$，在以 $0,x$ 为端点的闭区间上，$|t^{2n}|\leq|x|^{2n}$，由 $M$ 判别法可逐项积分，故
\[
    \int_0^x\frac{dt}{1-t^2}
       =\sum_{n=0}^\infty\int_0^xt^{2n}\,dt
       =\sum_{n=0}^\infty\frac{x^{2n+1}}{2n+1}.
\]
左边积分为 $\frac12\ln((1+x)/(1-x))$，得到所求公式。$x<0$ 时使用定积分的方向约定，推导仍然成立；这里用到的是内闭一致收敛，无须声称原几何级数在 $(-1,1)$ 上一致收敛。
\end{solve}

\begin{example}
\label{ex:euler-constant-integral-series}
计算 $\displaystyle I=\int_0^1\sum_{n=1}^\infty\frac{x}{n(n+x)}\,dx$。
\end{example}

\begin{solve}
在 $[0,1]$ 上，$0\leq x/(n(n+x))\leq1/n^2$，故由 $M$ 判别法及逐项积分定理，
\[
    I=\sum_{n=1}^\infty\int_0^1
          \left(\frac1n-\frac1{n+x}\right)\,dx
      =\lim_{N\to\infty}\left(\sum_{n=1}^N\frac1n-\ln(N+1)\right)
      =\gamma.
\]
这里 $\gamma=\lim_{N\to\infty}(H_N-\ln N)$ 为欧拉常数，$H_N=\sum_{n=1}^N1/n$。该极限存在，因为由 $1/(N+1)<\int_N^{N+1}dt/t$，$H_N-\ln N$ 单调下降；又 $\ln(N+1)=\int_1^{N+1}dt/t\leq H_N$，故它有下界。将 $\ln N$ 换为 $\ln(N+1)$ 只改变一个趋于零的量。
\end{solve}

\subsection{极限函数的导数}
\label{subsec:derivatives-of-limit-functions}

微分考察任意小尺度上的增量，因此对函数值的一致控制仍不足以控制导数；例~\ref{ex:pointwise-limits-and-operation-failures} 中的 $\sin(nx)/\sqrt n$ 已给出反例。另一方面，只知道导函数一致收敛也不够确定原函数，因为每一项还可加上任意常数。例如 $f_n(x)=n$ 的导函数恒为零，原函数却处处发散。因此需要同时控制导函数，并在至少一个点固定原函数的常数部分。

\begin{theorem}[导数与极限的交换]
\label{thm:differentiation-of-function-sequence}
设 $a<b$，每个 $f_n$ 在 $[a,b]$ 上可导，端点取有限单侧导数。若存在 $x_0\in[a,b]$，使数列 $\{f_n(x_0)\}$ 收敛，且 $f_n'\rightrightarrows g$ 于 $[a,b]$，则存在函数 $f$，使 $f_n\rightrightarrows f$ 于 $[a,b]$，$f$ 在该区间上可导，且
\begin{equation}
    f'(x)=g(x)=\lim_{n\to\infty}f_n'(x),
    \qquad x\in[a,b].
    \label{eq:uniform-derivative-limit}
\end{equation}
端点处的结论按相应单侧理解；本定理不要求各 $f_n'$ 连续。
\end{theorem}

\begin{proof}
先证原函数序列一致收敛。对 $f_n-f_m$ 在 $x_0$ 与任意 $x$ 之间使用拉格朗日中值定理，得
\[
    |f_n(x)-f_m(x)|
    \leq|f_n(x_0)-f_m(x_0)|
       +(b-a)\|f_n'-f_m'\|_\infty.
\]
右边由给定点处的柯西条件及导函数的一致柯西条件趋于零，故 $f_n\rightrightarrows f$。

再固定 $y\in[a,b]$，对 $x\neq y$ 记
\[
    Q_n(x)=\frac{f_n(x)-f_n(y)}{x-y},\qquad
    Q(x)=\frac{f(x)-f(y)}{x-y}.
\]
中值定理给出 $|Q_n(x)-Q_m(x)|\leq\|f_n'-f_m'\|_\infty$，所以 $\{Q_n\}$ 在 $[a,b]\setminus\{y\}$ 上一致柯西，并逐点收敛于 $Q$，因而一致收敛。另一方面，$\lim_{x\to y}Q_n(x)=f_n'(y)$。由定理~\ref{thm:exchange-limits-under-uniform-convergence} 可交换这两个极限，得到 $\lim_{x\to y}Q(x)=\lim_nf_n'(y)=g(y)$，即 $f'(y)=g(y)$。端点采用相应单侧极限，证明相同。
\end{proof}

\begin{corollary}[逐项求导]
\label{cor:termwise-differentiation-uniform-derivative-series}
设每个 $u_n$ 在 $[a,b]$ 上可导，端点按单侧理解。若 $\sum_{n=1}^\infty u_n(x_0)$ 在某个 $x_0\in[a,b]$ 处收敛，且导函数级数 $\sum_{n=1}^\infty u_n'$ 在 $[a,b]$ 上一致收敛，则原级数在该区间上一致收敛，和函数可导，并且
\begin{equation}
    \left(\sum_{n=1}^\infty u_n(x)\right)'
       =\sum_{n=1}^\infty u_n'(x).
    \label{eq:termwise-differentiation-function-series}
\end{equation}
这一交换称为\emph{逐项求导}（termwise differentiation）。
\end{corollary}

\begin{proof}
将前一定理应用于部分和序列。
\end{proof}

在开区间内，可导性同样是局部性质。因此，若 $f_n$ 已逐点收敛、$f_n'$ 内闭一致收敛，就能在每个闭子区间上应用上述定理，得到极限函数可导以及式~\eqref{eq:uniform-derivative-limit}。对级数也有相同结论。若只给出一个点的收敛，则对任意有限闭子区间，可先取同时包含它与该点的更大闭子区间，借此把常数部分的控制传递过去。上述结论保证有限闭子区间上的一致收敛，不能直接在无界区间上把中值定理中的长度替换为一个不存在的有限统一界。

\begin{example}
\label{ex:zeta-function-smoothness-by-series}
证明函数 $\displaystyle\zeta(x)=\sum_{n=1}^\infty n^{-x}$ 在 $(1,+\infty)$ 内具有任意阶连续导数，并写出其表达式。
\end{example}

\begin{solve}
固定 $\delta>0$ 和非负整数 $m$，通项的 $m$ 阶导数为 $(-\ln n)^m/n^x$，其中 $m=0$ 时将分子理解为 $1$。在 $x\geq1+\delta$ 上，
\[
    \left|\frac{(-\ln n)^m}{n^x}\right|
    \leq\frac{(\ln n)^m}{n^{1+\delta}}.
\]
对固定 $m,\delta$，对数增长慢于任意正幂，故充分大的 $n$ 满足 $(\ln n)^m\leq n^{\delta/2}$；右边因而可由收敛的 $p$ 级数 $\sum n^{-1-\delta/2}$ 控制。$M$ 判别法给出各阶导函数级数的一致收敛。先由连续性定理，再在任意有限闭子区间内反复逐项求导，得到
\[
    \zeta^{(m)}(x)=\sum_{n=1}^\infty\frac{(-\ln n)^m}{n^x},
    \qquad x>1.
\]
每个和函数均连续，故 $\zeta\in C^\infty((1,+\infty))$。这里 $\zeta$ 称为实变量上的\emph{黎曼 $\zeta$ 函数}（Riemann zeta function）；本例只讨论 $x>1$ 的级数定义。
\end{solve}

\begin{example}
\label{ex:series-with-prescribed-discontinuity-set}
设 $\{x_n\}\subset[0,1]$ 各项互不相同。构造一个有界函数，其在 $[0,1]$ 上的不连续点恰为全部 $x_n$；再构造一个可导函数，其导函数的不连续点恰为这些点。
\end{example}

\begin{solve}
定义符号函数 $\operatorname{sgn}(t)$ 在 $t>0$、$t=0$、$t<0$ 时分别为 $1,0,-1$，并令
\[
    F(x)=\sum_{n=1}^\infty\frac{\operatorname{sgn}(x-x_n)}{2^n},
    \qquad 0\leq x\leq1.
\]
每项绝对值不超过 $2^{-n}$，所以级数一致收敛，且 $|F|\leq1$。若 $x\notin\{x_n\}$，则每一项都在该点连续，故 $F$ 连续。若 $x=x_k$，去掉第 $k$ 项后，其余级数仍一致收敛，且每一项在 $x_k$ 连续，所以其和在该点连续；第 $k$ 项在该点不连续，因而 $F$ 也不连续。端点使用相对连续性，结论仍成立。

再令
\[
    h(t)=\begin{cases}t^2\sin(1/t),&t\neq0,\\0,&t=0,\end{cases}
    \qquad
    G(x)=\sum_{n=1}^\infty\frac{h(x-x_n)}{2^n}.
\]
$h$ 处处可导，$h'(0)=0$，而 $t\neq0$ 时 $h'(t)=2t\sin(1/t)-\cos(1/t)$；$h'$ 恰在零点不连续。在 $|t|\leq1$ 上有 $|h(t)|\leq1$、$|h'(t)|\leq3$，故原级数和导函数级数均由 $M$ 判别法一致收敛，逐项求导得到
\[
    G'(x)=\sum_{n=1}^\infty\frac{h'(x-x_n)}{2^n}.
\]
在 $x\notin\{x_n\}$ 处，各项连续，故 $G'$ 连续；在 $x=x_k$ 处，去掉第 $k$ 项后的和连续，而第 $k$ 项不连续，所以 $G'$ 不连续。端点导数按单侧理解。即使选取一个在 $[0,1]$ 中稠密的点列，这两个构造仍然成立；因此逐项求导定理本身并不要求导函数连续。
\end{solve}

教材最后还用三角波级数说明，连续函数甚至可以处处不可导。下面给出这一构造，并利用逐级网格上的割线斜率核查其性质。

\begin{example}
\label{ex:continuous-nowhere-differentiable-triangular-series}
设 $\varphi(x)=\operatorname{dist}(x,\mathbb Z)=\inf_{k\in\mathbb Z}|x-k|$，即把 $[-1/2,1/2]$ 上的 $|x|$ 作周期为 $1$ 的延拓。证明
\[
    W(x)=\sum_{n=1}^\infty\frac{\varphi(4^nx)}{4^n},
    \qquad x\in\mathbb R,
\]
连续，但没有任何一点存在有限导数。
\end{example}

\begin{solve}
由 $0\leq\varphi\leq1/2$ 及 $\sum4^{-n}<\infty$，$M$ 判别法保证级数在实轴上一致收敛；各项连续，所以 $W$ 连续。

固定 $x$，先设 $x$ 不属于任何网格 $4^{-p}\mathbb Z$。以 $\lfloor t\rfloor$ 表示不大于 $t$ 的最大整数，令
\[
    a_p=\frac{\lfloor4^px\rfloor}{4^p},\qquad b_p=a_p+4^{-p},
    \qquad D_p=\frac{W(b_p)-W(a_p)}{b_p-a_p}.
\]
有 $a_p<x<b_p$，二者均趋于 $x$。对 $n\geq p$，$4^na_p$、$4^nb_p$ 都是整数，所以相应项在两端均为零。对 $n<p$，第 $n$ 项的折点为 $k/(2\cdot4^n)$（$k\in\mathbb Z$），均落在 $4^{-p}\mathbb Z$ 中，因此该项在 $[a_p,b_p]$ 上为斜率等于 $1$ 或 $-1$ 的线性函数。由网格嵌套，进入下一层时，原有各项的斜率保持不变，仅新增第 $p$ 项的斜率，故
\[
    D_{p+1}-D_p\in\{-1,1\}.
\]
因此 $\{D_p\}$ 不可能收敛。若 $W'(x)=L\in\mathbb R$ 存在，由
\[
    W(t)=W(x)+L(t-x)+o(|t-x|)
\]
分别取 $t=a_p,b_p$，并利用 $|b_p-x|+|x-a_p|=b_p-a_p$，却会得到 $D_p\to L$，矛盾。

若 $x\in4^{-p_0}\mathbb Z$，改用 $[x,x+4^{-p}]$（$p\geq p_0$）的右侧割线。有限个 $n<p_0$ 的项在充分短的右邻域内斜率固定；对 $p_0\leq n<p$，$4^nx$ 为整数且 $0\leq4^n(t-x)\leq1/4$，相应项在该区间的斜率均为 $1$；对 $n\geq p$，端点值均为零。因此右侧割线斜率随 $p$ 趋于 $+\infty$，也不可能有有限导数。两种情形覆盖了全部实数。
\end{solve}

\section*{习题十}
\addcontentsline{toc}{section}{习题十}
\fi

\selectchapter{11}
\ifcompilechapter
\chapter{幂级数}
\label{chap:power-series}

% 本章正文依据数学分析-11.pdf 第 1--27 页，对应教材第 212--238 页；
% 第 27 页开始的习题十一不在本次整理范围内。
多项式容易计算、求导和积分，而幂级数的部分和都是多项式，因此幂级数提供了一条用简单函数研究复杂函数的途径。与一般函数项级数相比，幂级数的收敛域具有特别整齐的形状，并且在收敛区间内部可以任意次逐项求导和积分。本章先确定收敛范围，再研究和函数的性质以及初等函数的幂级数展开，最后讨论一个更一般的问题：即使函数不能展开成幂级数，是否仍能用多项式一致逼近？以下讨论实系数、实变量的幂级数，除另有说明外，级数的指标从 $0$ 开始。

\section{幂级数的收敛半径与收敛域}
\label{sec:radius-and-domain-of-convergence}

\subsection{幂级数的收敛半径与收敛域}
\label{subsec:radius-and-domain-of-convergence}

\begin{definition}
\label{def:real-power-series}
设 $x_0\in\mathbb R$，$a_n\in\mathbb R$。形如
\begin{equation}
    \sum_{n=0}^\infty a_n(x-x_0)^n
    =a_0+a_1(x-x_0)+a_2(x-x_0)^2+\cdots
    \label{eq:general-real-power-series}
\end{equation}
的函数项级数称为以 $x_0$ 为中心的\emph{幂级数}（power series），$a_n$ 称为系数。全部收敛点组成它的\emph{收敛域}（domain of convergence）；在收敛域上，级数的和定义其和函数。常数项在 $x=x_0$ 时取值为 $a_0$，故中心总是收敛点。
\end{definition}

作平移 $t=x-x_0$，即可把一般情形化为 $\sum a_nt^n$，所以先讨论中心为零的情形。以下几个级数说明，端点是否属于收敛域必须另外判断。

\begin{example}
\label{ex:power-series-basic-convergence-domains}
求 $\sum_{n=0}^\infty n!x^n$、$\sum_{n=0}^\infty x^n$、$\sum_{n=1}^\infty x^n/n$、$\sum_{n=1}^\infty x^n/n^2$ 及 $\sum_{n=0}^\infty x^n/n!$ 的收敛域，其中 $0!=1$。
\end{example}

\begin{solve}
当 $x\neq0$ 时，$n!|x|^n$ 的相邻项之比为 $(n+1)|x|\to+\infty$，故通项不趋于零，第一个级数仅在 $x=0$ 收敛。几何级数 $\sum x^n$ 在 $|x|<1$ 时绝对收敛，而在 $|x|\geq1$ 时通项不趋于零，所以收敛域为 $(-1,1)$。对于 $\sum x^n/n$，当 $|x|<1$ 时可用几何级数比较，当 $|x|>1$ 时通项不趋于零；在 $x=-1$ 处为收敛的交错级数，在 $x=1$ 处为发散的调和级数，因此收敛域为 $[-1,1)$。同理，$\sum x^n/n^2$ 的两个端点都绝对收敛，收敛域为 $[-1,1]$。最后，$x\neq0$ 时有
\[
    \frac{|x|^{n+1}/(n+1)!}{|x|^n/n!}
    =\frac{|x|}{n+1}\longrightarrow0,
\]
所以 $\sum x^n/n!$ 对每个实数都绝对收敛，收敛域为 $\mathbb R$。
\end{solve}

\begin{theorem}[由一个收敛点控制内部区间]
\label{thm:power-series-interior-absolute-uniform-convergence}
若 $\sum_{n=0}^\infty a_nx^n$ 在 $x=t\neq0$ 处收敛，则对每个 $0<r<|t|$，它在 $[-r,r]$ 上绝对一致收敛，即 $\sum |a_nx^n|$ 在该区间上一致收敛。
\end{theorem}

\begin{proof}
由收敛级数的通项趋于零，数列 $\{a_nt^n\}$ 有界，故存在 $M>0$，使所有 $n\geq0$ 都满足 $|a_nt^n|\leq M$。于是当 $|x|\leq r$ 时，
\[
    |a_nx^n|\leq M\left(\frac r{|t|}\right)^n.
\]
右边组成收敛的几何级数，由\crosschapterref{10}{定理~\ref{thm:weierstrass-uniform-m-test}}{魏尔斯特拉斯 $M$ 判别法} 即得结论。这里仅要求在 $t$ 处收敛，无须假设该点绝对收敛。
\end{proof}

\begin{theorem}[收敛半径与收敛域]
\label{thm:power-series-radius-and-domain}
对于任意幂级数 $\sum_{n=0}^\infty a_n(x-x_0)^n$，存在唯一的 $R\in[0,+\infty]$，使级数在 $|x-x_0|<R$ 时绝对收敛，并在这个开区间的每个有限闭子区间上绝对一致收敛；在 $|x-x_0|>R$ 时发散。$R$ 称为\emph{收敛半径}（radius of convergence）。当 $R=0$ 时，级数仅在 $x_0$ 处收敛；当 $R=+\infty$ 时，级数在整个实轴上收敛。当 $0<R<+\infty$ 时，两个端点 $x_0\pm R$ 的敛散性需要分别判断，故收敛域是
\[
    (x_0-R,x_0+R),\quad [x_0-R,x_0+R),\quad
    (x_0-R,x_0+R],\quad [x_0-R,x_0+R]
\]
四种区间中的一种。
\end{theorem}

\begin{proof}
平移后设 $x_0=0$，令
\[
    R=\sup\left\{|t|:\sum_{n=0}^\infty a_nt^n\text{ 收敛}\right\},
\]
集合无上界时规定 $R=+\infty$。若 $r<R$，由上确界的定义，可找到收敛点 $t$，使 $|t|>r$；前一定理保证级数在 $[-r,r]$ 上绝对一致收敛。任一有限闭子区间都可包含在这样的 $[-r,r]$ 内。若 $|x|>R$，按 $R$ 的定义该点不能收敛。$R=0$ 和 $R=+\infty$ 的结论也由定义及同一论证得到。唯一性来自内部收敛、外部发散这两个性质；有限正半径时，前例及将 $x$ 换成 $-x$ 得到的例子说明四种端点组合都可能出现。
\end{proof}

\emph[caution]{收敛半径控制的是开区间内部与闭区间外部，不能决定端点，也不能保证在整个开区间上一致收敛。}例如几何级数在 $(-1,1)$ 内闭绝对一致收敛，但在 $[0,1)$ 上的余项 $x^{N+1}/(1-x)$ 无法一致趋于零。

\subsection{收敛半径的求法}
\label{subsec:computing-the-radius-of-convergence}

\begin{theorem}[柯西--哈达玛定理]
\label{thm:cauchy-hadamard-radius}
设 $\sum a_n(x-x_0)^n$ 的收敛半径为 $R$，记
\[
    \rho=\limsup_{n\to\infty}\sqrt[n]{|a_n|}\in[0,+\infty].
\]
则
\begin{equation}
    R=\frac1\rho,
    \qquad \frac10:=+\infty,\quad \frac1{+\infty}:=0.
    \label{eq:cauchy-hadamard-radius}
\end{equation}
这称为\emph{柯西--哈达玛公式}（Cauchy--Hadamard formula）。
\end{theorem}

\begin{proof}
固定 $t=x-x_0\neq0$，有
\[
    \limsup_{n\to\infty}\sqrt[n]{|a_nt^n|}=|t|\rho.
\]
当 $0<\rho<+\infty$ 时，数项级数的根值判别法说明 $|t|\rho<1$ 时绝对收敛，$|t|\rho>1$ 时发散，因此 $R=1/\rho$。当 $\rho=0$ 时，每个有限 $t$ 都使根值上极限为零，故 $R=+\infty$；当 $\rho=+\infty$ 时，每个非零 $t$ 都使其为 $+\infty$，故 $R=0$。中心 $t=0$ 始终收敛，无须使用上述乘积式。
\end{proof}

\begin{corollary}[比值法]
\label{cor:power-series-ratio-radius}
若系数从某项起均不为零，且扩展实数极限
\[
    \lim_{n\to\infty}\left|\frac{a_{n+1}}{a_n}\right|=\rho
\]
存在，则仍有式~\eqref{eq:cauchy-hadamard-radius}。当 $0<\rho<+\infty$ 时，也可写成 $R=\lim_{n\to\infty}|a_n/a_{n+1}|$。
\end{corollary}

\begin{proof}
固定 $t\neq0$ 后，$|a_{n+1}t^{n+1}|/|a_nt^n|\to\rho|t|$，直接用比值判别法即可确定内部与外部的敛散性；$\rho=0,+\infty$ 时分别得到处处收敛与仅在中心收敛。
\end{proof}

比值极限不存在时，不能用比值法推出一个半径，但根值的上极限公式始终适用。对于\emph{缺项幂级数}（lacunary power series），必须按实际幂次确定系数，未出现的幂次系数为零；求根时不能把求和指标误当成幂次。

\begin{example}
\label{ex:lacunary-factorial-power-series}
求缺项幂级数 $\displaystyle\sum_{n=0}^\infty n!x^{n^2}$ 的收敛半径与收敛域。
\end{example}

\begin{solve}
在正幂次 $m=n^2$ 处系数为 $n!$，其余正幂次处为零。由
\[
    1\leq(n!)^{1/n^2}\leq(n^n)^{1/n^2}
    =n^{1/n}\longrightarrow1,
\]
可知完整系数序列的根值上极限为 $1$，故 $R=1$。在 $x=\pm1$ 处，通项的绝对值为 $n!$，不趋于零，所以收敛域为 $(-1,1)$。这里 $|x|<1$ 时，$x^{n^2}$ 的衰减能够压过 $n!$ 的增长。
\end{solve}

\begin{example}
\label{ex:shifted-and-parameter-power-series}
求 $\displaystyle\sum_{n=1}^\infty\frac{\ln(n+1)}{n^2}(x-3)^n$，以及带参数 $\alpha>0$ 的 $\displaystyle\sum_{n=1}^\infty\frac{(x+1)^n}{2^{n^\alpha}}$ 的收敛半径与收敛域。
\end{example}

\begin{solve}
第一个级数的中心为 $3$，且
\[
    \left|\frac{a_{n+1}}{a_n}\right|
    =\frac{n^2\ln(n+2)}{(n+1)^2\ln(n+1)}\longrightarrow1,
\]
故 $R=1$。在 $x=2,4$ 处，绝对值级数均为 $\sum\ln(n+1)/n^2$；因 $\ln(n+1)\leq\sqrt n$ 对充分大的 $n$ 成立，该级数收敛，所以收敛域为 $[2,4]$。

第二个级数的中心为 $-1$，系数满足
\[
    \sqrt[n]{a_n}=2^{-n^{\alpha-1}}
    \longrightarrow
    \begin{cases}
        1,&0<\alpha<1,\\
        1/2,&\alpha=1,\\
        0,&\alpha>1.
    \end{cases}
\]
因此半径依次为 $1,2,+\infty$。当 $0<\alpha<1$ 时，端点为 $-2,0$，绝对值级数为 $\sum2^{-n^\alpha}$；由 $n^\alpha/\ln n\to+\infty$，充分大的 $n$ 满足 $n^\alpha\ln2\geq2\ln n$，于是 $2^{-n^\alpha}\leq n^{-2}$，两个端点都绝对收敛。当 $\alpha=1$ 时，端点为 $-3,1$，通项绝对值恒为 $1$，两端均发散。当 $\alpha>1$ 时，没有有限端点需要检查。故其收敛域分别为 $[-2,0]$、$(-3,1)$ 与 $\mathbb R$。
\end{solve}

\section{幂级数的性质}
\label{sec:properties-of-power-series}

\begin{theorem}[阿贝尔定理]
\label{thm:abel-power-series-uniform-convergence}
设 $\sum a_n(x-x_0)^n$ 的收敛半径为 $R>0$。它在 $(x_0-R,x_0+R)$ 内闭绝对一致收敛，这一结论通常称为\emph{阿贝尔第一定理}（Abel's first theorem）。若 $R<+\infty$ 且级数在右端点 $x_0+R$ 收敛，则它在每个 $[c,x_0+R]$（$x_0-R<c<x_0+R$）上一致收敛；若在左端点 $x_0-R$ 收敛，则在每个 $[x_0-R,d]$（$x_0-R<d<x_0+R$）上一致收敛。端点处的这两个结论通常称为\emph{阿贝尔第二定理}（Abel's second theorem）。若两端都收敛，则在整个闭区间上一致收敛。
\end{theorem}

\begin{proof}
内部的结论已由定理~\ref{thm:power-series-radius-and-domain} 证明。若右端点收敛，在 $x\in[x_0,x_0+R]$ 上写成
\[
    a_n(x-x_0)^n=a_nR^n
        \left(\frac{x-x_0}{R}\right)^n.
\]
级数 $\sum a_nR^n$ 与 $x$ 无关，因而作为函数项级数一致收敛；后一因子关于 $n$ 单调不增，并一致有界于 $1$。由\crosschapterref{10}{定理~\ref{thm:uniform-abel-test}}{函数项级数的阿贝尔判别法}，乘积级数在该闭区间上一致收敛，再与内部的闭区间拼接即可。左端点的证明相同，只需写成
\[
    a_n(x-x_0)^n=a_n(-R)^n
        \left(\frac{x_0-x}{R}\right)^n,
    \qquad x_0-R\leq x\leq x_0.
\]
两端都收敛时，把左右两个闭区间合并，便得全区间的一致收敛。
\end{proof}

\begin{corollary}[和函数的连续性与端点极限]
\label{cor:power-series-continuity-and-abel-limits}
幂级数的和函数在其收敛域上连续，端点按相对连续性理解。特别地，当 $0<R<+\infty$ 时，若对应端点级数收敛，则
\begin{align}
    \lim_{x\to(x_0+R)-}\sum_{n=0}^\infty a_n(x-x_0)^n
        &=\sum_{n=0}^\infty a_nR^n,
        \label{eq:abel-right-endpoint-limit}\\
    \lim_{x\to(x_0-R)+}\sum_{n=0}^\infty a_n(x-x_0)^n
        &=\sum_{n=0}^\infty a_n(-R)^n.
        \label{eq:abel-left-endpoint-limit}
\end{align}
\end{corollary}

\begin{proof}
各项是连续的多项式，在内点附近及已经收敛的端点附近，都可使用前一定理得到闭区间上的一致收敛，再用连续函数一致极限的连续性即得。
\end{proof}

阿贝尔定理的逻辑方向需要区分：\emph[caution]{先有端点级数收敛，才能把内侧极限与端点求和交换；和函数在端点有有限内侧极限，并不足以反推端点级数收敛。}例如 $\sum_{n=0}^\infty(-1)^nx^n=1/(1+x)$ 在 $x\to1-$ 时趋于 $1/2$，但 $\sum(-1)^n$ 发散。

\begin{theorem}[逐项积分]
\label{thm:power-series-termwise-integration}
设 $f(x)=\sum_{n=0}^\infty a_n(x-x_0)^n$ 的收敛半径为 $R>0$，$\alpha,\beta$ 是其收敛域内任意两点，则
\begin{equation}
    \int_\alpha^\beta f(x)\,dx
    =\sum_{n=0}^\infty
       \frac{a_n}{n+1}
       \bigl[(\beta-x_0)^{n+1}-(\alpha-x_0)^{n+1}\bigr].
    \label{eq:power-series-integral-between-convergence-points}
\end{equation}
特别地，在收敛域内有
\begin{equation}
    \int_{x_0}^x f(t)\,dt
       =\sum_{n=0}^\infty\frac{a_n}{n+1}(x-x_0)^{n+1}.
    \label{eq:power-series-integral-from-center}
\end{equation}
\end{theorem}

\begin{proof}
先设 $\alpha\leq\beta$。若二者都在内部，由内闭一致收敛可逐项积分；若其中一个或两个为已经收敛的端点，则阿贝尔定理仍保证在 $[\alpha,\beta]$ 上一致收敛，由\crosschapterref{10}{推论~\ref{cor:termwise-integration-uniform-series}}{函数项级数的逐项积分推论} 得到同一结论。$\alpha>\beta$ 时由定积分的定向约定处理。
\end{proof}

\begin{theorem}[逐项求导及任意阶光滑性]
\label{thm:power-series-termwise-differentiation}
设 $f(x)=\sum a_n(x-x_0)^n$ 的收敛半径为 $R>0$，则 $f$ 在开区间 $(x_0-R,x_0+R)$ 上具有任意阶连续导数，且对每个正整数 $k$，
\begin{equation}
    f^{(k)}(x)=\sum_{n=k}^\infty
        n(n-1)\cdots(n-k+1)a_n(x-x_0)^{n-k}.
    \label{eq:power-series-arbitrary-derivatives}
\end{equation}
逐项求导或从中心逐项积分所得到的幂级数，都与原级数具有相同的收敛半径。
\end{theorem}

\begin{proof}
先说明一次求导与积分不改变半径。对求导级数记 $b_{n-1}=na_n$，对积分级数记 $c_{n+1}=a_n/(n+1)$，则对充分大的 $n$，
\[
    |b_{n-1}|^{1/(n-1)}
        =\bigl(n^{1/n}|a_n|^{1/n}\bigr)^{n/(n-1)},\qquad
    |c_{n+1}|^{1/(n+1)}
        =\left(\frac{|a_n|^{1/n}}{(n+1)^{1/n}}\right)^{n/(n+1)}.
\]
因 $n^{1/n},(n+1)^{1/n}\to1$，且 $n/(n\pm1)\to1$，这两列的上极限都等于 $\limsup|a_n|^{1/n}$，包括零与 $+\infty$ 两种情形。由柯西--哈达玛公式，它们的半径均为 $R$。反复应用可得任意阶求导级数也有半径 $R$。

任取内部的有限闭区间，并将它扩大为仍在内部且包含 $x_0$ 的闭区间。求导级数在其上一致收敛，原级数在 $x_0$ 处收敛，因此由\crosschapterref{10}{推论~\ref{cor:termwise-differentiation-uniform-derivative-series}}{函数项级数的逐项求导推论} 可逐项求导。各阶重复这一论证得到式~\eqref{eq:power-series-arbitrary-derivatives}；每个导函数又是连续项级数的内闭一致极限，故均连续。
\end{proof}

求导与积分保持收敛半径，但可能改变端点的敛散性。若原级数在某端点收敛，积分级数在该端点也收敛；若求导级数在某端点收敛，则原级数也在该端点收敛，因为可把原通项写成求导通项乘以该端点相对中心的距离再除以 $n$，使用阿贝尔判别法。反向推断一般不成立。例如从 $\sum_{n=0}^\infty x^n$ 连续两次从零积分，依次得到
\[
    \sum_{n=0}^\infty\frac{x^{n+1}}{n+1},\qquad
    \sum_{n=0}^\infty\frac{x^{n+2}}{(n+1)(n+2)},
\]
三者半径都为 $1$，收敛域却依次为 $(-1,1)$、$[-1,1)$ 和 $[-1,1]$。逐项求导公式首先是开区间内部的结论，不能仅凭原级数在端点收敛就延用到该点。

\begin{example}
\label{ex:power-series-logarithmic-endpoint-changes}
设 $p,q>0$，讨论
\[
    F(x)=\sum_{n=1}^\infty\frac{x^n}{n^p\ln^q(n+1)}
\]
及其从零逐项积分、逐项求导后所得级数的收敛域。
\end{example}

\begin{solve}
系数的 $n$ 次根趋于 $1$，故三个级数的半径均为 $1$。原级数在 $x=-1$ 处由交错级数判别法收敛；在 $x=1$ 处，由对数型正项级数的判别，$\sum1/[n^p\ln^q(n+1)]$ 当且仅当 $p>1$，或 $p=1$ 且 $q>1$ 时收敛。这里可用 $\ln(n+1)\sim\ln n$，再对 $1/(t^p\ln^qt)$ 使用积分判别法；$p=1$ 时换元 $u=\ln t$，$p\neq1$ 时用对数与幂的阶比较。因此原级数的收敛域为
\[
    \begin{cases}
        [-1,1],&p>1\text{，或 }p=1,\ q>1,\\
        [-1,1),&0<p<1\text{，或 }p=1,\ 0<q\leq1.
    \end{cases}
\]
积分级数为 $\sum x^{n+1}/[(n+1)n^p\ln^q(n+1)]$，因 $p>0$，在两个端点均绝对收敛，故收敛域为 $[-1,1]$。

求导级数为 $\sum x^{n-1}/[n^{p-1}\ln^q(n+1)]$。当 $p\geq1$ 时，端点 $-1$ 处的绝对值通项单调趋于零，因而交错收敛；当 $0<p<1$ 时，$n^{1-p}/\ln^q(n+1)\to+\infty$，两个端点的通项都不趋于零。端点 $1$ 的判别相当于把原级数的 $p$ 换成 $p-1$，故求导级数的收敛域为
\[
    \begin{cases}
        [-1,1],&p>2\text{，或 }p=2,\ q>1,\\
        [-1,1),&1\leq p<2\text{，或 }p=2,\ 0<q\leq1,\\
        (-1,1),&0<p<1.
    \end{cases}
\]
可见积分能够改善端点收敛性，而求导可能使原本收敛的端点消失。
\end{solve}

\begin{example}
\label{ex:power-series-summation-by-two-derivatives}
求 $\displaystyle f(x)=\sum_{n=1}^\infty\frac{x^{2n+1}}{n(2n+1)}$ 的和函数，并求 $\displaystyle\sum_{n=1}^\infty1/[n(2n+1)]$。
\end{example}

\begin{solve}
半径为 $1$，且两个端点均绝对收敛。在 $|x|<1$ 内连续两次逐项求导，得到
\[
    f'(x)=\sum_{n=1}^\infty\frac{x^{2n}}n,
    \qquad f''(x)=2\sum_{n=1}^\infty x^{2n-1}
       =\frac{2x}{1-x^2}.
\]
利用 $f'(0)=f(0)=0$，先后积分得
\[
    f'(x)=-\ln(1-x^2),\qquad
    f(x)=(1-x)\ln(1-x)+2x-(1+x)\ln(1+x).
\]
当 $x\to1-$ 时，用 $t\ln t\to0$ 及阿贝尔端点定理可得
\[
    \sum_{n=1}^\infty\frac1{n(2n+1)}=2-2\ln2.
\]
同理，$f(-1)=2\ln2-2$；闭区间上的和函数由上式作连续延拓得到。计算时先在开区间内部求导，再用已经核实的端点收敛性取极限，避免了直接在端点求导。
\end{solve}

\section{初等函数的幂级数展开}
\label{sec:power-series-expansions-of-elementary-functions}

\subsection{泰勒级数}
\label{subsec:taylor-series}

幂级数的性质回答了由系数研究和函数的问题；现在反过来，给定函数 $f$，希望找到系数，使幂级数的和恰好等于 $f$。逐项求导说明，如果这样的展开存在，其系数已经被 $f$ 在中心处的各阶导数完全确定。困难不在于写出一个形式上的级数，而在于证明它确实表示原函数。

\begin{definition}
\label{def:taylor-series-and-real-analyticity}
若 $f$ 在 $x_0$ 处具有任意阶导数，则幂级数
\begin{equation}
    \sum_{n=0}^\infty\frac{f^{(n)}(x_0)}{n!}(x-x_0)^n,
    \qquad f^{(0)}:=f,
    \label{eq:taylor-series-definition}
\end{equation}
称为 $f$ 在 $x_0$ 处的\emph{泰勒级数}（Taylor series）；$x_0=0$ 时称为\emph{麦克劳林级数}（Maclaurin series）。若存在 $\delta>0$，使该级数在 $|x-x_0|<\delta$ 内收敛且和等于 $f(x)$，则称这个等式为 $f$ 在 $x_0$ 附近的\emph{泰勒展式}（Taylor expansion），并称 $f$ 在 $x_0$ 处\emph{实解析}（real analytic）。若开区间 $I$ 的每一点都具有这样的局部表示，则称 $f$ 在 $I$ 上实解析。
\end{definition}

\begin{theorem}[幂级数展开的唯一性]
\label{thm:uniqueness-of-power-series-expansion}
若对 $|x-x_0|<R$（$R>0$）有 $f(x)=\sum a_n(x-x_0)^n$，则
\begin{equation}
    a_n=\frac{f^{(n)}(x_0)}{n!},\qquad n=0,1,2,\ldots.
    \label{eq:power-series-coefficients-by-derivatives}
\end{equation}
因此，同一中心处的幂级数展开若存在，必唯一，且必为该函数的泰勒展式。
\end{theorem}

\begin{proof}
由定理~\ref{thm:power-series-termwise-differentiation}，可在内部逐项求任意阶导数。令 $x=x_0$ 时，$n$ 阶导数级数中只有原第 $n$ 项留下，得到 $f^{(n)}(x_0)=n!a_n$；$n=0$ 时直接代入。
\end{proof}

\begin{example}
\label{ex:smooth-flat-function-not-real-analytic}
设
\[
    f(x)=\begin{cases}e^{-1/x^2},&x\neq0,\\0,&x=0.\end{cases}
\]
证明 $f\in C^\infty(\mathbb R)$，但它在零点附近不能展开成幂级数。
\end{example}

\begin{solve}
先注意，对每个非负整数 $m$，令 $u=1/x^2\to+\infty$，由指数增长快于幂增长可知
\[
    |x|^{-m}e^{-1/x^2}=u^{m/2}e^{-u}\longrightarrow0.
\]
在 $x\neq0$ 处反复求导，每一阶导数都具有 $P_n(1/x)e^{-1/x^2}$ 的形式，其中 $P_n$ 为多项式。这可由归纳法核实：若当前多项式为 $P(t)$，下一阶对应多项式为 $-t^2P'(t)+2t^3P(t)$。

归纳证明零点处各阶导数存在且等于零。已知第 $n$ 阶导数在零点为零、在非零点具有上述形式，则
\[
    \lim_{h\to0}\frac{f^{(n)}(h)-f^{(n)}(0)}h
       =\lim_{h\to0}\frac{P_n(1/h)}h e^{-1/h^2}=0.
\]
所以第 $n+1$ 阶导数在零点存在且为零；它在非零点的表达式也趋于零，故在零点连续。由 $n=0$ 起归纳，得到 $f\in C^\infty(\mathbb R)$ 及 $f^{(n)}(0)=0$。于是其麦克劳林级数恒为零，而任意非零 $x$ 都有 $f(x)>0$；由展开的唯一性，$f$ 不可能在零点的任何邻域内由幂级数表示。
\end{solve}

此例区分了三个层次：有任意阶导数，泰勒级数具有正收敛半径，泰勒级数的和等于原函数。它们不能仅凭形式上的系数计算互相替代；本例的泰勒级数甚至在整个实轴收敛，却仍不能在零点附近表示 $f$。

\subsection{初等函数的泰勒展式}
\label{subsec:taylor-expansions-of-elementary-functions}

\begin{theorem}[泰勒展开的余项判据]
\label{thm:taylor-expansion-remainder-criterion}
设 $f\in C^\infty((x_0-\delta,x_0+\delta))$，记
\[
    T_N(x)=\sum_{n=0}^N\frac{f^{(n)}(x_0)}{n!}(x-x_0)^n,
    \qquad R_N(x)=f(x)-T_N(x).
\]
则 $f$ 在这个区间内等于其泰勒级数，当且仅当对每个固定的 $x$ 都有 $R_N(x)\to0$。若进一步在某个集合 $E$ 上有 $\sup_{x\in E}|R_N(x)|\to0$，则泰勒部分和在 $E$ 上一致收敛于 $f$。
\end{theorem}

\begin{proof}
由 $T_N(x)=f(x)-R_N(x)$，逐点收敛与一致收敛的两个结论分别直接来自定义。应用时可使用\crosschapterref{5}{第~\ref{sec:taylor-formula} 节}{泰勒公式一节}的拉格朗日余项或\crosschapterref{7}{式~\eqref{eq:taylor-integral-remainder}}{积分余项公式} 中的积分余项估计 $R_N$，但不能在没有估计的情况下把有限阶泰勒公式直接改写为无穷级数。
\end{proof}

\begin{proposition}[指数、正弦与余弦的展开]
\label{prop:exponential-sine-cosine-power-expansions}
对每个 $x\in\mathbb R$，有
\begin{align}
    e^x&=\sum_{n=0}^\infty\frac{x^n}{n!},
        \label{eq:exponential-maclaurin-expansion}\\
    \sin x&=\sum_{n=0}^\infty
        \frac{(-1)^nx^{2n+1}}{(2n+1)!},
        \label{eq:sine-maclaurin-expansion}\\
    \cos x&=\sum_{n=0}^\infty
        \frac{(-1)^nx^{2n}}{(2n)!}.
        \label{eq:cosine-maclaurin-expansion}
\end{align}
三者的收敛半径均为 $+\infty$，并在每个有限闭区间上绝对一致收敛。
\end{proposition}

\begin{proof}
指数函数各阶导数仍为自身；正弦与余弦的各阶导数由四步循环得到，绝对值不超过 $1$，故泰勒系数分别如上。固定 $x$，拉格朗日余项满足
\[
    |R_N^{\exp}(x)|\leq e^{|x|}\frac{|x|^{N+1}}{(N+1)!},
    \qquad
    |R_N^{\sin}(x)|,\ |R_N^{\cos}(x)|
       \leq\frac{|x|^{N+1}}{(N+1)!}.
\]
这里上标只标明所属函数，不表示求导。对固定 $x$，右边的阶乘分母使其趋于零，由余项判据得到展开；在 $|x|\leq M$ 上以 $M$ 替代 $|x|$ 还得到一致控制。绝对收敛及收敛半径也可直接用比值法核实。
\end{proof}

\begin{theorem}[广义二项式展开]
\label{thm:generalized-binomial-series}
对 $\alpha\in\mathbb R$，定义\emph{广义二项式系数}（generalized binomial coefficient）为
\[
    \binom\alpha0=1,\qquad
    \binom\alpha n=\frac{\alpha(\alpha-1)\cdots(\alpha-n+1)}{n!}
    \quad(n\geq1).
\]
当 $|x|<1$ 时，成立
\begin{equation}
    (1+x)^\alpha=\sum_{n=0}^\infty\binom\alpha n x^n.
    \label{eq:generalized-binomial-expansion}
\end{equation}
若 $\alpha$ 是非负整数，右边为有限和，即通常的二项式定理，半径为 $+\infty$，并按多项式意义对所有实数成立。其余情形半径为 $1$，端点的结论如下：当 $\alpha>0$ 且非整数时，两端均绝对收敛，式~\eqref{eq:generalized-binomial-expansion} 可连续延拓到 $[-1,1]$；当 $-1<\alpha<0$ 时，只在右端点条件收敛，等式成立于 $(-1,1]$；当 $\alpha\leq-1$ 时，两端都发散，收敛域为 $(-1,1)$。
\end{theorem}

\begin{proof}
非负整数情形由有限二项式定理得到，以下设 $\alpha$ 不是非负整数，记 $c_n=\binom\alpha n$。相邻项满足
\[
    \left|\frac{c_{n+1}}{c_n}\right|
       =\frac{|\alpha-n|}{n+1}\longrightarrow1,
\]
故半径为 $1$。为了证明和确为 $(1+x)^\alpha$，使用积分余项：
\[
    R_N(x)=\frac{\alpha(\alpha-1)\cdots(\alpha-N)}{N!}
       \int_0^x(x-t)^N(1+t)^{\alpha-N-1}\,dt.
\]
固定 $0<|x|<1$，当 $t$ 位于 $0$ 与 $x$ 之间时，$1+t>0$，并有
\[
    \left|\frac{x-t}{1+t}\right|\leq|x|.
\]
当 $x>0$ 时可用 $0\leq x-t\leq x(1+t)$；当 $x<0$ 时可用 $0\leq t-x\leq(-x)(1+t)$，后一不等式等价于 $t(1+x)\leq0$。令
\[
    C_{\alpha,x}=\int_{\min(0,x)}^{\max(0,x)}(1+t)^{\alpha-1}\,dt<+\infty,
\]
于是
\[
    |R_N(x)|\leq (N+1)|c_{N+1}|\,|x|^N C_{\alpha,x}
       \longrightarrow0.
\]
最后的极限可由右侧数列相邻项的比值趋于 $|x|<1$ 得到；$x=0$ 时余项恒为零。这证明了内部的等式。

再判断端点。若 $\alpha>0$，对充分大的 $n$ 有
\[
    n\left(\frac{|c_n|}{|c_{n+1}|}-1\right)
       =\frac{n(\alpha+1)}{n-\alpha}
       \longrightarrow\alpha+1>1.
\]
由拉贝判别法，$\sum|c_n|$ 收敛，两个端点都绝对收敛。阿贝尔定理将内部等式延拓到端点，和分别为 $0$ 与 $2^\alpha$。

若 $-1<\alpha<0$，则 $c_n=(-1)^nb_n$，$b_n>0$，且
\[
    \frac{b_{n+1}}{b_n}=1-\frac{1+\alpha}{n+1}\in(0,1).
\]
由 $1-u\leq e^{-u}$ 及调和级数发散，乘积表达式说明 $b_n\to0$，故 $x=1$ 处交错收敛，再由阿贝尔定理知和为 $2^\alpha$。在 $x=-1$ 处，通项为 $b_n>0$；若该级数收敛，阿贝尔定理会推出 $(1+x)^\alpha$ 在 $x\to-1+$ 时有有限极限，与其趋于 $+\infty$ 矛盾，故发散，右端点也不是绝对收敛。

若 $\alpha\leq-1$，则 $b_n=|c_n|$ 满足 $b_0=1$、$b_{n+1}/b_n=(n-\alpha)/(n+1)\geq1$，故两端的通项都不趋于零。各情形至此全部确定。
\end{proof}

\begin{proposition}[反正弦与对数的展开]
\label{prop:arcsine-and-logarithm-power-expansions}
有
\begin{align}
    \arcsin x
       &=\sum_{n=0}^\infty
         \frac{(2n-1)!!}{(2n)!!}\frac{x^{2n+1}}{2n+1},
       &&-1\leq x\leq1,
       \label{eq:arcsine-maclaurin-expansion}\\
    \ln(1+x)
       &=\sum_{n=1}^\infty\frac{(-1)^{n+1}}n x^n,
       &&-1<x\leq1.
       \label{eq:logarithm-maclaurin-expansion}
\end{align}
其中 $m!!$ 表示同奇偶性的正整数连乘，约定 $(-1)!!=0!!=1$。这两个级数的收敛半径均为 $1$；反正弦级数在两个端点都绝对收敛，对数级数在右端点条件收敛、左端点发散。
\end{proposition}

\begin{proof}
先在二项式展开中取 $\alpha=-1/2$ 并将变量换为 $-x^2$，得到
\[
    \frac1{\sqrt{1-x^2}}
       =\sum_{n=0}^\infty\frac{(2n-1)!!}{(2n)!!}x^{2n},
    \qquad |x|<1.
\]
在 $0$ 与任一内部点 $x$ 之间逐项积分，由 $\arcsin0=0$ 得到式~\eqref{eq:arcsine-maclaurin-expansion} 的内部等式。记
\[
    d_n=\frac{(2n-1)!!}{(2n)!!(2n+1)}.
\]
有 $n(d_n/d_{n+1}-1)\to3/2>1$，拉贝判别法说明 $\sum d_n$ 收敛，故两个端点均绝对收敛，再由阿贝尔定理及 $\arcsin$ 的连续性得到端点等式。特别地，
\begin{equation}
    \frac\pi2=\sum_{n=0}^\infty
       \frac{(2n-1)!!}{(2n)!!(2n+1)}.
    \label{eq:pi-by-arcsine-series}
\end{equation}

另一方面，在 $|x|<1$ 内由几何级数有 $1/(1+x)=\sum_{n=0}^\infty(-1)^nx^n$，从 $0$ 到 $x$ 逐项积分即得对数展开。在 $x=1$ 处，级数交错收敛，由阿贝尔定理知其和为 $\ln2$；在 $x=-1$ 处，级数为 $-\sum_{n=1}^\infty1/n$，故发散。
\end{proof}

求较复杂函数的泰勒展式时，应先寻找已知展式之间的联系。变量平移与代换、逐项积分或求导以及绝对收敛级数的乘积，常比直接计算一般的高阶导数更有效；每次变换仍须说明收敛范围。若所得展开以另一点为中心，要先把自变量改写成该中心的增量。

\begin{example}
\label{ex:shifted-sine-composition-taylor-expansion}
求 $f(x)=\sin(x^2-2x)$ 在 $x=1$ 处的泰勒展式。
\end{example}

\begin{solve}
令 $h=x-1$，则 $x^2-2x=h^2-1$，利用正弦差角公式和已知展开得到
\[
    \sin(x^2-2x)
       =\cos1\sum_{n=0}^\infty
           \frac{(-1)^n(x-1)^{4n+2}}{(2n+1)!}
        -\sin1\sum_{n=0}^\infty
           \frac{(-1)^n(x-1)^{4n}}{(2n)!}.
\]
两个级数对所有实数 $x$ 都绝对收敛，其合并形式是以 $1$ 为中心的幂级数；由唯一性，它正是所求泰勒展式。
\end{solve}

\begin{example}
\label{ex:finite-geometric-denominator-expansion}
设整数 $k\geq2$，求 $\displaystyle1/(1+x+\cdots+x^{k-1})$ 的麦克劳林展式。
\end{example}

\begin{solve}
对 $|x|<1$，有限几何和公式及无穷几何级数给出
\[
    \frac1{1+x+\cdots+x^{k-1}}
       =\frac{1-x}{1-x^k}
       =\sum_{n=0}^\infty x^{kn}-\sum_{n=0}^\infty x^{kn+1}.
\]
因 $k\geq2$，两列幂次不重合，合并后的系数仅为 $0,1,-1$，且有无穷多个非零系数，所以系数根值的上极限为 $1$，半径为 $1$。在 $x=\pm1$ 处，合并后级数的通项均不趋于零，故收敛域为 $(-1,1)$。即使原有理函数在某个端点有定义，也不能据此推断其麦克劳林级数在该点收敛。
\end{solve}

\begin{proposition}[幂级数的乘积]
\label{prop:product-of-power-series}
设 $f(x)=\sum a_n(x-x_0)^n$ 与 $g(x)=\sum b_n(x-x_0)^n$ 的半径分别为 $R_1,R_2>0$，令 $r=\min(R_1,R_2)$。当 $|x-x_0|<r$ 时，
\begin{equation}
    f(x)g(x)=\sum_{n=0}^\infty c_n(x-x_0)^n,
    \qquad c_n=\sum_{j=0}^na_jb_{n-j}.
    \label{eq:cauchy-product-of-power-series}
\end{equation}
乘积级数的半径至少为 $r$，可能因系数抵消而更大。
\end{proposition}

\begin{proof}
固定内部点，两个数项级数都绝对收敛，故可使用柯西乘积并重排，按总幂次合并即得公式。对任意 $0<s<r$，还有
\[
    \sum_{n=0}^\infty|c_n|s^n
       \leq\left(\sum_{n=0}^\infty|a_n|s^n\right)
          \left(\sum_{n=0}^\infty|b_n|s^n\right)<+\infty,
\]
故新级数半径至少为 $r$。
\end{proof}

\begin{example}
\label{ex:square-logarithm-maclaurin-expansion}
求 $\ln^2(1+x)$ 的麦克劳林展式，并确定其收敛域。
\end{example}

\begin{solve}
当 $|x|<1$ 时，对式~\eqref{eq:logarithm-maclaurin-expansion} 作柯西乘积，第 $m$ 次幂（$m\geq2$）的系数为
\[
    (-1)^m\sum_{j=1}^{m-1}\frac1{j(m-j)}
       =\frac{2(-1)^m}{m}\sum_{j=1}^{m-1}\frac1j.
\]
记调和数 $H_n=\sum_{j=1}^n1/j$，得到
\begin{equation}
    \ln^2(1+x)=2\sum_{m=2}^\infty
        \frac{(-1)^mH_{m-1}}m x^m,
    \qquad -1<x\leq1.
    \label{eq:square-logarithm-maclaurin-expansion}
\end{equation}
先在 $|x|<1$ 内获得这个等式，再单独检查端点：$H_{m-1}/m\to0$，且 $H_m/(m+1)\leq H_{m-1}/m$（$m\geq2$），故 $x=1$ 时交错收敛，阿贝尔定理说明和为 $\ln^22$。在 $x=-1$ 时，各项为 $2H_{m-1}/m\geq2/m$，级数发散。因此收敛域为 $(-1,1]$，右端点为条件收敛。端点结论来自独立判别，不能仅由内部的乘积公式直接推出。
\end{solve}

\section{连续函数的多项式逼近}
\label{sec:polynomial-approximation-of-continuous-functions}

若函数在某点附近有幂级数展开，取部分和就能在收敛区间内部的闭子区间上一致逼近它。但是，这要求函数在相应区域具有很强的局部性质；例~\ref{ex:smooth-flat-function-not-real-analytic} 说明，甚至无限次可微也不足以保证泰勒展开。若只要求逼近误差一致趋于零，而允许每次重新选择多项式的全部系数，则可以把函数类扩大到任意有限闭区间上的连续函数。

\begin{definition}
\label{def:uniform-polynomial-approximation}
设 $f:I\to\mathbb R$，$I$ 为非空区间。若对每个 $\varepsilon>0$，都存在实系数多项式 $P$，使
\begin{equation}
    \|f-P\|_{\infty,I}:=\sup_{x\in I}|f(x)-P(x)|<\varepsilon,
    \label{eq:uniform-polynomial-approximation}
\end{equation}
则称 $f$ 在 $I$ 上\emph{可被多项式一致逼近}（uniformly approximable by polynomials），也简称可被多项式逼近。多项式的次数与系数允许依赖于 $\varepsilon$。
\end{definition}

这个定义等价于存在多项式序列 $\{P_n\}$ 在 $I$ 上一致收敛于 $f$：一个方向可对误差 $1/n$ 逐次选取 $P_n$，另一个方向直接由一致收敛的定义得到。由于多项式连续，连续函数一致极限的定理说明，连续性是可被多项式一致逼近的必要条件。以下定理证明，在有限闭区间上，这也是充分条件。

\begin{theorem}[魏尔斯特拉斯逼近定理]
\label{thm:weierstrass-polynomial-approximation}
设 $a<b$ 为有限实数，$f\in C([a,b])$，则 $f$ 在 $[a,b]$ 上可被多项式一致逼近，即存在多项式序列 $\{P_n\}$，使
\[
    \sup_{x\in[a,b]}|P_n(x)-f(x)|\longrightarrow0.
\]
这称为\emph{魏尔斯特拉斯逼近定理}（Weierstrass approximation theorem）。
\end{theorem}

为证明定理，先考察多项式一致逼近在有限线性组合与一致极限下的稳定性，再从绝对值函数构造分段线性函数的逼近。

\begin{lemma}[有限线性组合与一致极限]
\label{lem:polynomial-approximation-closure}
在固定的 $[a,b]$ 上，若 $f_1,\ldots,f_m$ 都可被多项式一致逼近，则对任意实数 $c_1,\ldots,c_m$，$\sum_{j=1}^mc_jf_j$ 也可被多项式一致逼近。若每个 $f_n$ 都可被多项式一致逼近，且 $f_n\rightrightarrows f$ 于 $[a,b]$，则 $f$ 也可被多项式一致逼近。
\end{lemma}

\begin{proof}
给定 $\varepsilon>0$，令 $C=\sum_{j=1}^m|c_j|$，对每个 $j$ 选多项式 $Q_j$，使 $\|f_j-Q_j\|_{\infty,[a,b]}<\varepsilon/(C+1)$。则 $\sum c_jQ_j$ 是多项式，且
\[
    \left\|\sum_{j=1}^mc_jf_j-\sum_{j=1}^mc_jQ_j\right\|_{\infty,[a,b]}
       \leq\sum_{j=1}^m|c_j|\|f_j-Q_j\|_{\infty,[a,b]}
       <\varepsilon.
\]
再证第二个结论。先取某个 $n$，使 $\|f-f_n\|_{\infty,[a,b]}<\varepsilon/2$，再选多项式 $P$，使 $\|f_n-P\|_{\infty,[a,b]}<\varepsilon/2$。三角不等式给出 $\|f-P\|_{\infty,[a,b]}<\varepsilon$。
\end{proof}

\begin{lemma}[绝对值函数的多项式逼近]
\label{lem:absolute-value-polynomial-approximation}
对每个固定的 $c\in\mathbb R$，存在多项式序列 $\{P_{n,c}\}$，在任意有限闭区间上一致收敛于 $|x-c|$。
\end{lemma}

\begin{proof}
先在 $[-1,1]$ 上逼近 $|x|$。广义二项式定理在 $\alpha=1/2$ 时给出
\[
    \sqrt{1+t}=\sum_{j=0}^\infty\binom{1/2}j t^j,
    \qquad -1\leq t\leq1.
\]
级数两端均绝对收敛，故由阿贝尔定理在整个 $[-1,1]$ 上一致收敛。记部分和为 $S_N(t)$，将 $t=x^2-1\in[-1,0]$ 代入，得到
\[
    S_N(x^2-1)\rightrightarrows\sqrt{1+x^2-1}=|x|,
    \qquad -1\leq x\leq1.
\]
$S_N(x^2-1)$ 是关于 $x$ 的多项式，所以 $|x|$ 在该区间可被多项式一致逼近。

对每个正整数 $n$，选多项式 $Q_n$，使 $\sup_{|t|\leq1}|Q_n(t)-|t||<1/n^2$，再令
\[
    P_{n,c}(x)=nQ_n\left(\frac{x-c}{n}\right).
\]
固定有限闭区间 $[A,B]$，当 $n\geq\max_{x\in[A,B]}|x-c|$ 时，所有代入点均属于 $[-1,1]$，因此
\[
    \sup_{x\in[A,B]}|P_{n,c}(x)-|x-c||<\frac1n\longrightarrow0.
\]
同一个序列由此在每个有限闭区间上一致收敛；这里没有声称在整个实轴上一致收敛。
\end{proof}

\begin{definition}
\label{def:continuous-piecewise-linear-function}
设 $g\in C([a,b])$。若存在分割 $a=t_0<t_1<\cdots<t_m=b$，使 $g$ 在每个 $[t_{i-1},t_i]$ 上都是一次或常数函数，则称 $g$ 为\emph{连续分段线性函数}（continuous piecewise linear function）。另记\emph{正部函数}（positive-part function）为
\[
    H(x)=\max\{x,0\}=\frac{x+|x|}{2},\qquad x\in\mathbb R.
\]
\end{definition}

\begin{lemma}[分段线性函数的表示与逼近]
\label{lem:piecewise-linear-hinge-representation}
设 $g$ 对分割 $a=t_0<\cdots<t_m=b$ 连续分段线性，各段斜率为
\[
    s_i=\frac{g(t_i)-g(t_{i-1})}{t_i-t_{i-1}},\qquad 1\leq i\leq m.
\]
则对 $x\in[a,b]$，有
\begin{equation}
    g(x)=g(a)+s_1H(x-a)
       +\sum_{j=1}^{m-1}(s_{j+1}-s_j)H(x-t_j).
    \label{eq:piecewise-linear-positive-part-representation}
\end{equation}
因此，每个连续分段线性函数都可在 $[a,b]$ 上被多项式一致逼近。
\end{lemma}

\begin{proof}
式~\eqref{eq:piecewise-linear-positive-part-representation} 右边是连续分段线性函数，在 $x=a$ 处等于 $g(a)$。对 $x\in(t_{i-1},t_i)$，只有分点在 $x$ 左侧的正部函数有斜率 $1$，所以右边在该段的斜率为
\[
    s_1+\sum_{j=1}^{i-1}(s_{j+1}-s_j)=s_i.
\]
从首段起逐段比较函数值与斜率，即知两边相等。又因 $H(x-c)=(x-c+|x-c|)/2$，由前一引理及有限线性组合的稳定性，每个 $H(x-c)$ 都可被多项式一致逼近，再对表示式使用同一稳定性即得结论。
\end{proof}

\begin{proof}[魏尔斯特拉斯逼近定理的证明]
由\crosschapterref{3}{定理~\ref{thm:heine-cantor-uniform-continuity}}{康托尔定理}，$f$ 在 $[a,b]$ 上一致连续。给定 $\varepsilon>0$，选 $\delta>0$，使 $|x-y|<\delta$ 时有 $|f(x)-f(y)|<\varepsilon/2$。取网格宽度小于 $\delta$ 的分割 $a=t_0<\cdots<t_m=b$，令 $g$ 为通过各点 $(t_i,f(t_i))$ 的连续分段线性插值函数。若 $x\in[t_{i-1},t_i]$，写成
\[
    x=(1-\theta)t_{i-1}+\theta t_i,\qquad
    g(x)=(1-\theta)f(t_{i-1})+\theta f(t_i),
    \qquad 0\leq\theta\leq1,
\]
便有
\[
    |g(x)-f(x)|
       \leq(1-\theta)|f(t_{i-1})-f(x)|
          +\theta|f(t_i)-f(x)|<\frac\varepsilon2.
\]
由于 $g-f$ 在紧区间上连续并取到最大值，上述逐点严格不等式也给出 $\|g-f\|_{\infty,[a,b]}<\varepsilon/2$。由引理~\ref{lem:piecewise-linear-hinge-representation}，可再选多项式 $P$，使 $\|P-g\|_{\infty,[a,b]}<\varepsilon/2$。于是
\[
    \|P-f\|_{\infty,[a,b]}
       \leq\|P-g\|_{\infty,[a,b]}+\|g-f\|_{\infty,[a,b]}
       <\varepsilon.
\]
对每个误差 $1/n$ 这样选取多项式，即得所需的多项式序列。等价地，也可先取一列网格宽度趋于零的分段线性插值函数，再使用引理~\ref{lem:polynomial-approximation-closure} 中关于一致极限的结论。
\end{proof}

这个证明把一般连续函数的逼近分为两步：由一致连续性控制线性插值误差，再把有限个绝对值函数的逼近组合成一个多项式。它不要求 $f$ 可导，也不要求这些多项式是某点的泰勒部分和；例如 $|x|$ 在零点不可导，仍可在每个有限闭区间上被多项式一致逼近。反过来，多项式一致收敛也不能保证导数收敛。区间条件同样重要，例如 $1/x$ 虽在 $(0,1)$ 连续，却不能在整个 $(0,1)$ 上被多项式一致逼近，因为每个多项式在 $[0,1]$ 上有界，而 $1/x$ 在零点附近无界。

\section*{习题十一}
\addcontentsline{toc}{section}{习题十一}
\fi

\selectchapter{12}
\ifcompilechapter
\chapter{傅里叶级数}
\label{chap:fourier-series}

幂级数用 $1,x,x^2,\ldots$ 的无穷线性组合表示函数，适合研究函数在某一点附近的性质。对于周期现象，更自然的基本函数是具有共同周期的正弦、余弦。本章研究\emph{傅里叶级数}（Fourier series）：先由三角函数的正交性确定展开系数，再考察级数在给定点的收敛、一致逼近与均方逼近。系数能由积分确定，并不意味着相应级数已经收敛，更不意味着其和处处等于原函数；这些问题需要分别证明。

\section{函数的傅里叶级数}
\label{sec:fourier-series-of-functions}

\subsection{基本三角函数系}
\label{subsec:basic-trigonometric-system}

函数系
\[
    1,\cos x,\sin x,\cos2x,\sin2x,\ldots,\cos nx,\sin nx,\ldots
\]
称为\emph{基本三角函数系}（basic trigonometric system），其中各函数均以 $2\pi$ 为周期。它们的有限线性组合称为\emph{三角多项式}（trigonometric polynomial）；形如
\[
    T_N(x)=\frac{\alpha_0}{2}
      +\sum_{n=1}^N(\alpha_n\cos nx+\beta_n\sin nx)
\]
的函数称为至多 $N$ 阶的三角多项式。常数项写成 $\alpha_0/2$，是为了使以后余弦系数的积分公式在 $n=0$ 时仍然适用。

对有限区间 $[a,b]$（$a<b$）上的实连续函数，定义\emph{内积}（inner product）$\langle f,g\rangle=\int_a^bf(x)g(x)\,dx$；若该积分为零，则称两函数\emph{正交}（orthogonal）。对于黎曼可积函数，同一积分也有意义，但 $\int f^2=0$ 不保证 $f$ 处处为零，因而要把零均方距离的函数视为同一个对象，才能得到通常的内积空间。以下正交计算只用到积分本身，不需要预先建立这一空间。

\begin{proposition}[基本三角函数系的正交关系]
\label{prop:trigonometric-orthogonality}
对任意 $c\in\mathbb R$ 及正整数 $m,n$，有
\begin{align}
    \int_c^{c+2\pi}\cos mx\cos nx\,dx
       &=\begin{cases}\pi,&m=n,\\0,&m\neq n,\end{cases}\notag\\
    \int_c^{c+2\pi}\sin mx\sin nx\,dx
       &=\begin{cases}\pi,&m=n,\\0,&m\neq n,\end{cases}\notag\\
    \int_c^{c+2\pi}\cos mx\sin nx\,dx&=0.
    \label{eq:trigonometric-orthogonality}
\end{align}
此外，$\int_c^{c+2\pi}1\,dx=2\pi$，而任意非恒定基本三角函数在该区间上的积分为零。
\end{proposition}

\begin{proof}
由周期性，只须在 $[-\pi,\pi]$ 上验证。使用积化和差公式，例如
\[
    \sin mx\sin nx
      =\frac12\bigl(\cos(m-n)x-\cos(m+n)x\bigr),
\]
当 $m\neq n$ 时，右边两项积分均为零；当 $m=n$ 时，右边为 $(1-\cos2nx)/2$，积分为 $\pi$。余弦与混合乘积同理，常数项的积分直接计算即可。
\end{proof}

正交性使不同频率之间的积分交叉项消失。特别地，若一个三角多项式恒为零，依次乘以 $1,\cos nx,\sin nx$ 并积分，就得到其全部系数为零，因此基本三角函数系线性无关。

\subsection{周期为 \texorpdfstring{$2\pi$}{2pi} 的函数的傅里叶级数}
\label{subsec:fourier-series-of-two-pi-periodic-functions}

设 $f$ 为实值 $2\pi$ 周期函数，在一个周期内黎曼可积。若它已有一致收敛的三角展开，则将展开式分别乘以 $1,\cos nx,\sin nx$，由逐项积分定理与正交关系就能求出各系数。这一计算提示我们：即使尚未知道能否展开，也可先用积分定义系数。

\begin{definition}
\label{def:fourier-coefficients-and-series}
函数 $f$ 的\emph{傅里叶系数}（Fourier coefficients）定义为
\begin{align}
    a_n&=\frac1\pi\int_{-\pi}^{\pi}f(x)\cos nx\,dx,
          &&n=0,1,2,\ldots,\notag\\
    b_n&=\frac1\pi\int_{-\pi}^{\pi}f(x)\sin nx\,dx,
          &&n=1,2,\ldots.
    \label{eq:fourier-coefficients-two-pi}
\end{align}
相应的傅里叶级数及其部分和记为
\begin{equation}
    f(x)\sim\frac{a_0}{2}
          +\sum_{n=1}^\infty(a_n\cos nx+b_n\sin nx),
    \qquad
    S_N(f;x)=\frac{a_0}{2}
          +\sum_{n=1}^N(a_n\cos nx+b_n\sin nx).
    \label{eq:fourier-series-and-partial-sums}
\end{equation}
记号 $\sim$ 只表示右边的系数由左边函数按式~\eqref{eq:fourier-coefficients-two-pi} 确定，不预设级数的敛散性或和函数。
\end{definition}

被积函数均为 $2\pi$ 周期函数，故系数积分可以在任意一个完整周期上计算。改变 $f$ 在有限个点的值，不改变其傅里叶系数；因此傅里叶级数不可能在这些点同时恢复各种不同的赋值。若 $f$ 在对称区间内为奇函数，则所有 $a_n=0$；若为偶函数，则所有 $b_n=0$。这里的奇偶性应针对给定区间及其周期延拓判断，端点的有限个赋值不影响计算。

\begin{theorem}[连续周期函数的傅里叶系数唯一性]
\label{thm:fourier-uniqueness-continuous}
若连续的 $2\pi$ 周期函数 $f$ 的全部傅里叶系数均为零，则 $f\equiv0$。因此，两个连续的 $2\pi$ 周期函数若傅里叶系数相同，它们就处处相等。
\end{theorem}

\begin{proof}
全部系数为零意味着对任意三角多项式 $T$，都有 $\int_{x_0-\pi}^{x_0+\pi}f(x)T(x)\,dx=0$。若 $f$ 不恒为零，必要时改用 $-f$，可设 $f(x_0)>0$。取 $0<\delta<\pi/2$ 与 $M>0$，使 $|x-x_0|\leq\delta$ 时 $f(x)\geq M$，并记 $B=\max|f|$。令
\[
    H(x)=1+\cos(x-x_0)-\cos\delta,
    \qquad q=1+\cos(\delta/2)-\cos\delta>1.
\]
在 $|x-x_0|\leq\delta/2$ 上有 $H\geq q$；在 $\delta\leq|x-x_0|\leq\pi$ 上有 $|H|\leq1$。$H^{2m}$ 是三角多项式，在 $|x-x_0|<\delta$ 上非负，故
\[
    0=\int_{x_0-\pi}^{x_0+\pi}f(x)H(x)^{2m}\,dx
       \geq M\delta q^{2m}-2\pi B.
\]
右边在 $m\to\infty$ 时趋于 $+\infty$，矛盾。第二个结论应用于两函数之差。
\end{proof}

这个唯一性结论并未断言连续函数的傅里叶部分和处处收敛；它只说明系数作为整体能够唯一确定连续周期函数。后文的费叶平均将给出从系数恢复该函数的一种一致逼近方法。

\begin{example}
\label{ex:fourier-sawtooth-and-period-choice}
分别求在 $[-\pi,\pi)$ 上取 $f(x)=x$、在 $[0,2\pi)$ 上取 $g(x)=x$ 的两个 $2\pi$ 周期函数的傅里叶级数。
\end{example}

\begin{solve}
$f$ 在 $(-\pi,\pi)$ 内为奇函数，故 $a_n=0$，分部积分得
\[
    b_n=\frac2\pi\int_0^\pi x\sin nx\,dx
        =\frac{2(-1)^{n+1}}n,
    \qquad
    f(x)\sim2\sum_{n=1}^\infty\frac{(-1)^{n+1}}n\sin nx.
\]
对 $g$ 在 $[0,2\pi]$ 上积分，则 $a_0=2\pi$、$a_n=0$（$n\geq1$）、$b_n=-2/n$，所以
\[
    g(x)\sim\pi-2\sum_{n=1}^\infty\frac{\sin nx}{n}.
\]
两种周期延拓的跳跃点位置不同，定义的是两个不同函数。另若在 $[0,2\pi)$ 上取 $h(x)=\pi-x$，其周期延拓为奇函数（跳跃点赋值除外），相应级数为 $2\sum_{n\geq1}\sin nx/n$；不能仅凭区间内的解析式判断延拓后的奇偶性。
\end{solve}

\begin{example}
\label{ex:fourier-square-on-shifted-period}
求在 $[0,2\pi)$ 上取值 $q(x)=x^2$ 的 $2\pi$ 周期函数的傅里叶级数。
\end{example}

\begin{solve}
在给定周期上连续两次分部积分，得到
\[
    a_0=\frac{8\pi^2}{3},\qquad
    a_n=\frac4{n^2},\qquad b_n=-\frac{4\pi}{n},\quad n\geq1.
\]
因此
\begin{equation}
    q(x)\sim\frac{4\pi^2}{3}
       +4\sum_{n=1}^\infty\frac{\cos nx}{n^2}
       -4\pi\sum_{n=1}^\infty\frac{\sin nx}{n}.
    \label{eq:fourier-square-shifted-period}
\end{equation}
虽然 $x^2$ 在实轴上为偶函数，但这里的周期函数由 $[0,2\pi)$ 上的取值规定，并非将 $[-\pi,\pi)$ 上的 $x^2$ 周期延拓，因此正弦系数不为零。
\end{solve}

\subsection{正弦级数与余弦级数}
\label{subsec:sine-and-cosine-series}

只在半个周期上给定函数时，需要先选择如何延拓。设 $f$ 在 $(0,\pi)$ 上有定义，并能通过端点赋值成为 $[0,\pi]$ 上的黎曼可积函数。将它在 $(-\pi,0)$ 上分别定义为 $-f(-x)$ 或 $f(-x)$，就得到\emph{奇延拓}（odd extension）或\emph{偶延拓}（even extension），再作 $2\pi$ 周期延拓。奇延拓在零点及周期拼接点可取零，偶延拓可取任意有限端点值，若有端点极限也可按该极限赋值；有限个点的选择均不影响系数。

\begin{definition}
\label{def:half-range-fourier-expansions}
$f$ 在 $(0,\pi)$ 上的\emph{正弦级数}（sine series）和\emph{余弦级数}（cosine series）分别为
\begin{align}
    f(x)&\sim\sum_{n=1}^\infty b_n\sin nx,
       & b_n&=\frac2\pi\int_0^\pi f(x)\sin nx\,dx,\notag\\
    f(x)&\sim\frac{a_0}{2}+\sum_{n=1}^\infty a_n\cos nx,
       & a_n&=\frac2\pi\int_0^\pi f(x)\cos nx\,dx\quad(n\geq0).
    \label{eq:half-range-fourier-coefficients}
\end{align}
它们分别是上述奇延拓、偶延拓的傅里叶级数，称为\emph{半区间展开}（half-range expansion）。
\end{definition}

\begin{example}
\label{ex:half-range-fourier-square}
求 $f(x)=x^2$（$0<x<\pi$）的正弦级数与余弦级数。
\end{example}

\begin{solve}
分部积分给出
\[
    b_n=\frac{2\pi(-1)^{n+1}}n
           +\frac{4\bigl((-1)^n-1\bigr)}{\pi n^3},
    \qquad a_0=\frac{2\pi^2}{3},\qquad
    a_n=\frac{4(-1)^n}{n^2}\quad(n\geq1).
\]
因此两种展开分别为
\begin{align}
    x^2&\sim2\pi\sum_{n=1}^\infty\frac{(-1)^{n+1}}n\sin nx
       -\frac8\pi\sum_{n=1}^\infty
                    \frac{\sin((2n-1)x)}{(2n-1)^3},
       \label{eq:half-range-square-sine}\\
    x^2&\sim\frac{\pi^2}{3}
          +4\sum_{n=1}^\infty\frac{(-1)^n}{n^2}\cos nx.
       \label{eq:half-range-square-cosine}
\end{align}
两个级数对应不同的全周期函数，故在区间端点的行为也可能不同；其确切收敛值将在点态收敛判别之后确定。
\end{solve}

\subsection{周期为 \texorpdfstring{$2T$}{2T} 的函数的傅里叶级数}
\label{subsec:fourier-series-of-two-t-periodic-functions}

设 $T>0$，$f$ 为 $2T$ 周期函数且在一个周期内可积。作变量代换 $t=\pi x/T$，则 $F(t)=f(Tt/\pi)$ 为 $2\pi$ 周期函数，把其傅里叶级数换回 $x$ 就得到
\begin{equation}
    f(x)\sim\frac{a_0}{2}
       +\sum_{n=1}^\infty\left(
          a_n\cos\frac{n\pi x}{T}
          +b_n\sin\frac{n\pi x}{T}\right),
    \label{eq:fourier-series-general-period}
\end{equation}
其中
\begin{equation}
    a_n=\frac1T\int_c^{c+2T}f(x)\cos\frac{n\pi x}{T}\,dx\quad(n\geq0),
    \qquad
    b_n=\frac1T\int_c^{c+2T}f(x)\sin\frac{n\pi x}{T}\,dx\quad(n\geq1).
    \label{eq:fourier-coefficients-general-period}
\end{equation}
$c$ 任意。基本角频率为 $\pi/T$，不必是 $f$ 的最小正周期所对应的角频率；只要给定的 $2T$ 确为一个周期，公式就成立。相应地，在 $(0,T)$ 上的正弦、余弦半区间展开，把式~\eqref{eq:half-range-fourier-coefficients} 中的 $2/\pi$ 改为 $2/T$，并把 $nx$ 改为 $n\pi x/T$。

\begin{example}
\label{ex:fourier-x-cos-x-period-pi}
将 $f(x)=x\cos x$（$-\pi/2\leq x<\pi/2$）作 $\pi$ 周期延拓，求其傅里叶级数。
\end{example}

\begin{solve}
此时 $T=\pi/2$，延拓后的函数为奇函数，故只含 $\sin2nx$ 项。由积化和差及分部积分，
\[
    b_n=\frac4\pi\int_0^{\pi/2}x\cos x\sin2nx\,dx
       =\frac{16n(-1)^{n+1}}{\pi(4n^2-1)^2},
\]
因此
\begin{equation}
    f(x)\sim\frac{16}{\pi}\sum_{n=1}^\infty
             \frac{(-1)^{n+1}n}{(4n^2-1)^2}\sin2nx.
    \label{eq:fourier-x-cos-x-period-pi}
\end{equation}
在端点 $x=\pm\pi/2$，左右延拓的函数值都趋于零，因此该周期函数在拼接点连续。后文的收敛结论通过相同的变量代换适用于这一展开。
\end{solve}

\section{傅里叶级数的敛散性}
\label{sec:convergence-of-fourier-series}

从本节起，除黎曼可积函数外，也允许在一个周期内有有限个瑕点的函数：它在避开瑕点的闭子区间上黎曼可积，且 $\int|f|$ 的相应瑕积分收敛。下文简称这两类情形为\emph{一个周期内绝对可积}，所有涉及瑕点的积分按绝对收敛的广义积分理解。该条件保证傅里叶系数存在，但尚不保证傅里叶级数在每点收敛。

\subsection{狄利克雷积分}
\label{subsec:dirichlet-integral}

\begin{definition}
\label{def:dirichlet-kernel}
对非负整数 $N$，定义\emph{狄利克雷核}（Dirichlet kernel）
\begin{equation}
    D_N(t)=1+2\sum_{n=1}^N\cos nt
       =\frac{\sin((N+\tfrac12)t)}{\sin(t/2)}
       \quad(t\notin2\pi\mathbb Z).
    \label{eq:dirichlet-kernel}
\end{equation}
在 $t\in2\pi\mathbb Z$ 处按连续延拓取 $D_N(t)=2N+1$。本章把归一化常数 $1/(2\pi)$ 放在积分公式外。
\end{definition}

利用 $2\sin(t/2)\cos nt=\sin((n+1/2)t)-\sin((n-1/2)t)$ 求和，就得到式~\eqref{eq:dirichlet-kernel}。将系数积分代入有限和，利用余弦差角公式及周期性，可得\emph{狄利克雷积分表示}（Dirichlet integral representation）
\begin{align}
    S_N(f;x)&=\frac1{2\pi}\int_{-\pi}^{\pi}f(x+t)D_N(t)\,dt\notag\\
       &=\frac1{2\pi}\int_0^\pi
         \bigl(f(x+t)+f(x-t)\bigr)
         \frac{\sin((N+\tfrac12)t)}{\sin(t/2)}\,dt.
    \label{eq:dirichlet-integral-representation}
\end{align}
这里交换的仅是有限和与积分。核为偶函数，并满足
\begin{equation}
    \int_{-\pi}^{\pi}D_N(t)\,dt=2\pi,
    \qquad \int_0^\pi D_N(t)\,dt=\pi.
    \label{eq:dirichlet-kernel-mass}
\end{equation}
因此，对于希望得到的极限 $s\in\mathbb R$，若记
\[
    \phi_{x,s}(t)=f(x+t)+f(x-t)-2s,
\]
就有
\begin{equation}
    S_N(f;x)-s=\frac1\pi\int_0^\pi
       \frac{\phi_{x,s}(t)}{2\sin(t/2)}
                    \sin((N+\tfrac12)t)\,dt.
    \label{eq:fourier-pointwise-error-integral}
\end{equation}
傅里叶级数的点态收敛，由此转化为一个振荡积分趋于零的问题。

\begin{lemma}[黎曼--勒贝格引理]
\label{lem:riemann-lebesgue-finite-integral}
若 $h$ 在有限区间 $[a,b]$ 上黎曼可积，或有有限个瑕点且绝对可积，则
\begin{equation}
    \lim_{\lambda\to+\infty}\int_a^bh(t)\sin\lambda t\,dt=0,
    \qquad
    \lim_{\lambda\to+\infty}\int_a^bh(t)\cos\lambda t\,dt=0.
    \label{eq:riemann-lebesgue-finite-integral}
\end{equation}
\end{lemma}

\begin{proof}
先设 $h$ 黎曼可积。由达布可积准则，对任意 $\varepsilon>0$，可选一个固定分割及其上的阶梯函数 $p$，使 $\int_a^b|h-p|<\varepsilon/2$。若 $p=c_j$ 于第 $j$ 个子区间，则
\[
    \left|\int_a^bp(t)\sin\lambda t\,dt\right|
        \leq\frac2\lambda\sum_j|c_j|\longrightarrow0.
\]
而 $|\int(h-p)\sin\lambda t|\leq\int|h-p|<\varepsilon/2$，即得正弦的结论，余弦同理。若存在瑕点，先取它们的有限个小邻域，使 $\int|h|$ 在这些邻域上的总和小于 $\varepsilon/2$；余下部分是有限个闭区间，对各段使用已证结论。瑕点邻域上的误差估计与 $\lambda$ 无关，故两种极限仍成立。
\end{proof}

把引理用于 $h=f$，得到 $a_n\to0$、$b_n\to0$。这只是傅里叶系数的必要性质，不能单独推出相应三角级数收敛。

\begin{theorem}[黎曼局部化定理]
\label{thm:riemann-localization-fourier}
设 $f,g$ 均为一个周期内绝对可积的 $2\pi$ 周期函数。若存在 $\delta>0$，使它们在 $0<|t|<\delta$ 时满足 $f(x_0+t)=g(x_0+t)$，则
\[
    S_N(f;x_0)-S_N(g;x_0)\longrightarrow0.
\]
因而两级数在 $x_0$ 处同时收敛或同时发散；若收敛，其和相同。
\end{theorem}

\begin{proof}
取 $0<\delta<\pi$。在部分和之差的积分表示中，$|t|<\delta$ 的积分为零，单点 $t=0$ 的取值不影响积分；而在 $\delta\leq|t|\leq\pi$ 上，$1/\sin(t/2)$ 有界。将
\[
    \frac{f(x_0+t)-g(x_0+t)}{\sin(t/2)}
\]
分别限制在这两个闭区间，使用黎曼--勒贝格引理即可。
\end{proof}

系数取决于一个完整周期的取值，某一点的收敛却只受该点任意小邻域内的行为控制。进一步，函数
\[
    r(t)=\frac1{2\sin(t/2)}-\frac1t\quad(t>0)
\]
在 $t\to0$ 时趋于零，可补成 $[0,\pi]$ 上的连续函数。对误差积分中由 $r(t)$ 产生的项应用黎曼--勒贝格引理，再去掉远离零点的一段，便得到：对任意固定的 $0<\delta<\pi$，
\begin{equation}
    S_N(f;x_0)-s
       =\frac1\pi\int_0^\delta
           \frac{\phi_{x_0,s}(t)}t
                  \sin((N+\tfrac12)t)\,dt+o(1).
    \label{eq:localized-fourier-error}
\end{equation}
这里 $o(1)$ 随 $N\to\infty$ 趋于零；即使 $\phi(t)/t$ 尚未绝对可积，对每个固定 $N$，$\sin((N+1/2)t)/t$ 在零点仍有有限极限，因此该积分存在。

\begin{example}
\label{ex:dirichlet-sine-integral-value}
利用狄利克雷核计算 $\displaystyle\int_0^{+\infty}\frac{\sin u}{u}\,du$。
\end{example}

\begin{solve}
由无穷积分的狄利克雷判别法，该积分收敛。在式~\eqref{eq:dirichlet-kernel-mass} 中，用 $1/t$ 代替 $1/(2\sin(t/2))$，二者之差乘以振荡因子后的积分由黎曼--勒贝格引理趋于零，故
\[
    1=\lim_{N\to\infty}\frac2\pi\int_0^\pi
                  \frac{\sin((N+\tfrac12)t)}t\,dt
      =\frac2\pi\int_0^{+\infty}\frac{\sin u}{u}\,du.
\]
最后一步作 $u=(N+1/2)t$，再使用已经证明的广义积分收敛性，得到积分值为 $\pi/2$。
\end{solve}

\subsection{傅里叶级数的收敛判别法}
\label{subsec:convergence-tests-for-fourier-series}

式~\eqref{eq:localized-fourier-error} 给出了收敛到 $s$ 的充要条件，但具体应用还需把它转化为便于检查的局部条件。以下始终假定 $f$ 为一个周期内绝对可积的 $2\pi$ 周期函数；在周期拼接点，单侧极限按周期延拓后的函数计算。

\begin{theorem}[傅里叶级数的狄尼判别法]
\label{thm:dini-fourier-convergence}
若存在 $s\in\mathbb R$ 与 $\delta>0$，使
\begin{equation}
    \int_0^\delta
       \frac{|f(x_0+t)+f(x_0-t)-2s|}{t}\,dt<+\infty,
    \label{eq:dini-fourier-condition}
\end{equation}
则 $S_N(f;x_0)\to s$。
\end{theorem}

\begin{proof}
对 $\phi_{x_0,s}(t)/t$ 在 $(0,\delta]$ 上使用黎曼--勒贝格引理，再用式~\eqref{eq:localized-fourier-error}。
\end{proof}

这里的狄尼判别法（Dini test）是振荡积分的局部可积条件，与前章关于单调函数序列一致收敛的狄尼定理是不同结论。若左右极限存在，通常取 $s=(f(x_0+0)+f(x_0-0))/2$；此时只要两项 $|f(x_0\pm t)-f(x_0\pm0)|/t$ 分别可积，就满足式~\eqref{eq:dini-fourier-condition}。

\begin{corollary}[可导情形与分段可微情形]
\label{cor:piecewise-differentiable-fourier-convergence}
若 $f$ 在 $x_0$ 可导，则 $S_N(f;x_0)\to f(x_0)$。更一般地，若有限的左右极限存在，且下列两个极限均为有限实数：
\[
    \lim_{t\downarrow0}\frac{f(x_0+t)-f(x_0+0)}t,
    \qquad
    \lim_{t\downarrow0}\frac{f(x_0-t)-f(x_0-0)}{-t},
\]
则
\begin{equation}
    S_N(f;x_0)\longrightarrow
          \frac{f(x_0+0)+f(x_0-0)}2.
    \label{eq:fourier-midpoint-limit}
\end{equation}
特别地，一个周期内分段可微且各分段端点具有上述有限单侧导数的函数，傅里叶级数处处收敛到左右极限的平均值。
\end{corollary}

\begin{proof}
两个差商在零点附近有界，故取式~\eqref{eq:fourier-midpoint-limit} 右边为 $s$ 时，$\phi_{x_0,s}(t)/t$ 在零点附近有界，满足狄尼条件。可导情形中两个单侧函数值均为 $f(x_0)$。
\end{proof}

这里\emph{分段可微}（piecewise differentiable）指存在有限分割，使函数在各开子区间可导，分点处具有有限左右极限及上述以单侧极限为基准的单侧导数；分点的实际赋值可以不同于相应极限。常用的分段 $C^1$ 函数满足此条件，但无需要求整个周期内处处连续。

\begin{corollary}[李普西茨判别法]
\label{cor:holder-fourier-convergence}
若存在 $L>0$、$\delta>0$、$\alpha>0$，使 $|t|<\delta$ 时
\[
    |f(x_0+t)-f(x_0)|\leq L|t|^\alpha,
\]
即 $f$ 在 $x_0$ 满足\emph{赫尔德条件}（H\"older condition），则傅里叶级数在该点收敛到 $f(x_0)$。$\alpha=1$ 时称为\emph{李普西茨条件}（Lipschitz condition）。
\end{corollary}

\begin{proof}
取 $s=f(x_0)$，则 $|\phi_{x_0,s}(t)|/t\leq2Lt^{\alpha-1}$，其在零点附近可积。左右两侧若分别以各自的极限满足正指数的赫尔德估计，同样得到式~\eqref{eq:fourier-midpoint-limit}。
\end{proof}

\begin{theorem}[狄利克雷判别法]
\label{thm:dirichlet-fourier-convergence}
若 $f$ 在 $(x_0-\delta_0,x_0)$ 和 $(x_0,x_0+\delta_0)$ 内分别单调，且 $f(x_0-0),f(x_0+0)$ 均为有限实数，则式~\eqref{eq:fourier-midpoint-limit} 成立。因而在一个周期内具有有限个单调段、各分点具有有限单侧极限的函数，其傅里叶级数处处收敛到左右极限的平均值。
\end{theorem}

\begin{proof}
先考虑右侧，令 $h(t)=f(x_0+t)-f(x_0+0)$。$h$ 单调且 $h(t)\to0$（$t\downarrow0$）。广义积分 $\int_0^{+\infty}\sin u/u\,du$ 收敛，所以存在常数 $C$，使任意 $0\leq A<B$ 均满足
\[
    \left|\int_A^B\frac{\sin u}{u}\,du\right|\leq C.
\]
给定 $\varepsilon>0$，选 $0<\delta<\delta_0$，使 $|h(t)|<\varepsilon$ 于 $(0,\delta]$。对任意 $0<\eta<\delta$，使用定积分第二中值定理的端点形式，$\int_\eta^\delta h(t)\sin(\lambda t)/t\,dt$ 可写成 $h(\eta)$、$h(\delta)$ 与两个子区间上 $\sin(\lambda t)/t$ 的积分之线性组合；换元 $u=\lambda t$ 后，各积分绝对值至多为 $C$。故令 $\eta\downarrow0$，得到
\[
    \left|\int_0^\delta h(t)\frac{\sin\lambda t}{t}\,dt\right|
        \leq2C\varepsilon.
\]
在 $[\delta,\delta_0]$ 上，$h(t)/t$ 可积，其振荡积分由黎曼--勒贝格引理趋于零。因此先令 $\lambda\to\infty$，再令 $\varepsilon\downarrow0$，右侧积分趋于零。左侧对 $f(x_0-t)-f(x_0-0)$ 作同样论证，再在式~\eqref{eq:localized-fourier-error} 中取 $s$ 为左右极限的平均值。
\end{proof}

\emph[caution]{在连续点，左右极限平均值等于函数值；在跳跃点，收敛值由两侧极限决定，与该点的单独赋值无关。}函数在一点连续本身不足以套用以上判别法；应检查可微性、赫尔德条件、狄尼条件或局部单调性中的适用条件。

\begin{example}
\label{ex:fourier-endpoint-sums}
确定前面各个多项式型例子的傅里叶级数在区间内部与端点处的和。
\end{example}

\begin{solve}
这些周期函数都分段可微。在例~\ref{ex:fourier-sawtooth-and-period-choice} 中，$f(x)=x$ 的级数在 $(-\pi,\pi)$ 内等于 $x$，在 $x=\pm\pi$ 处等于零；$g(x)=x$ 的级数在 $(0,2\pi)$ 内等于 $x$，在 $x=0,2\pi$ 处等于 $\pi$。相应地，
\begin{equation}
    2\sum_{n=1}^\infty\frac{\sin nx}{n}
      =\begin{cases}\pi-x,&0<x<2\pi,\\0,&x=0,2\pi.\end{cases}
    \label{eq:fourier-pi-minus-x-sum}
\end{equation}
式~\eqref{eq:fourier-square-shifted-period} 在 $(0,2\pi)$ 内的和为 $x^2$，在 $0,2\pi$ 处为 $(0+4\pi^2)/2=2\pi^2$。式~\eqref{eq:half-range-square-sine} 在 $0\leq x<\pi$ 的和为 $x^2$，在 $x=\pi$ 处为零；式~\eqref{eq:half-range-square-cosine} 则在整个 $[0,\pi]$ 上等于 $x^2$。奇延拓在 $\pi$ 处产生跳跃，而偶延拓在该点连续，这解释了两种端点结果的差异。式~\eqref{eq:fourier-x-cos-x-period-pi} 则在 $[-\pi/2,\pi/2]$ 上均等于 $x\cos x$。
\end{solve}

\begin{example}
\label{ex:fourier-cos-noninteger-frequency}
设 $0<\alpha<1$，将 $f(x)=\cos\alpha x$（$-\pi\leq x\leq\pi$）作 $2\pi$ 周期延拓，求其傅里叶级数及和函数。
\end{example}

\begin{solve}
由偶性，$b_n=0$。积化和差给出
\[
    a_0=\frac{2\sin\alpha\pi}{\alpha\pi},\qquad
    a_n=\frac{2\alpha(-1)^n\sin\alpha\pi}
                       {\pi(\alpha^2-n^2)}\quad(n\geq1).
\]
该周期延拓连续且分段可微，故
\begin{equation}
    \cos\alpha x=\frac{\sin\alpha\pi}{\alpha\pi}
       +\frac{2\alpha\sin\alpha\pi}{\pi}
          \sum_{n=1}^\infty
                    \frac{(-1)^n\cos nx}{\alpha^2-n^2},
    \qquad -\pi\leq x\leq\pi.
    \label{eq:fourier-cos-noninteger-frequency}
\end{equation}
区间外的和是左边在这一周期上的周期延拓，而不一定等于实轴上原来的 $\cos\alpha x$。端点左右函数值相接，因此端点处也成立等式。
\end{solve}

\section{傅里叶级数的其他收敛性}
\label{sec:other-convergence-properties-of-fourier-series}

点态收敛考察每个固定点的部分和极限，一致收敛要求对全体点给出统一的误差控制；均方收敛则用误差平方的积分衡量逼近。后两者研究的对象和条件不同。本节先通过部分和的算术平均构造一致逼近，再利用正交性证明部分和的均方最佳逼近，最后讨论在何种条件下原傅里叶级数本身一致收敛及允许逐项运算。

\subsection{连续函数的三角多项式一致逼近}
\label{subsec:uniform-trigonometric-approximation}

若数列 $s_N\to s$，则它的算术平均 $(s_0+\cdots+s_N)/(N+1)$ 也趋于 $s$，反向推论一般不成立。平均能够减弱部分和的振荡，因此即使尚不能证明 $S_N(f;x)$ 收敛，也可以研究其平均。

\begin{definition}
\label{def:fejer-means-and-kernel}
$f$ 的傅里叶部分和的\emph{费叶平均}（Fej\'er mean）定义为
\begin{equation}
    \sigma_N(f;x)=\frac1{N+1}\sum_{k=0}^NS_k(f;x)
       =\frac{a_0}{2}+\sum_{n=1}^N
           \left(1-\frac n{N+1}\right)
                        (a_n\cos nx+b_n\sin nx).
    \label{eq:fejer-means}
\end{equation}
相应的\emph{费叶核}（Fej\'er kernel）为
\begin{equation}
    F_N(t)=\frac1{N+1}\sum_{k=0}^ND_k(t)
       =\frac1{N+1}
             \left(\frac{\sin((N+1)t/2)}{\sin(t/2)}\right)^2,
    \label{eq:fejer-kernel}
\end{equation}
在 $2\pi\mathbb Z$ 上取连续延拓值 $N+1$。
\end{definition}

求和公式可由式~\eqref{eq:dirichlet-kernel} 再次用积化和差得到。费叶核具有三个关键性质：非负性 $F_N\geq0$，归一化 $\int_{-\pi}^{\pi}F_N=2\pi$，以及对任意固定的 $0<\delta<\pi$，
\begin{equation}
    \sup_{\delta\leq|t|\leq\pi}F_N(t)
       \leq\frac1{(N+1)\sin^2(\delta/2)}\longrightarrow0.
    \label{eq:fejer-kernel-concentration}
\end{equation}
归一化来自各 $D_k$ 的积分，最后的估计来自分子不超过 $1$。它们说明核的总质量保持不变，而远离零点的质量逐渐消失。对有限个部分和的积分表示求平均，有
\begin{equation}
    \sigma_N(f;x)=\frac1{2\pi}\int_{-\pi}^{\pi}
                                     f(x+t)F_N(t)\,dt.
    \label{eq:fejer-convolution}
\end{equation}

\begin{theorem}[费叶定理与魏尔斯特拉斯第二逼近定理]
\label{thm:fejer-uniform-approximation}
若 $f$ 为连续的 $2\pi$ 周期函数，则 $\sigma_N(f;\cdot)$ 在整个实轴上一致收敛到 $f$。因此，连续周期函数可由三角多项式一致逼近。
\end{theorem}

\begin{proof}
连续周期函数在实轴上有界且一致连续。记 $M=\max|f|$，给定 $\varepsilon>0$，取 $0<\delta<\pi$，使 $|t|<\delta$ 时对一切 $x$ 有 $|f(x+t)-f(x)|<\varepsilon/2$。由核的非负性和归一化，
\[
    |\sigma_N(f;x)-f(x)|
       \leq\frac1{2\pi}\int_{-\pi}^{\pi}
                  |f(x+t)-f(x)|F_N(t)\,dt.
\]
在 $|t|<\delta$ 上，积分的贡献至多为 $\varepsilon/2$；在其余部分，由式~\eqref{eq:fejer-kernel-concentration}，贡献至多为
\[
    \frac{2M}{(N+1)\sin^2(\delta/2)}.
\]
充分大的 $N$ 使此项小于 $\varepsilon/2$，且这个选择与 $x$ 无关。$\sigma_N$ 是至多 $N$ 阶的三角多项式，故一致逼近结论同时成立。
\end{proof}

此定理中的逼近多项式是 $\sigma_N$，它对第 $n$ 个频率的系数乘以 $1-n/(N+1)$，并非原来的傅里叶部分和 $S_N$。证明依赖费叶核的非负性；狄利克雷核会变号，不能照搬上述估计来证明连续函数的傅里叶部分和一致收敛。另一方面，若 $S_N(f;x_0)$ 在某一点已经收敛到 $\ell$，则其算术平均也趋于 $\ell$；费叶定理又给出 $\sigma_N(f;x_0)\to f(x_0)$，所以 $\ell=f(x_0)$。由此得到教材中的结论：连续周期函数的傅里叶级数若处处收敛，其和必处处为原函数。

费叶定理也给出定理~\ref{thm:fourier-uniqueness-continuous} 的另一证明：全部系数为零时，每个 $\sigma_N$ 均为零，而它们一致趋于 $f$。任意周期 $2T$ 的情形只须作 $t=\pi x/T$ 的变量代换。

\subsection{傅里叶级数的均方收敛}
\label{subsec:mean-square-convergence-of-fourier-series}

\begin{definition}
\label{def:mean-square-convergence}
设 $f$ 在有限区间 $[a,b]$ 上黎曼可积，或仅有有限个瑕点、在避开瑕点的闭子区间可积，且相应广义积分 $\int_a^bf(x)^2\,dx<+\infty$，则称其\emph{平方可积}（square integrable）。对平方可积的实函数 $f_n,f$，若
\[
    \int_a^b|f_n(x)-f(x)|^2\,dx\longrightarrow0,
\]
则称 $f_n$ \emph{均方收敛}（mean-square convergence）到 $f$。在 $[-\pi,\pi]$ 上，记
\begin{equation}
    \langle f,g\rangle_2=\frac1\pi\int_{-\pi}^{\pi}f(x)g(x)\,dx,
    \qquad \|f\|_2=\left(\frac1\pi\int_{-\pi}^{\pi}|f(x)|^2\,dx\right)^{1/2}.
    \label{eq:normalized-square-integral}
\end{equation}
后一个量称为归一化均方范数；严格地说，在未把零均方距离的函数等同之前，它是半范数。
\end{definition}

本节的平方可积首先按上述黎曼积分与有限瑕点的框架理解；更一般的 $L^2$ 理论需要勒贝格积分，此处不预先使用它。黎曼可积函数有界，故平方可积；平方可积函数绝对可积，因为 $2|f|\leq1+f^2$。平方可积比绝对可积要求更强，例如 $x^{-1/2}$ 在 $(0,1]$ 上绝对可积，而其平方的积分发散。

\begin{proposition}[积分型柯西不等式与三角不等式]
\label{prop:square-integral-inequalities}
若 $f,g$ 在 $[a,b]$ 上平方可积，则 $fg$ 绝对可积，$f+g$ 平方可积，且
\begin{align}
    \left|\int_a^bf(x)g(x)\,dx\right|
       &\leq\left(\int_a^bf^2\right)^{1/2}
              \left(\int_a^bg^2\right)^{1/2},
       \label{eq:integral-cauchy-schwarz-fourier}\\
    \left(\int_a^b|f+g|^2\right)^{1/2}
       &\leq\left(\int_a^bf^2\right)^{1/2}
                 +\left(\int_a^bg^2\right)^{1/2}.
       \label{eq:integral-triangle-fourier}
\end{align}
\end{proposition}

\begin{proof}
由 $2|fg|\leq f^2+g^2$ 与 $(f+g)^2\leq2f^2+2g^2$，先得到相应可积性。对每个实数 $t$，$\int_a^b(f+tg)^2\geq0$；若 $\int g^2>0$，此二次式的判别式非正，得到第一个不等式。若 $\int g^2=0$，其一次项也必须为零，结论仍成立。展开 $\int(f+g)^2$，再对交叉项使用第一个不等式，即得第二个不等式。
\end{proof}

在有限区间上，一致收敛推出均方收敛，因为 $\int|f_n-f|^2\leq(b-a)\|f_n-f\|_\infty^2$。反向推论不成立，而且均方极限不决定个别点的取值：在一个点取 $1$、其余点取 $0$ 的函数与零函数均方距离为零。更具体地，令 $h_n(x)=\max\{1-nx,0\}$ 于 $[0,1]$，则 $\int_0^1h_n^2=1/(3n)\to0$，但 $h_n(0)=1$，既不一致趋于零，也不在每点趋于零。

\begin{theorem}[傅里叶部分和的最佳均方逼近]
\label{thm:fourier-best-square-approximation}
设 $f$ 在 $[-\pi,\pi]$ 上平方可积，其傅里叶系数为 $a_n,b_n$。对任意至多 $N$ 阶的三角多项式
\[
    T_N(x)=\frac{\alpha_0}{2}
             +\sum_{n=1}^N(\alpha_n\cos nx+\beta_n\sin nx),
\]
有
\begin{equation}
    \|f-T_N\|_2^2=\|f-S_N(f)\|_2^2
            +\frac{(\alpha_0-a_0)^2}{2}
            +\sum_{n=1}^N\bigl((\alpha_n-a_n)^2+(\beta_n-b_n)^2\bigr).
    \label{eq:fourier-square-error-decomposition}
\end{equation}
因此 $S_N(f)$ 是这一类多项式中唯一使均方误差最小者。
\end{theorem}

\begin{proof}
由系数定义，$f-S_N(f)$ 与 $1,\cos nx,\sin nx$（$1\leq n\leq N$）均正交，故也与 $S_N(f)-T_N$ 正交。把
\[
    f-T_N=(f-S_N(f))+(S_N(f)-T_N)
\]
平方积分，交叉项为零，便得到两项范数平方的和；再对第二项应用正交关系即得式~\eqref{eq:fourier-square-error-decomposition}。等号达到最小值当且仅当每一个系数差均为零，即 $T_N=S_N(f)$。
\end{proof}

特别地，直接展开部分和的误差得到
\begin{equation}
    \|f-S_N(f)\|_2^2
       =\|f\|_2^2-\frac{a_0^2}{2}
                          -\sum_{n=1}^N(a_n^2+b_n^2)\geq0.
    \label{eq:fourier-residual-energy}
\end{equation}
令 $N\to\infty$，便得到\emph{贝塞尔不等式}（Bessel's inequality）
\begin{equation}
    \frac{a_0^2}{2}+\sum_{n=1}^\infty(a_n^2+b_n^2)\leq\|f\|_2^2.
    \label{eq:bessel-fourier-inequality}
\end{equation}
它还说明均方误差随 $N$ 单调不增。不过，误差非负且递减只保证存在极限，要证明极限为零，还需三角多项式对函数类具有充分的逼近能力。

\begin{lemma}[平方可积函数的三角多项式逼近]
\label{lem:trigonometric-density-square-integral}
对本节定义的任意平方可积函数 $f$ 及任意 $\varepsilon>0$，存在三角多项式 $T$，使 $\|f-T\|_2<\varepsilon$。
\end{lemma}

\begin{proof}
先设 $f$ 黎曼可积，$|f|\leq M$；$M=0$ 时无需证明。取分割，使 $\sum_j\omega_j(f)\Delta x_j$ 任意小，并在各段取 $f$ 的一个值组成阶梯函数 $p$，则
\[
    \int_{-\pi}^{\pi}|f-p|^2
       \leq2M\sum_j\omega_j(f)\Delta x_j.
\]
阶梯函数只有有限个跳跃点。在这些点附近的小区间内用线段连接左右常值，在两端的小区间内再使端点值都变为零，便得到连续且 $g(-\pi)=g(\pi)$ 的函数 $g$，并保持 $|g|\leq M$。$p,g$ 只在总长度可任意小的集合内不同，且差的绝对值至多为 $2M$，所以 $\|p-g\|_2$ 可任意小。将 $g$ 作连续周期延拓，费叶定理给出三角多项式 $T$ 使 $\|g-T\|_\infty$ 任意小，故 $\|g-T\|_2\leq\sqrt2\|g-T\|_\infty$ 也任意小。用三角不等式合并三项误差，就得到结论。

若 $f$ 有有限个瑕点，由 $\int f^2$ 收敛，可取这些点的小邻域，使其上 $f^2$ 的积分任意小。把这些邻域内的函数值改为零，得到一个有界黎曼可积函数 $f_0$，并使 $\|f-f_0\|_2$ 任意小。对 $f_0$ 用前面的结论，再次使用三角不等式即可。
\end{proof}

\begin{theorem}[傅里叶级数的均方收敛与帕塞瓦尔等式]
\label{thm:fourier-mean-square-and-parseval}
若 $f$ 在 $[-\pi,\pi]$ 上平方可积，则 $S_N(f)$ 均方收敛到 $f$，并且成立\emph{帕塞瓦尔等式}（Parseval's identity）
\begin{equation}
    \frac1\pi\int_{-\pi}^{\pi}f(x)^2\,dx
       =\frac{a_0^2}{2}+\sum_{n=1}^\infty(a_n^2+b_n^2).
    \label{eq:parseval-fourier-identity}
\end{equation}
\end{theorem}

\begin{proof}
给定 $\varepsilon>0$，由前一引理取三角多项式 $T$，使 $\|f-T\|_2<\varepsilon$。设其阶数为 $m$，则对一切 $N\geq m$，最佳逼近定理给出 $\|f-S_N(f)\|_2\leq\|f-T\|_2<\varepsilon$，从而均方收敛。再在式~\eqref{eq:fourier-residual-energy} 中令 $N\to\infty$ 即得帕塞瓦尔等式。
\end{proof}

等式说明函数的平方积分可以分解为各频率系数的平方和，这常称为能量分解。它在没有点态收敛保证时仍成立，因此证明中没有把一个未经核实的无穷展开直接平方并逐项积分。反过来，并非任意三角级数都是某个平方可积函数的傅里叶级数；例如 $\sum_{n\geq1}\cos nx/\sqrt n$ 的系数平方和发散，不能来自这样的函数。

\begin{corollary}[广义帕塞瓦尔等式]
\label{cor:generalized-parseval-fourier}
若 $f,g$ 平方可积，傅里叶系数分别为 $a_n,b_n$ 和 $\alpha_n,\beta_n$，则
\begin{equation}
    \frac1\pi\int_{-\pi}^{\pi}f(x)g(x)\,dx
       =\frac{a_0\alpha_0}{2}
            +\sum_{n=1}^\infty(a_n\alpha_n+b_n\beta_n).
    \label{eq:generalized-parseval-fourier}
\end{equation}
右边级数绝对收敛。
\end{corollary}

\begin{proof}
由贝塞尔不等式及数列的柯西不等式，$\sum|a_n\alpha_n|$ 和 $\sum|b_n\beta_n|$ 收敛。分别对 $f+g$、$f-g$ 使用帕塞瓦尔等式，二式相减后除以 $4$，即得结论。
\end{proof}

\begin{example}
\label{ex:parseval-sums-second-fourth-powers}
用帕塞瓦尔等式计算 $\sum_{n=1}^\infty n^{-2}$ 和 $\sum_{n=1}^\infty n^{-4}$。
\end{example}

\begin{solve}
对 $[-\pi,\pi]$ 上的 $f(x)=x$，由 $b_n=2(-1)^{n+1}/n$，得到
\[
    \frac1\pi\int_{-\pi}^{\pi}x^2\,dx
       =4\sum_{n=1}^\infty\frac1{n^2},
    \qquad \sum_{n=1}^\infty\frac1{n^2}=\frac{\pi^2}{6}.
\]
对同一区间上的 $g(x)=x^2$，式~\eqref{eq:half-range-square-cosine} 给出 $a_0=2\pi^2/3$、$a_n=4(-1)^n/n^2$，故
\[
    \frac{2\pi^4}{5}=\frac{2\pi^4}{9}
                   +16\sum_{n=1}^\infty\frac1{n^4},
    \qquad \sum_{n=1}^\infty\frac1{n^4}=\frac{\pi^4}{90}.
\]
两次运用只要求相应函数平方可积，不依赖它们的傅里叶级数先有一致收敛性。
\end{solve}

\subsection{傅里叶级数的一致收敛性}
\label{subsec:uniform-convergence-of-fourier-series}

光滑性使高频系数衰减得更快，而系数的绝对可和性又能给出一致收敛。联系这两者的关键是对系数积分作分部积分，再用导函数的贝塞尔不等式。

\begin{theorem}[一个一致收敛的充分条件]
\label{thm:fourier-absolute-uniform-convergence}
设 $f$ 为连续的 $2\pi$ 周期函数，在一个周期内有限分段可导，且导函数可在分点处补值成为黎曼可积函数，则
\begin{equation}
    \sum_{n=1}^\infty(|a_n|+|b_n|)<+\infty,
    \label{eq:fourier-absolute-coefficient-summability}
\end{equation}
其傅里叶级数在整个实轴上绝对且一致收敛到 $f$。特别地，教材中 $f$ 在整个周期上可导、$f'$ 黎曼可积的条件满足本定理；连续的周期分段 $C^1$ 函数也满足本定理。
\end{theorem}

\begin{proof}
设 $f'$ 的傅里叶系数为 $A_n,B_n$。在各段上分部积分，因 $f$ 连续，内部分点的边界项相消；周期端点满足 $f(-\pi)=f(\pi)$，所以
\begin{equation}
    A_0=0,\qquad A_n=nb_n,\qquad B_n=-na_n\quad(n\geq1).
    \label{eq:fourier-derivative-coefficients}
\end{equation}
这里可积导数的牛顿--莱布尼茨公式保证各段的分部积分成立。由于 $f'$ 黎曼可积，从而平方可积，贝塞尔不等式给出
\[
    \sum_{n=1}^\infty n^2(a_n^2+b_n^2)
       =\sum_{n=1}^\infty(A_n^2+B_n^2)\leq\|f'\|_2^2.
\]
对有限和使用柯西不等式，再令上限趋于无穷，得到
\[
    \sum_{n=1}^\infty(|a_n|+|b_n|)
       \leq\left(\sum_{n=1}^\infty n^2(a_n^2+b_n^2)\right)^{1/2}
          \left(2\sum_{n=1}^\infty\frac1{n^2}\right)^{1/2}<+\infty.
\]
由 $M$ 判别法，级数在实轴上绝对且一致收敛。其连续和函数必为 $f$，这可由费叶定理关于已收敛部分和的结论得到。
\end{proof}

周期拼接处的连续性是这里的实质条件。例如，$x$ 在 $(-\pi,\pi)$ 内光滑，但它的周期延拓有跳跃，系数仅为 $1/n$ 量级；其部分和都是连续函数，不可能在整个周期上一致收敛到跳跃的和函数。另一方面，式~\eqref{eq:fourier-derivative-coefficients} 只是导函数的系数公式，不能仅凭形式微分就断言导数级数收敛。

\begin{theorem}[傅里叶级数的逐项求导]
\label{thm:fourier-termwise-differentiation}
设 $f$ 为 $2\pi$ 周期函数，$f\in C^1(\mathbb R)$，且在一个周期内二阶可导、$f''$ 黎曼可积。则 $f$ 的傅里叶级数可逐项求导，并有
\begin{equation}
    f'(x)=\sum_{n=1}^\infty
                       (nb_n\cos nx-na_n\sin nx),
    \qquad x\in\mathbb R.
    \label{eq:fourier-termwise-differentiation}
\end{equation}
原级数与上述导数级数都在实轴上一致收敛；特别地，$2\pi$ 周期的 $C^2$ 函数满足这些条件。
\end{theorem}

\begin{proof}
$f'$ 连续且为 $2\pi$ 周期函数，对它使用前一定理，得到其傅里叶级数一致收敛到 $f'$。由式~\eqref{eq:fourier-derivative-coefficients}，这个级数正是原级数形式微分后的级数。原级数由前一定理收敛到 $f$，故在任一有限闭区间上，由\crosschapterref{10}{推论~\ref{cor:termwise-differentiation-uniform-derivative-series}}{函数项级数的逐项求导推论} 可逐项求导，得到式~\eqref{eq:fourier-termwise-differentiation}。导数级数本身的实轴一致收敛由前述对 $f'$ 的应用保证。
\end{proof}

逐项积分需要的条件弱得多：原函数在一个周期内黎曼可积即可，原傅里叶级数本身无需先有逐点或一致收敛性。

\begin{theorem}[傅里叶级数的逐项积分]
\label{thm:fourier-termwise-integration}
设 $f$ 为在一个周期内黎曼可积的 $2\pi$ 周期函数，其傅里叶系数为 $a_n,b_n$。则对 $0\leq x\leq2\pi$，
\begin{equation}
    \int_0^xf(t)\,dt=\frac{a_0x}{2}
       +\sum_{n=1}^\infty
          \left(\frac{a_n}{n}\sin nx
                  +\frac{b_n}{n}(1-\cos nx)\right).
    \label{eq:fourier-termwise-integration}
\end{equation}
右边级数绝对且一致收敛。更一般地，对任意实数 $u,v$，
\begin{equation}
    \int_u^vf(t)\,dt=\frac{a_0(v-u)}2
       +\sum_{n=1}^\infty\left(
          \frac{a_n}{n}(\sin nv-\sin nu)
           +\frac{b_n}{n}(\cos nu-\cos nv)\right).
    \label{eq:fourier-integral-arbitrary-endpoints}
\end{equation}
\end{theorem}

\begin{proof}
因 $f$ 黎曼可积，所以平方可积。固定 $x\in[0,2\pi]$，在 $[0,2\pi]$ 上定义 $g_x(t)$ 为 $(0,x)$ 的示性函数，并作周期延拓；端点赋值任意。其系数为
\[
    \alpha_0=\frac x\pi,\qquad
    \alpha_n=\frac{\sin nx}{\pi n},\qquad
    \beta_n=\frac{1-\cos nx}{\pi n}.
\]
广义帕塞瓦尔等式的积分可移到 $[0,2\pi]$ 上，对 $f,g_x$ 应用该式，再乘以 $\pi$，就得到式~\eqref{eq:fourier-termwise-integration}。由贝塞尔不等式与柯西不等式，
\[
    \sum_{n=1}^\infty\frac{|a_n|+2|b_n|}{n}
       \leq\left(\sum_{n=1}^\infty(a_n^2+b_n^2)\right)^{1/2}
           \left(5\sum_{n=1}^\infty\frac1{n^2}\right)^{1/2}<+\infty,
\]
故积分后的级数绝对且一致收敛。式~\eqref{eq:fourier-termwise-integration} 两边从 $x$ 到 $x+2\pi$ 都增加 $\pi a_0$，因而等式通过周期性延伸到所有实数 $x$；取 $x=v$ 与 $x=u$ 的差即得一般端点公式，$u>v$ 时按定积分的方向约定理解。
\end{proof}

这一证明先由均方理论确定积分结果，再证明积分后级数的一致收敛，没有把一般函数项级数的逐项积分定理用在一个尚未证明一致收敛的原级数上。常数项积分为 $a_0x/2$，除非 $a_0=0$，原函数的积分通常不再是周期函数。

\begin{example}
\label{ex:fourier-integrated-sawtooth}
由 $\pi-x$ 在 $(0,2\pi)$ 上的正弦展开，求 $x^2/2-\pi x$ 的傅里叶级数及和函数。
\end{example}

\begin{solve}
由式~\eqref{eq:fourier-pi-minus-x-sum} 及逐项积分定理，
\[
    \pi x-\frac{x^2}{2}
       =2\sum_{n=1}^\infty\frac{1-\cos nx}{n^2},
       \qquad 0\leq x\leq2\pi.
\]
使用 $\sum n^{-2}=\pi^2/6$，得到
\begin{equation}
    \frac{x^2}{2}-\pi x
       =-\frac{\pi^2}{3}
           +2\sum_{n=1}^\infty\frac{\cos nx}{n^2},
       \qquad 0\leq x\leq2\pi.
    \label{eq:fourier-integrated-sawtooth}
\end{equation}
右边绝对且一致收敛，可逐项积分验证它的傅里叶系数恰为常数项 $-\pi^2/3$ 与余弦系数 $2/n^2$，故这是所求傅里叶级数。左边在 $0,2\pi$ 处均为零，其周期延拓连续，右边的和正是这一连续周期函数。原来的 $\pi-x$ 的周期延拓有跳跃，积分后则具有更好的连续性与系数衰减。
\end{solve}

\section*{习题十二}
\addcontentsline{toc}{section}{习题十二}
\fi

\selectchapter{13}
\ifcompilechapter
\chapter{多元函数的极限和连续}
\label{chap:multivariable-limits-and-continuity}

% 依据教材《数学分析》第十三章分章原文 PDF 第 1--34 页正文整理，不含章末习题。

一个量往往同时依赖若干个变量，例如温度可依赖空间位置与时间。研究这类函数，需要把实直线上的极限与连续理论推广到 $\mathbb R^n$。推广的关键不是逐个增加自变量，而是用\emph{距离}统一描述自变量的接近程度：在高维空间中，一个点可以从不同方向、沿不同曲线趋近另一个点，极限必须同时控制所有允许的趋近方式。

本章先建立欧氏空间中点列、点集和紧性的基础，再讨论实值与向量值函数的极限，最后把连续函数的局部性质推广为紧集或道路连通集上的整体性质。以下 $n,m$ 均为正整数；一般将 $x=(x_1,\ldots,x_n)$ 看作点或向量，用 $k$ 标记点列的项，用 $i$ 标记坐标，以免混淆维数与序号。对于二元、三元函数，仍沿用 $(x,y)$、$(x,y,z)$ 表示自变量。

\section{欧氏空间 \texorpdfstring{$\mathbb{R}^n$}{R^n}}
\label{sec:euclidean-space}

\subsection{欧氏空间 \texorpdfstring{$\mathbb{R}^n$}{R^n}}
\label{subsec:euclidean-space}

集合 $\mathbb R^n=\{(x_1,\ldots,x_n):x_i\in\mathbb R\}$ 配备按坐标进行的加法与数乘
\[
    x+y=(x_1+y_1,\ldots,x_n+y_n),\qquad
    \alpha x=(\alpha x_1,\ldots,\alpha x_n),\quad \alpha\in\mathbb R,
\]
构成 $n$ 维实向量空间。零向量为 $0=(0,\ldots,0)$，单位坐标向量为 $e_i=(0,\ldots,0,1,0,\ldots,0)$，其中第 $i$ 个坐标为 $1$，因而 $x=\sum_{i=1}^n x_i e_i$。在线性运算之外，定义\emph{内积}（inner product）
\[
    x\cdot y=\sum_{i=1}^n x_i y_i.
\]
它满足对称性、对每个变量的线性及正定性：$x\cdot y=y\cdot x$，$(\alpha x+\beta y)\cdot z=\alpha(x\cdot z)+\beta(y\cdot z)$，且 $x\cdot x\geq0$，等号当且仅当 $x=0$。赋予这一内积后的空间称为\emph{欧几里得空间}（Euclidean space），简称欧氏空间。

\begin{definition}
\label{def:euclidean-norm-and-distance}
向量 $x\in\mathbb R^n$ 的\emph{欧氏范数}（Euclidean norm），也称模，定义为
\[
    \|x\|=\sqrt{x\cdot x}=\left(\sum_{i=1}^n x_i^2\right)^{1/2}.
\]
两点 $x,y$ 的\emph{距离}（distance）定义为 $d(x,y)=\|x-y\|$。以下也用 $|x|$、$|x-y|$ 分别表示向量的模与两点间距离；当对象为实数时，竖线仍表示绝对值。
\end{definition}

\crosschapterref{5}{例~\ref{ex:holder-inequality-from-concavity}}{赫尔德不等式例题} 中的柯西--施瓦茨不等式给出 $|x\cdot y|\leq|x||y|$。由此
\[
    |x+y|^2=|x|^2+2x\cdot y+|y|^2\leq(|x|+|y|)^2,
\]
开方得到\emph{三角不等式} $|x+y|\leq|x|+|y|$。同时有 $|\alpha x|=|\alpha||x|$ 及反三角不等式 $\bigl||x|-|y|\bigr|\leq|x-y|$。因此距离非负且仅在 $x=y$ 时为零，对称，并满足 $|x-z|\leq|x-y|+|y-z|$。对于非零向量，夹角 $\theta\in[0,\pi]$ 可由 $\cos\theta=(x\cdot y)/(|x||y|)$ 定义；在二维和三维中，这与通常的几何夹角一致。特别地，非零向量也可以满足 $x\cdot y=0$，这时称它们\emph{正交}（orthogonal）。

欧氏范数与各坐标之间有以下反复使用的估计：
\begin{equation}
    \max_{1\leq i\leq n}|x_i|
    \leq |x|\leq\sum_{i=1}^n|x_i|
    \leq\sqrt n\,|x|,
    \qquad
    |x|\leq\sqrt n\max_{1\leq i\leq n}|x_i|.
    \label{eq:euclidean-coordinate-norm-bounds}
\end{equation}
其中最后一个不等式来自 $\sum x_i^2\leq n\max_i|x_i|^2$，而 $\sum|x_i|\leq\sqrt n\,|x|$ 来自柯西--施瓦茨不等式。这些估计说明，在固定的有限维空间中，控制所有坐标与控制距离可以相互转化。

\begin{example}
\label{ex:matrix-frobenius-vector-bound}
设 $A=(a_{ji})$ 为 $m\times n$ 实矩阵，定义其\emph{弗罗贝尼乌斯范数}（Frobenius norm）
\[
    \|A\|_{\mathrm F}=\left(\sum_{j=1}^m\sum_{i=1}^n a_{ji}^2\right)^{1/2}.
\]
证明对任意列向量 $x\in\mathbb R^n$，都有 $|Ax|\leq\|A\|_{\mathrm F}|x|$。
\end{example}

\begin{solve}
分别对各行使用柯西--施瓦茨不等式，得
\[
    |Ax|^2=\sum_{j=1}^m\left(\sum_{i=1}^n a_{ji}x_i\right)^2
    \leq\sum_{j=1}^m\left(\sum_{i=1}^n a_{ji}^2\right)
                            \left(\sum_{i=1}^n x_i^2\right)
    =\|A\|_{\mathrm F}^2|x|^2.
\]
两边开方即得结论。这里的范数明确指各矩阵元平方和的平方根，估计对任意 $m,n$ 均成立。
\end{solve}

\subsection{点列极限}
\label{subsec:limits-of-point-sequences}

\begin{definition}
\label{def:euclidean-neighborhoods}
给定 $a=(a_1,\ldots,a_n)\in\mathbb R^n$ 和 $\delta>0$，称
\[
    U(a,\delta)=\{x:|x-a|<\delta\},\qquad
    U_0(a,\delta)=U(a,\delta)\setminus\{a\}
\]
为 $a$ 的\emph{球形邻域}（ball neighborhood）及\emph{去心邻域}（punctured neighborhood）；称
\[
    N(a,\delta)=\{x:|x_i-a_i|<\delta,\ i=1,\ldots,n\},\qquad
    N_0(a,\delta)=N(a,\delta)\setminus\{a\}
\]
为 $a$ 的\emph{方形邻域}（cubical neighborhood）及方形去心邻域。
\end{definition}

由式~\eqref{eq:euclidean-coordinate-norm-bounds}，有
\begin{equation}
    U(a,\delta)\subseteq N(a,\delta)\subseteq U(a,\sqrt n\,\delta).
    \label{eq:ball-cube-neighborhood-inclusions}
\end{equation}
因此，在用任意小邻域定义的概念中，两类邻域可以互换。在 $\mathbb R^2$ 中，它们分别是开圆盘与开正方形；在 $\mathbb R^3$ 中，分别是开球与开立方体。这里“方形”在高维时也指各边平行于坐标轴的开立方体。

\begin{definition}
\label{def:convergent-euclidean-sequence}
设 $\{x_k\}\subseteq\mathbb R^n$。若存在 $a\in\mathbb R^n$，使得
\[
    \forall\varepsilon>0\ \exists K\in\mathbb N\ \forall k>K,
    \qquad |x_k-a|<\varepsilon,
\]
则称点列 $\{x_k\}$ \emph{收敛}（convergent）于 $a$，记为 $x_k\to a$ 或 $\lim_{k\to\infty}x_k=a$；若不存在这样的 $a$，则称其发散。等价地，$x_k\to a$ 就是实数列 $|x_k-a|\to0$。
\end{definition}

\begin{theorem}[按分量判定收敛]
\label{thm:euclidean-sequence-componentwise-convergence}
设 $x_k=(x_{k,1},\ldots,x_{k,n})$，$a=(a_1,\ldots,a_n)$，则 $x_k\to a$ 当且仅当每个坐标数列 $x_{k,i}\to a_i$。
\end{theorem}

\begin{proof}
若 $x_k\to a$，由 $|x_{k,i}-a_i|\leq|x_k-a|$ 得到各坐标收敛。反之，给定 $\varepsilon>0$，对每个 $i$ 取 $K_i$，使 $k>K_i$ 时 $|x_{k,i}-a_i|<\varepsilon/\sqrt n$。因维数有限，可取 $K=\max_i K_i$，从而 $k>K$ 时
\[
    |x_k-a|=\left(\sum_{i=1}^n|x_{k,i}-a_i|^2\right)^{1/2}<\varepsilon.
    \qedhere
\]
\end{proof}

集合 $E\subseteq\mathbb R^n$ 称为\emph{有界集}（bounded set），是指存在 $M>0$ 使 $|x|\leq M$ 对所有 $x\in E$ 成立；点列有界是指其各项组成的集合有界。按坐标判别可知，点列收敛蕴含有界，但有界并不蕴含收敛，例如 $x_k=(-1)^k e_1$。

\begin{proposition}
\label{prop:euclidean-sequence-limit-rules}
收敛点列的极限唯一，任意子列与原点列有相同极限。若 $x_k\to a$、$y_k\to b$，且 $\alpha,\beta\in\mathbb R$，则
\[
    \alpha x_k+\beta y_k\to\alpha a+\beta b,\qquad
    x_k\cdot y_k\to a\cdot b,\qquad |x_k|\to|a|.
\]
\end{proposition}

\begin{proof}
若同时有 $x_k\to a,b$，则 $|a-b|\leq|a-x_k|+|x_k-b|\to0$，故 $a=b$。收敛点列的尾项落在 $U(a,1)$ 中，再合并有限个首项即可得有界性。子列结论直接来自极限定义；线性运算与内积极限由按分量判别及有限次实数运算得到，范数极限由反三角不等式得到。
\end{proof}

\begin{example}
\label{ex:four-dimensional-point-sequence-limit}
求点列
\[
    x_k=\left(\frac{\sqrt k\sin k}{k},\ k^4e^{-k^2},
              \cos k\ln\left(1+\frac1k\right),\ k\tan\frac1k\right)
    \in\mathbb R^4
\]
的极限。
\end{example}

\begin{solve}
第一、三坐标的绝对值分别不超过 $k^{-1/2}$、$\ln(1+1/k)$，均趋于零；指数衰减快于固定次幂，故第二坐标也趋于零；第四坐标由 $\tan t/t\to1$ 得到极限 $1$。由定理~\ref{thm:euclidean-sequence-componentwise-convergence}，$x_k\to(0,0,0,1)$。向量非零与内积非零是不同的条件，例如两个恒定点列 $e_1,e_2$ 都有非零极限，其内积却恒为零。
\end{solve}

\subsection{聚点}
\label{subsec:accumulation-points-in-euclidean-space}

\begin{definition}
\label{def:euclidean-accumulation-and-isolated-points}
设 $E\subseteq\mathbb R^n$。若对任意 $\delta>0$ 都有 $U_0(a,\delta)\cap E\neq\varnothing$，则称 $a$ 为 $E$ 的\emph{聚点}或\emph{极限点}（accumulation point, limit point）。全体聚点组成的集合记为 $E'$，称为\emph{导集}（derived set）。若 $a\in E$ 而 $a\notin E'$，则称 $a$ 是 $E$ 的\emph{孤立点}（isolated point）；这时存在 $\delta>0$ 使 $U(a,\delta)\cap E=\{a\}$。
\end{definition}

聚点未必属于原集合，孤立点则必须属于原集合。聚点的每个邻域中实际含有无穷多个 $E$ 中的点：若某个去心邻域只含有限个这样的点，便可再缩小半径，避开这有限个点，产生与聚点定义矛盾的空去心邻域。由式~\eqref{eq:ball-cube-neighborhood-inclusions}，也可用方形去心邻域定义聚点。

\begin{theorem}[聚点的点列刻画]
\label{thm:euclidean-accumulation-sequence-criterion}
$a$ 是 $E\subseteq\mathbb R^n$ 的聚点，当且仅当存在 $E\setminus\{a\}$ 中两两不同的点 $x_k$，使 $x_k\to a$。
\end{theorem}

\begin{proof}
若 $a\in E'$，可从 $U_0(a,1/k)\cap E$ 中逐次选取 $x_k$，每次避开先前已选的有限个点；上述无限性保证选择可行，且 $|x_k-a|<1/k$。反之，若这样的点列存在，则 $a$ 的任意去心邻域都包含充分靠后的 $x_k$，故 $a\in E'$。
\end{proof}

若 $B\subseteq A\subseteq\mathbb R^n$，且 $A$ 中每个点的每个邻域都与 $B$ 相交，则称 $B$ 在 $A$ 中\emph{稠密}（dense）。这里不要求交点异于邻域中心，因此“稠密”与“每点都是聚点”在定义上应加以区分。

\begin{example}
\label{ex:rational-coordinate-points-dense}
求有理坐标点集 $\mathbb Q^n=\{(x_1,\ldots,x_n):x_i\in\mathbb Q\}$ 的导集，并说明其稠密性。
\end{example}

\begin{solve}
任取 $a\in\mathbb R^n$ 和 $\delta>0$，在各区间 $(a_i-\delta/\sqrt n,a_i+\delta/\sqrt n)$ 中选取有理数 $q_i$，并使 $q_1\neq a_1$。于是 $q=(q_1,\ldots,q_n)\in\mathbb Q^n$，$q\neq a$ 且 $|q-a|<\delta$。故 $(\mathbb Q^n)'=\mathbb R^n$，从而 $\mathbb Q^n$ 在 $\mathbb R^n$ 中稠密。此例也说明，一个集合的聚点可以遍布整个空间，即使原集合只含有理坐标点。
\end{solve}

\subsection{开集与闭集}
\label{subsec:open-and-closed-sets}

\begin{definition}
\label{def:euclidean-interior-exterior-boundary}
设 $E\subseteq\mathbb R^n$，$E^c=\mathbb R^n\setminus E$。若存在 $\delta>0$ 使 $U(a,\delta)\subseteq E$，则称 $a$ 为 $E$ 的\emph{内点}（interior point），全体内点组成\emph{内部} $E^\circ$。若存在 $\delta>0$ 使 $U(a,\delta)\cap E=\varnothing$，则称 $a$ 为\emph{外点}（exterior point），全体外点组成 $(E^c)^\circ$。若每个 $U(a,\delta)$ 都同时与 $E$ 和 $E^c$ 相交，则称 $a$ 为\emph{边界点}（boundary point），全体边界点组成\emph{边界} $\partial E$。
\end{definition}

内点一定属于 $E$，外点一定不属于 $E$，边界点则可能属于或不属于 $E$。每个点恰属于内部、外部、边界三者之一，因而
\[
    \mathbb R^n=E^\circ\cup(E^c)^\circ\cup\partial E
\]
是两两不交的分解。\emph[caution]{边界点与聚点不是同一个概念}：孤立点属于边界，却不是聚点；开球的内点都是聚点，却不是边界点。

\begin{definition}
\label{def:euclidean-open-closed-closure}
若 $E=E^\circ$，即每一点都是内点，则称 $E$ 为\emph{开集}（open set）；若 $E^c$ 为开集，则称 $E$ 为\emph{闭集}（closed set）。集合 $\overline E=E\cup E'$ 称为 $E$ 的\emph{闭包}（closure）。空集按定义既开又闭。
\end{definition}

球形邻域是开集：若 $x\in U(a,r)$，则 $U(x,r-|x-a|)\subseteq U(a,r)$。方形邻域同样是开集，可取各坐标到边界距离的最小值作为更小邻域的半径。$\mathbb R^n$ 也既开又闭；因此开与闭并非互为否定，一个集合还可能既不开也不闭。

\begin{theorem}[开闭集的运算]
\label{thm:euclidean-open-closed-set-operations}
任意一族开集的并为开集，有限个开集的交为开集；任意一族闭集的交为闭集，有限个闭集的并为闭集。
\end{theorem}

\begin{proof}
若一点属于一族开集的并，它属于其中某一个开集，该开集提供一个包含于并集的邻域。若 $x\in\bigcap_{j=1}^s O_j$ 且各 $O_j$ 开，对每个 $j$ 选 $r_j>0$ 使 $U(x,r_j)\subseteq O_j$，再取 $r=\min_jr_j>0$，就有 $U(x,r)\subseteq\bigcap_jO_j$。闭集结论由补集及德摩根公式
\[
    \left(\bigcup_{\lambda\in\Lambda}E_\lambda\right)^c
       =\bigcap_{\lambda\in\Lambda}E_\lambda^c,\qquad
    \left(\bigcap_{\lambda\in\Lambda}E_\lambda\right)^c
       =\bigcup_{\lambda\in\Lambda}E_\lambda^c
\]
得到，其中指标集 $\Lambda$ 不必有限或可数。
\end{proof}

“有限个”的限制不可随意删除。例如
\[
    \bigcap_{k=1}^\infty U(a,1/k)=\{a\}
\]
不是开集；在 $\mathbb R$ 中，$\bigcup_{k=2}^\infty[1/k,1-1/k]=(0,1)$ 不是闭集。开集也未必只有一块，例如 $\bigcup_{k\in\mathbb Z}U((0,k),1/2)\subseteq\mathbb R^2$ 是互不相交的开圆盘之并。在实直线上，每个开集可分解为至多可数个互不相交的开区间，允许无界区间：对开集中的点，合并所有包含该点并且包含于开集的开区间，得到一个极大开区间；这些极大区间相交即相同，每个又含一个有理数，故至多可数。高维开集的边界与形状可以复杂得多，不能套用一维区间的直观。

\begin{theorem}[闭集的聚点与点列刻画]
\label{thm:euclidean-closed-set-criteria}
对 $E\subseteq\mathbb R^n$，下列条件等价：$E$ 为闭集；$E'\subseteq E$；$\overline E=E$；$E$ 中任一收敛点列的极限仍属于 $E$。
\end{theorem}

\begin{proof}
若 $E$ 闭，则每个 $a\notin E$ 都有一个包含于 $E^c$ 的邻域，故 $a\notin E'$，得到 $E'\subseteq E$。反之，若 $E'\subseteq E$，则 $a\notin E$ 时 $a\notin E'$，有一个去心邻域与 $E$ 不交；由于 $a\notin E$，整个邻域都与 $E$ 不交，故 $E^c$ 开。$E'\subseteq E$ 与 $E\cup E'=E$ 的等价性直接来自集合运算。

若 $E$ 闭且 $x_k\in E$、$x_k\to a\notin E$，则 $a$ 在 $E^c$ 中的邻域最终包含 $x_k$，矛盾。反之，若收敛点列的极限均留在 $E$，则对 $a\in E'$ 应用定理~\ref{thm:euclidean-accumulation-sequence-criterion}，即可得到 $a\in E$，从而 $E$ 闭。
\end{proof}

闭包还可由邻域描述：$a\in\overline E$ 当且仅当 $a$ 的每个邻域都与 $E$ 相交。若 $a\notin\overline E$，可取一个与 $E$ 不交的邻域，其内每点都有更小邻域与 $E$ 不交，故 $(\overline E)^c$ 开，$\overline E$ 闭。由此及点的分类，有
\begin{equation}
    \overline E=E\cup\partial E,
    \qquad \partial E=\overline E\setminus E^\circ
                      =\overline E\cap\overline{E^c}.
    \label{eq:euclidean-closure-boundary-identities}
\end{equation}
例如，若 $a\in\overline E\setminus E^\circ$，则每个邻域既与 $E$ 相交，又不包含于 $E$，所以 $a\in\partial E$；其余等式同样由邻域描述得到。特别地，$\overline E$ 是包含 $E$ 的最小闭集：若闭集 $F\supseteq E$，则 $E'\subseteq F'\subseteq F$，所以 $\overline E\subseteq F$。

\begin{example}
\label{ex:half-open-disk-interior-boundary}
设 $E=\{(x,y):x^2+y^2<1,\ y\geq x\}$，求其内部、外部和边界。
\end{example}

\begin{solve}
点在圆盘内部且严格位于直线 $y=x$ 上方时，可取足够小的圆盘保留两个严格不等式，故
\[
    E^\circ=\{(x,y):x^2+y^2<1,\ y>x\}.
\]
直线 $y=x$ 上位于圆盘内的点，其任意邻域都跨越直线；满足 $x^2+y^2=1$、$y\geq x$ 的圆周点，其任意邻域也同时含有 $E$ 与 $E^c$ 中的点，包括两个交点。因此
\[
    \overline E=\{(x,y):x^2+y^2\leq1,\ y\geq x\},\qquad
    (E^c)^\circ=\mathbb R^2\setminus\overline E,
\]
\[
    \partial E=\{(x,y):x^2+y^2=1,\ y\geq x\}
          \cup\{(x,y):x^2+y^2<1,\ y=x\}.
\]
直线段的内部点属于 $E$，圆弧上的点不属于 $E$，所以本例同时展示了两类边界点。
\end{solve}

\subsection{欧氏空间 \texorpdfstring{$\mathbb{R}^n$}{R^n} 中的基本定理}
\label{subsec:fundamental-theorems-of-euclidean-space}

实数系中完备性、区间套、聚点原理与有限覆盖定理的证明，大多可以借助距离和有限个坐标推广。本小节按这一依赖关系建立相应结论，它们也是连续函数整体性质的基础。

\begin{definition}
\label{def:euclidean-cauchy-sequence-completeness}
点列 $\{x_k\}\subseteq\mathbb R^n$ 称为\emph{柯西点列}（Cauchy sequence），是指
\[
    \forall\varepsilon>0\ \exists K\in\mathbb N\ \forall p,q>K,
    \qquad |x_p-x_q|<\varepsilon.
\]
若子集 $X\subseteq\mathbb R^n$ 中每个柯西点列都收敛，且极限仍属于 $X$，则称 $X$ 在所继承的距离下\emph{完备}（complete）。
\end{definition}

\begin{theorem}[欧氏空间的完备性]
\label{thm:euclidean-completeness}
$\mathbb R^n$ 中的点列收敛，当且仅当它是柯西点列。特别地，$\mathbb R^n$ 完备。
\end{theorem}

\begin{proof}
若 $x_k\to a$，三角不等式给出 $|x_p-x_q|\leq|x_p-a|+|x_q-a|$，因此满足柯西条件。反之，若 $\{x_k\}$ 是柯西点列，则 $|x_{p,i}-x_{q,i}|\leq|x_p-x_q|$ 说明每个坐标数列都是实数系中的柯西数列。由实数系的完备性，存在 $a_i=\lim_kx_{k,i}$；按分量判定收敛得到 $x_k\to(a_1,\ldots,a_n)$。
\end{proof}

完备性还要求极限留在所讨论的空间内，例如 $(0,1)$ 中的柯西数列 $1/(k+1)$ 的极限为 $0\notin(0,1)$，所以 $(0,1)$ 不完备；由闭集的点列刻画，$\mathbb R^n$ 的每个闭子集都是完备的。

\begin{definition}
\label{def:euclidean-set-diameter}
非空集合 $E\subseteq\mathbb R^n$ 的\emph{直径}（diameter）定义为
\[
    \operatorname{diam}E=\sup\{|x-y|:x,y\in E\},
\]
无界时允许其值为 $+\infty$。
\end{definition}

\begin{theorem}[闭集套定理]
\label{thm:euclidean-nested-closed-sets}
设 $F_k\subseteq\mathbb R^n$ 为非空闭集，$F_{k+1}\subseteq F_k$，且 $\operatorname{diam}F_k\to0$。则存在唯一的 $a\in\mathbb R^n$，使
\[
    \bigcap_{k=1}^\infty F_k=\{a\}.
\]
\end{theorem}

\begin{proof}
任取 $x_k\in F_k$。给定 $\varepsilon>0$，取 $K$ 使 $\operatorname{diam}F_K<\varepsilon$。当 $p,q\geq K$ 时，嵌套性保证 $x_p,x_q\in F_K$，故 $|x_p-x_q|<\varepsilon$。由完备性，$x_k\to a$。对每个固定的 $j$，尾列均在闭集 $F_j$ 中，所以 $a\in F_j$，即 $a$ 属于全部 $F_j$。若 $b$ 也是公共点，则 $|a-b|\leq\operatorname{diam}F_k$ 对所有 $k$ 成立，令 $k\to\infty$ 得 $a=b$。
\end{proof}

闭性保证所构造的极限不被排除，直径趋于零保证唯一性，并使所选点列满足柯西条件。若把闭集换成 $(0,1/k)$，公共交为空；若仅保留非空、闭性和嵌套性，$[k,+\infty)$ 的公共交仍可为空。教材中的闭区间套定理是上述结论在一维的特例。

\begin{theorem}[波尔查诺--魏尔斯特拉斯定理（Bolzano--Weierstrass theorem）]
\label{thm:euclidean-bolzano-weierstrass}
$\mathbb R^n$ 中任一有界点列都有收敛子列；任一有界无限集至少有一个聚点。
\end{theorem}

\begin{proof}
设 $\{x_k\}$ 有界。先从第一坐标数列选取收敛子列，再在该子列中使第二坐标收敛，依次进行。每次再取子列不破坏之前坐标的收敛性，经过有限的 $n$ 次选取，得到所有坐标都收敛的同一个子列，因而得到收敛点列。若 $E$ 是有界无限集，先在其中选取两两不同的点 $x_k$，再取收敛子列；由聚点的点列刻画，其极限是 $E$ 的聚点，但未必属于 $E$。

反过来，有界无限集存在聚点也能推出点列结论：若点列有无限个不同取值，聚点的每个邻域含无穷多个不同取值，便能递增选取原点列的下标，使相应项与聚点的距离小于 $1/j$；若只有有限个取值，则至少一个值出现无穷多次，提供常值子列。因此两种表述等价。
\end{proof}

\begin{definition}
\label{def:euclidean-compact-set}
设 $E\subseteq\mathbb R^n$。若一族 $\mathbb R^n$ 中的开集 $\{O_\lambda\}_{\lambda\in\Lambda}$ 满足 $E\subseteq\bigcup_{\lambda\in\Lambda}O_\lambda$，则称其为 $E$ 的\emph{开覆盖}（open cover）。若 $E$ 的每个开覆盖都能从原覆盖中选出有限个成员，仍覆盖 $E$，则称 $E$ 为\emph{紧集}（compact set），所选成员构成\emph{有限子覆盖}（finite subcover）。
\end{definition}

\emph[caution]{紧性要求对任意给定的开覆盖都能选出有限子覆盖。}仅仅找到一个有限开覆盖毫无判别力，因为单个开集 $\mathbb R^n$ 就能覆盖任何 $E$。空集也是紧集；下文涉及最大、最小值时则必须另加非空条件。

\begin{theorem}[海涅--博雷尔定理（Heine--Borel theorem）]
\label{thm:euclidean-heine-borel}
集合 $E\subseteq\mathbb R^n$ 为紧集，当且仅当 $E$ 是有界闭集。
\end{theorem}

\begin{proof}
空集的情形直接成立，以下设 $E\neq\varnothing$。若 $E$ 紧，则开覆盖 $\{U(0,k)\}_{k\geq1}$ 有有限子覆盖；这些球逐个包含，故 $E\subseteq U(0,K)$ 对某个 $K$ 成立，$E$ 有界。若 $a\in E'\setminus E$，对每个 $x\in E$ 令 $r_x=|x-a|/2>0$。球 $U(x,r_x)$ 覆盖 $E$，且与 $U(a,r_x)$ 不交。选出有限子覆盖 $U(x_j,r_{x_j})$，取 $r=\min_jr_{x_j}>0$，就有 $U(a,r)\cap E=\varnothing$，与 $a\in E'$ 矛盾。故 $E'\subseteq E$，$E$ 闭。

反之，设 $E$ 是非空有界闭集。若某个开覆盖 $\{O_\lambda\}$ 无有限子覆盖，取闭立方体 $Q_0=[-a,a]^n\supseteq E$。把每条边二等分，将 $Q_0$ 分为 $2^n$ 个闭立方体，其中至少一个 $Q_1$ 满足 $E\cap Q_1$ 不能由原覆盖的有限个成员覆盖，否则合并各部分的有限子覆盖就能覆盖 $E$。反复二分并选取，可得嵌套的 $Q_k$，使每个 $E\cap Q_k$ 均无有限子覆盖，且 $\operatorname{diam}Q_k=2a\sqrt n/2^k\to0$。

令 $F_k=E\cap Q_k$。它们是非空嵌套闭集，直径趋于零。由闭集套定理，存在唯一公共点 $x_0\in E$。选覆盖中包含 $x_0$ 的 $O_{\lambda_0}$，再取 $r>0$ 使 $U(x_0,r)\subseteq O_{\lambda_0}$。当 $k$ 充分大时，对所有 $x\in F_k$ 都有 $|x-x_0|\leq\operatorname{diam}F_k<r$，故 $F_k\subseteq O_{\lambda_0}$，这与 $F_k$ 没有有限子覆盖矛盾。因此 $E$ 紧。
\end{proof}

\begin{corollary}[紧性的点列刻画]
\label{cor:euclidean-sequential-compactness}
$E\subseteq\mathbb R^n$ 紧，当且仅当 $E$ 中每个点列都有一个极限仍在 $E$ 中的收敛子列。
\end{corollary}

\begin{proof}
若 $E$ 紧，则有界性提供收敛子列，闭性保证极限属于 $E$。反之，若点列条件成立而 $E$ 无界，可选 $x_k\in E$ 使 $|x_k|>k$，它不可能有收敛子列，矛盾。若 $x_k\in E$ 收敛于 $a$，再取一个极限属于 $E$ 的收敛子列，子列仍收敛于 $a$，故 $a\in E$；于是 $E$ 闭，再由定理~\ref{thm:euclidean-heine-borel} 得紧性。
\end{proof}

例如 $\{1/k:k\in\mathbb N\}\subseteq\mathbb R$ 有界却不闭，因遗漏聚点 $0$ 而不紧；补上 $0$ 后就是紧集。有限维欧氏空间中“紧”等价于“有界且闭”，这一等价关系依赖当前的空间，不能不加条件地作为任意空间中紧性的定义。

\section{多元函数与向量函数的极限}
\label{sec:limits-of-multivariable-and-vector-functions}

\subsection{多元函数的概念}
\label{subsec:concept-of-multivariable-functions}

\begin{definition}
\label{def:multivariable-real-valued-function}
设 $\varnothing\neq E\subseteq\mathbb R^n$。若对应法则 $f$ 为每个 $x=(x_1,\ldots,x_n)\in E$ 指定唯一实数 $u$，则称 $f:E\to\mathbb R$ 为 $E$ 上的一个\emph{$n$ 元实值函数}（real-valued function of $n$ variables），记为 $u=f(x)=f(x_1,\ldots,x_n)$。$E$ 为定义域，$f(E)=\{f(x):x\in E\}$ 为值域；当 $n\geq2$ 时，统称\emph{多元函数}（multivariable function）。
\end{definition}

对二元函数 $z=f(x,y)$，其图像是 $\mathbb R^3$ 中的点集 $\{(x,y,f(x,y)):(x,y)\in E\}$，具有适当正则性时可看作曲面。一般 $n$ 元实值函数的图像位于 $\mathbb R^{n+1}$，定义域则位于 $\mathbb R^n$，两者所在空间应当区分。函数有界是指存在 $M>0$ 使 $|f(x)|\leq M$ 对所有 $x\in E$ 成立；其上、下确界分别是实数集 $f(E)$ 的上、下确界。

\begin{example}
\label{ex:sphere-two-hemisphere-graphs}
说明单位球面 $x^2+y^2+z^2=1$ 与二元函数图像的关系。
\end{example}

\begin{solve}
在闭单位圆盘 $\overline\Delta=\{(x,y):x^2+y^2\leq1\}$ 上，两个函数
\[
    z=\sqrt{1-x^2-y^2},\qquad z=-\sqrt{1-x^2-y^2}
\]
分别表示上、下半球面。整个球面不能写成同一 $(x,y)$ 定义域上的单个 $z=f(x,y)$ 的图像，因为圆盘内部每个 $(x,y)$ 对应两个不同的 $z$；函数的单值性需要在选择上、下分支时保证。
\end{solve}

\subsection{多元函数的极限}
\label{subsec:limits-of-multivariable-functions}

\begin{definition}
\label{def:multivariable-function-limit}
设 $f:E\to\mathbb R$，$E\subseteq\mathbb R^n$，$a\in E'$。若存在 $A\in\mathbb R$，使
\[
    \forall\varepsilon>0\ \exists\delta>0\ \forall x\in E,
    \qquad 0<|x-a|<\delta\ \Longrightarrow\ |f(x)-A|<\varepsilon,
\]
则称 $A$ 为 $x$ 在 $E$ 中趋于 $a$ 时 $f(x)$ 的\emph{极限}，记为 $\lim_{\substack{x\to a\\ x\in E}}f(x)=A$，也写成 $\lim_{E\ni x\to a}f(x)=A$。若 $E$ 包含 $a$ 的某个去心邻域，可省略 $x\in E$。
\end{definition}

条件 $a\in E'$ 保证任意小的去心邻域中都有可检验的自变量；极限不要求 $a\in E$，即使 $f(a)$ 有定义也不参与该极限的判定。对二元函数，$|(x,y)-(a,b)|=\sqrt{(x-a)^2+(y-b)^2}$，可等价地用方形去心邻域描述接近。\emph[caution]{方形去心邻域只删去中心点，并不删去穿过中心的两条坐标直线。}相对于 $E$ 的定义也统一包含一元函数的单侧极限以及从高维区域内部趋近边界的极限。

\begin{theorem}[多元函数极限的点列准则]
\label{thm:multivariable-heine-criterion}
设 $f:E\to\mathbb R$，$a\in E'$，$A\in\mathbb R$。则 $\lim_{E\ni x\to a}f(x)=A$ 当且仅当对每个 $E\setminus\{a\}$ 中满足 $x_k\to a$ 的点列，都有 $f(x_k)\to A$。
\end{theorem}

\begin{proof}
若函数极限存在，给定 $\varepsilon>0$ 后先取定义中的 $\delta$，再用 $x_k\to a$ 保证尾项满足 $0<|x_k-a|<\delta$，即可得 $f(x_k)\to A$。反之，若函数极限不等于 $A$，存在 $\varepsilon_0>0$，使对每个 $k$ 都能选取 $x_k\in E$ 满足
\[
    0<|x_k-a|<1/k,\qquad |f(x_k)-A|\geq\varepsilon_0.
\]
于是 $x_k\to a$，但 $f(x_k)$ 不趋于 $A$，与点列条件矛盾。
\end{proof}

与一元情形相同，有限极限唯一，存在时函数在某个相对去心邻域内有界；若 $A>0$，则在足够小的相对去心邻域内有 $f(x)>A/2$。唯一性可取一个趋于 $a$ 的点列，应用实数列极限的唯一性；后两个结论分别在定义中取 $\varepsilon=1$、$\varepsilon=A/2$ 即得。若同一定义域上的 $f,g$ 分别趋于有限的 $A,B$，则 $\alpha f+\beta g\to\alpha A+\beta B$、$fg\to AB$，且 $B\neq0$ 时 $f/g\to A/B$，最后一个商只需在足够小邻域内定义。实值函数的夹逼定理同样成立。这些结论均可由点列准则归结为实数列的对应定理。

无穷极限也可用距离定义：$f(x)\to+\infty$ 是指对每个 $M>0$，存在 $\delta>0$，使 $x\in E$、$0<|x-a|<\delta$ 时 $f(x)>M$；$f(x)\to-\infty$ 类似，$|f(x)|\to+\infty$ 则用 $|f(x)|>M$ 描述。若定义域无界，还可把 $0<|x-a|<\delta$ 换成 $|x|>R$，定义 $x$ 趋于无穷远时的函数极限。点列准则对这些情形也成立，证明中的 $\varepsilon_0$ 条件相应换成固定阈值条件即可。

\begin{example}
\label{ex:multivariable-tangent-cubic-limit}
求 $\displaystyle\lim_{(x,y)\to(0,0)}\frac{\tan(x^3+y^3)}{x^2+y^2}$。
\end{example}

\begin{solve}
令 $r=\sqrt{x^2+y^2}$，则 $|x^3+y^3|\leq|x|^3+|y|^3\leq r(x^2+y^2)=r^3$。由 $\tan t/t\to1$，存在 $\eta>0$ 使 $|t|<\eta$ 时 $|\tan t|\leq2|t|$，其中 $t=0$ 也成立。因此对足够小的 $r>0$，
\[
    \left|\frac{\tan(x^3+y^3)}{x^2+y^2}\right|
       \leq2\frac{|x^3+y^3|}{r^2}\leq2r\to0.
\]
所以所求极限为 $0$。这里得到的界只依赖距离 $r$，因而同时控制了所有趋近方向，也避免了在 $x^3+y^3=0$ 时使用未定义的比值。
\end{solve}

\begin{example}
\label{ex:multivariable-xy-direction-dependent-limit}
判断 $\displaystyle f(x,y)=\frac{xy}{x^2+y^2}$ 在 $(x,y)\to(0,0)$ 时是否有极限，其中 $(x,y)\neq(0,0)$。
\end{example}

\begin{solve}
沿点列 $(0,1/k)$，函数值恒为 $0$；沿 $(1/k,1/k)$，函数值恒为 $1/2$。由点列准则，联合极限不存在。更一般地，沿直线 $y=cx$ 有 $f(x,cx)=c/(1+c^2)$，结果依赖斜率。用极坐标 $x=r\cos\theta$、$y=r\sin\theta$，则 $f=(1/2)\sin2\theta$，角度变化也显示了这种方向依赖。
\end{solve}

\begin{example}
\label{ex:all-line-limits-insufficient}
判断 $\displaystyle f(x,y)=\frac{x^2y}{x^4+y^2}$ 在原点是否有极限，其中 $(x,y)\neq(0,0)$。
\end{example}

\begin{solve}
沿 $y=cx$，若 $c\neq0$，则 $f(x,cx)=cx/(x^2+c^2)\to0$；若 $c=0$，函数值恒为零，沿竖直线 $x=0$ 也为零。但沿抛物线 $y=x^2$，有 $f(x,x^2)=1/2$，所以联合极限不存在。\emph[caution]{沿所有直线所得极限相同，仍不足以推出多元极限存在。}若使用极坐标证明存在性，必须对所有角度给出统一控制；只对每个固定角度令 $r\to0$，仍然只检查了直线方向。
\end{solve}

\begin{example}
\label{ex:ball-boundary-radial-logarithm-limit}
设 $E=\{(x,y,z):x^2+y^2+z^2<1\}$，$a\in\partial E$，求
\[
    \lim_{E\ni(x,y,z)\to a}
    [\sin1-\sin(x^2+y^2+z^2)]\ln[1-(x^2+y^2+z^2)].
\]
\end{example}

\begin{solve}
令 $t=x^2+y^2+z^2$。当自变量从 $E$ 趋于 $a$ 时，由范数的连续估计有 $t\to1^-$。一元中值定理给出 $|\sin1-\sin t|\leq1-t$，故原式绝对值不超过 $(1-t)|\ln(1-t)|\to0$。所求极限为 $0$，且对任意球面点 $a$ 结论相同。这里函数只在开球内定义，极限从该定义域取值，不涉及球外使对数无意义的点。
\end{solve}

\subsection{累次极限}
\label{subsec:iterated-limits}

联合极限同时约束全部自变量，而\emph{累次极限}（iterated limit）先固定一部分变量，对其余变量取极限，再对所得到的函数取极限；两者的量词和存在条件不同。以下以二元函数说明。

\begin{definition}
\label{def:iterated-multivariable-limits}
设 $f$ 在 $N_0((a,b),\delta_0)$ 上有定义。若对每个 $0<|y-b|<\delta_0$，极限 $\varphi(y)=\lim_{x\to a}f(x,y)$ 存在，且 $\lim_{y\to b}\varphi(y)=A$，则称
\[
    \lim_{y\to b}\lim_{x\to a}f(x,y)=A
\]
为先 $x$ 后 $y$ 的累次极限。先 $y$ 后 $x$ 的累次极限 $\lim_{x\to a}\lim_{y\to b}f(x,y)$ 类似定义。
\end{definition}

两个累次极限都不涉及直线 $x=a$、$y=b$ 上的函数值：内层极限删去相应中心值，外层极限也删去另一变量的中心值；联合极限则仍须考虑这些直线上异于中心的点。例如
\[
    u(x,y)=\begin{cases}0,&xy=0,\\1,&xy\neq0\end{cases}
\]
的两个累次极限在原点都为 $1$，但沿坐标轴函数值为 $0$，沿 $y=x\neq0$ 为 $1$，联合极限不存在。即使避开坐标轴，这两类极限仍不等价：例~\ref{ex:multivariable-xy-direction-dependent-limit} 的两个累次极限都为 $0$，联合极限却不存在。

联合极限存在也不保证内层极限存在。令
\[
    v(x,y)=\begin{cases}
        (x+y)\sin(1/x)\sin(1/y),&xy\neq0,\\
        0,&xy=0.
    \end{cases}
\]
因为 $|v(x,y)|\leq|x|+|y|$，其联合极限为 $0$。但固定满足 $y\sin(1/y)\neq0$ 的 $y$ 时，写成
\[
    v(x,y)=y\sin(1/y)\sin(1/x)+x\sin(1/y)\sin(1/x),
\]
右边第二项趋于零，第一项振荡而没有极限。任意小的去心 $y$ 邻域都存在这样的 $y$，所以先 $x$ 后 $y$ 的累次极限不满足上述定义；对称地，另一顺序也不存在。

\begin{theorem}[联合极限与已有的内层极限]
\label{thm:joint-limit-implies-existing-iterated-limit}
设 $f$ 在 $N_0((a,b),\delta_0)$ 上有定义，联合极限为有限数 $A$。若对每个 $0<|y-b|<\delta_0$ 都存在有限的 $\varphi(y)=\lim_{x\to a}f(x,y)$，则
\[
    \lim_{y\to b}\varphi(y)
       =\lim_{y\to b}\lim_{x\to a}f(x,y)=A.
\]
交换 $x,y$ 有对应结论。因此，在联合极限及两组内层极限均存在的条件下，两个累次极限存在且等于联合极限。
\end{theorem}

\begin{proof}
给定 $\varepsilon>0$，由联合极限的方形邻域描述，取 $0<\delta<\delta_0$，使 $0<\max\{|x-a|,|y-b|\}<\delta$ 时 $|f(x,y)-A|<\varepsilon/2$。固定任意 $0<|y-b|<\delta$，在此不等式中令 $x\to a$，得到 $|\varphi(y)-A|\leq\varepsilon/2<\varepsilon$，所以 $\varphi(y)\to A$。关键在于，联合极限给出的同一个 $\delta$ 对所有这些 $y$ 都有效。
\end{proof}

若联合极限为 $+\infty$，而各内层极限存在，则先用任意阈值 $M+1$ 代替上述误差界，再令 $x\to a$，得到 $\varphi(y)\geq M+1>M$，故累次极限也为 $+\infty$；$-\infty$ 的情形相同。不能跳过内层极限的存在性，直接把联合极限改写为累次极限。

\subsection{向量函数的定义与极限}
\label{subsec:definition-and-limits-of-vector-functions}

一般地，\emph{映射}（mapping）$f:A\to B$ 把集合 $A$ 中每个元素指定为 $B$ 中唯一的一个元素；$f(A)\subseteq B$ 是像集。若不同元素有不同像，则为\emph{单射}（injection）；若 $f(A)=B$，则为\emph{满射}（surjection）；两者兼具时为\emph{双射}（bijection），这时存在逆映射 $f^{-1}:B\to A$。区分值域 $f(A)$ 与指定的目标集合 $B$，有助于准确判断满射性。

\begin{definition}
\label{def:multivariable-vector-valued-function}
映射 $f:E\subseteq\mathbb R^n\to\mathbb R^m$ 称为一个 $m$ 维 $n$ 元\emph{向量函数}（vector-valued function）。它唯一表示为
\[
    f(x)=(f_1(x),\ldots,f_m(x)),
\]
其中 $f_j:E\to\mathbb R$ 是第 $j$ 个分量函数。若存在 $M>0$ 使 $|f(x)|\leq M$ 对所有 $x\in E$ 成立，则称 $f$ 有界。
\end{definition}

\begin{definition}
\label{def:vector-valued-function-limit}
设 $a\in E'$，$A\in\mathbb R^m$。若对任意 $\varepsilon>0$，存在 $\delta>0$，使 $x\in E$、$0<|x-a|<\delta$ 时 $|f(x)-A|<\varepsilon$，则称 $f(x)$ 趋于 $A$。由坐标与范数的估计，这等价于
\[
    \lim_{E\ni x\to a}f_j(x)=A_j\quad(j=1,\ldots,m),\qquad
    A=(A_1,\ldots,A_m).
\]
\end{definition}

必要性来自 $|f_j(x)-A_j|\leq|f(x)-A|$；充分性则对有限个分量分别取误差 $\varepsilon/\sqrt m$，再取相应 $\delta_j$ 的最小值。任一分量没有有限极限，就没有 $\mathbb R^m$ 中的向量极限，不能把实数的 $+\infty$ 当作向量的一个普通坐标。点列准则、极限唯一性及局部有界性仍成立；同维向量函数的线性组合、内积和范数的极限由各分量的实数运算得到。

\begin{example}
\label{ex:polar-map-boundary-limits}
设 $P(r,\theta)=(r\cos\theta,r\sin\theta)$，定义域为 $D=(0,+\infty)\times[0,2\pi)$。说明此映射的单射性，并求固定 $r_0>0$ 时 $(r,\theta)\to(r_0,0)$ 与 $(r,\theta)\to(r_0,2\pi)$ 的极限。
\end{example}

\begin{solve}
对每个非零平面点，半径 $r>0$ 及取在 $[0,2\pi)$ 内的角度唯一，所以 $P:D\to\mathbb R^2\setminus\{0\}$ 是双射。若允许 $r=0$，则不同角度都映成原点，不再单射。按分量取极限，有
\[
    \lim_{D\ni(r,\theta)\to(r_0,0)}P(r,\theta)
       =\lim_{D\ni(r,\theta)\to(r_0,2\pi)}P(r,\theta)=(r_0,0).
\]
第二个参数点不属于 $D$，却是其聚点，因而相对定义域的极限有意义。两端角度在像中汇合的现象，将在同胚映射部分解释逆映射为何可能不连续。
\end{solve}

\section{多元连续函数}
\label{sec:continuous-multivariable-functions}

\subsection{多元连续函数}
\label{subsec:continuous-multivariable-functions}

\begin{definition}
\label{def:relative-multivariable-continuity}
设 $f:E\subseteq\mathbb R^n\to\mathbb R$，$a\in E$。若
\[
    \forall\varepsilon>0\ \exists\delta>0\ \forall x\in E,
    \qquad |x-a|<\delta\ \Longrightarrow\ |f(x)-f(a)|<\varepsilon,
\]
则称 $f$ 在 $a$ 处\emph{连续}（continuous）；若在 $E$ 的每一点都连续，则称 $f$ 在 $E$ 上连续。不连续的点称为\emph{间断点}（point of discontinuity）。这里的连续性始终相对于定义域 $E$。
\end{definition}

若 $a\in E'$，连续等价于 $\lim_{E\ni x\to a}f(x)=f(a)$；若 $a$ 是 $E$ 的孤立点，取 $U(a,\delta)\cap E=\{a\}$，上述条件自动成立，所以任何函数在定义域的孤立点都连续。与函数极限不同，连续性要求 $a\in E$，并且比较的目标值固定为 $f(a)$。由点列准则，连续性还等价于：每个 $E$ 中收敛于 $a$ 的点列均满足 $f(x_k)\to f(a)$，其中允许 $x_k=a$。

连续函数在每一点附近相对有界；若 $f(a)>0$，则在足够小的相对邻域内 $f(x)>f(a)/2$。同一定义域上的连续实值函数经过和、差、积仍连续，作商时必须保证分母在所考查点非零。这些局部结论不能自动升级为整个定义域上的有界性或统一的正下界；后者需要紧性等整体条件。

\begin{example}
\label{ex:oscillatory-multivariable-continuity-set}
讨论函数
\[
    f(x,y)=\begin{cases}
        x\sin(1/y),&xy\neq0,\\
        0,&xy=0
    \end{cases}
\]
在 $\mathbb R^2$ 上的连续性。
\end{example}

\begin{solve}
当 $y_0\neq0$ 时，在该点附近有 $y\neq0$，两种写法均与 $x\sin(1/y)$ 一致，故函数连续。对任意 $y_0\in\mathbb R$，又有 $|f(x,y)|\leq|x|$，所以 $(x,y)\to(0,y_0)$ 时 $f(x,y)\to0=f(0,y_0)$，函数在整个 $y$ 轴上连续。

若 $x_0\neq0$，固定 $x=x_0$，取
\[
    y_k=\frac1{2k\pi+\pi/2},\qquad
    \widetilde y_k=\frac1{2k\pi+3\pi/2},
\]
则 $(x_0,y_k)$ 与 $(x_0,\widetilde y_k)$ 均趋于 $(x_0,0)$，函数值却分别为 $x_0$ 与 $-x_0$。所以间断点集恰为 $\{(x,0):x\neq0\}$。决定连续性的不是三角函数是否出现振荡，而是振荡的幅度是否被趋于零的因子统一控制。
\end{solve}

分别固定其余变量所得的一元连续性，称为\emph{分别连续}（separately continuous）。联合连续必然蕴含分别连续，但反向一般不成立。例如把例~\ref{ex:multivariable-xy-direction-dependent-limit} 的函数在原点补定义为 $0$，固定任一 $x$ 或 $y$ 后均得到连续的一元函数，它在原点却没有联合极限。

\subsection{多元连续向量函数}
\label{subsec:continuous-multivariable-vector-functions}

\begin{definition}
\label{def:relative-vector-function-continuity}
设 $f:E\subseteq\mathbb R^n\to\mathbb R^m$，$a\in E$。若对任意 $\varepsilon>0$ 存在 $\delta>0$，使 $x\in E$、$|x-a|<\delta$ 时 $|f(x)-f(a)|<\varepsilon$，则称 $f$ 在 $a$ 处连续。若该条件在 $E$ 的每点成立，则称 $f$ 在 $E$ 上连续。
\end{definition}

按分量判别向量极限可知，$f=(f_1,\ldots,f_m)$ 连续当且仅当每个分量 $f_j$ 连续，包括定义域的孤立点情形。向量连续性同样有点列刻画。同维连续向量函数的和、差、常数倍仍连续；内积是连续实值函数，范数也连续，因为 $\bigl||f(x)|-|f(a)|\bigr|\leq|f(x)-f(a)|$。

\begin{theorem}[复合向量函数的连续性]
\label{thm:continuous-vector-function-composition}
设 $f:E\subseteq\mathbb R^n\to\mathbb R^m$ 在 $a\in E$ 处连续，$g:F\subseteq\mathbb R^m\to\mathbb R^\ell$ 在 $f(a)$ 处连续，且 $f(E)\subseteq F$。则 $g\circ f:E\to\mathbb R^\ell$ 在 $a$ 处连续。特别地，若两映射分别在整个相应定义域上连续，则复合映射也连续。
\end{theorem}

\begin{proof}
给定 $\varepsilon>0$，由 $g$ 在 $f(a)$ 连续取 $\eta>0$，使 $y\in F$、$|y-f(a)|<\eta$ 时 $|g(y)-g(f(a))|<\varepsilon$。再由 $f$ 连续取 $\delta>0$，使 $x\in E$、$|x-a|<\delta$ 时 $|f(x)-f(a)|<\eta$。代入 $y=f(x)$ 即得结论。这里允许 $f(x)=f(a)$，因此不存在复合函数极限中可能遇到的去心条件问题。
\end{proof}

\begin{example}
\label{ex:disk-logarithm-continuous-extension}
讨论定义在闭单位圆盘上的函数
\[
    f(x,y)=\begin{cases}
        \sqrt{1-x^2-y^2}\ln(1-x^2-y^2),&x^2+y^2<1,\\
        0,&x^2+y^2=1
    \end{cases}
\]
的连续性。
\end{example}

\begin{solve}
令 $H(s)=\sqrt s\ln s$（$0<s\leq1$），并补定义 $H(0)=0$。当 $s=e^{-t}\to0^+$ 时，$H(s)=-te^{-t/2}\to0$，所以 $H$ 在 $[0,1]$ 上连续。而闭圆盘上 $g(x,y)=1-x^2-y^2$ 连续，值落在 $[0,1]$，原函数正是 $H\circ g$。由复合函数连续性，$f$ 在整个闭圆盘上连续。这一论证同时涵盖从圆盘内部及沿圆周趋近边界点的情形；边界上连续并不要求将函数定义到圆盘外部。
\end{solve}

\subsection{集合的连通性}
\label{subsec:connectedness-of-sets}

连续映射 $h:[\alpha,\beta]\to\mathbb R^n$（$\alpha<\beta$）称为一条\emph{道路}或\emph{路径}（path），也称参数化的连续曲线。通过参数的线性变换，可把参数区间统一为 $[0,1]$；道路允许自交，也允许端点重合。

\begin{definition}
\label{def:path-connected-set-and-domain}
若非空集合 $E\subseteq\mathbb R^n$ 中任意两点 $a,b$ 都能由一条完全位于 $E$ 内的道路连接，即存在连续的 $h:[0,1]\to E$ 满足 $h(0)=a$、$h(1)=b$，则称 $E$ \emph{道路连通}（path-connected）。非空道路连通开集称为\emph{区域}（domain），区域 $D$ 的闭包 $\overline D=D\cup\partial D$ 称为\emph{闭区域}。
\end{definition}

\begin{notealert}[type=caution,title={本章的连通性约定}]
教材在本章及后续微积分讨论中，将“道路连通”简称为“连通”，以下沿用这一约定。一般拓扑学中的连通性有更广泛的定义，与道路连通并非对任意集合都等价；引用本章结论时应保留这里的约定。
\end{notealert}

在 $\mathbb R$ 中，非空集合道路连通当且仅当它是区间，允许单点区间及无界区间。若两点属于同一区间，线段 $h(t)=(1-t)a+tb$ 提供道路；反之，连接实数 $a<b$ 的连续道路由介值定理必须经过它们之间的每个数，所以道路连通集必为区间。在高维空间中，开圆盘、开球以及开圆盘闭包的外部都可成为区域；区域可以有孔，也未必有界。闭区域虽为闭集，却不必为紧集，还需有界性。

若 $E\subseteq\mathbb R^n$ 对任意 $a,b\in E$、$t\in[0,1]$ 都有 $(1-t)a+tb\in E$，则称 $E$ 为\emph{凸集}（convex set）；非空开凸集称为\emph{凸域}（convex domain）。上述线段直接证明非空凸集道路连通，但道路连通不要求线段留在集合内。例如圆环 $\{(x,y):1<x^2+y^2<4\}$ 中两点可由径向线段与圆弧连接，却有经过孔洞的连线，因而不凸。

\subsection{连续函数的性质}
\label{subsec:properties-of-continuous-multivariable-functions}

紧性控制点列的极限位置，道路连通性控制两点间能否由连续路径联系。它们分别对应连续函数的有界性、取到最值与一致连续性，以及介值性质，不能相互替代。

\begin{theorem}[连续像保持紧性]
\label{thm:continuous-image-of-euclidean-compact-set}
若 $E\subseteq\mathbb R^n$ 为紧集，$f:E\to\mathbb R^m$ 连续，则 $f(E)$ 为 $\mathbb R^m$ 中的紧集。
\end{theorem}

\begin{proof}
若 $f(E)$ 无界，可取 $x_k\in E$ 使 $|f(x_k)|>k$。由 $E$ 的紧性，存在子列 $x_{k_j}\to a\in E$，连续性给出 $f(x_{k_j})\to f(a)$，与其范数趋于无穷矛盾，所以 $f(E)$ 有界。

若 $y\in(f(E))'$，取 $y_k\in f(E)$ 使 $y_k\to y$，并为每项选一个原像 $x_k\in E$，使 $f(x_k)=y_k$。再取 $x_{k_j}\to a\in E$，由连续性及极限唯一性得到 $y=\lim_jy_{k_j}=f(a)\in f(E)$。因此 $f(E)$ 包含其全部聚点，是闭集；由海涅--博雷尔定理，它是紧集。若 $E$ 为空，像集为空，结论同样成立。
\end{proof}

\begin{corollary}[有界性与最值定理]
\label{cor:multivariable-extreme-value-theorem}
设 $\varnothing\neq E\subseteq\mathbb R^n$ 紧，$f:E\to\mathbb R$ 连续。则 $f$ 有界，且存在 $a,b\in E$，使
\[
    f(a)=\min_{x\in E}f(x)=\inf f(E),\qquad
    f(b)=\max_{x\in E}f(x)=\sup f(E).
\]
\end{corollary}

\begin{proof}
$f(E)$ 是非空有界闭实数集。令 $M=\sup f(E)$，对每个 $k$ 取 $u_k\in f(E)$ 使 $M-1/k<u_k\leq M$；则 $u_k\to M$，闭性保证 $M\in f(E)$。下确界同理。因此两个确界都有原像，最大、最小值均能取到。
\end{proof}

向量值函数的像集也有界，但 $\mathbb R^m$ 中没有本章所用的实数全序，不能笼统地说向量函数“取得最大值”。可以对每个实值分量或范数应用最值定理。例如连续且处处为正的实值函数在非空紧集上有严格正的最小值，因为取到最小值的点仍属于该定义域。

\begin{theorem}[连续像保持道路连通性]
\label{thm:continuous-image-path-connected}
若非空集合 $E\subseteq\mathbb R^n$ 道路连通，$f:E\to\mathbb R^m$ 连续，则 $f(E)$ 道路连通。
\end{theorem}

\begin{proof}
任取 $u,v\in f(E)$，选 $a,b\in E$ 使 $f(a)=u$、$f(b)=v$。取 $E$ 中连接 $a,b$ 的道路 $h:[0,1]\to E$，则 $f\circ h$ 连续，像包含于 $f(E)$，两端分别为 $u,v$，因而就是所需道路。
\end{proof}

\begin{corollary}[多元连续函数的介值定理]
\label{cor:multivariable-intermediate-value-theorem}
若 $E\subseteq\mathbb R^n$ 道路连通，$f:E\to\mathbb R$ 连续，且 $f(a)<c<f(b)$，则存在 $\xi\in E$ 使 $f(\xi)=c$。因此 $f(E)$ 是一个区间。若 $E$ 还非空且紧，则
\[
    f(E)=[\min_E f,\max_E f],
\]
允许两端点相同。
\end{corollary}

\begin{proof}
在连接 $a,b$ 的道路 $h$ 上，实值函数 $f\circ h$ 连续，由一元介值定理存在 $t_0\in(0,1)$ 使 $f(h(t_0))=c$，取 $\xi=h(t_0)$。再结合最值定理，即得非空紧且道路连通情形的像集表达式。
\end{proof}

介值定理只保证中间值能取到，不保证其原像唯一。若定义域紧却不连通，函数的像也不必是区间；例如在紧集 $\{-1,1\}$ 上，连续函数 $f(x)=x$ 的像只含两个点。反之，连通而不紧的定义域上，连续函数也未必有界或取到其确界。

\begin{definition}
\label{def:uniform-continuity-vector-function}
设 $f:E\subseteq\mathbb R^n\to\mathbb R^m$。若
\[
    \forall\varepsilon>0\ \exists\delta>0\ \forall x,y\in E,
    \qquad |x-y|<\delta\ \Longrightarrow\ |f(x)-f(y)|<\varepsilon,
\]
则称 $f$ 在 $E$ 上\emph{一致连续}（uniformly continuous）。同一个 $\delta$ 对所有点对有效，不能依赖 $x,y$。有限维情形下，这等价于各实值分量都在 $E$ 上一致连续。
\end{definition}

\begin{theorem}[紧集上的一致连续性]
\label{thm:multivariable-heine-cantor}
若 $E\subseteq\mathbb R^n$ 紧，$f:E\to\mathbb R^m$ 连续，则 $f$ 在 $E$ 上一致连续。
\end{theorem}

\begin{proof}
若不一致连续，则存在 $\varepsilon_0>0$，使每个 $k$ 都能选出 $x_k,y_k\in E$ 满足 $|x_k-y_k|<1/k$ 而 $|f(x_k)-f(y_k)|\geq\varepsilon_0$。紧性提供子列 $x_{k_j}\to a\in E$；由 $|y_{k_j}-a|\leq|y_{k_j}-x_{k_j}|+|x_{k_j}-a|$，还有 $y_{k_j}\to a$。连续性于是给出 $f(x_{k_j}),f(y_{k_j})\to f(a)$，两者之差的范数趋于零，矛盾。
\end{proof}

\begin{example}
\label{ex:sine-product-not-uniformly-continuous}
证明 $f(x,y)=\sin(xy)$ 在 $\mathbb R^2$ 的每个紧子集上一致连续，但在整个 $\mathbb R^2$ 上不一致连续。
\end{example}

\begin{solve}
函数处处连续，故由定理~\ref{thm:multivariable-heine-cantor}，它在任意紧子集上一致连续。取
\[
    p_k=(k,1/k),\qquad q_k=(k,2/k),
\]
则 $|p_k-q_k|=1/k\to0$，但 $|f(p_k)-f(q_k)|=|\sin1-\sin2|>0$，不满足一致连续性。严格地取 $\varepsilon_0=|\sin1-\sin2|/2$，任意 $\delta>0$ 下都可选 $k>1/\delta$ 得到违反定义的点对。该例表明，“在每个紧子集上一致连续”不等于“在整个空间上一致连续”。
\end{solve}

\subsection{同胚映射}
\label{subsec:homeomorphisms}

\begin{definition}
\label{def:euclidean-homeomorphism}
设 $E\subseteq\mathbb R^n$，$F\subseteq\mathbb R^m$。若 $f:E\to F$ 是双射，且 $f$ 与逆映射 $f^{-1}:F\to E$ 都连续，则称 $f$ 为 $E$ 到 $F$ 的\emph{同胚映射}（homeomorphism），称 $E,F$ 同胚。两映射的连续性均相对于各自定义域；教材也将同胚映射称为变换，将其逆映射称为逆变换。
\end{definition}

同胚在两个方向都保持连续性。由连续像的性质，同胚集合同时为紧集或同时不为紧集，也同时道路连通或同时不道路连通；同胚的限制仍是相应子集之间的同胚。这些性质可用于排除某些看似可能存在的变换。

\begin{example}
\label{ex:real-line-plane-not-homeomorphic}
证明 $\mathbb R$ 与 $\mathbb R^2$ 不同胚。
\end{example}

\begin{solve}
若存在同胚 $f:\mathbb R\to\mathbb R^2$，删去点 $0$ 及其像 $p=f(0)$，得到同胚
\[
    f:\mathbb R\setminus\{0\}\longrightarrow\mathbb R^2\setminus\{p\}.
\]
平面删去一点后仍道路连通：给定 $u,v\neq p$，若线段 $[u,v]$ 不经过 $p$，可直接连接；否则选不在直线 $pu$ 与 $pv$ 上的点 $w$，折线 $u\to w\to v$ 就避开 $p$。但实直线删去零点后不道路连通，因为任何从 $-1$ 到 $1$ 的连续道路由介值定理都要经过 $0$。于是连续逆映射会把道路连通的 $\mathbb R^2\setminus\{p\}$ 映成非道路连通的 $\mathbb R\setminus\{0\}$，与定理~\ref{thm:continuous-image-path-connected} 矛盾。
\end{solve}

区间上的连续单射实值函数严格单调，因而其反函数连续；一般集合间的连续双射则未必是同胚。考虑例~\ref{ex:polar-map-boundary-limits} 的极坐标映射
\[
    P:D=(0,+\infty)\times[0,2\pi)\longrightarrow
          \mathbb R^2\setminus\{0\},\qquad
    P(r,\theta)=(r\cos\theta,r\sin\theta).
\]
它连续且为双射，逆映射的半径分量 $r=\sqrt{x^2+y^2}$ 连续，但角度取在 $[0,2\pi)$ 时为
\[
    \theta(x,y)=\begin{cases}
        \arctan(y/x),&x>0,\ y\geq0,\\
        \pi/2,&x=0,\ y>0,\\
        \pi+\arctan(y/x),&x<0,\\
        3\pi/2,&x=0,\ y<0,\\
        2\pi+\arctan(y/x),&x>0,\ y<0.
    \end{cases}
\]
固定 $r_0>0$，取 $z_k=(r_0\cos(1/k),-r_0\sin(1/k))\to(r_0,0)$，则
\[
    P^{-1}(z_k)=(r_0,2\pi-1/k)\to(r_0,2\pi)
       \neq P^{-1}(r_0,0)=(r_0,0).
\]
这里左侧收敛是在环境空间 $\mathbb R^2$ 中讨论的，极限参数点不属于 $D$；它仍足以说明逆映射在正实轴上不连续，因此这个连续双射不是同胚。

若改取开参数域 $D_0=(0,+\infty)\times(0,2\pi)$，则其像为
\[
    \Omega_0=\mathbb R^2\setminus\{(x,y):x\geq0,\ y=0\}.
\]
删去这一条射线后，每个像点附近都可选取无跳跃的角度分支：$x<0$ 时用 $\pi+\arctan(y/x)$，$x>0$ 时按 $y$ 的正负使用上述对应分支；在 $y>0$ 或 $y<0$ 且靠近 $x=0$ 时，也可分别使用 $\pi/2-\arctan(x/y)$ 或 $3\pi/2-\arctan(x/y)$。因此逆映射的角度与半径均连续，$P:D_0\to\Omega_0$ 就是同胚。需要检查的是角度分支在所选定义域内是否连续，不能仅凭坐标公式可逆便断言同胚。

\section*{习题十三}
\addcontentsline{toc}{section}{习题十三}
\fi

\selectchapter{14}
\ifcompilechapter
\chapter{多元微分学}
\label{chap:multivariable-differential-calculus}

% 本章依据教材第 40--120 页（分章 PDF 第 1--81 页）整理；第 120 页起的习题十四不在本次范围内。
一元函数的导数同时描述变化率与局部线性近似；自变量具有多个分量时，沿一条直线观察到的变化还不足以控制所有趋近方式，因此必须区分偏导数、方向导数与全微分。本章先用线性映射建立可微性的统一语言，再由链式法则和泰勒公式研究隐函数、局部逆映射、极值与切空间。除另有说明外，$D\subseteq\mathbb R^n$ 是开集，所论点均为其内点，向量采用列向量，$\|\cdot\|$ 表示欧氏范数，$\boldsymbol e_i$ 表示第 $i$ 个标准单位向量，$B(\boldsymbol a,r)=\{\boldsymbol x:\|\boldsymbol x-\boldsymbol a\|<r\}$ 表示开球；函数在区域上为常数的结论还要求该区域连通。

\section{偏导数与全微分}
\label{sec:partial-derivatives-and-total-differentials}

% 教材第 40--57 页；分章 PDF 第 1--18 页。

\subsection{偏导数}
\label{subsec:partial-derivatives}

\begin{definition}
\label{def:multivariable-partial-derivative}
设 $f:D\to\mathbb R$，$\boldsymbol x_0=(x_{01},\ldots,x_{0n})^{\mathsf T}\in D$。若有限极限
\[
    \frac{\partial f}{\partial x_i}(\boldsymbol x_0)
    =\lim_{t\to0}\frac{f(\boldsymbol x_0+t\boldsymbol e_i)-f(\boldsymbol x_0)}{t}
\]
存在，则称 $f$ 在 $\boldsymbol x_0$ 处关于 $x_i$ 可偏导，该极限称为相应的\emph{偏导数}（partial derivative），也记为 $f_{x_i}(\boldsymbol x_0)$。若它在 $D$ 的每一点都存在，就得到定义在 $D$ 上的偏导函数 $f_{x_i}$。
\end{definition}

求偏导数就是固定其余自变量，只对一个自变量求一元导数。对于曲面 $z=f(x,y)$，$f_x(x_0,y_0)$ 是截线 $z=f(x,y_0)$ 在相应点的切线斜率，$f_y(x_0,y_0)$ 则对应另一条坐标截线。\emph[caution]{各个偏导数只考察有限个坐标方向，其存在不能保证函数连续。}例如，令 $f(x,y)$ 在 $xy\ne0$ 时为 $1$，在 $xy=0$ 时为 $0$，则 $f_x(0,0)=f_y(0,0)=0$，但沿 $y=x\ne0$ 趋近原点时函数值恒为 $1$。

\begin{example}
\label{ex:partial-derivative-computation-ch14}
求 $f(x,y)=\tan(x^2/y)$ 在 $(0,1)$ 处的两个偏导数，并求 $g(x,y)=\arcsin\bigl(x/\sqrt{x^2+y^2}\bigr)$ 在 $y\ne0$ 时的偏导数。
\end{example}
\begin{solve}
第一个函数在 $y\ne0$ 且正切有定义的点附近可按一元法则计算：
\[
    f_x=\frac{2x}{y}\sec^2\frac{x^2}{y},
    \qquad f_y=-\frac{x^2}{y^2}\sec^2\frac{x^2}{y}.
\]
故 $f_x(0,1)=f_y(0,1)=0$。第二个函数中，令 $r=\sqrt{x^2+y^2}$。由于 $\sqrt{1-x^2/r^2}=|y|/r$，在 $y\ne0$ 时有
\[
    g_x=\frac{r}{|y|}\frac{y^2}{r^3}
        =\frac{|y|}{x^2+y^2},
    \qquad
    g_y=-\frac{r}{|y|}\frac{xy}{r^3}
        =-\frac{x\operatorname{sgn}y}{x^2+y^2}.
\]
这里绝对值来自平方根，不能将 $\sqrt{y^2}$ 直接写成 $y$。
\end{solve}

\subsection{方向导数}
\label{subsec:directional-derivatives}

\begin{definition}
\label{def:one-sided-directional-derivative-ch14}
设 $f$ 在 $\boldsymbol x_0$ 的邻域内有定义，$\boldsymbol v=(v_1,\ldots,v_n)^{\mathsf T}$ 为单位向量。若有限极限
\begin{equation}
\label{eq:one-sided-directional-derivative-ch14}
    D_{\boldsymbol v}f(\boldsymbol x_0)
    =\lim_{t\to0+}\frac{f(\boldsymbol x_0+t\boldsymbol v)-f(\boldsymbol x_0)}{t}
\end{equation}
存在，则称其为 $f$ 在 $\boldsymbol x_0$ 沿 $\boldsymbol v$ 的\emph{方向导数}（directional derivative），也记为 $\partial f/\partial\boldsymbol v$。教材采用 $t\to0+$ 的约定，因此它表示沿指定射线、按单位距离计算的变化率。$v_i=\cos\theta_i$ 称为该方向的方向余弦，满足 $\sum_i\cos^2\theta_i=1$。
\end{definition}

有些文献改用 $t\to0$ 定义方向导数；两种约定在可微函数上相同，在一般函数上须加区分。若 $f_{x_i}(\boldsymbol x_0)$ 存在，则
\[
    D_{\boldsymbol e_i}f(\boldsymbol x_0)=f_{x_i}(\boldsymbol x_0),
    \qquad D_{-\boldsymbol e_i}f(\boldsymbol x_0)=-f_{x_i}(\boldsymbol x_0).
\]
反过来，两个单侧方向导数都存在且互为相反数，才可推出相应的双侧偏导数存在。对于非单位向量 $\boldsymbol w\ne0$，差商 $[f(\boldsymbol x_0+t\boldsymbol w)-f(\boldsymbol x_0)]/t$ 还包含速度因子 $\|\boldsymbol w\|$，与按单位距离定义的方向导数相差这一倍数。

\begin{example}
\label{ex:directional-derivatives-without-continuity-ch14}
说明所有方向导数存在既不保证各个偏导数存在，也不保证函数连续。
\end{example}
\begin{solve}
对于 $f(x,y)=\sqrt{x^2+y^2}$，任意单位向量 $\boldsymbol v$ 都满足
\[
    D_{\boldsymbol v}f(0,0)=\lim_{t\to0+}\frac{t}{t}=1.
\]
然而沿 $x$ 轴的双侧差商为 $|t|/t$，故 $f_x(0,0)$ 不存在，$f_y(0,0)$ 同样不存在。

再令 $g(x,y)$ 在 $y=x^2$ 且 $x\ne0$ 时为 $1$，其余点为 $0$。任一固定射线在充分接近原点时不再与这条去掉原点的抛物线相交：若 $\boldsymbol v=(a,b)$ 且 $a\ne0$，非零交点的参数必须满足 $t=b/a^2$，至多一个；若 $a=0$，则没有非零交点。因此 $D_{\boldsymbol v}g(0,0)=0$ 对所有方向成立。但 $g(s,s^2)=1$（$s\ne0$），故 $g$ 在原点不连续。这里每个方向所需的邻域可以不同，不能将逐方向的信息当成对所有趋近方式的统一控制。
\end{solve}

\subsection{全微分}
\label{subsec:total-differentials}

\begin{definition}
\label{def:total-differentiability-ch14}
设 $f:D\to\mathbb R$。若存在与增量 $\boldsymbol h$ 无关的常数 $a_1,\ldots,a_n$，使
\begin{equation}
\label{eq:total-differentiability-ch14}
    f(\boldsymbol x_0+\boldsymbol h)-f(\boldsymbol x_0)
    =\sum_{i=1}^n a_i h_i+r(\boldsymbol h),
    \qquad \frac{r(\boldsymbol h)}{\|\boldsymbol h\|}\longrightarrow0
    \quad(\boldsymbol h\to\boldsymbol0),
\end{equation}
则称 $f$ 在 $\boldsymbol x_0$ \emph{可微}（differentiable）。线性映射
$\mathrm df_{\boldsymbol x_0}(\boldsymbol h)=\sum_i a_i h_i$ 称为其\emph{全微分}（total differential）；令 $\mathrm dx_i=h_i$，也写为 $\mathrm df=\sum_i a_i\,\mathrm dx_i$。它是函数全增量的一阶线性主要部分。
\end{definition}

\begin{theorem}
\label{thm:differentiability-implies-partials-ch14}
若 $f$ 在 $\boldsymbol x_0$ 可微，则它在该点连续，各个偏导数都存在，且式~\eqref{eq:total-differentiability-ch14} 中的系数唯一，满足 $a_i=f_{x_i}(\boldsymbol x_0)$。
\end{theorem}
\begin{proof}
由柯西不等式，$|\sum_i a_i h_i|\leq(\sum_i a_i^2)^{1/2}\|\boldsymbol h\|$，故全增量趋于零。再令 $\boldsymbol h=t\boldsymbol e_i$，在增量公式两边除以 $t\ne0$；由于 $r(t\boldsymbol e_i)/t\to0$，取极限得 $f_{x_i}(\boldsymbol x_0)=a_i$，也就证明了系数的唯一性。
\end{proof}

若各个偏导数已存在，判定可微性仍须验证
\[
    \frac{f(\boldsymbol x_0+\boldsymbol h)-f(\boldsymbol x_0)
       -\sum_i f_{x_i}(\boldsymbol x_0)h_i}{\|\boldsymbol h\|}\longrightarrow0.
\]
仅仅写出形式上的 $\sum_i f_{x_i}\,\mathrm dx_i$ 不等于已经证明它是全微分。实际计算中，以下充分条件最常使用。

\begin{theorem}
\label{thm:continuous-partials-imply-differentiability-ch14}
若 $f$ 在 $\boldsymbol x_0$ 的某邻域内存在各个偏导数，且这些偏导数都在 $\boldsymbol x_0$ 连续，则 $f$ 在 $\boldsymbol x_0$ 可微。若所有一阶偏导数在开集 $D$ 上连续，称 $f$ 在 $D$ 上\emph{连续可微}（continuously differentiable），记为 $f\in C^1(D)$。
\end{theorem}
\begin{proof}
取完全包含在 $D$ 内的一个小球，使下面依次改变坐标的线段都在其中。令
$\boldsymbol p_0=\boldsymbol x_0$，$\boldsymbol p_i=\boldsymbol x_0+\sum_{j=1}^i h_j\boldsymbol e_j$。逐段应用一元拉格朗日中值定理，有
\[
    f(\boldsymbol x_0+\boldsymbol h)-f(\boldsymbol x_0)
    =\sum_{i=1}^n f_{x_i}(\boldsymbol\xi_i)h_i,
\]
其中 $\boldsymbol\xi_i$ 在 $\boldsymbol p_{i-1}$ 与 $\boldsymbol p_i$ 之间；当 $h_i=0$ 时相应项直接取零。每个 $\boldsymbol\xi_i\to\boldsymbol x_0$，故
\[
    \left|\sum_i\bigl(f_{x_i}(\boldsymbol\xi_i)-f_{x_i}(\boldsymbol x_0)\bigr)h_i\right|
    \leq\sqrt n\max_i|f_{x_i}(\boldsymbol\xi_i)-f_{x_i}(\boldsymbol x_0)|\,
       \|\boldsymbol h\|
    =o(\|\boldsymbol h\|).
\]
这正是可微性的定义。
\end{proof}

\begin{example}
\label{ex:differentiable-discontinuous-partials-ch14}
可微函数的偏导数不必连续。令
\[
    f(x,y)=
    \begin{cases}
        (x^2+y^2)\sin\dfrac1{x^2+y^2},&(x,y)\ne(0,0),\\
        0,&(x,y)=(0,0).
    \end{cases}
\]
证明 $f$ 在原点可微，但两个偏导数在该点都不连续。
\end{example}
\begin{solve}
记 $r=\sqrt{x^2+y^2}$，则 $|f(x,y)|/r\leq r\to0$，所以原点处的全微分为零。原点外
\[
    f_x=2x\sin\frac1{r^2}-\frac{2x}{r^2}\cos\frac1{r^2}.
\]
取 $x_k=(2\pi k)^{-1/2}$、$y_k=0$，有 $f_x(x_k,0)=-2/x_k\to-\infty$，而 $f_x(0,0)=0$。交换 $x,y$ 即得 $f_y$ 的不连续性。
\end{solve}

\begin{proposition}
\label{prop:directional-derivative-of-differentiable-function-ch14}
若 $f$ 在 $\boldsymbol x_0$ 可微，则每个单位方向的方向导数都存在，且
\begin{equation}
\label{eq:directional-derivative-formula-ch14}
    D_{\boldsymbol v}f(\boldsymbol x_0)
    =\sum_{i=1}^n f_{x_i}(\boldsymbol x_0)v_i.
\end{equation}
\end{proposition}
\begin{proof}
在可微增量式中代入 $\boldsymbol h=t\boldsymbol v$，除以 $t>0$，再令 $t\to0+$ 即得结论，因为 $\|t\boldsymbol v\|=t$。
\end{proof}

\begin{example}
\label{ex:total-differential-approximation-ch14}
利用全微分估计 $2.99^2\times1.02^3$。
\end{example}
\begin{solve}
取 $f(x,y)=x^2y^3$，在 $(3,1)$ 处有 $f=9$、$f_x=6$、$f_y=27$。因此
\[
    f(2.99,1.02)\approx9+6(-0.01)+27(0.02)=9.48.
\]
这是舍去高于一阶增量所得的近似；可微性说明误差为 $o(\sqrt{(\Delta x)^2+(\Delta y)^2})$，给定增量下若要具体误差界，还需利用高阶导数估计余项。
\end{solve}

\subsection{梯度}
\label{subsec:gradient}

\begin{definition}
\label{def:gradient-ch14}
若 $f$ 在 $\boldsymbol x_0$ 可微，称列向量
\[
    \nabla f(\boldsymbol x_0)
    =\operatorname{grad}f(\boldsymbol x_0)
    =\bigl(f_{x_1}(\boldsymbol x_0),\ldots,f_{x_n}(\boldsymbol x_0)\bigr)^{\mathsf T}
\]
为其\emph{梯度}（gradient）。于是 $\mathrm df_{\boldsymbol x_0}(\boldsymbol h)=\nabla f(\boldsymbol x_0)^{\mathsf T}\boldsymbol h$，且 $D_{\boldsymbol v}f=\nabla f^{\mathsf T}\boldsymbol v$。
\end{definition}

由柯西不等式，当 $\nabla f(\boldsymbol x_0)\ne\boldsymbol0$ 时，方向导数的最大值为 $\|\nabla f(\boldsymbol x_0)\|$，唯一取最大值的单位方向是 $\nabla f/\|\nabla f\|$；最小值为其相反数，对应负梯度方向。与梯度正交的单位方向恰为方向导数为零的方向。若梯度为零，则所有方向导数都为零，却不能据此断言函数在该点取得极值。对于可微的 $f,g$ 与常数 $a,b$，逐分量使用一元求导法则可得
\[
    \nabla(af+bg)=a\nabla f+b\nabla g,\qquad
    \nabla(fg)=f\nabla g+g\nabla f,\qquad
    \nabla\frac fg=\frac{g\nabla f-f\nabla g}{g^2}\quad(g\ne0).
\]

\begin{example}
\label{ex:continuous-all-directional-not-differentiable-ch14}
设整数 $k\geq2$，用极坐标定义 $f(0,0)=0$，$f(r\cos\theta,r\sin\theta)=r\sin(k\theta)$（$r>0$）。证明该函数在原点连续，所有方向导数都存在，但不可微。
\end{example}
\begin{solve}
由 $|f(x,y)|\leq\sqrt{x^2+y^2}$，函数在原点连续；沿方向 $\boldsymbol v=(\cos\theta,\sin\theta)$ 有 $D_{\boldsymbol v}f(0,0)=\sin(k\theta)$。若函数在原点可微，此方向导数应等于某个固定向量与 $\boldsymbol v$ 的内积。该向量不为零，因为方向导数取到 $1$，而非零向量的内积在单位圆上只有一个最大值方向；$\sin(k\theta)$ 却有 $k$ 个不同的最大值方向，产生矛盾。特别地，$k=2$ 时原点外 $f(x,y)=2xy/\sqrt{x^2+y^2}$，原点的两个偏导数都为零；它同时说明连续且各个偏导数存在仍不足以保证可微。
\end{solve}

\subsection{向量函数的导数与全微分}
\label{subsec:derivatives-and-total-differentials-of-vector-functions}

\begin{definition}
\label{def:frechet-derivative-and-jacobian-ch14}
设 $F:D\subseteq\mathbb R^n\to\mathbb R^m$。若存在 $m\times n$ 矩阵 $A$，使
\[
    F(\boldsymbol x_0+\boldsymbol h)-F(\boldsymbol x_0)
    =A\boldsymbol h+R(\boldsymbol h),\qquad
    \frac{\|R(\boldsymbol h)\|}{\|\boldsymbol h\|}\longrightarrow0,
\]
则称 $F$ 在 $\boldsymbol x_0$ 可微或可导，$A$ 为其\emph{弗雷歇导数}（Fr\'echet derivative），记为 $DF(\boldsymbol x_0)$ 或 $F'(\boldsymbol x_0)$；其全微分为 $\mathrm dF_{\boldsymbol x_0}(\boldsymbol h)=DF(\boldsymbol x_0)\boldsymbol h$。标量函数是 $m=1$ 的特例，此时 $Df=\nabla f^{\mathsf T}$ 为行矩阵。
\end{definition}

\begin{theorem}
\label{thm:componentwise-vector-differentiability-ch14}
向量函数 $F=(F_1,\ldots,F_m)^{\mathsf T}$ 在一点可微，当且仅当其每个分量函数在该点可微。此时
\begin{equation}
\label{eq:jacobian-matrix-ch14}
    DF(\boldsymbol x_0)=J_F(\boldsymbol x_0)
    =\left(\frac{\partial F_i}{\partial x_j}(\boldsymbol x_0)\right)_{
        \substack{1\leq i\leq m\\1\leq j\leq n}}.
\end{equation}
该矩阵称为\emph{雅可比矩阵}（Jacobian matrix）；当 $m=n$ 时，$\det DF$ 称为\emph{雅可比行列式}（Jacobian determinant），也记为 $\partial(F_1,\ldots,F_n)/\partial(x_1,\ldots,x_n)$。
\end{theorem}
\begin{proof}
若向量余项为 $o(\|\boldsymbol h\|)$，则每个分量的绝对值不超过其范数，故各分量可微。反之，若各分量的余项为 $r_i(\boldsymbol h)=o(\|\boldsymbol h\|)$，则
\[
    \frac{\|R(\boldsymbol h)\|}{\|\boldsymbol h\|}
    =\left(\sum_{i=1}^m\frac{|r_i(\boldsymbol h)|^2}{\|\boldsymbol h\|^2}\right)^{1/2}
    \longrightarrow0.
\]
将各分量全微分按行排列即得矩阵表达式。
\end{proof}

所有偏导数存在时可以形式地写出 $J_F$，但只有可微时它才表示函数增量的线性主要部分。若各分量都属于 $C^1(D)$，就记 $F\in C^1(D,\mathbb R^m)$。当 $F$ 是两个开集之间的双射，且 $F$ 与 $F^{-1}$ 都连续可微时，称为 $C^1$ 微分同胚（$C^1$ diffeomorphism）；教材也称之为 $C^1$ 变换，单有 $F\in C^1$ 不包含逆映射的存在。

\begin{example}
\label{ex:spherical-jacobian-ch14}
对球坐标映射 $F(r,\varphi,\theta)=(r\sin\varphi\cos\theta,r\sin\varphi\sin\theta,r\cos\varphi)^{\mathsf T}$，求其导数与雅可比行列式，其中 $r>0$、$0<\varphi<\pi$、$0<\theta<2\pi$。
\end{example}
\begin{solve}
各分量均光滑，故
\[
    DF=
    \begin{pmatrix}
    \sin\varphi\cos\theta&r\cos\varphi\cos\theta&-r\sin\varphi\sin\theta\\
    \sin\varphi\sin\theta&r\cos\varphi\sin\theta&r\sin\varphi\cos\theta\\
    \cos\varphi&-r\sin\varphi&0
    \end{pmatrix},\qquad
    \det DF=r^2\sin\varphi.
\]
三列分别对应 $r,\varphi,\theta$ 的变化；行列式的符号依赖于这一次序，交换两个变量会改变符号。
\end{solve}

\section{多元函数求导法}
\label{sec:multivariable-differentiation-rules}

% 教材第 57--74 页；分章 PDF 第 18--35 页。

\subsection{导数的四则运算}
\label{subsec:arithmetic-rules-for-multivariable-derivatives}

若标量函数 $f,g$ 在同一点可微，则和、差、积仍可微，且在 $g\ne0$ 时商也可微：
\[
    D(f\pm g)=Df\pm Dg,\qquad
    D(fg)=f\,Dg+g\,Df,\qquad
    D\!\left(\frac fg\right)=\frac{g\,Df-f\,Dg}{g^2}.
\]
这些结论既可由增量公式证明，也可先逐分量计算偏导数，再核实余项。以积为例，展开 $(f+\Delta f)(g+\Delta g)-fg$，其中 $\Delta f\Delta g=O(\|\boldsymbol h\|^2)=o(\|\boldsymbol h\|)$，其余项便给出乘积法则；商法则由 $D(1/g)=-Dg/g^2$ 与乘积法则得到，分母在所论点非零及其连续性保证商在一个邻域内有定义。

对于标量函数 $a$ 与列向量函数 $F:D\to\mathbb R^m$，有
\[
    D(aF)=a\,DF+F\,Da.
\]
这里 $F$ 为 $m\times1$ 矩阵、$Da$ 为 $1\times n$ 矩阵，$F\,Da$ 为外积。若 $F,G:D\to\mathbb R^m$ 均可微，则内积的求导法则为
\[
    D(F^{\mathsf T}G)=G^{\mathsf T}DF+F^{\mathsf T}DG.
\]
这可由 $F^{\mathsf T}G=\sum_iF_iG_i$ 逐项求导得到；矩阵乘法须保持次序，不能直接照抄标量乘法中的交换操作。

\subsection{复合函数的求导法}
\label{subsec:multivariable-chain-rule}

\begin{theorem}[链式法则]
\label{thm:multivariable-chain-rule-ch14}
设 $U:D\subseteq\mathbb R^n\to\mathbb R^m$ 在 $\boldsymbol x_0$ 可微，$F:\Omega\subseteq\mathbb R^m\to\mathbb R^p$ 在 $\boldsymbol u_0=U(\boldsymbol x_0)$ 可微，且 $U$ 将 $\boldsymbol x_0$ 的某邻域映入 $\Omega$。则 $F\circ U$ 在 $\boldsymbol x_0$ 可微，且
\begin{equation}
\label{eq:multivariable-chain-rule-ch14}
    D(F\circ U)(\boldsymbol x_0)
    =DF(\boldsymbol u_0)\,DU(\boldsymbol x_0).
\end{equation}
右边是 $p\times m$ 矩阵乘以 $m\times n$ 矩阵，结果为 $p\times n$ 矩阵。
\end{theorem}
\begin{proof}
记 $A=DU(\boldsymbol x_0)$、$B=DF(\boldsymbol u_0)$。由可微性，
\[
    \boldsymbol k=U(\boldsymbol x_0+\boldsymbol h)-U(\boldsymbol x_0)
        =A\boldsymbol h+r(\boldsymbol h),\qquad
    F(\boldsymbol u_0+\boldsymbol k)-F(\boldsymbol u_0)
        =B\boldsymbol k+s(\boldsymbol k),
\]
其中 $r(\boldsymbol h)=o(\|\boldsymbol h\|)$、$s(\boldsymbol k)=o(\|\boldsymbol k\|)$，规定 $s(\boldsymbol0)=\boldsymbol0$。第一个等式给出 $\|\boldsymbol k\|=O(\|\boldsymbol h\|)$，故即使部分增量使 $\boldsymbol k=\boldsymbol0$，仍有 $s(\boldsymbol k)=o(\|\boldsymbol h\|)$。代入第二式得
\[
    F(U(\boldsymbol x_0+\boldsymbol h))-F(U(\boldsymbol x_0))
    =BA\boldsymbol h+Br(\boldsymbol h)+s(\boldsymbol k)
    =BA\boldsymbol h+o(\|\boldsymbol h\|),
\]
从而得到结论。
\end{proof}

当外层 $f$ 为标量函数时，链式法则的分量形式为
\begin{equation}
\label{eq:partial-chain-rule-ch14}
    \frac{\partial(f\circ U)}{\partial x_j}(\boldsymbol x_0)
    =\sum_{i=1}^m
        f_{u_i}(\boldsymbol u_0)\frac{\partial U_i}{\partial x_j}(\boldsymbol x_0).
\end{equation}
若只求关于 $x_j$ 的一个偏导数，可以将内层的可微性减弱为 $U$ 关于 $x_j$ 的各个分量偏导数存在：将 $U$ 限制在直线 $\boldsymbol x_0+t\boldsymbol e_j$ 上，再应用一元参数的链式法则即可。\emph[caution]{外层函数可微这一条件不能仅用各个偏导数存在代替。}例如，取前述在 $uv\ne0$ 时等于 $1$、其余点等于 $0$ 的函数 $f(u,v)$，虽然原点处的两个偏导数都为零，但 $f(t,t)$ 在 $t=0$ 不连续，不能套用链式法则。

若 $F$ 为 $C^1$ 微分同胚，令 $\boldsymbol y=F(\boldsymbol x)$，对 $F^{-1}\circ F$ 求导可得
\begin{equation}
\label{eq:inverse-derivative-ch14}
    D(F^{-1})(\boldsymbol y)=[DF(\boldsymbol x)]^{-1},\qquad
    \det D(F^{-1})(\boldsymbol y)\det DF(\boldsymbol x)=1.
\end{equation}
这是已知局部逆映射可微之后的求导公式，逆映射是否存在还需后面的存在定理。

\begin{example}
\label{ex:chain-rule-explicit-abstract-ch14}
设 $u=e^{x^2+y^2+z^2}$，$z=x^2\sin y$，求复合后的 $u_x,u_y$；再设 $w=f(xy,y/x,z^2)$（$x\ne0$），其中 $f$ 可微，写出 $w$ 的各个偏导数。
\end{example}
\begin{solve}
第一个问题中，求导既要计入 $x,y$ 的直接影响，也要计入经 $z$ 传递的影响：
\[
    u_x=(2x+2z\cdot2x\sin y)e^{x^2+y^2+z^2}
        =2x(1+2x^2\sin^2y)e^{x^2+y^2+z^2},
\]
\[
    u_y=(2y+2z\cdot x^2\cos y)e^{x^2+y^2+z^2}
        =(2y+x^4\sin2y)e^{x^2+y^2+z^2},
\]
式中的 $z$ 均须取 $x^2\sin y$。第二个问题记 $f_i$ 为外层 $f$ 关于第 $i$ 个自变量的偏导数，均在 $(xy,y/x,z^2)$ 处取值，则
\[
    w_x=yf_1-\frac{y}{x^2}f_2,\qquad
    w_y=xf_1+\frac1x f_2,\qquad
    w_z=2zf_3.
\]
按中间变量逐条列出影响路径，可以避免遗漏直接项或间接项。
\end{solve}

\begin{example}
\label{ex:polar-gradient-and-transport-ch14}
用极坐标表示平面函数梯度的模，并用变量替换求解整个 $(x,t)$ 平面上的一阶方程 $u_t+\lambda u_x=0$，其中 $\lambda\in\mathbb R$ 为常数，$u\in C^1(\mathbb R^2)$。
\end{example}
\begin{solve}
令 $\widetilde f(r,\theta)=f(r\cos\theta,r\sin\theta)$，$r>0$。链式法则给出
\[
    \widetilde f_r=f_x\cos\theta+f_y\sin\theta,\qquad
    \widetilde f_\theta=-rf_x\sin\theta+rf_y\cos\theta.
\]
解此线性方程组得
\[
    f_x=\cos\theta\,\widetilde f_r-\frac{\sin\theta}{r}\widetilde f_\theta,
    \qquad
    f_y=\sin\theta\,\widetilde f_r+\frac{\cos\theta}{r}\widetilde f_\theta,
\]
从而 $\|\nabla f\|^2=\widetilde f_r^2+r^{-2}\widetilde f_\theta^2$。角度变量只需在局部选择连续分支，上述公式不用于 $r=0$。

对于偏微分方程，令 $\xi=x-\lambda t$、$\tau=t$，并定义 $v(\xi,\tau)=u(\xi+\lambda\tau,\tau)$。于是 $v_\tau=\lambda u_x+u_t=0$。对每个固定 $\xi$，一元导数为零说明 $v$ 与 $\tau$ 无关，故 $u(x,t)=g(x-\lambda t)$，其中 $g\in C^1(\mathbb R)$；直接求导又可验证每个这样的函数都是解。若定义域只是一般开集，该结论首先是局部结论，不能忽略同一特征直线在定义域内可能有多个连通部分。
\end{solve}

\begin{example}
\label{ex:hyperspherical-jacobian-ch14}
设 $n\geq2$。在 $n$ 维球坐标
\[
    x_j=r\left(\prod_{\ell=1}^{j-1}\sin\theta_\ell\right)\cos\theta_j
       \quad(1\leq j\leq n-1),\qquad
    x_n=r\prod_{\ell=1}^{n-1}\sin\theta_\ell
\]
中，证明按 $(r,\theta_1,\ldots,\theta_{n-1})$ 排列的雅可比行列式为
\[
    \frac{\partial(x_1,\ldots,x_n)}{\partial(r,\theta_1,\ldots,\theta_{n-1})}
    =r^{n-1}\prod_{j=1}^{n-2}\sin^{n-1-j}\theta_j,
\]
其中空乘积规定为 $1$。
\end{example}
\begin{solve}
$n=2$ 时直接计算得 $r$。对 $n>2$，先令 $x_1=r\cos\theta_1$、$s=r\sin\theta_1$，再以 $s,\theta_2,\ldots,\theta_{n-1}$ 作其余 $n-1$ 个坐标的球坐标变换。第一步的雅可比行列式为 $r$，第二步由归纳假设为 $s^{n-2}\prod_{j=2}^{n-2}\sin^{n-1-j}\theta_j$，由链式法则取行列式并相乘即得公式。坐标计算通常取 $r>0$，前 $n-2$ 个角在 $(0,\pi)$ 内，最后一个角在一个长度为 $2\pi$ 的开区间内，以避开坐标退化与角度重复。
\end{solve}

\subsection{高阶偏导数}
\label{subsec:higher-partial-derivatives}

偏导函数仍是多元函数，继续求偏导便得到高阶偏导数。以下约定 $f_{x_ix_j}=\partial_{x_j}(\partial_{x_i}f)$，即先对下标左边的变量求导，再对右边求导；用分式写作 $\partial^2f/(\partial x_j\partial x_i)$。在尚未证明可以交换次序前，这两个次序必须区别。若 $f$ 的各阶偏导数直到 $k$ 阶都存在且连续，记为 $f\in C^k(D)$。

\begin{example}
\label{ex:unequal-mixed-partials-ch14}
令 $f(0,0)=0$，原点外 $f(x,y)=xy(x^2-y^2)/(x^2+y^2)$，说明混合偏导数可能不相等。
\end{example}
\begin{solve}
由定义直接计算轴线上的一阶偏导数，得 $f_x(0,y)=-y$，$f_y(x,0)=x$，包括原点处的零值。因此
\[
    \partial_y(\partial_xf)(0,0)=-1,\qquad
    \partial_x(\partial_yf)(0,0)=1.
\]
原点外的表达式可以常规求导，但原点处仍应回到定义，不能把只在原点外成立的有理式直接代入原点。
\end{solve}

\begin{theorem}[混合偏导数交换定理]
\label{thm:schwarz-mixed-partials-ch14}
若 $f$ 在 $\boldsymbol x_0$ 的某邻域内存在 $f_{x_ix_j}$ 与 $f_{x_jx_i}$，且二者在 $\boldsymbol x_0$ 连续，则
$f_{x_ix_j}(\boldsymbol x_0)=f_{x_jx_i}(\boldsymbol x_0)$。特别地，$C^k$ 函数的 $k$ 阶偏导数与求导顺序无关，只依赖各变量被求导的次数。
\end{theorem}
\begin{proof}
固定其他坐标，只证二元情形。对 $h,k\ne0$，令
\[
    Q(h,k)=\frac{f(a+h,b+k)-f(a+h,b)-f(a,b+k)+f(a,b)}{hk}.
\]
先沿 $x$、再沿 $y$ 对差分应用一元中值定理，可将 $Q$ 写成小矩形内某点的 $f_{xy}$；反过来应用，则写成另一个点的 $f_{yx}$。两点都随 $(h,k)\to(0,0)$ 趋于 $(a,b)$，由连续性得到二者在该点相等。对高阶偏导数，可将两个相邻求导操作应用于相应的低阶偏导函数，再利用任意排列都由相邻交换组成，得到最后的结论。
\end{proof}

例如，$C^2$ 函数的二阶偏导数原本有 $n^2$ 个有序写法，交换定理使其中只有 $n(n+1)/2$ 类需要分别计算；这里说的是不同的求导类型，其数值还可能偶然相同。

\subsection{复合函数的高阶偏导数}
\label{subsec:higher-partial-derivatives-of-composite-functions}

高阶链式法则要对一阶公式逐次求导，每次既要微分外层导数，也要微分内层导数。设 $U:D\subseteq\mathbb R^n\to\mathbb R^m$、$f:\Omega\to\mathbb R$ 均为 $C^2$，$u=f\circ U$，则
\begin{equation}
\label{eq:second-order-chain-rule-ch14}
    u_{x_jx_k}
    =\sum_{a,b=1}^m f_{u_au_b}(U)\,(U_a)_{x_j}(U_b)_{x_k}
      +\sum_{a=1}^m f_{u_a}(U)\,(U_a)_{x_jx_k}.
\end{equation}
例如，当 $u(s,t)=f(x(s,t),y(s,t))$ 时，
\[
    u_{ss}=f_{xx}x_s^2+2f_{xy}x_sy_s+f_{yy}y_s^2
        +f_x\,x_{ss}+f_y\,y_{ss}.
\]
最后两项反映内层变换的非线性，只有当相应内层二阶导数为零时才能省略。

\begin{example}
\label{ex:radial-harmonic-function-ch14}
设 $u(x,y,z)=g(r)$，$r=\sqrt{x^2+y^2+z^2}>0$，$g$ 在 $(0,\infty)$ 上二阶连续可微。求使 $u_{xx}+u_{yy}+u_{zz}=0$ 成立的所有 $g$。
\end{example}
\begin{solve}
由链式法则，
\[
    u_x=g'(r)\frac{x}{r},\qquad
    u_{xx}=g''(r)\frac{x^2}{r^2}
        +g'(r)\left(\frac1r-\frac{x^2}{r^3}\right).
\]
将 $y,z$ 对应的两式相加，得到\emph{拉普拉斯算子}（Laplacian）作用于径向函数的表达式
\[
    \Delta u=g''(r)+\frac2r g'(r).
\]
因 $r>0$，方程等价于 $(r^2g'(r))'=0$。积分得 $g(r)=A+B/r$，其中 $A,B$ 为任意实常数。若另外要求该函数在整个原点邻域连续，则必须 $B=0$；原题在 $r>0$ 上并不要求这一点。
\end{solve}

\subsection{一阶微分的形式不变性与高阶微分}
\label{subsec:invariance-and-higher-multivariable-differentials}

设 $f\in C^k(D)$，各 $x_i$ 为独立变量，并将增量 $\mathrm dx_i$ 视为固定量。连续微分得到
\begin{equation}
\label{eq:higher-total-differentials-ch14}
    \mathrm d^k f(\boldsymbol x)
    =\left(\sum_{i=1}^n\mathrm dx_i\frac{\partial}{\partial x_i}\right)^k
       f(\boldsymbol x).
\end{equation}
二阶时 $\mathrm d^2f=\sum_{i,j}f_{x_ix_j}\,\mathrm dx_i\mathrm dx_j$，二元情形则为
\[
    \mathrm d^kf(x,y)=\sum_{j=0}^k\binom kj
       \frac{\partial^kf}{\partial x^{k-j}\partial y^j}
       (\mathrm dx)^{k-j}(\mathrm dy)^j.
\]
这里使用了混合偏导数交换定理，符号上的乘方表示固定增量方向的微分算子反复作用，不是把函数值求幂。

若 $\boldsymbol u=U(\boldsymbol x)$ 是可微的中间变量，则链式法则给出
\[
    \mathrm d(f\circ U)=Df(U)\,DU\,\mathrm d\boldsymbol x
        =Df(\boldsymbol u)\,\mathrm d\boldsymbol u.
\]
无论 $\boldsymbol u$ 是独立变量还是中间变量，最后的形式相同，这称为\emph{一阶微分的形式不变性}（invariance of the first differential）。二阶时若 $f,U\in C^2$，则
\begin{equation}
\label{eq:noninvariance-second-differential-ch14}
    \mathrm d^2(f\circ U)
    =\sum_{a,b}f_{u_au_b}\,\mathrm du_a\mathrm du_b
       +\sum_a f_{u_a}\,\mathrm d^2u_a.
\end{equation}
中间变量的二阶微分通常不为零，所以高阶微分一般不具有上述形式不变性；若 $U$ 为仿射变换，则这些额外项消失。

\begin{example}
\label{ex:second-partials-by-differentials-ch14}
设 $w(x,y)=f(xy,y/x)$，$x\ne0$，$f\in C^2$，利用二阶微分求 $w$ 的全部二阶偏导数。
\end{example}
\begin{solve}
令 $u=xy$、$v=y/x$，则
\[
    \mathrm du=y\,\mathrm dx+x\,\mathrm dy,\quad
    \mathrm dv=-\frac{y}{x^2}\,\mathrm dx+\frac1x\,\mathrm dy,
    \quad\mathrm d^2u=2\,\mathrm dx\mathrm dy,
\]
\[
    \mathrm d^2v=\frac{2y}{x^3}(\mathrm dx)^2
       -\frac2{x^2}\,\mathrm dx\mathrm dy.
\]
代入式~\eqref{eq:noninvariance-second-differential-ch14}，比较 $(\mathrm dx)^2$、$2\,\mathrm dx\mathrm dy$、$(\mathrm dy)^2$ 的系数，得到
\[
    w_{xx}=y^2f_{11}-\frac{2y^2}{x^2}f_{12}
       +\frac{y^2}{x^4}f_{22}+\frac{2y}{x^3}f_2,
\]
\[
    w_{xy}=xyf_{11}-\frac{y}{x^3}f_{22}+f_1-\frac1{x^2}f_2,
    \qquad
    w_{yy}=x^2f_{11}+2f_{12}+\frac1{x^2}f_{22},
\]
其中外层导数均在 $(xy,y/x)$ 处取值，且 $w_{yx}=w_{xy}$。二阶微分中混合项的系数是混合偏导数的两倍，比较系数时应保留这一因子。
\end{solve}

\section{泰勒公式}
\label{sec:multivariable-taylor-formula}

% 教材第 74--79 页；分章 PDF 第 35--40 页。

多元泰勒公式的基本思路是固定起点和增量，将函数限制在连接两点的线段上，再应用一元泰勒公式。为简化记号，对固定增量 $\boldsymbol h=(h_1,\ldots,h_n)^{\mathsf T}$ 记
\[
    L_{\boldsymbol h}=\sum_{i=1}^n h_i\frac{\partial}{\partial x_i},\qquad
    L_{\boldsymbol h}^{\ell}f(\boldsymbol a)
    =\sum_{i_1,\ldots,i_\ell=1}^n
       h_{i_1}\cdots h_{i_\ell}
       \partial_{x_{i_\ell}}\cdots\partial_{x_{i_1}}f(\boldsymbol a),
\]
并规定 $L_{\boldsymbol h}^0f=f$。该算子中各个 $h_i$ 始终视为常数。

\begin{theorem}[带拉格朗日余项的多元泰勒公式]
\label{thm:multivariable-taylor-lagrange-ch14}
设 $K\geq0$ 为整数，$f\in C^{K+1}(D)$，且线段 $\{\boldsymbol a+t\boldsymbol h:0\leq t\leq1\}$ 包含在 $D$ 内。则存在依赖于 $\boldsymbol a,\boldsymbol h$ 的 $\theta\in(0,1)$，使
\begin{equation}
\label{eq:multivariable-taylor-lagrange-ch14}
    f(\boldsymbol a+\boldsymbol h)
    =\sum_{\ell=0}^K\frac1{\ell!}L_{\boldsymbol h}^{\ell}f(\boldsymbol a)
      +\frac1{(K+1)!}L_{\boldsymbol h}^{K+1}f(\boldsymbol a+\theta\boldsymbol h).
\end{equation}
最后一项称为\emph{拉格朗日余项}（Lagrange remainder）。
\end{theorem}
\begin{proof}
构造 $\phi(t)=f(\boldsymbol a+t\boldsymbol h)$。沿线段的链式法则给出 $\phi^{(\ell)}(t)=L_{\boldsymbol h}^{\ell}f(\boldsymbol a+t\boldsymbol h)$。将一元泰勒公式用于 $\phi(1)$ 在 $t=0$ 处的展开，再代入这些导数即可。
\end{proof}

\begin{corollary}[带佩亚诺余项的多元泰勒公式]
\label{cor:multivariable-taylor-peano-ch14}
若 $K\geq1$ 且 $f\in C^K$ 于 $\boldsymbol a$ 的某邻域，则
\begin{equation}
\label{eq:multivariable-taylor-peano-ch14}
    f(\boldsymbol a+\boldsymbol h)
    =\sum_{\ell=0}^K\frac1{\ell!}L_{\boldsymbol h}^{\ell}f(\boldsymbol a)
       +o(\|\boldsymbol h\|^K)\qquad(\boldsymbol h\to\boldsymbol0).
\end{equation}
最后一项称为\emph{佩亚诺余项}（Peano remainder）。
\end{corollary}
\begin{proof}
将定理用于 $K-1$ 阶展开，再加减 $L_{\boldsymbol h}^Kf(\boldsymbol a)/K!$。所得余项为
\[
    \frac1{K!}\left(
       L_{\boldsymbol h}^Kf(\boldsymbol a+\theta\boldsymbol h)
       -L_{\boldsymbol h}^Kf(\boldsymbol a)\right).
\]
展开后每项都是 $K$ 个 $h_i$ 的乘积乘以一个趋于零的 $K$ 阶偏导数差；因项数有限且 $|h_i|\leq\|\boldsymbol h\|$，余项即为 $o(\|\boldsymbol h\|^K)$。
\end{proof}

泰勒多项式在相应阶数内是唯一的：若两个次数不超过 $K$ 的多项式都给出余项为 $o(\|\boldsymbol h\|^K)$ 的展开，它们的差也满足这一估计；若差有最低次非零齐次部分，沿一个使该部分非零的固定方向代入 $\boldsymbol h=t\boldsymbol v$，并除以相应最低次幂，即产生矛盾。因此计算展开时可以使用代数变形和已知的一元展开，不必逐个求高阶偏导数，但必须控制总次数与余项。

\begin{corollary}
\label{cor:multivariable-mean-value-ch14}
若 $f\in C^1(D)$，连接 $\boldsymbol a,\boldsymbol b$ 的线段在 $D$ 内，则存在 $\theta\in(0,1)$，使
\[
    f(\boldsymbol b)-f(\boldsymbol a)
    =\nabla f(\boldsymbol a+\theta(\boldsymbol b-\boldsymbol a))^{\mathsf T}
       (\boldsymbol b-\boldsymbol a).
\]
若 $D$ 还是连通开集，且各个偏导数在 $D$ 上恒为零，则 $f$ 在 $D$ 上为常数。
\end{corollary}
\begin{proof}
第一式是式~\eqref{eq:multivariable-taylor-lagrange-ch14} 中 $K=0$ 的情形。对于第二个结论，先在 $D$ 内每个小球上应用第一式，得到 $f$ 局部为常数。固定 $\boldsymbol a\in D$ 后，集合 $\{\boldsymbol x\in D:f(\boldsymbol x)=f(\boldsymbol a)\}$ 与其补集都相对开；由连通性，前者只能是整个 $D$。
\end{proof}

上述共同中值点的等式适用于实值函数，不能直接推广为向量函数所有分量共用同一个 $\theta$。例如 $F(t)=(\cos t,\sin t)^{\mathsf T}$ 在 $[0,2\pi]$ 的端点增量为零，但 $F'(t)$ 处处非零，不存在使 $F(2\pi)-F(0)=2\pi F'(\theta)$ 成立的点。

对二阶展开，定义\emph{海塞矩阵}（Hessian matrix，教材称海色矩阵）
\[
    H_f(\boldsymbol a)=\bigl(f_{x_ix_j}(\boldsymbol a)\bigr)_{1\leq i,j\leq n}.
\]
当 $f\in C^2$ 时它是对称矩阵，且
\begin{equation}
\label{eq:second-order-taylor-peano-ch14}
    f(\boldsymbol a+\boldsymbol h)
    =f(\boldsymbol a)+\nabla f(\boldsymbol a)^{\mathsf T}\boldsymbol h
       +\frac12\boldsymbol h^{\mathsf T}H_f(\boldsymbol a)\boldsymbol h
       +o(\|\boldsymbol h\|^2).
\end{equation}
这说明在一阶项消失时，局部增量首先由一个二次型决定，是后面极值判别的依据。

\begin{example}
\label{ex:multivariable-taylor-by-expansion-ch14}
求 $f(x,y)=(x-y)/(x+y)$ 在 $(1,0)$ 附近到二次项的展开，并求 $g(x,y)=xe^{x+y}$ 在原点的所有四阶偏导数。
\end{example}
\begin{solve}
记 $a=x-1$、$b=y$，利用 $(1+s)^{-1}=1-s+s^2+O(s^3)$，得到
\[
    \frac{1+a-b}{1+a+b}
    =(1+a-b)\bigl(1-a-b+(a+b)^2\bigr)
       +o(a^2+b^2)
    =1-2b+2ab+2b^2+o(a^2+b^2).
\]
这里 $|a+b|\leq\sqrt2\sqrt{a^2+b^2}$，保证舍去项确为所需的小量。

第二个函数的四阶泰勒多项式可由
\[
    xe^{x+y}
    =x+x(x+y)+\frac{x}{2}(x+y)^2+\frac{x}{6}(x+y)^3
       +o\bigl((x^2+y^2)^2\bigr)
\]
直接读出。四次齐次部分为 $(x^4+3x^3y+3x^2y^2+xy^3)/6$，故按 $x^4,x^3y,x^2y^2,xy^3,y^4$ 的顺序，相应四阶偏导数在原点的值依次为 $4,3,2,1,0$。读系数时应乘以各变量次数的阶乘，例如 $x^3y$ 的系数应乘 $3!1!$，而非统一乘 $4!$。
\end{solve}

\section{隐函数存在定理}
\label{sec:implicit-function-existence-theorem}

% 教材第 79--94 页；分章 PDF 第 40--55 页。

\subsection{单个方程的情形}
\label{subsec:implicit-function-for-one-equation}

方程 $F(x,y)=0$ 描述一个零点集合。所谓由它确定\emph{隐函数}（implicit function），是指在指定零点附近，该零点集合恰好是一个函数 $y=\varphi(x)$ 的图像。圆周 $x^2+y^2=1$ 在 $(0,1)$ 附近可写成 $y=\sqrt{1-x^2}$，但在 $(1,0)$ 附近不能将所有邻近圆周点唯一地写成 $y$ 关于 $x$ 的函数；能否选定某个变量作为因变量，是一个局部问题。

\begin{theorem}[单个方程的隐函数存在定理]
\label{thm:scalar-implicit-function-ch14}
设 $F(\boldsymbol x,y)$ 在 $(\boldsymbol x_0,y_0)\in\mathbb R^n\times\mathbb R$ 的某邻域内连续，$F_y$ 在该邻域内存在且连续，并且
\[
    F(\boldsymbol x_0,y_0)=0,\qquad F_y(\boldsymbol x_0,y_0)\ne0.
\]
则存在 $\boldsymbol x_0$ 的邻域 $U$ 与 $y_0$ 的开区间邻域 $I$，使对每个 $\boldsymbol x\in U$，方程在 $I$ 中恰有一个解 $y=\varphi(\boldsymbol x)$。函数 $\varphi:U\to I$ 连续，$\varphi(\boldsymbol x_0)=y_0$。若进一步有 $F\in C^1$，则 $\varphi\in C^1(U)$，且
\begin{equation}
\label{eq:scalar-implicit-derivative-ch14}
    \frac{\partial\varphi}{\partial x_i}(\boldsymbol x)
    =-\frac{F_{x_i}(\boldsymbol x,\varphi(\boldsymbol x))}
             {F_y(\boldsymbol x,\varphi(\boldsymbol x))},
    \qquad i=1,\ldots,n.
\end{equation}
\end{theorem}
\begin{proof}
不妨设 $F_y(\boldsymbol x_0,y_0)>0$。由连续性，在一个足够小的乘积邻域内 $F_y$ 恒正。选取 $b>0$，使上下端点仍在其中，则一元中值定理给出
\[
    F(\boldsymbol x_0,y_0-b)<0<F(\boldsymbol x_0,y_0+b).
\]
缩小 $\boldsymbol x$ 的邻域 $U$，由 $F$ 的连续性可使这两个严格不等式对所有 $\boldsymbol x\in U$ 保持成立。对固定的 $\boldsymbol x$，函数 $y\mapsto F(\boldsymbol x,y)$ 在 $I=(y_0-b,y_0+b)$ 上连续且严格增加，由介值定理存在唯一零点 $\varphi(\boldsymbol x)$。

为证连续性，固定 $\overline{\boldsymbol x}\in U$，记 $\overline y=\varphi(\overline{\boldsymbol x})$。对任意充分小的 $\varepsilon>0$，有
\[
    F(\overline{\boldsymbol x},\overline y-\varepsilon)<0
       <F(\overline{\boldsymbol x},\overline y+\varepsilon).
\]
当 $\boldsymbol x$ 充分接近 $\overline{\boldsymbol x}$ 时这两个符号不变，严格单调性便将唯一零点夹在 $(\overline y-\varepsilon,\overline y+\varepsilon)$ 内，得到连续性。

若 $F\in C^1$，固定 $\boldsymbol x\in U$，沿第 $i$ 个坐标作增量 $h$，记 $y=\varphi(\boldsymbol x)$、$\Delta y=\varphi(\boldsymbol x+h\boldsymbol e_i)-\varphi(\boldsymbol x)$。将零点等式之差拆为先改变 $\boldsymbol x$、再改变 $y$ 的两项，并分别应用一元中值定理，得到
\[
    0=F_{x_i}(\boldsymbol x+\theta h\boldsymbol e_i,y+\Delta y)h
       +F_y(\boldsymbol x,y+\eta\Delta y)\Delta y,
       \qquad 0<\theta,\eta<1.
\]
当某个增量为零时相应项直接为零。由已证连续性，$\Delta y\to0$；除以 $h\ne0$ 后取极限，即得式~\eqref{eq:scalar-implicit-derivative-ch14}。右边随 $\boldsymbol x$ 连续，所以所有偏导数连续，$\varphi\in C^1(U)$。
\end{proof}

上述唯一性限于选定的乘积邻域，不排除方程在远处另有解。非零偏导数是一个充分条件，并非隐函数存在的必要条件：例如 $(y-x)^3=0$ 仍唯一确定光滑函数 $y=x$，但在零点集合上 $F_y=0$，定理的导数商式不能直接使用。若 $F\in C^k$（$k\geq1$）且满足非退化条件，则所确定的隐函数也属于 $C^k$；可由式~\eqref{eq:scalar-implicit-derivative-ch14} 反复求导归纳证明，分母在足够小的邻域内始终不为零。

例如，开普勒方程 $E-e\sin E=M$（固定 $0<e<1$）中，关于 $E$ 的导数为 $1-e\cos E\geq1-e>0$，左边严格增加且趋向两端无穷，故每个 $M\in\mathbb R$ 都对应唯一 $E$。局部隐函数拼合成全局光滑函数，且 $\mathrm dE/\mathrm dM=(1-e\cos E)^{-1}$；这里无须先显式解出 $E$。

\begin{example}
\label{ex:cyclic-partial-derivative-ch14}
设 $F(x,y,z)=0$ 在一点附近可分别解出三个变量，且 $F\in C^1$、$F_xF_yF_z\ne0$。解释偏导数的循环关系。
\end{example}
\begin{solve}
在同一曲面点处，由隐函数求导公式，
\[
    \left(\frac{\partial x}{\partial y}\right)_z=-\frac{F_y}{F_x},\quad
    \left(\frac{\partial y}{\partial z}\right)_x=-\frac{F_z}{F_y},\quad
    \left(\frac{\partial z}{\partial x}\right)_y=-\frac{F_x}{F_z}.
\]
所以三个因子之积为 $-1$。此外，保持同一个第三变量不变时，$\left(\partial x/\partial y\right)_z\left(\partial y/\partial x\right)_z=1$。下标标明保持不变的变量；偏导数符号不能脱离这些条件当作普通分数任意约分，否则会错误地将循环乘积算成 $1$。
\end{solve}

\begin{example}
\label{ex:implicit-taylor-expansion-ch14}
证明方程 $-2x+y-x^2+y^2+z+\sin z=0$ 在原点附近唯一确定光滑函数 $z=\varphi(x,y)$，并求其到二次项的泰勒展开。
\end{example}
\begin{solve}
令左边为 $F(x,y,z)$，则 $F(0,0,0)=0$、$F_z(0,0,0)=2\ne0$，故隐函数存在且光滑。因 $\varphi(0,0)=0$ 且可微，令 $r=\sqrt{x^2+y^2}$，有 $\varphi=O(r)$，所以 $\sin\varphi=\varphi+O(r^3)$。代回原方程得
\[
    \varphi(x,y)=x-\frac12y+\frac12x^2-\frac12y^2+o(x^2+y^2).
\]
先由存在定理保证光滑性，再比较多项式系数，才为这种隐式展开提供了依据。
\end{solve}

\subsection{方程组的情形}
\label{subsec:implicit-functions-for-equation-systems}

\begin{theorem}[隐函数组存在定理]
\label{thm:implicit-function-system-ch14}
设 $F=(F_1,\ldots,F_m)^{\mathsf T}$ 是 $(\boldsymbol x_0,\boldsymbol u_0)\in\mathbb R^n\times\mathbb R^m$ 某邻域内的 $C^1$ 映射，且
\[
    F(\boldsymbol x_0,\boldsymbol u_0)=\boldsymbol0,\qquad
    \det D_{\boldsymbol u}F(\boldsymbol x_0,\boldsymbol u_0)\ne0.
\]
则存在邻域 $U\ni\boldsymbol x_0$、$V\ni\boldsymbol u_0$ 及唯一的 $C^1$ 映射 $\varphi:U\to V$，使 $\varphi(\boldsymbol x_0)=\boldsymbol u_0$，且对 $(\boldsymbol x,\boldsymbol u)\in U\times V$，
\[
    F(\boldsymbol x,\boldsymbol u)=\boldsymbol0
       \quad\Longleftrightarrow\quad\boldsymbol u=\varphi(\boldsymbol x).
\]
缩小邻域使 $D_{\boldsymbol u}F$ 处处可逆后，导数为
\begin{equation}
\label{eq:implicit-system-derivative-ch14}
    D\varphi(\boldsymbol x)
    =-[D_{\boldsymbol u}F(\boldsymbol x,\varphi(\boldsymbol x))]^{-1}
       D_{\boldsymbol x}F(\boldsymbol x,\varphi(\boldsymbol x)).
\end{equation}
其中 $D_{\boldsymbol u}F$ 为 $m\times m$ 矩阵，$D_{\boldsymbol x}F$ 与 $D\varphi$ 均为 $m\times n$ 矩阵。
\end{theorem}
\begin{proof}
对方程个数 $m$ 归纳。$m=1$ 已证。设 $m>1$，可逆矩阵至少有一个非零元素，重新排列方程与因变量后可设 $(F_m)_{u_m}(\boldsymbol x_0,\boldsymbol u_0)\ne0$。由单个方程的定理，$F_m=0$ 在局部唯一确定
$u_m=g(\boldsymbol x,\boldsymbol v)$，其中 $\boldsymbol v=(u_1,\ldots,u_{m-1})^{\mathsf T}$，$g\in C^1$。将它代入前 $m-1$ 个方程，得到约化方程组 $H(\boldsymbol x,\boldsymbol v)=\boldsymbol0$。

在基点将原来的因变量雅可比矩阵分块写为
\[
    D_{\boldsymbol u}F=
    \begin{pmatrix}A&\boldsymbol b\\\boldsymbol c^{\mathsf T}&d\end{pmatrix},
    \qquad d\ne0.
\]
由隐函数求导公式 $D_{\boldsymbol v}g=-d^{-1}\boldsymbol c^{\mathsf T}$，可得
\[
    D_{\boldsymbol v}H=A-\boldsymbol b d^{-1}\boldsymbol c^{\mathsf T},\qquad
    \det D_{\boldsymbol u}F=d\det D_{\boldsymbol v}H.
\]
后一等式可用分块消元验证，故约化雅可比行列式非零。由归纳假设，$H=0$ 唯一确定 $C^1$ 函数 $\boldsymbol v=\psi(\boldsymbol x)$，再代入 $g$ 就构造出 $\varphi$。适当缩小乘积邻域后，任何原方程组的解都必须先满足唯一的 $u_m=g(\boldsymbol x,\boldsymbol v)$，进而满足约化方程组，故只能是构造的解，这也证明局部唯一性。

最后对恒等式 $F(\boldsymbol x,\varphi(\boldsymbol x))=\boldsymbol0$ 应用链式法则，得
\[
    D_{\boldsymbol x}F+D_{\boldsymbol u}F\,D\varphi=0.
\]
左乘可逆矩阵 $[D_{\boldsymbol u}F]^{-1}$，即得式~\eqref{eq:implicit-system-derivative-ch14}。
\end{proof}

定理中的行列式对应拟解出的因变量，选择不同变量时必须重新检查相应子矩阵。$F\in C^k$ 时同样有 $\varphi\in C^k$，依据仍是导数公式和可逆矩阵元素的光滑性。计算时通常直接微分原方程组，再解关于未知偏导数的线性方程组，比先求出隐函数的显式表达式更方便。

\begin{example}
\label{ex:implicit-system-computation-ch14}
证明方程组
\[
    x^2+y^2+u^2-v^2=2,\qquad u+v+x+y=0
\]
在 $(1,1,-1,-1)$ 附近确定 $u=u(x,y)$、$v=v(x,y)$，并求其一阶偏导数。
\end{example}
\begin{solve}
基点满足方程，关于 $(u,v)$ 的雅可比行列式为
\[
    \det\begin{pmatrix}2u&-2v\\1&1\end{pmatrix}
    =2(u+v),
\]
在基点等于 $-4$，所以所求隐函数存在且光滑。对原方程组微分并将第一式除以 $2$，得到
\[
    u\,\mathrm du-v\,\mathrm dv=-x\,\mathrm dx-y\,\mathrm dy,
    \qquad \mathrm du+\mathrm dv=-\mathrm dx-\mathrm dy.
\]
分别比较 $\mathrm dx,\mathrm dy$ 的系数并解线性方程组，得
\[
    u_x=-\frac{x+v}{u+v},\quad v_x=\frac{x-u}{u+v},\qquad
    u_y=-\frac{y+v}{u+v},\quad v_y=\frac{y-u}{u+v}.
\]
这些公式中的 $u,v$ 是所确定的隐函数；分母在足够小的邻域内非零，基点处的四个偏导数依次为 $0,-1,0,-1$。
\end{solve}

\subsection{逆映射存在定理}
\label{subsec:inverse-mapping-existence-theorem}

\begin{theorem}[逆映射存在定理]
\label{thm:inverse-function-ch14}
设 $F:D\subseteq\mathbb R^n\to\mathbb R^n$ 属于 $C^1$，且 $\det DF(\boldsymbol x_0)\ne0$。令 $\boldsymbol y_0=F(\boldsymbol x_0)$，则存在开邻域 $U\ni\boldsymbol x_0$、$V\ni\boldsymbol y_0$，使 $F:U\to V$ 是 $C^1$ 微分同胚，其逆映射 $G$ 满足
\[
    DG(\boldsymbol y)=[DF(G(\boldsymbol y))]^{-1}.
\]
若 $F\in C^k$，则 $G\in C^k$。
\end{theorem}
\begin{proof}
将 $F(\boldsymbol x)-\boldsymbol y=\boldsymbol0$ 看作以 $\boldsymbol x$ 为因变量、$\boldsymbol y$ 为自变量的方程组，关于因变量的雅可比矩阵为 $DF(\boldsymbol x_0)$。由隐函数组定理，在乘积邻域 $U_0\times V_0$ 内，方程的解恰为 $\boldsymbol x=G(\boldsymbol y)$，且 $G\in C^1$。缩小 $V_0$ 为开邻域 $V$，使 $G(V)\subseteq U_0$，并令 $U=U_0\cap F^{-1}(V)$。连续性使 $U$ 为开集；方程组的局部唯一性给出 $G(F(\boldsymbol x))=\boldsymbol x$（$\boldsymbol x\in U$），而构造又给出 $F(G(\boldsymbol y))=\boldsymbol y$（$\boldsymbol y\in V$），故二者互为逆映射。对后一恒等式求导便得逆映射导数公式；更高光滑性由隐函数组定理继承。
\end{proof}

若 $\det DF$ 在 $D$ 内处处非零，则 $F$ 是局部微分同胚，因而将开集映为开集；$D$ 连通时，$F(D)$ 还因连续性而连通。但局部可逆不蕴含整体单射。

\begin{example}
\label{ex:local-versus-global-invertibility-ch14}
说明雅可比行列式处处非零不能保证全局逆映射存在，而连续逆映射的存在也不要求雅可比行列式处处非零。
\end{example}
\begin{solve}
映射 $F(x,y)=(e^x\cos y,e^x\sin y)^{\mathsf T}$ 的雅可比行列式为 $e^{2x}>0$，但 $F(x,y+2\pi)=F(x,y)$，所以它虽处处局部可逆，却不是从 $\mathbb R^2$ 到其像的双射。另一方面，$H(x,y)=(x^3,y)^{\mathsf T}$ 是 $\mathbb R^2$ 到自身的同胚，逆映射为 $(u,v)\mapsto(\sqrt[3]u,v)$，但 $\det DH(0,y)=0$，逆映射在 $u=0$ 处不可微。非零雅可比是局部连续可微逆映射的条件，不能与只有连续性的逆映射混为一谈。
\end{solve}

\section{多元函数的极值}
\label{sec:extrema-of-multivariable-functions}

% 教材第 95--109 页；分章 PDF 第 56--70 页。

\subsection{通常极值问题}
\label{subsec:unconstrained-extrema}

\begin{definition}
\label{def:multivariable-local-extrema-ch14}
设 $f:D\to\mathbb R$，$\boldsymbol x_0\in D$。若存在 $\delta>0$，使 $B(\boldsymbol x_0,\delta)\subseteq D$，且该球内恒有 $f(\boldsymbol x)\geq f(\boldsymbol x_0)$，则称 $\boldsymbol x_0$ 为\emph{局部极小值点}（local minimum point）；不等号反向则称为\emph{局部极大值点}（local maximum point）。若对球内 $\boldsymbol x\ne\boldsymbol x_0$ 不等式严格成立，则称为严格局部极值点。若 $f$ 可微且 $\nabla f(\boldsymbol x_0)=\boldsymbol0$，称 $\boldsymbol x_0$ 为\emph{驻点}（stationary point）。
\end{definition}

局部极值只与一个邻域有关，全局最值则要求相应不等式在整个指定集合上成立。驻点附近若既有比 $f(\boldsymbol x_0)$ 大的值，也有比它小的值，称为\emph{鞍点}（saddle point），例如 $f(x,y)=x^2-y^2$ 的原点。闭集上的最值还可能在边界取得，不能只求内点驻点。

\begin{theorem}[极值的必要条件与二阶判别]
\label{thm:multivariable-extremum-tests-ch14}
若 $\boldsymbol x_0$ 为内点极值点，且 $f$ 在该点存在各个偏导数，则所有偏导数都为零。进一步设 $f\in C^2$ 于该点的某邻域，且 $\nabla f(\boldsymbol x_0)=\boldsymbol0$，则当 $H_f(\boldsymbol x_0)$ 正定时，$\boldsymbol x_0$ 为严格局部极小值点；负定时为严格局部极大值点；不定时为鞍点。反过来，局部极小值点的海塞矩阵必须半正定，局部极大值点的海塞矩阵必须半负定。
\end{theorem}
\begin{proof}
将 $f$ 限制在过 $\boldsymbol x_0$ 的每条坐标直线上，所得一元函数仍在参数零点取得极值，由费马定理得到偏导数为零。对于二阶结论，由式~\eqref{eq:second-order-taylor-peano-ch14}，
\[
    f(\boldsymbol x_0+\boldsymbol h)-f(\boldsymbol x_0)
    =\frac12\boldsymbol h^{\mathsf T}H_f(\boldsymbol x_0)\boldsymbol h
       +o(\|\boldsymbol h\|^2).
\]
若矩阵正定，则二次型在单位球面上连续且处处正，由紧性存在 $c>0$，使其对所有 $\boldsymbol h$ 不小于 $c\|\boldsymbol h\|^2$。充分小时，余项的绝对值不超过 $c\|\boldsymbol h\|^2/4$，故非零增量使函数值严格增加。对 $-f$ 应用这一结论即得负定情形。若二次型不定，选取使其分别为正、负的两个固定方向，沿这两个方向取充分小的增量，便得到正、负两种函数增量。

若 $\boldsymbol x_0$ 为局部极小值点，对每个固定向量 $\boldsymbol v$ 将上式代入 $\boldsymbol h=t\boldsymbol v$，除以 $t^2$ 后令 $t\to0$，即得 $\boldsymbol v^{\mathsf T}H_f(\boldsymbol x_0)\boldsymbol v\geq0$。极大值情形同理。
\end{proof}

若海塞矩阵只是半正定或半负定，二阶信息一般不能决定极值。例如 $x^4+y^4$、$-x^4-y^4$、$x^4-y^4$ 在原点具有相同的零海塞矩阵，却分别给出严格极小、严格极大与鞍点。对二元函数，记 $A=f_{xx}$、$B=f_{xy}$、$C=f_{yy}$，均在驻点取值，$\Delta=AC-B^2$，则 $\Delta>0,A>0$ 判为严格极小，$\Delta>0,A<0$ 判为严格极大，$\Delta<0$ 判为鞍点，$\Delta=0$ 时须另作分析。正定与负定可用顺序主子式等线性代数判据检验。

\begin{example}
\label{ex:quadratic-global-minimum-ch14}
求 $f(x,y)=x^2+3xy+3y^2-6x-3y$ 的极值。
\end{example}
\begin{solve}
驻点方程为 $2x+3y-6=0$、$3x+6y-3=0$，唯一解为 $(9,-4)$。海塞矩阵为 $\begin{pmatrix}2&3\\3&6\end{pmatrix}$，其首阶顺序主子式为 $2$，行列式为 $3$，故正定。再令 $a=x-9$、$b=y+4$，可直接得到
\[
    f(x,y)=-21+a^2+3ab+3b^2
       =-21+\left(a+\frac32b\right)^2+\frac34b^2.
\]
所以 $-21$ 还是唯一的全局最小值，函数没有局部极大值。
\end{solve}

\begin{example}
\label{ex:open-box-material-minimum-ch14}
无盖长方体盒子的容积为 $1\,\mathrm m^3$，底板厚度为侧板厚度的三倍。采用忽略接缝与板厚对容积影响的薄壁模型，求最省材料的长、宽、高。
\end{example}
\begin{solve}
设长、宽、高的米数分别为 $x,y,z>0$，共同的侧板厚度因子不影响最优尺寸。用料与 $3xy+2xz+2yz$ 成正比，约束为 $xyz=1$。消去 $z$ 后，目标函数为
\[
    S(x,y)=3xy+\frac2x+\frac2y,\qquad x,y>0.
\]
驻点方程 $3y-2/x^2=0$、$3x-2/y^2=0$ 给出 $x=y=a=(2/3)^{1/3}$。为确认全局最小值，使用三个正数的算术--几何平均不等式：
\[
    S(x,y)\geq3\sqrt[3]{(3xy)(2/x)(2/y)}=3\sqrt[3]{12},
\]
且等号当且仅当 $3xy=2/x=2/y$，恰对应 $x=y=a$。因此最优底面是正方形，高为 $z=a^{-2}=(3/2)^{2/3}=3a/2$，上述各长度的单位均为米。唯一驻点本身并不能代替这个全局最优性的论证。
\end{solve}

\begin{example}
\label{ex:least-squares-line-ch14}
给定平面上的 $N\geq2$ 个点 $(x_i,y_i)$，且 $x_i$ 不全相等，求使平方残差和 $Q(a,b)=\sum_{i=1}^N(ax_i+b-y_i)^2$ 最小的直线 $y=ax+b$。
\end{example}
\begin{solve}
令 $\bar x=N^{-1}\sum_i x_i$、$\bar y=N^{-1}\sum_i y_i$。驻点方程为
\[
    a\sum_i x_i^2+b\sum_i x_i=\sum_i x_iy_i,\qquad
    a\sum_i x_i+Nb=\sum_i y_i.
\]
因 $\sum_i(x_i-\bar x)^2>0$，它有唯一解
\[
    a_* =\frac{\sum_i(x_i-\bar x)(y_i-\bar y)}{\sum_i(x_i-\bar x)^2},
    \qquad b_* =\bar y-a_*\bar x.
\]
记 $r_i=a_*x_i+b_*-y_i$。正规方程使 $\sum_i r_i=\sum_i x_ir_i=0$，故任意增量 $s,t$ 满足
\[
    Q(a_*+s,b_*+t)-Q(a_*,b_*)=\sum_i(sx_i+t)^2.
\]
若右边为零，则 $sx_i+t=0$ 对所有 $i$ 成立；由于至少有两个不同的 $x_i$，必有 $s=t=0$。因此该解是唯一的全局最小值点。这称为直线拟合的\emph{最小二乘法}（least squares method），其结论是一个确定性的优化结果，无须预设观测值的概率模型。
\end{solve}

\subsection{条件极值问题}
\label{subsec:constrained-extrema}

在等式约束 $\Phi(\boldsymbol x)=\boldsymbol0$ 下求 $f$ 的极值，只允许在可行集合内比较函数值。若能将约束解为参数函数，可以代入目标函数化为通常极值问题；隐函数定理则允许我们在不显式求出参数函数的情况下消去其导数，得到拉格朗日乘数法。

\begin{theorem}[拉格朗日乘数法的必要条件]
\label{thm:lagrange-multiplier-ch14}
设 $f:D\to\mathbb R$、$\Phi=(\phi_1,\ldots,\phi_m)^{\mathsf T}:D\to\mathbb R^m$ 均为 $C^1$，$m<n$。若 $\boldsymbol x_0$ 是 $f$ 在约束 $\Phi=\boldsymbol0$ 下的局部极值点，且
\[
    \operatorname{rank}D\Phi(\boldsymbol x_0)=m,
\]
则存在唯一的乘数向量 $\boldsymbol\lambda=(\lambda_1,\ldots,\lambda_m)^{\mathsf T}$，使
\begin{equation}
\label{eq:lagrange-multiplier-ch14}
    \nabla f(\boldsymbol x_0)+D\Phi(\boldsymbol x_0)^{\mathsf T}\boldsymbol\lambda
       =\boldsymbol0,\qquad
    \Phi(\boldsymbol x_0)=\boldsymbol0.
\end{equation}
等价地，令\emph{拉格朗日函数}（Lagrangian）
$L(\boldsymbol x,\boldsymbol\lambda)=f(\boldsymbol x)+\sum_{j=1}^m\lambda_j\phi_j(\boldsymbol x)$，则候选点须满足 $\nabla_{\boldsymbol x}L=\boldsymbol0$、$\nabla_{\boldsymbol\lambda}L=\boldsymbol0$。
\end{theorem}
\begin{proof}
由满秩条件，重新排列坐标后可写 $\boldsymbol x=(\boldsymbol s,\boldsymbol t)$，其中 $\boldsymbol t\in\mathbb R^m$，且 $B=D_{\boldsymbol t}\Phi(\boldsymbol x_0)$ 可逆。局部约束唯一写为 $\boldsymbol t=g(\boldsymbol s)$；令 $A=D_{\boldsymbol s}\Phi(\boldsymbol x_0)$，则 $Dg=-B^{-1}A$。将约束代入 $f$ 后，其关于 $\boldsymbol s$ 的导数在极值点为零，即
\[
    D_{\boldsymbol s}f-D_{\boldsymbol t}f\,B^{-1}A=0.
\]
取 $\boldsymbol\lambda^{\mathsf T}=-D_{\boldsymbol t}f\,B^{-1}$，上式与这一取法分别给出关于 $\boldsymbol s$、$\boldsymbol t$ 的乘数方程，合并即得结论。若两组乘数都满足条件，其差在 $D\Phi^{\mathsf T}$ 的核中；满行秩使该核为零，故乘数唯一。
\end{proof}

满秩条件表示约束梯度线性无关，不能省略。例如约束 $x^2+y^2=0$ 的可行集合只有原点，$f(x,y)=x$ 在其中既取最大值也取最小值，但原点约束梯度为零，不可能满足乘数方程。即使满秩，式~\eqref{eq:lagrange-multiplier-ch14} 仍只是必要条件；它没有把原问题变为对 $(\boldsymbol x,\boldsymbol\lambda)$ 无约束求极值的问题，乘数方向上的线性结构通常也不会产生这样的极值。要求全局最值时，还须处理约束退化点、边界以及最值存在性。

若 $f,\Phi\in C^2$，在满足乘数方程的正则可行点定义约束的\emph{切空间}（tangent space）
\[
    T=\{\boldsymbol v\in\mathbb R^n:D\Phi(\boldsymbol x_0)\boldsymbol v=\boldsymbol0\}.
\]
二阶判别应使用固定乘数 $\boldsymbol\lambda_0$ 后的矩阵
\[
    H_{\boldsymbol x}L
       =H_f+\sum_j\lambda_{0j}H_{\phi_j}
\]
在 $T$ 上的二次型：若它在 $T\setminus\{\boldsymbol0\}$ 上正定，则为严格条件极小值；负定则为严格条件极大值。条件极小值要求它在 $T$ 上半正定，条件极大值要求半负定。

这些结论可由前面局部参数化验证。令可行参数化为 $\boldsymbol x=r(\boldsymbol s)$，则 $Dr$ 将参数空间一一映到 $T$。对 $\Phi\circ r=\boldsymbol0$ 求二阶导，并利用 $\nabla f=-D\Phi^{\mathsf T}\boldsymbol\lambda_0$，可得约化目标 $f\circ r$ 的海塞矩阵为
\[
    H_{f\circ r}=Dr^{\mathsf T}(H_{\boldsymbol x}L)Dr.
\]
于是通常极值的二阶判别直接给出上述结论。整个空间上正定或负定当然也是充分条件，但并非必要：例如 $f(x,y)=x^2-y^2$ 在约束 $y=0$ 下于原点取严格极小值，虽然相应海塞矩阵在整个平面上不定。

\begin{example}
\label{ex:ellipsoid-inscribed-box-ch14}
求椭球面 $16x^2+4y^2+9z^2=144$ 内接、各面平行于坐标平面的长方体的最大体积。
\end{example}
\begin{solve}
沿长方体任一条平行于 $x$ 轴的棱，其两端都在椭球面上；将两端的曲面方程相减可知，两端 $x$ 坐标的平方相等，故为互相反号的数。对另外两组棱同理，长方体的八个顶点必为 $(\pm x,\pm y,\pm z)$，其中 $x,y,z>0$，目标为 $V=8xyz$。先允许 $x,y,z\geq0$，则可行集合闭且有界，连续目标在其上取到最大值；有坐标为零的边缘部分满足 $V=0$，而三坐标均正时有 $V>0$，故最大值点的三个坐标均正，且约束梯度非零。乘数方程为
\[
    8yz+32\lambda x=0,\quad
    8xz+8\lambda y=0,\quad
    8xy+18\lambda z=0.
\]
分别乘以 $x,y,z$，由三个坐标均正可知 $\lambda\ne0$，进而
\[
    16x^2=4y^2=9z^2=48.
\]
因此唯一的正坐标候选点为 $x=\sqrt3$、$y=2\sqrt3$、$z=4\sqrt3/3$。存在性与边界分析已排除其他最大值位置，故 $V_{\max}=64\sqrt3$。
\end{solve}

\begin{example}
\label{ex:product-maximum-and-amgm-ch14}
固定 $n\geq2$ 个正数的和为 $c>0$，求乘积的最大值，并得到算术--几何平均不等式。
\end{example}
\begin{solve}
将问题先扩充到紧集 $E=\{\boldsymbol x:x_i\geq0,\ \sum_i x_i=c\}$ 上。乘积 $P=\prod_i x_i$ 连续，故取到最大值；边界乘积为零，而等分点的乘积为正，所以最大值点满足每个 $x_i>0$。约束梯度为 $(1,\ldots,1)^{\mathsf T}$，满秩条件成立。乘数方程为 $P/x_i+\lambda=0$，从而所有 $x_i$ 相等。于是
\[
    \max_{\substack{x_i>0\\\sum_i x_i=c}}\prod_{i=1}^n x_i
       =\left(\frac cn\right)^n,
\]
且等号仅在 $x_1=\cdots=x_n=c/n$ 时成立。对给定正数 $a_1,\ldots,a_n$ 取 $c=\sum_i a_i$，便得到
\[
    \sqrt[n]{a_1\cdots a_n}\leq\frac{a_1+\cdots+a_n}{n},
\]
等号当且仅当各数相等。先在闭包上证明最大值存在，再排除边界，避免了直接把开单纯形误当作紧集。
\end{solve}

\section{多元微分学的几何应用}
\label{sec:geometric-applications-of-multivariable-differentiation}

% 教材第 109--120 页；分章 PDF 第 70--81 页，止于习题十四之前。

\subsection{空间曲线的切线与法平面}
\label{subsec:tangent-lines-and-normal-planes-of-space-curves}

连续映射 $\gamma:[a,b]\to\mathbb R^n$ 描述一条参数曲线。若它在 $[a,b]$ 上单射，称为\emph{简单曲线}（simple curve）；若在 $[a,b)$ 上单射且 $\gamma(a)=\gamma(b)$，称为\emph{简单闭曲线}（simple closed curve），也称\emph{约当曲线}（Jordan curve）。仅有连续性不能保证曲线具有切线；本节主要讨论 $\gamma\in C^1$ 且 $\gamma'(t)\ne\boldsymbol0$ 的\emph{正则曲线}（regular curve），教材称为光滑曲线。每一点附近至少有一个坐标分量的导数保持非零，因此该坐标沿参数严格单调，曲线局部为简单曲线。

设空间曲线为 $\gamma(t)=(x(t),y(t),z(t))^{\mathsf T}$，$t_0$ 为参数区间的内点。由
\[
    \frac{\gamma(t)-\gamma(t_0)}{t-t_0}\longrightarrow\gamma'(t_0)\ne\boldsymbol0,
\]
割线的方向直线趋于 $\gamma'(t_0)$ 张成的直线，因此 $\boldsymbol p=\gamma(t_0)$ 处的\emph{切线}（tangent line）为
\begin{equation}
\label{eq:space-curve-tangent-ch14}
    \boldsymbol x=\boldsymbol p+s\gamma'(t_0),\qquad s\in\mathbb R.
\end{equation}
过 $\boldsymbol p$ 且以该切向量为法向量的\emph{法平面}（normal plane）为
\begin{equation}
\label{eq:space-curve-normal-plane-ch14}
    \gamma'(t_0)^{\mathsf T}(\boldsymbol x-\boldsymbol p)=0.
\end{equation}
切线的对称式 $(x-x_0)/x'(t_0)=(y-y_0)/y'(t_0)=(z-z_0)/z'(t_0)$ 只在各分母非零时直接使用；参数式始终适用。若 $\gamma'(t_0)=\boldsymbol0$，这些公式不能确定切线，但也不能据此断言几何曲线必无切线，例如 $\gamma(t)=(t^3,0,0)$。

若空间曲线由 $F(x,y,z)=G(x,y,z)=0$ 给出，$F,G\in C^1$，且在公共零点 $\boldsymbol p$ 处两个梯度线性无关，则隐函数组定理保证它在局部是一条正则曲线。沿曲线求导得
\[
    \nabla F(\boldsymbol p)^{\mathsf T}\gamma'(t_0)=0,\qquad
    \nabla G(\boldsymbol p)^{\mathsf T}\gamma'(t_0)=0.
\]
因此可取非零切向量
\begin{equation}
\label{eq:implicit-curve-tangent-vector-ch14}
    \boldsymbol T=\nabla F(\boldsymbol p)\times\nabla G(\boldsymbol p)
      =(A,B,C)^{\mathsf T},\quad
    A=\frac{\partial(F,G)}{\partial(y,z)},\quad
    B=\frac{\partial(F,G)}{\partial(z,x)},\quad
    C=\frac{\partial(F,G)}{\partial(x,y)},
\end{equation}
各行列式均在 $\boldsymbol p$ 取值。以 $\boldsymbol T$ 代替前两式中的 $\gamma'(t_0)$ 即得切线和法平面；切线也就是两个方程各自的切平面的交线。若两个梯度相关，叉积为零，则该方法失效，需要重新分析零点集合，不能直接认定其仍为一条正则曲线。

\begin{example}
\label{ex:implicit-space-curve-tangent-ch14}
求曲线 $x^2-y^2+2z^2=2$、$x+y+z=3$ 在 $(1,1,1)$ 处的切线与法平面。
\end{example}
\begin{solve}
令 $F=x^2-y^2+2z^2-2$、$G=x+y+z-3$，则在该点
\[
    \nabla F=(2,-2,4)^{\mathsf T},\quad
    \nabla G=(1,1,1)^{\mathsf T},\quad
    \nabla F\times\nabla G=(-6,2,4)^{\mathsf T}.
\]
约去共同因子 $2$，切线可写为 $(x,y,z)=(1,1,1)+s(-3,1,2)$，法平面为 $-3(x-1)+(y-1)+2(z-1)=0$，即 $3x-y-2z=0$。另一种求切线的方法是取两个切平面的交线：
\[
    x-y+2z=2,\qquad x+y+z=3.
\]
两个平面法向量独立，其交线方向与上面求得的叉积一致。
\end{solve}

\subsection{曲面的切平面与法线}
\label{subsec:tangent-planes-and-normal-lines-of-surfaces}

设曲面在 $\boldsymbol p=(x_0,y_0,z_0)^{\mathsf T}$ 附近由 $F(x,y,z)=0$ 确定，其中 $F\in C^1$、$F(\boldsymbol p)=0$、$\nabla F(\boldsymbol p)\ne\boldsymbol0$。对曲面上任一过 $\boldsymbol p$ 的可微曲线 $\gamma$，链式法则给出 $\nabla F(\boldsymbol p)^{\mathsf T}\gamma'(t_0)=0$。反过来，隐函数定理给出曲面的局部参数化，其导数的像正好是这个二维核空间，故每个这样的切向量都可由曲面内曲线实现。因此曲面在该点的\emph{切平面}（tangent plane）与\emph{法线}（normal line）分别为
\begin{equation}
\label{eq:implicit-surface-tangent-normal-ch14}
    \nabla F(\boldsymbol p)^{\mathsf T}(\boldsymbol x-\boldsymbol p)=0,
    \qquad \boldsymbol x=\boldsymbol p+s\nabla F(\boldsymbol p),\quad s\in\mathbb R.
\end{equation}
具体地，切平面方程为 $F_x(\boldsymbol p)(x-x_0)+F_y(\boldsymbol p)(y-y_0)+F_z(\boldsymbol p)(z-z_0)=0$。对于图像 $z=f(x,y)$，取 $F=z-f(x,y)$，得到
\[
    z-z_0=f_x(x_0,y_0)(x-x_0)+f_y(x_0,y_0)(y-y_0),\qquad
    \boldsymbol n=(-f_x,-f_y,1)^{\mathsf T}.
\]
这正是可微性的局部线性近似在几何上的表达。

若曲面以 $\boldsymbol r(u,v)=(x(u,v),y(u,v),z(u,v))^{\mathsf T}$ 参数化，$\boldsymbol r\in C^1$，且 $D\boldsymbol r(u_0,v_0)$ 的秩为 $2$，则两条坐标曲线的切向量 $\boldsymbol r_u,\boldsymbol r_v$ 线性无关。对任意参数曲线应用链式法则，其切向量均在二者张成的平面内，因此可取
\begin{equation}
\label{eq:parametric-surface-normal-ch14}
    \boldsymbol n=\boldsymbol r_u\times\boldsymbol r_v
       =\left(\frac{\partial(y,z)}{\partial(u,v)},
              \frac{\partial(z,x)}{\partial(u,v)},
              \frac{\partial(x,y)}{\partial(u,v)}\right)^{\mathsf T}
\end{equation}
作为法向量，各项在 $(u_0,v_0)$ 计算。令 $\boldsymbol p=\boldsymbol r(u_0,v_0)$，则切平面为 $\boldsymbol n^{\mathsf T}(\boldsymbol x-\boldsymbol p)=0$，法线为 $\boldsymbol x=\boldsymbol p+s\boldsymbol n$。交换参数次序会改变法向量的符号，但不改变切平面与法线本身。

两个正则曲面在公共点处的夹角由法向量决定，通常取两法线的锐角或直角，满足
\[
    \cos\alpha=\frac{|\boldsymbol n_1^{\mathsf T}\boldsymbol n_2|}
       {\|\boldsymbol n_1\|\|\boldsymbol n_2\|},\qquad 0\leq\alpha\leq\frac\pi2.
\]
法向量正交时称曲面正交，相互平行时切平面相同而称相切；公共点处法向量均非零是这些判定的前提。推广到 $\mathbb R^n$，正则曲线的法平面相应成为 $\gamma'(t_0)^{\mathsf T}(\boldsymbol x-\boldsymbol p)=0$ 的法超平面；正则水平超曲面 $F=0$ 的切超平面仍由 $\nabla F(\boldsymbol p)^{\mathsf T}(\boldsymbol x-\boldsymbol p)=0$ 给出，其法线沿梯度方向。

\begin{example}
\label{ex:graph-and-parametric-surface-tangent-ch14}
求抛物面 $z=x^2+y^2$ 在 $(1,1,2)$ 处的切平面与法线，并核对隐式与参数表示的结果。
\end{example}
\begin{solve}
由图像公式，切平面为 $z-2=2(x-1)+2(y-1)$，即 $z=2x+2y-2$，法线为
\[
    (x,y,z)=(1,1,2)+s(-2,-2,1).
\]
取参数化 $\boldsymbol r(u,v)=(u,v,u^2+v^2)^{\mathsf T}$，在 $(1,1)$ 处 $\boldsymbol r_u=(1,0,2)^{\mathsf T}$、$\boldsymbol r_v=(0,1,2)^{\mathsf T}$，其叉积为 $(-2,-2,1)^{\mathsf T}$。这与隐式函数 $F=z-x^2-y^2$ 的梯度相同，两种方法给出一致结果。
\end{solve}

\subsection{多元凸函数}
\label{subsec:multivariable-convex-functions}

\begin{definition}
\label{def:multivariable-convex-function-ch14}
设 $D\subseteq\mathbb R^n$ 为凸集，即任意两点之间的整条线段仍在 $D$ 内。若 $f:D\to\mathbb R$ 满足
\[
    f((1-t)\boldsymbol x+t\boldsymbol y)
       \leq(1-t)f(\boldsymbol x)+tf(\boldsymbol y),\qquad
    \boldsymbol x,\boldsymbol y\in D,\quad 0<t<1,
\]
则称 $f$ 为\emph{凸函数}（convex function）。若对 $\boldsymbol x\ne\boldsymbol y$ 总成立严格不等式，称为\emph{严格凸函数}（strictly convex function）。$-f$ 凸时称 $f$ 为凹函数。凸性本身不要求可微，以下导数判别则在凸开集上讨论。
\end{definition}

\begin{theorem}[凸函数的一阶与二阶判别]
\label{thm:multivariable-convexity-tests-ch14}
设 $D$ 是凸开集，$f\in C^2(D)$，则下列条件等价：$f$ 在 $D$ 上凸；任意 $\boldsymbol x,\boldsymbol y\in D$ 满足
\begin{equation}
\label{eq:convex-first-order-inequality-ch14}
    f(\boldsymbol y)\geq f(\boldsymbol x)
       +\nabla f(\boldsymbol x)^{\mathsf T}(\boldsymbol y-\boldsymbol x);
\end{equation}
$H_f(\boldsymbol x)$ 在 $D$ 的每一点半正定。其中凸性与式~\eqref{eq:convex-first-order-inequality-ch14} 的等价只需 $f$ 可微。
\end{theorem}
\begin{proof}
若 $f$ 凸，令 $\boldsymbol x_t=\boldsymbol x+t(\boldsymbol y-\boldsymbol x)$，则
\[
    \frac{f(\boldsymbol x_t)-f(\boldsymbol x)}t
       \leq f(\boldsymbol y)-f(\boldsymbol x).
\]
令 $t\to0+$，由可微性得到式~\eqref{eq:convex-first-order-inequality-ch14}。反过来，对 $\boldsymbol z=(1-t)\boldsymbol x+t\boldsymbol y$，分别将一阶不等式用于 $f(\boldsymbol x)$ 与 $f(\boldsymbol y)$ 在 $\boldsymbol z$ 处的比较，再乘以 $1-t,t$ 相加，梯度项相消，即得凸性。这也证明了一阶判别只需可微。

若一阶不等式成立，固定 $\boldsymbol x$ 和 $\boldsymbol v$，在足够小的 $s$ 下代入 $\boldsymbol y=\boldsymbol x+s\boldsymbol v$，并用二阶泰勒公式，得
\[
    \frac12s^2\boldsymbol v^{\mathsf T}H_f(\boldsymbol x)\boldsymbol v+o(s^2)\geq0.
\]
除以 $s^2$ 后取极限可知海塞矩阵半正定。反之，若海塞矩阵处处半正定，由沿线段的二阶泰勒公式，
\[
    f(\boldsymbol y)-f(\boldsymbol x)-\nabla f(\boldsymbol x)^{\mathsf T}(\boldsymbol y-\boldsymbol x)
    =\frac12(\boldsymbol y-\boldsymbol x)^{\mathsf T}
       H_f(\boldsymbol x+\theta(\boldsymbol y-\boldsymbol x))
       (\boldsymbol y-\boldsymbol x)\geq0,
\]
故一阶不等式成立。
\end{proof}

几何上，一阶不等式说明函数图像位于任一点的切超平面上方；二阶判别说明沿每条直线的弯曲程度都非负。海塞矩阵处处正定足以保证严格凸性，因为限制到任意非平凡线段后，其二阶导数严格为正；但这不是必要条件，例如 $f(\boldsymbol x)=\sum_i x_i^4$ 严格凸，在原点的海塞矩阵却为零。对可微凸函数，$\nabla f(\boldsymbol x_0)=\boldsymbol0$ 由一阶不等式直接推出全局最小值；反之，开集上的最小值点必为驻点。严格凸函数的最小值点若存在则唯一，否则两个不同最小值点的中点会有更小的函数值。

\begin{example}
\label{ex:separable-convex-function-ch14}
设一元函数 $g\in C^2((a,b))$ 为凸函数，证明 $f(x_1,\ldots,x_n)=\sum_{i=1}^n g(x_i)$ 在 $(a,b)^n$ 上凸。
\end{example}
\begin{solve}
一元凸性给出 $g''\geq0$，而 $f$ 的海塞矩阵为对角矩阵
\[
    H_f(\boldsymbol x)=\operatorname{diag}(g''(x_1),\ldots,g''(x_n)).
\]
故对任意向量 $\boldsymbol v$，有 $\boldsymbol v^{\mathsf T}H_f\boldsymbol v=\sum_i g''(x_i)v_i^2\geq0$，由二阶判别即得凸性。若 $g''>0$，则该矩阵正定，$f$ 严格凸；更一般地，若 $g$ 本身严格凸，直接逐项应用定义，也可得到 $f$ 严格凸，而不必要求 $g''$ 处处严格为正。
\end{solve}


\section*{习题十四}
\addcontentsline{toc}{section}{习题十四}
\fi

\selectchapter{15}
\ifcompilechapter
\chapter{重积分}
\label{chap:multiple-integrals}

一元定积分把区间上的局部贡献累加起来；多元函数的\emph{重积分}（multiple integral）把同样的思想用于平面或高维空间中的区域。例如，物体占据区域 $D\subset\mathbb R^3$，密度为 $\rho(x,y,z)$，把 $D$ 分成小块后，以每块的密度乘体积近似其质量，便得到 $\sum_k\rho(\xi_k)\Delta V_k$。重积分所研究的是这些和式何时具有与分割、取点无关的极限，以及怎样计算这一极限。与一元情形相比，主要的新问题是区域的形状：必须先定义体积，再处理分割、切片和坐标变换，最后区分无界区域或无界函数的广义积分。

% 本章依据《数学分析（伍胜健）》第十五章正文整理；对应数学分析-15.pdf 第 1--51 页，不含习题十五。

\section{重积分的定义}
\label{sec:definition-of-multiple-integrals}

\subsection{\texorpdfstring{$\mathbb{R}^n$}{R^n} 空间中集合的体积}
\label{subsec:volumes-of-sets-in-euclidean-space}

对 $x=(x_1,\ldots,x_n)$，采用欧氏距离 $|x-y|=(\sum_{j=1}^n(x_j-y_j)^2)^{1/2}$。坐标长方体 $Q=\prod_{j=1}^n[a_j,b_j]$ 的体积规定为 $V(Q)=\prod_{j=1}^n(b_j-a_j)$；一般集合的体积用有限个这种长方体从内外逼近。这里的体积是\emph{若尔当体积}（Jordan content），只讨论有界集合，不能预先赋予它对任意可数并的可加性。

\begin{definition}
\label{def:multiple-jordan-volume}
设 $E\subset\mathbb R^n$ 有界，取包含 $E$ 的长方体 $Q$，用有限组平行于坐标超平面的平面把 $Q$ 分成小长方体，所得分割记为 $\mathcal P$。定义内、外体积和
\[
    v_*(E,\mathcal P)=\sum_{Q_k\subset E}V(Q_k),\qquad
    v^*(E,\mathcal P)=\sum_{Q_k\cap E\ne\varnothing}V(Q_k),
\]
以及内、外体积
\[
    v_*(E)=\sup_{\mathcal P}v_*(E,\mathcal P),\qquad
    v^*(E)=\inf_{\mathcal P}v^*(E,\mathcal P).
\]
若 $v_*(E)=v^*(E)$，则称 $E$ \emph{可求体积}（Jordan measurable），公共值记为 $V(E)$。当 $V(E)=0$ 时，称 $E$ 为\emph{零体积集}（set of Jordan content zero）；二维时也称可求面积、零面积集，面积有时记为 $\sigma(E)$。
\end{definition}

细分使内和不减、外和不增；任意两个网格有公共细分，因而总有 $v_*(E)\leq v^*(E)$。定义不依赖包含 $E$ 的大长方体：扩大长方体只增加不与 $E$ 相交的部分，两个候选长方体又可放进同一个更大的长方体比较。由定义，$E$ 为零体积集，当且仅当对每个 $\varepsilon>0$，都能用有限个长方体覆盖 $E$，且这些长方体的体积总和小于 $\varepsilon$。

\begin{theorem}
\label{thm:multiple-jordan-boundary}
有界集合 $E\subset\mathbb R^n$ 可求体积，当且仅当其边界 $\partial E$ 是零体积集。
\end{theorem}
\begin{proof}
先设 $E$ 可求体积。给定 $\varepsilon>0$，选网格使外和与内和之差小于 $\varepsilon$。除网格平面上的点外，$\partial E$ 的点所在的小长方体必既接触 $E$ 又不完全含于 $E$，故这些长方体的总体积小于 $\varepsilon$。有限个网格平面在大长方体内的部分，可用总厚度任意小的薄长方体覆盖，体积也可任意小。因此 $\partial E$ 为零体积集。

反过来，设 $\partial E$ 为零体积集。将它覆盖在有限个开长方体的并 $G$ 中，使其体积总和小于 $\varepsilon$；这是把零体积覆盖略微扩大所得的。由于 $\partial E$ 紧，存在 $\delta>0$，使 $\partial E$ 的 $\delta$ 邻域包含在 $G$ 中。取直径小于 $\delta$ 的网格。凡既与 $E$ 相交又不完全含于 $E$ 的小长方体，都接触 $\partial E$，从而含于 $G$。这些小长方体内部互不相交，体积总和小于 $\varepsilon$，故 $v^*(E)-v_*(E)<\varepsilon$。令 $\varepsilon\downarrow0$ 即得结论。
\end{proof}

\begin{corollary}
\label{cor:multiple-jordan-set-operations}
若有界集合 $A,B$ 可求体积，则 $A\cup B$、$A\cap B$、$A\setminus B$ 均可求体积。可求体积的有界区域 $U$ 与其闭包 $\overline U$ 具有相同体积；若另有 $U=\operatorname{int}\overline U$，则 $U$ 可求体积当且仅当 $\overline U$ 可求体积。
\end{corollary}
\begin{proof}
前三个集合的边界均包含在 $\partial A\cup\partial B$ 中，有限个零体积集的并仍为零体积集。若开区域 $U$ 可求体积，则 $\partial\overline U\subset\partial U$，从而 $\overline U$ 可求体积，两者只相差零体积边界。另有 $U=\operatorname{int}\overline U$ 时，$\partial U=\partial\overline U$，逆向结论也成立。
\end{proof}

\begin{example}
\label{ex:multiple-rectifiable-boundary}
证明：若平面有界区域的边界由有限条可求长曲线组成，则它可求面积。
\end{example}
\begin{solve}
设其中一条曲线长为 $L>0$，按弧长分成 $N$ 段，每段的直径不超过 $L/N$，可包含在一个边长 $2L/N$ 的正方形内。全部覆盖的面积不超过 $4L^2/N\to0$，故该曲线面积为零；长度为零的曲线退化为一点，也有相同结论。有限条曲线的并面积为零，再用定理~\ref{thm:multiple-jordan-boundary} 即得所述区域可求面积。特别地，有限条分段光滑曲线围成的有界区域符合这一条件。
\end{solve}

\begin{notealert}[type=caution,title={有界与可求体积是不同条件}]
有界集合未必可求体积，判定时仍须检查边界条件。有限零体积覆盖不能任意换成可数覆盖；本章始终使用若尔当体积。
\end{notealert}

\subsection{重积分的定义}
\label{subsec:definition-of-multiple-integrals}

以下把\emph{有界闭区域}理解为某个有界区域的闭包，因而 $D=\overline{D^\circ}$；另行要求 $D$ 可求体积。这一约定排除了在主体外另添孤立点等情形，有助于准确表述可积函数必有界等结论。

\begin{definition}
\label{def:multiple-riemann-integral}
设 $D\subset\mathbb R^n$ 为可求体积的有界闭区域。它的\emph{分割} $\mathcal P=\{D_1,\ldots,D_N\}$ 是有限个可求体积闭子集，满足
\[
    D=\bigcup_{k=1}^ND_k,\qquad V(D_j\cap D_k)=0\quad(j\ne k).
\]
记 $\Delta V_k=V(D_k)$，网格宽度 $\|\mathcal P\|=\max_k\operatorname{diam}D_k$，并在每个非空 $D_k$ 任取标记点 $\xi_k$。若存在 $I\in\mathbb R$，使对每个 $\varepsilon>0$，都有 $\delta>0$，且对任何满足 $\|\mathcal P\|<\delta$ 的分割及任何标记点，均有
\[
    \left|\sum_{k=1}^Nf(\xi_k)\Delta V_k-I\right|<\varepsilon,
\]
则称 $f$ 在 $D$ 上\emph{黎曼可积}，$I$ 称为其 $n$ 重积分，记为
\begin{equation}
    \int_D f\,dV=\lim_{\|\mathcal P\|\to0}
    \sum_{k=1}^Nf(\xi_k)\Delta V_k,
    \qquad dV=dx_1\cdots dx_n.
    \label{eq:multiple-riemann-definition}
\end{equation}
\end{definition}

二维、三维分别写成 $\iint_D f(x,y)\,dx\,dy$、$\iiint_D f(x,y,z)\,dx\,dy\,dz$，也可用 $d\sigma$、$dV$ 表示面积元、体积元。这里的 $\int_D$ 统一表示维数由区域决定的重积分，$\iint$、$\iiint$ 则把维数明确写出。极限必须同时独立于分割方式和标记点；一列特殊黎曼和收敛只能在已经证明可积后用于计算积分值。与一元定积分的有向上下限不同，本章区域体积总非负，交换坐标名称或改变坐标方向本身不会产生积分的负号。

若 $f\geq0$ 连续，$\iint_D f\,d\sigma$ 表示曲顶柱体 $\{(x,y,z):(x,y)\in D,\ 0\leq z\leq f(x,y)\}$ 的体积；若 $f$ 变号，它表示正、负两部分体积之差。特别地，$\int_D1\,dV=V(D)$。若 $\rho\geq0$ 是物体的连续密度，则质量为 $m=\iiint_D\rho\,dV$；在 $m>0$ 时，质心坐标为
\begin{equation}
    (\overline x,\overline y,\overline z)
    =\frac1m\left(\iiint_Dx\rho\,dV,
    \iiint_Dy\rho\,dV,\iiint_Dz\rho\,dV\right).
    \label{eq:multiple-mass-centroid}
\end{equation}

\section{多元函数的可积性理论与重积分的性质}
\label{sec:multivariable-integrability-and-multiple-integral-properties}

\subsection{达布理论}
\label{subsec:multivariable-darboux-theory}

\begin{theorem}
\label{thm:multiple-integrable-bounded}
可求体积的有界闭区域上的黎曼可积函数必有界。
\end{theorem}
\begin{proof}
若 $f$ 无界，紧性保证存在 $x_0\in D$，使 $f$ 在 $x_0$ 的每个相对邻域内无界；否则有限个局部有界邻域便给出全局界。给定任意小的网格宽度，先取足够小的闭球与 $D$ 的交作为一个分块，再把其余部分用网格细分。此分块可求体积，且因 $D=\overline{D^\circ}$ 而具有正体积，$f$ 在其中仍无界。固定其他标记点，仅改变该块的取点，就能使黎曼和的绝对值任意大，这与足够细的所有黎曼和都接近同一有限数矛盾。
\end{proof}

设 $f$ 有界，对分割 $\mathcal P$ 记
\[
    m_k=\inf_{D_k}f,\quad M_k=\sup_{D_k}f,\quad
    \omega_k=M_k-m_k,
\]
并定义\emph{达布下和、上和}（lower and upper Darboux sums）及下、上积分
\[
    L(f,\mathcal P)=\sum_km_k\Delta V_k,\qquad
    U(f,\mathcal P)=\sum_kM_k\Delta V_k,
\]
\[
    \underline{\int_D}f\,dV=\sup_{\mathcal P}L(f,\mathcal P),\qquad
    \overline{\int_D}f\,dV=\inf_{\mathcal P}U(f,\mathcal P).
\]
任一黎曼和介于对应的上下和之间；细分使下和不减、上和不增。两个分割可由所有非空交集构成公共细分，故任一下和不大于任一上和。

\begin{theorem}
\label{thm:multiple-darboux-criterion}
设 $f$ 在可求体积的有界闭区域 $D$ 上有界，则当 $\|\mathcal P\|\to0$ 时，达布下和、上和分别趋于下积分、上积分。并且下列条件等价：$f$ 可积；上、下积分相等；对每个 $\varepsilon>0$，存在一个分割 $\mathcal P$ 使
\begin{equation}
    U(f,\mathcal P)-L(f,\mathcal P)
    =\sum_k\omega_k\Delta V_k<\varepsilon.
    \label{eq:multiple-darboux-criterion}
\end{equation}
\end{theorem}
\begin{proof}
固定一个分割 $\mathcal P_0$。它的各块边界均为零体积集，故这些边界的有限并可覆盖在体积任意小的开集内。对任意足够细的另一个分割 $\mathcal P$，凡不能归入 $\mathcal P_0$ 某一块的分块，都落在这一边界邻域内；若 $|f|\leq M$，这些块在比较上下和时造成的误差至多为 $2M$ 乘其总体积。其余块可直接用细分单调性比较。因此当 $\|\mathcal P\|\to0$ 时，$U(f,\mathcal P)$ 不超过 $U(f,\mathcal P_0)$ 加任意小的误差；再对 $\mathcal P_0$ 取下确界，便得上和极限等于上积分。下和同理。

若上下积分相等，则所有充分细的黎曼和被夹在趋于同一数的上下和之间，故可积。若 $f$ 可积，在每块选取函数值任意接近上、下确界的点，再用黎曼和极限的任意性，得到上下和具有同一极限。最后，若式~\eqref{eq:multiple-darboux-criterion} 成立，则上下积分之差介于 $0$ 与任意 $\varepsilon$ 之间；反之，分别选取接近公共积分值的上和与下和，再取公共细分即可。
\end{proof}

实际计算常只用小坐标长方体。取包含 $D$ 的长方体 $Q$，将 $f$ 在 $Q\setminus D$ 上延拓为 $0$，得到 $\widetilde f$。由于 $\partial D$ 为零体积集，接触边界的网格块的体积总和趋于 $0$，故 $f$ 在 $D$ 上可积等价于 $\widetilde f$ 在 $Q$ 上可积，积分值也相等。同理，在足够细的坐标网格中可以省略接触 $\partial D$ 的块，只对完全落在 $D^\circ$ 的小长方体作黎曼和；省去部分的绝对值不超过 $\sup_D|f|$ 乘边界块总体积，因而仍趋于同一积分。

\begin{theorem}
\label{thm:multiple-continuous-integrable}
可求体积的有界闭区域上的连续函数可积。更一般地，若 $f$ 有界，且其间断点集包含于一个零体积集，则 $f$ 可积。
\end{theorem}
\begin{proof}
连续情形由紧集上的一致连续性得出：给定 $\varepsilon>0$，令各块直径充分小，就有 $\omega_k<\varepsilon/V(D)$，于是式~\eqref{eq:multiple-darboux-criterion} 成立。若 $V(D)=0$，任何有界函数的上下和均为 $0$，无需此除法。

一般情形中，先把间断点包含在一个总体积任意小的有限开长方体并 $G$ 中。在紧集 $D\setminus G$ 上，每点都是 $f$ 相对于整个 $D$ 的连续点；由连续性邻域的有限覆盖，可令所有与该紧集相交的充分小分块的振幅一致地小。其余分块全部落在 $G$ 中，振幅至多为 $2\sup_D|f|$。两类分块的振幅体积和分别由连续性和 $G$ 的小体积控制，故可积。
\end{proof}

\begin{example}
\label{ex:multiple-dense-line-discontinuities}
将 $[0,1]$ 中有理数无重复地枚举为 $q_1,q_2,\ldots$，在 $D=[0,1]^2$ 上定义
\[
    h(x,y)=\begin{cases}1/k,&x=q_k,\\0,&x\notin\mathbb Q.\end{cases}
\]
证明 $h$ 可积且积分为 $0$。
\end{example}
\begin{solve}
给定 $\varepsilon>0$，先取 $K$ 使 $1/(K+1)<\varepsilon/2$。用总宽度小于 $\varepsilon/2$ 的有限竖条覆盖前 $K$ 条直线 $x=q_k$，把竖条边界加入网格。竖条内 $h\leq1$，其余部分 $h\leq1/(K+1)$，故可使达布上和小于 $\varepsilon$；每个有正面积的小矩形都含横坐标无理的点，达布下和为 $0$。因此积分为 $0$。这个例子在每条 $x=q_k$ 上间断，间断点集在正方形内稠密，因而并非若尔当零面积集；它说明上一充分条件不是必要条件，不能把有限零面积集并的性质直接用于可数并。
\end{solve}

\subsection{重积分的性质}
\label{subsec:properties-of-multiple-integrals}

\begin{theorem}
\label{thm:multiple-integral-algebra-order}
设 $f,g$ 在可求体积的有界闭区域 $D$ 上可积，$a,b\in\mathbb R$。则 $af+bg$、$fg$、$|f|$ 均可积，且
\[
    \int_D(af+bg)\,dV=a\int_Df\,dV+b\int_Dg\,dV,
    \qquad
    \left|\int_Df\,dV\right|\leq\int_D|f|\,dV.
\]
若 $f\leq g$，则 $\int_Df\,dV\leq\int_Dg\,dV$；特别地，当 $m\leq f\leq M$ 时，
\begin{equation}
    mV(D)\leq\int_Df\,dV\leq MV(D).
    \label{eq:multiple-integral-bounds}
\end{equation}
\end{theorem}
\begin{proof}
线性和保序性直接由黎曼和取极限得到。对任一分块，$|f|$ 的振幅不超过 $f$ 的振幅；若 $|f|\leq A$、$|g|\leq B$，则 $fg$ 的振幅至多为 $A\omega_g+B\omega_f$。将分别控制两函数振幅和的分割取公共细分，达布判据即给出乘积和绝对值的可积性。由 $-|f|\leq f\leq|f|$ 得积分的绝对值估计，再以常数函数比较即得式~\eqref{eq:multiple-integral-bounds}。
\end{proof}

\begin{proposition}
\label{prop:multiple-domain-additivity}
若 $D=\bigcup_{j=1}^mD_j$，各 $D_j$ 可求体积且两两交集体积为零，则有界函数 $f$ 在 $D$ 上可积当且仅当各限制 $f|_{D_j}$ 可积，并且
\begin{equation}
    \int_Df\,dV=\sum_{j=1}^m\int_{D_j}f\,dV.
    \label{eq:multiple-domain-additivity}
\end{equation}
这里也据此定义有限个互不相接的闭区域之并上的积分。
\end{proposition}
\begin{proof}
用所有 $\partial D_j$ 对分割作公共细分，上下和分别相加，零体积交集不贡献体积项。各区域上达布上下和之差均非负，故其和能任意小，等价于每一项都能任意小；积分的等式随之成立。
\end{proof}

在零体积集上有界地改变函数值不改变可积性及积分值：用总体积任意小的有限长方体覆盖该集合，差函数的振幅和只在这些覆盖中贡献。这里的有界性不能省略，否则所讨论的已不再是通常的黎曼积分。

\begin{theorem}[重积分第一中值定理]
\label{thm:multiple-weighted-mean-value}
设 $D$ 是连通的可求体积有界闭区域，$f\in C(D)$，$g$ 可积且不变号，则存在 $\xi\in D$，使
\begin{equation}
    \int_Dfg\,dV=f(\xi)\int_Dg\,dV.
    \label{eq:multiple-weighted-mean-value}
\end{equation}
特别地，$\int_Df\,dV=f(\xi)V(D)$。
\end{theorem}
\begin{proof}
改变 $g$ 的符号可设 $g\geq0$。令 $m=\min_Df$、$M=\max_Df$，由 $mg\leq fg\leq Mg$ 得到相应积分不等式。若 $\int_Dg\,dV=0$，则 $\int_Dfg\,dV=0$，任取 $\xi$ 即可；否则两边除以 $\int_Dg\,dV$，所得比值位于 $[m,M]$ 内。连续函数把连通集映成区间，故该比值等于某个 $f(\xi)$。
\end{proof}

连通性保证均值确由某一点取到。若积分区域有互不相接的两个分量，而 $f$ 在两分量上分别恒为 $0$、$1$，平均值可能严格介于 $0$ 与 $1$ 之间，却不是任何函数值；积分估计仍成立，中值形式则需相应条件。

\section{化重积分为累次积分}
\label{sec:reduction-of-multiple-integrals-to-iterated-integrals}

\subsection{化二重积分为累次积分}
\label{subsec:reduction-of-double-integrals-to-iterated-integrals}

重积分的定义同时细分全部坐标方向，\emph{累次积分}（iterated integral）则先固定部分坐标，把余下坐标积分后再继续积分。两者相等需要切片积分存在等条件，不能仅凭积分号的书写顺序推断。

\begin{theorem}[矩形上的累次积分公式]
\label{thm:multiple-rectangle-fubini}
设 $f$ 在矩形 $R=[a,b]\times[c,d]$ 上黎曼可积，且对每个 $x\in[a,b]$，一元积分 $F(x)=\int_c^df(x,y)\,dy$ 都存在。则 $F$ 在 $[a,b]$ 上可积，且
\begin{equation}
    \iint_Rf(x,y)\,dx\,dy
    =\int_a^b\left(\int_c^df(x,y)\,dy\right)dx.
    \label{eq:multiple-rectangle-fubini}
\end{equation}
交换 $x,y$ 后有对应结论；若两个方向的所有切片积分都存在，则两种累次积分相等。连续函数满足这些条件。
\end{theorem}
\begin{proof}
分别分割 $[a,b]$、$[c,d]$，所得网格小矩形为 $R_{ij}$，其上 $f$ 的下、上确界记为 $m_{ij},M_{ij}$。当 $x$ 属于第 $i$ 个横坐标小区间时，切片积分满足
\[
    \sum_jm_{ij}\Delta y_j\leq F(x)
    \leq\sum_jM_{ij}\Delta y_j.
\]
因此 $F$ 的下、上和被夹在 $f$ 关于该矩形网格的下、上和之间。先让二维网格足够细，$f$ 的上下和之差任意小，便知 $F$ 的上、下积分相等；其公共值又夹在趋于 $\iint_Rf$ 的上下和之间，故得公式。
\end{proof}

\begin{example}
\label{ex:multiple-integrable-bad-slices}
二重积分存在，却有一条切片不可积。取 $D=[0,1]^2$，定义
\[
    f(x,y)=\begin{cases}1,&x=1/2\text{ 且 }y\in\mathbb Q,\\0,&\text{其他}。\end{cases}
\]
\end{example}
\begin{solve}
$f$ 有界，且只在零面积线段 $x=1/2$ 上可能非零，故它在 $D$ 上可积、积分为 $0$。但 $f(1/2,y)$ 是 $[0,1]$ 上的\emph{狄利克雷函数}（Dirichlet function），每个小区间内上下确界分别为 $1,0$，其一元黎曼积分不存在。因此不能直接把该二重积分写成所有切片均存在的 $\int_0^1(\int_0^1f(x,y)\,dy)dx$。本章采用黎曼积分，这里不以忽略个别切片的其他积分理论替代其假设。
\end{solve}

\begin{example}
\label{ex:multiple-square-polynomial}
计算 $\iint_{[0,1]^2}x(x-y)^2\,dx\,dy$。
\end{example}
\begin{solve}
被积函数连续，故
\[
    \int_0^1x\left(\int_0^1(x-y)^2\,dy\right)dx
    =\int_0^1\left(x^3-x^2+\frac{x}{3}\right)dx
    =\frac1{12}.
\]
\end{solve}

\begin{definition}
\label{def:multiple-normal-domains}
设 $\varphi_1,\varphi_2\in C([a,b])$，且在 $(a,b)$ 内 $\varphi_1<\varphi_2$，则
\[
    D=\{(x,y):a\leq x\leq b,\ \varphi_1(x)\leq y\leq\varphi_2(x)\}
\]
称为 $X$ 型区域，即对 $x$ 的\emph{简单区域}（$x$-simple region）。相应地，
\[
    D=\{(x,y):c\leq y\leq d,\ \psi_1(y)\leq x\leq\psi_2(y)\}
\]
在边界函数连续、内部严格有序时称为 $Y$ 型区域。这里 $X$ 型意指先固定 $x$ 作竖直切片，计算时内层变量是 $y$。
\end{definition}

这些区域可求面积。以连续图像 $y=\varphi(x)$ 为例，把 $[a,b]$ 分成足够短的小区间，各段图像覆盖矩形的总面积不超过 $(b-a)$ 乘各段最大振幅，再加任意小的扩大覆盖误差；一致连续性使其趋于零。区域边界由两条这种图像及端点竖线组成，故面积为零。

\begin{theorem}
\label{thm:multiple-normal-domain-fubini}
若 $f$ 在上述 $X$ 型区域 $D$ 上可积，且每条竖直切片上的一元积分均存在，则
\begin{equation}
    \iint_Df(x,y)\,dx\,dy
    =\int_a^b\left(\int_{\varphi_1(x)}^{\varphi_2(x)}f(x,y)\,dy\right)dx.
    \label{eq:multiple-x-simple-integral}
\end{equation}
在 $Y$ 型区域上，相应公式为
\begin{equation}
    \iint_Df(x,y)\,dx\,dy
    =\int_c^d\left(\int_{\psi_1(y)}^{\psi_2(y)}f(x,y)\,dx\right)dy.
    \label{eq:multiple-y-simple-integral}
\end{equation}
\end{theorem}
\begin{proof}
取包含 $D$ 的矩形，将 $f$ 在矩形内的 $D$ 外部置为 $0$。由边界零面积性，此延拓仍可积；每条切片的延拓仅在相应一元区间外取零，其积分就是原切片积分。应用定理~\ref{thm:multiple-rectangle-fubini} 即得结论。
\end{proof}

若一般区域能用有限条零面积边界分成若干简单区域，就先用区域可加性分块。确定积分界限时，先把区域写成不等式，再说明固定外层变量后内层变量的起点和终点；换序改变的是同一区域的描述，被积函数仍取同一函数。

\begin{example}
\label{ex:multiple-cubic-line-domain}
设 $D$ 由 $y=0$、$y=x^3$、$x+y=2$ 围成，$f\in C(D)$，写出两种累次积分。
\end{example}
\begin{solve}
曲线与斜线的交点为 $(1,1)$。固定 $y\in[0,1]$ 时，$y^{1/3}\leq x\leq2-y$；固定 $x$ 时，上边界在 $x=1$ 处改变。因此
\[
    \iint_Df\,d\sigma
    =\int_0^1\left(\int_{y^{1/3}}^{2-y}f(x,y)\,dx\right)dy
\]
\[
    =\int_0^1\left(\int_0^{x^3}f(x,y)\,dy\right)dx
    +\int_1^2\left(\int_0^{2-x}f(x,y)\,dy\right)dx.
\]
一种顺序只需一个区域，另一种顺序必须分段，这也是选择积分次序时应先考虑的因素。
\end{solve}

\begin{example}
\label{ex:multiple-change-order-sinc}
计算 $\displaystyle I=\int_0^1\left(\int_x^{\sqrt x}\frac{\sin y}{y}\,dy\right)dx$，其中 $\sin y/y$ 在 $y=0$ 处连续延拓为 $1$。
\end{example}
\begin{solve}
积分区域为 $0\leq x\leq1$、$x\leq y\leq\sqrt x$，等价于 $0\leq y\leq1$、$y^2\leq x\leq y$。被积函数连续，换序得到
\[
    I=\int_0^1\left(\int_{y^2}^y\frac{\sin y}{y}\,dx\right)dy
    =\int_0^1(1-y)\sin y\,dy=1-\sin1.
\]
这里先对 $x$ 积分消去了分母 $y$，避免直接寻找 $\sin y/y$ 的初等原函数。
\end{solve}

\begin{example}
\label{ex:multiple-two-cylinders-volume}
求由 $x^2+y^2\leq a^2$、$x^2+z^2\leq a^2$ 确定的相交圆柱体的体积，其中 $a>0$。
\end{example}
\begin{solve}
八个卦限的体积相等。第一卦限的底面为 $0\leq x\leq a$、$0\leq y\leq\sqrt{a^2-x^2}$，顶面为 $z=\sqrt{a^2-x^2}$，故
\[
    V=8\int_0^a\left(\int_0^{\sqrt{a^2-x^2}}\sqrt{a^2-x^2}\,dy\right)dx
    =8\int_0^a(a^2-x^2)\,dx=\frac{16a^3}{3}.
\]
\end{solve}

\subsection{化三重积分为累次积分}
\label{subsec:reduction-of-triple-integrals-to-iterated-integrals}

三重积分有两种常用的切片方式。第一种先固定平面坐标 $(x,y)$，沿竖直线积分；第二种先固定高度 $z$，在整个水平截面上作二重积分。它们分别适合上下表面易描述、截面形状易描述的立体。

\begin{theorem}
\label{thm:multiple-triple-slicing}
设 $\Omega\subset\mathbb R^2$ 是可求面积的有界闭区域，$\varphi,\psi\in C(\Omega)$，在 $\Omega^\circ$ 内 $\varphi<\psi$，并令
\[
    D=\{(x,y,z):(x,y)\in\Omega,\ \varphi(x,y)\leq z\leq\psi(x,y)\}.
\]
则 $D$ 可求体积。若 $f\in C(D)$，则
\begin{equation}
    \iiint_Df\,dV
    =\iint_\Omega\left(\int_{\varphi(x,y)}^{\psi(x,y)}f(x,y,z)\,dz\right)dx\,dy.
    \label{eq:multiple-triple-line-slicing}
\end{equation}
另一种情形中，设 $D$ 是可求体积的有界闭区域，$a\leq z\leq b$，每个截面 $D_z=\{(x,y):(x,y,z)\in D\}$ 均可求面积。若 $f\in C(D)$，则
\begin{equation}
    \iiint_Df\,dV
    =\int_a^b\left(\iint_{D_z}f(x,y,z)\,dx\,dy\right)dz.
    \label{eq:multiple-triple-plane-slicing}
\end{equation}
坐标互换后有相应公式。
\end{theorem}
\begin{proof}
连续函数的图像具有三维零体积：用底面小矩形乘图像在其上的振幅作覆盖，再用一致连续性控制总体积；接触 $\partial\Omega$ 的底面块总体积趋于零，高度又有界。故第一种立体的边界，即上下图像和底面边界上方的侧面，均为零体积集，$D$ 可求体积。

将 $f$ 零延拓至包含 $D$ 的长方体。第一种情形中，对每个 $(x,y)$ 的竖直切片可积；相应切片积分是 $\Omega$ 上的连续函数，因为边界函数连续且 $f$ 一致连续，函数变化和上下限变化造成的积分误差都趋于零。将矩形情形的达布上下和夹逼推广到两个外层变量，即得式~\eqref{eq:multiple-triple-line-slicing}。第二种情形的每个水平切片也可积，同样的上下和比较表明截面积分作为 $z$ 的函数可积，并给出式~\eqref{eq:multiple-triple-plane-slicing}。截面退化为零面积集时，其积分取 $0$。
\end{proof}

\begin{example}
\label{ex:multiple-tetrahedron-integral}
在四面体 $D=\{(x,y,z):x,y,z\geq0,\ x+y+z\leq1\}$ 上计算 $\displaystyle I=\iiint_D(1+x+y+z)^{-3}\,dV$。
\end{example}
\begin{solve}
依次固定 $x,y$，得到 $0\leq z\leq1-x-y$，因此
\[
    I=\int_0^1\int_0^{1-x}\int_0^{1-x-y}
    \frac{dz\,dy\,dx}{(1+x+y+z)^3}
    =\frac12\int_0^1\int_0^{1-x}
    \left(\frac1{(1+x+y)^2}-\frac14\right)dy\,dx
\]
\[
    =\frac12\int_0^1\left(\frac1{1+x}-\frac12-\frac{1-x}{4}\right)dx
    =\frac12\log2-\frac5{16}.
\]
\end{solve}

\begin{example}
\label{ex:multiple-ellipsoid-second-moment}
设 $D=\{(x,y,z):x^2/a^2+y^2/b^2+z^2/c^2\leq1\}$，$a,b,c>0$，求 $\iiint_D(x+y+z)^2\,dV$。
\end{example}
\begin{solve}
区域关于各坐标平面反射对称，含 $xy,xz,yz$ 的混合项积分均为 $0$。固定 $x$ 时，椭圆截面的面积为 $\pi bc(1-x^2/a^2)$，所以
\[
    \iiint_Dx^2\,dV
    =\pi bc\int_{-a}^ax^2\left(1-\frac{x^2}{a^2}\right)dx
    =\frac{4\pi a^3bc}{15}.
\]
其余两项交换坐标即得，故原积分为
\[
    \frac{4\pi abc}{15}(a^2+b^2+c^2).
\]
利用对称性前应同时检查区域和被积函数；单凭区域对称不能使偶函数的积分消失。
\end{solve}

\section{重积分的变量替换}
\label{sec:change-of-variables-in-multiple-integrals}

\subsection{重积分的变量替换公式}
\label{subsec:change-of-variables-formula-for-multiple-integrals}

累次积分改变积分的先后次序，\emph{变量替换}（change of variables）则用新坐标重新描述区域和局部体积。对平面变换 $T(u,v)=(x(u,v),y(u,v))$，线性近似把边向量 $(du,0)$、$(0,dv)$ 分别映为 $T_u\,du$、$T_v\,dv$，它们张成的平行四边形面积为
\[
    |J_T(u,v)|\,du\,dv,\qquad
    J_T(u,v)=\frac{\partial(x,y)}{\partial(u,v)}
    =\begin{vmatrix}x_u&x_v\\y_u&y_v\end{vmatrix}.
\]
$J_T$ 称为\emph{雅可比行列式}（Jacobian determinant）。行列式的正负反映定向，而重积分中的体积非负，所以换元因子必须取绝对值。

\begin{theorem}[重积分的变量替换公式]
\label{thm:multiple-change-of-variables}
设 $D\subset\mathbb R^n$ 为可求体积的有界闭区域，$T:U\to T(U)$ 是包含 $D$ 的开集 $U$ 上的 $C^1$ \emph{微分同胚}（diffeomorphism），即 $T$ 一一对应、$T,T^{-1}$ 均为 $C^1$ 映射。记
\[
    T(u)=\bigl(x_1(u),\ldots,x_n(u)\bigr),\qquad
    J_T(u)=\det DT(u)=\frac{\partial(x_1,\ldots,x_n)}{\partial(u_1,\ldots,u_n)},
    \qquad \Omega=T(D).
\]
则 $\Omega$ 可求体积；若 $f$ 在 $\Omega$ 上黎曼可积，则 $(f\circ T)|J_T|$ 在 $D$ 上可积，且
\begin{equation}
    \int_\Omega f(x)\,dx_1\cdots dx_n
    =\int_D f(T(u))|J_T(u)|\,du_1\cdots du_n.
    \label{eq:multiple-change-of-variables}
\end{equation}
特别地，$V(\Omega)=\int_D|J_T(u)|\,du$。由逆函数定理，也可把上述微分同胚条件表述为：$T$ 在 $U$ 上为 $C^1$、一一且 $J_T$ 处处非零。
\end{theorem}
\begin{proof}
先说明局部体积变化。紧集附近 $DT$ 有界，最小伸缩率又有正下界；因此 $T$ 与逆映射在足够小的邻域内均满足统一的\emph{利普希茨估计}（Lipschitz estimate），即像点距离不超过原像距离乘某个固定常数。此类映射把零体积集映为零体积集：把原集合用总体积任意小的小立方体覆盖，每个像可放进边长至多为原边长固定倍数的立方体中，总体积仍可任意小。因而 $\partial\Omega=T(\partial D)$ 为零体积集，$\Omega$ 可求体积。

设 $Q\subset D^\circ$ 是边长为 $h$、中心为 $u_Q$ 的小立方体，$A=DT(u_Q)$。导数的一致连续性给出
\[
    T(u)=T(u_Q)+A(u-u_Q)+r_Q(u),\qquad
    \sup_{u\in Q}|r_Q(u)|\leq\eta(h)h,\quad \eta(h)\to0,
\]
其中估计对全部这样的 $Q$ 一致成立。记线性近似像为平行多面体 $P=T(u_Q)+A(Q-u_Q)$。由此 $T(Q)$ 包含在 $P$ 的 $\eta(h)h$ 邻域内，而其边界 $T(\partial Q)$ 落在 $\partial P$ 的同样小邻域内。$P$ 中离边界超过 $\eta(h)h$ 的内部部分为连通集，含中心 $T(u_Q)$ 且不接触 $\partial T(Q)$，所以也含于 $T(Q)$。这些内外薄层的体积至多为 $C\eta(h)h^n$，这里 $C$ 只依赖导数及逆导数的统一界。故
\begin{equation}
    V(T(Q))=|J_T(u_Q)|V(Q)+e_QV(Q),
    \qquad \sup_Q|e_Q|\to0\quad(h\to0).
    \label{eq:multiple-local-volume-change}
\end{equation}

以细立方体网格分割 $D$。接触 $\partial D$ 的部分总体积趋于零，其像的总体积也由局部利普希茨估计趋于零，故两边都可暂时舍去这些块。内部立方体的像构成 $\Omega$ 的分割主体，其直径一致趋于零，且由式~\eqref{eq:multiple-local-volume-change} 有 $V(T(Q))\geq cV(Q)$，其中 $c>0$ 与小立方体无关。记 $\operatorname{osc}_Af=\sup_Af-\inf_Af$ 为振幅，则
\[
    \sum_Q\operatorname{osc}_Q(f\circ T)V(Q)
    \leq c^{-1}\sum_Q\operatorname{osc}_{T(Q)}f\,V(T(Q))\to0.
\]
达布判据说明 $f\circ T$ 可积；连续函数 $|J_T|$ 与它的乘积也可积。

在像块中取标记点 $T(u_Q)$，将 $\Omega$ 上的黎曼和与 $D$ 上的黎曼和比较，差的绝对值至多为
\[
    \sup_\Omega|f|\,\sup_Q|e_Q|\,V(D)
\]
再加两边边界块的积分误差，故差趋于零。两种黎曼和的极限分别为式~\eqref{eq:multiple-change-of-variables} 两边，公式得证。
\end{proof}

此定理用开邻域上的微分同胚陈述，便于把导数、逆导数和边界条件说清楚。计算中还常见在坐标轴、角度接缝等零体积集上退化的变换。应先在去掉这些集合的小邻域后应用定理，再令被去掉部分趋于零；对于本章的有界可积函数，误差可由函数上界乘被去掉部分的体积控制。不能仅因某个雅可比行列式写得出来，就忽略映射的一一性或重复覆盖次数。

\subsection{利用变量替换计算重积分}
\label{subsec:evaluating-multiple-integrals-by-change-of-variables}

平面\emph{极坐标}（polar coordinates）为
\begin{equation}
    x=r\cos\theta,\qquad y=r\sin\theta,\qquad
    \left|\frac{\partial(x,y)}{\partial(r,\theta)}\right|=r,
    \qquad dx\,dy=r\,dr\,d\theta.
    \label{eq:multiple-polar-coordinates}
\end{equation}
通常取 $r\geq0$，角度选取长度不超过 $2\pi$ 的区间，内部避免重复覆盖。$r=0$ 时雅可比为零，$\theta$ 相差 $2\pi$ 时像点相同；在有界区域中可先排除一个小圆和一条窄角带，最后令其面积趋于零。半径为 $a$ 的圆盘面积于是为 $\int_0^{2\pi}\int_0^ar\,dr\,d\theta=\pi a^2$。

\begin{example}
\label{ex:multiple-lemniscate-integral}
设 $a>0$，$D$ 为双纽线 $(x^2+y^2)^2=a^2(x^2-y^2)$ 的右叶所围闭区域，计算 $\iint_D(x^2+y^2)\,dx\,dy$。
\end{example}
\begin{solve}
极坐标下边界为 $r^2=a^2\cos2\theta$，右叶对应 $-\pi/4\leq\theta\leq\pi/4$、$0\leq r\leq a\sqrt{\cos2\theta}$。故
\[
    \iint_D(x^2+y^2)\,dx\,dy
    =\int_{-\pi/4}^{\pi/4}\int_0^{a\sqrt{\cos2\theta}}r^3\,dr\,d\theta
    =\frac{a^4}{4}\int_{-\pi/4}^{\pi/4}\cos^22\theta\,d\theta
    =\frac{\pi a^4}{16}.
\]
\end{solve}

\begin{example}
\label{ex:multiple-hyperbolic-substitution}
设 $D$ 是第一象限内由 $x^2-y^2=1,9$ 和 $xy=2,4$ 围成的闭区域，计算 $\iint_D(x^2+y^2)\,dx\,dy$。
\end{example}
\begin{solve}
取 $u=x^2-y^2$、$v=2xy$，则 $D$ 对应矩形 $1\leq u\leq9$、$4\leq v\leq8$，并有
\[
    \frac{\partial(u,v)}{\partial(x,y)}=4(x^2+y^2)>0,\qquad
    x^2+y^2=\sqrt{u^2+v^2}.
\]
第一象限内取正平方根可唯一恢复 $x,y$，所以该变换一一。原积分化为
\[
    \int_1^9\int_4^8\frac14\,dv\,du=8.
\]
选择第一象限这一分支保证了逆映射唯一；若把第三象限的另一分量也纳入积分区域，需另计它的贡献。
\end{solve}

\begin{example}
\label{ex:multiple-ellipsoid-volume}
求椭球体 $x^2/a^2+y^2/b^2+z^2/c^2\leq1$ 的体积，其中 $a,b,c>0$。
\end{example}
\begin{solve}
以 $xy$ 平面上的椭圆为底，体积为 $2c\iint_{x^2/a^2+y^2/b^2\leq1}\sqrt{1-x^2/a^2-y^2/b^2}\,dx\,dy$。取广义极坐标 $x=ar\cos\theta$、$y=br\sin\theta$，雅可比为 $abr$，于是
\[
    V=2abc\int_0^{2\pi}\int_0^1r\sqrt{1-r^2}\,dr\,d\theta
    =\frac{4\pi abc}{3}.
\]
\end{solve}

\begin{example}
\label{ex:multiple-simplex-squared-angle}
在 $D=\{(x,y):x,y\geq0,\ x+y\leq1\}$ 上计算 $\displaystyle\iint_D\frac{xy}{\sqrt{x+y}}\,dx\,dy$，原点处取连续延拓值 $0$。
\end{example}
\begin{solve}
为使 $x+y$ 只依赖一个变量，取 $x=r\cos^2\theta$、$y=r\sin^2\theta$，其中 $0\leq r\leq1$、$0\leq\theta\leq\pi/2$。内部变换一一，雅可比绝对值为 $r\sin2\theta$。因此
\[
    \iint_D\frac{xy}{\sqrt{x+y}}\,dx\,dy
    =2\int_0^1r^{5/2}\,dr\int_0^{\pi/2}\cos^3\theta\sin^3\theta\,d\theta
    =2\cdot\frac27\cdot\frac1{12}=\frac1{21}.
\]
被积函数在原点的延拓连续，因为 $0\leq xy/\sqrt{x+y}\leq(x+y)^{3/2}/4\to0$。
\end{solve}

三维\emph{柱坐标}（cylindrical coordinates）仅把水平投影改写为极坐标：
\begin{equation}
    x=r\cos\theta,\qquad y=r\sin\theta,\qquad z=z,
    \qquad dV=r\,dr\,d\theta\,dz.
    \label{eq:multiple-cylindrical-coordinates}
\end{equation}
\emph{球坐标}（spherical coordinates）则同时使用到原点的距离和两个角：
\begin{equation}
    x=r\sin\varphi\cos\theta,\quad
    y=r\sin\varphi\sin\theta,\quad z=r\cos\varphi,
    \qquad dV=r^2\sin\varphi\,dr\,d\varphi\,d\theta,
    \label{eq:multiple-spherical-coordinates}
\end{equation}
其中 $r\geq0$、$0\leq\varphi\leq\pi$，$\theta$ 取一周。$\varphi$ 是从正 $z$ 轴量起的极角，$\theta$ 是水平投影的方位角；本式雅可比按 $(r,\varphi,\theta)$ 的顺序为 $r^2\sin\varphi$。球面给出 $r$ 的界，圆锥面常给出 $\varphi$ 的界，绕 $z$ 轴的半平面给出 $\theta$ 的界。

\begin{example}
\label{ex:multiple-paraboloid-volume}
求由 $z=25-x^2-y^2$ 与 $z=0$ 围成的立体体积。
\end{example}
\begin{solve}
柱坐标中 $0\leq\theta\leq2\pi$、$0\leq r\leq5$、$0\leq z\leq25-r^2$，故
\[
    V=\int_0^{2\pi}\int_0^5\int_0^{25-r^2}r\,dz\,dr\,d\theta
    =2\pi\int_0^5r(25-r^2)\,dr=\frac{625\pi}{2}.
\]
\end{solve}

\begin{example}
\label{ex:multiple-spherical-cone-centroid}
求均匀物体 $D=\{(x,y,z):x^2+y^2+z^2\leq1,\ z\geq\sqrt{x^2+y^2}\}$ 的质心。
\end{example}
\begin{solve}
设密度为 $1$，轴对称性给出 $\overline x=\overline y=0$。球坐标中 $0\leq r\leq1$、$0\leq\varphi\leq\pi/4$、$0\leq\theta\leq2\pi$，故
\[
    V=2\pi\int_0^{\pi/4}\sin\varphi\,d\varphi\int_0^1r^2\,dr
    =\frac{2\pi}{3}\left(1-\frac1{\sqrt2}\right),
\]
\[
    M_{xy}=\iiint_Dz\,dV
    =2\pi\int_0^{\pi/4}\cos\varphi\sin\varphi\,d\varphi\int_0^1r^3\,dr
    =\frac\pi8.
\]
由式~\eqref{eq:multiple-mass-centroid}，质心为 $\bigl(0,0,3(2+\sqrt2)/16\bigr)$。
\end{solve}

\section{广义重积分}
\label{sec:improper-multiple-integrals}

当积分区域无界，或被积函数在有界区域内无界时，通常的黎曼重积分不再直接适用。前者称\emph{无穷重积分}（improper multiple integral over an unbounded domain），后者称\emph{瑕重积分}（improper multiple integral with singularities），统称\emph{广义重积分}（improper multiple integral）。关键是明确以哪些有界子区域逼近原区域；指定一种特殊逼近方式可能产生另一个概念。

\subsection{无穷重积分的基本概念}
\label{subsec:basic-concepts-of-infinite-multiple-integrals}

先讨论平面情形，记 $B_R=\{(x,y):x^2+y^2\leq R^2\}$。设 $D$ 是无界闭区域，并假定充分大的圆盘截断 $D\cap B_R$、圆环截断 $D\cap(B_S\setminus B_R^\circ)$ 都是可求面积的有界闭区域或有限个这类区域之并。这样的假设保证以下截断和差集上的积分有意义。称此类位于 $D$ 内的有界闭子区域及其有限并为允许的截断区域，积分按区域可加性计算。

\begin{definition}
\label{def:multiple-improper-at-infinity}
设 $f$ 在 $D$ 的每个允许的有界闭子区域上黎曼可积。若存在 $I\in\mathbb R$，使对每个 $\varepsilon>0$ 都有 $R>0$，对任何允许的截断区域 $K\subset D$，只要 $D\cap B_R\subset K$，就有
\begin{equation}
    \left|\iint_Kf\,d\sigma-I\right|<\varepsilon,
    \label{eq:multiple-improper-unrestricted-limit}
\end{equation}
则称 $f$ 在 $D$ 上的无穷重积分\emph{收敛}，记作 $\iint_Df\,d\sigma=I$；否则称其发散。高维情形把圆盘、面积分别换成闭球、体积即可。
\end{definition}

定义中的量词是\emph[caution]{对任何包含充分大中心部分的允许截断区域}。因此收敛时圆盘截断积分必趋于 $I$，但仅有 $\lim_{R\to\infty}\iint_{D\cap B_R}f$ 存在，尚不足以证明收敛。只按对称圆盘截断取得的极限有时称为相应的\emph{柯西主值}（Cauchy principal value），必须明确所用的截断方式。

\begin{example}
\label{ex:multiple-principal-value-not-convergence}
讨论 $\displaystyle\iint_{\mathbb R^2}\frac{x}{1+x^2+y^2}\,dx\,dy$。
\end{example}
\begin{solve}
每个中心圆盘上的积分由关于 $y$ 轴的奇对称性等于 $0$，所以圆盘主值为 $0$。但取
\[
    Q_R=[R,2R]\times[-R,R],\qquad K_R=B_R\cup Q_R,
\]
则 $K_R$ 包含 $B_R$，两部分的交集面积为零，而在 $Q_R$ 上 $x\geq R$、$1+x^2+y^2\leq1+5R^2$。因此
\[
    \iint_{K_R}\frac{x}{1+x^2+y^2}\,d\sigma
    =\iint_{Q_R}\frac{x}{1+x^2+y^2}\,d\sigma
    \geq\frac{2R^3}{1+5R^2}\to+\infty.
\]
这与圆盘截断的值不一致，故无穷重积分发散。若要求每个截断区域连通，可把矩形向左稍微延长，使它与圆盘内部相接；新增部分位于 $x>0$，仍保留同一发散下界。
\end{solve}

\subsection{无穷重积分敛散性的判定}
\label{subsec:convergence-of-infinite-multiple-integrals}

\begin{definition}
\label{def:multiple-exhaustion}
若允许的有界闭子区域列 $K_1\subset K_2\subset\cdots\subset D$ 满足：对每个 $R>0$，从某项起都有 $D\cap B_R\subset K_k$，则称 $\{K_k\}$ 为 $D$ 的\emph{穷尽列}（exhaustion）。这一要求比单纯 $\bigcup_kK_k=D$ 更明确，它要求每个固定的中心截断最终整体包含于 $K_k$。
\end{definition}

\begin{theorem}
\label{thm:multiple-improper-exhaustion}
在定义~\ref{def:multiple-improper-at-infinity} 的局部可积性条件下，无穷重积分收敛，当且仅当对每个穷尽列 $\{K_k\}$，序列 $\iint_{K_k}f\,d\sigma$ 都收敛；此时所有极限相同，均等于无穷重积分值。
\end{theorem}
\begin{proof}
必要性直接由定义得到。反过来，若每个穷尽列的积分都收敛，先取圆盘截断列，记其极限为 $I$。任意两个穷尽列都可交替抽取子列，使其集合依次包含，因为每个集合有界，而另一个穷尽列最终包含它。交替所得仍是穷尽列，所以两列的极限必须一致。

若定义中的一致截断要求对 $I$ 不成立，则存在 $\varepsilon_0>0$，对任意大的 $R$ 都有包含 $D\cap B_R$ 的允许 $K$，其积分与 $I$ 相差至少 $\varepsilon_0$。逐步把 $R$ 取得大到包含上一选出的有界集合，便得到一个嵌套的坏截断穷尽列；在其间再交替插入包含前一集合的圆盘截断，会得到一个积分不能收敛的穷尽列，矛盾。
\end{proof}

\begin{theorem}[非负函数的有界性判据]
\label{thm:multiple-nonnegative-exhaustion}
若 $f\geq0$ 且满足上述局部可积性条件，则无穷重积分收敛，当且仅当存在一个穷尽列 $\{K_k\}$，使 $\{\iint_{K_k}f\,d\sigma\}$ 有上界。此时任一穷尽列均给出相同积分值；若该单调序列无界，则积分发散到 $+\infty$。
\end{theorem}
\begin{proof}
若积分收敛，则任一穷尽列积分收敛，因而有界。反过来，给定一个积分有界的穷尽列，其积分由非负性单调增加，故趋于某个有限 $I$。任一有界允许区域都包含在该列的某个集合内，所以其积分不超过 $I$。又可选 $K_k$ 使其积分大于 $I-\varepsilon$，再取 $R$ 使 $K_k\subset D\cap B_R$；于是任何包含此中心部分的截断积分都在 $[I-\varepsilon,I]$ 内，满足定义。无界情形按同一比较得到所有足够大截断积分都超过任意给定常数。
\end{proof}

\begin{proposition}[柯西尾部判据]
\label{prop:multiple-improper-cauchy}
无穷重积分收敛，当且仅当对每个 $\varepsilon>0$，存在 $R>0$，使任何位于 $D\setminus B_R$ 的允许有界子区域 $H$ 都满足
\begin{equation}
    \left|\iint_Hf\,d\sigma\right|<\varepsilon.
    \label{eq:multiple-improper-cauchy}
\end{equation}
\end{proposition}
\begin{proof}
若积分收敛，把充分大的中心截断 $C=D\cap B_R$ 与 $C\cup H$ 的积分相减，即得尾部估计。反过来，取 $H$ 为两个大圆盘截断之间的部分，知圆盘截断积分是柯西族，故趋于有限 $I$。一般截断 $K\supset D\cap B_R$ 与中心截断的积分之差，又是外部区域 $K\setminus B_R$ 上的积分，按假设可任意小；让中心半径充分大即得定义。边界圆周上的取值不影响任何有界积分。
\end{proof}

\begin{theorem}[收敛与绝对收敛等价]
\label{thm:multiple-improper-absolute}
在上述局部可积性条件下，
\begin{equation}
    \iint_Df\,d\sigma\text{ 收敛}
    \quad\Longleftrightarrow\quad
    \iint_D|f|\,d\sigma<+\infty.
    \label{eq:multiple-improper-absolute}
\end{equation}
因此，在这里采用的任意区域截断定义下，不存在条件收敛的无穷重积分。
\end{theorem}
\begin{proof}
若 $|f|$ 的积分收敛，由有界区域上的 $|\iint_Hf|\leq\iint_H|f|$ 及尾部判据立即得 $f$ 的积分收敛。

反过来，设 $f$ 的积分收敛。给定 $\varepsilon>0$，取 $R$ 使式~\eqref{eq:multiple-improper-cauchy} 对所有外部允许区域 $H$ 成立。令 $f_+=\max(f,0)$、$f_-=\max(-f,0)$。在任一有界外部区域 $A$ 内，达布判据允许将 $f_+$ 的积分从下方用有限个小长方体上的正下和逼近：舍去接触边界的块，并只保留 $\inf_Q f>0$ 的块。设这些块的并为 $H\subset A$，则对任意 $\delta>0$ 可使
\[
    \iint_Af_+\,d\sigma-\delta
    <\sum_{Q\subset H}(\inf_Q f)V(Q)
    \leq\iint_Hf\,d\sigma<\varepsilon.
\]
这里舍去边界块的误差由局部有界性与零面积边界控制，正下和的逼近由 $f_+$ 的可积性保证。令 $\delta\downarrow0$，得 $\iint_Af_+\leq\varepsilon$；对 $-f$ 同理得 $\iint_Af_-\leq\varepsilon$。故 $\iint_A|f|\leq2\varepsilon$。取 $A$ 为任意大的外部圆环截断，加上固定中心部分的有限积分，便使非负函数 $|f|$ 的圆盘截断积分有上界；由定理~\ref{thm:multiple-nonnegative-exhaustion}，$|f|$ 的积分收敛。
\end{proof}

以上论证说明，任意形状的外部小块可以分别收集正、负贡献，因而仅靠相消不能保证收敛。这与一元广义积分通常只让区间端点趋于无穷的定义具有不同的截断自由度。

\begin{proposition}[比较判别法]
\label{prop:multiple-improper-comparison}
设局部可积函数在充分远处满足 $0\leq f\leq g$。若 $\iint_Dg<+\infty$，则 $\iint_Df<+\infty$；若 $\iint_Df=+\infty$，则 $\iint_Dg=+\infty$。若 $f,g>0$ 且在 $|x|\to\infty$、$x\in D$ 时 $f(x)/g(x)\to c\in(0,+\infty)$，则二者同敛散。若 $|f|\leq g$ 且 $g$ 可积，则 $f$ 的广义重积分收敛。
\end{proposition}
\begin{proof}
在共同的充分大中心截断之外逐个比较非负截断积分，再用有界性判据。极限比较中，充分远处有 $c_1g\leq f\leq c_2g$，其中 $0<c_1<c<c_2$，两次比较即得结论。最后一项先比较 $|f|$，再用绝对收敛判据。
\end{proof}

\begin{example}
\label{ex:multiple-infinity-sector-strip}
比较 $r^{-\alpha}=(x^2+y^2)^{-\alpha/2}$ 在无界扇形与条带上的敛散性，其中 $\alpha\in\mathbb R$。
\end{example}
\begin{solve}
在 $D_1=\{(r\cos\theta,r\sin\theta):r\geq1,\ \theta_1\leq\theta\leq\theta_2\}$ 上，设 $0<\theta_2-\theta_1\leq2\pi$。非负性允许只检查扇形截断，极坐标给出
\[
    \iint_{D_1}r^{-\alpha}\,d\sigma
    =(\theta_2-\theta_1)\int_1^\infty r^{1-\alpha}\,dr.
\]
因此恰在 $\alpha>2$ 时收敛。另一方面，在 $D_2=\{(x,y):x\geq1,\ 0\leq y\leq1\}$ 上，$r/x\in[1,\sqrt2]$，所以对每个固定 $\alpha$，$r^{-\alpha}$ 与 $x^{-\alpha}$ 之比被两个正常数夹住。条带积分遂与 $\int_1^\infty x^{-\alpha}\,dx$ 同敛散，恰在 $\alpha>1$ 时收敛。相同函数的收敛阈值取决于区域在无穷远处的宽度，不能只看环境空间维数。
\end{solve}

\begin{theorem}[广义重积分的累次积分与换元]
\label{thm:multiple-improper-fubini-change}
设 $f$ 在 $D=[a,+\infty)\times[b,+\infty)$ 上连续。若
\begin{equation}
    \int_a^\infty\left(\int_b^\infty|f(x,y)|\,dy\right)dx<+\infty,
    \label{eq:multiple-improper-fubini-absolute}
\end{equation}
其中内层广义积分对每个 $x$ 存在，外层被积函数在每个有限区间上可积，则 $\iint_Df$ 收敛，且在相应内、外层积分都存在时
\[
    \iint_Df\,d\sigma
    =\int_a^\infty\left(\int_b^\infty f(x,y)\,dy\right)dx.
\]
交换积分次序后有相应结论。对于无界区域之间的 $C^1$ 微分同胚 $T:D\to\Omega$，在有限截断上满足换元条件、且截断像与逆像构成允许穷尽的条件下，广义换元公式
\begin{equation}
    \int_\Omega f(x)\,dx=\int_Df(T(u))|J_T(u)|\,du
    \label{eq:multiple-improper-change-of-variables}
\end{equation}
表示一边收敛当且仅当另一边收敛，收敛时两值相等；非负情形也允许两边同时为 $+\infty$。
\end{theorem}

累次积分部分采用一组便于直接核对的充分条件，\emph{此处仅给证明思路}：先在有限矩形上对 $|f|$ 用累次积分公式，式~\eqref{eq:multiple-improper-fubini-absolute} 给出这些非负积分的统一上界，因而 $\iint_D|f|$ 收敛。再以绝对值积分控制被截去的尾部，得到所写累次积分与重积分一致。换元部分则可直接证明：在每个允许截断上使用式~\eqref{eq:multiple-change-of-variables}，先对 $|f|$ 用非负有界性判据比较两边的穷尽积分，再对 $f$ 取极限。这样不依赖未经证明的正负相消，也不把两个已存在的条件累次积分当成绝对收敛的依据。

\begin{example}
\label{ex:multiple-two-iterates-divergent}
设 $D=[1,+\infty)^2$，$\displaystyle f(x,y)=\frac{x^2-y^2}{(x^2+y^2)^2}$。说明两个累次积分都存在仍不能保证广义重积分收敛。
\end{example}
\begin{solve}
因为
\[
    \frac{\partial}{\partial y}\frac{y}{x^2+y^2}=f(x,y),
    \qquad
    \frac{\partial}{\partial x}\frac{x}{x^2+y^2}=-f(x,y),
\]
故
\[
    \int_1^\infty\left(\int_1^\infty f(x,y)\,dy\right)dx=-\frac\pi4,
    \qquad
    \int_1^\infty\left(\int_1^\infty f(x,y)\,dx\right)dy=\frac\pi4.
\]
严格的发散判定可直接检查绝对值。取 $0<\theta_1<\theta_2<\pi/8$，半径足够大的这一扇形落在 $D$ 内，其上
\[
    |f|=\frac{\cos2\theta}{r^2},\qquad
    \int_{\theta_1}^{\theta_2}\int_{R_0}^R|f|r\,dr\,d\theta
    =\left(\int_{\theta_1}^{\theta_2}\cos2\theta\,d\theta\right)\log\frac R{R_0}\to+\infty.
\]
由定理~\ref{thm:multiple-improper-absolute}，原无穷重积分发散。
\end{solve}

\begin{example}
\label{ex:multiple-gaussian-integral}
利用二重积分计算 $\displaystyle\int_0^\infty e^{-x^2}\,dx$。
\end{example}
\begin{solve}
先对非负函数 $e^{-(x^2+y^2)}$ 用圆盘截断和极坐标计算：
\[
    \iint_{B_R}e^{-(x^2+y^2)}\,d\sigma
    =2\pi\int_0^Re^{-r^2}r\,dr
    =\pi(1-e^{-R^2})\to\pi.
\]
由非负有界性判据，全平面广义积分收敛且等于 $\pi$。再取正方形穷尽 $[-R,R]^2$，有限区域上的累次积分公式给出
\[
    \left(\int_{-R}^Re^{-x^2}\,dx\right)^2
    =\iint_{[-R,R]^2}e^{-(x^2+y^2)}\,d\sigma\to\pi.
\]
括号内非负，故其极限为 $\sqrt\pi$；由偶性
\[
    \int_0^\infty e^{-x^2}\,dx=\frac{\sqrt\pi}{2}.
\]
这一过程先证明非负二重积分有限，再用两个穷尽列的共同极限，避免预先交换尚未证明收敛的无穷积分。
\end{solve}

\subsection{瑕重积分}
\label{subsec:singular-multiple-integrals}

设 $D\subset\mathbb R^2$ 是可求面积的有界闭区域，$x_0\in D$，$f$ 在 $D\setminus\{x_0\}$ 上有定义，并在每个避开 $x_0$ 的允许闭子区域上黎曼可积。若 $f$ 在 $x_0$ 附近无界，称 $x_0$ 为\emph{瑕点}（singular point）。假定充分小的去心截断 $D\setminus B_\delta(x_0)^\circ$ 以及环形部分都是允许的有限区域并，以保证下述逼近有意义。

\begin{definition}
\label{def:multiple-singular-integral}
若存在 $I\in\mathbb R$，使对每个 $\varepsilon>0$，都有 $\delta>0$，对任何满足
\[
    D\setminus B_\delta(x_0)^\circ\subset K\subset D\setminus\{x_0\}
\]
的允许闭子区域 $K$，均有 $|\iint_Kf\,d\sigma-I|<\varepsilon$，则称瑕重积分收敛，记为 $\iint_Df\,d\sigma=I$；否则称其发散。瑕点本身的函数值不参与这一定义。
\end{definition}

相应的穷尽列是嵌套的允许闭子区域 $K_k\subset D\setminus\{x_0\}$，要求每个固定的去心截断最终包含在 $K_k$ 中。与无穷重积分一样，必须控制所有允许的挖去方式；只对以瑕点为中心的圆孔检查极限，对变号函数一般不够。

\begin{theorem}
\label{thm:multiple-singular-criteria}
在上述局部可积与区域条件下，瑕重积分具有下列判据：它收敛当且仅当任一穷尽列的积分都收敛，且此时各极限相同；对非负函数，只须存在一个穷尽列使其积分有上界；对一般实值函数，收敛当且仅当 $|f|$ 的瑕重积分有限。非负函数的比较、极限比较判别法仍成立，只须把比较范围改为瑕点的充分小去心邻域。
\end{theorem}
\begin{proof}
用去心截断代替中心圆盘，将定理~\ref{thm:multiple-improper-exhaustion}、定理~\ref{thm:multiple-nonnegative-exhaustion} 中的包含关系逐项对应，分别得到穷尽列判据及非负判据。柯西条件相应变为：对每个 $\varepsilon>0$，存在 $\delta>0$，使任何位于 $D\cap(B_\delta(x_0)\setminus\{x_0\})$ 内的允许闭子区域 $H$ 都有 $|\iint_Hf|<\varepsilon$。在这样的局部尾部内，对 $f_+$、$f_-$ 仍可分别以正下和选取有限小长方体，完全重复定理~\ref{thm:multiple-improper-absolute} 的估计，得 $|f|$ 在瑕点附近的积分有限；反向则由绝对值不等式立即成立。比较判据由非负积分的单调性得到。
\end{proof}

\begin{example}
\label{ex:multiple-singular-radial-power}
讨论单位圆盘 $D=\{(x,y):x^2+y^2\leq1\}$ 上 $\displaystyle\iint_D(x^2+y^2)^{-\alpha/2}\,dx\,dy$ 的敛散性，其中 $\alpha\in\mathbb R$。
\end{example}
\begin{solve}
对去心圆环用极坐标，有
\[
    \iint_{\delta\leq\sqrt{x^2+y^2}\leq1}
    (x^2+y^2)^{-\alpha/2}\,d\sigma
    =2\pi\int_\delta^1r^{1-\alpha}\,dr.
\]
被积函数非负，所以只须检查这一穷尽方式。$\alpha<2$ 时极限为 $2\pi/(2-\alpha)$；$\alpha=2$ 时为 $2\pi\log(1/\delta)\to+\infty$；$\alpha>2$ 时也趋于 $+\infty$。故收敛当且仅当 $\alpha<2$。当 $\alpha\leq0$ 时，原点处可有界延拓，这时结论也可作为通常积分理解。

也可用反演把瑕点移到无穷远。取
\[
    x=\frac{s}{s^2+t^2},\qquad y=-\frac{t}{s^2+t^2},\qquad
    \left|\frac{\partial(x,y)}{\partial(s,t)}\right|
    =\frac1{(s^2+t^2)^2}.
\]
它将 $s^2+t^2\geq1$ 映到去心单位圆盘，且原积分化为
\[
    \iint_{s^2+t^2\geq1}(s^2+t^2)^{\alpha/2-2}\,ds\,dt.
\]
由例~\ref{ex:multiple-infinity-sector-strip} 的扇形判据，它收敛当且仅当 $4-\alpha>2$，仍为 $\alpha<2$。
\end{solve}

因此 $r^{-\alpha}$ 在全平面上的广义积分对任何 $\alpha$ 都发散：原点附近要求 $\alpha<2$，无穷远处要求 $\alpha>2$，两条件不能同时满足。对应地，在 $n$ 维完整球邻域及完整球外部，径向体积增长为常数乘 $r^{n-1}\,dr$，分别得到原点附近 $\alpha<n$、无穷远处 $\alpha>n$ 的阈值；区域若在瑕点附近变窄，则须重新比较其实际宽度。

\begin{example}
\label{ex:multiple-singular-thin-domain}
说明在单位圆盘内的尖角区域上，$r^{-\alpha}$ 即使在 $2\leq\alpha<3$ 时仍可能可积。
\end{example}
\begin{solve}
取条带 $G=\{(s,t):s\geq1,\ 0\leq t\leq1\}$，用上一例的反演 $T$ 定义 $\Omega=T(G)\cup\{(0,0)\}$。因 $|T(s,t)|=(s^2+t^2)^{-1/2}\leq1$，$\Omega$ 位于单位圆盘内，原点在其边界。两条长边的像在原点相接，构成逐渐变窄的区域；其边界为有限条光滑弧及原点，故可求面积。换元后
\[
    \iint_\Omega r^{-\alpha}\,d\sigma
    =\iint_G(s^2+t^2)^{-(4-\alpha)/2}\,ds\,dt.
\]
由条带判据，右边恰在 $4-\alpha>1$ 即 $\alpha<3$ 时收敛，因而 $2\leq\alpha<3$ 给出所求例子。完整圆盘的阈值不能不加条件地套用于所有以原点为边界点的子区域。
\end{solve}

多个瑕点可用互不相交的小邻域分别处理；在去掉这些邻域的紧区域上，若函数连续，则只有通常积分问题。更一般地，若瑕点集是零面积闭集 $S$，例如一条连续函数图像，定义中可用 $S$ 的任意充分小邻域替代点的圆邻域，并要求对所有允许的避开 $S$ 的闭子区域有同一极限。判别时应考察奇性沿整条集合的累积，不能只把某一个点的径向判据照搬过去。

\begin{example}
\label{ex:multiple-singular-curve}
设 $a\in C([0,1])$ 且 $a(x)>0$，$D=\{(x,y):0\leq x\leq1,\ 0\leq y\leq a(x)\}$。讨论 $\displaystyle\iint_D|y-a(x)|^{-\alpha}\,dx\,dy$，其中瑕点集为上边界图像。
\end{example}
\begin{solve}
由紧性 $m=\min a>0$。对 $0<\varepsilon<m$，取截断区域 $D_\varepsilon=\{0\leq x\leq1,\ 0\leq y\leq a(x)-\varepsilon\}$，被积函数在其上连续，故可化为累次积分。令 $t=a(x)-y$，得
\[
    \iint_{D_\varepsilon}|y-a(x)|^{-\alpha}\,d\sigma
    =\int_0^1\left(\int_\varepsilon^{a(x)}t^{-\alpha}\,dt\right)dx.
\]
这些区域随 $\varepsilon\downarrow0$ 构成避开上边界的穷尽，且函数非负。$\alpha<1$ 时，右边等于
\[
    \frac1{1-\alpha}\int_0^1a(x)^{1-\alpha}\,dx
    -\frac{\varepsilon^{1-\alpha}}{1-\alpha},
\]
极限有限；$\alpha=1$ 时积分为 $\int_0^1\log a(x)\,dx-\log\varepsilon\to+\infty$；$\alpha>1$ 时 $\varepsilon^{1-\alpha}/(\alpha-1)$ 使积分趋于 $+\infty$。因此恰在 $\alpha<1$ 时收敛，积分值为 $\frac1{1-\alpha}\int_0^1a(x)^{1-\alpha}\,dx$。这里上边界的高度只要求连续，一元换元是在每条固定 $x$ 的切片内进行，未对 $a$ 求导。
\end{solve}

\section*{习题十五}
\addcontentsline{toc}{section}{习题十五}
\fi

\selectchapter{16}
\ifcompilechapter
\chapter{曲线积分与曲面积分}
\label{chap:line-and-surface-integrals}

% 本章正文依据教材第 188--261 页（分章 PDF 第 1--74 页，止于习题十六之前）。
重积分是在平面或空间区域上累积一个量，曲线积分与曲面积分则把积分的范围换成曲线和曲面。细线的质量、薄壳的质量属于沿长度或面积累积数量的模型；变力做功、流体通量则还要考虑切向或法向。这一区别产生了两类积分：\emph{第一型积分}（integral of the first kind）以弧长或面积为权重，\emph{第二型积分}（integral of the second kind）以有向位移或有向面积为权重。本章先分别建立四种积分，再用格林、高斯和斯托克斯公式联系区域内部与边界，最后讨论微分形式、路径无关性及场的局部性质。

以下曲线、曲面通常由有限个光滑部分组成，接合处只在端点或边界上重合；被积函数取实值，向量场取值于相应的欧氏空间。光滑曲线的正则参数表示满足 $\boldsymbol r\in C^1$ 且 $\boldsymbol r'\neq0$，光滑曲面片的正则参数表示满足 $\boldsymbol r\in C^1$ 且 $\boldsymbol r_u\times\boldsymbol r_v\neq0$。闭曲线的首尾重合以及曲面参数图之间的边界重合均允许，但计算时不能重复计入一段弧或一块曲面。

\section{第一型曲线积分}
\label{sec:line-integrals-of-the-first-kind}

\subsection{第一型曲线积分的定义}
\label{subsec:definition-of-line-integrals-of-the-first-kind}

设 $\Gamma$ 是端点为 $A,B$ 的连续简单曲线，按曲线上的先后次序取分点 $A=A_0,A_1,\ldots,A_N=B$。连接相邻分点所得折线的长度为 $\sum_{k=1}^N|A_k-A_{k-1}|$；若这些长度的上确界有限，则称 $\Gamma$ \emph{可求长}（rectifiable），该上确界就是弧长 $L$。可求长曲线的每一段子弧也可求长，故可用小弧段长度 $\Delta s_k$ 作为分割的权重。若细线在点 $M$ 的线密度为 $\rho(M)$，则 $\rho(M_k)\Delta s_k$ 近似于第 $k$ 段细线的质量，从而得到以下定义。

\begin{definition}
\label{def:line-integral-first-kind}
设 $\Gamma$ 为可求长简单曲线，$f:\Gamma\to\mathbb R$。对任意弧段分割 $\mathcal P$，在第 $k$ 段子弧 $\Gamma_k$ 上任取 $M_k$，并记 $\lambda(\mathcal P)=\max_k\Delta s_k$。若存在 $I\in\mathbb R$，使对每个 $\varepsilon>0$，都存在 $\delta>0$，当 $\lambda(\mathcal P)<\delta$ 时，对任意这样的分割及取点均有
\[
    \left|\sum_{k=1}^N f(M_k)\Delta s_k-I\right|<\varepsilon,
\]
则称 $I$ 为 $f$ 沿 $\Gamma$ 的\emph{第一型曲线积分}，记为
\[
    I=\int_\Gamma f\,\mathrm ds
      =\int_\Gamma f(x,y,z)\,\mathrm ds.
\]
平面曲线只需省去坐标 $z$；有限条子弧组成的分段曲线按各段积分之和定义。
\end{definition}

由定义，$\int_\Gamma1\,\mathrm ds=L$；若 $\rho\geq0$ 为线密度，则质量为 $m=\int_\Gamma\rho\,\mathrm ds$。对于平面曲线上的非负函数 $f$，还可把积分理解为以该曲线为底、沿垂直于曲线所在平面的方向立起、高为 $f$ 的柱形曲面面积。这里的长度权重始终非负，所以把曲线方向反过来并不改变积分。

\begin{proposition}
\label{prop:first-line-integral-properties}
第一型曲线积分对被积函数具有线性性，对只在端点相交的有限条积分子弧具有可加性。在相关积分存在时，$f\leq g$ 蕴含 $\int_\Gamma f\,\mathrm ds\leq\int_\Gamma g\,\mathrm ds$，且
\[
    \left|\int_\Gamma f\,\mathrm ds\right|
    \leq\int_\Gamma|f|\,\mathrm ds
    \leq L\sup_{M\in\Gamma}|f(M)|.
\]
其中第一项比较要求 $|f|$ 可积，最后一项要求 $f$ 有界。特别地，$\int_{-\Gamma}f\,\mathrm ds=\int_\Gamma f\,\mathrm ds$，其中 $-\Gamma$ 表示反向曲线。
\end{proposition}

\begin{proof}
线性性来自有限和及极限的线性性；把分割在接合点细分后，可分别在各子弧取极限，得到可加性。由于 $\Delta s_k\geq0$，和式保持大小关系，且其绝对值不超过 $\sum_k|f(M_k)|\Delta s_k$；令网格趋于零即得估计。反向只把子弧次序颠倒，既不改变取值也不改变长度，故积分不变。
\end{proof}

\subsection{第一型曲线积分的存在性与计算公式}
\label{subsec:existence-and-evaluation-of-line-integrals-of-the-first-kind}

\begin{theorem}
\label{thm:first-line-integral-arclength}
设 $\Gamma$ 为长度为 $L$ 的可求长简单曲线，$f$ 在 $\Gamma$ 上连续。以从端点 $A$ 起算的弧长 $s$ 为参数，记曲线为 $\boldsymbol r=\boldsymbol r(s)$，$0\leq s\leq L$，则第一型曲线积分存在，且
\begin{equation}
\label{eq:first-line-integral-arclength}
    \int_\Gamma f\,\mathrm ds=\int_0^L f(\boldsymbol r(s))\,\mathrm ds.
\end{equation}
\end{theorem}

\begin{proof}
弧长参数表示连续，故 $f(\boldsymbol r(s))$ 在 $[0,L]$ 上连续。曲线分割对应于 $0=s_0<\cdots<s_N=L$，且 $\Delta s_k=s_k-s_{k-1}$；在子弧上取 $M_k$ 等价于取 $s_k^*\in[s_{k-1},s_k]$，使 $M_k=\boldsymbol r(s_k^*)$。定义中的和式于是恰为连续函数 $f\circ\boldsymbol r$ 的黎曼和，网格也相同，故收敛到式~\eqref{eq:first-line-integral-arclength} 的右端。
\end{proof}

\begin{theorem}
\label{thm:first-line-integral-parametric}
设 $\Gamma$ 有正则参数表示 $\boldsymbol r(t)=(x(t),y(t),z(t))$，$a\leq t\leq b$，且 $f$ 在 $\Gamma$ 上连续，则
\begin{equation}
\label{eq:first-line-integral-parametric}
    \int_\Gamma f\,\mathrm ds
    =\int_a^b f(\boldsymbol r(t))\,|\boldsymbol r'(t)|\,\mathrm dt,
    \qquad
    \mathrm ds=\sqrt{x'(t)^2+y'(t)^2+z'(t)^2}\,\mathrm dt.
\end{equation}
对分段光滑曲线逐段使用此式。若平面曲线为 $y=\varphi(x)$，$a\leq x\leq b$，则 $\mathrm ds=\sqrt{1+\varphi'(x)^2}\,\mathrm dx$。
\end{theorem}

\begin{proof}
令 $s(t)=\int_a^t|\boldsymbol r'(\tau)|\,\mathrm d\tau$，则 $s'(t)=|\boldsymbol r'(t)|>0$。在式~\eqref{eq:first-line-integral-arclength} 中作一元换元即得结论。若改用反向参数，换元中的方向变化由速度的绝对值抵消，因此公式与第一型积分的方向无关性一致。
\end{proof}

\begin{example}
\label{ex:first-line-integral-offset-circle}
设 $\Gamma$ 为圆周 $x^2+y^2=x$，计算 $\int_\Gamma\sqrt{1-x^2-y^2}\,\mathrm ds$。
\end{example}
\begin{solve}
取 $x=(1+\cos t)/2$、$y=\sin t/2$，$0\leq t\leq2\pi$，则 $\mathrm ds=\mathrm dt/2$，而圆周上 $1-x^2-y^2=1-x=(1-\cos t)/2$。所以
\[
    \int_\Gamma\sqrt{1-x^2-y^2}\,\mathrm ds
    =\frac12\int_0^{2\pi}\left|\sin\frac t2\right|\,\mathrm dt=2.
\]
也可利用关于 $x$ 轴的对称性，将上半圆写成 $y=\sqrt{x-x^2}$，得到 $2\int_0^1\sqrt{1-x}\,[2\sqrt{x-x^2}]^{-1}\,\mathrm dx=2$。后一写法在端点处是瑕积分，需要按极限理解，不能把参数图在端点的导数发散误认为原曲线不光滑。
\end{solve}

\begin{example}
\label{ex:first-line-integral-semicircle-boundary}
设 $\Gamma$ 为上半单位圆盘的整个边界，计算 $\int_\Gamma(x+y)^2\,\mathrm ds$。
\end{example}
\begin{solve}
边界由上半圆弧与直径组成。圆弧上 $x^2+y^2=1$，而 $xy$ 关于 $y$ 轴反号，故其贡献为 $\int_{\text{圆弧}}(1+2xy)\,\mathrm ds=\pi$。直径上 $y=0$，其贡献为 $\int_{-1}^1x^2\,\mathrm dx=2/3$，因此积分为 $\pi+2/3$。对称性只在被积函数、弧长权重和积分曲线同时相容时才能使用。
\end{solve}

\section{第二型曲线积分}
\label{sec:line-integrals-of-the-second-kind}

\subsection{第二型曲线积分的定义}
\label{subsec:definition-of-line-integrals-of-the-second-kind}

设力场 $\boldsymbol F=(P,Q,R)$ 作用下的质点沿曲线从 $A$ 移到 $B$。在一小段子弧上用常力和弦位移近似，所做的功为 $\boldsymbol F(M_k)\cdot(A_k-A_{k-1})$。它不仅取决于弧段长度，还取决于位移方向，这就引出第二型曲线积分。

\begin{definition}
\label{def:line-integral-second-kind}
设 $\Gamma$ 为取定方向的分段光滑曲线，$\boldsymbol F=(P,Q,R)$ 在其上有定义。沿指定方向依次取分点 $A_0,\ldots,A_N$，在各子弧上任取 $M_k$，记 $A_k-A_{k-1}=(\Delta x_k,\Delta y_k,\Delta z_k)$，$\lambda(\mathcal P)=\max_k\Delta s_k$。若下列极限存在，且与分割及取点无关，则称为\emph{第二型曲线积分}：
\begin{equation}
\label{eq:second-line-integral-definition}
    \int_\Gamma\boldsymbol F\cdot\mathrm d\boldsymbol r
    =\int_\Gamma P\,\mathrm dx+Q\,\mathrm dy+R\,\mathrm dz
    =\lim_{\lambda(\mathcal P)\to0}
      \sum_{k=1}^N\bigl(P(M_k)\Delta x_k+Q(M_k)\Delta y_k+R(M_k)\Delta z_k\bigr).
\end{equation}
这里的极限同样要求对所有充分细分割和任意取点一致成立。闭曲线仍须另行指定方向，其积分常写成 $\oint_\Gamma$。
\end{definition}

积分对向量场具有线性性，对沿途相接的有向曲线具有可加性；反向时每个坐标增量均反号，所以
\[
    \int_{-\Gamma}\boldsymbol F\cdot\mathrm d\boldsymbol r
    =-\int_\Gamma\boldsymbol F\cdot\mathrm d\boldsymbol r.
\]
若同一路段走过去再沿原路返回，两段积分相消。做功模型中的 $\boldsymbol F$ 是力，积分给出该力做的功；沿闭曲线的积分也称为向量场的\emph{环量}（circulation）。

\subsection{第二型曲线积分的存在性与计算公式}
\label{subsec:existence-and-evaluation-of-line-integrals-of-the-second-kind}

\begin{theorem}
\label{thm:second-line-integral-parametric}
设 $\boldsymbol r:[a,b]\to\mathbb R^3$ 是 $\Gamma$ 的分段正则参数表示，$t$ 增大时沿 $\Gamma$ 的指定方向运动，$\boldsymbol F$ 在 $\Gamma$ 上连续，则第二型曲线积分存在，且
\begin{equation}
\label{eq:second-line-integral-parametric}
    \int_\Gamma\boldsymbol F\cdot\mathrm d\boldsymbol r
    =\int_a^b\boldsymbol F(\boldsymbol r(t))\cdot\boldsymbol r'(t)\,\mathrm dt
    =\int_a^b\bigl(Px'+Qy'+Rz'\bigr)\,\mathrm dt,
\end{equation}
其中 $P,Q,R$ 均在 $\boldsymbol r(t)$ 处取值。在光滑部分令 $\boldsymbol\tau=\boldsymbol r'/|\boldsymbol r'|$ 为沿指定方向的单位切向量，则两类曲线积分满足
\begin{equation}
\label{eq:relation-two-line-integrals}
    \int_\Gamma\boldsymbol F\cdot\mathrm d\boldsymbol r
    =\int_\Gamma\boldsymbol F\cdot\boldsymbol\tau\,\mathrm ds.
\end{equation}
\end{theorem}

\begin{proof}
先在一个光滑参数段上讨论。将 $\Delta\boldsymbol r_k=\int_{t_{k-1}}^{t_k}\boldsymbol r'(t)\,\mathrm dt$ 代入定义中的和式，该和式与式~\eqref{eq:second-line-integral-parametric} 右端的差为
\[
    \sum_k\int_{t_{k-1}}^{t_k}
    \bigl[\boldsymbol F(M_k)-\boldsymbol F(\boldsymbol r(t))\bigr]
    \cdot\boldsymbol r'(t)\,\mathrm dt.
\]
同一子弧上两点的距离不超过该弧长度。由 $\boldsymbol F$ 在紧曲线上的一致连续性，记上述距离不超过 $\eta$ 时函数值差的模的上界为 $\omega(\eta)$，则差的绝对值不超过 $L\omega(\lambda(\mathcal P))\to0$。这证明积分存在及计算公式；再用 $\boldsymbol r'\,\mathrm dt=\boldsymbol\tau\,\mathrm ds$ 得到式~\eqref{eq:relation-two-line-integrals}。分段情形用可加性处理。
\end{proof}

\begin{notealert}[type=caution,title={长度微元与有向位移}]
第一型积分代入 $\mathrm ds=|\boldsymbol r'|\,\mathrm dt$，第二型积分代入 $\mathrm d\boldsymbol r=\boldsymbol r'\,\mathrm dt$。若所选参数的增大方向与题目相反，应反转积分上下限或在结果前加负号，不能对 $x',y',z'$ 分别取绝对值。第二型积分化为第一型积分时，被积函数也变为 $\boldsymbol F\cdot\boldsymbol\tau$，反向会使它反号。
\end{notealert}

\begin{example}
\label{ex:second-line-integral-semicircle}
设 $\Gamma$ 为上半单位圆周，从 $(1,0)$ 到 $(-1,0)$，求 $\int_\Gamma-y\,\mathrm dx+x\,\mathrm dy$。
\end{example}
\begin{solve}
取 $x=\cos t$、$y=\sin t$，$0\leq t\leq\pi$，方向与题意一致，故积分为
\[
    \int_0^\pi\bigl[(-\sin t)(-\sin t)+\cos^2t\bigr]\,\mathrm dt=\pi.
\]
若反向走同一半圆，结果为 $-\pi$。若把被积形式改成 $x\,\mathrm dx+y\,\mathrm dy$，则它在圆周上恒为零，因为位置向量与切向量正交。
\end{solve}

\begin{example}
\label{ex:second-line-integral-variable-path}
沿曲线 $y=x^a$（$a>0$）从 $(0,0)$ 到 $(1,1)$，计算 $\int_\Gamma(x^2-y^2)\,\mathrm dx-2xy\,\mathrm dy$。
\end{example}
\begin{solve}
以 $x$ 为参数，代入 $y=x^a$、$\mathrm dy=ax^{a-1}\,\mathrm dx$，得到
\[
    \int_0^1\bigl[x^2-(1+2a)x^{2a}\bigr]\,\mathrm dx
    =\frac13-1=-\frac23.
\]
当 $0<a<1$ 时，$y'(x)$ 在 $0$ 附近无界，可先在 $[\varepsilon,1]$ 上计算再取极限，或改用 $y$ 为参数。结果与 $a$ 无关；事实上该形式是 $x^3/3-xy^2$ 的全微分，其一般原理见第~\ref{sec:path-independence-of-line-integrals} 节。
\end{solve}

\begin{example}
\label{ex:second-line-integral-sphere-circle}
设 $\Gamma$ 是球面 $x^2+y^2+z^2=1$ 与平面 $x+y+z=0$ 的交线，计算 $\oint_\Gamma x\,\mathrm dx+y\,\mathrm dy+z\,\mathrm dz$。
\end{example}
\begin{solve}
球面在 $\boldsymbol r$ 处的法向量平行于 $\boldsymbol r=(x,y,z)$，而交线切向量位于球面的切平面中，故 $\boldsymbol r\cdot\boldsymbol\tau=0$，由式~\eqref{eq:relation-two-line-integrals} 得积分为 $0$。也可沿曲线对 $x^2+y^2+z^2=1$ 求微分，直接得到被积形式恒为零；两种方向的结果均为零。
\end{solve}

\section{第一型曲面积分}
\label{sec:surface-integrals-of-the-first-kind}

\subsection{曲面的面积}
\label{subsec:surface-area}

曲面面积不能直接照搬弧长的定义，任意内接三角网的面积未必趋于同一极限。以半径和高均为 $1$ 的圆柱侧面为例，将高度分成 $m$ 层，每圈取 $n$ 个等分点，相邻两圈的分点错开半格，连成的内接三角网面积为
\[
    A_{m,n}=2n\sin\frac\pi n
    \sqrt{1+m^2\left(1-\cos\frac\pi n\right)^2}.
\]
当 $m=n\to\infty$ 时，其极限为 $2\pi$；当 $m=n^3\to\infty$ 时，其极限为 $+\infty$，而两种情形中的小三角形直径都趋于零。这就是\emph{施瓦茨灯笼}（Schwarz lantern）现象：过于扁长的三角形所在平面未必逼近曲面切平面，只有网格变细还不足以保证面积逼近正确。

设曲面片由 $\boldsymbol r(u,v)=(x(u,v),y(u,v),z(u,v))$，$(u,v)\in D$ 给出，其中 $D$ 是边界分段光滑的有界闭平面区域，$\boldsymbol r$ 在其邻域内为 $C^1$，在内部一一对应且正则。小参数矩形的两条邻边分别近似为 $\boldsymbol r_u\,\mathrm du$ 与 $\boldsymbol r_v\,\mathrm dv$，它们张成的切平行四边形面积为 $|\boldsymbol r_u\times\boldsymbol r_v|\,\mathrm du\,\mathrm dv$。因此用切平面的局部面积伸缩率定义曲面面积，而不取任意内接三角网面积的极限。

\begin{definition}
\label{def:surface-area-parametric}
上述正则曲面片的面积定义为
\begin{equation}
\label{eq:surface-area-parametric}
    |S|=\iint_D J_S(u,v)\,\mathrm du\,\mathrm dv,
    \qquad J_S=|\boldsymbol r_u\times\boldsymbol r_v|.
\end{equation}
分片光滑曲面的面积为各片面积之和。参数图接缝只在边界上重合时，不影响面积积分。
\end{definition}

记
\begin{equation}
\label{eq:surface-normal-minors}
    A=\frac{\partial(y,z)}{\partial(u,v)},\qquad
    B=\frac{\partial(z,x)}{\partial(u,v)},\qquad
    C=\frac{\partial(x,y)}{\partial(u,v)},
\end{equation}
则 $\boldsymbol r_u\times\boldsymbol r_v=(A,B,C)$。再令 $E=\boldsymbol r_u\cdot\boldsymbol r_u$、$F=\boldsymbol r_u\cdot\boldsymbol r_v$、$G=\boldsymbol r_v\cdot\boldsymbol r_v$，由叉积的模长恒等式得到
\begin{equation}
\label{eq:surface-area-element}
    \mathrm dS=\sqrt{A^2+B^2+C^2}\,\mathrm du\,\mathrm dv
    =\sqrt{EG-F^2}\,\mathrm du\,\mathrm dv.
\end{equation}
这里 $F$ 是第一基本形式中的一个系数，与向量场 $\boldsymbol F$ 有别。若作正则参数替换 $(u,v)=\phi(\xi,\eta)$，则新叉积等于原叉积乘以 $\det D\phi$；模长乘以 $|\det D\phi|$，再由二重积分换元公式，面积保持不变。这也说明上述定义与正则参数的选择无关。对一般光滑曲面，先用有限个参数图分片，再应用同一公式。

对于图形曲面 $z=g(x,y)$，取 $\boldsymbol r(x,y)=(x,y,g(x,y))$，得到
\begin{equation}
\label{eq:surface-area-graph}
    \mathrm dS=\sqrt{1+g_x^2+g_y^2}\,\mathrm dx\,\mathrm dy.
\end{equation}
相应地，$x=g(y,z)$ 或 $y=g(z,x)$ 也有循环换位后的公式。球坐标等参数在少数点失去正则性时，应分片计算，或先避开这些点再取极限，不能直接声称参数在整个闭参数域上处处正则。

\begin{example}
\label{ex:surface-area-cylinder-and-sphere}
求单位圆柱面在 $0\leq z\leq1$ 之间的侧面积，并求单位球面被柱面 $x^2+y^2=x$ 截在柱内的部分的面积。
\end{example}
\begin{solve}
圆柱取参数 $\boldsymbol r(\theta,z)=(\cos\theta,\sin\theta,z)$，面积伸缩率为 $1$，所以侧面积为 $\int_0^{2\pi}\int_0^1\mathrm dz\,\mathrm d\theta=2\pi$。球面所求部分关于 $z=0$ 对称，上半球写成 $z=\sqrt{1-x^2-y^2}$，其投影满足 $(x-1/2)^2+y^2\leq1/4$。用极坐标 $x=r\cos\theta$、$y=r\sin\theta$，有 $-\pi/2\leq\theta\leq\pi/2$、$0\leq r\leq\cos\theta$，故面积为
\[
    2\int_{-\pi/2}^{\pi/2}\int_0^{\cos\theta}
    \frac r{\sqrt{1-r^2}}\,\mathrm dr\,\mathrm d\theta
    =4\int_0^{\pi/2}(1-\sin\theta)\,\mathrm d\theta=2\pi-4.
\]
球面的图形表示在赤道处出现无界导数，上式按非负瑕积分理解，亦可改用分片球面参数避免这一问题。
\end{solve}

\subsection{第一型曲面积分的定义}
\label{subsec:definition-of-surface-integrals-of-the-first-kind}

\begin{definition}
\label{def:surface-integral-first-kind}
设 $S$ 为紧的分片光滑曲面，$f:S\to\mathbb R$。将 $S$ 分成只在边界相交的小曲面 $S_k$，各片面积为 $\Delta S_k$，令 $\lambda(\mathcal P)=\max_k\operatorname{diam}(S_k)$，并在各片上任取 $M_k$。若当 $\lambda(\mathcal P)\to0$ 时，和式 $\sum_k f(M_k)\Delta S_k$ 对任意分割及取点都趋于同一有限值，则称此值为 $f$ 在 $S$ 上的\emph{第一型曲面积分}，记为
\[
    \iint_S f\,\mathrm dS.
\]
分割需在各正则参数片内对应于可求面积的参数子区域，公共边界不重复计入面积。
\end{definition}

若薄壳面密度为 $\rho\geq0$，则质量为 $m=\iint_S\rho\,\mathrm dS$；取 $f=1$ 即得曲面面积。第一型曲面积分与曲面取哪一侧无关，具有线性性、曲面可加性、保序性和估计
\[
    \left|\iint_S f\,\mathrm dS\right|
    \leq\iint_S|f|\,\mathrm dS
    \leq |S|\sup_S|f|,
\]
条件与第一型曲线积分的对应性质相同。若 $S$ 是平面区域，面积权重就退化为平面面积元，因此相应积分为通常的二重积分。

\subsection{第一型曲面积分的存在性与计算公式}
\label{subsec:existence-and-evaluation-of-surface-integrals-of-the-first-kind}

\begin{theorem}
\label{thm:first-surface-integral-parametric}
设 $S$ 是正则参数曲面片 $\boldsymbol r(D)$，$f$ 在 $S$ 上连续，则第一型曲面积分存在，且
\begin{equation}
\label{eq:first-surface-integral-parametric}
    \iint_S f\,\mathrm dS
    =\iint_D f(\boldsymbol r(u,v))J_S(u,v)\,\mathrm du\,\mathrm dv.
\end{equation}
特别地，对 $z=g(x,y)$ 有
\[
    \iint_S f\,\mathrm dS
    =\iint_D f(x,y,g(x,y))\sqrt{1+g_x^2+g_y^2}\,\mathrm dx\,\mathrm dy.
\]
\end{theorem}

\begin{proof}
记参数子区域为 $D_k$，$M_k=\boldsymbol r(u_k,v_k)$，并令 $h=f\circ\boldsymbol r$。由面积定义，$\Delta S_k=\iint_{D_k}J_S$，所以积分和与 $\iint_DhJ_S$ 的差的绝对值不超过
\[
    \sum_k\iint_{D_k}
    |h(u_k,v_k)-h(u,v)|J_S(u,v)\,\mathrm du\,\mathrm dv
    \leq |S|\,\omega_f(\lambda(\mathcal P)),
\]
其中 $\omega_f$ 为 $f$ 在紧曲面上的连续模，且 $\omega_f(\eta)\to0$。该估计对所有取点成立，故定义中的极限存在且等于右端。有限个参数片再用可加性合并。
\end{proof}

\begin{example}
\label{ex:first-surface-integral-sphere-oddness}
计算 $\iint_Sx^2z\,\mathrm dS$，其中 $S$ 分别为上半单位球面和整个单位球面。
\end{example}
\begin{solve}
在上半球令 $D=\{x^2+y^2\leq1\}$，则 $\mathrm dS=\mathrm dx\,\mathrm dy/\sqrt{1-x^2-y^2}$，故 $z$ 与面积元中的分母相消，得到
\[
    \iint_Sx^2z\,\mathrm dS
    =\iint_Dx^2\,\mathrm dx\,\mathrm dy
    =\int_0^{2\pi}\cos^2\theta\,\mathrm d\theta\int_0^1r^3\,\mathrm dr
    =\frac\pi4.
\]
反射 $(x,y,z)\mapsto(x,y,-z)$ 保持面积而使被积函数反号，所以整个球面的积分为 $0$。
\end{solve}

\begin{example}
\label{ex:first-surface-integral-paraboloid-centroid}
设薄壳 $S$ 为抛物面 $z=2-x^2-y^2$ 在 $z\geq0$ 的部分，面密度均匀，求其质心。
\end{example}
\begin{solve}
由对称性，$\bar x=\bar y=0$，而 $\bar z=(\iint_Sz\,\mathrm dS)/|S|$。投影为 $r\leq\sqrt2$，面积元为 $\sqrt{1+4r^2}\,r\,\mathrm dr\,\mathrm d\theta$，因此
\[
    |S|=2\pi\int_0^{\sqrt2}r\sqrt{1+4r^2}\,\mathrm dr=\frac{13\pi}3,
    \qquad
    \iint_Sz\,\mathrm dS
    =2\pi\int_0^{\sqrt2}(2-r^2)r\sqrt{1+4r^2}\,\mathrm dr
    =\frac{37\pi}{10}.
\]
故质心为 $(0,0,111/130)$。这里必须按曲面面积加权，不能用其平面投影上的平均高度代替薄壳质心高度。
\end{solve}

\section{第二型曲面积分}
\label{sec:surface-integrals-of-the-second-kind}

\subsection{曲面的侧}
\label{subsec:sides-of-surfaces}

\begin{definition}
\label{def:orientable-surface}
若在连通光滑曲面 $S$ 上可以连续选取单位法向量场 $\boldsymbol n$，则称 $S$ 为\emph{可定向曲面}（orientable surface），也称双侧曲面；$\boldsymbol n$ 与 $-\boldsymbol n$ 确定其两个定向或两侧。不存在这种全局连续法向选择时，称为不可定向曲面或单侧曲面。分片光滑曲面要求各光滑片的定向在接缝处相容。
\end{definition}

等价地，固定一点的法向朝向，沿曲面内任意闭路径连续移动法向，回到原点时应始终保持原朝向。莫比乌斯带（M\"obius strip）沿中线绕行一周会把法向翻转，因而不可定向；局部能选法向并不保证全局能选法向。若曲面由一个正则且一一对应的参数图给出，则 $\boldsymbol r_u\times\boldsymbol r_v$ 提供一个定向；多个图拼接时须检查相容性。

图形曲面 $z=g(x,y)$ 的上侧取
\[
    \boldsymbol n=\frac{(-g_x,-g_y,1)}{\sqrt{1+g_x^2+g_y^2}},
\]
下侧取其相反数。所谓上侧是法向 $z$ 分量为正，并非曲面本身位于 $z>0$。闭曲面围成一个有界区域时，通常指定从区域内部指向外部的外侧；若区域有内腔，内腔边界的区域外法向指入空腔。

\subsection{第二型曲面积分的定义}
\label{subsec:definition-of-surface-integrals-of-the-second-kind}

设流体速度场为 $\boldsymbol v$，单位法向为 $\boldsymbol n$，则单位时间通过面积 $\Delta S$ 的有向体积近似为 $(\boldsymbol v\cdot\boldsymbol n)\Delta S$。正值表示流体沿所选法向通过曲面，负值表示反向通过；若要计算质量流率，应把速度场换成质量通量密度 $\rho\boldsymbol v$。

\begin{definition}
\label{def:surface-integral-second-kind}
设 $S$ 为已定向的分片光滑曲面，单位法向为 $\boldsymbol n=(n_x,n_y,n_z)$，$\boldsymbol F=(P,Q,R)$ 在其上有定义。若下列积分和在曲面分割网格趋于零时，对所有分割及任意取点都趋于同一有限值，则称其为\emph{第二型曲面积分}或\emph{通量积分}（flux integral）：
\begin{equation}
\label{eq:second-surface-integral-definition}
    \iint_S\boldsymbol F\cdot\boldsymbol n\,\mathrm dS
    =\lim_{\lambda(\mathcal P)\to0}
      \sum_k\boldsymbol F(M_k)\cdot\boldsymbol n(M_k)\Delta S_k
    =\iint_S P\,\mathrm dy\,\mathrm dz+Q\,\mathrm dz\,\mathrm dx+R\,\mathrm dx\,\mathrm dy.
\end{equation}
分割在各光滑片内进行，接缝上可由相邻光滑片取法向的单侧值，不影响积分。
\end{definition}

第二型曲面积分就是数量函数 $\boldsymbol F\cdot\boldsymbol n$ 的第一型曲面积分。记号的含义为
\begin{equation}
\label{eq:oriented-surface-elements}
    \mathrm dy\,\mathrm dz=n_x\,\mathrm dS,\qquad
    \mathrm dz\,\mathrm dx=n_y\,\mathrm dS,\qquad
    \mathrm dx\,\mathrm dy=n_z\,\mathrm dS;
\end{equation}
它们是有向投影面积元，可能为负，不能一律取成投影面积的绝对值。积分具有线性性和定向相容的曲面可加性；把曲面定向反过来后积分反号。曲面有向面积元也记为 $\mathrm d\boldsymbol S=\boldsymbol n\,\mathrm dS$，从而通量可写成 $\iint_S\boldsymbol F\cdot\mathrm d\boldsymbol S$。

\subsection{第二型曲面积分的存在性与计算公式}
\label{subsec:existence-and-evaluation-of-surface-integrals-of-the-second-kind}

\begin{theorem}
\label{thm:second-surface-integral-parametric}
设 $S=\boldsymbol r(D)$ 为正则曲面片，取定方向与 $\boldsymbol r_u\times\boldsymbol r_v$ 一致，且 $\boldsymbol F=(P,Q,R)$ 在 $S$ 上连续，则第二型曲面积分存在，且
\begin{equation}
\label{eq:second-surface-integral-parametric}
    \iint_S\boldsymbol F\cdot\boldsymbol n\,\mathrm dS
    =\iint_D\boldsymbol F(\boldsymbol r(u,v))\cdot
      (\boldsymbol r_u\times\boldsymbol r_v)\,\mathrm du\,\mathrm dv
    =\iint_D(PA+QB+RC)\,\mathrm du\,\mathrm dv,
\end{equation}
其中 $A,B,C$ 见式~\eqref{eq:surface-normal-minors}。若指定方向与叉积相反，右端整体取负号。
\end{theorem}
\begin{proof}
单位法向 $\boldsymbol n=(\boldsymbol r_u\times\boldsymbol r_v)/J_S$ 连续，所以 $\boldsymbol F\cdot\boldsymbol n$ 连续，由定理~\ref{thm:first-surface-integral-parametric} 可积。代入第一型曲面积分公式，约去 $J_S$，即得结论。
\end{proof}

若 $S$ 为 $z=g(x,y)$ 并取上侧，公式化为
\begin{equation}
\label{eq:second-surface-integral-graph}
    \iint_S P\,\mathrm dy\,\mathrm dz+Q\,\mathrm dz\,\mathrm dx+R\,\mathrm dx\,\mathrm dy
    =\iint_D[-Pg_x-Qg_y+R]_{z=g(x,y)}\,\mathrm dx\,\mathrm dy.
\end{equation}
取下侧则整体反号。若以 $x=g(y,z)$ 为图形并取法向 $x$ 分量为正的一侧，则用 $[P-Qg_y-Rg_z]_{x=g(y,z)}\,\mathrm dy\,\mathrm dz$；若以 $y=g(z,x)$ 为图形并取法向 $y$ 分量为正的一侧，则用 $[-Pg_x+Q-Rg_z]_{y=g(z,x)}\,\mathrm dz\,\mathrm dx$。计算时先固定法向，再决定符号，比根据曲面位置猜正负更可靠。

\begin{example}
\label{ex:second-surface-integral-unit-sphere}
设 $S$ 为单位球面的外侧，计算向量场 $(x,y,z)$ 与 $(x^2,y^2,z^2)$ 的通量。
\end{example}
\begin{solve}
单位球面的外单位法向为 $\boldsymbol n=(x,y,z)$，因此第一个通量为 $\iint_S1\,\mathrm dS=4\pi$；第二个为 $\iint_S(x^3+y^3+z^3)\,\mathrm dS=0$，因为中心反射保持面积而使被积函数反号。这说明第二型曲面积分中的对称性必须同时考察场与法向，不能直接套用对坐标函数奇偶性的第一型积分判断。
\end{solve}

\begin{example}
\label{ex:second-surface-integral-tetrahedron}
设 $S$ 为三个坐标平面与平面 $x+y+z=1$ 围成的四面体边界，取外侧，求 $(x,y,z)$ 的通量。
\end{example}
\begin{solve}
在三个坐标面上，位置向量与外法向正交，通量均为零。斜面 $z=1-x-y$ 的外侧为上侧，故由式~\eqref{eq:second-surface-integral-graph} 得其通量为
\[
    \iint_{\{x\geq0,\,y\geq0,\,x+y\leq1\}}
    [x+y+(1-x-y)]\,\mathrm dx\,\mathrm dy=\frac12.
\]
总通量为 $1/2$；这也将在高斯公式中化为常数散度 $3$ 与体积 $1/6$ 的乘积。
\end{solve}

\section{各类积分之间的联系}
\label{sec:relations-among-types-of-integrals}

牛顿--莱布尼茨公式 $\int_a^bf'(x)\,\mathrm dx=f(b)-f(a)$ 将区间内部导数的积分化为端点函数值之差。格林、高斯和斯托克斯公式延续这一思想：区域或曲面上某种导数的积分，等于其有向边界上原来函数或向量场的积分。运用这些公式时，边界必须完整，定向必须匹配，被积场必须在所涉及的整个区域内满足光滑性条件。

\subsection{格林公式}
\label{subsec:green-formula}

设 $D\subset\mathbb R^2$ 为有界闭区域，边界由有限条互不相交的分段光滑简单闭曲线组成。边界的\emph{正向}（positive orientation）规定为沿边界前进时区域位于左侧，所以外边界为逆时针方向，内孔边界为顺时针方向。只有一条外边界时区域没有平面孔洞；多条边界的积分必须按各自相对于 $D$ 的正向相加。

\begin{theorem}[格林公式]
\label{thm:green-formula}
设 $D$ 如上，$P,Q$ 在包含 $D$ 的某个开集内具有连续一阶偏导数，则
\begin{equation}
\label{eq:green-formula}
    \oint_{\partial D}P\,\mathrm dx+Q\,\mathrm dy
    =\iint_D\left(\frac{\partial Q}{\partial x}-\frac{\partial P}{\partial y}\right)
      \,\mathrm dx\,\mathrm dy.
\end{equation}
这称为\emph{格林公式}（Green's theorem），左端取 $D$ 的全部正向边界。
\end{theorem}

\begin{proof}[对标准简单区域及其有限剖分的证明]
若 $D=\{a\leq x\leq b,\ \varphi_1(x)\leq y\leq\varphi_2(x)\}$，则正向边界的下边从左向右、上边从右向左，竖直边上的 $\mathrm dx=0$，所以
\[
    \oint_{\partial D}P\,\mathrm dx
    =\int_a^b[P(x,\varphi_1(x))-P(x,\varphi_2(x))]\,\mathrm dx
    =-\iint_D P_y\,\mathrm dx\,\mathrm dy.
\]
最后一步是在固定 $x$ 时应用牛顿--莱布尼茨公式。若 $D=\{c\leq y\leq d,\ \psi_1(y)\leq x\leq\psi_2(y)\}$，同理得
\[
    \oint_{\partial D}Q\,\mathrm dy
    =\int_c^d[Q(\psi_2(y),y)-Q(\psi_1(y),y)]\,\mathrm dy
    =\iint_D Q_x\,\mathrm dx\,\mathrm dy.
\]
当区域同时具有两种表示时，两式相加即得结论。若区域能有限剖分为这类小区域，则在每块上应用公式并相加；内部公共边界在相邻两块中方向相反，积分恰好相消，只剩原区域的正向边界。带孔区域的内边界也由同样的定向规则保留下来。
\end{proof}

对定理所容许的更一般分段光滑区域，可经局部分片与区域逼近把上面的证明推广；控制误差需要用到边界的分段光滑性以及场和偏导数在紧邻域上的连续有界性。这里与教材一样不展开该一般化步骤，上面的直接证明限于可作相应有限剖分的区域，不能仅以“任意区域都可这样剖分”代替论证。

若 $\boldsymbol\tau=(\tau_x,\tau_y)$ 为正向单位切向量，则区域外单位法向为 $\boldsymbol n=(\tau_y,-\tau_x)$，从而 $n_x\,\mathrm ds=\mathrm dy$、$n_y\,\mathrm ds=-\mathrm dx$。把式~\eqref{eq:green-formula} 中的 $(P,Q)$ 换为 $(-Q,P)$，得到二维通量形式
\begin{equation}
\label{eq:green-flux-formula}
    \oint_{\partial D}(P,Q)\cdot\boldsymbol n\,\mathrm ds
    =\oint_{\partial D}-Q\,\mathrm dx+P\,\mathrm dy
    =\iint_D(P_x+Q_y)\,\mathrm dx\,\mathrm dy.
\end{equation}
分别取 $(P,Q)=(0,x)$、$(-y,0)$ 或 $(-y/2,x/2)$，还可得到面积公式
\begin{equation}
\label{eq:green-area-formula}
    |D|=\oint_{\partial D}x\,\mathrm dy
       =-\oint_{\partial D}y\,\mathrm dx
       =\frac12\oint_{\partial D}x\,\mathrm dy-y\,\mathrm dx.
\end{equation}

\begin{example}
\label{ex:green-exact-integral}
设 $\Gamma$ 从 $(0,0)$ 到 $(0,1)$，计算 $\int_\Gamma2xy\,\mathrm dx+(x^2-y^2)\,\mathrm dy$。
\end{example}
\begin{solve}
这里 $Q_x-P_y=2x-2x=0$，可利用格林公式把合适闭合边界上的积分化为零，再改走 $y$ 轴线段。更直接地，令 $u=x^2y-y^3/3$，则被积形式就是 $\mathrm du$。沿分段光滑路径使用链式法则，积分为 $u(0,1)-u(0,0)=-1/3$，故对任意这样的路径结果相同。全微分的计算也避免了任意原路径不一定与补上线段围成简单区域的问题。
\end{solve}

\begin{example}
\label{ex:green-angular-integral}
设 $\Gamma$ 是不经过原点的分段光滑简单闭曲线，按其内部区域的正向定向，求
\[
    \oint_\Gamma\frac{x\,\mathrm dy-y\,\mathrm dx}{x^2+y^2}.
\]
\end{example}
\begin{solve}
令 $P=-y/(x^2+y^2)$、$Q=x/(x^2+y^2)$，则在原点以外 $Q_x-P_y=0$。若原点在曲线外部，格林公式给出积分为 $0$；若原点在内部，则必须挖去一个完全位于内部的小圆盘，不能直接对含原点的区域应用公式。圆环形区域上的格林公式表明原曲线积分等于小圆周的逆时针积分。令 $x=\varepsilon\cos t$、$y=\varepsilon\sin t$，得到小圆周积分为 $\int_0^{2\pi}\mathrm dt=2\pi$，故原积分为 $2\pi$。反向定向时结果整体反号。
\end{solve}

\subsection{高斯公式}
\label{subsec:gauss-formula}

\begin{theorem}[高斯公式]
\label{thm:gauss-formula}
设 $V\subset\mathbb R^3$ 为有界闭区域，其边界由有限块相容拼接的光滑曲面组成，取区域外侧的单位法向 $\boldsymbol n$。若 $\boldsymbol F=(P,Q,R)$ 在包含 $V$ 的某个开集内为 $C^1$，则
\begin{equation}
\label{eq:gauss-formula}
    \iint_{\partial V}\boldsymbol F\cdot\boldsymbol n\,\mathrm dS
    =\iiint_V(P_x+Q_y+R_z)\,\mathrm dx\,\mathrm dy\,\mathrm dz.
\end{equation}
这称为\emph{高斯公式}（Gauss's theorem），也称散度定理（divergence theorem）。有内腔时，应计入指向空腔的内边界通量。
\end{theorem}

\begin{proof}[对标准简单区域及其有限剖分的证明]
先设 $V=\{(x,y,z):(x,y)\in D,\ \varphi_1(x,y)\leq z\leq\varphi_2(x,y)\}$，上下图形光滑，侧边为相应的竖直曲面。上边界的外法向向上，下边界向下；竖直侧边上 $n_z=0$，故
\[
    \iint_{\partial V}R\,\mathrm dx\,\mathrm dy
    =\iint_D[R(x,y,\varphi_2)-R(x,y,\varphi_1)]\,\mathrm dx\,\mathrm dy
    =\iiint_V R_z\,\mathrm dV.
\]
对沿 $x$ 和 $y$ 方向的简单区域同理有 $\iint_{\partial V}P\,\mathrm dy\,\mathrm dz=\iiint_VP_x\,\mathrm dV$ 与 $\iint_{\partial V}Q\,\mathrm dz\,\mathrm dx=\iiint_VQ_y\,\mathrm dV$。若区域同时具有三种表示，相加即得高斯公式。若可作相应有限剖分，相邻小区域在公共曲面上的外法向相反，内部通量相消，结论推广到整个区域。
\end{proof}

一般分片光滑区域的证明仍需局部分片与逼近步骤，这里只给出上述直接证明及推广思路。公式的实质是牛顿--莱布尼茨公式在各坐标方向上的累次应用，公共边界相消则解释了为什么最终只出现整体边界。

\begin{example}
\label{ex:gauss-close-hemisphere}
设 $S$ 为上半单位球面，取上侧，求向量场
\[
    \boldsymbol F=(\sin(yz)+x,\ e^{xz}+y,\ xy+z)
\]
通过 $S$ 的通量。
\end{example}
\begin{solve}
给 $S$ 补上底面圆盘 $D_0=\{z=0,\ x^2+y^2\leq1\}$ 并取下侧，则组成上半球体 $V$ 的完整外边界。由于 $P_x+Q_y+R_z=3$，总通量为 $3\operatorname{Vol}(V)=3(2\pi/3)=2\pi$。底面通量为 $-\iint_{x^2+y^2\leq1}xy\,\mathrm dx\,\mathrm dy=0$，故所求为 $2\pi$。补面法必须同时检查补面朝向并从总通量中减去补面的贡献。
\end{solve}

\begin{example}
\label{ex:gauss-point-source-flux}
设 $S=\partial V$ 为封闭分片光滑曲面的外侧，$\boldsymbol a\notin S$，求点源场
\[
    \boldsymbol F(\boldsymbol x)=\frac{\boldsymbol x-\boldsymbol a}{|\boldsymbol x-\boldsymbol a|^3}
\]
通过 $S$ 的通量。
\end{example}
\begin{solve}
令 $r=|\boldsymbol x-\boldsymbol a|$，在 $r>0$ 处直接求导得
\[
    P_x+Q_y+R_z=\frac3{r^3}-\frac{3|\boldsymbol x-\boldsymbol a|^2}{r^5}=0.
\]
若 $\boldsymbol a\notin V$，则场在 $V$ 的邻域光滑，由高斯公式通量为 $0$。若 $\boldsymbol a$ 在 $V$ 内部，挖去以它为球心、半径充分小的球 $B_\varepsilon$，在剩余区域应用高斯公式，得 $S$ 的通量等于小球面的向外通量。在小球面上 $\boldsymbol F\cdot\boldsymbol n=\varepsilon^{-2}$，故通量为 $\varepsilon^{-2}(4\pi\varepsilon^2)=4\pi$。因此结果为内部点源 $4\pi$、外部点源 $0$；奇点处不能使用“散度为零”来忽略点源贡献。
\end{solve}

\subsection{斯托克斯公式}
\label{subsec:stokes-formula}

有向曲面 $S$ 的法向决定边界 $\partial S$ 的正向：沿选定法向站立并沿边界前进时，曲面应位于左侧。这与右手规则一致；对上侧平面圆盘，边界从 $z$ 轴正向看为逆时针。曲面有孔时，各内边界方向与外边界相反；反转曲面定向时，全部边界方向也随之反转。

\begin{theorem}[斯托克斯公式]
\label{thm:stokes-formula}
设 $S$ 为紧的已定向分片 $C^2$ 正则曲面，边界由有限条分段光滑闭曲线组成，并取与曲面相容的正向。若 $\boldsymbol F=(P,Q,R)$ 在 $S$ 的某个开邻域内为 $C^1$，则
\begin{equation}
\label{eq:stokes-formula}
    \oint_{\partial S}P\,\mathrm dx+Q\,\mathrm dy+R\,\mathrm dz
    =\iint_S(R_y-Q_z)\,\mathrm dy\,\mathrm dz
      +(P_z-R_x)\,\mathrm dz\,\mathrm dx
      +(Q_x-P_y)\,\mathrm dx\,\mathrm dy.
\end{equation}
这称为\emph{斯托克斯公式}（Stokes' theorem）。记
\[
    \operatorname{curl}\boldsymbol F=(R_y-Q_z,\ P_z-R_x,\ Q_x-P_y),
\]
则右端为 $\iint_S(\operatorname{curl}\boldsymbol F)\cdot\boldsymbol n\,\mathrm dS$。
\end{theorem}

\begin{proof}
先设 $S=\boldsymbol r(D)$ 为一个 $C^2$ 正则参数片，取定向 $\boldsymbol r_u\times\boldsymbol r_v$。令
\[
    p(u,v)=\boldsymbol F(\boldsymbol r(u,v))\cdot\boldsymbol r_u,
    \qquad q(u,v)=\boldsymbol F(\boldsymbol r(u,v))\cdot\boldsymbol r_v.
\]
参数曲线积分公式及格林公式给出
\[
    \oint_{\partial S}\boldsymbol F\cdot\mathrm d\boldsymbol r
    =\oint_{\partial D}p\,\mathrm du+q\,\mathrm dv
    =\iint_D(q_u-p_v)\,\mathrm du\,\mathrm dv.
\]
链式法则中含 $\boldsymbol r_{uv}$ 与 $\boldsymbol r_{vu}$ 的两项相消，剩下
\[
    q_u-p_v
    =[\operatorname{curl}\boldsymbol F](\boldsymbol r(u,v))
      \cdot(\boldsymbol r_u\times\boldsymbol r_v).
\]
代入式~\eqref{eq:second-surface-integral-parametric} 即得结论。一般紧曲面作有限参数分片后相加，公共边界因诱导方向相反而相消。反向定向时等式两边同时反号。
\end{proof}

上述 $C^2$ 条件便于直接使用混合偏导数相等；斯托克斯公式也适用于适当的分片 $C^1$ 正则曲面，该推广可用光滑逼近处理，这里不展开。应用时可换用具有相同有向边界的另一曲面，只要场在所用曲面的邻域内光滑，就可得到相同的旋度通量；边界相同并不许可让替换曲面穿过场的奇点。

\begin{example}
\label{ex:stokes-sphere-plane-circle}
设 $\Gamma$ 为单位球面与平面 $x+y+z=\sqrt3h$（$|h|<1$）的交线，从 $z$ 轴正向看取逆时针方向，求
\[
    \oint_\Gamma(y^2-z^2)\,\mathrm dx+(z^2-x^2)\,\mathrm dy+(x^2-y^2)\,\mathrm dz.
\]
\end{example}
\begin{solve}
取平面圆盘 $S_h$，法向为 $(1,1,1)/\sqrt3$，与指定边界方向相容。向量场的旋度为 $-2(y+z,z+x,x+y)$，故
\[
    (\operatorname{curl}\boldsymbol F)\cdot\boldsymbol n
    =-\frac4{\sqrt3}(x+y+z)=-4h.
\]
圆心到原点的距离为 $|h|$，圆盘面积为 $\pi(1-h^2)$，由斯托克斯公式得到积分为 $-4\pi h(1-h^2)$。若将平面写成 $x+y+z=c$，则应同时把距离换成 $|c|/\sqrt3$，不能把平面方程中的常数直接当作到原点的距离。
\end{solve}

\section{微分形式简介}
\label{sec:introduction-to-differential-forms}

四种积分中，$P\,\mathrm dx+Q\,\mathrm dy+R\,\mathrm dz$ 自然沿曲线积分，$P\,\mathrm dy\,\mathrm dz+Q\,\mathrm dz\,\mathrm dx+R\,\mathrm dx\,\mathrm dy$ 自然在有向曲面上积分。微分形式把这些表达式的次数、反对称性和定向规律统一起来。第一型积分则使用非负长度或面积密度，不能仅靠把其中的乘积符号改成外积，就把它当成第二型积分。

\subsection{微分形式}
\label{subsec:differential-forms}

设 $U\subset\mathbb R^n$ 为开区域，坐标为 $x_1,\ldots,x_n$。在每一点，$\mathrm dx_i$ 是从切向量 $\boldsymbol v=(v_1,\ldots,v_n)$ 中取出第 $i$ 个分量的线性函数，即 $\mathrm dx_i(\boldsymbol v)=v_i$。因而 $\mathrm df=\sum_i f_{x_i}\,\mathrm dx_i$ 对向量的作用就是 $f$ 的方向微分。

\begin{definition}
\label{def:differential-k-form}
区域 $U$ 上的一个\emph{$k$ 次微分形式}（differential $k$-form）是在每点对 $k$ 个切向量作交替多线性作用的对象，其坐标表达式为
\begin{equation}
\label{eq:differential-form-standard}
    \omega=\sum_{1\leq i_1<\cdots<i_k\leq n}
    a_{i_1\cdots i_k}(x)\,
    \mathrm dx_{i_1}\wedge\cdots\wedge\mathrm dx_{i_k}.
\end{equation}
称其为 $C^r$ 形式，是指全部系数函数均为 $C^r$。$0$ 次形式就是函数；$k>n$ 时交替性迫使该形式为零。$1$ 次形式写成 $\sum_i a_i\,\mathrm dx_i$，$n$ 次形式只有 $a(x)\,\mathrm dx_1\wedge\cdots\wedge\mathrm dx_n$ 一项。
\end{definition}

基本反对称规则是交换相邻两个 $\mathrm dx_i$ 后变号，特别地 $\mathrm dx_i\wedge\mathrm dx_i=0$。在一点处，基元素在 $k$ 个向量上的取值为
\[
    (\mathrm dx_{i_1}\wedge\cdots\wedge\mathrm dx_{i_k})
    (\boldsymbol v_1,\ldots,\boldsymbol v_k)
    =\det\bigl((v_j)_{i_\ell}\bigr)_{\ell,j=1}^k,
\]
所以交换两个向量也会使结果反号。这是形式能反映有向长度、面积和体积的代数原因。在 $\mathbb R^3$ 中常按循环顺序写 $2$ 次形式为
\[
    \omega=P\,\mathrm dy\wedge\mathrm dz+Q\,\mathrm dz\wedge\mathrm dx
           +R\,\mathrm dx\wedge\mathrm dy,
\]
其中 $\mathrm dz\wedge\mathrm dx=-\mathrm dx\wedge\mathrm dz$，需与升序标准式的符号区别。

形式可以逐项相加、乘实数，也可以乘相应光滑度的函数。\emph[caution]{“有 $\binom nk$ 个基元素”是指逐点的线性空间维数，或把系数函数作为标量时的模的秩。}若把所有系数可变的形式看成实线性空间，它通常是无限维的；特别地，全部 $0$ 次形式就是函数空间，不能因为只出现基元素 $1$ 就称它为一维实线性空间。

\subsection{微分形式的外积}
\label{subsec:exterior-products-of-differential-forms}

\begin{definition}
\label{def:exterior-product}
若 $\omega=\sum_Ia_I\,\mathrm dx_I$ 为 $p$ 次形式，$\eta=\sum_Jb_J\,\mathrm dx_J$ 为 $q$ 次形式，则它们的\emph{外积}（exterior product）为
\[
    \omega\wedge\eta=\sum_{I,J}a_Ib_J\,
    \mathrm dx_{i_1}\wedge\cdots\wedge\mathrm dx_{i_p}
    \wedge\mathrm dx_{j_1}\wedge\cdots\wedge\mathrm dx_{j_q}.
\]
运算按双线性展开，再用相邻交换变号和重复指标项为零的规则化为标准形式；结果为 $p+q$ 次形式，$p+q>n$ 时恒为零。函数与形式的外积就是通常的函数乘法。
\end{definition}

\begin{proposition}
\label{prop:exterior-product-rules}
外积满足结合律、对加法的分配律及分次交换律
\[
    \omega\wedge\eta=(-1)^{pq}\eta\wedge\omega,
    \qquad \deg\omega=p,\quad\deg\eta=q.
\]
因此 $1$ 次形式的自外积为零；更一般地，奇数次形式的自外积为零，偶数次形式则未必。
\end{proposition}
\begin{proof}
结合律和分配律由基元素拼接及双线性展开得到。把 $p$ 个微分移过另外 $q$ 个微分需要 $pq$ 次相邻交换，故出现因子 $(-1)^{pq}$；逐项相加即得交换律。取 $\eta=\omega$ 且 $p$ 为奇数时，$\omega\wedge\omega=-\omega\wedge\omega$，所以为零。
\end{proof}

外积与坐标变换中的雅可比行列式直接相连。设 $(x,y)=(x(u,v),y(u,v))$ 为 $C^1$ 坐标变换，则
\begin{align*}
    \mathrm dx\wedge\mathrm dy
    &=(x_u\,\mathrm du+x_v\,\mathrm dv)\wedge(y_u\,\mathrm du+y_v\,\mathrm dv)\\
    &=(x_uy_v-x_vy_u)\,\mathrm du\wedge\mathrm dv
     =\frac{\partial(x,y)}{\partial(u,v)}\,\mathrm du\wedge\mathrm dv.
\end{align*}
同理，在三维及更高维有
\begin{equation}
\label{eq:exterior-jacobian}
    \mathrm dx_1\wedge\cdots\wedge\mathrm dx_n
    =\det D\phi\,\mathrm du_1\wedge\cdots\wedge\mathrm du_n,
    \qquad x=\phi(u).
\end{equation}
对微分形式作变量代入并把 $\mathrm dx_i$ 换成其全微分，称为把形式\emph{拉回}（pullback）到参数域。式~\eqref{eq:exterior-jacobian} 正是体积形式的拉回公式；曲面上 $\mathrm dy\wedge\mathrm dz$ 的拉回则为 $A\,\mathrm du\wedge\mathrm dv$，与式~\eqref{eq:surface-normal-minors} 一致。

若 $\phi$ 为连通区域间的 $C^1$ 微分同胚，则 $\det D\phi$ 连续且处处非零，所以符号恒定；正号表示保向，负号表示反向。有向形式换元保留带符号行列式，目标区域与参数域的定向要相容；普通重积分按非负体积计算，仍须使用 $|\det D\phi|$。外积记录了定向，并没有取消普通重积分中的绝对值要求。

\subsection{外微分}
\label{subsec:exterior-derivative}

\begin{definition}
\label{def:exterior-derivative}
对系数为 $C^1$ 的 $k$ 次形式 $\omega=\sum_Ia_I\,\mathrm dx_I$，定义其\emph{外微分}（exterior derivative）为
\begin{equation}
\label{eq:exterior-derivative-definition}
    \mathrm d\omega=\sum_I\mathrm da_I\wedge\mathrm dx_I
    =\sum_I\sum_{j=1}^n\frac{\partial a_I}{\partial x_j}
      \,\mathrm dx_j\wedge\mathrm dx_I.
\end{equation}
它是系数连续的 $k+1$ 次形式。对于 $0$ 次形式，这就是函数的全微分；对于 $n$ 次形式，外微分为零。
\end{definition}

\begin{proposition}
\label{prop:exterior-derivative-rules}
外微分具有线性性。若 $\deg\omega=p$ 且 $\omega,\eta$ 为 $C^1$ 形式，则
\begin{equation}
\label{eq:exterior-product-derivative}
    \mathrm d(\omega\wedge\eta)
    =\mathrm d\omega\wedge\eta+(-1)^p\omega\wedge\mathrm d\eta.
\end{equation}
若 $\omega$ 为 $C^2$ 形式，则 $\mathrm d^2\omega=0$。
\end{proposition}
\begin{proof}
线性性来自系数微分的线性性。对单项 $a\,\mathrm dx_I$ 和 $b\,\mathrm dx_J$ 使用 $\mathrm d(ab)=b\,\mathrm da+a\,\mathrm db$，把 $\mathrm db$ 移过 $p$ 个基微分时产生 $(-1)^p$，逐项相加即得式~\eqref{eq:exterior-product-derivative}。再次外微分时，系数中出现 $a_{x_ix_j}$；交换 $i,j$ 后混合偏导数相等，但 $\mathrm dx_i\wedge\mathrm dx_j$ 反号，两项相消，而 $i=j$ 时外积为零，故 $\mathrm d^2\omega=0$。
\end{proof}

平面上 $\omega=P\,\mathrm dx+Q\,\mathrm dy$ 满足
\[
    \mathrm d\omega=(Q_x-P_y)\,\mathrm dx\wedge\mathrm dy.
\]
在三维，若 $\eta=P\,\mathrm dx+Q\,\mathrm dy+R\,\mathrm dz$，则
\[
    \mathrm d\eta=(R_y-Q_z)\,\mathrm dy\wedge\mathrm dz
      +(P_z-R_x)\,\mathrm dz\wedge\mathrm dx
      +(Q_x-P_y)\,\mathrm dx\wedge\mathrm dy;
\]
若 $\theta=P\,\mathrm dy\wedge\mathrm dz+Q\,\mathrm dz\wedge\mathrm dx+R\,\mathrm dx\wedge\mathrm dy$，则
\[
    \mathrm d\theta=(P_x+Q_y+R_z)\,\mathrm dx\wedge\mathrm dy\wedge\mathrm dz.
\]
因此三个积分公式都可写成同一个表达式
\begin{equation}
\label{eq:generalized-stokes-formula}
    \int_M\mathrm d\omega=\int_{\partial M}\omega,
\end{equation}
其中 $M$ 是相应的紧有向区域或曲面，$\omega$ 的次数比 $M$ 的维数小 $1$，并满足对应定理中的光滑性条件，边界取诱导定向。对区间 $[a,b]$，规定 $\int_{\partial[a,b]}f=f(b)-f(a)$，则同一表达式也包括牛顿--莱布尼茨公式。这个统一记号并不替代各定理的条件和证明；它说明边界积分与外微分之间具有共同结构。

\section{曲线积分与路径的无关性}
\label{sec:path-independence-of-line-integrals}

设 $U\subset\mathbb R^3$ 为开连通区域，$\boldsymbol F=(P,Q,R)$ 在 $U$ 内连续。开连通性保证任意两点之间可用位于 $U$ 内的有限折线连接：固定一点，把可由折线到达的点组成集合，该集合及其补集在 $U$ 内都是开的，连通性便迫使前者等于 $U$。因此下面关于不同路径的比较都有实际意义。

\begin{definition}
\label{def:path-independent-and-potential}
若对于 $U$ 中任意两点 $A,B$，积分 $\int_\Gamma\boldsymbol F\cdot\mathrm d\boldsymbol r$ 对所有从 $A$ 到 $B$ 的分段光滑路径均相等，则称该积分在 $U$ 内\emph{与路径无关}（path independent）。若存在 $u\in C^1(U)$，使
\[
    \mathrm du=P\,\mathrm dx+Q\,\mathrm dy+R\,\mathrm dz,
    \qquad\text{即}\qquad \nabla u=\boldsymbol F,
\]
则称 $u$ 为这一微分形式的\emph{原函数}或\emph{势函数}（potential function），并称该形式为\emph{恰当形式}（exact form）。此时 $\boldsymbol F$ 称为保守场（conservative field）。物理学若约定力为势能的负梯度，则势能为 $-u$，应与这里的原函数符号区分。
\end{definition}

\begin{theorem}
\label{thm:path-independence-equivalences}
对开连通区域 $U$ 内的连续向量场 $\boldsymbol F$，以下三个条件等价：曲线积分在 $U$ 内与路径无关；沿 $U$ 内任意分段光滑闭路径的积分为零；存在 $u\in C^1(U)$ 使 $\nabla u=\boldsymbol F$。条件成立时，对从 $A$ 到 $B$ 的任意路径都有
\begin{equation}
\label{eq:gradient-line-integral-theorem}
    \int_A^B\boldsymbol F\cdot\mathrm d\boldsymbol r=u(B)-u(A),
\end{equation}
且势函数在连通区域上只相差一个常数。
\end{theorem}

\begin{proof}
若积分与路径无关，把闭路径与同点不动的常路径比较，即得闭路径积分为零。反过来，若所有闭路径积分均为零，任取从 $A$ 到 $B$ 的两条路径 $\Gamma_1,\Gamma_2$，则沿 $\Gamma_1$ 再沿 $-\Gamma_2$ 构成闭路径，故两积分之差为零。

现在固定 $A_0\in U$，在路径无关条件下定义
\[
    u(A)=\int_{A_0}^A\boldsymbol F\cdot\mathrm d\boldsymbol r.
\]
对 $A=(x,y,z)$ 及充分小的 $h$，连接 $A$ 与 $(x+h,y,z)$ 的线段在 $U$ 内，因而
\[
    u(x+h,y,z)-u(x,y,z)=\int_x^{x+h}P(t,y,z)\,\mathrm dt.
\]
除以 $h$ 后令 $h\to0$，由 $P$ 的连续性得到 $u_x=P$；同理 $u_y=Q$、$u_z=R$。三个偏导数连续，故 $u\in C^1(U)$ 且为所求势函数。

若反过来已知 $\nabla u=\boldsymbol F$，沿路径 $\boldsymbol r(t)$ 用链式法则，有 $\boldsymbol F(\boldsymbol r(t))\cdot\boldsymbol r'(t)=\frac{\mathrm d}{\mathrm dt}u(\boldsymbol r(t))$，逐段积分即得式~\eqref{eq:gradient-line-integral-theorem}，所以与路径无关。若 $u,v$ 是两个势函数，则 $\nabla(u-v)=0$，沿连接任意两点的折线积分可知 $u-v$ 为常数。
\end{proof}

在平面区域内，完全相同的结论适用于 $P\,\mathrm dx+Q\,\mathrm dy$，势函数条件为 $u_x=P$、$u_y=Q$。这里闭路径允许自交；采用所有闭路径的表述，可以直接用“去程加反向回程”证明等价性，而不必假定两条任意路径之并总是简单闭曲线。

\begin{definition}
\label{def:simply-connected-and-closed-form}
若开连通区域 $U$ 内的每条闭路径都能在 $U$ 内连续收缩为一点，则称 $U$ \emph{单连通}（simply connected）。对 $C^1$ 的 $1$ 次形式 $\omega$，若 $\mathrm d\omega=0$，则称其为\emph{闭形式}（closed form）。在三维，这等价于 $\operatorname{curl}\boldsymbol F=0$；在二维，等价于 $Q_x=P_y$。
\end{definition}

单连通是关于闭路径能否收缩的条件，不能在三维直接解释成“没有任何空腔”。例如球壳区域 $a<|\boldsymbol x|<b$（$0<a<b$）是单连通的，而实心环面内部不是单连通的；前者的空腔阻碍某些封闭曲面收缩，却不阻碍闭路径收缩。凸区域及星形区域都是单连通的，因为可沿连向一个固定点的线段直接收缩。

\begin{theorem}
\label{thm:simply-connected-curl-free}
设 $U\subset\mathbb R^3$ 为开连通区域，$\boldsymbol F=(P,Q,R)\in C^1(U)$。若曲线积分与路径无关，则必有
\begin{equation}
\label{eq:curl-free-conditions}
    R_y=Q_z,\qquad P_z=R_x,\qquad Q_x=P_y.
\end{equation}
若 $U$ 进一步单连通，则式~\eqref{eq:curl-free-conditions} 也是路径无关的充分条件。对平面单连通区域，充分必要条件相应为 $Q_x=P_y$。
\end{theorem}

\begin{proof}
必要性由定理~\ref{thm:path-independence-equivalences} 得到 $u_x=P$、$u_y=Q$、$u_z=R$；因 $P,Q,R\in C^1$，有 $u\in C^2$，混合偏导数相等给出式~\eqref{eq:curl-free-conditions}。

为证明充分性，先在以 $\boldsymbol a$ 为中心、包含于 $U$ 的一个球内构造局部势函数
\[
    v(\boldsymbol x)=\int_0^1
    \boldsymbol F(\boldsymbol a+t(\boldsymbol x-\boldsymbol a))
    \cdot(\boldsymbol x-\boldsymbol a)\,\mathrm dt.
\]
记 $F_i$ 为第 $i$ 个分量。连续偏导数允许在积分号下求导，而闭形式条件为 $\partial_iF_j=\partial_jF_i$，所以
\begin{align*}
    \partial_i v
    &=\int_0^1\left[F_i(\boldsymbol a+t(\boldsymbol x-\boldsymbol a))
      +t\sum_j(\partial_iF_j)(\boldsymbol a+t(\boldsymbol x-\boldsymbol a))(x_j-a_j)\right]\,\mathrm dt\\
    &=\int_0^1\frac{\mathrm d}{\mathrm dt}
      \left[tF_i(\boldsymbol a+t(\boldsymbol x-\boldsymbol a))\right]\,\mathrm dt
      =F_i(\boldsymbol x).
\end{align*}
故每个充分小球内的积分均与路径无关。

任取 $U$ 内一条分段光滑闭路径，由单连通性，存在将它连续收缩到一点的同伦。这个同伦是紧参数正方形到 $U$ 的连续映射。用上述小球覆盖其像，并把参数正方形分得足够细，使每个小格的像落在某一小球内。把各格点的像用直线段连接；共享边的端点同时属于相邻两格对应的小球，由球的凸性，该直线段在两球内均可用。每个小格对应的闭折线积分为零，相加后内部边相消；收缩端及两条侧边的贡献也相消，故原路径所对应的边界折线积分为零。原路径每小段与替代线段在同一小球内，局部势函数保证二者积分相同，因此原闭路径积分为零。由定理~\ref{thm:path-independence-equivalences} 得到全局路径无关。二维证明完全相同。
\end{proof}

\begin{notealert}[type=caution,title={局部微分条件与全局结论}]
无旋条件是局部条件，路径无关是全局条件。例~\ref{ex:green-angular-integral} 中的角形式在 $\mathbb R^2\setminus\{(0,0)\}$ 上满足 $Q_x=P_y$，却在绕原点一周时积分为 $2\pi$，所以不是全局恰当形式。单连通是使所有闭 $1$ 次形式都恰当的充分条件，但某一个具体形式恰当，并不要求其定义域单连通。
\end{notealert}

\begin{example}
\label{ex:path-independent-difference-angular-forms}
设 $E=\{(x,0):0\leq x\leq1\}$，$U=\mathbb R^2\setminus E$，证明
\[
    \omega=\frac{x\,\mathrm dy-y\,\mathrm dx}{x^2+y^2}
     -\frac{(x-1)\,\mathrm dy-y\,\mathrm dx}{(x-1)^2+y^2}
\]
在 $U$ 内与路径无关。
\end{example}
\begin{solve}
两项分别描述相对于 $(0,0)$ 和 $(1,0)$ 的角变化；闭曲线若不与线段 $E$ 相交，就不能只围住其中一个端点。为避免依赖图形判断，直接构造一个全局势函数。令
\[
    a=x(x-1)+y^2,\qquad b=-y,\qquad q=\sqrt{a^2+b^2},
    \qquad u=2\arctan\frac b{q+a}.
\]
在 $U$ 内 $q>0$ 且 $q+a>0$：若 $b=0$，则 $y=0$ 且 $x\notin[0,1]$，故 $a=x(x-1)>0$；若 $b\neq0$，则 $q>|a|$。因此 $u$ 在整个 $U$ 内光滑。用反正切求导公式有
\[
    \mathrm du=\frac{a\,\mathrm db-b\,\mathrm da}{a^2+b^2}.
\]
再由 $a^2+b^2=(x^2+y^2)((x-1)^2+y^2)$ 化简即得 $\mathrm du=\omega$，所以积分与路径无关。若只删去两个端点而保留线段内部，则小圆可仅绕一个端点，积分不再与路径无关；决定结论的是定义域和形式共同的性质。
\end{solve}

\begin{example}
\label{ex:path-independent-radial-field}
设 $f:(0,+\infty)\to\mathbb R$ 连续，$r=\sqrt{x^2+y^2+z^2}$。证明 $f(r)(x\,\mathrm dx+y\,\mathrm dy+z\,\mathrm dz)$ 在任意不含原点的开区域内有原函数。
\end{example}
\begin{solve}
令 $H(r)=\int_1^r t f(t)\,\mathrm dt$，则 $H'(r)=rf(r)$，而 $\mathrm dr=(x\,\mathrm dx+y\,\mathrm dy+z\,\mathrm dz)/r$，故
\[
    \mathrm dH(r)=f(r)(x\,\mathrm dx+y\,\mathrm dy+z\,\mathrm dz).
\]
因此 $H(r)$ 就是势函数。这里只要求 $f$ 连续，不能为了套用无旋判别而额外假设其可导；也不能漏掉构造 $H$ 时的因子 $t$。若被积形式本来是 $f(r)\,\mathrm dr$，才可直接取 $f$ 的一元原函数。
\end{solve}

求势函数通常有两种方法。第一种先把 $y,z$ 看作常数，对 $u_x=P$ 积分，写成 $u=\int P\,\mathrm dx+\varphi(y,z)$，再用 $u_y=Q$ 和 $u_z=R$ 逐步确定余下函数；“积分常数”在此必须允许依赖尚未积分的变量。第二种固定基点 $A_0$，选择位于定义域内的便于计算的路径，取 $u(A)=\int_{A_0}^A\boldsymbol F\cdot\mathrm d\boldsymbol r$。若区域允许依次沿坐标轴方向走折线，则可写成
\begin{align*}
    u(x,y,z)={}&\int_{x_0}^xP(t,y_0,z_0)\,\mathrm dt
    +\int_{y_0}^yQ(x,t,z_0)\,\mathrm dt\\
    &+\int_{z_0}^zR(x,y,t)\,\mathrm dt+C.
\end{align*}
这条折线必须完整地位于定义域内；一般带孔区域不能不加检查地使用坐标折线路径。

\begin{example}
\label{ex:potential-exponential-polynomial}
在 $\mathbb R^3$ 中求原函数，使
\[
    \mathrm du=(yze^{xyz}+3x^2)\,\mathrm dx
      +(xze^{xyz}+\sin y)\,\mathrm dy
      +(xye^{xyz}+2z)\,\mathrm dz.
\]
\end{example}
\begin{solve}
先对 $x$ 积分得 $u=e^{xyz}+x^3+\varphi(y,z)$。与 $u_y=xze^{xyz}+\sin y$ 比较，得 $\varphi=-\cos y+\psi(z)$；再与 $u_z=xye^{xyz}+2z$ 比较，得 $\psi=z^2+C$。因此
\[
    u=e^{xyz}+x^3-\cos y+z^2+C.
\]
沿 $(0,0,0)\to(x,0,0)\to(x,y,0)\to(x,y,z)$ 的坐标折线积分也得到同一结果。最后逐项求偏导数，便可直接核查原函数，且看出所有原函数只差一个常数。
\end{solve}

\section{场论简介}
\label{sec:introduction-to-field-theory}

空间区域内每点所对应的物理量称为\emph{场}（field）。温度、密度可用标量函数描述，称为\emph{数量场}或标量场（scalar field）；速度、力等还需要大小和方向，称为\emph{向量场}（vector field）。不显含时间的场称为稳定场（steady field）；把实际现象看作稳定场是一种模型假设，是否适用取决于所研究的时间尺度。以下在固定的欧氏直角坐标系中讨论，只写空间变量。

\subsection{数量场的梯度}
\label{subsec:gradient-of-scalar-fields}

\begin{definition}
\label{def:scalar-field-gradient}
设 $u:U\subset\mathbb R^3\to\mathbb R$ 在一点可微，定义其\emph{梯度}（gradient）为
\[
    \operatorname{grad}u=\nabla u=(u_x,u_y,u_z).
\]
对任意单位向量 $\boldsymbol e$，该点沿 $\boldsymbol e$ 的方向导数为 $D_{\boldsymbol e}u=\nabla u\cdot\boldsymbol e$。
\end{definition}

由柯西--施瓦茨不等式，$D_{\boldsymbol e}u\leq|\nabla u|$；当 $\nabla u\neq0$ 时，等号恰在 $\boldsymbol e=\nabla u/|\nabla u|$ 时成立。因此梯度方向为最快上升方向，最大上升率为梯度模长；反方向为最快下降方向。若梯度为零，所有方向导数都为零，此时没有唯一的最快上升方向，也不能据此断定该点为极值点。

集合 $\{\boldsymbol x:u(\boldsymbol x)=c\}$ 称为\emph{等位集}或等值集（level set）。若 $u\in C^1$ 且某点处 $\nabla u\neq0$，隐函数定理保证该点附近的等位集是光滑曲面，称为等位面。沿其中任一光滑曲线求导，得 $\nabla u\cdot\boldsymbol r'=0$，所以梯度垂直于等位面。仅有 $u\in C^1$ 并不足以保证每个等位集都是曲面，梯度非零是这一局部正则结论的关键条件。

\subsection{向量场的向量线}
\label{subsec:field-lines-of-vector-fields}

\begin{definition}
\label{def:vector-field-lines}
设 $\boldsymbol F=(P,Q,R)\in C^1(U)$。在 $\boldsymbol F\neq0$ 的区域，一条曲线若每点切线均与该点场向量平行，就称为\emph{向量线}或场线（field line）。沿场方向选择参数后，可将它表示为微分方程
\begin{equation}
\label{eq:field-line-ode}
    \boldsymbol r'(t)=\boldsymbol F(\boldsymbol r(t)),
    \qquad\text{即}\qquad x'=P,\quad y'=Q,\quad z'=R
\end{equation}
的非平衡解轨迹。
\end{definition}

常见写法 $\mathrm dx/P=\mathrm dy/Q=\mathrm dz/R$ 只是上述方程的简记；某个分量为零时不能真的除以该分量，应回到参数方程处理。由于 $C^1$ 向量场局部满足利普希茨条件，常微分方程的局部存在唯一性定理保证过每个初始点的解唯一，并可延拓为唯一的最大解。因此非零点处的最大场线不会分叉，相交的两条解轨迹实际上是同一条轨迹。这里借用常微分方程理论的结论，不在本章证明；最大解也未必对所有实数时间存在。

若 $\boldsymbol F(\boldsymbol a)=0$，则 $\boldsymbol r(t)\equiv\boldsymbol a$ 为平衡解，不能据此称它为具有非零切向量的光滑场线。场线若以零点为极限，可在无限时间内接近它，但在唯一性条件下不能在有限时间到达该点后再沿另一条轨迹离开。

\begin{example}
\label{ex:radial-vector-field-lines}
求 $\boldsymbol F(x,y,z)=(x,y,z)$ 在 $\mathbb R^3\setminus\{0\}$ 中的向量线。
\end{example}
\begin{solve}
初值为 $\boldsymbol r(t_0)=\boldsymbol a\neq0$ 的解为 $\boldsymbol r(t)=e^{t-t_0}\boldsymbol a$，故每条场线是一条从原点出发但不包含原点的射线，方向朝外。用参数解能同时处理某些坐标为零的情形；只写 $y=C_1x$、$z=C_2x$ 会遗漏 $x=0$ 的场线，也会把同一直线上的两条相反射线误并为一条。
\end{solve}

\subsection{向量场的散度}
\label{subsec:divergence-of-vector-fields}

\begin{definition}
\label{def:vector-field-divergence}
对 $C^1$ 向量场 $\boldsymbol F=(P,Q,R)$，定义其\emph{散度}（divergence）为数量场
\[
    \operatorname{div}\boldsymbol F=P_x+Q_y+R_z.
\]
\end{definition}

散度是单位体积净外流通量的局部极限。具体地，设闭球 $\overline{B_\varepsilon(\boldsymbol a)}\subset U$，球面取外侧，由高斯公式有
\begin{equation}
\label{eq:divergence-local-flux}
    \operatorname{div}\boldsymbol F(\boldsymbol a)
    =\lim_{\varepsilon\to0+}
      \frac1{\operatorname{Vol}(B_\varepsilon)}
      \iint_{\partial B_\varepsilon(\boldsymbol a)}
      \boldsymbol F\cdot\boldsymbol n\,\mathrm dS.
\end{equation}
因为右端在取极限前等于球内散度的体积平均，而连续函数在收缩球上的平均趋于中心值。通量和体积都具有几何意义，所以散度不依赖于欧氏直角坐标系的选取。

散度在一点为正时，其邻域内充分小区域具有正的净外流，称该点具有源的性质；为负时具有汇的性质。若散度在整个区域恒为零，称为\emph{无源场}（solenoidal field）。仅在一点散度为零，只能说明该点的一阶净通量密度为零，不能推出它邻近的每个小闭曲面的通量都恰为零。例~\ref{ex:gauss-point-source-flux} 中的点源场在点源位置以外无源，含奇点的闭曲面却可有非零通量，二者没有矛盾。

\subsection{向量场的旋度}
\label{subsec:curl-of-vector-fields}

\begin{definition}
\label{def:vector-field-curl}
对 $C^1$ 向量场 $\boldsymbol F=(P,Q,R)$，定义其\emph{旋度}（curl）为向量场
\[
    \operatorname{curl}\boldsymbol F
    =\operatorname{rot}\boldsymbol F
    =(R_y-Q_z,\ P_z-R_x,\ Q_x-P_y).
\]
恒满足 $\operatorname{curl}\boldsymbol F=0$ 的场称为\emph{无旋场}（irrotational field）。
\end{definition}

旋度在单位方向 $\boldsymbol n$ 上的分量，是垂直于该方向的小圆盘边界环量与圆盘面积之比的极限。设 $D_\varepsilon$ 是以 $\boldsymbol a$ 为中心、法向为 $\boldsymbol n$、半径为 $\varepsilon$ 的圆盘，其边界按右手规则定向。由斯托克斯公式及旋度的连续性，
\begin{equation}
\label{eq:curl-local-circulation}
    (\operatorname{curl}\boldsymbol F)(\boldsymbol a)\cdot\boldsymbol n
    =\lim_{\varepsilon\to0+}
      \frac1{\pi\varepsilon^2}
      \oint_{\partial D_\varepsilon}\boldsymbol F\cdot\mathrm d\boldsymbol r.
\end{equation}
不同方向的分量由此共同确定旋度；固定空间定向后，这一描述与所用的正交坐标无关。局部环量密度非零体现旋转趋势，但旋度并不表示每个质点都沿闭合圆周运动。

例如 $\boldsymbol F=(-\Omega y,\Omega x,0)$ 表示绕 $z$ 轴的刚性转动速度场，其中 $\Omega$ 为常数。半径 $\varepsilon$ 的水平圆周环量为 $2\pi\Omega\varepsilon^2$，故 $\operatorname{curl}\boldsymbol F=(0,0,2\Omega)$；旋度模长是该模型角速度绝对值的两倍。路径无关必然推出无旋，反向结论还需第~\ref{sec:path-independence-of-line-integrals} 节中的定义域条件。

\subsection{一些重要算子}
\label{subsec:important-field-operators}

把偏导数形式地排成向量，得到\emph{哈密尔顿算子}（Hamilton operator），又称 nabla 或 del 算子：
\[
    \nabla=\left(\frac{\partial}{\partial x_1},\ldots,
    \frac{\partial}{\partial x_n}\right).
\]
它是作用于函数的微分算子，不能在乘积中把它当成普通常向量。对可微函数 $f,g$ 及常数 $c$，有 $\nabla(cf)=c\nabla f$、$\nabla(f+g)=\nabla f+\nabla g$ 和 $\nabla(fg)=f\nabla g+g\nabla f$。在三维，梯度、散度和旋度分别记为
\[
    \operatorname{grad}f=\nabla f,\qquad
    \operatorname{div}\boldsymbol F=\nabla\cdot\boldsymbol F,\qquad
    \operatorname{curl}\boldsymbol F=\nabla\times\boldsymbol F.
\]
叉积的形式行列式按第一行展开理解为各分量的偏导数组合，并不是把微分算子当作数的任意行列式运算。

\begin{definition}
\label{def:laplacian-and-harmonic-function}
在 $\mathbb R^n$ 中，\emph{拉普拉斯算子}（Laplacian）定义为
\begin{equation}
\label{eq:laplace-operator}
    \Delta=\nabla\cdot\nabla
    =\sum_{i=1}^n\frac{\partial^2}{\partial x_i^2}.
\end{equation}
若 $u\in C^2(U)$ 且 $\Delta u=0$，则称 $u$ 为 $U$ 上的\emph{调和函数}（harmonic function）。一维连通区间内调和函数恰为仿射函数 $u(x)=ax+b$；二维时 $e^x\cos y$、$2xy$ 都是调和函数。
\end{definition}

\begin{proposition}
\label{prop:basic-vector-calculus-identities}
若 $u\in C^2(U)$、$\boldsymbol F\in C^2(U;\mathbb R^3)$，则
\begin{equation}
\label{eq:vector-calculus-identities}
    \operatorname{curl}(\nabla u)=0,\qquad
    \operatorname{div}(\operatorname{curl}\boldsymbol F)=0,
    \qquad \operatorname{div}(\nabla u)=\Delta u.
\end{equation}
此外，逐分量定义 $\Delta\boldsymbol F=(\Delta P,\Delta Q,\Delta R)$ 时，有
\[
    \operatorname{curl}(\operatorname{curl}\boldsymbol F)
    =\nabla(\operatorname{div}\boldsymbol F)-\Delta\boldsymbol F.
\]
\end{proposition}
\begin{proof}
第一式各分量均为 $u$ 的两个相等混合偏导数之差；第二式展开为六项，按混合偏导数相等两两相消；第三式由定义得到。最后一式的第一分量为
\[
    \partial_y(Q_x-P_y)-\partial_z(P_z-R_x)
    =\partial_x(P_x+Q_y+R_z)-(P_{xx}+P_{yy}+P_{zz}),
\]
其余两分量由循环换位得到。前两式也分别对应微分形式中的 $\mathrm d^2u=0$ 和 $\mathrm d^2\eta=0$。
\end{proof}

因此调和函数的梯度场既无旋又无源。但从“无源”反推某个场必为梯度场没有依据；无源只是一条标量偏微分条件，梯度场还须满足无旋及适当的全局条件。

\begin{example}
\label{ex:harmonic-conjugate-potential}
设 $U\subset\mathbb R^2$ 为单连通区域，$u\in C^2(U)$ 且 $u_{xx}+u_{yy}=0$。证明 $u_y\,\mathrm dx-u_x\,\mathrm dy$ 有原函数，且其原函数也是调和函数。
\end{example}
\begin{solve}
令 $P=u_y$、$Q=-u_x$，则 $P,Q\in C^1$，且 $Q_x-P_y=-u_{xx}-u_{yy}=0$。由定理~\ref{thm:simply-connected-curl-free}，存在 $v$ 满足 $v_x=u_y$、$v_y=-u_x$。右端为 $C^1$ 函数，故 $v\in C^2$，而
\[
    \Delta v=(u_y)_x-(u_x)_y=0.
\]
所以 $v$ 也是调和函数。沿教材的符号约定可称其为共轭调和函数；若采用复函数 $u+\mathrm iv$ 的通常柯西--黎曼约定，则相应共轭函数应取这里的 $-v$。例如 $u=x^2-y^2$ 时，上述原函数为 $v=-2xy+C$，而通常与 $u$ 配成 $u+\mathrm iv$ 的共轭函数为 $2xy+C$。
\end{solve}



\section*{习题十六}
\addcontentsline{toc}{section}{习题十六}
\fi

\selectchapter{17}
\ifcompilechapter
\chapter{含参变量积分}
\label{chap:parameter-dependent-integrals}

% 本章依据“数学分析/textbooks/数学分析（分章原文）/数学分析-17.pdf”第 1--28 页习题十七之前的正文整理。
积分中除积分变量外还可以含有参数。例如，随时间变化的向量场通过一个固定曲面的瞬时通量，是对空间变量积分后得到的时间函数；求一段时间内的总通量，还需再对时间积分。本章研究的核心是：积分得到的参数函数是否连续、可导，以及极限、微分、积分能否与原积分交换。有限区间上的联合连续性提供一致控制；积分区间无穷或被积函数有瑕点时，还必须控制截断误差。本章以黎曼积分及相应广义积分为基础，先讨论一般理论，再用它研究 $\Gamma$ 函数与 $B$ 函数。

\section{含参变量定积分}
\label{sec:parameter-dependent-definite-integrals}

设 $a<b$、$c<d$，函数 $f(x,y)$ 定义在矩形 $D=[a,b]\times[c,d]$ 上。若对每个固定的 $x\in[a,b]$，$f(x,\cdot)$ 在 $[c,d]$ 上可积，则
\[
    I(x)=\int_c^d f(x,y)\,\mathrm{d}y
\]
定义了一个以 $x$ 为自变量的函数，称为\emph{含参变量定积分}（parameter-dependent definite integral）。这里 $y$ 是积分变量，$x$ 是参变量；交换二者的角色也得到同类问题。以下端点处的连续性与导数按相应单侧理解。

\begin{theorem}[连续性与极限交换]
\label{thm:parameter-finite-integral-continuity}
若 $f\in C(D)$，则 $I\in C([a,b])$。特别地，对任意 $x_0\in[a,b]$，
\begin{equation}
    \lim_{x\to x_0}\int_c^d f(x,y)\,\mathrm{d}y
    =\int_c^d f(x_0,y)\,\mathrm{d}y
    =\int_c^d\lim_{x\to x_0}f(x,y)\,\mathrm{d}y.
    \label{eq:parameter-finite-limit-exchange}
\end{equation}
\end{theorem}

\begin{proof}
$D$ 紧且 $f$ 连续，故 $f$ 在 $D$ 上一致连续。当 $x\to x_0$ 时，$\sup_{y\in[c,d]}|f(x,y)-f(x_0,y)|\to0$，于是
\[
    |I(x)-I(x_0)|\leq(d-c)
    \sup_{y\in[c,d]}|f(x,y)-f(x_0,y)|\longrightarrow0.
\]
这里将逐点连续性加强为对所有积分变量统一有效的估计，正是能够交换极限与积分的原因。
\end{proof}

同样，$F(x,u)=\int_c^u f(x,y)\,\mathrm{d}y$ 在 $D$ 上联合连续。事实上，若 $M=\max_D|f|$，则
\[
    |F(x,u)-F(x_0,u_0)|
    \leq(d-c)\sup_{y\in[c,d]}|f(x,y)-f(x_0,y)|+M|u-u_0|.
\]
联合连续只是积分函数连续的充分条件。例如 $f(x,y)=h(x)g(y)$，若 $g$ 可积且 $\int_c^d g=0$，则无论 $h$ 是否连续，都有 $I(x)=0$。

\begin{theorem}[积分次序交换]
\label{thm:parameter-finite-integration-exchange}
若 $f\in C(D)$，则两个含参变量定积分分别是相应闭区间上的连续函数，且
\begin{equation}
    \int_a^b\left(\int_c^d f(x,y)\,\mathrm{d}y\right)\mathrm{d}x
    =\int_c^d\left(\int_a^b f(x,y)\,\mathrm{d}x\right)\mathrm{d}y.
    \label{eq:parameter-finite-integration-exchange}
\end{equation}
\end{theorem}

\begin{proof}
连续性由定理~\ref{thm:parameter-finite-integral-continuity} 得到。连续函数在矩形上重可积，两个累次积分都等于 $\iint_D f(x,y)\,\mathrm{d}x\,\mathrm{d}y$，故相等。
\end{proof}

\begin{theorem}[积分号下求导]
\label{thm:parameter-finite-differentiation}
若 $f$ 与偏导数 $f_x=\partial f/\partial x$ 在 $D$ 上连续，则 $I\in C^1([a,b])$，且
\begin{equation}
    I'(x)=\frac{\mathrm{d}}{\mathrm{d}x}\int_c^d f(x,y)\,\mathrm{d}y
    =\int_c^d f_x(x,y)\,\mathrm{d}y.
    \label{eq:parameter-finite-differentiation}
\end{equation}
\end{theorem}

\begin{proof}
令 $g(x)=\int_c^d f_x(x,y)\,\mathrm{d}y$，由连续性定理知 $g$ 连续。对固定的 $y$ 使用微积分基本定理，再交换有限区间上的积分，得
\[
    I(x)-I(a)=\int_c^d\int_a^x f_x(t,y)\,\mathrm{d}t\,\mathrm{d}y
    =\int_a^x g(t)\,\mathrm{d}t.
\]
两端对 $x$ 求导，便得 $I'=g$。
\end{proof}

\begin{theorem}[变限积分的莱布尼茨公式]
\label{thm:parameter-leibniz-variable-limits}
设 $f,f_x\in C(D)$，$\alpha,\beta:[a,b]\to[c,d]$ 可导，则
\begin{equation}
    \frac{\mathrm{d}}{\mathrm{d}x}
      \int_{\alpha(x)}^{\beta(x)}f(x,y)\,\mathrm{d}y
    =\int_{\alpha(x)}^{\beta(x)}f_x(x,y)\,\mathrm{d}y
      +f(x,\beta(x))\beta'(x)-f(x,\alpha(x))\alpha'(x).
    \label{eq:parameter-leibniz-variable-limits}
\end{equation}
积分按定向区间理解，不要求始终有 $\alpha(x)\leq\beta(x)$。
\end{theorem}

\begin{proof}
对 $F(x,u)=\int_c^u f(x,y)\,\mathrm{d}y$，前一定理与微积分基本定理分别给出
\[
    F_x(x,u)=\int_c^u f_x(x,y)\,\mathrm{d}y,
    \qquad F_u(x,u)=f(x,u).
\]
这两个偏导数联合连续，故 $F$ 可微。把原积分写成 $F(x,\beta(x))-F(x,\alpha(x))$，用复合函数求导法则即得结论。若积分限在端点取到 $c$ 或 $d$，相应增量沿矩形内取值，仍有同样的单侧微分估计。
\end{proof}

\emph[tip]{变限求导包含两类贡献：被积函数随参数的变化，以及上下限移动所增减的积分。}上下限固定时只保留积分项；被积函数不显含 $x$ 时只保留边界项。

\begin{example}
\label{ex:parameter-finite-limit-and-logarithm}
求 $\displaystyle\lim_{x\to0+}\int_0^1(1+y)^x e^{xy}\sin(x+y)\,\mathrm{d}y$，并计算
\[
    \int_0^1\frac{x^{2e-1}-x^{e-1}}{\ln x}\,\mathrm{d}x.
\]
\end{example}

\begin{solve}
第一个被积函数在 $[0,1]\times[0,1]$ 上连续，故极限等于
\[
    \int_0^1\sin y\,\mathrm{d}y=1-\cos1.
\]
第二个积分利用参数积分表示
\[
    \frac{x^{2e-1}-x^{e-1}}{\ln x}
    =\int_{e-1}^{2e-1}x^y\,\mathrm{d}y\qquad(0<x<1).
\]
因为 $y\geq e-1>0$，把 $x^y$ 在 $x=0$ 处补为零后，它在 $[0,1]\times[e-1,2e-1]$ 上联合连续；在 $x=1$ 处原商的连续延拓值为 $e$。因此交换积分次序合法，并有
\[
    \int_{e-1}^{2e-1}\int_0^1 x^y\,\mathrm{d}x\,\mathrm{d}y
    =\int_{e-1}^{2e-1}\frac{\mathrm{d}y}{1+y}
    =\ln2.
\]
参数积分把原来不易直接求原函数的商，化成了幂函数的简单积分。
\end{solve}

\begin{example}
\label{ex:parameter-nested-variable-limits}
设 $f\in C(\mathbb R^2)$，令
\[
    F(x)=\int_0^{x^2}\left(\int_0^x f(s,t)\,\mathrm{d}t\right)\mathrm{d}s.
\]
求 $F'(x)$。
\end{example}

\begin{solve}
记 $g(s,x)=\int_0^x f(s,t)\,\mathrm{d}t$，则 $g$ 联合连续，且 $g_x(s,x)=f(s,x)$ 联合连续。对任意固定 $x$，在包含所有相关积分限的有限矩形上应用式~\eqref{eq:parameter-leibniz-variable-limits}，得
\[
    F'(x)=\int_0^{x^2}f(s,x)\,\mathrm{d}s
      +2x\int_0^x f(x^2,t)\,\mathrm{d}t.
\]
这里只有 $f$ 的连续性，不需要 $f$ 本身有偏导数；所需的参数导数来自内层变上限积分。
\end{solve}

\section{含参变量广义积分}
\label{sec:parameter-dependent-improper-integrals}

把有限积分区间改为无穷区间，或允许被积函数在端点无界后，积分还包含一个截断极限。因此，有限截断积分的连续性或可导性，只解决了截断之前的问题；要把性质传给广义积分，还需控制参数变化时的尾部误差。

\subsection{含参变量无穷积分}
\label{subsec:parameter-dependent-infinite-integrals}

以下设 $E\subseteq\mathbb R$ 非空，$f$ 定义在 $E\times[c,+\infty)$ 上，且对每个 $x\in E$，$f(x,\cdot)$ 在每个有限区间 $[c,A]$ 上黎曼可积。若
\[
    I(x)=\int_c^{+\infty}f(x,y)\,\mathrm{d}y
    =\lim_{A\to+\infty}F_A(x),\qquad
    F_A(x)=\int_c^A f(x,y)\,\mathrm{d}y
\]
对每个 $x\in E$ 均存在且有限，就得到一个\emph{含参变量无穷积分}（parameter-dependent improper integral over an infinite interval）。

\begin{definition}[一致收敛与绝对一致收敛]
\label{def:parameter-infinite-uniform-convergence}
若上述积分在 $E$ 上逐点收敛，且
\begin{equation}
    \forall\varepsilon>0\ \exists A_0>c\ \forall A>A_0\ \forall x\in E,
    \quad\left|\int_A^{+\infty}f(x,y)\,\mathrm{d}y\right|<\varepsilon,
    \label{eq:parameter-infinite-uniform-tail}
\end{equation}
则称它在 $E$ 上\emph{一致收敛}（uniform convergence）。若 $\int_c^{+\infty}|f(x,y)|\,\mathrm{d}y$ 在 $E$ 上一致收敛，则称原积分在 $E$ 上\emph{绝对一致收敛}（uniform absolute convergence）。
\end{definition}

逐点收敛允许截断阈值 $A_0$ 随 $x$ 改变，一致收敛要求同一个阈值适用于所有参数。绝对一致收敛蕴含一致收敛，也蕴含每个参数处的绝对收敛；每个参数处绝对收敛本身不能保证一致收敛。与函数序列不同，截断参数 $A$ 是连续变化的实数，不能未经论证就只检验整数截断。

\begin{theorem}[一致收敛的柯西准则]
\label{thm:parameter-infinite-uniform-cauchy}
含参变量无穷积分在 $E$ 上一致收敛，当且仅当
\begin{equation}
    \forall\varepsilon>0\ \exists A_0>c\ \forall B>A>A_0\ \forall x\in E,
    \quad\left|\int_A^B f(x,y)\,\mathrm{d}y\right|<\varepsilon.
    \label{eq:parameter-infinite-uniform-cauchy}
\end{equation}
\end{theorem}

\begin{proof}
必要性由 $\int_A^B f=\int_A^{+\infty}f-\int_B^{+\infty}f$ 及三角不等式得到。反之，固定 $x$ 后，式~\eqref{eq:parameter-infinite-uniform-cauchy} 是普通无穷积分的柯西条件，故 $I(x)$ 存在。在该式中先把误差取为 $\varepsilon/2$，再令 $B\to+\infty$，可得对所有 $x$ 统一成立的尾积分估计 $|\int_A^{+\infty}f|\leq\varepsilon/2<\varepsilon$。
\end{proof}

\begin{theorem}[魏尔斯特拉斯判别法]
\label{thm:parameter-infinite-weierstrass}
若存在与 $x$ 无关的非负局部可积函数 $g$，使
\[
    |f(x,y)|\leq g(y)\quad(x\in E,\ y\geq c),\qquad
    \int_c^{+\infty}g(y)\,\mathrm{d}y<+\infty,
\]
则 $\int_c^{+\infty}f(x,y)\,\mathrm{d}y$ 在 $E$ 上绝对一致收敛。
\end{theorem}

\begin{proof}
对任意 $B>A\geq c$ 与 $x\in E$，$\int_A^B|f(x,y)|\,\mathrm{d}y\leq\int_A^B g(y)\,\mathrm{d}y$，右端满足与 $x$ 无关的柯西估计，应用前一定理即可。
\end{proof}

\begin{theorem}[狄利克雷与阿贝尔判别法]
\label{thm:parameter-infinite-dirichlet-abel}
设 $f(x,\cdot),g(x,\cdot)$ 在每个有限截断区间上可积，并且对每个固定 $x\in E$，$g(x,y)$ 关于 $y$ 单调。以下各组条件均能保证 $\int_c^{+\infty}f(x,y)g(x,y)\,\mathrm{d}y$ 在 $E$ 上一致收敛。

狄利克雷条件（Dirichlet test）：存在 $M>0$，使 $|\int_c^A f(x,y)\,\mathrm{d}y|\leq M$ 对所有 $A>c$、$x\in E$ 成立，且 $g(x,y)$ 在 $y\to+\infty$ 时关于 $x\in E$ 一致趋于零。阿贝尔条件（Abel test）：$\int_c^{+\infty}f(x,y)\,\mathrm{d}y$ 在 $E$ 上一致收敛，且存在 $L>0$，使 $|g(x,y)|\leq L$ 对所有 $x\in E$、$y\geq c$ 成立。
\end{theorem}

\begin{proof}
由定积分第二中值定理，对任意 $B>A\geq c$，存在 $\xi\in[A,B]$，使
\[
    \int_A^B fg\,\mathrm{d}y
    =g(x,A)\int_A^\xi f\,\mathrm{d}y
      +g(x,B)\int_\xi^B f\,\mathrm{d}y.
\]
在狄利克雷条件下，右侧两个 $f$ 的积分绝对值均不超过 $2M$，故原积分绝对值不超过 $2M(|g(x,A)|+|g(x,B)|)$，并一致趋于零。在阿贝尔条件下，一致柯西准则保证右侧两个积分对足够大的 $A$ 均小于任意指定的 $\delta>0$，故乘积积分绝对值不超过 $2L\delta$。两种情形都满足式~\eqref{eq:parameter-infinite-uniform-cauchy}。
\end{proof}

\begin{example}
\label{ex:parameter-exponential-tail-domain}
讨论 $\displaystyle\int_1^{+\infty}y^x e^{-y}\,\mathrm{d}y$ 在 $[0,M]$（$0<M<+\infty$）及 $[0,+\infty)$ 上的一致收敛性。
\end{example}

\begin{solve}
在 $[0,M]$ 上，因 $y\geq1$，有 $0\leq y^x e^{-y}\leq y^M e^{-y}$。后者可积，因为充分大的 $y$ 满足 $y^M\leq e^{y/2}$，故由魏尔斯特拉斯判别法绝对一致收敛。在 $[0,+\infty)$ 上，任给 $A>1$，都有
\[
    \int_A^{A+1}y^x e^{-y}\,\mathrm{d}y
    \geq A^x e^{-(A+1)}.
\]
选取足够大的 $x$ 可使右端至少为 $1$，所以一致柯西条件失败。每个固定 $x$ 处的绝对收敛，并不能抵消参数无界增长造成的尾部失控。
\end{solve}

\begin{example}
\label{ex:parameter-sine-tail-uniformity}
证明 $\displaystyle\int_1^{+\infty}\frac{\sin(xy)}{y}\,\mathrm{d}y$ 在每个 $[a_0,+\infty)$（$a_0>0$）上一致收敛，但在 $(0,+\infty)$ 上不一致收敛。
\end{example}

\begin{solve}
当 $x\geq a_0$ 时，$|\int_1^A\sin(xy)\,\mathrm{d}y|\leq2/a_0$，而 $1/y$ 单调趋于零且与参数无关，故一致狄利克雷判别法适用。为证后一结论，取 $A_n=n$、$B_n=2n$、$x_n=1/n$，换元 $t=y/n$ 得
\[
    \int_{A_n}^{B_n}\frac{\sin(x_ny)}{y}\,\mathrm{d}y
    =\int_1^2\frac{\sin t}{t}\,\mathrm{d}t>0.
\]
右端是不随 $n$ 改变的正常数，故尾积分不能一致趋于零。参数靠近零时，振荡抵消所需的积分长度不断增大。
\end{solve}

\begin{example}
\label{ex:parameter-abel-damping-equivalence}
设 $h\in C([0,+\infty))$。证明 $\int_0^{+\infty}e^{-\alpha t}h(t)\,\mathrm{d}t$ 对 $\alpha\in(0,+\infty)$ 一致收敛，当且仅当 $\int_0^{+\infty}h(t)\,\mathrm{d}t$ 收敛。
\end{example}

\begin{solve}
若无阻尼积分收敛，把 $h(t)$ 看成与 $\alpha$ 无关的函数，则它的积分关于 $\alpha$ 一致收敛；$e^{-\alpha t}$ 关于 $t$ 单调且统一有界于 $1$，阿贝尔判别法给出充分性。反之，若带阻尼积分一致收敛，则对任意 $\varepsilon>0$，存在 $A_0$，使所有 $B>A>A_0$ 及 $\alpha>0$ 都满足
\[
    \left|\int_A^B e^{-\alpha t}h(t)\,\mathrm{d}t\right|<\varepsilon/2.
\]
先固定有限的 $A,B$，再由有限区间积分的连续性令 $\alpha\to0+$，得到 $|\int_A^B h(t)\,\mathrm{d}t|\leq\varepsilon/2<\varepsilon$。普通无穷积分的柯西准则给出必要性。
\end{solve}

\subsection{含参变量无穷积分的性质}
\label{subsec:properties-of-parameter-dependent-infinite-integrals}

\begin{theorem}[连续截断与序列截断的联系]
\label{thm:parameter-infinite-sequential-criterion}
在局部可积的基本假设下，$\int_c^{+\infty}f(x,y)\,\mathrm{d}y$ 在 $E$ 上一致收敛，当且仅当对于每个严格递增且趋于 $+\infty$ 的实数序列 $\{t_n\}$（$t_1>c$），函数序列
\[
    F_{t_n}(x)=\int_c^{t_n}f(x,y)\,\mathrm{d}y
\]
在 $E$ 上一致收敛。此时各序列的极限都是同一个积分函数 $I(x)$。
\end{theorem}

\begin{proof}
必要性由定义直接得到。若积分不一致收敛，则一致柯西条件失败，存在 $\varepsilon_0>0$，可以逐次选择
\[
    c<A_1<B_1<A_2<B_2<\cdots,\qquad A_n>n,
\]
以及 $x_n\in E$，使 $|\int_{A_n}^{B_n}f(x_n,y)\,\mathrm{d}y|\geq\varepsilon_0$。将所有 $A_n,B_n$ 依次排成 $\{t_n\}$，所得截断函数序列不满足一致柯西准则，与假设矛盾。
\end{proof}

\emph[caution]{这里要求每个趋于无穷的递增截断序列都适用。}例如与参数无关的 $f(y)=\sin(2\pi y)$ 在整数截断处的积分恒为零，但 $\int_0^A f(y)\,\mathrm{d}y=(1-\cos(2\pi A))/(2\pi)$ 在 $A\to+\infty$ 时无极限。单独检验一列截断值会漏掉相邻截断之间的变化。

\begin{theorem}[连续性与有限区间上的积分交换]
\label{thm:parameter-infinite-continuity-integration}
设 $E$ 为区间，$f\in C(E\times[c,+\infty))$。若 $I(x)=\int_c^{+\infty}f(x,y)\,\mathrm{d}y$ 在 $E$ 上一致收敛，则 $I$ 在 $E$ 上连续。若 $E=[a,b]$，还有
\begin{equation}
    \int_a^b\left(\int_c^{+\infty}f(x,y)\,\mathrm{d}y\right)\mathrm{d}x
    =\int_c^{+\infty}\left(\int_a^b f(x,y)\,\mathrm{d}x\right)\mathrm{d}y,
    \label{eq:parameter-infinite-integration-exchange}
\end{equation}
其中右侧无穷积分也收敛。
\end{theorem}

\begin{proof}
每个有限截断积分 $F_A(x)$ 在 $E$ 上连续：对任意参数点，只需在其有限闭邻域与 $[c,A]$ 的乘积上应用定积分连续性定理。由一致极限保持连续性，$I$ 连续。在 $E=[a,b]$ 时，
\[
    \left|\int_a^b I(x)\,\mathrm{d}x-\int_a^bF_A(x)\,\mathrm{d}x\right|
    \leq(b-a)\sup_{x\in[a,b]}|I(x)-F_A(x)|\longrightarrow0.
\]
对 $F_A$ 的积分使用有限矩形上的次序交换，再令实数 $A\to+\infty$，便得到右侧广义积分的存在性及公式。
\end{proof}

若只要求 $I$ 在开区间内部连续，只需对每个闭子区间证明一致收敛，即\emph{内闭一致收敛}（locally uniform convergence）。但式~\eqref{eq:parameter-infinite-integration-exchange} 在整个 $[a,b]$ 上的结论，不能仅凭 $(a,b)$ 内部的局部一致收敛推出。

\begin{theorem}[无穷积分的积分号下求导]
\label{thm:parameter-infinite-differentiation}
设 $E$ 为非退化区间，$f,f_x\in C(E\times[c,+\infty))$。假设某个 $x_0\in E$ 处的积分 $\int_c^{+\infty}f(x_0,y)\,\mathrm{d}y$ 收敛，并且 $\int_c^{+\infty}f_x(x,y)\,\mathrm{d}y$ 在 $E$ 的每个闭子区间上一致收敛，则原积分在 $E$ 上处处收敛且内闭一致收敛，积分函数 $I$ 连续可导，并有
\begin{equation}
    I'(x)=\int_c^{+\infty}f_x(x,y)\,\mathrm{d}y.
    \label{eq:parameter-infinite-differentiation}
\end{equation}
若 $E$ 本身有界，且导数积分在整个 $E$ 上一致收敛，则原积分在整个 $E$ 上也一致收敛。
\end{theorem}

\begin{proof}
取任意闭子区间 $K\subseteq E$，必要时把它扩大为仍在 $E$ 内且包含 $x_0$ 的有限闭区间 $K_0$。由有限截断积分的求导公式与微积分基本定理，对 $B>A>c$ 及 $x\in K_0$，有
\[
    \int_A^B f(x,y)\,\mathrm{d}y
    =\int_A^B f(x_0,y)\,\mathrm{d}y
      +\int_{x_0}^x\left(\int_A^B f_x(t,y)\,\mathrm{d}y\right)\mathrm{d}t.
\]
因此其绝对值不超过
\[
    \left|\int_A^B f(x_0,y)\,\mathrm{d}y\right|
    +|K_0|\sup_{t\in K_0}
       \left|\int_A^B f_x(t,y)\,\mathrm{d}y\right|,
\]
其中 $|K_0|$ 表示区间长度。两项均趋于零，故原积分在 $K_0$ 上一致收敛。再任取 $t_n\uparrow+\infty$，截断函数 $F_{t_n}$ 在 $x_0$ 处收敛，其导函数在 $K_0$ 上一致收敛。由\crosschapterref{10}{定理~\ref{thm:differentiation-of-function-sequence}}{导数与极限的交换定理}，其极限可导且导数等于导函数的极限，即式~\eqref{eq:parameter-infinite-differentiation}。导数的连续性由前一连续性定理得到。若 $E$ 有界且导数积分在 $E$ 上一致收敛，上述估计中可用 $E$ 的有限长度统一控制 $|x-x_0|$，从而原积分在整个 $E$ 上一致收敛。
\end{proof}

\begin{remark}
\label{rem:parameter-derivative-unbounded-domain}
求导公式是局部结论，不能在无界参数区间上无条件加强为原积分整体一致收敛。例如 $f(x,y)=xe^{-y}$，$x\geq0$、$y\geq0$，则 $\int_0^{+\infty}f_x(x,y)\,\mathrm{d}y=1$ 关于 $x$ 一致收敛，$x=0$ 处原积分也收敛；但原积分的尾部为 $xe^{-A}$，在 $[0,+\infty)$ 上没有统一趋于零。此时 $I(x)=x$ 仍然可导，求导公式完全成立。
\end{remark}

\begin{proposition}[非负积分的狄尼定理]
\label{prop:parameter-integral-dini}
设 $f\in C([a,b]\times[c,+\infty))$，且 $f\geq0$。若 $I(x)=\int_c^{+\infty}f(x,y)\,\mathrm{d}y$ 对每个 $x\in[a,b]$ 收敛并在 $[a,b]$ 上连续，则该积分在 $[a,b]$ 上一致收敛。$f\leq0$ 时结论相同。
\end{proposition}

\begin{proof}
对任意 $t_n\uparrow+\infty$，连续函数 $F_{t_n}$ 单调增加并逐点趋于连续函数 $I$。由函数序列的狄尼定理\crosschapterref{10}{，即定理~\ref{thm:dini-uniform-convergence}}{}，它一致收敛。也可固定一列 $\{t_n\}$，由 $A\geq t_n$ 时 $0\leq I(x)-F_A(x)\leq I(x)-F_{t_n}(x)$，直接得到所有实数截断的一致估计。非正情形对 $-f$ 应用同一结论。
\end{proof}

\begin{example}
\label{ex:parameter-dirichlet-integral-evaluation}
利用含参变量积分计算狄利克雷积分 $\displaystyle\int_0^{+\infty}\frac{\sin t}{t}\,\mathrm{d}t$。
\end{example}

\begin{solve}
先确认积分存在：$\sin t/t$ 在零点可连续延拓为 $1$，而无穷尾部由普通狄利克雷判别法收敛。引入阻尼参数
\[
    I(\alpha)=\int_0^{+\infty}e^{-\alpha t}\frac{\sin t}{t}\,\mathrm{d}t,
    \qquad\alpha\geq0.
\]
无阻尼积分收敛，且 $e^{-\alpha t}$ 对 $t$ 单调、关于 $\alpha\geq0$ 统一有界于 $1$，故一致阿贝尔判别法保证 $I$ 的积分在 $[0,+\infty)$ 上一致收敛。被积函数在 $t=0$ 处按连续延拓理解，因此 $I$ 在 $\alpha=0$ 处右连续。

对任意 $\delta>0$，当 $\alpha\geq\delta$ 时，参数偏导数为 $-e^{-\alpha t}\sin t$，绝对值不超过可积函数 $e^{-\delta t}$。故对所有 $\alpha>0$ 可在积分号下求导，得到
\[
    I'(\alpha)=-\int_0^{+\infty}e^{-\alpha t}\sin t\,\mathrm{d}t
    =-\frac1{1+\alpha^2}.
\]
最后一步可用原函数 $-e^{-\alpha t}(\alpha\sin t+\cos t)/(1+\alpha^2)$ 核对。于是 $I(\alpha)=C-\arctan\alpha$。又由 $|\sin t|\leq t$，
\[
    |I(\alpha)|\leq\int_0^{+\infty}e^{-\alpha t}\,\mathrm{d}t
    =\frac1\alpha\longrightarrow0\qquad(\alpha\to+\infty),
\]
从而 $C=\pi/2$。利用零点处已经证明的右连续性，得
\begin{equation}
    \int_0^{+\infty}\frac{\sin t}{t}\,\mathrm{d}t
    =\lim_{\alpha\to0+}I(\alpha)=\frac\pi2.
    \label{eq:parameter-dirichlet-integral-value}
\end{equation}
这里的极限交换由一致收敛支持；仅算出 $\alpha>0$ 时的积分表达式，还不足以识别无阻尼积分。
\end{solve}

\subsection{含参变量瑕积分}
\label{subsec:parameter-dependent-singular-integrals}

设 $f\in C(E\times(c,d])$，其中 $c<d$ 为有限数，允许 $y=c$ 为瑕点。若每个 $x\in E$ 处都有有限极限
\[
    I(x)=\int_c^d f(x,y)\,\mathrm{d}y
    :=\lim_{\eta\to0+}\int_{c+\eta}^d f(x,y)\,\mathrm{d}y,
\]
则称其为\emph{含参变量瑕积分}（parameter-dependent improper integral with an endpoint singularity）。其一致收敛要求
\begin{equation}
    \forall\varepsilon>0\ \exists\eta_0\in(0,d-c)
    \forall\eta\in(0,\eta_0)\ \forall x\in E,
    \quad\left|\int_c^{c+\eta}f(x,y)\,\mathrm{d}y\right|<\varepsilon.
    \label{eq:parameter-singular-uniform-tail}
\end{equation}
等价的一致柯西条件为：对充分小的 $0<\eta_1<\eta_2$，$|\int_{c+\eta_1}^{c+\eta_2}f(x,y)\,\mathrm{d}y|$ 对全部 $x\in E$ 统一地小。

令 $t=1/(y-c)$，则 $y=c+1/t$，$\mathrm{d}y=-t^{-2}\mathrm{d}t$，从而
\begin{equation}
    \int_c^d f(x,y)\,\mathrm{d}y
    =\int_{1/(d-c)}^{+\infty}
        \frac{f(x,c+1/t)}{t^2}\,\mathrm{d}t.
    \label{eq:parameter-singular-to-infinite}
\end{equation}
两个截断误差完全对应，因此前面的一致柯西准则、魏尔斯特拉斯判别法、连续性、积分交换及积分号下求导定理均可转移过来。例如，若 $|f(x,y)|\leq g(y)$ 且非负函数 $g$ 在 $(c,d]$ 上瑕可积，就有绝对一致收敛；若 $f,f_x$ 联合连续、某个参数处原积分收敛，且 $\int_c^d f_x(x,y)\,\mathrm{d}y$ 内闭一致收敛，则 $I'(x)=\int_c^d f_x(x,y)\,\mathrm{d}y$。若两端都有瑕点，应固定一个内部点分别讨论两端，不能用同一个趋近参数使两端的发散部分互相抵消。

\begin{example}
\label{ex:parameter-power-singularity-uniformity}
讨论 $\displaystyle\int_0^1y^{s-1}\,\mathrm{d}y$ 的参数范围及一致收敛性。
\end{example}

\begin{solve}
当且仅当 $s>0$ 时积分收敛，值为 $1/s$。对任意 $\delta>0$、$0<\eta<1$ 及 $s\geq\delta$，
\[
    \int_0^\eta y^{s-1}\,\mathrm{d}y
    =\frac{\eta^s}{s}\leq\frac{\eta^\delta}{\delta}\longrightarrow0.
\]
因此在每个 $[\delta,+\infty)$ 上一致收敛。但在 $(0,+\infty)$ 上，固定任何 $\eta\in(0,1)$，$\eta^s/s\to+\infty$（$s\to0+$），故不一致收敛。参数逼近定义域边界时，可积奇性的强弱可能失去统一控制；这也解释了后面讨论 $\Gamma$ 与 $B$ 函数时为何把参数限制在正半轴的闭子区间上。
\end{solve}

含参变量积分还可推广到向量参数以及重积分、曲线积分和曲面积分。例如，设 $S$ 是固定的紧致、分片光滑、有指定侧的曲面，面积为 $|S|$，$\boldsymbol n$ 为其指定单位法向量；若 $\boldsymbol F(\boldsymbol r,t)$ 在 $S\times[a,b]$ 上连续，则瞬时有向通量
\[
    q(t)=\iint_S\boldsymbol F(\boldsymbol r,t)\cdot\boldsymbol n(\boldsymbol r)\,\mathrm{d}S
\]
连续，因为
\[
    |q(t)-q(t_0)|\leq |S|\sup_{\boldsymbol r\in S}
       |\boldsymbol F(\boldsymbol r,t)-\boldsymbol F(\boldsymbol r,t_0)|\longrightarrow0.
\]
这里即使分片交界处的法向量不连续，估计仍成立，因为法向量长度为 $1$，交界线不影响曲面积分。在流体模型中，若 $\boldsymbol F$ 表示体积通量密度，则 $[t_0,t_1]$ 内的净通过量为 $Q=\int_{t_0}^{t_1}q(t)\,\mathrm{d}t$；稳定场的 $q$ 与时间无关，因而 $Q=(t_1-t_0)q$。这一例子回到本章的动机，也说明关键在于积分区域上能够统一控制参数变化。

\section{\texorpdfstring{$\Gamma$}{Gamma} 函数与 \texorpdfstring{$B$}{B} 函数}
\label{sec:gamma-and-beta-functions}

前两节建立了由积分定义函数并研究其参数性质的方法。本节把它用于两类不能一般地用初等函数表示的积分：$\Gamma$ 函数把整数阶乘延伸到正实数参数，$B$ 函数统一处理两个端点的幂函数因子。先核实各自的收敛范围，再建立递推、光滑性与二者之间的关系，便能把一批定积分和几何量的计算归约到这些函数的特殊值。

\subsection{\texorpdfstring{$\Gamma$}{Gamma} 函数}
\label{subsec:gamma-function}

\begin{definition}[伽马函数]
\label{def:gamma-euler-integral}
对 $s>0$，定义\emph{伽马函数}（gamma function）
\begin{equation}
    \Gamma(s)=\int_0^{+\infty}t^{s-1}e^{-t}\,\mathrm{d}t.
    \label{eq:gamma-euler-integral}
\end{equation}
本章只讨论实参数的积分定义。
\end{definition}

在 $(0,1]$ 上，$e^{-1}t^{s-1}\leq t^{s-1}e^{-t}\leq t^{s-1}$，故零点附近当且仅当 $s>0$ 时可积；在 $[1,+\infty)$ 上，对任意固定实数 $s$，幂函数增长都可被指数衰减控制，因而无穷尾部收敛。因此式~\eqref{eq:gamma-euler-integral} 的实数定义域恰为 $(0,+\infty)$，且 $\Gamma(s)>0$。当 $0<s<1$ 时，它同时包含零点瑕积分与无穷积分，应在 $1$ 处分开理解。

\begin{proposition}[递推公式与阶乘]
\label{prop:gamma-recurrence-factorials}
对 $s>0$，有
\begin{equation}
    \Gamma(s+1)=s\Gamma(s),\qquad
    \Gamma(1)=1,\qquad \Gamma(n+1)=n!\quad(n=0,1,2,\ldots).
    \label{eq:gamma-recurrence-factorials}
\end{equation}
此外，换元 $t=u^2$ 得到
\begin{equation}
    \Gamma(s)=2\int_0^{+\infty}u^{2s-1}e^{-u^2}\,\mathrm{d}u,
    \qquad s>0.
    \label{eq:gamma-gaussian-form}
\end{equation}
\end{proposition}

\begin{proof}
先在 $[\eta,A]$ 上分部积分，再令 $\eta\to0+$、$A\to+\infty$。因为 $s>0$ 时 $t^se^{-t}$ 在两端均趋于零，故
\[
    \Gamma(s+1)=\bigl[-t^se^{-t}\bigr]_0^{+\infty}
       +s\int_0^{+\infty}t^{s-1}e^{-t}\,\mathrm{d}t
    =s\Gamma(s).
\]
由 $\Gamma(1)=\int_0^{+\infty}e^{-t}\,\mathrm{d}t=1$，归纳即得阶乘关系。平方换元也先在有限正常区间上进行，再取端点极限，得到式~\eqref{eq:gamma-gaussian-form}。
\end{proof}

\begin{theorem}[光滑性与各阶导数]
\label{thm:gamma-smoothness-derivatives}
$\Gamma\in C^\infty((0,+\infty))$，且对每个非负整数 $k$，
\begin{equation}
    \Gamma^{(k)}(s)=\int_0^{+\infty}
       t^{s-1}(\ln t)^k e^{-t}\,\mathrm{d}t,\qquad s>0.
    \label{eq:gamma-derivative-integrals}
\end{equation}
这些积分在每个闭区间 $[\delta,M]\subset(0,+\infty)$ 上绝对一致收敛。
\end{theorem}

\begin{proof}
固定 $0<\delta<M$ 及 $k\geq0$。对 $s\in[\delta,M]$，有
\[
    |t^{s-1}(\ln t)^k e^{-t}|
    \leq\begin{cases}
        t^{\delta-1}|\ln t|^k,&0<t\leq1,\\
        t^{M-1}(\ln t)^k e^{-t},&t\geq1.
    \end{cases}
\]
第一个控制函数经 $u=-\ln t$ 换元，其积分为 $\int_0^{+\infty}e^{-\delta u}u^k\,\mathrm{d}u<+\infty$；第二个控制函数在无穷远处可由常数倍的 $e^{-t/2}$ 控制，因而也可积。魏尔斯特拉斯判别法分别控制两端，保证每阶参数偏导数的积分绝对一致收敛。反复应用含参变量积分的求导定理及连续性定理，即得所有阶导数及其连续性。由于任意 $s>0$ 都包含在某个这样的闭区间内部，结论在整个正半轴成立。
\end{proof}

\begin{proposition}[严格凸性与严格对数凸性]
\label{prop:gamma-strict-log-convexity}
$\Gamma(s)$ 与 $\ln\Gamma(s)$ 都是 $(0,+\infty)$ 上的严格凸函数。
\end{proposition}

\begin{proof}
由式~\eqref{eq:gamma-derivative-integrals}，
\[
    \Gamma''(s)=\int_0^{+\infty}t^{s-1}(\ln t)^2e^{-t}\,\mathrm{d}t>0,
\]
因为被积函数非负，且在不包含 $1$ 的任一非退化正区间上为正。又因 $\Gamma>0$，
\[
    (\ln\Gamma(s))''
    =\frac{\Gamma(s)\Gamma''(s)-(\Gamma'(s))^2}{\Gamma(s)^2}.
\]
对 $u(t)=t^{(s-1)/2}e^{-t/2}$ 与 $v(t)=u(t)\ln t$ 应用积分形式的柯西--施瓦茨不等式，得 $(\Gamma'(s))^2\leq\Gamma(s)\Gamma''(s)$。为核查严格性，令 $m=\Gamma'(s)/\Gamma(s)$，则
\[
    \Gamma''(s)-\frac{(\Gamma'(s))^2}{\Gamma(s)}
    =\int_0^{+\infty}t^{s-1}e^{-t}(\ln t-m)^2\,\mathrm{d}t>0.
\]
右端的连续非负被积函数只在 $t=e^m$ 处为零，在其他点附近均给出正积分，故 $(\ln\Gamma)''>0$。
\end{proof}

\subsection{\texorpdfstring{$B$}{B} 函数}
\label{subsec:beta-function}

\begin{definition}[贝塔函数]
\label{def:beta-euler-integral}
对 $p,q>0$，定义\emph{贝塔函数}（beta function）
\begin{equation}
    B(p,q)=\int_0^1 t^{p-1}(1-t)^{q-1}\,\mathrm{d}t.
    \label{eq:beta-euler-integral}
\end{equation}
它是含两个参数的广义积分。
\end{definition}

在零点附近，$(1-t)^{q-1}$ 是有正下界的连续函数，故可积性要求 $p>0$；在 $1$ 附近同理要求 $q>0$。所以该积分的实数定义域恰为 $(0,+\infty)^2$，且在此域内恒正。固定任意内部分点，例如 $1/2$，分别取两个端点的瑕积分极限，便得到严格的定义；当某端点没有奇性时，该部分就是普通定积分。

\begin{proposition}[对称性与递推公式]
\label{prop:beta-symmetry-recurrence}
对 $p,q>0$，有
\begin{align}
    B(p,q)&=B(q,p),\label{eq:beta-symmetry}\\
    B(p+1,q)&=\frac{p}{p+q}B(p,q),\qquad
    B(p,q+1)=\frac{q}{p+q}B(p,q).
    \label{eq:beta-forward-recurrence}
\end{align}
特别地，
\begin{align}
    B(p,q)&=\frac{p-1}{p+q-1}B(p-1,q)
        &&(p>1,\ q>0),\label{eq:beta-backward-recurrence}\\
    B(p,q)&=\frac{(p-1)(q-1)}{(p+q-1)(p+q-2)}B(p-1,q-1)
        &&(p>1,\ q>1).\label{eq:beta-double-recurrence}
\end{align}
\end{proposition}

\begin{proof}
在定义式中令 $t=1-u$，即得对称性。由 $1=t+(1-t)$，有
\[
    B(p,q)=B(p+1,q)+B(p,q+1).
\]
另一方面，$t^p(1-t)^q$ 在两端均趋于零，其导数的两个分项在 $p,q>0$ 时分别可积。先在 $[\eta,1-\eta]$ 上积分导数再令 $\eta\to0+$，得到
\[
    0=pB(p,q+1)-qB(p+1,q).
\]
联立即得向前递推式；对参数平移并重复使用，即得后两式。向后递推时必须保留 $p>1$、$q>1$ 等条件，保证右侧的积分仍在定义域内。
\end{proof}

\begin{proposition}[三角形式与无穷积分形式]
\label{prop:beta-alternative-integrals}
对 $p,q>0$，有
\begin{align}
    B(p,q)&=2\int_0^{\pi/2}
       (\cos\theta)^{2p-1}(\sin\theta)^{2q-1}\,\mathrm{d}\theta,
       \label{eq:beta-trigonometric-form}\\
    B(p,q)&=\int_0^{+\infty}
       \frac{u^{p-1}}{(1+u)^{p+q}}\,\mathrm{d}u.
       \label{eq:beta-infinite-form}
\end{align}
\end{proposition}

\begin{proof}
第一式在定义式中取 $t=\cos^2\theta$，注意积分方向反转；第二式取 $t=u/(1+u)$，此时 $\mathrm{d}t=(1+u)^{-2}\mathrm{d}u$。两次换元先在避开端点的闭区间上进行，再分别取端点极限。由对称性，式~\eqref{eq:beta-trigonometric-form} 中的正弦、余弦幂次也可以互换。
\end{proof}

\subsection{\texorpdfstring{$\Gamma$}{Gamma} 函数与 \texorpdfstring{$B$}{B} 函数的关系}
\label{subsec:relation-between-gamma-and-beta-functions}

\begin{theorem}[贝塔--伽马关系]
\label{thm:beta-gamma-relation}
对 $p,q>0$，有
\begin{equation}
    B(p,q)=\frac{\Gamma(p)\Gamma(q)}{\Gamma(p+q)}.
    \label{eq:beta-gamma-relation}
\end{equation}
从而 $B$ 在 $(0,+\infty)^2$ 内具有任意阶连续偏导数。
\end{theorem}

\begin{proof}
由式~\eqref{eq:gamma-gaussian-form}，两个一维正积分的乘积可以写成第一象限上的广义二重积分
\[
    \Gamma(p)\Gamma(q)
    =4\iint_{x>0,\,y>0}x^{2p-1}y^{2q-1}e^{-(x^2+y^2)}
       \,\mathrm{d}x\,\mathrm{d}y.
\]
为说明换元的合法性，先在矩形 $[\varepsilon,R]^2$ 上积分，再令 $\varepsilon\downarrow0$、$R\uparrow+\infty$，矩形积分等于两个截断一维积分的乘积，故极限存在。被积函数非负，也可改用极坐标截断区域 $\delta\leq r\leq L$、$\eta\leq\theta\leq\pi/2-\eta$：每个这样的区域包含于某个矩形截断区域，每个矩形截断区域也包含于某个这样的区域，所以两族积分的上确界相同。这为改换穷竭区域以及随后的极坐标变换提供依据，即使坐标轴附近有瑕点也成立。

令 $x=r\cos\theta$、$y=r\sin\theta$，雅可比因子为 $r$，于是
\begin{align*}
    \Gamma(p)\Gamma(q)
    &=4\left(\int_0^{+\infty}r^{2(p+q)-1}e^{-r^2}\,\mathrm{d}r\right)
       \left(\int_0^{\pi/2}\cos^{2p-1}\theta\sin^{2q-1}\theta
          \,\mathrm{d}\theta\right)\\
    &=\Gamma(p+q)B(p,q).
\end{align*}
两个因子都收敛，且 $\Gamma(p+q)>0$，故可相除。最后，$\Gamma$ 光滑且在正半轴不为零，关系式便给出 $B$ 的光滑性。
\end{proof}

\begin{theorem}[余元公式]
\label{thm:gamma-reflection-formula}
若 $0<p<1$，则
\begin{equation}
    \Gamma(p)\Gamma(1-p)=B(p,1-p)=\frac{\pi}{\sin(\pi p)}.
    \label{eq:gamma-reflection-formula}
\end{equation}
该恒等式也称\emph{反射公式}（reflection formula）；这里的参数范围保证两边的 $\Gamma$ 函数都由本章的正参数积分定义。
\end{theorem}

\begin{proof}
贝塔--伽马关系与 $\Gamma(1)=1$ 给出第一个等号。由式~\eqref{eq:beta-infinite-form}，在 $[1,+\infty)$ 部分取 $u=1/t$，得
\[
    B(p,1-p)=\int_0^{+\infty}\frac{t^{p-1}}{1+t}\,\mathrm{d}t
    =\int_0^1\frac{t^{p-1}+t^{-p}}{1+t}\,\mathrm{d}t.
\]
使用有限几何恒等式
\[
    \frac1{1+t}=\sum_{k=0}^N(-1)^kt^k
       +\frac{(-1)^{N+1}t^{N+1}}{1+t}.
\]
余项乘以 $t^{p-1}+t^{-p}$ 后的积分绝对值不超过
$1/(N+p+1)+1/(N+2-p)$，故趋于零，无须假定几何级数在 $[0,1]$ 上一致收敛。逐项计算有限和并在极限中合并相邻指标，得到
\begin{equation}
    B(p,1-p)
    =\lim_{N\to\infty}\sum_{k=0}^N(-1)^k
       \left(\frac1{k+p}+\frac1{k+1-p}\right)
    =\frac1p+\sum_{k=1}^{\infty}(-1)^k\frac{2p}{p^2-k^2}.
    \label{eq:beta-reflection-series}
\end{equation}
合并时只差一个趋于零的末项；合并后的级数绝对收敛。

下面按教材思路用 $\cos(px)$ 的傅里叶系数计算这个级数，并补出在 $x=0$ 处取值的依据。直接积分得
\[
    \frac1\pi\int_0^\pi\cos(pt)\,\mathrm{d}t
       =\frac{\sin(\pi p)}{\pi p},\qquad
    \frac2\pi\int_0^\pi\cos(pt)\cos(kt)\,\mathrm{d}t
       =\frac{\sin(\pi p)}\pi\frac{2p(-1)^k}{p^2-k^2}.
\]
因此对应傅里叶部分和在零点的值为
\[
    S_N=\frac{\sin(\pi p)}\pi
       \left(\frac1p+\sum_{k=1}^N\frac{2p(-1)^k}{p^2-k^2}\right).
\]
利用有限三角恒等式 $1+2\sum_{k=1}^N\cos(kt)=\sin((N+1/2)t)/\sin(t/2)$，以及常数函数的相应积分为 $1$，有
\[
    S_N-1=\frac1\pi\int_0^\pi
      \frac{\cos(pt)-1}{\sin(t/2)}\sin((N+1/2)t)\,\mathrm{d}t.
\]
函数 $h(t)=(\cos(pt)-1)/\sin(t/2)$ 在 $t=0$ 处补为零后属于 $C^1([0,\pi])$。对右侧分部积分，其绝对值不超过
\[
    \frac{|h(0)|+|h(\pi)|+\int_0^\pi|h'(t)|\,\mathrm{d}t}
         {\pi(N+1/2)}\longrightarrow0.
\]
故 $S_N\to1$。结合式~\eqref{eq:beta-reflection-series} 即得 $B(p,1-p)=\pi/\sin(\pi p)$。
\end{proof}

取 $p=1/2$，由 $\Gamma(1/2)>0$ 得
\begin{equation}
    \Gamma\left(\frac12\right)=\sqrt\pi,\qquad
    \Gamma\left(n+\frac12\right)
       =\frac{(2n)!}{4^n n!}\sqrt\pi\quad(n=0,1,2,\ldots).
    \label{eq:gamma-half-integer-values}
\end{equation}
其中第二式由递推公式得到。再由式~\eqref{eq:gamma-gaussian-form}，得到高斯积分
\begin{equation}
    \int_0^{+\infty}e^{-t^2}\,\mathrm{d}t=\frac{\sqrt\pi}{2},
    \qquad\int_{-\infty}^{+\infty}e^{-t^2}\,\mathrm{d}t=\sqrt\pi.
    \label{eq:gamma-gaussian-integral}
\end{equation}
递推公式将正参数处的值归约到 $(0,1]$ 上的值，余元公式又联系 $p$ 与 $1-p$；这些公式的使用都应遵守各自的参数范围。

\begin{example}
\label{ex:beta-sine-power-integrals}
计算 $\displaystyle I_n=\int_0^{\pi/2}\sin^n x\,\mathrm{d}x$（$n=0,1,2,\ldots$）。
\end{example}

\begin{solve}
由贝塔函数的三角形式及其与伽马函数的关系，
\begin{equation}
    I_n=\frac12B\left(\frac12,\frac{n+1}{2}\right)
       =\frac{\sqrt\pi\,\Gamma((n+1)/2)}{2\Gamma((n+2)/2)}
       =\begin{cases}
           \dfrac{(n-1)!!}{n!!}\dfrac\pi2,&n\text{ 为偶数},\\
           \dfrac{(n-1)!!}{n!!},&n\text{ 为奇数}.
         \end{cases}
    \label{eq:beta-sine-power-integrals}
\end{equation}
这里 $n!!=n(n-2)\cdots$，约定 $0!!=(-1)!!=1$；故 $n=0$ 时仍给出 $I_0=\pi/2$。换元 $x\mapsto\pi/2-x$ 可知余弦同次幂的积分相等。中间的伽马函数表达式还对任意实数 $n>-1$ 成立，此范围正是端点处可积的条件。
\end{solve}

\begin{example}
\label{ex:gamma-unit-ball-volume}
求 $\mathbb R^n$ 中单位球体 $D_n=\{(x_1,\ldots,x_n):x_1^2+\cdots+x_n^2\leq1\}$ 的体积 $V_n$。
\end{example}

\begin{solve}
先用截面递推使计算中的维数关系清楚。$V_1=2$；$n\geq2$ 时，在 $x_n=t$ 处的截面是半径为 $\sqrt{1-t^2}$ 的 $(n-1)$ 维球体，其体积由缩放变换等于 $V_{n-1}(1-t^2)^{(n-1)/2}$。球面作为有限个局部光滑图形之并具有零 $n$ 维体积，因而球体的示性函数可积，截面积分公式适用。由 $u=t^2$ 换元，
\begin{align*}
    V_n&=V_{n-1}\int_{-1}^1(1-t^2)^{(n-1)/2}\,\mathrm{d}t\\
       &=V_{n-1}B\left(\frac12,\frac{n+1}{2}\right)
        =V_{n-1}\frac{\sqrt\pi\,\Gamma((n+1)/2)}{\Gamma((n+2)/2)}.
\end{align*}
由归纳法及 $\Gamma(3/2)=\sqrt\pi/2$，得到
\begin{equation}
    V_n=\frac{\pi^{n/2}}{\Gamma(n/2+1)}
    =\begin{cases}
       \dfrac{\pi^m}{m!},&n=2m,\ m\geq1,\\
       \dfrac{2^{m+1}\pi^m}{(2m+1)!!},&n=2m+1,\ m\geq0.
     \end{cases}
    \label{eq:gamma-unit-ball-volume}
\end{equation}
例如 $V_2=\pi$、$V_3=4\pi/3$；半径为 $R>0$ 的球体体积再乘 $R^n$。

教材使用的高维球坐标给出同一结果：在 $n\geq2$ 时，径向变量 $r\in(0,1)$，前 $n-2$ 个角属于 $(0,\pi)$，最后一个角属于 $(0,2\pi)$；逐次作平面极坐标变换，雅可比绝对值为
\[
    r^{n-1}\sin^{n-2}\theta_1\sin^{n-3}\theta_2\cdots\sin\theta_{n-2}.
\]
坐标退化集及角度分界集为零体积集，不影响积分。于是
\[
    V_n=\frac{2\pi}{n}\prod_{j=1}^{n-2}
         \int_0^\pi\sin^j\theta\,\mathrm{d}\theta
       =\frac{2\pi}{n}\prod_{j=1}^{n-2}
         \frac{\sqrt\pi\,\Gamma((j+1)/2)}{\Gamma((j+2)/2)}
       =\frac{\pi^{n/2}}{\Gamma(n/2+1)}.
\]
$n=2$ 时乘积为空，按 $1$ 计算；一般情形的伽马函数因子逐项消去，与截面递推结果一致。
\end{solve}

\section*{习题十七}
\addcontentsline{toc}{section}{习题十七}
\fi

\end{document}
